% Phu luc 06: Trang tom tat tieng Anh. Khoang 360 tu, toi da 4 tu khoa.
% F1: du 7 y (context, problem, motivation, methodology, methods, results toi da 3, conclusions).
\IfFileExists{generated/rq-numbers.tex}{% SINH TU DONG boi writing/scripts/gen-rq-macros.py. KHONG SUA TAY.
% Bang tra cuu: docs/proposal/rq-macro-map.md. ?? = thieu nguon; n/a = mau so 0.
% Repo commit: a8b0b65a41f772c806ef4cc799c5ca7dd8831df7; harness root commit: a8b0b65a41f772c806ef4cc799c5ca7dd8831df7; split: harness 7f30afb5e9e7c3fed1f144e5af16ed1d57efb4b3, seeds [0, 1, 2], expect p_rep
% Population: P0 khong pilot
% P0, P_acc, M, R0 theo giai doan: cung nguon Bang 3.1 (tab-m-population.tex). P_rep, nguyen nhan truoc P_rep (PreRep*), phan cua R0_test (RZeroTest*): Le chap nhan khong kiem tay 2026-11-12.
% Moi so tu metrics (RQ1, RQ2, RQ3, B/N/U, robustness): kiem tren mau bang ba luot gan nhan doc lap va Le phan xu (R9: exp 31, 32; RQ2: exp 34), sai so o robust-numbers.tex (rqRNine*, rqRTwoCheck*); decision-log 2026-11-03.
% sha256 3af1f18bb0ab41b64fcf0715de4231f8bca06359955ad1160efdd39d16c03540  /tmp/claude-0/-home-claude/aa483f75-7ea5-5ae5-b102-034cb2fcbdee/scratchpad/rqin/metrics.json
% sha256 bbbdfb491cbfe83237f7237daf81e45729f8a7de5852d217b02b2263f8dc7cb0  /tmp/claude-0/-home-claude/aa483f75-7ea5-5ae5-b102-034cb2fcbdee/scratchpad/rqin/metrics-split.json
% sha256 2ad972bb0395a7ad37484fa7a6ad225908916125c4f8b9b19e46683e24b71c11  /tmp/claude-0/-home-claude/aa483f75-7ea5-5ae5-b102-034cb2fcbdee/scratchpad/rqin/analysis-exe.json
% sha256 447e678cd65f52c8a4c22e991af294519892d6cca84521504f1ba47dbc7e9b02  /tmp/claude-0/-home-claude/aa483f75-7ea5-5ae5-b102-034cb2fcbdee/scratchpad/rqin/analysis-lib.json
% sha256 f3536eabed6576568d3423a7a628b35cc7d28498f1119cbaa3ea98ace408bed3  experiments/02-harvest-repro/data/pub-decision-exe.tsv
% sha256 b92ec2b2e0b81ce26cd4dba4b185881499282ba09b4894193f1373b33f3f48f1  experiments/02-harvest-repro/results/census-exe-ec2.jsonl
% sha256 1c4635a05c54f35cd5530b5e9158f71b3d18edd5a0f6d708bc46f930c8159c28  experiments/02-harvest-repro/results/recheck-timeout-ec2.jsonl
% sha256 888aed06c374ea709d16696fc3f5840954e4555c0353a4efb8602b06dcf69afa  experiments/06-library-census/results/census-lib-c7i.jsonl
% sha256 54b87990ba21dfaa8162b7e465caa5d4c4a8e1af0a9359d27e5b5aa39017bca3  experiments/06-library-census/results/census-lib.jsonl
% sha256 6e9886225bc329fc1836d3b53d62a2a58e2ab544c128d5d17a3a3a767fe69ef4  /tmp/claude-0/-home-claude/aa483f75-7ea5-5ae5-b102-034cb2fcbdee/scratchpad/c3/ci-out/census-lib3/census-lib.jsonl
% sha256 4d3f2d366c938ac1fbfeb5b1de34a2c83be4c5ef5884ea75735d7ccaf7f86243  /tmp/claude-0/-home-claude/aa483f75-7ea5-5ae5-b102-034cb2fcbdee/scratchpad/c3/ci-out/census-lib3/recheck-loss.jsonl
% sha256 e0a559fa251f47f1096a5de31040db67c43f69c31041b6453c46305981d4f458  experiments/09-bnu-rules/scripts/bnu_v0.py

\providecommand{\rqSet}[2]{\expandafter\def\csname #1\endcsname{#2}}
\rqSet{rqPzeroExeAllNum}{1060}
\rqSet{rqPzeroExeAllDen}{1060}
\rqSet{rqPzeroExeAllKiroNum}{106}
\rqSet{rqPzeroExeAllKiroDen}{106}
\rqSet{rqPzeroExeAllClaudeNum}{106}
\rqSet{rqPzeroExeAllClaudeDen}{106}
\rqSet{rqPzeroExeAllCodexNum}{106}
\rqSet{rqPzeroExeAllCodexDen}{106}
\rqSet{rqPzeroExeAllCTwoRustNum}{106}
\rqSet{rqPzeroExeAllCTwoRustDen}{106}
\rqSet{rqPzeroExeAllLaertesNum}{106}
\rqSet{rqPzeroExeAllLaertesDen}{106}
\rqSet{rqPzeroExeAllCTwoSaferRustNum}{106}
\rqSet{rqPzeroExeAllCTwoSaferRustDen}{106}
\rqSet{rqPzeroExeAllSmartCTwoRustNum}{106}
\rqSet{rqPzeroExeAllSmartCTwoRustDen}{106}
\rqSet{rqPzeroExeAllKimiNum}{106}
\rqSet{rqPzeroExeAllKimiDen}{106}
\rqSet{rqPzeroExeAllGptFiveNum}{106}
\rqSet{rqPzeroExeAllGptFiveDen}{106}
\rqSet{rqPzeroExeAllGeminiNum}{106}
\rqSet{rqPzeroExeAllGeminiDen}{106}
\rqSet{rqPzeroExeBOneNum}{430}
\rqSet{rqPzeroExeBOneDen}{430}
\rqSet{rqPzeroExeBOneKiroNum}{43}
\rqSet{rqPzeroExeBOneKiroDen}{43}
\rqSet{rqPzeroExeBOneClaudeNum}{43}
\rqSet{rqPzeroExeBOneClaudeDen}{43}
\rqSet{rqPzeroExeBOneCodexNum}{43}
\rqSet{rqPzeroExeBOneCodexDen}{43}
\rqSet{rqPzeroExeBOneCTwoRustNum}{43}
\rqSet{rqPzeroExeBOneCTwoRustDen}{43}
\rqSet{rqPzeroExeBOneLaertesNum}{43}
\rqSet{rqPzeroExeBOneLaertesDen}{43}
\rqSet{rqPzeroExeBOneCTwoSaferRustNum}{43}
\rqSet{rqPzeroExeBOneCTwoSaferRustDen}{43}
\rqSet{rqPzeroExeBOneSmartCTwoRustNum}{43}
\rqSet{rqPzeroExeBOneSmartCTwoRustDen}{43}
\rqSet{rqPzeroExeBOneKimiNum}{43}
\rqSet{rqPzeroExeBOneKimiDen}{43}
\rqSet{rqPzeroExeBOneGptFiveNum}{43}
\rqSet{rqPzeroExeBOneGptFiveDen}{43}
\rqSet{rqPzeroExeBOneGeminiNum}{43}
\rqSet{rqPzeroExeBOneGeminiDen}{43}
\rqSet{rqPzeroExeBTwoNum}{140}
\rqSet{rqPzeroExeBTwoDen}{140}
\rqSet{rqPzeroExeBTwoKiroNum}{14}
\rqSet{rqPzeroExeBTwoKiroDen}{14}
\rqSet{rqPzeroExeBTwoClaudeNum}{14}
\rqSet{rqPzeroExeBTwoClaudeDen}{14}
\rqSet{rqPzeroExeBTwoCodexNum}{14}
\rqSet{rqPzeroExeBTwoCodexDen}{14}
\rqSet{rqPzeroExeBTwoCTwoRustNum}{14}
\rqSet{rqPzeroExeBTwoCTwoRustDen}{14}
\rqSet{rqPzeroExeBTwoLaertesNum}{14}
\rqSet{rqPzeroExeBTwoLaertesDen}{14}
\rqSet{rqPzeroExeBTwoCTwoSaferRustNum}{14}
\rqSet{rqPzeroExeBTwoCTwoSaferRustDen}{14}
\rqSet{rqPzeroExeBTwoSmartCTwoRustNum}{14}
\rqSet{rqPzeroExeBTwoSmartCTwoRustDen}{14}
\rqSet{rqPzeroExeBTwoKimiNum}{14}
\rqSet{rqPzeroExeBTwoKimiDen}{14}
\rqSet{rqPzeroExeBTwoGptFiveNum}{14}
\rqSet{rqPzeroExeBTwoGptFiveDen}{14}
\rqSet{rqPzeroExeBTwoGeminiNum}{14}
\rqSet{rqPzeroExeBTwoGeminiDen}{14}
\rqSet{rqPzeroExePZeroNum}{10}
\rqSet{rqPzeroExePZeroDen}{10}
\rqSet{rqPzeroExePZeroKiroNum}{1}
\rqSet{rqPzeroExePZeroKiroDen}{1}
\rqSet{rqPzeroExePZeroClaudeNum}{1}
\rqSet{rqPzeroExePZeroClaudeDen}{1}
\rqSet{rqPzeroExePZeroCodexNum}{1}
\rqSet{rqPzeroExePZeroCodexDen}{1}
\rqSet{rqPzeroExePZeroCTwoRustNum}{1}
\rqSet{rqPzeroExePZeroCTwoRustDen}{1}
\rqSet{rqPzeroExePZeroLaertesNum}{1}
\rqSet{rqPzeroExePZeroLaertesDen}{1}
\rqSet{rqPzeroExePZeroCTwoSaferRustNum}{1}
\rqSet{rqPzeroExePZeroCTwoSaferRustDen}{1}
\rqSet{rqPzeroExePZeroSmartCTwoRustNum}{1}
\rqSet{rqPzeroExePZeroSmartCTwoRustDen}{1}
\rqSet{rqPzeroExePZeroKimiNum}{1}
\rqSet{rqPzeroExePZeroKimiDen}{1}
\rqSet{rqPzeroExePZeroGptFiveNum}{1}
\rqSet{rqPzeroExePZeroGptFiveDen}{1}
\rqSet{rqPzeroExePZeroGeminiNum}{1}
\rqSet{rqPzeroExePZeroGeminiDen}{1}
\rqSet{rqPzeroExePOneNum}{480}
\rqSet{rqPzeroExePOneDen}{480}
\rqSet{rqPzeroExePOneKiroNum}{48}
\rqSet{rqPzeroExePOneKiroDen}{48}
\rqSet{rqPzeroExePOneClaudeNum}{48}
\rqSet{rqPzeroExePOneClaudeDen}{48}
\rqSet{rqPzeroExePOneCodexNum}{48}
\rqSet{rqPzeroExePOneCodexDen}{48}
\rqSet{rqPzeroExePOneCTwoRustNum}{48}
\rqSet{rqPzeroExePOneCTwoRustDen}{48}
\rqSet{rqPzeroExePOneLaertesNum}{48}
\rqSet{rqPzeroExePOneLaertesDen}{48}
\rqSet{rqPzeroExePOneCTwoSaferRustNum}{48}
\rqSet{rqPzeroExePOneCTwoSaferRustDen}{48}
\rqSet{rqPzeroExePOneSmartCTwoRustNum}{48}
\rqSet{rqPzeroExePOneSmartCTwoRustDen}{48}
\rqSet{rqPzeroExePOneKimiNum}{48}
\rqSet{rqPzeroExePOneKimiDen}{48}
\rqSet{rqPzeroExePOneGptFiveNum}{48}
\rqSet{rqPzeroExePOneGptFiveDen}{48}
\rqSet{rqPzeroExePOneGeminiNum}{48}
\rqSet{rqPzeroExePOneGeminiDen}{48}
\rqSet{rqPaccExeAllNum}{559}
\rqSet{rqPaccExeAllDen}{1060}
\rqSet{rqPaccExeAllPct}{52.7}
\rqSet{rqPaccExeAllKiroNum}{96}
\rqSet{rqPaccExeAllKiroDen}{106}
\rqSet{rqPaccExeAllKiroPct}{90.6}
\rqSet{rqPaccExeAllClaudeNum}{101}
\rqSet{rqPaccExeAllClaudeDen}{106}
\rqSet{rqPaccExeAllClaudePct}{95.3}
\rqSet{rqPaccExeAllCodexNum}{63}
\rqSet{rqPaccExeAllCodexDen}{106}
\rqSet{rqPaccExeAllCodexPct}{59.4}
\rqSet{rqPaccExeAllCTwoRustNum}{54}
\rqSet{rqPaccExeAllCTwoRustDen}{106}
\rqSet{rqPaccExeAllCTwoRustPct}{50.9}
\rqSet{rqPaccExeAllLaertesNum}{56}
\rqSet{rqPaccExeAllLaertesDen}{106}
\rqSet{rqPaccExeAllLaertesPct}{52.8}
\rqSet{rqPaccExeAllCTwoSaferRustNum}{42}
\rqSet{rqPaccExeAllCTwoSaferRustDen}{106}
\rqSet{rqPaccExeAllCTwoSaferRustPct}{39.6}
\rqSet{rqPaccExeAllSmartCTwoRustNum}{40}
\rqSet{rqPaccExeAllSmartCTwoRustDen}{106}
\rqSet{rqPaccExeAllSmartCTwoRustPct}{37.7}
\rqSet{rqPaccExeAllKimiNum}{28}
\rqSet{rqPaccExeAllKimiDen}{106}
\rqSet{rqPaccExeAllKimiPct}{26.4}
\rqSet{rqPaccExeAllGptFiveNum}{37}
\rqSet{rqPaccExeAllGptFiveDen}{106}
\rqSet{rqPaccExeAllGptFivePct}{34.9}
\rqSet{rqPaccExeAllGeminiNum}{42}
\rqSet{rqPaccExeAllGeminiDen}{106}
\rqSet{rqPaccExeAllGeminiPct}{39.6}
\rqSet{rqPaccExeBOneNum}{379}
\rqSet{rqPaccExeBOneDen}{430}
\rqSet{rqPaccExeBOnePct}{88.1}
\rqSet{rqPaccExeBOneKiroNum}{43}
\rqSet{rqPaccExeBOneKiroDen}{43}
\rqSet{rqPaccExeBOneKiroPct}{100.0}
\rqSet{rqPaccExeBOneClaudeNum}{43}
\rqSet{rqPaccExeBOneClaudeDen}{43}
\rqSet{rqPaccExeBOneClaudePct}{100.0}
\rqSet{rqPaccExeBOneCodexNum}{42}
\rqSet{rqPaccExeBOneCodexDen}{43}
\rqSet{rqPaccExeBOneCodexPct}{97.7}
\rqSet{rqPaccExeBOneCTwoRustNum}{43}
\rqSet{rqPaccExeBOneCTwoRustDen}{43}
\rqSet{rqPaccExeBOneCTwoRustPct}{100.0}
\rqSet{rqPaccExeBOneLaertesNum}{43}
\rqSet{rqPaccExeBOneLaertesDen}{43}
\rqSet{rqPaccExeBOneLaertesPct}{100.0}
\rqSet{rqPaccExeBOneCTwoSaferRustNum}{32}
\rqSet{rqPaccExeBOneCTwoSaferRustDen}{43}
\rqSet{rqPaccExeBOneCTwoSaferRustPct}{74.4}
\rqSet{rqPaccExeBOneSmartCTwoRustNum}{38}
\rqSet{rqPaccExeBOneSmartCTwoRustDen}{43}
\rqSet{rqPaccExeBOneSmartCTwoRustPct}{88.4}
\rqSet{rqPaccExeBOneKimiNum}{26}
\rqSet{rqPaccExeBOneKimiDen}{43}
\rqSet{rqPaccExeBOneKimiPct}{60.5}
\rqSet{rqPaccExeBOneGptFiveNum}{33}
\rqSet{rqPaccExeBOneGptFiveDen}{43}
\rqSet{rqPaccExeBOneGptFivePct}{76.7}
\rqSet{rqPaccExeBOneGeminiNum}{36}
\rqSet{rqPaccExeBOneGeminiDen}{43}
\rqSet{rqPaccExeBOneGeminiPct}{83.7}
\rqSet{rqPaccExeBTwoNum}{66}
\rqSet{rqPaccExeBTwoDen}{140}
\rqSet{rqPaccExeBTwoPct}{47.1}
\rqSet{rqPaccExeBTwoKiroNum}{4}
\rqSet{rqPaccExeBTwoKiroDen}{14}
\rqSet{rqPaccExeBTwoKiroPct}{28.6}
\rqSet{rqPaccExeBTwoClaudeNum}{9}
\rqSet{rqPaccExeBTwoClaudeDen}{14}
\rqSet{rqPaccExeBTwoClaudePct}{64.3}
\rqSet{rqPaccExeBTwoCodexNum}{8}
\rqSet{rqPaccExeBTwoCodexDen}{14}
\rqSet{rqPaccExeBTwoCodexPct}{57.1}
\rqSet{rqPaccExeBTwoCTwoRustNum}{10}
\rqSet{rqPaccExeBTwoCTwoRustDen}{14}
\rqSet{rqPaccExeBTwoCTwoRustPct}{71.4}
\rqSet{rqPaccExeBTwoLaertesNum}{12}
\rqSet{rqPaccExeBTwoLaertesDen}{14}
\rqSet{rqPaccExeBTwoLaertesPct}{85.7}
\rqSet{rqPaccExeBTwoCTwoSaferRustNum}{9}
\rqSet{rqPaccExeBTwoCTwoSaferRustDen}{14}
\rqSet{rqPaccExeBTwoCTwoSaferRustPct}{64.3}
\rqSet{rqPaccExeBTwoSmartCTwoRustNum}{2}
\rqSet{rqPaccExeBTwoSmartCTwoRustDen}{14}
\rqSet{rqPaccExeBTwoSmartCTwoRustPct}{14.3}
\rqSet{rqPaccExeBTwoKimiNum}{2}
\rqSet{rqPaccExeBTwoKimiDen}{14}
\rqSet{rqPaccExeBTwoKimiPct}{14.3}
\rqSet{rqPaccExeBTwoGptFiveNum}{4}
\rqSet{rqPaccExeBTwoGptFiveDen}{14}
\rqSet{rqPaccExeBTwoGptFivePct}{28.6}
\rqSet{rqPaccExeBTwoGeminiNum}{6}
\rqSet{rqPaccExeBTwoGeminiDen}{14}
\rqSet{rqPaccExeBTwoGeminiPct}{42.9}
\rqSet{rqPaccExePZeroNum}{6}
\rqSet{rqPaccExePZeroDen}{10}
\rqSet{rqPaccExePZeroPct}{60.0}
\rqSet{rqPaccExePZeroKiroNum}{1}
\rqSet{rqPaccExePZeroKiroDen}{1}
\rqSet{rqPaccExePZeroKiroPct}{100.0}
\rqSet{rqPaccExePZeroClaudeNum}{1}
\rqSet{rqPaccExePZeroClaudeDen}{1}
\rqSet{rqPaccExePZeroClaudePct}{100.0}
\rqSet{rqPaccExePZeroCodexNum}{1}
\rqSet{rqPaccExePZeroCodexDen}{1}
\rqSet{rqPaccExePZeroCodexPct}{100.0}
\rqSet{rqPaccExePZeroCTwoRustNum}{1}
\rqSet{rqPaccExePZeroCTwoRustDen}{1}
\rqSet{rqPaccExePZeroCTwoRustPct}{100.0}
\rqSet{rqPaccExePZeroLaertesNum}{1}
\rqSet{rqPaccExePZeroLaertesDen}{1}
\rqSet{rqPaccExePZeroLaertesPct}{100.0}
\rqSet{rqPaccExePZeroCTwoSaferRustNum}{1}
\rqSet{rqPaccExePZeroCTwoSaferRustDen}{1}
\rqSet{rqPaccExePZeroCTwoSaferRustPct}{100.0}
\rqSet{rqPaccExePZeroSmartCTwoRustNum}{0}
\rqSet{rqPaccExePZeroSmartCTwoRustDen}{1}
\rqSet{rqPaccExePZeroSmartCTwoRustPct}{0.0}
\rqSet{rqPaccExePZeroKimiNum}{0}
\rqSet{rqPaccExePZeroKimiDen}{1}
\rqSet{rqPaccExePZeroKimiPct}{0.0}
\rqSet{rqPaccExePZeroGptFiveNum}{0}
\rqSet{rqPaccExePZeroGptFiveDen}{1}
\rqSet{rqPaccExePZeroGptFivePct}{0.0}
\rqSet{rqPaccExePZeroGeminiNum}{0}
\rqSet{rqPaccExePZeroGeminiDen}{1}
\rqSet{rqPaccExePZeroGeminiPct}{0.0}
\rqSet{rqPaccExePOneNum}{108}
\rqSet{rqPaccExePOneDen}{480}
\rqSet{rqPaccExePOnePct}{22.5}
\rqSet{rqPaccExePOneKiroNum}{48}
\rqSet{rqPaccExePOneKiroDen}{48}
\rqSet{rqPaccExePOneKiroPct}{100.0}
\rqSet{rqPaccExePOneClaudeNum}{48}
\rqSet{rqPaccExePOneClaudeDen}{48}
\rqSet{rqPaccExePOneClaudePct}{100.0}
\rqSet{rqPaccExePOneCodexNum}{12}
\rqSet{rqPaccExePOneCodexDen}{48}
\rqSet{rqPaccExePOneCodexPct}{25.0}
\rqSet{rqPaccExePOneCTwoRustNum}{0}
\rqSet{rqPaccExePOneCTwoRustDen}{48}
\rqSet{rqPaccExePOneCTwoRustPct}{0.0}
\rqSet{rqPaccExePOneLaertesNum}{0}
\rqSet{rqPaccExePOneLaertesDen}{48}
\rqSet{rqPaccExePOneLaertesPct}{0.0}
\rqSet{rqPaccExePOneCTwoSaferRustNum}{0}
\rqSet{rqPaccExePOneCTwoSaferRustDen}{48}
\rqSet{rqPaccExePOneCTwoSaferRustPct}{0.0}
\rqSet{rqPaccExePOneSmartCTwoRustNum}{0}
\rqSet{rqPaccExePOneSmartCTwoRustDen}{48}
\rqSet{rqPaccExePOneSmartCTwoRustPct}{0.0}
\rqSet{rqPaccExePOneKimiNum}{0}
\rqSet{rqPaccExePOneKimiDen}{48}
\rqSet{rqPaccExePOneKimiPct}{0.0}
\rqSet{rqPaccExePOneGptFiveNum}{0}
\rqSet{rqPaccExePOneGptFiveDen}{48}
\rqSet{rqPaccExePOneGptFivePct}{0.0}
\rqSet{rqPaccExePOneGeminiNum}{0}
\rqSet{rqPaccExePOneGeminiDen}{48}
\rqSet{rqPaccExePOneGeminiPct}{0.0}
\rqSet{rqMissingExeAllNum}{47}
\rqSet{rqMissingExeAllDen}{1060}
\rqSet{rqMissingExeAllPct}{4.4}
\rqSet{rqMissingExeAllKiroNum}{0}
\rqSet{rqMissingExeAllKiroDen}{106}
\rqSet{rqMissingExeAllKiroPct}{0.0}
\rqSet{rqMissingExeAllClaudeNum}{0}
\rqSet{rqMissingExeAllClaudeDen}{106}
\rqSet{rqMissingExeAllClaudePct}{0.0}
\rqSet{rqMissingExeAllCodexNum}{0}
\rqSet{rqMissingExeAllCodexDen}{106}
\rqSet{rqMissingExeAllCodexPct}{0.0}
\rqSet{rqMissingExeAllCTwoRustNum}{0}
\rqSet{rqMissingExeAllCTwoRustDen}{106}
\rqSet{rqMissingExeAllCTwoRustPct}{0.0}
\rqSet{rqMissingExeAllLaertesNum}{0}
\rqSet{rqMissingExeAllLaertesDen}{106}
\rqSet{rqMissingExeAllLaertesPct}{0.0}
\rqSet{rqMissingExeAllCTwoSaferRustNum}{0}
\rqSet{rqMissingExeAllCTwoSaferRustDen}{106}
\rqSet{rqMissingExeAllCTwoSaferRustPct}{0.0}
\rqSet{rqMissingExeAllSmartCTwoRustNum}{0}
\rqSet{rqMissingExeAllSmartCTwoRustDen}{106}
\rqSet{rqMissingExeAllSmartCTwoRustPct}{0.0}
\rqSet{rqMissingExeAllKimiNum}{47}
\rqSet{rqMissingExeAllKimiDen}{106}
\rqSet{rqMissingExeAllKimiPct}{44.3}
\rqSet{rqMissingExeAllGptFiveNum}{0}
\rqSet{rqMissingExeAllGptFiveDen}{106}
\rqSet{rqMissingExeAllGptFivePct}{0.0}
\rqSet{rqMissingExeAllGeminiNum}{0}
\rqSet{rqMissingExeAllGeminiDen}{106}
\rqSet{rqMissingExeAllGeminiPct}{0.0}
\rqSet{rqMissingExeBOneNum}{0}
\rqSet{rqMissingExeBOneDen}{430}
\rqSet{rqMissingExeBOnePct}{0.0}
\rqSet{rqMissingExeBOneKiroNum}{0}
\rqSet{rqMissingExeBOneKiroDen}{43}
\rqSet{rqMissingExeBOneKiroPct}{0.0}
\rqSet{rqMissingExeBOneClaudeNum}{0}
\rqSet{rqMissingExeBOneClaudeDen}{43}
\rqSet{rqMissingExeBOneClaudePct}{0.0}
\rqSet{rqMissingExeBOneCodexNum}{0}
\rqSet{rqMissingExeBOneCodexDen}{43}
\rqSet{rqMissingExeBOneCodexPct}{0.0}
\rqSet{rqMissingExeBOneCTwoRustNum}{0}
\rqSet{rqMissingExeBOneCTwoRustDen}{43}
\rqSet{rqMissingExeBOneCTwoRustPct}{0.0}
\rqSet{rqMissingExeBOneLaertesNum}{0}
\rqSet{rqMissingExeBOneLaertesDen}{43}
\rqSet{rqMissingExeBOneLaertesPct}{0.0}
\rqSet{rqMissingExeBOneCTwoSaferRustNum}{0}
\rqSet{rqMissingExeBOneCTwoSaferRustDen}{43}
\rqSet{rqMissingExeBOneCTwoSaferRustPct}{0.0}
\rqSet{rqMissingExeBOneSmartCTwoRustNum}{0}
\rqSet{rqMissingExeBOneSmartCTwoRustDen}{43}
\rqSet{rqMissingExeBOneSmartCTwoRustPct}{0.0}
\rqSet{rqMissingExeBOneKimiNum}{0}
\rqSet{rqMissingExeBOneKimiDen}{43}
\rqSet{rqMissingExeBOneKimiPct}{0.0}
\rqSet{rqMissingExeBOneGptFiveNum}{0}
\rqSet{rqMissingExeBOneGptFiveDen}{43}
\rqSet{rqMissingExeBOneGptFivePct}{0.0}
\rqSet{rqMissingExeBOneGeminiNum}{0}
\rqSet{rqMissingExeBOneGeminiDen}{43}
\rqSet{rqMissingExeBOneGeminiPct}{0.0}
\rqSet{rqMissingExeBTwoNum}{0}
\rqSet{rqMissingExeBTwoDen}{140}
\rqSet{rqMissingExeBTwoPct}{0.0}
\rqSet{rqMissingExeBTwoKiroNum}{0}
\rqSet{rqMissingExeBTwoKiroDen}{14}
\rqSet{rqMissingExeBTwoKiroPct}{0.0}
\rqSet{rqMissingExeBTwoClaudeNum}{0}
\rqSet{rqMissingExeBTwoClaudeDen}{14}
\rqSet{rqMissingExeBTwoClaudePct}{0.0}
\rqSet{rqMissingExeBTwoCodexNum}{0}
\rqSet{rqMissingExeBTwoCodexDen}{14}
\rqSet{rqMissingExeBTwoCodexPct}{0.0}
\rqSet{rqMissingExeBTwoCTwoRustNum}{0}
\rqSet{rqMissingExeBTwoCTwoRustDen}{14}
\rqSet{rqMissingExeBTwoCTwoRustPct}{0.0}
\rqSet{rqMissingExeBTwoLaertesNum}{0}
\rqSet{rqMissingExeBTwoLaertesDen}{14}
\rqSet{rqMissingExeBTwoLaertesPct}{0.0}
\rqSet{rqMissingExeBTwoCTwoSaferRustNum}{0}
\rqSet{rqMissingExeBTwoCTwoSaferRustDen}{14}
\rqSet{rqMissingExeBTwoCTwoSaferRustPct}{0.0}
\rqSet{rqMissingExeBTwoSmartCTwoRustNum}{0}
\rqSet{rqMissingExeBTwoSmartCTwoRustDen}{14}
\rqSet{rqMissingExeBTwoSmartCTwoRustPct}{0.0}
\rqSet{rqMissingExeBTwoKimiNum}{0}
\rqSet{rqMissingExeBTwoKimiDen}{14}
\rqSet{rqMissingExeBTwoKimiPct}{0.0}
\rqSet{rqMissingExeBTwoGptFiveNum}{0}
\rqSet{rqMissingExeBTwoGptFiveDen}{14}
\rqSet{rqMissingExeBTwoGptFivePct}{0.0}
\rqSet{rqMissingExeBTwoGeminiNum}{0}
\rqSet{rqMissingExeBTwoGeminiDen}{14}
\rqSet{rqMissingExeBTwoGeminiPct}{0.0}
\rqSet{rqMissingExePZeroNum}{0}
\rqSet{rqMissingExePZeroDen}{10}
\rqSet{rqMissingExePZeroPct}{0.0}
\rqSet{rqMissingExePZeroKiroNum}{0}
\rqSet{rqMissingExePZeroKiroDen}{1}
\rqSet{rqMissingExePZeroKiroPct}{0.0}
\rqSet{rqMissingExePZeroClaudeNum}{0}
\rqSet{rqMissingExePZeroClaudeDen}{1}
\rqSet{rqMissingExePZeroClaudePct}{0.0}
\rqSet{rqMissingExePZeroCodexNum}{0}
\rqSet{rqMissingExePZeroCodexDen}{1}
\rqSet{rqMissingExePZeroCodexPct}{0.0}
\rqSet{rqMissingExePZeroCTwoRustNum}{0}
\rqSet{rqMissingExePZeroCTwoRustDen}{1}
\rqSet{rqMissingExePZeroCTwoRustPct}{0.0}
\rqSet{rqMissingExePZeroLaertesNum}{0}
\rqSet{rqMissingExePZeroLaertesDen}{1}
\rqSet{rqMissingExePZeroLaertesPct}{0.0}
\rqSet{rqMissingExePZeroCTwoSaferRustNum}{0}
\rqSet{rqMissingExePZeroCTwoSaferRustDen}{1}
\rqSet{rqMissingExePZeroCTwoSaferRustPct}{0.0}
\rqSet{rqMissingExePZeroSmartCTwoRustNum}{0}
\rqSet{rqMissingExePZeroSmartCTwoRustDen}{1}
\rqSet{rqMissingExePZeroSmartCTwoRustPct}{0.0}
\rqSet{rqMissingExePZeroKimiNum}{0}
\rqSet{rqMissingExePZeroKimiDen}{1}
\rqSet{rqMissingExePZeroKimiPct}{0.0}
\rqSet{rqMissingExePZeroGptFiveNum}{0}
\rqSet{rqMissingExePZeroGptFiveDen}{1}
\rqSet{rqMissingExePZeroGptFivePct}{0.0}
\rqSet{rqMissingExePZeroGeminiNum}{0}
\rqSet{rqMissingExePZeroGeminiDen}{1}
\rqSet{rqMissingExePZeroGeminiPct}{0.0}
\rqSet{rqMissingExePOneNum}{47}
\rqSet{rqMissingExePOneDen}{480}
\rqSet{rqMissingExePOnePct}{9.8}
\rqSet{rqMissingExePOneKiroNum}{0}
\rqSet{rqMissingExePOneKiroDen}{48}
\rqSet{rqMissingExePOneKiroPct}{0.0}
\rqSet{rqMissingExePOneClaudeNum}{0}
\rqSet{rqMissingExePOneClaudeDen}{48}
\rqSet{rqMissingExePOneClaudePct}{0.0}
\rqSet{rqMissingExePOneCodexNum}{0}
\rqSet{rqMissingExePOneCodexDen}{48}
\rqSet{rqMissingExePOneCodexPct}{0.0}
\rqSet{rqMissingExePOneCTwoRustNum}{0}
\rqSet{rqMissingExePOneCTwoRustDen}{48}
\rqSet{rqMissingExePOneCTwoRustPct}{0.0}
\rqSet{rqMissingExePOneLaertesNum}{0}
\rqSet{rqMissingExePOneLaertesDen}{48}
\rqSet{rqMissingExePOneLaertesPct}{0.0}
\rqSet{rqMissingExePOneCTwoSaferRustNum}{0}
\rqSet{rqMissingExePOneCTwoSaferRustDen}{48}
\rqSet{rqMissingExePOneCTwoSaferRustPct}{0.0}
\rqSet{rqMissingExePOneSmartCTwoRustNum}{0}
\rqSet{rqMissingExePOneSmartCTwoRustDen}{48}
\rqSet{rqMissingExePOneSmartCTwoRustPct}{0.0}
\rqSet{rqMissingExePOneKimiNum}{47}
\rqSet{rqMissingExePOneKimiDen}{48}
\rqSet{rqMissingExePOneKimiPct}{97.9}
\rqSet{rqMissingExePOneGptFiveNum}{0}
\rqSet{rqMissingExePOneGptFiveDen}{48}
\rqSet{rqMissingExePOneGptFivePct}{0.0}
\rqSet{rqMissingExePOneGeminiNum}{0}
\rqSet{rqMissingExePOneGeminiDen}{48}
\rqSet{rqMissingExePOneGeminiPct}{0.0}
\rqSet{rqRZeroBuildExeAllNum}{326}
\rqSet{rqRZeroBuildExeAllDen}{1060}
\rqSet{rqRZeroBuildExeAllPct}{30.8}
\rqSet{rqRZeroBuildExeAllKiroNum}{0}
\rqSet{rqRZeroBuildExeAllKiroDen}{106}
\rqSet{rqRZeroBuildExeAllKiroPct}{0.0}
\rqSet{rqRZeroBuildExeAllClaudeNum}{0}
\rqSet{rqRZeroBuildExeAllClaudeDen}{106}
\rqSet{rqRZeroBuildExeAllClaudePct}{0.0}
\rqSet{rqRZeroBuildExeAllCodexNum}{0}
\rqSet{rqRZeroBuildExeAllCodexDen}{106}
\rqSet{rqRZeroBuildExeAllCodexPct}{0.0}
\rqSet{rqRZeroBuildExeAllCTwoRustNum}{52}
\rqSet{rqRZeroBuildExeAllCTwoRustDen}{106}
\rqSet{rqRZeroBuildExeAllCTwoRustPct}{49.1}
\rqSet{rqRZeroBuildExeAllLaertesNum}{50}
\rqSet{rqRZeroBuildExeAllLaertesDen}{106}
\rqSet{rqRZeroBuildExeAllLaertesPct}{47.2}
\rqSet{rqRZeroBuildExeAllCTwoSaferRustNum}{52}
\rqSet{rqRZeroBuildExeAllCTwoSaferRustDen}{106}
\rqSet{rqRZeroBuildExeAllCTwoSaferRustPct}{49.1}
\rqSet{rqRZeroBuildExeAllSmartCTwoRustNum}{58}
\rqSet{rqRZeroBuildExeAllSmartCTwoRustDen}{106}
\rqSet{rqRZeroBuildExeAllSmartCTwoRustPct}{54.7}
\rqSet{rqRZeroBuildExeAllKimiNum}{14}
\rqSet{rqRZeroBuildExeAllKimiDen}{106}
\rqSet{rqRZeroBuildExeAllKimiPct}{13.2}
\rqSet{rqRZeroBuildExeAllGptFiveNum}{51}
\rqSet{rqRZeroBuildExeAllGptFiveDen}{106}
\rqSet{rqRZeroBuildExeAllGptFivePct}{48.1}
\rqSet{rqRZeroBuildExeAllGeminiNum}{49}
\rqSet{rqRZeroBuildExeAllGeminiDen}{106}
\rqSet{rqRZeroBuildExeAllGeminiPct}{46.2}
\rqSet{rqRZeroBuildExeBOneNum}{10}
\rqSet{rqRZeroBuildExeBOneDen}{430}
\rqSet{rqRZeroBuildExeBOnePct}{2.3}
\rqSet{rqRZeroBuildExeBOneKiroNum}{0}
\rqSet{rqRZeroBuildExeBOneKiroDen}{43}
\rqSet{rqRZeroBuildExeBOneKiroPct}{0.0}
\rqSet{rqRZeroBuildExeBOneClaudeNum}{0}
\rqSet{rqRZeroBuildExeBOneClaudeDen}{43}
\rqSet{rqRZeroBuildExeBOneClaudePct}{0.0}
\rqSet{rqRZeroBuildExeBOneCodexNum}{0}
\rqSet{rqRZeroBuildExeBOneCodexDen}{43}
\rqSet{rqRZeroBuildExeBOneCodexPct}{0.0}
\rqSet{rqRZeroBuildExeBOneCTwoRustNum}{0}
\rqSet{rqRZeroBuildExeBOneCTwoRustDen}{43}
\rqSet{rqRZeroBuildExeBOneCTwoRustPct}{0.0}
\rqSet{rqRZeroBuildExeBOneLaertesNum}{0}
\rqSet{rqRZeroBuildExeBOneLaertesDen}{43}
\rqSet{rqRZeroBuildExeBOneLaertesPct}{0.0}
\rqSet{rqRZeroBuildExeBOneCTwoSaferRustNum}{0}
\rqSet{rqRZeroBuildExeBOneCTwoSaferRustDen}{43}
\rqSet{rqRZeroBuildExeBOneCTwoSaferRustPct}{0.0}
\rqSet{rqRZeroBuildExeBOneSmartCTwoRustNum}{2}
\rqSet{rqRZeroBuildExeBOneSmartCTwoRustDen}{43}
\rqSet{rqRZeroBuildExeBOneSmartCTwoRustPct}{4.7}
\rqSet{rqRZeroBuildExeBOneKimiNum}{5}
\rqSet{rqRZeroBuildExeBOneKimiDen}{43}
\rqSet{rqRZeroBuildExeBOneKimiPct}{11.6}
\rqSet{rqRZeroBuildExeBOneGptFiveNum}{2}
\rqSet{rqRZeroBuildExeBOneGptFiveDen}{43}
\rqSet{rqRZeroBuildExeBOneGptFivePct}{4.7}
\rqSet{rqRZeroBuildExeBOneGeminiNum}{1}
\rqSet{rqRZeroBuildExeBOneGeminiDen}{43}
\rqSet{rqRZeroBuildExeBOneGeminiPct}{2.3}
\rqSet{rqRZeroBuildExeBTwoNum}{27}
\rqSet{rqRZeroBuildExeBTwoDen}{140}
\rqSet{rqRZeroBuildExeBTwoPct}{19.3}
\rqSet{rqRZeroBuildExeBTwoKiroNum}{0}
\rqSet{rqRZeroBuildExeBTwoKiroDen}{14}
\rqSet{rqRZeroBuildExeBTwoKiroPct}{0.0}
\rqSet{rqRZeroBuildExeBTwoClaudeNum}{0}
\rqSet{rqRZeroBuildExeBTwoClaudeDen}{14}
\rqSet{rqRZeroBuildExeBTwoClaudePct}{0.0}
\rqSet{rqRZeroBuildExeBTwoCodexNum}{0}
\rqSet{rqRZeroBuildExeBTwoCodexDen}{14}
\rqSet{rqRZeroBuildExeBTwoCodexPct}{0.0}
\rqSet{rqRZeroBuildExeBTwoCTwoRustNum}{4}
\rqSet{rqRZeroBuildExeBTwoCTwoRustDen}{14}
\rqSet{rqRZeroBuildExeBTwoCTwoRustPct}{28.6}
\rqSet{rqRZeroBuildExeBTwoLaertesNum}{2}
\rqSet{rqRZeroBuildExeBTwoLaertesDen}{14}
\rqSet{rqRZeroBuildExeBTwoLaertesPct}{14.3}
\rqSet{rqRZeroBuildExeBTwoCTwoSaferRustNum}{4}
\rqSet{rqRZeroBuildExeBTwoCTwoSaferRustDen}{14}
\rqSet{rqRZeroBuildExeBTwoCTwoSaferRustPct}{28.6}
\rqSet{rqRZeroBuildExeBTwoSmartCTwoRustNum}{8}
\rqSet{rqRZeroBuildExeBTwoSmartCTwoRustDen}{14}
\rqSet{rqRZeroBuildExeBTwoSmartCTwoRustPct}{57.1}
\rqSet{rqRZeroBuildExeBTwoKimiNum}{8}
\rqSet{rqRZeroBuildExeBTwoKimiDen}{14}
\rqSet{rqRZeroBuildExeBTwoKimiPct}{57.1}
\rqSet{rqRZeroBuildExeBTwoGptFiveNum}{1}
\rqSet{rqRZeroBuildExeBTwoGptFiveDen}{14}
\rqSet{rqRZeroBuildExeBTwoGptFivePct}{7.1}
\rqSet{rqRZeroBuildExeBTwoGeminiNum}{0}
\rqSet{rqRZeroBuildExeBTwoGeminiDen}{14}
\rqSet{rqRZeroBuildExeBTwoGeminiPct}{0.0}
\rqSet{rqRZeroBuildExePZeroNum}{0}
\rqSet{rqRZeroBuildExePZeroDen}{10}
\rqSet{rqRZeroBuildExePZeroPct}{0.0}
\rqSet{rqRZeroBuildExePZeroKiroNum}{0}
\rqSet{rqRZeroBuildExePZeroKiroDen}{1}
\rqSet{rqRZeroBuildExePZeroKiroPct}{0.0}
\rqSet{rqRZeroBuildExePZeroClaudeNum}{0}
\rqSet{rqRZeroBuildExePZeroClaudeDen}{1}
\rqSet{rqRZeroBuildExePZeroClaudePct}{0.0}
\rqSet{rqRZeroBuildExePZeroCodexNum}{0}
\rqSet{rqRZeroBuildExePZeroCodexDen}{1}
\rqSet{rqRZeroBuildExePZeroCodexPct}{0.0}
\rqSet{rqRZeroBuildExePZeroCTwoRustNum}{0}
\rqSet{rqRZeroBuildExePZeroCTwoRustDen}{1}
\rqSet{rqRZeroBuildExePZeroCTwoRustPct}{0.0}
\rqSet{rqRZeroBuildExePZeroLaertesNum}{0}
\rqSet{rqRZeroBuildExePZeroLaertesDen}{1}
\rqSet{rqRZeroBuildExePZeroLaertesPct}{0.0}
\rqSet{rqRZeroBuildExePZeroCTwoSaferRustNum}{0}
\rqSet{rqRZeroBuildExePZeroCTwoSaferRustDen}{1}
\rqSet{rqRZeroBuildExePZeroCTwoSaferRustPct}{0.0}
\rqSet{rqRZeroBuildExePZeroSmartCTwoRustNum}{0}
\rqSet{rqRZeroBuildExePZeroSmartCTwoRustDen}{1}
\rqSet{rqRZeroBuildExePZeroSmartCTwoRustPct}{0.0}
\rqSet{rqRZeroBuildExePZeroKimiNum}{0}
\rqSet{rqRZeroBuildExePZeroKimiDen}{1}
\rqSet{rqRZeroBuildExePZeroKimiPct}{0.0}
\rqSet{rqRZeroBuildExePZeroGptFiveNum}{0}
\rqSet{rqRZeroBuildExePZeroGptFiveDen}{1}
\rqSet{rqRZeroBuildExePZeroGptFivePct}{0.0}
\rqSet{rqRZeroBuildExePZeroGeminiNum}{0}
\rqSet{rqRZeroBuildExePZeroGeminiDen}{1}
\rqSet{rqRZeroBuildExePZeroGeminiPct}{0.0}
\rqSet{rqRZeroBuildExePOneNum}{289}
\rqSet{rqRZeroBuildExePOneDen}{480}
\rqSet{rqRZeroBuildExePOnePct}{60.2}
\rqSet{rqRZeroBuildExePOneKiroNum}{0}
\rqSet{rqRZeroBuildExePOneKiroDen}{48}
\rqSet{rqRZeroBuildExePOneKiroPct}{0.0}
\rqSet{rqRZeroBuildExePOneClaudeNum}{0}
\rqSet{rqRZeroBuildExePOneClaudeDen}{48}
\rqSet{rqRZeroBuildExePOneClaudePct}{0.0}
\rqSet{rqRZeroBuildExePOneCodexNum}{0}
\rqSet{rqRZeroBuildExePOneCodexDen}{48}
\rqSet{rqRZeroBuildExePOneCodexPct}{0.0}
\rqSet{rqRZeroBuildExePOneCTwoRustNum}{48}
\rqSet{rqRZeroBuildExePOneCTwoRustDen}{48}
\rqSet{rqRZeroBuildExePOneCTwoRustPct}{100.0}
\rqSet{rqRZeroBuildExePOneLaertesNum}{48}
\rqSet{rqRZeroBuildExePOneLaertesDen}{48}
\rqSet{rqRZeroBuildExePOneLaertesPct}{100.0}
\rqSet{rqRZeroBuildExePOneCTwoSaferRustNum}{48}
\rqSet{rqRZeroBuildExePOneCTwoSaferRustDen}{48}
\rqSet{rqRZeroBuildExePOneCTwoSaferRustPct}{100.0}
\rqSet{rqRZeroBuildExePOneSmartCTwoRustNum}{48}
\rqSet{rqRZeroBuildExePOneSmartCTwoRustDen}{48}
\rqSet{rqRZeroBuildExePOneSmartCTwoRustPct}{100.0}
\rqSet{rqRZeroBuildExePOneKimiNum}{1}
\rqSet{rqRZeroBuildExePOneKimiDen}{48}
\rqSet{rqRZeroBuildExePOneKimiPct}{2.1}
\rqSet{rqRZeroBuildExePOneGptFiveNum}{48}
\rqSet{rqRZeroBuildExePOneGptFiveDen}{48}
\rqSet{rqRZeroBuildExePOneGptFivePct}{100.0}
\rqSet{rqRZeroBuildExePOneGeminiNum}{48}
\rqSet{rqRZeroBuildExePOneGeminiDen}{48}
\rqSet{rqRZeroBuildExePOneGeminiPct}{100.0}
\rqSet{rqRZeroTestExeAllNum}{128}
\rqSet{rqRZeroTestExeAllDen}{1060}
\rqSet{rqRZeroTestExeAllPct}{12.1}
\rqSet{rqRZeroTestExeAllKiroNum}{10}
\rqSet{rqRZeroTestExeAllKiroDen}{106}
\rqSet{rqRZeroTestExeAllKiroPct}{9.4}
\rqSet{rqRZeroTestExeAllClaudeNum}{5}
\rqSet{rqRZeroTestExeAllClaudeDen}{106}
\rqSet{rqRZeroTestExeAllClaudePct}{4.7}
\rqSet{rqRZeroTestExeAllCodexNum}{43}
\rqSet{rqRZeroTestExeAllCodexDen}{106}
\rqSet{rqRZeroTestExeAllCodexPct}{40.6}
\rqSet{rqRZeroTestExeAllCTwoRustNum}{0}
\rqSet{rqRZeroTestExeAllCTwoRustDen}{106}
\rqSet{rqRZeroTestExeAllCTwoRustPct}{0.0}
\rqSet{rqRZeroTestExeAllLaertesNum}{0}
\rqSet{rqRZeroTestExeAllLaertesDen}{106}
\rqSet{rqRZeroTestExeAllLaertesPct}{0.0}
\rqSet{rqRZeroTestExeAllCTwoSaferRustNum}{12}
\rqSet{rqRZeroTestExeAllCTwoSaferRustDen}{106}
\rqSet{rqRZeroTestExeAllCTwoSaferRustPct}{11.3}
\rqSet{rqRZeroTestExeAllSmartCTwoRustNum}{8}
\rqSet{rqRZeroTestExeAllSmartCTwoRustDen}{106}
\rqSet{rqRZeroTestExeAllSmartCTwoRustPct}{7.5}
\rqSet{rqRZeroTestExeAllKimiNum}{17}
\rqSet{rqRZeroTestExeAllKimiDen}{106}
\rqSet{rqRZeroTestExeAllKimiPct}{16.0}
\rqSet{rqRZeroTestExeAllGptFiveNum}{18}
\rqSet{rqRZeroTestExeAllGptFiveDen}{106}
\rqSet{rqRZeroTestExeAllGptFivePct}{17.0}
\rqSet{rqRZeroTestExeAllGeminiNum}{15}
\rqSet{rqRZeroTestExeAllGeminiDen}{106}
\rqSet{rqRZeroTestExeAllGeminiPct}{14.2}
\rqSet{rqRZeroTestExeBOneNum}{41}
\rqSet{rqRZeroTestExeBOneDen}{430}
\rqSet{rqRZeroTestExeBOnePct}{9.5}
\rqSet{rqRZeroTestExeBOneKiroNum}{0}
\rqSet{rqRZeroTestExeBOneKiroDen}{43}
\rqSet{rqRZeroTestExeBOneKiroPct}{0.0}
\rqSet{rqRZeroTestExeBOneClaudeNum}{0}
\rqSet{rqRZeroTestExeBOneClaudeDen}{43}
\rqSet{rqRZeroTestExeBOneClaudePct}{0.0}
\rqSet{rqRZeroTestExeBOneCodexNum}{1}
\rqSet{rqRZeroTestExeBOneCodexDen}{43}
\rqSet{rqRZeroTestExeBOneCodexPct}{2.3}
\rqSet{rqRZeroTestExeBOneCTwoRustNum}{0}
\rqSet{rqRZeroTestExeBOneCTwoRustDen}{43}
\rqSet{rqRZeroTestExeBOneCTwoRustPct}{0.0}
\rqSet{rqRZeroTestExeBOneLaertesNum}{0}
\rqSet{rqRZeroTestExeBOneLaertesDen}{43}
\rqSet{rqRZeroTestExeBOneLaertesPct}{0.0}
\rqSet{rqRZeroTestExeBOneCTwoSaferRustNum}{11}
\rqSet{rqRZeroTestExeBOneCTwoSaferRustDen}{43}
\rqSet{rqRZeroTestExeBOneCTwoSaferRustPct}{25.6}
\rqSet{rqRZeroTestExeBOneSmartCTwoRustNum}{3}
\rqSet{rqRZeroTestExeBOneSmartCTwoRustDen}{43}
\rqSet{rqRZeroTestExeBOneSmartCTwoRustPct}{7.0}
\rqSet{rqRZeroTestExeBOneKimiNum}{12}
\rqSet{rqRZeroTestExeBOneKimiDen}{43}
\rqSet{rqRZeroTestExeBOneKimiPct}{27.9}
\rqSet{rqRZeroTestExeBOneGptFiveNum}{8}
\rqSet{rqRZeroTestExeBOneGptFiveDen}{43}
\rqSet{rqRZeroTestExeBOneGptFivePct}{18.6}
\rqSet{rqRZeroTestExeBOneGeminiNum}{6}
\rqSet{rqRZeroTestExeBOneGeminiDen}{43}
\rqSet{rqRZeroTestExeBOneGeminiPct}{14.0}
\rqSet{rqRZeroTestExeBTwoNum}{47}
\rqSet{rqRZeroTestExeBTwoDen}{140}
\rqSet{rqRZeroTestExeBTwoPct}{33.6}
\rqSet{rqRZeroTestExeBTwoKiroNum}{10}
\rqSet{rqRZeroTestExeBTwoKiroDen}{14}
\rqSet{rqRZeroTestExeBTwoKiroPct}{71.4}
\rqSet{rqRZeroTestExeBTwoClaudeNum}{5}
\rqSet{rqRZeroTestExeBTwoClaudeDen}{14}
\rqSet{rqRZeroTestExeBTwoClaudePct}{35.7}
\rqSet{rqRZeroTestExeBTwoCodexNum}{6}
\rqSet{rqRZeroTestExeBTwoCodexDen}{14}
\rqSet{rqRZeroTestExeBTwoCodexPct}{42.9}
\rqSet{rqRZeroTestExeBTwoCTwoRustNum}{0}
\rqSet{rqRZeroTestExeBTwoCTwoRustDen}{14}
\rqSet{rqRZeroTestExeBTwoCTwoRustPct}{0.0}
\rqSet{rqRZeroTestExeBTwoLaertesNum}{0}
\rqSet{rqRZeroTestExeBTwoLaertesDen}{14}
\rqSet{rqRZeroTestExeBTwoLaertesPct}{0.0}
\rqSet{rqRZeroTestExeBTwoCTwoSaferRustNum}{1}
\rqSet{rqRZeroTestExeBTwoCTwoSaferRustDen}{14}
\rqSet{rqRZeroTestExeBTwoCTwoSaferRustPct}{7.1}
\rqSet{rqRZeroTestExeBTwoSmartCTwoRustNum}{4}
\rqSet{rqRZeroTestExeBTwoSmartCTwoRustDen}{14}
\rqSet{rqRZeroTestExeBTwoSmartCTwoRustPct}{28.6}
\rqSet{rqRZeroTestExeBTwoKimiNum}{4}
\rqSet{rqRZeroTestExeBTwoKimiDen}{14}
\rqSet{rqRZeroTestExeBTwoKimiPct}{28.6}
\rqSet{rqRZeroTestExeBTwoGptFiveNum}{9}
\rqSet{rqRZeroTestExeBTwoGptFiveDen}{14}
\rqSet{rqRZeroTestExeBTwoGptFivePct}{64.3}
\rqSet{rqRZeroTestExeBTwoGeminiNum}{8}
\rqSet{rqRZeroTestExeBTwoGeminiDen}{14}
\rqSet{rqRZeroTestExeBTwoGeminiPct}{57.1}
\rqSet{rqRZeroTestExePZeroNum}{4}
\rqSet{rqRZeroTestExePZeroDen}{10}
\rqSet{rqRZeroTestExePZeroPct}{40.0}
\rqSet{rqRZeroTestExePZeroKiroNum}{0}
\rqSet{rqRZeroTestExePZeroKiroDen}{1}
\rqSet{rqRZeroTestExePZeroKiroPct}{0.0}
\rqSet{rqRZeroTestExePZeroClaudeNum}{0}
\rqSet{rqRZeroTestExePZeroClaudeDen}{1}
\rqSet{rqRZeroTestExePZeroClaudePct}{0.0}
\rqSet{rqRZeroTestExePZeroCodexNum}{0}
\rqSet{rqRZeroTestExePZeroCodexDen}{1}
\rqSet{rqRZeroTestExePZeroCodexPct}{0.0}
\rqSet{rqRZeroTestExePZeroCTwoRustNum}{0}
\rqSet{rqRZeroTestExePZeroCTwoRustDen}{1}
\rqSet{rqRZeroTestExePZeroCTwoRustPct}{0.0}
\rqSet{rqRZeroTestExePZeroLaertesNum}{0}
\rqSet{rqRZeroTestExePZeroLaertesDen}{1}
\rqSet{rqRZeroTestExePZeroLaertesPct}{0.0}
\rqSet{rqRZeroTestExePZeroCTwoSaferRustNum}{0}
\rqSet{rqRZeroTestExePZeroCTwoSaferRustDen}{1}
\rqSet{rqRZeroTestExePZeroCTwoSaferRustPct}{0.0}
\rqSet{rqRZeroTestExePZeroSmartCTwoRustNum}{1}
\rqSet{rqRZeroTestExePZeroSmartCTwoRustDen}{1}
\rqSet{rqRZeroTestExePZeroSmartCTwoRustPct}{100.0}
\rqSet{rqRZeroTestExePZeroKimiNum}{1}
\rqSet{rqRZeroTestExePZeroKimiDen}{1}
\rqSet{rqRZeroTestExePZeroKimiPct}{100.0}
\rqSet{rqRZeroTestExePZeroGptFiveNum}{1}
\rqSet{rqRZeroTestExePZeroGptFiveDen}{1}
\rqSet{rqRZeroTestExePZeroGptFivePct}{100.0}
\rqSet{rqRZeroTestExePZeroGeminiNum}{1}
\rqSet{rqRZeroTestExePZeroGeminiDen}{1}
\rqSet{rqRZeroTestExePZeroGeminiPct}{100.0}
\rqSet{rqRZeroTestExePOneNum}{36}
\rqSet{rqRZeroTestExePOneDen}{480}
\rqSet{rqRZeroTestExePOnePct}{7.5}
\rqSet{rqRZeroTestExePOneKiroNum}{0}
\rqSet{rqRZeroTestExePOneKiroDen}{48}
\rqSet{rqRZeroTestExePOneKiroPct}{0.0}
\rqSet{rqRZeroTestExePOneClaudeNum}{0}
\rqSet{rqRZeroTestExePOneClaudeDen}{48}
\rqSet{rqRZeroTestExePOneClaudePct}{0.0}
\rqSet{rqRZeroTestExePOneCodexNum}{36}
\rqSet{rqRZeroTestExePOneCodexDen}{48}
\rqSet{rqRZeroTestExePOneCodexPct}{75.0}
\rqSet{rqRZeroTestExePOneCTwoRustNum}{0}
\rqSet{rqRZeroTestExePOneCTwoRustDen}{48}
\rqSet{rqRZeroTestExePOneCTwoRustPct}{0.0}
\rqSet{rqRZeroTestExePOneLaertesNum}{0}
\rqSet{rqRZeroTestExePOneLaertesDen}{48}
\rqSet{rqRZeroTestExePOneLaertesPct}{0.0}
\rqSet{rqRZeroTestExePOneCTwoSaferRustNum}{0}
\rqSet{rqRZeroTestExePOneCTwoSaferRustDen}{48}
\rqSet{rqRZeroTestExePOneCTwoSaferRustPct}{0.0}
\rqSet{rqRZeroTestExePOneSmartCTwoRustNum}{0}
\rqSet{rqRZeroTestExePOneSmartCTwoRustDen}{48}
\rqSet{rqRZeroTestExePOneSmartCTwoRustPct}{0.0}
\rqSet{rqRZeroTestExePOneKimiNum}{0}
\rqSet{rqRZeroTestExePOneKimiDen}{48}
\rqSet{rqRZeroTestExePOneKimiPct}{0.0}
\rqSet{rqRZeroTestExePOneGptFiveNum}{0}
\rqSet{rqRZeroTestExePOneGptFiveDen}{48}
\rqSet{rqRZeroTestExePOneGptFivePct}{0.0}
\rqSet{rqRZeroTestExePOneGeminiNum}{0}
\rqSet{rqRZeroTestExePOneGeminiDen}{48}
\rqSet{rqRZeroTestExePOneGeminiPct}{0.0}
\rqSet{rqPrepExeAllNum}{548}
\rqSet{rqPrepExeAllDen}{559}
\rqSet{rqPrepExeAllPct}{98.0}
\rqSet{rqPrepExeAllKiroNum}{96}
\rqSet{rqPrepExeAllKiroDen}{96}
\rqSet{rqPrepExeAllKiroPct}{100.0}
\rqSet{rqPrepExeAllClaudeNum}{101}
\rqSet{rqPrepExeAllClaudeDen}{101}
\rqSet{rqPrepExeAllClaudePct}{100.0}
\rqSet{rqPrepExeAllCodexNum}{62}
\rqSet{rqPrepExeAllCodexDen}{63}
\rqSet{rqPrepExeAllCodexPct}{98.4}
\rqSet{rqPrepExeAllCTwoRustNum}{53}
\rqSet{rqPrepExeAllCTwoRustDen}{54}
\rqSet{rqPrepExeAllCTwoRustPct}{98.1}
\rqSet{rqPrepExeAllLaertesNum}{53}
\rqSet{rqPrepExeAllLaertesDen}{56}
\rqSet{rqPrepExeAllLaertesPct}{94.6}
\rqSet{rqPrepExeAllCTwoSaferRustNum}{39}
\rqSet{rqPrepExeAllCTwoSaferRustDen}{42}
\rqSet{rqPrepExeAllCTwoSaferRustPct}{92.9}
\rqSet{rqPrepExeAllSmartCTwoRustNum}{40}
\rqSet{rqPrepExeAllSmartCTwoRustDen}{40}
\rqSet{rqPrepExeAllSmartCTwoRustPct}{100.0}
\rqSet{rqPrepExeAllKimiNum}{27}
\rqSet{rqPrepExeAllKimiDen}{28}
\rqSet{rqPrepExeAllKimiPct}{96.4}
\rqSet{rqPrepExeAllGptFiveNum}{36}
\rqSet{rqPrepExeAllGptFiveDen}{37}
\rqSet{rqPrepExeAllGptFivePct}{97.3}
\rqSet{rqPrepExeAllGeminiNum}{41}
\rqSet{rqPrepExeAllGeminiDen}{42}
\rqSet{rqPrepExeAllGeminiPct}{97.6}
\rqSet{rqPrepExeBOneNum}{378}
\rqSet{rqPrepExeBOneDen}{379}
\rqSet{rqPrepExeBOnePct}{99.7}
\rqSet{rqPrepExeBOneKiroNum}{43}
\rqSet{rqPrepExeBOneKiroDen}{43}
\rqSet{rqPrepExeBOneKiroPct}{100.0}
\rqSet{rqPrepExeBOneClaudeNum}{43}
\rqSet{rqPrepExeBOneClaudeDen}{43}
\rqSet{rqPrepExeBOneClaudePct}{100.0}
\rqSet{rqPrepExeBOneCodexNum}{41}
\rqSet{rqPrepExeBOneCodexDen}{42}
\rqSet{rqPrepExeBOneCodexPct}{97.6}
\rqSet{rqPrepExeBOneCTwoRustNum}{43}
\rqSet{rqPrepExeBOneCTwoRustDen}{43}
\rqSet{rqPrepExeBOneCTwoRustPct}{100.0}
\rqSet{rqPrepExeBOneLaertesNum}{43}
\rqSet{rqPrepExeBOneLaertesDen}{43}
\rqSet{rqPrepExeBOneLaertesPct}{100.0}
\rqSet{rqPrepExeBOneCTwoSaferRustNum}{32}
\rqSet{rqPrepExeBOneCTwoSaferRustDen}{32}
\rqSet{rqPrepExeBOneCTwoSaferRustPct}{100.0}
\rqSet{rqPrepExeBOneSmartCTwoRustNum}{38}
\rqSet{rqPrepExeBOneSmartCTwoRustDen}{38}
\rqSet{rqPrepExeBOneSmartCTwoRustPct}{100.0}
\rqSet{rqPrepExeBOneKimiNum}{26}
\rqSet{rqPrepExeBOneKimiDen}{26}
\rqSet{rqPrepExeBOneKimiPct}{100.0}
\rqSet{rqPrepExeBOneGptFiveNum}{33}
\rqSet{rqPrepExeBOneGptFiveDen}{33}
\rqSet{rqPrepExeBOneGptFivePct}{100.0}
\rqSet{rqPrepExeBOneGeminiNum}{36}
\rqSet{rqPrepExeBOneGeminiDen}{36}
\rqSet{rqPrepExeBOneGeminiPct}{100.0}
\rqSet{rqPrepExeBTwoNum}{56}
\rqSet{rqPrepExeBTwoDen}{66}
\rqSet{rqPrepExeBTwoPct}{84.8}
\rqSet{rqPrepExeBTwoKiroNum}{4}
\rqSet{rqPrepExeBTwoKiroDen}{4}
\rqSet{rqPrepExeBTwoKiroPct}{100.0}
\rqSet{rqPrepExeBTwoClaudeNum}{9}
\rqSet{rqPrepExeBTwoClaudeDen}{9}
\rqSet{rqPrepExeBTwoClaudePct}{100.0}
\rqSet{rqPrepExeBTwoCodexNum}{8}
\rqSet{rqPrepExeBTwoCodexDen}{8}
\rqSet{rqPrepExeBTwoCodexPct}{100.0}
\rqSet{rqPrepExeBTwoCTwoRustNum}{9}
\rqSet{rqPrepExeBTwoCTwoRustDen}{10}
\rqSet{rqPrepExeBTwoCTwoRustPct}{90.0}
\rqSet{rqPrepExeBTwoLaertesNum}{9}
\rqSet{rqPrepExeBTwoLaertesDen}{12}
\rqSet{rqPrepExeBTwoLaertesPct}{75.0}
\rqSet{rqPrepExeBTwoCTwoSaferRustNum}{6}
\rqSet{rqPrepExeBTwoCTwoSaferRustDen}{9}
\rqSet{rqPrepExeBTwoCTwoSaferRustPct}{66.7}
\rqSet{rqPrepExeBTwoSmartCTwoRustNum}{2}
\rqSet{rqPrepExeBTwoSmartCTwoRustDen}{2}
\rqSet{rqPrepExeBTwoSmartCTwoRustPct}{100.0}
\rqSet{rqPrepExeBTwoKimiNum}{1}
\rqSet{rqPrepExeBTwoKimiDen}{2}
\rqSet{rqPrepExeBTwoKimiPct}{50.0}
\rqSet{rqPrepExeBTwoGptFiveNum}{3}
\rqSet{rqPrepExeBTwoGptFiveDen}{4}
\rqSet{rqPrepExeBTwoGptFivePct}{75.0}
\rqSet{rqPrepExeBTwoGeminiNum}{5}
\rqSet{rqPrepExeBTwoGeminiDen}{6}
\rqSet{rqPrepExeBTwoGeminiPct}{83.3}
\rqSet{rqPrepExePZeroNum}{6}
\rqSet{rqPrepExePZeroDen}{6}
\rqSet{rqPrepExePZeroPct}{100.0}
\rqSet{rqPrepExePZeroKiroNum}{1}
\rqSet{rqPrepExePZeroKiroDen}{1}
\rqSet{rqPrepExePZeroKiroPct}{100.0}
\rqSet{rqPrepExePZeroClaudeNum}{1}
\rqSet{rqPrepExePZeroClaudeDen}{1}
\rqSet{rqPrepExePZeroClaudePct}{100.0}
\rqSet{rqPrepExePZeroCodexNum}{1}
\rqSet{rqPrepExePZeroCodexDen}{1}
\rqSet{rqPrepExePZeroCodexPct}{100.0}
\rqSet{rqPrepExePZeroCTwoRustNum}{1}
\rqSet{rqPrepExePZeroCTwoRustDen}{1}
\rqSet{rqPrepExePZeroCTwoRustPct}{100.0}
\rqSet{rqPrepExePZeroLaertesNum}{1}
\rqSet{rqPrepExePZeroLaertesDen}{1}
\rqSet{rqPrepExePZeroLaertesPct}{100.0}
\rqSet{rqPrepExePZeroCTwoSaferRustNum}{1}
\rqSet{rqPrepExePZeroCTwoSaferRustDen}{1}
\rqSet{rqPrepExePZeroCTwoSaferRustPct}{100.0}
\rqSet{rqPrepExePZeroSmartCTwoRustNum}{0}
\rqSet{rqPrepExePZeroSmartCTwoRustDen}{0}
\rqSet{rqPrepExePZeroSmartCTwoRustPct}{n/a}
\rqSet{rqPrepExePZeroKimiNum}{0}
\rqSet{rqPrepExePZeroKimiDen}{0}
\rqSet{rqPrepExePZeroKimiPct}{n/a}
\rqSet{rqPrepExePZeroGptFiveNum}{0}
\rqSet{rqPrepExePZeroGptFiveDen}{0}
\rqSet{rqPrepExePZeroGptFivePct}{n/a}
\rqSet{rqPrepExePZeroGeminiNum}{0}
\rqSet{rqPrepExePZeroGeminiDen}{0}
\rqSet{rqPrepExePZeroGeminiPct}{n/a}
\rqSet{rqPrepExePOneNum}{108}
\rqSet{rqPrepExePOneDen}{108}
\rqSet{rqPrepExePOnePct}{100.0}
\rqSet{rqPrepExePOneKiroNum}{48}
\rqSet{rqPrepExePOneKiroDen}{48}
\rqSet{rqPrepExePOneKiroPct}{100.0}
\rqSet{rqPrepExePOneClaudeNum}{48}
\rqSet{rqPrepExePOneClaudeDen}{48}
\rqSet{rqPrepExePOneClaudePct}{100.0}
\rqSet{rqPrepExePOneCodexNum}{12}
\rqSet{rqPrepExePOneCodexDen}{12}
\rqSet{rqPrepExePOneCodexPct}{100.0}
\rqSet{rqPrepExePOneCTwoRustNum}{0}
\rqSet{rqPrepExePOneCTwoRustDen}{0}
\rqSet{rqPrepExePOneCTwoRustPct}{n/a}
\rqSet{rqPrepExePOneLaertesNum}{0}
\rqSet{rqPrepExePOneLaertesDen}{0}
\rqSet{rqPrepExePOneLaertesPct}{n/a}
\rqSet{rqPrepExePOneCTwoSaferRustNum}{0}
\rqSet{rqPrepExePOneCTwoSaferRustDen}{0}
\rqSet{rqPrepExePOneCTwoSaferRustPct}{n/a}
\rqSet{rqPrepExePOneSmartCTwoRustNum}{0}
\rqSet{rqPrepExePOneSmartCTwoRustDen}{0}
\rqSet{rqPrepExePOneSmartCTwoRustPct}{n/a}
\rqSet{rqPrepExePOneKimiNum}{0}
\rqSet{rqPrepExePOneKimiDen}{0}
\rqSet{rqPrepExePOneKimiPct}{n/a}
\rqSet{rqPrepExePOneGptFiveNum}{0}
\rqSet{rqPrepExePOneGptFiveDen}{0}
\rqSet{rqPrepExePOneGptFivePct}{n/a}
\rqSet{rqPrepExePOneGeminiNum}{0}
\rqSet{rqPrepExePOneGeminiDen}{0}
\rqSet{rqPrepExePOneGeminiPct}{n/a}
\rqSet{rqPtestExeAllNum}{476}
\rqSet{rqPtestExeAllDen}{559}
\rqSet{rqPtestExeAllPct}{85.2}
\rqSet{rqPtestExeAllKiroNum}{88}
\rqSet{rqPtestExeAllKiroDen}{96}
\rqSet{rqPtestExeAllKiroPct}{91.7}
\rqSet{rqPtestExeAllClaudeNum}{93}
\rqSet{rqPtestExeAllClaudeDen}{101}
\rqSet{rqPtestExeAllClaudePct}{92.1}
\rqSet{rqPtestExeAllCodexNum}{54}
\rqSet{rqPtestExeAllCodexDen}{63}
\rqSet{rqPtestExeAllCodexPct}{85.7}
\rqSet{rqPtestExeAllCTwoRustNum}{44}
\rqSet{rqPtestExeAllCTwoRustDen}{54}
\rqSet{rqPtestExeAllCTwoRustPct}{81.5}
\rqSet{rqPtestExeAllLaertesNum}{44}
\rqSet{rqPtestExeAllLaertesDen}{56}
\rqSet{rqPtestExeAllLaertesPct}{78.6}
\rqSet{rqPtestExeAllCTwoSaferRustNum}{32}
\rqSet{rqPtestExeAllCTwoSaferRustDen}{42}
\rqSet{rqPtestExeAllCTwoSaferRustPct}{76.2}
\rqSet{rqPtestExeAllSmartCTwoRustNum}{33}
\rqSet{rqPtestExeAllSmartCTwoRustDen}{40}
\rqSet{rqPtestExeAllSmartCTwoRustPct}{82.5}
\rqSet{rqPtestExeAllKimiNum}{24}
\rqSet{rqPtestExeAllKimiDen}{28}
\rqSet{rqPtestExeAllKimiPct}{85.7}
\rqSet{rqPtestExeAllGptFiveNum}{30}
\rqSet{rqPtestExeAllGptFiveDen}{37}
\rqSet{rqPtestExeAllGptFivePct}{81.1}
\rqSet{rqPtestExeAllGeminiNum}{34}
\rqSet{rqPtestExeAllGeminiDen}{42}
\rqSet{rqPtestExeAllGeminiPct}{81.0}
\rqSet{rqPtestExeBOneNum}{315}
\rqSet{rqPtestExeBOneDen}{379}
\rqSet{rqPtestExeBOnePct}{83.1}
\rqSet{rqPtestExeBOneKiroNum}{36}
\rqSet{rqPtestExeBOneKiroDen}{43}
\rqSet{rqPtestExeBOneKiroPct}{83.7}
\rqSet{rqPtestExeBOneClaudeNum}{36}
\rqSet{rqPtestExeBOneClaudeDen}{43}
\rqSet{rqPtestExeBOneClaudePct}{83.7}
\rqSet{rqPtestExeBOneCodexNum}{34}
\rqSet{rqPtestExeBOneCodexDen}{42}
\rqSet{rqPtestExeBOneCodexPct}{81.0}
\rqSet{rqPtestExeBOneCTwoRustNum}{36}
\rqSet{rqPtestExeBOneCTwoRustDen}{43}
\rqSet{rqPtestExeBOneCTwoRustPct}{83.7}
\rqSet{rqPtestExeBOneLaertesNum}{36}
\rqSet{rqPtestExeBOneLaertesDen}{43}
\rqSet{rqPtestExeBOneLaertesPct}{83.7}
\rqSet{rqPtestExeBOneCTwoSaferRustNum}{27}
\rqSet{rqPtestExeBOneCTwoSaferRustDen}{32}
\rqSet{rqPtestExeBOneCTwoSaferRustPct}{84.4}
\rqSet{rqPtestExeBOneSmartCTwoRustNum}{31}
\rqSet{rqPtestExeBOneSmartCTwoRustDen}{38}
\rqSet{rqPtestExeBOneSmartCTwoRustPct}{81.6}
\rqSet{rqPtestExeBOneKimiNum}{23}
\rqSet{rqPtestExeBOneKimiDen}{26}
\rqSet{rqPtestExeBOneKimiPct}{88.5}
\rqSet{rqPtestExeBOneGptFiveNum}{27}
\rqSet{rqPtestExeBOneGptFiveDen}{33}
\rqSet{rqPtestExeBOneGptFivePct}{81.8}
\rqSet{rqPtestExeBOneGeminiNum}{29}
\rqSet{rqPtestExeBOneGeminiDen}{36}
\rqSet{rqPtestExeBOneGeminiPct}{80.6}
\rqSet{rqPtestExeBTwoNum}{53}
\rqSet{rqPtestExeBTwoDen}{66}
\rqSet{rqPtestExeBTwoPct}{80.3}
\rqSet{rqPtestExeBTwoKiroNum}{4}
\rqSet{rqPtestExeBTwoKiroDen}{4}
\rqSet{rqPtestExeBTwoKiroPct}{100.0}
\rqSet{rqPtestExeBTwoClaudeNum}{9}
\rqSet{rqPtestExeBTwoClaudeDen}{9}
\rqSet{rqPtestExeBTwoClaudePct}{100.0}
\rqSet{rqPtestExeBTwoCodexNum}{8}
\rqSet{rqPtestExeBTwoCodexDen}{8}
\rqSet{rqPtestExeBTwoCodexPct}{100.0}
\rqSet{rqPtestExeBTwoCTwoRustNum}{8}
\rqSet{rqPtestExeBTwoCTwoRustDen}{10}
\rqSet{rqPtestExeBTwoCTwoRustPct}{80.0}
\rqSet{rqPtestExeBTwoLaertesNum}{8}
\rqSet{rqPtestExeBTwoLaertesDen}{12}
\rqSet{rqPtestExeBTwoLaertesPct}{66.7}
\rqSet{rqPtestExeBTwoCTwoSaferRustNum}{5}
\rqSet{rqPtestExeBTwoCTwoSaferRustDen}{9}
\rqSet{rqPtestExeBTwoCTwoSaferRustPct}{55.6}
\rqSet{rqPtestExeBTwoSmartCTwoRustNum}{2}
\rqSet{rqPtestExeBTwoSmartCTwoRustDen}{2}
\rqSet{rqPtestExeBTwoSmartCTwoRustPct}{100.0}
\rqSet{rqPtestExeBTwoKimiNum}{1}
\rqSet{rqPtestExeBTwoKimiDen}{2}
\rqSet{rqPtestExeBTwoKimiPct}{50.0}
\rqSet{rqPtestExeBTwoGptFiveNum}{3}
\rqSet{rqPtestExeBTwoGptFiveDen}{4}
\rqSet{rqPtestExeBTwoGptFivePct}{75.0}
\rqSet{rqPtestExeBTwoGeminiNum}{5}
\rqSet{rqPtestExeBTwoGeminiDen}{6}
\rqSet{rqPtestExeBTwoGeminiPct}{83.3}
\rqSet{rqPtestExePZeroNum}{0}
\rqSet{rqPtestExePZeroDen}{6}
\rqSet{rqPtestExePZeroPct}{0.0}
\rqSet{rqPtestExePZeroKiroNum}{0}
\rqSet{rqPtestExePZeroKiroDen}{1}
\rqSet{rqPtestExePZeroKiroPct}{0.0}
\rqSet{rqPtestExePZeroClaudeNum}{0}
\rqSet{rqPtestExePZeroClaudeDen}{1}
\rqSet{rqPtestExePZeroClaudePct}{0.0}
\rqSet{rqPtestExePZeroCodexNum}{0}
\rqSet{rqPtestExePZeroCodexDen}{1}
\rqSet{rqPtestExePZeroCodexPct}{0.0}
\rqSet{rqPtestExePZeroCTwoRustNum}{0}
\rqSet{rqPtestExePZeroCTwoRustDen}{1}
\rqSet{rqPtestExePZeroCTwoRustPct}{0.0}
\rqSet{rqPtestExePZeroLaertesNum}{0}
\rqSet{rqPtestExePZeroLaertesDen}{1}
\rqSet{rqPtestExePZeroLaertesPct}{0.0}
\rqSet{rqPtestExePZeroCTwoSaferRustNum}{0}
\rqSet{rqPtestExePZeroCTwoSaferRustDen}{1}
\rqSet{rqPtestExePZeroCTwoSaferRustPct}{0.0}
\rqSet{rqPtestExePZeroSmartCTwoRustNum}{0}
\rqSet{rqPtestExePZeroSmartCTwoRustDen}{0}
\rqSet{rqPtestExePZeroSmartCTwoRustPct}{n/a}
\rqSet{rqPtestExePZeroKimiNum}{0}
\rqSet{rqPtestExePZeroKimiDen}{0}
\rqSet{rqPtestExePZeroKimiPct}{n/a}
\rqSet{rqPtestExePZeroGptFiveNum}{0}
\rqSet{rqPtestExePZeroGptFiveDen}{0}
\rqSet{rqPtestExePZeroGptFivePct}{n/a}
\rqSet{rqPtestExePZeroGeminiNum}{0}
\rqSet{rqPtestExePZeroGeminiDen}{0}
\rqSet{rqPtestExePZeroGeminiPct}{n/a}
\rqSet{rqPtestExePOneNum}{108}
\rqSet{rqPtestExePOneDen}{108}
\rqSet{rqPtestExePOnePct}{100.0}
\rqSet{rqPtestExePOneKiroNum}{48}
\rqSet{rqPtestExePOneKiroDen}{48}
\rqSet{rqPtestExePOneKiroPct}{100.0}
\rqSet{rqPtestExePOneClaudeNum}{48}
\rqSet{rqPtestExePOneClaudeDen}{48}
\rqSet{rqPtestExePOneClaudePct}{100.0}
\rqSet{rqPtestExePOneCodexNum}{12}
\rqSet{rqPtestExePOneCodexDen}{12}
\rqSet{rqPtestExePOneCodexPct}{100.0}
\rqSet{rqPtestExePOneCTwoRustNum}{0}
\rqSet{rqPtestExePOneCTwoRustDen}{0}
\rqSet{rqPtestExePOneCTwoRustPct}{n/a}
\rqSet{rqPtestExePOneLaertesNum}{0}
\rqSet{rqPtestExePOneLaertesDen}{0}
\rqSet{rqPtestExePOneLaertesPct}{n/a}
\rqSet{rqPtestExePOneCTwoSaferRustNum}{0}
\rqSet{rqPtestExePOneCTwoSaferRustDen}{0}
\rqSet{rqPtestExePOneCTwoSaferRustPct}{n/a}
\rqSet{rqPtestExePOneSmartCTwoRustNum}{0}
\rqSet{rqPtestExePOneSmartCTwoRustDen}{0}
\rqSet{rqPtestExePOneSmartCTwoRustPct}{n/a}
\rqSet{rqPtestExePOneKimiNum}{0}
\rqSet{rqPtestExePOneKimiDen}{0}
\rqSet{rqPtestExePOneKimiPct}{n/a}
\rqSet{rqPtestExePOneGptFiveNum}{0}
\rqSet{rqPtestExePOneGptFiveDen}{0}
\rqSet{rqPtestExePOneGptFivePct}{n/a}
\rqSet{rqPtestExePOneGeminiNum}{0}
\rqSet{rqPtestExePOneGeminiDen}{0}
\rqSet{rqPtestExePOneGeminiPct}{n/a}
\rqSet{rqLpExeAllNum}{83}
\rqSet{rqLpExeAllDen}{559}
\rqSet{rqLpExeAllPct}{14.8}
\rqSet{rqLpExeAllKiroNum}{8}
\rqSet{rqLpExeAllKiroDen}{96}
\rqSet{rqLpExeAllKiroPct}{8.3}
\rqSet{rqLpExeAllClaudeNum}{8}
\rqSet{rqLpExeAllClaudeDen}{101}
\rqSet{rqLpExeAllClaudePct}{7.9}
\rqSet{rqLpExeAllCodexNum}{9}
\rqSet{rqLpExeAllCodexDen}{63}
\rqSet{rqLpExeAllCodexPct}{14.3}
\rqSet{rqLpExeAllCTwoRustNum}{10}
\rqSet{rqLpExeAllCTwoRustDen}{54}
\rqSet{rqLpExeAllCTwoRustPct}{18.5}
\rqSet{rqLpExeAllLaertesNum}{12}
\rqSet{rqLpExeAllLaertesDen}{56}
\rqSet{rqLpExeAllLaertesPct}{21.4}
\rqSet{rqLpExeAllCTwoSaferRustNum}{10}
\rqSet{rqLpExeAllCTwoSaferRustDen}{42}
\rqSet{rqLpExeAllCTwoSaferRustPct}{23.8}
\rqSet{rqLpExeAllSmartCTwoRustNum}{7}
\rqSet{rqLpExeAllSmartCTwoRustDen}{40}
\rqSet{rqLpExeAllSmartCTwoRustPct}{17.5}
\rqSet{rqLpExeAllKimiNum}{4}
\rqSet{rqLpExeAllKimiDen}{28}
\rqSet{rqLpExeAllKimiPct}{14.3}
\rqSet{rqLpExeAllGptFiveNum}{7}
\rqSet{rqLpExeAllGptFiveDen}{37}
\rqSet{rqLpExeAllGptFivePct}{18.9}
\rqSet{rqLpExeAllGeminiNum}{8}
\rqSet{rqLpExeAllGeminiDen}{42}
\rqSet{rqLpExeAllGeminiPct}{19.0}
\rqSet{rqLpExeBOneNum}{64}
\rqSet{rqLpExeBOneDen}{379}
\rqSet{rqLpExeBOnePct}{16.9}
\rqSet{rqLpExeBOneKiroNum}{7}
\rqSet{rqLpExeBOneKiroDen}{43}
\rqSet{rqLpExeBOneKiroPct}{16.3}
\rqSet{rqLpExeBOneClaudeNum}{7}
\rqSet{rqLpExeBOneClaudeDen}{43}
\rqSet{rqLpExeBOneClaudePct}{16.3}
\rqSet{rqLpExeBOneCodexNum}{8}
\rqSet{rqLpExeBOneCodexDen}{42}
\rqSet{rqLpExeBOneCodexPct}{19.0}
\rqSet{rqLpExeBOneCTwoRustNum}{7}
\rqSet{rqLpExeBOneCTwoRustDen}{43}
\rqSet{rqLpExeBOneCTwoRustPct}{16.3}
\rqSet{rqLpExeBOneLaertesNum}{7}
\rqSet{rqLpExeBOneLaertesDen}{43}
\rqSet{rqLpExeBOneLaertesPct}{16.3}
\rqSet{rqLpExeBOneCTwoSaferRustNum}{5}
\rqSet{rqLpExeBOneCTwoSaferRustDen}{32}
\rqSet{rqLpExeBOneCTwoSaferRustPct}{15.6}
\rqSet{rqLpExeBOneSmartCTwoRustNum}{7}
\rqSet{rqLpExeBOneSmartCTwoRustDen}{38}
\rqSet{rqLpExeBOneSmartCTwoRustPct}{18.4}
\rqSet{rqLpExeBOneKimiNum}{3}
\rqSet{rqLpExeBOneKimiDen}{26}
\rqSet{rqLpExeBOneKimiPct}{11.5}
\rqSet{rqLpExeBOneGptFiveNum}{6}
\rqSet{rqLpExeBOneGptFiveDen}{33}
\rqSet{rqLpExeBOneGptFivePct}{18.2}
\rqSet{rqLpExeBOneGeminiNum}{7}
\rqSet{rqLpExeBOneGeminiDen}{36}
\rqSet{rqLpExeBOneGeminiPct}{19.4}
\rqSet{rqLpExeBTwoNum}{13}
\rqSet{rqLpExeBTwoDen}{66}
\rqSet{rqLpExeBTwoPct}{19.7}
\rqSet{rqLpExeBTwoKiroNum}{0}
\rqSet{rqLpExeBTwoKiroDen}{4}
\rqSet{rqLpExeBTwoKiroPct}{0.0}
\rqSet{rqLpExeBTwoClaudeNum}{0}
\rqSet{rqLpExeBTwoClaudeDen}{9}
\rqSet{rqLpExeBTwoClaudePct}{0.0}
\rqSet{rqLpExeBTwoCodexNum}{0}
\rqSet{rqLpExeBTwoCodexDen}{8}
\rqSet{rqLpExeBTwoCodexPct}{0.0}
\rqSet{rqLpExeBTwoCTwoRustNum}{2}
\rqSet{rqLpExeBTwoCTwoRustDen}{10}
\rqSet{rqLpExeBTwoCTwoRustPct}{20.0}
\rqSet{rqLpExeBTwoLaertesNum}{4}
\rqSet{rqLpExeBTwoLaertesDen}{12}
\rqSet{rqLpExeBTwoLaertesPct}{33.3}
\rqSet{rqLpExeBTwoCTwoSaferRustNum}{4}
\rqSet{rqLpExeBTwoCTwoSaferRustDen}{9}
\rqSet{rqLpExeBTwoCTwoSaferRustPct}{44.4}
\rqSet{rqLpExeBTwoSmartCTwoRustNum}{0}
\rqSet{rqLpExeBTwoSmartCTwoRustDen}{2}
\rqSet{rqLpExeBTwoSmartCTwoRustPct}{0.0}
\rqSet{rqLpExeBTwoKimiNum}{1}
\rqSet{rqLpExeBTwoKimiDen}{2}
\rqSet{rqLpExeBTwoKimiPct}{50.0}
\rqSet{rqLpExeBTwoGptFiveNum}{1}
\rqSet{rqLpExeBTwoGptFiveDen}{4}
\rqSet{rqLpExeBTwoGptFivePct}{25.0}
\rqSet{rqLpExeBTwoGeminiNum}{1}
\rqSet{rqLpExeBTwoGeminiDen}{6}
\rqSet{rqLpExeBTwoGeminiPct}{16.7}
\rqSet{rqLpExePZeroNum}{6}
\rqSet{rqLpExePZeroDen}{6}
\rqSet{rqLpExePZeroPct}{100.0}
\rqSet{rqLpExePZeroKiroNum}{1}
\rqSet{rqLpExePZeroKiroDen}{1}
\rqSet{rqLpExePZeroKiroPct}{100.0}
\rqSet{rqLpExePZeroClaudeNum}{1}
\rqSet{rqLpExePZeroClaudeDen}{1}
\rqSet{rqLpExePZeroClaudePct}{100.0}
\rqSet{rqLpExePZeroCodexNum}{1}
\rqSet{rqLpExePZeroCodexDen}{1}
\rqSet{rqLpExePZeroCodexPct}{100.0}
\rqSet{rqLpExePZeroCTwoRustNum}{1}
\rqSet{rqLpExePZeroCTwoRustDen}{1}
\rqSet{rqLpExePZeroCTwoRustPct}{100.0}
\rqSet{rqLpExePZeroLaertesNum}{1}
\rqSet{rqLpExePZeroLaertesDen}{1}
\rqSet{rqLpExePZeroLaertesPct}{100.0}
\rqSet{rqLpExePZeroCTwoSaferRustNum}{1}
\rqSet{rqLpExePZeroCTwoSaferRustDen}{1}
\rqSet{rqLpExePZeroCTwoSaferRustPct}{100.0}
\rqSet{rqLpExePZeroSmartCTwoRustNum}{0}
\rqSet{rqLpExePZeroSmartCTwoRustDen}{0}
\rqSet{rqLpExePZeroSmartCTwoRustPct}{n/a}
\rqSet{rqLpExePZeroKimiNum}{0}
\rqSet{rqLpExePZeroKimiDen}{0}
\rqSet{rqLpExePZeroKimiPct}{n/a}
\rqSet{rqLpExePZeroGptFiveNum}{0}
\rqSet{rqLpExePZeroGptFiveDen}{0}
\rqSet{rqLpExePZeroGptFivePct}{n/a}
\rqSet{rqLpExePZeroGeminiNum}{0}
\rqSet{rqLpExePZeroGeminiDen}{0}
\rqSet{rqLpExePZeroGeminiPct}{n/a}
\rqSet{rqLpExePOneNum}{0}
\rqSet{rqLpExePOneDen}{108}
\rqSet{rqLpExePOnePct}{0.0}
\rqSet{rqLpExePOneKiroNum}{0}
\rqSet{rqLpExePOneKiroDen}{48}
\rqSet{rqLpExePOneKiroPct}{0.0}
\rqSet{rqLpExePOneClaudeNum}{0}
\rqSet{rqLpExePOneClaudeDen}{48}
\rqSet{rqLpExePOneClaudePct}{0.0}
\rqSet{rqLpExePOneCodexNum}{0}
\rqSet{rqLpExePOneCodexDen}{12}
\rqSet{rqLpExePOneCodexPct}{0.0}
\rqSet{rqLpExePOneCTwoRustNum}{0}
\rqSet{rqLpExePOneCTwoRustDen}{0}
\rqSet{rqLpExePOneCTwoRustPct}{n/a}
\rqSet{rqLpExePOneLaertesNum}{0}
\rqSet{rqLpExePOneLaertesDen}{0}
\rqSet{rqLpExePOneLaertesPct}{n/a}
\rqSet{rqLpExePOneCTwoSaferRustNum}{0}
\rqSet{rqLpExePOneCTwoSaferRustDen}{0}
\rqSet{rqLpExePOneCTwoSaferRustPct}{n/a}
\rqSet{rqLpExePOneSmartCTwoRustNum}{0}
\rqSet{rqLpExePOneSmartCTwoRustDen}{0}
\rqSet{rqLpExePOneSmartCTwoRustPct}{n/a}
\rqSet{rqLpExePOneKimiNum}{0}
\rqSet{rqLpExePOneKimiDen}{0}
\rqSet{rqLpExePOneKimiPct}{n/a}
\rqSet{rqLpExePOneGptFiveNum}{0}
\rqSet{rqLpExePOneGptFiveDen}{0}
\rqSet{rqLpExePOneGptFivePct}{n/a}
\rqSet{rqLpExePOneGeminiNum}{0}
\rqSet{rqLpExePOneGeminiDen}{0}
\rqSet{rqLpExePOneGeminiPct}{n/a}
\rqSet{rqPreRepBuildExeAllNum}{7}
\rqSet{rqPreRepBuildExeAllDen}{559}
\rqSet{rqPreRepBuildExeAllPct}{1.3}
\rqSet{rqPreRepBuildExeAllKiroNum}{0}
\rqSet{rqPreRepBuildExeAllKiroDen}{96}
\rqSet{rqPreRepBuildExeAllKiroPct}{0.0}
\rqSet{rqPreRepBuildExeAllClaudeNum}{0}
\rqSet{rqPreRepBuildExeAllClaudeDen}{101}
\rqSet{rqPreRepBuildExeAllClaudePct}{0.0}
\rqSet{rqPreRepBuildExeAllCodexNum}{0}
\rqSet{rqPreRepBuildExeAllCodexDen}{63}
\rqSet{rqPreRepBuildExeAllCodexPct}{0.0}
\rqSet{rqPreRepBuildExeAllCTwoRustNum}{1}
\rqSet{rqPreRepBuildExeAllCTwoRustDen}{54}
\rqSet{rqPreRepBuildExeAllCTwoRustPct}{1.9}
\rqSet{rqPreRepBuildExeAllLaertesNum}{3}
\rqSet{rqPreRepBuildExeAllLaertesDen}{56}
\rqSet{rqPreRepBuildExeAllLaertesPct}{5.4}
\rqSet{rqPreRepBuildExeAllCTwoSaferRustNum}{3}
\rqSet{rqPreRepBuildExeAllCTwoSaferRustDen}{42}
\rqSet{rqPreRepBuildExeAllCTwoSaferRustPct}{7.1}
\rqSet{rqPreRepBuildExeAllSmartCTwoRustNum}{0}
\rqSet{rqPreRepBuildExeAllSmartCTwoRustDen}{40}
\rqSet{rqPreRepBuildExeAllSmartCTwoRustPct}{0.0}
\rqSet{rqPreRepBuildExeAllKimiNum}{0}
\rqSet{rqPreRepBuildExeAllKimiDen}{28}
\rqSet{rqPreRepBuildExeAllKimiPct}{0.0}
\rqSet{rqPreRepBuildExeAllGptFiveNum}{0}
\rqSet{rqPreRepBuildExeAllGptFiveDen}{37}
\rqSet{rqPreRepBuildExeAllGptFivePct}{0.0}
\rqSet{rqPreRepBuildExeAllGeminiNum}{0}
\rqSet{rqPreRepBuildExeAllGeminiDen}{42}
\rqSet{rqPreRepBuildExeAllGeminiPct}{0.0}
\rqSet{rqPreRepBuildExeBOneNum}{0}
\rqSet{rqPreRepBuildExeBOneDen}{379}
\rqSet{rqPreRepBuildExeBOnePct}{0.0}
\rqSet{rqPreRepBuildExeBTwoNum}{7}
\rqSet{rqPreRepBuildExeBTwoDen}{66}
\rqSet{rqPreRepBuildExeBTwoPct}{10.6}
\rqSet{rqPreRepBuildExePZeroNum}{0}
\rqSet{rqPreRepBuildExePZeroDen}{6}
\rqSet{rqPreRepBuildExePZeroPct}{0.0}
\rqSet{rqPreRepBuildExePOneNum}{0}
\rqSet{rqPreRepBuildExePOneDen}{108}
\rqSet{rqPreRepBuildExePOnePct}{0.0}
\rqSet{rqPreRepTimeoutExeAllNum}{1}
\rqSet{rqPreRepTimeoutExeAllDen}{559}
\rqSet{rqPreRepTimeoutExeAllPct}{0.2}
\rqSet{rqPreRepTimeoutExeAllKiroNum}{0}
\rqSet{rqPreRepTimeoutExeAllKiroDen}{96}
\rqSet{rqPreRepTimeoutExeAllKiroPct}{0.0}
\rqSet{rqPreRepTimeoutExeAllClaudeNum}{0}
\rqSet{rqPreRepTimeoutExeAllClaudeDen}{101}
\rqSet{rqPreRepTimeoutExeAllClaudePct}{0.0}
\rqSet{rqPreRepTimeoutExeAllCodexNum}{1}
\rqSet{rqPreRepTimeoutExeAllCodexDen}{63}
\rqSet{rqPreRepTimeoutExeAllCodexPct}{1.6}
\rqSet{rqPreRepTimeoutExeAllCTwoRustNum}{0}
\rqSet{rqPreRepTimeoutExeAllCTwoRustDen}{54}
\rqSet{rqPreRepTimeoutExeAllCTwoRustPct}{0.0}
\rqSet{rqPreRepTimeoutExeAllLaertesNum}{0}
\rqSet{rqPreRepTimeoutExeAllLaertesDen}{56}
\rqSet{rqPreRepTimeoutExeAllLaertesPct}{0.0}
\rqSet{rqPreRepTimeoutExeAllCTwoSaferRustNum}{0}
\rqSet{rqPreRepTimeoutExeAllCTwoSaferRustDen}{42}
\rqSet{rqPreRepTimeoutExeAllCTwoSaferRustPct}{0.0}
\rqSet{rqPreRepTimeoutExeAllSmartCTwoRustNum}{0}
\rqSet{rqPreRepTimeoutExeAllSmartCTwoRustDen}{40}
\rqSet{rqPreRepTimeoutExeAllSmartCTwoRustPct}{0.0}
\rqSet{rqPreRepTimeoutExeAllKimiNum}{0}
\rqSet{rqPreRepTimeoutExeAllKimiDen}{28}
\rqSet{rqPreRepTimeoutExeAllKimiPct}{0.0}
\rqSet{rqPreRepTimeoutExeAllGptFiveNum}{0}
\rqSet{rqPreRepTimeoutExeAllGptFiveDen}{37}
\rqSet{rqPreRepTimeoutExeAllGptFivePct}{0.0}
\rqSet{rqPreRepTimeoutExeAllGeminiNum}{0}
\rqSet{rqPreRepTimeoutExeAllGeminiDen}{42}
\rqSet{rqPreRepTimeoutExeAllGeminiPct}{0.0}
\rqSet{rqPreRepTimeoutExeBOneNum}{1}
\rqSet{rqPreRepTimeoutExeBOneDen}{379}
\rqSet{rqPreRepTimeoutExeBOnePct}{0.3}
\rqSet{rqPreRepTimeoutExeBTwoNum}{0}
\rqSet{rqPreRepTimeoutExeBTwoDen}{66}
\rqSet{rqPreRepTimeoutExeBTwoPct}{0.0}
\rqSet{rqPreRepTimeoutExePZeroNum}{0}
\rqSet{rqPreRepTimeoutExePZeroDen}{6}
\rqSet{rqPreRepTimeoutExePZeroPct}{0.0}
\rqSet{rqPreRepTimeoutExePOneNum}{0}
\rqSet{rqPreRepTimeoutExePOneDen}{108}
\rqSet{rqPreRepTimeoutExePOnePct}{0.0}
\rqSet{rqPreRepUnstableExeAllNum}{0}
\rqSet{rqPreRepUnstableExeAllDen}{559}
\rqSet{rqPreRepUnstableExeAllPct}{0.0}
\rqSet{rqPreRepUnstableExeAllKiroNum}{0}
\rqSet{rqPreRepUnstableExeAllKiroDen}{96}
\rqSet{rqPreRepUnstableExeAllKiroPct}{0.0}
\rqSet{rqPreRepUnstableExeAllClaudeNum}{0}
\rqSet{rqPreRepUnstableExeAllClaudeDen}{101}
\rqSet{rqPreRepUnstableExeAllClaudePct}{0.0}
\rqSet{rqPreRepUnstableExeAllCodexNum}{0}
\rqSet{rqPreRepUnstableExeAllCodexDen}{63}
\rqSet{rqPreRepUnstableExeAllCodexPct}{0.0}
\rqSet{rqPreRepUnstableExeAllCTwoRustNum}{0}
\rqSet{rqPreRepUnstableExeAllCTwoRustDen}{54}
\rqSet{rqPreRepUnstableExeAllCTwoRustPct}{0.0}
\rqSet{rqPreRepUnstableExeAllLaertesNum}{0}
\rqSet{rqPreRepUnstableExeAllLaertesDen}{56}
\rqSet{rqPreRepUnstableExeAllLaertesPct}{0.0}
\rqSet{rqPreRepUnstableExeAllCTwoSaferRustNum}{0}
\rqSet{rqPreRepUnstableExeAllCTwoSaferRustDen}{42}
\rqSet{rqPreRepUnstableExeAllCTwoSaferRustPct}{0.0}
\rqSet{rqPreRepUnstableExeAllSmartCTwoRustNum}{0}
\rqSet{rqPreRepUnstableExeAllSmartCTwoRustDen}{40}
\rqSet{rqPreRepUnstableExeAllSmartCTwoRustPct}{0.0}
\rqSet{rqPreRepUnstableExeAllKimiNum}{0}
\rqSet{rqPreRepUnstableExeAllKimiDen}{28}
\rqSet{rqPreRepUnstableExeAllKimiPct}{0.0}
\rqSet{rqPreRepUnstableExeAllGptFiveNum}{0}
\rqSet{rqPreRepUnstableExeAllGptFiveDen}{37}
\rqSet{rqPreRepUnstableExeAllGptFivePct}{0.0}
\rqSet{rqPreRepUnstableExeAllGeminiNum}{0}
\rqSet{rqPreRepUnstableExeAllGeminiDen}{42}
\rqSet{rqPreRepUnstableExeAllGeminiPct}{0.0}
\rqSet{rqPreRepUnstableExeBOneNum}{0}
\rqSet{rqPreRepUnstableExeBOneDen}{379}
\rqSet{rqPreRepUnstableExeBOnePct}{0.0}
\rqSet{rqPreRepUnstableExeBTwoNum}{0}
\rqSet{rqPreRepUnstableExeBTwoDen}{66}
\rqSet{rqPreRepUnstableExeBTwoPct}{0.0}
\rqSet{rqPreRepUnstableExePZeroNum}{0}
\rqSet{rqPreRepUnstableExePZeroDen}{6}
\rqSet{rqPreRepUnstableExePZeroPct}{0.0}
\rqSet{rqPreRepUnstableExePOneNum}{0}
\rqSet{rqPreRepUnstableExePOneDen}{108}
\rqSet{rqPreRepUnstableExePOnePct}{0.0}
\rqSet{rqPreRepNonReproducedExeAllNum}{3}
\rqSet{rqPreRepNonReproducedExeAllDen}{559}
\rqSet{rqPreRepNonReproducedExeAllPct}{0.5}
\rqSet{rqPreRepNonReproducedExeAllKiroNum}{0}
\rqSet{rqPreRepNonReproducedExeAllKiroDen}{96}
\rqSet{rqPreRepNonReproducedExeAllKiroPct}{0.0}
\rqSet{rqPreRepNonReproducedExeAllClaudeNum}{0}
\rqSet{rqPreRepNonReproducedExeAllClaudeDen}{101}
\rqSet{rqPreRepNonReproducedExeAllClaudePct}{0.0}
\rqSet{rqPreRepNonReproducedExeAllCodexNum}{0}
\rqSet{rqPreRepNonReproducedExeAllCodexDen}{63}
\rqSet{rqPreRepNonReproducedExeAllCodexPct}{0.0}
\rqSet{rqPreRepNonReproducedExeAllCTwoRustNum}{0}
\rqSet{rqPreRepNonReproducedExeAllCTwoRustDen}{54}
\rqSet{rqPreRepNonReproducedExeAllCTwoRustPct}{0.0}
\rqSet{rqPreRepNonReproducedExeAllLaertesNum}{0}
\rqSet{rqPreRepNonReproducedExeAllLaertesDen}{56}
\rqSet{rqPreRepNonReproducedExeAllLaertesPct}{0.0}
\rqSet{rqPreRepNonReproducedExeAllCTwoSaferRustNum}{0}
\rqSet{rqPreRepNonReproducedExeAllCTwoSaferRustDen}{42}
\rqSet{rqPreRepNonReproducedExeAllCTwoSaferRustPct}{0.0}
\rqSet{rqPreRepNonReproducedExeAllSmartCTwoRustNum}{0}
\rqSet{rqPreRepNonReproducedExeAllSmartCTwoRustDen}{40}
\rqSet{rqPreRepNonReproducedExeAllSmartCTwoRustPct}{0.0}
\rqSet{rqPreRepNonReproducedExeAllKimiNum}{1}
\rqSet{rqPreRepNonReproducedExeAllKimiDen}{28}
\rqSet{rqPreRepNonReproducedExeAllKimiPct}{3.6}
\rqSet{rqPreRepNonReproducedExeAllGptFiveNum}{1}
\rqSet{rqPreRepNonReproducedExeAllGptFiveDen}{37}
\rqSet{rqPreRepNonReproducedExeAllGptFivePct}{2.7}
\rqSet{rqPreRepNonReproducedExeAllGeminiNum}{1}
\rqSet{rqPreRepNonReproducedExeAllGeminiDen}{42}
\rqSet{rqPreRepNonReproducedExeAllGeminiPct}{2.4}
\rqSet{rqPreRepNonReproducedExeBOneNum}{0}
\rqSet{rqPreRepNonReproducedExeBOneDen}{379}
\rqSet{rqPreRepNonReproducedExeBOnePct}{0.0}
\rqSet{rqPreRepNonReproducedExeBTwoNum}{3}
\rqSet{rqPreRepNonReproducedExeBTwoDen}{66}
\rqSet{rqPreRepNonReproducedExeBTwoPct}{4.5}
\rqSet{rqPreRepNonReproducedExePZeroNum}{0}
\rqSet{rqPreRepNonReproducedExePZeroDen}{6}
\rqSet{rqPreRepNonReproducedExePZeroPct}{0.0}
\rqSet{rqPreRepNonReproducedExePOneNum}{0}
\rqSet{rqPreRepNonReproducedExePOneDen}{108}
\rqSet{rqPreRepNonReproducedExePOnePct}{0.0}
\rqSet{rqLeaveTotalExeAllNum}{72}
\rqSet{rqLeaveTotalExeAllDen}{548}
\rqSet{rqLeaveTotalExeAllPct}{13.1}
\rqSet{rqLeaveTotalExeAllKiroNum}{8}
\rqSet{rqLeaveTotalExeAllKiroDen}{96}
\rqSet{rqLeaveTotalExeAllKiroPct}{8.3}
\rqSet{rqLeaveTotalExeAllClaudeNum}{8}
\rqSet{rqLeaveTotalExeAllClaudeDen}{101}
\rqSet{rqLeaveTotalExeAllClaudePct}{7.9}
\rqSet{rqLeaveTotalExeAllCodexNum}{8}
\rqSet{rqLeaveTotalExeAllCodexDen}{62}
\rqSet{rqLeaveTotalExeAllCodexPct}{12.9}
\rqSet{rqLeaveTotalExeAllCTwoRustNum}{9}
\rqSet{rqLeaveTotalExeAllCTwoRustDen}{53}
\rqSet{rqLeaveTotalExeAllCTwoRustPct}{17.0}
\rqSet{rqLeaveTotalExeAllLaertesNum}{9}
\rqSet{rqLeaveTotalExeAllLaertesDen}{53}
\rqSet{rqLeaveTotalExeAllLaertesPct}{17.0}
\rqSet{rqLeaveTotalExeAllCTwoSaferRustNum}{7}
\rqSet{rqLeaveTotalExeAllCTwoSaferRustDen}{39}
\rqSet{rqLeaveTotalExeAllCTwoSaferRustPct}{17.9}
\rqSet{rqLeaveTotalExeAllSmartCTwoRustNum}{7}
\rqSet{rqLeaveTotalExeAllSmartCTwoRustDen}{40}
\rqSet{rqLeaveTotalExeAllSmartCTwoRustPct}{17.5}
\rqSet{rqLeaveTotalExeAllKimiNum}{3}
\rqSet{rqLeaveTotalExeAllKimiDen}{27}
\rqSet{rqLeaveTotalExeAllKimiPct}{11.1}
\rqSet{rqLeaveTotalExeAllGptFiveNum}{6}
\rqSet{rqLeaveTotalExeAllGptFiveDen}{36}
\rqSet{rqLeaveTotalExeAllGptFivePct}{16.7}
\rqSet{rqLeaveTotalExeAllGeminiNum}{7}
\rqSet{rqLeaveTotalExeAllGeminiDen}{41}
\rqSet{rqLeaveTotalExeAllGeminiPct}{17.1}
\rqSet{rqLeaveTotalExeBOneNum}{63}
\rqSet{rqLeaveTotalExeBOneDen}{378}
\rqSet{rqLeaveTotalExeBOnePct}{16.7}
\rqSet{rqLeaveTotalExeBTwoNum}{3}
\rqSet{rqLeaveTotalExeBTwoDen}{56}
\rqSet{rqLeaveTotalExeBTwoPct}{5.4}
\rqSet{rqLeaveTotalExePZeroNum}{6}
\rqSet{rqLeaveTotalExePZeroDen}{6}
\rqSet{rqLeaveTotalExePZeroPct}{100.0}
\rqSet{rqLeaveTotalExePOneNum}{0}
\rqSet{rqLeaveTotalExePOneDen}{108}
\rqSet{rqLeaveTotalExePOnePct}{0.0}
\rqSet{rqLeaveNoInputFieldsExeAllNum}{0}
\rqSet{rqLeaveNoInputFieldsExeAllDen}{548}
\rqSet{rqLeaveNoInputFieldsExeAllPct}{0.0}
\rqSet{rqLeaveNoInputFieldsExeAllKiroNum}{0}
\rqSet{rqLeaveNoInputFieldsExeAllKiroDen}{96}
\rqSet{rqLeaveNoInputFieldsExeAllKiroPct}{0.0}
\rqSet{rqLeaveNoInputFieldsExeAllClaudeNum}{0}
\rqSet{rqLeaveNoInputFieldsExeAllClaudeDen}{101}
\rqSet{rqLeaveNoInputFieldsExeAllClaudePct}{0.0}
\rqSet{rqLeaveNoInputFieldsExeAllCodexNum}{0}
\rqSet{rqLeaveNoInputFieldsExeAllCodexDen}{62}
\rqSet{rqLeaveNoInputFieldsExeAllCodexPct}{0.0}
\rqSet{rqLeaveNoInputFieldsExeAllCTwoRustNum}{0}
\rqSet{rqLeaveNoInputFieldsExeAllCTwoRustDen}{53}
\rqSet{rqLeaveNoInputFieldsExeAllCTwoRustPct}{0.0}
\rqSet{rqLeaveNoInputFieldsExeAllLaertesNum}{0}
\rqSet{rqLeaveNoInputFieldsExeAllLaertesDen}{53}
\rqSet{rqLeaveNoInputFieldsExeAllLaertesPct}{0.0}
\rqSet{rqLeaveNoInputFieldsExeAllCTwoSaferRustNum}{0}
\rqSet{rqLeaveNoInputFieldsExeAllCTwoSaferRustDen}{39}
\rqSet{rqLeaveNoInputFieldsExeAllCTwoSaferRustPct}{0.0}
\rqSet{rqLeaveNoInputFieldsExeAllSmartCTwoRustNum}{0}
\rqSet{rqLeaveNoInputFieldsExeAllSmartCTwoRustDen}{40}
\rqSet{rqLeaveNoInputFieldsExeAllSmartCTwoRustPct}{0.0}
\rqSet{rqLeaveNoInputFieldsExeAllKimiNum}{0}
\rqSet{rqLeaveNoInputFieldsExeAllKimiDen}{27}
\rqSet{rqLeaveNoInputFieldsExeAllKimiPct}{0.0}
\rqSet{rqLeaveNoInputFieldsExeAllGptFiveNum}{0}
\rqSet{rqLeaveNoInputFieldsExeAllGptFiveDen}{36}
\rqSet{rqLeaveNoInputFieldsExeAllGptFivePct}{0.0}
\rqSet{rqLeaveNoInputFieldsExeAllGeminiNum}{0}
\rqSet{rqLeaveNoInputFieldsExeAllGeminiDen}{41}
\rqSet{rqLeaveNoInputFieldsExeAllGeminiPct}{0.0}
\rqSet{rqLeaveNoInputFieldsExeBOneNum}{0}
\rqSet{rqLeaveNoInputFieldsExeBOneDen}{378}
\rqSet{rqLeaveNoInputFieldsExeBOnePct}{0.0}
\rqSet{rqLeaveNoInputFieldsExeBTwoNum}{0}
\rqSet{rqLeaveNoInputFieldsExeBTwoDen}{56}
\rqSet{rqLeaveNoInputFieldsExeBTwoPct}{0.0}
\rqSet{rqLeaveNoInputFieldsExePZeroNum}{0}
\rqSet{rqLeaveNoInputFieldsExePZeroDen}{6}
\rqSet{rqLeaveNoInputFieldsExePZeroPct}{0.0}
\rqSet{rqLeaveNoInputFieldsExePOneNum}{0}
\rqSet{rqLeaveNoInputFieldsExePOneDen}{108}
\rqSet{rqLeaveNoInputFieldsExePOnePct}{0.0}
\rqSet{rqLeaveCBuildFailedExeAllNum}{0}
\rqSet{rqLeaveCBuildFailedExeAllDen}{548}
\rqSet{rqLeaveCBuildFailedExeAllPct}{0.0}
\rqSet{rqLeaveCBuildFailedExeAllKiroNum}{0}
\rqSet{rqLeaveCBuildFailedExeAllKiroDen}{96}
\rqSet{rqLeaveCBuildFailedExeAllKiroPct}{0.0}
\rqSet{rqLeaveCBuildFailedExeAllClaudeNum}{0}
\rqSet{rqLeaveCBuildFailedExeAllClaudeDen}{101}
\rqSet{rqLeaveCBuildFailedExeAllClaudePct}{0.0}
\rqSet{rqLeaveCBuildFailedExeAllCodexNum}{0}
\rqSet{rqLeaveCBuildFailedExeAllCodexDen}{62}
\rqSet{rqLeaveCBuildFailedExeAllCodexPct}{0.0}
\rqSet{rqLeaveCBuildFailedExeAllCTwoRustNum}{0}
\rqSet{rqLeaveCBuildFailedExeAllCTwoRustDen}{53}
\rqSet{rqLeaveCBuildFailedExeAllCTwoRustPct}{0.0}
\rqSet{rqLeaveCBuildFailedExeAllLaertesNum}{0}
\rqSet{rqLeaveCBuildFailedExeAllLaertesDen}{53}
\rqSet{rqLeaveCBuildFailedExeAllLaertesPct}{0.0}
\rqSet{rqLeaveCBuildFailedExeAllCTwoSaferRustNum}{0}
\rqSet{rqLeaveCBuildFailedExeAllCTwoSaferRustDen}{39}
\rqSet{rqLeaveCBuildFailedExeAllCTwoSaferRustPct}{0.0}
\rqSet{rqLeaveCBuildFailedExeAllSmartCTwoRustNum}{0}
\rqSet{rqLeaveCBuildFailedExeAllSmartCTwoRustDen}{40}
\rqSet{rqLeaveCBuildFailedExeAllSmartCTwoRustPct}{0.0}
\rqSet{rqLeaveCBuildFailedExeAllKimiNum}{0}
\rqSet{rqLeaveCBuildFailedExeAllKimiDen}{27}
\rqSet{rqLeaveCBuildFailedExeAllKimiPct}{0.0}
\rqSet{rqLeaveCBuildFailedExeAllGptFiveNum}{0}
\rqSet{rqLeaveCBuildFailedExeAllGptFiveDen}{36}
\rqSet{rqLeaveCBuildFailedExeAllGptFivePct}{0.0}
\rqSet{rqLeaveCBuildFailedExeAllGeminiNum}{0}
\rqSet{rqLeaveCBuildFailedExeAllGeminiDen}{41}
\rqSet{rqLeaveCBuildFailedExeAllGeminiPct}{0.0}
\rqSet{rqLeaveCBuildFailedExeBOneNum}{0}
\rqSet{rqLeaveCBuildFailedExeBOneDen}{378}
\rqSet{rqLeaveCBuildFailedExeBOnePct}{0.0}
\rqSet{rqLeaveCBuildFailedExeBTwoNum}{0}
\rqSet{rqLeaveCBuildFailedExeBTwoDen}{56}
\rqSet{rqLeaveCBuildFailedExeBTwoPct}{0.0}
\rqSet{rqLeaveCBuildFailedExePZeroNum}{0}
\rqSet{rqLeaveCBuildFailedExePZeroDen}{6}
\rqSet{rqLeaveCBuildFailedExePZeroPct}{0.0}
\rqSet{rqLeaveCBuildFailedExePOneNum}{0}
\rqSet{rqLeaveCBuildFailedExePOneDen}{108}
\rqSet{rqLeaveCBuildFailedExePOnePct}{0.0}
\rqSet{rqLeaveRustBuildFailedExeAllNum}{0}
\rqSet{rqLeaveRustBuildFailedExeAllDen}{548}
\rqSet{rqLeaveRustBuildFailedExeAllPct}{0.0}
\rqSet{rqLeaveRustBuildFailedExeAllKiroNum}{0}
\rqSet{rqLeaveRustBuildFailedExeAllKiroDen}{96}
\rqSet{rqLeaveRustBuildFailedExeAllKiroPct}{0.0}
\rqSet{rqLeaveRustBuildFailedExeAllClaudeNum}{0}
\rqSet{rqLeaveRustBuildFailedExeAllClaudeDen}{101}
\rqSet{rqLeaveRustBuildFailedExeAllClaudePct}{0.0}
\rqSet{rqLeaveRustBuildFailedExeAllCodexNum}{0}
\rqSet{rqLeaveRustBuildFailedExeAllCodexDen}{62}
\rqSet{rqLeaveRustBuildFailedExeAllCodexPct}{0.0}
\rqSet{rqLeaveRustBuildFailedExeAllCTwoRustNum}{0}
\rqSet{rqLeaveRustBuildFailedExeAllCTwoRustDen}{53}
\rqSet{rqLeaveRustBuildFailedExeAllCTwoRustPct}{0.0}
\rqSet{rqLeaveRustBuildFailedExeAllLaertesNum}{0}
\rqSet{rqLeaveRustBuildFailedExeAllLaertesDen}{53}
\rqSet{rqLeaveRustBuildFailedExeAllLaertesPct}{0.0}
\rqSet{rqLeaveRustBuildFailedExeAllCTwoSaferRustNum}{0}
\rqSet{rqLeaveRustBuildFailedExeAllCTwoSaferRustDen}{39}
\rqSet{rqLeaveRustBuildFailedExeAllCTwoSaferRustPct}{0.0}
\rqSet{rqLeaveRustBuildFailedExeAllSmartCTwoRustNum}{0}
\rqSet{rqLeaveRustBuildFailedExeAllSmartCTwoRustDen}{40}
\rqSet{rqLeaveRustBuildFailedExeAllSmartCTwoRustPct}{0.0}
\rqSet{rqLeaveRustBuildFailedExeAllKimiNum}{0}
\rqSet{rqLeaveRustBuildFailedExeAllKimiDen}{27}
\rqSet{rqLeaveRustBuildFailedExeAllKimiPct}{0.0}
\rqSet{rqLeaveRustBuildFailedExeAllGptFiveNum}{0}
\rqSet{rqLeaveRustBuildFailedExeAllGptFiveDen}{36}
\rqSet{rqLeaveRustBuildFailedExeAllGptFivePct}{0.0}
\rqSet{rqLeaveRustBuildFailedExeAllGeminiNum}{0}
\rqSet{rqLeaveRustBuildFailedExeAllGeminiDen}{41}
\rqSet{rqLeaveRustBuildFailedExeAllGeminiPct}{0.0}
\rqSet{rqLeaveRustBuildFailedExeBOneNum}{0}
\rqSet{rqLeaveRustBuildFailedExeBOneDen}{378}
\rqSet{rqLeaveRustBuildFailedExeBOnePct}{0.0}
\rqSet{rqLeaveRustBuildFailedExeBTwoNum}{0}
\rqSet{rqLeaveRustBuildFailedExeBTwoDen}{56}
\rqSet{rqLeaveRustBuildFailedExeBTwoPct}{0.0}
\rqSet{rqLeaveRustBuildFailedExePZeroNum}{0}
\rqSet{rqLeaveRustBuildFailedExePZeroDen}{6}
\rqSet{rqLeaveRustBuildFailedExePZeroPct}{0.0}
\rqSet{rqLeaveRustBuildFailedExePOneNum}{0}
\rqSet{rqLeaveRustBuildFailedExePOneDen}{108}
\rqSet{rqLeaveRustBuildFailedExePOnePct}{0.0}
\rqSet{rqLeaveHarnessErrorExeAllNum}{0}
\rqSet{rqLeaveHarnessErrorExeAllDen}{548}
\rqSet{rqLeaveHarnessErrorExeAllPct}{0.0}
\rqSet{rqLeaveHarnessErrorExeAllKiroNum}{0}
\rqSet{rqLeaveHarnessErrorExeAllKiroDen}{96}
\rqSet{rqLeaveHarnessErrorExeAllKiroPct}{0.0}
\rqSet{rqLeaveHarnessErrorExeAllClaudeNum}{0}
\rqSet{rqLeaveHarnessErrorExeAllClaudeDen}{101}
\rqSet{rqLeaveHarnessErrorExeAllClaudePct}{0.0}
\rqSet{rqLeaveHarnessErrorExeAllCodexNum}{0}
\rqSet{rqLeaveHarnessErrorExeAllCodexDen}{62}
\rqSet{rqLeaveHarnessErrorExeAllCodexPct}{0.0}
\rqSet{rqLeaveHarnessErrorExeAllCTwoRustNum}{0}
\rqSet{rqLeaveHarnessErrorExeAllCTwoRustDen}{53}
\rqSet{rqLeaveHarnessErrorExeAllCTwoRustPct}{0.0}
\rqSet{rqLeaveHarnessErrorExeAllLaertesNum}{0}
\rqSet{rqLeaveHarnessErrorExeAllLaertesDen}{53}
\rqSet{rqLeaveHarnessErrorExeAllLaertesPct}{0.0}
\rqSet{rqLeaveHarnessErrorExeAllCTwoSaferRustNum}{0}
\rqSet{rqLeaveHarnessErrorExeAllCTwoSaferRustDen}{39}
\rqSet{rqLeaveHarnessErrorExeAllCTwoSaferRustPct}{0.0}
\rqSet{rqLeaveHarnessErrorExeAllSmartCTwoRustNum}{0}
\rqSet{rqLeaveHarnessErrorExeAllSmartCTwoRustDen}{40}
\rqSet{rqLeaveHarnessErrorExeAllSmartCTwoRustPct}{0.0}
\rqSet{rqLeaveHarnessErrorExeAllKimiNum}{0}
\rqSet{rqLeaveHarnessErrorExeAllKimiDen}{27}
\rqSet{rqLeaveHarnessErrorExeAllKimiPct}{0.0}
\rqSet{rqLeaveHarnessErrorExeAllGptFiveNum}{0}
\rqSet{rqLeaveHarnessErrorExeAllGptFiveDen}{36}
\rqSet{rqLeaveHarnessErrorExeAllGptFivePct}{0.0}
\rqSet{rqLeaveHarnessErrorExeAllGeminiNum}{0}
\rqSet{rqLeaveHarnessErrorExeAllGeminiDen}{41}
\rqSet{rqLeaveHarnessErrorExeAllGeminiPct}{0.0}
\rqSet{rqLeaveHarnessErrorExeBOneNum}{0}
\rqSet{rqLeaveHarnessErrorExeBOneDen}{378}
\rqSet{rqLeaveHarnessErrorExeBOnePct}{0.0}
\rqSet{rqLeaveHarnessErrorExeBTwoNum}{0}
\rqSet{rqLeaveHarnessErrorExeBTwoDen}{56}
\rqSet{rqLeaveHarnessErrorExeBTwoPct}{0.0}
\rqSet{rqLeaveHarnessErrorExePZeroNum}{0}
\rqSet{rqLeaveHarnessErrorExePZeroDen}{6}
\rqSet{rqLeaveHarnessErrorExePZeroPct}{0.0}
\rqSet{rqLeaveHarnessErrorExePOneNum}{0}
\rqSet{rqLeaveHarnessErrorExePOneDen}{108}
\rqSet{rqLeaveHarnessErrorExePOnePct}{0.0}
\rqSet{rqLeaveNoInputsExeAllNum}{0}
\rqSet{rqLeaveNoInputsExeAllDen}{548}
\rqSet{rqLeaveNoInputsExeAllPct}{0.0}
\rqSet{rqLeaveNoInputsExeAllKiroNum}{0}
\rqSet{rqLeaveNoInputsExeAllKiroDen}{96}
\rqSet{rqLeaveNoInputsExeAllKiroPct}{0.0}
\rqSet{rqLeaveNoInputsExeAllClaudeNum}{0}
\rqSet{rqLeaveNoInputsExeAllClaudeDen}{101}
\rqSet{rqLeaveNoInputsExeAllClaudePct}{0.0}
\rqSet{rqLeaveNoInputsExeAllCodexNum}{0}
\rqSet{rqLeaveNoInputsExeAllCodexDen}{62}
\rqSet{rqLeaveNoInputsExeAllCodexPct}{0.0}
\rqSet{rqLeaveNoInputsExeAllCTwoRustNum}{0}
\rqSet{rqLeaveNoInputsExeAllCTwoRustDen}{53}
\rqSet{rqLeaveNoInputsExeAllCTwoRustPct}{0.0}
\rqSet{rqLeaveNoInputsExeAllLaertesNum}{0}
\rqSet{rqLeaveNoInputsExeAllLaertesDen}{53}
\rqSet{rqLeaveNoInputsExeAllLaertesPct}{0.0}
\rqSet{rqLeaveNoInputsExeAllCTwoSaferRustNum}{0}
\rqSet{rqLeaveNoInputsExeAllCTwoSaferRustDen}{39}
\rqSet{rqLeaveNoInputsExeAllCTwoSaferRustPct}{0.0}
\rqSet{rqLeaveNoInputsExeAllSmartCTwoRustNum}{0}
\rqSet{rqLeaveNoInputsExeAllSmartCTwoRustDen}{40}
\rqSet{rqLeaveNoInputsExeAllSmartCTwoRustPct}{0.0}
\rqSet{rqLeaveNoInputsExeAllKimiNum}{0}
\rqSet{rqLeaveNoInputsExeAllKimiDen}{27}
\rqSet{rqLeaveNoInputsExeAllKimiPct}{0.0}
\rqSet{rqLeaveNoInputsExeAllGptFiveNum}{0}
\rqSet{rqLeaveNoInputsExeAllGptFiveDen}{36}
\rqSet{rqLeaveNoInputsExeAllGptFivePct}{0.0}
\rqSet{rqLeaveNoInputsExeAllGeminiNum}{0}
\rqSet{rqLeaveNoInputsExeAllGeminiDen}{41}
\rqSet{rqLeaveNoInputsExeAllGeminiPct}{0.0}
\rqSet{rqLeaveNoInputsExeBOneNum}{0}
\rqSet{rqLeaveNoInputsExeBOneDen}{378}
\rqSet{rqLeaveNoInputsExeBOnePct}{0.0}
\rqSet{rqLeaveNoInputsExeBTwoNum}{0}
\rqSet{rqLeaveNoInputsExeBTwoDen}{56}
\rqSet{rqLeaveNoInputsExeBTwoPct}{0.0}
\rqSet{rqLeaveNoInputsExePZeroNum}{0}
\rqSet{rqLeaveNoInputsExePZeroDen}{6}
\rqSet{rqLeaveNoInputsExePZeroPct}{0.0}
\rqSet{rqLeaveNoInputsExePOneNum}{0}
\rqSet{rqLeaveNoInputsExePOneDen}{108}
\rqSet{rqLeaveNoInputsExePOnePct}{0.0}
\rqSet{rqLeaveIncompleteExeAllNum}{0}
\rqSet{rqLeaveIncompleteExeAllDen}{548}
\rqSet{rqLeaveIncompleteExeAllPct}{0.0}
\rqSet{rqLeaveIncompleteExeAllKiroNum}{0}
\rqSet{rqLeaveIncompleteExeAllKiroDen}{96}
\rqSet{rqLeaveIncompleteExeAllKiroPct}{0.0}
\rqSet{rqLeaveIncompleteExeAllClaudeNum}{0}
\rqSet{rqLeaveIncompleteExeAllClaudeDen}{101}
\rqSet{rqLeaveIncompleteExeAllClaudePct}{0.0}
\rqSet{rqLeaveIncompleteExeAllCodexNum}{0}
\rqSet{rqLeaveIncompleteExeAllCodexDen}{62}
\rqSet{rqLeaveIncompleteExeAllCodexPct}{0.0}
\rqSet{rqLeaveIncompleteExeAllCTwoRustNum}{0}
\rqSet{rqLeaveIncompleteExeAllCTwoRustDen}{53}
\rqSet{rqLeaveIncompleteExeAllCTwoRustPct}{0.0}
\rqSet{rqLeaveIncompleteExeAllLaertesNum}{0}
\rqSet{rqLeaveIncompleteExeAllLaertesDen}{53}
\rqSet{rqLeaveIncompleteExeAllLaertesPct}{0.0}
\rqSet{rqLeaveIncompleteExeAllCTwoSaferRustNum}{0}
\rqSet{rqLeaveIncompleteExeAllCTwoSaferRustDen}{39}
\rqSet{rqLeaveIncompleteExeAllCTwoSaferRustPct}{0.0}
\rqSet{rqLeaveIncompleteExeAllSmartCTwoRustNum}{0}
\rqSet{rqLeaveIncompleteExeAllSmartCTwoRustDen}{40}
\rqSet{rqLeaveIncompleteExeAllSmartCTwoRustPct}{0.0}
\rqSet{rqLeaveIncompleteExeAllKimiNum}{0}
\rqSet{rqLeaveIncompleteExeAllKimiDen}{27}
\rqSet{rqLeaveIncompleteExeAllKimiPct}{0.0}
\rqSet{rqLeaveIncompleteExeAllGptFiveNum}{0}
\rqSet{rqLeaveIncompleteExeAllGptFiveDen}{36}
\rqSet{rqLeaveIncompleteExeAllGptFivePct}{0.0}
\rqSet{rqLeaveIncompleteExeAllGeminiNum}{0}
\rqSet{rqLeaveIncompleteExeAllGeminiDen}{41}
\rqSet{rqLeaveIncompleteExeAllGeminiPct}{0.0}
\rqSet{rqLeaveIncompleteExeBOneNum}{0}
\rqSet{rqLeaveIncompleteExeBOneDen}{378}
\rqSet{rqLeaveIncompleteExeBOnePct}{0.0}
\rqSet{rqLeaveIncompleteExeBTwoNum}{0}
\rqSet{rqLeaveIncompleteExeBTwoDen}{56}
\rqSet{rqLeaveIncompleteExeBTwoPct}{0.0}
\rqSet{rqLeaveIncompleteExePZeroNum}{0}
\rqSet{rqLeaveIncompleteExePZeroDen}{6}
\rqSet{rqLeaveIncompleteExePZeroPct}{0.0}
\rqSet{rqLeaveIncompleteExePOneNum}{0}
\rqSet{rqLeaveIncompleteExePOneDen}{108}
\rqSet{rqLeaveIncompleteExePOnePct}{0.0}
\rqSet{rqLeaveSanInfraExeAllNum}{0}
\rqSet{rqLeaveSanInfraExeAllDen}{548}
\rqSet{rqLeaveSanInfraExeAllPct}{0.0}
\rqSet{rqLeaveSanInfraExeAllKiroNum}{0}
\rqSet{rqLeaveSanInfraExeAllKiroDen}{96}
\rqSet{rqLeaveSanInfraExeAllKiroPct}{0.0}
\rqSet{rqLeaveSanInfraExeAllClaudeNum}{0}
\rqSet{rqLeaveSanInfraExeAllClaudeDen}{101}
\rqSet{rqLeaveSanInfraExeAllClaudePct}{0.0}
\rqSet{rqLeaveSanInfraExeAllCodexNum}{0}
\rqSet{rqLeaveSanInfraExeAllCodexDen}{62}
\rqSet{rqLeaveSanInfraExeAllCodexPct}{0.0}
\rqSet{rqLeaveSanInfraExeAllCTwoRustNum}{0}
\rqSet{rqLeaveSanInfraExeAllCTwoRustDen}{53}
\rqSet{rqLeaveSanInfraExeAllCTwoRustPct}{0.0}
\rqSet{rqLeaveSanInfraExeAllLaertesNum}{0}
\rqSet{rqLeaveSanInfraExeAllLaertesDen}{53}
\rqSet{rqLeaveSanInfraExeAllLaertesPct}{0.0}
\rqSet{rqLeaveSanInfraExeAllCTwoSaferRustNum}{0}
\rqSet{rqLeaveSanInfraExeAllCTwoSaferRustDen}{39}
\rqSet{rqLeaveSanInfraExeAllCTwoSaferRustPct}{0.0}
\rqSet{rqLeaveSanInfraExeAllSmartCTwoRustNum}{0}
\rqSet{rqLeaveSanInfraExeAllSmartCTwoRustDen}{40}
\rqSet{rqLeaveSanInfraExeAllSmartCTwoRustPct}{0.0}
\rqSet{rqLeaveSanInfraExeAllKimiNum}{0}
\rqSet{rqLeaveSanInfraExeAllKimiDen}{27}
\rqSet{rqLeaveSanInfraExeAllKimiPct}{0.0}
\rqSet{rqLeaveSanInfraExeAllGptFiveNum}{0}
\rqSet{rqLeaveSanInfraExeAllGptFiveDen}{36}
\rqSet{rqLeaveSanInfraExeAllGptFivePct}{0.0}
\rqSet{rqLeaveSanInfraExeAllGeminiNum}{0}
\rqSet{rqLeaveSanInfraExeAllGeminiDen}{41}
\rqSet{rqLeaveSanInfraExeAllGeminiPct}{0.0}
\rqSet{rqLeaveSanInfraExeBOneNum}{0}
\rqSet{rqLeaveSanInfraExeBOneDen}{378}
\rqSet{rqLeaveSanInfraExeBOnePct}{0.0}
\rqSet{rqLeaveSanInfraExeBTwoNum}{0}
\rqSet{rqLeaveSanInfraExeBTwoDen}{56}
\rqSet{rqLeaveSanInfraExeBTwoPct}{0.0}
\rqSet{rqLeaveSanInfraExePZeroNum}{0}
\rqSet{rqLeaveSanInfraExePZeroDen}{6}
\rqSet{rqLeaveSanInfraExePZeroPct}{0.0}
\rqSet{rqLeaveSanInfraExePOneNum}{0}
\rqSet{rqLeaveSanInfraExePOneDen}{108}
\rqSet{rqLeaveSanInfraExePOnePct}{0.0}
\rqSet{rqLeaveMissingSeedExeAllNum}{28}
\rqSet{rqLeaveMissingSeedExeAllDen}{548}
\rqSet{rqLeaveMissingSeedExeAllPct}{5.1}
\rqSet{rqLeaveMissingSeedExeAllKiroNum}{3}
\rqSet{rqLeaveMissingSeedExeAllKiroDen}{96}
\rqSet{rqLeaveMissingSeedExeAllKiroPct}{3.1}
\rqSet{rqLeaveMissingSeedExeAllClaudeNum}{3}
\rqSet{rqLeaveMissingSeedExeAllClaudeDen}{101}
\rqSet{rqLeaveMissingSeedExeAllClaudePct}{3.0}
\rqSet{rqLeaveMissingSeedExeAllCodexNum}{3}
\rqSet{rqLeaveMissingSeedExeAllCodexDen}{62}
\rqSet{rqLeaveMissingSeedExeAllCodexPct}{4.8}
\rqSet{rqLeaveMissingSeedExeAllCTwoRustNum}{4}
\rqSet{rqLeaveMissingSeedExeAllCTwoRustDen}{53}
\rqSet{rqLeaveMissingSeedExeAllCTwoRustPct}{7.5}
\rqSet{rqLeaveMissingSeedExeAllLaertesNum}{4}
\rqSet{rqLeaveMissingSeedExeAllLaertesDen}{53}
\rqSet{rqLeaveMissingSeedExeAllLaertesPct}{7.5}
\rqSet{rqLeaveMissingSeedExeAllCTwoSaferRustNum}{4}
\rqSet{rqLeaveMissingSeedExeAllCTwoSaferRustDen}{39}
\rqSet{rqLeaveMissingSeedExeAllCTwoSaferRustPct}{10.3}
\rqSet{rqLeaveMissingSeedExeAllSmartCTwoRustNum}{2}
\rqSet{rqLeaveMissingSeedExeAllSmartCTwoRustDen}{40}
\rqSet{rqLeaveMissingSeedExeAllSmartCTwoRustPct}{5.0}
\rqSet{rqLeaveMissingSeedExeAllKimiNum}{1}
\rqSet{rqLeaveMissingSeedExeAllKimiDen}{27}
\rqSet{rqLeaveMissingSeedExeAllKimiPct}{3.7}
\rqSet{rqLeaveMissingSeedExeAllGptFiveNum}{2}
\rqSet{rqLeaveMissingSeedExeAllGptFiveDen}{36}
\rqSet{rqLeaveMissingSeedExeAllGptFivePct}{5.6}
\rqSet{rqLeaveMissingSeedExeAllGeminiNum}{2}
\rqSet{rqLeaveMissingSeedExeAllGeminiDen}{41}
\rqSet{rqLeaveMissingSeedExeAllGeminiPct}{4.9}
\rqSet{rqLeaveMissingSeedExeBOneNum}{19}
\rqSet{rqLeaveMissingSeedExeBOneDen}{378}
\rqSet{rqLeaveMissingSeedExeBOnePct}{5.0}
\rqSet{rqLeaveMissingSeedExeBTwoNum}{3}
\rqSet{rqLeaveMissingSeedExeBTwoDen}{56}
\rqSet{rqLeaveMissingSeedExeBTwoPct}{5.4}
\rqSet{rqLeaveMissingSeedExePZeroNum}{6}
\rqSet{rqLeaveMissingSeedExePZeroDen}{6}
\rqSet{rqLeaveMissingSeedExePZeroPct}{100.0}
\rqSet{rqLeaveMissingSeedExePOneNum}{0}
\rqSet{rqLeaveMissingSeedExePOneDen}{108}
\rqSet{rqLeaveMissingSeedExePOnePct}{0.0}
\rqSet{rqLeaveMissingRecordExeAllNum}{44}
\rqSet{rqLeaveMissingRecordExeAllDen}{548}
\rqSet{rqLeaveMissingRecordExeAllPct}{8.0}
\rqSet{rqLeaveMissingRecordExeAllKiroNum}{5}
\rqSet{rqLeaveMissingRecordExeAllKiroDen}{96}
\rqSet{rqLeaveMissingRecordExeAllKiroPct}{5.2}
\rqSet{rqLeaveMissingRecordExeAllClaudeNum}{5}
\rqSet{rqLeaveMissingRecordExeAllClaudeDen}{101}
\rqSet{rqLeaveMissingRecordExeAllClaudePct}{5.0}
\rqSet{rqLeaveMissingRecordExeAllCodexNum}{5}
\rqSet{rqLeaveMissingRecordExeAllCodexDen}{62}
\rqSet{rqLeaveMissingRecordExeAllCodexPct}{8.1}
\rqSet{rqLeaveMissingRecordExeAllCTwoRustNum}{5}
\rqSet{rqLeaveMissingRecordExeAllCTwoRustDen}{53}
\rqSet{rqLeaveMissingRecordExeAllCTwoRustPct}{9.4}
\rqSet{rqLeaveMissingRecordExeAllLaertesNum}{5}
\rqSet{rqLeaveMissingRecordExeAllLaertesDen}{53}
\rqSet{rqLeaveMissingRecordExeAllLaertesPct}{9.4}
\rqSet{rqLeaveMissingRecordExeAllCTwoSaferRustNum}{3}
\rqSet{rqLeaveMissingRecordExeAllCTwoSaferRustDen}{39}
\rqSet{rqLeaveMissingRecordExeAllCTwoSaferRustPct}{7.7}
\rqSet{rqLeaveMissingRecordExeAllSmartCTwoRustNum}{5}
\rqSet{rqLeaveMissingRecordExeAllSmartCTwoRustDen}{40}
\rqSet{rqLeaveMissingRecordExeAllSmartCTwoRustPct}{12.5}
\rqSet{rqLeaveMissingRecordExeAllKimiNum}{2}
\rqSet{rqLeaveMissingRecordExeAllKimiDen}{27}
\rqSet{rqLeaveMissingRecordExeAllKimiPct}{7.4}
\rqSet{rqLeaveMissingRecordExeAllGptFiveNum}{4}
\rqSet{rqLeaveMissingRecordExeAllGptFiveDen}{36}
\rqSet{rqLeaveMissingRecordExeAllGptFivePct}{11.1}
\rqSet{rqLeaveMissingRecordExeAllGeminiNum}{5}
\rqSet{rqLeaveMissingRecordExeAllGeminiDen}{41}
\rqSet{rqLeaveMissingRecordExeAllGeminiPct}{12.2}
\rqSet{rqLeaveMissingRecordExeBOneNum}{44}
\rqSet{rqLeaveMissingRecordExeBOneDen}{378}
\rqSet{rqLeaveMissingRecordExeBOnePct}{11.6}
\rqSet{rqLeaveMissingRecordExeBTwoNum}{0}
\rqSet{rqLeaveMissingRecordExeBTwoDen}{56}
\rqSet{rqLeaveMissingRecordExeBTwoPct}{0.0}
\rqSet{rqLeaveMissingRecordExePZeroNum}{0}
\rqSet{rqLeaveMissingRecordExePZeroDen}{6}
\rqSet{rqLeaveMissingRecordExePZeroPct}{0.0}
\rqSet{rqLeaveMissingRecordExePOneNum}{0}
\rqSet{rqLeaveMissingRecordExePOneDen}{108}
\rqSet{rqLeaveMissingRecordExePOnePct}{0.0}
\rqSet{rqLeaveOtherExeAllNum}{0}
\rqSet{rqLeaveOtherExeAllDen}{548}
\rqSet{rqLeaveOtherExeAllPct}{0.0}
\rqSet{rqLeaveOtherExeAllKiroNum}{0}
\rqSet{rqLeaveOtherExeAllKiroDen}{96}
\rqSet{rqLeaveOtherExeAllKiroPct}{0.0}
\rqSet{rqLeaveOtherExeAllClaudeNum}{0}
\rqSet{rqLeaveOtherExeAllClaudeDen}{101}
\rqSet{rqLeaveOtherExeAllClaudePct}{0.0}
\rqSet{rqLeaveOtherExeAllCodexNum}{0}
\rqSet{rqLeaveOtherExeAllCodexDen}{62}
\rqSet{rqLeaveOtherExeAllCodexPct}{0.0}
\rqSet{rqLeaveOtherExeAllCTwoRustNum}{0}
\rqSet{rqLeaveOtherExeAllCTwoRustDen}{53}
\rqSet{rqLeaveOtherExeAllCTwoRustPct}{0.0}
\rqSet{rqLeaveOtherExeAllLaertesNum}{0}
\rqSet{rqLeaveOtherExeAllLaertesDen}{53}
\rqSet{rqLeaveOtherExeAllLaertesPct}{0.0}
\rqSet{rqLeaveOtherExeAllCTwoSaferRustNum}{0}
\rqSet{rqLeaveOtherExeAllCTwoSaferRustDen}{39}
\rqSet{rqLeaveOtherExeAllCTwoSaferRustPct}{0.0}
\rqSet{rqLeaveOtherExeAllSmartCTwoRustNum}{0}
\rqSet{rqLeaveOtherExeAllSmartCTwoRustDen}{40}
\rqSet{rqLeaveOtherExeAllSmartCTwoRustPct}{0.0}
\rqSet{rqLeaveOtherExeAllKimiNum}{0}
\rqSet{rqLeaveOtherExeAllKimiDen}{27}
\rqSet{rqLeaveOtherExeAllKimiPct}{0.0}
\rqSet{rqLeaveOtherExeAllGptFiveNum}{0}
\rqSet{rqLeaveOtherExeAllGptFiveDen}{36}
\rqSet{rqLeaveOtherExeAllGptFivePct}{0.0}
\rqSet{rqLeaveOtherExeAllGeminiNum}{0}
\rqSet{rqLeaveOtherExeAllGeminiDen}{41}
\rqSet{rqLeaveOtherExeAllGeminiPct}{0.0}
\rqSet{rqLeaveOtherExeBOneNum}{0}
\rqSet{rqLeaveOtherExeBOneDen}{378}
\rqSet{rqLeaveOtherExeBOnePct}{0.0}
\rqSet{rqLeaveOtherExeBTwoNum}{0}
\rqSet{rqLeaveOtherExeBTwoDen}{56}
\rqSet{rqLeaveOtherExeBTwoPct}{0.0}
\rqSet{rqLeaveOtherExePZeroNum}{0}
\rqSet{rqLeaveOtherExePZeroDen}{6}
\rqSet{rqLeaveOtherExePZeroPct}{0.0}
\rqSet{rqLeaveOtherExePOneNum}{0}
\rqSet{rqLeaveOtherExePOneDen}{108}
\rqSet{rqLeaveOtherExePOnePct}{0.0}
\rqSet{rqUnexpectedExeAllNum}{0}
\rqSet{rqUnexpectedExeAllDen}{548}
\rqSet{rqUnexpectedExeBOneNum}{0}
\rqSet{rqUnexpectedExeBOneDen}{378}
\rqSet{rqUnexpectedExeBTwoNum}{0}
\rqSet{rqUnexpectedExeBTwoDen}{56}
\rqSet{rqUnexpectedExePZeroNum}{0}
\rqSet{rqUnexpectedExePZeroDen}{6}
\rqSet{rqUnexpectedExePOneNum}{0}
\rqSet{rqUnexpectedExePOneDen}{108}
\rqSet{rqFPExeAllNum}{176}
\rqSet{rqFPExeAllDen}{476}
\rqSet{rqFPExeAllPct}{37.0}
\rqSet{rqFPExeAllKiroNum}{29}
\rqSet{rqFPExeAllKiroDen}{88}
\rqSet{rqFPExeAllKiroPct}{33.0}
\rqSet{rqFPExeAllClaudeNum}{24}
\rqSet{rqFPExeAllClaudeDen}{93}
\rqSet{rqFPExeAllClaudePct}{25.8}
\rqSet{rqFPExeAllCodexNum}{9}
\rqSet{rqFPExeAllCodexDen}{54}
\rqSet{rqFPExeAllCodexPct}{16.7}
\rqSet{rqFPExeAllCTwoRustNum}{3}
\rqSet{rqFPExeAllCTwoRustDen}{44}
\rqSet{rqFPExeAllCTwoRustPct}{6.8}
\rqSet{rqFPExeAllLaertesNum}{3}
\rqSet{rqFPExeAllLaertesDen}{44}
\rqSet{rqFPExeAllLaertesPct}{6.8}
\rqSet{rqFPExeAllCTwoSaferRustNum}{12}
\rqSet{rqFPExeAllCTwoSaferRustDen}{32}
\rqSet{rqFPExeAllCTwoSaferRustPct}{37.5}
\rqSet{rqFPExeAllSmartCTwoRustNum}{27}
\rqSet{rqFPExeAllSmartCTwoRustDen}{33}
\rqSet{rqFPExeAllSmartCTwoRustPct}{81.8}
\rqSet{rqFPExeAllKimiNum}{18}
\rqSet{rqFPExeAllKimiDen}{24}
\rqSet{rqFPExeAllKimiPct}{75.0}
\rqSet{rqFPExeAllGptFiveNum}{25}
\rqSet{rqFPExeAllGptFiveDen}{30}
\rqSet{rqFPExeAllGptFivePct}{83.3}
\rqSet{rqFPExeAllGeminiNum}{26}
\rqSet{rqFPExeAllGeminiDen}{34}
\rqSet{rqFPExeAllGeminiPct}{76.5}
\rqSet{rqFPExeBOneNum}{152}
\rqSet{rqFPExeBOneDen}{315}
\rqSet{rqFPExeBOnePct}{48.3}
\rqSet{rqFPExeBOneKiroNum}{26}
\rqSet{rqFPExeBOneKiroDen}{36}
\rqSet{rqFPExeBOneKiroPct}{72.2}
\rqSet{rqFPExeBOneClaudeNum}{18}
\rqSet{rqFPExeBOneClaudeDen}{36}
\rqSet{rqFPExeBOneClaudePct}{50.0}
\rqSet{rqFPExeBOneCodexNum}{6}
\rqSet{rqFPExeBOneCodexDen}{34}
\rqSet{rqFPExeBOneCodexPct}{17.6}
\rqSet{rqFPExeBOneCTwoRustNum}{1}
\rqSet{rqFPExeBOneCTwoRustDen}{36}
\rqSet{rqFPExeBOneCTwoRustPct}{2.8}
\rqSet{rqFPExeBOneLaertesNum}{1}
\rqSet{rqFPExeBOneLaertesDen}{36}
\rqSet{rqFPExeBOneLaertesPct}{2.8}
\rqSet{rqFPExeBOneCTwoSaferRustNum}{10}
\rqSet{rqFPExeBOneCTwoSaferRustDen}{27}
\rqSet{rqFPExeBOneCTwoSaferRustPct}{37.0}
\rqSet{rqFPExeBOneSmartCTwoRustNum}{26}
\rqSet{rqFPExeBOneSmartCTwoRustDen}{31}
\rqSet{rqFPExeBOneSmartCTwoRustPct}{83.9}
\rqSet{rqFPExeBOneKimiNum}{18}
\rqSet{rqFPExeBOneKimiDen}{23}
\rqSet{rqFPExeBOneKimiPct}{78.3}
\rqSet{rqFPExeBOneGptFiveNum}{23}
\rqSet{rqFPExeBOneGptFiveDen}{27}
\rqSet{rqFPExeBOneGptFivePct}{85.2}
\rqSet{rqFPExeBOneGeminiNum}{23}
\rqSet{rqFPExeBOneGeminiDen}{29}
\rqSet{rqFPExeBOneGeminiPct}{79.3}
\rqSet{rqFPExeBTwoNum}{24}
\rqSet{rqFPExeBTwoDen}{53}
\rqSet{rqFPExeBTwoPct}{45.3}
\rqSet{rqFPExeBTwoKiroNum}{3}
\rqSet{rqFPExeBTwoKiroDen}{4}
\rqSet{rqFPExeBTwoKiroPct}{75.0}
\rqSet{rqFPExeBTwoClaudeNum}{6}
\rqSet{rqFPExeBTwoClaudeDen}{9}
\rqSet{rqFPExeBTwoClaudePct}{66.7}
\rqSet{rqFPExeBTwoCodexNum}{3}
\rqSet{rqFPExeBTwoCodexDen}{8}
\rqSet{rqFPExeBTwoCodexPct}{37.5}
\rqSet{rqFPExeBTwoCTwoRustNum}{2}
\rqSet{rqFPExeBTwoCTwoRustDen}{8}
\rqSet{rqFPExeBTwoCTwoRustPct}{25.0}
\rqSet{rqFPExeBTwoLaertesNum}{2}
\rqSet{rqFPExeBTwoLaertesDen}{8}
\rqSet{rqFPExeBTwoLaertesPct}{25.0}
\rqSet{rqFPExeBTwoCTwoSaferRustNum}{2}
\rqSet{rqFPExeBTwoCTwoSaferRustDen}{5}
\rqSet{rqFPExeBTwoCTwoSaferRustPct}{40.0}
\rqSet{rqFPExeBTwoSmartCTwoRustNum}{1}
\rqSet{rqFPExeBTwoSmartCTwoRustDen}{2}
\rqSet{rqFPExeBTwoSmartCTwoRustPct}{50.0}
\rqSet{rqFPExeBTwoKimiNum}{0}
\rqSet{rqFPExeBTwoKimiDen}{1}
\rqSet{rqFPExeBTwoKimiPct}{0.0}
\rqSet{rqFPExeBTwoGptFiveNum}{2}
\rqSet{rqFPExeBTwoGptFiveDen}{3}
\rqSet{rqFPExeBTwoGptFivePct}{66.7}
\rqSet{rqFPExeBTwoGeminiNum}{3}
\rqSet{rqFPExeBTwoGeminiDen}{5}
\rqSet{rqFPExeBTwoGeminiPct}{60.0}
\rqSet{rqFPExePZeroNum}{0}
\rqSet{rqFPExePZeroDen}{0}
\rqSet{rqFPExePZeroPct}{n/a}
\rqSet{rqFPExePZeroKiroNum}{0}
\rqSet{rqFPExePZeroKiroDen}{0}
\rqSet{rqFPExePZeroKiroPct}{n/a}
\rqSet{rqFPExePZeroClaudeNum}{0}
\rqSet{rqFPExePZeroClaudeDen}{0}
\rqSet{rqFPExePZeroClaudePct}{n/a}
\rqSet{rqFPExePZeroCodexNum}{0}
\rqSet{rqFPExePZeroCodexDen}{0}
\rqSet{rqFPExePZeroCodexPct}{n/a}
\rqSet{rqFPExePZeroCTwoRustNum}{0}
\rqSet{rqFPExePZeroCTwoRustDen}{0}
\rqSet{rqFPExePZeroCTwoRustPct}{n/a}
\rqSet{rqFPExePZeroLaertesNum}{0}
\rqSet{rqFPExePZeroLaertesDen}{0}
\rqSet{rqFPExePZeroLaertesPct}{n/a}
\rqSet{rqFPExePZeroCTwoSaferRustNum}{0}
\rqSet{rqFPExePZeroCTwoSaferRustDen}{0}
\rqSet{rqFPExePZeroCTwoSaferRustPct}{n/a}
\rqSet{rqFPExePZeroSmartCTwoRustNum}{0}
\rqSet{rqFPExePZeroSmartCTwoRustDen}{0}
\rqSet{rqFPExePZeroSmartCTwoRustPct}{n/a}
\rqSet{rqFPExePZeroKimiNum}{0}
\rqSet{rqFPExePZeroKimiDen}{0}
\rqSet{rqFPExePZeroKimiPct}{n/a}
\rqSet{rqFPExePZeroGptFiveNum}{0}
\rqSet{rqFPExePZeroGptFiveDen}{0}
\rqSet{rqFPExePZeroGptFivePct}{n/a}
\rqSet{rqFPExePZeroGeminiNum}{0}
\rqSet{rqFPExePZeroGeminiDen}{0}
\rqSet{rqFPExePZeroGeminiPct}{n/a}
\rqSet{rqFPExePOneNum}{0}
\rqSet{rqFPExePOneDen}{108}
\rqSet{rqFPExePOnePct}{0.0}
\rqSet{rqFPExePOneKiroNum}{0}
\rqSet{rqFPExePOneKiroDen}{48}
\rqSet{rqFPExePOneKiroPct}{0.0}
\rqSet{rqFPExePOneClaudeNum}{0}
\rqSet{rqFPExePOneClaudeDen}{48}
\rqSet{rqFPExePOneClaudePct}{0.0}
\rqSet{rqFPExePOneCodexNum}{0}
\rqSet{rqFPExePOneCodexDen}{12}
\rqSet{rqFPExePOneCodexPct}{0.0}
\rqSet{rqFPExePOneCTwoRustNum}{0}
\rqSet{rqFPExePOneCTwoRustDen}{0}
\rqSet{rqFPExePOneCTwoRustPct}{n/a}
\rqSet{rqFPExePOneLaertesNum}{0}
\rqSet{rqFPExePOneLaertesDen}{0}
\rqSet{rqFPExePOneLaertesPct}{n/a}
\rqSet{rqFPExePOneCTwoSaferRustNum}{0}
\rqSet{rqFPExePOneCTwoSaferRustDen}{0}
\rqSet{rqFPExePOneCTwoSaferRustPct}{n/a}
\rqSet{rqFPExePOneSmartCTwoRustNum}{0}
\rqSet{rqFPExePOneSmartCTwoRustDen}{0}
\rqSet{rqFPExePOneSmartCTwoRustPct}{n/a}
\rqSet{rqFPExePOneKimiNum}{0}
\rqSet{rqFPExePOneKimiDen}{0}
\rqSet{rqFPExePOneKimiPct}{n/a}
\rqSet{rqFPExePOneGptFiveNum}{0}
\rqSet{rqFPExePOneGptFiveDen}{0}
\rqSet{rqFPExePOneGptFivePct}{n/a}
\rqSet{rqFPExePOneGeminiNum}{0}
\rqSet{rqFPExePOneGeminiDen}{0}
\rqSet{rqFPExePOneGeminiPct}{n/a}
\rqSet{rqFPBoundLoExeAllNum}{176}
\rqSet{rqFPBoundLoExeAllDen}{559}
\rqSet{rqFPBoundLoExeAllPct}{31.5}
\rqSet{rqFPBoundLoExeAllKiroNum}{29}
\rqSet{rqFPBoundLoExeAllKiroDen}{96}
\rqSet{rqFPBoundLoExeAllKiroPct}{30.2}
\rqSet{rqFPBoundLoExeAllClaudeNum}{24}
\rqSet{rqFPBoundLoExeAllClaudeDen}{101}
\rqSet{rqFPBoundLoExeAllClaudePct}{23.8}
\rqSet{rqFPBoundLoExeAllCodexNum}{9}
\rqSet{rqFPBoundLoExeAllCodexDen}{63}
\rqSet{rqFPBoundLoExeAllCodexPct}{14.3}
\rqSet{rqFPBoundLoExeAllCTwoRustNum}{3}
\rqSet{rqFPBoundLoExeAllCTwoRustDen}{54}
\rqSet{rqFPBoundLoExeAllCTwoRustPct}{5.6}
\rqSet{rqFPBoundLoExeAllLaertesNum}{3}
\rqSet{rqFPBoundLoExeAllLaertesDen}{56}
\rqSet{rqFPBoundLoExeAllLaertesPct}{5.4}
\rqSet{rqFPBoundLoExeAllCTwoSaferRustNum}{12}
\rqSet{rqFPBoundLoExeAllCTwoSaferRustDen}{42}
\rqSet{rqFPBoundLoExeAllCTwoSaferRustPct}{28.6}
\rqSet{rqFPBoundLoExeAllSmartCTwoRustNum}{27}
\rqSet{rqFPBoundLoExeAllSmartCTwoRustDen}{40}
\rqSet{rqFPBoundLoExeAllSmartCTwoRustPct}{67.5}
\rqSet{rqFPBoundLoExeAllKimiNum}{18}
\rqSet{rqFPBoundLoExeAllKimiDen}{28}
\rqSet{rqFPBoundLoExeAllKimiPct}{64.3}
\rqSet{rqFPBoundLoExeAllGptFiveNum}{25}
\rqSet{rqFPBoundLoExeAllGptFiveDen}{37}
\rqSet{rqFPBoundLoExeAllGptFivePct}{67.6}
\rqSet{rqFPBoundLoExeAllGeminiNum}{26}
\rqSet{rqFPBoundLoExeAllGeminiDen}{42}
\rqSet{rqFPBoundLoExeAllGeminiPct}{61.9}
\rqSet{rqFPBoundLoExeBOneNum}{152}
\rqSet{rqFPBoundLoExeBOneDen}{379}
\rqSet{rqFPBoundLoExeBOnePct}{40.1}
\rqSet{rqFPBoundLoExeBOneKiroNum}{26}
\rqSet{rqFPBoundLoExeBOneKiroDen}{43}
\rqSet{rqFPBoundLoExeBOneKiroPct}{60.5}
\rqSet{rqFPBoundLoExeBOneClaudeNum}{18}
\rqSet{rqFPBoundLoExeBOneClaudeDen}{43}
\rqSet{rqFPBoundLoExeBOneClaudePct}{41.9}
\rqSet{rqFPBoundLoExeBOneCodexNum}{6}
\rqSet{rqFPBoundLoExeBOneCodexDen}{42}
\rqSet{rqFPBoundLoExeBOneCodexPct}{14.3}
\rqSet{rqFPBoundLoExeBOneCTwoRustNum}{1}
\rqSet{rqFPBoundLoExeBOneCTwoRustDen}{43}
\rqSet{rqFPBoundLoExeBOneCTwoRustPct}{2.3}
\rqSet{rqFPBoundLoExeBOneLaertesNum}{1}
\rqSet{rqFPBoundLoExeBOneLaertesDen}{43}
\rqSet{rqFPBoundLoExeBOneLaertesPct}{2.3}
\rqSet{rqFPBoundLoExeBOneCTwoSaferRustNum}{10}
\rqSet{rqFPBoundLoExeBOneCTwoSaferRustDen}{32}
\rqSet{rqFPBoundLoExeBOneCTwoSaferRustPct}{31.2}
\rqSet{rqFPBoundLoExeBOneSmartCTwoRustNum}{26}
\rqSet{rqFPBoundLoExeBOneSmartCTwoRustDen}{38}
\rqSet{rqFPBoundLoExeBOneSmartCTwoRustPct}{68.4}
\rqSet{rqFPBoundLoExeBOneKimiNum}{18}
\rqSet{rqFPBoundLoExeBOneKimiDen}{26}
\rqSet{rqFPBoundLoExeBOneKimiPct}{69.2}
\rqSet{rqFPBoundLoExeBOneGptFiveNum}{23}
\rqSet{rqFPBoundLoExeBOneGptFiveDen}{33}
\rqSet{rqFPBoundLoExeBOneGptFivePct}{69.7}
\rqSet{rqFPBoundLoExeBOneGeminiNum}{23}
\rqSet{rqFPBoundLoExeBOneGeminiDen}{36}
\rqSet{rqFPBoundLoExeBOneGeminiPct}{63.9}
\rqSet{rqFPBoundLoExeBTwoNum}{24}
\rqSet{rqFPBoundLoExeBTwoDen}{66}
\rqSet{rqFPBoundLoExeBTwoPct}{36.4}
\rqSet{rqFPBoundLoExeBTwoKiroNum}{3}
\rqSet{rqFPBoundLoExeBTwoKiroDen}{4}
\rqSet{rqFPBoundLoExeBTwoKiroPct}{75.0}
\rqSet{rqFPBoundLoExeBTwoClaudeNum}{6}
\rqSet{rqFPBoundLoExeBTwoClaudeDen}{9}
\rqSet{rqFPBoundLoExeBTwoClaudePct}{66.7}
\rqSet{rqFPBoundLoExeBTwoCodexNum}{3}
\rqSet{rqFPBoundLoExeBTwoCodexDen}{8}
\rqSet{rqFPBoundLoExeBTwoCodexPct}{37.5}
\rqSet{rqFPBoundLoExeBTwoCTwoRustNum}{2}
\rqSet{rqFPBoundLoExeBTwoCTwoRustDen}{10}
\rqSet{rqFPBoundLoExeBTwoCTwoRustPct}{20.0}
\rqSet{rqFPBoundLoExeBTwoLaertesNum}{2}
\rqSet{rqFPBoundLoExeBTwoLaertesDen}{12}
\rqSet{rqFPBoundLoExeBTwoLaertesPct}{16.7}
\rqSet{rqFPBoundLoExeBTwoCTwoSaferRustNum}{2}
\rqSet{rqFPBoundLoExeBTwoCTwoSaferRustDen}{9}
\rqSet{rqFPBoundLoExeBTwoCTwoSaferRustPct}{22.2}
\rqSet{rqFPBoundLoExeBTwoSmartCTwoRustNum}{1}
\rqSet{rqFPBoundLoExeBTwoSmartCTwoRustDen}{2}
\rqSet{rqFPBoundLoExeBTwoSmartCTwoRustPct}{50.0}
\rqSet{rqFPBoundLoExeBTwoKimiNum}{0}
\rqSet{rqFPBoundLoExeBTwoKimiDen}{2}
\rqSet{rqFPBoundLoExeBTwoKimiPct}{0.0}
\rqSet{rqFPBoundLoExeBTwoGptFiveNum}{2}
\rqSet{rqFPBoundLoExeBTwoGptFiveDen}{4}
\rqSet{rqFPBoundLoExeBTwoGptFivePct}{50.0}
\rqSet{rqFPBoundLoExeBTwoGeminiNum}{3}
\rqSet{rqFPBoundLoExeBTwoGeminiDen}{6}
\rqSet{rqFPBoundLoExeBTwoGeminiPct}{50.0}
\rqSet{rqFPBoundLoExePZeroNum}{0}
\rqSet{rqFPBoundLoExePZeroDen}{6}
\rqSet{rqFPBoundLoExePZeroPct}{0.0}
\rqSet{rqFPBoundLoExePZeroKiroNum}{0}
\rqSet{rqFPBoundLoExePZeroKiroDen}{1}
\rqSet{rqFPBoundLoExePZeroKiroPct}{0.0}
\rqSet{rqFPBoundLoExePZeroClaudeNum}{0}
\rqSet{rqFPBoundLoExePZeroClaudeDen}{1}
\rqSet{rqFPBoundLoExePZeroClaudePct}{0.0}
\rqSet{rqFPBoundLoExePZeroCodexNum}{0}
\rqSet{rqFPBoundLoExePZeroCodexDen}{1}
\rqSet{rqFPBoundLoExePZeroCodexPct}{0.0}
\rqSet{rqFPBoundLoExePZeroCTwoRustNum}{0}
\rqSet{rqFPBoundLoExePZeroCTwoRustDen}{1}
\rqSet{rqFPBoundLoExePZeroCTwoRustPct}{0.0}
\rqSet{rqFPBoundLoExePZeroLaertesNum}{0}
\rqSet{rqFPBoundLoExePZeroLaertesDen}{1}
\rqSet{rqFPBoundLoExePZeroLaertesPct}{0.0}
\rqSet{rqFPBoundLoExePZeroCTwoSaferRustNum}{0}
\rqSet{rqFPBoundLoExePZeroCTwoSaferRustDen}{1}
\rqSet{rqFPBoundLoExePZeroCTwoSaferRustPct}{0.0}
\rqSet{rqFPBoundLoExePZeroSmartCTwoRustNum}{0}
\rqSet{rqFPBoundLoExePZeroSmartCTwoRustDen}{0}
\rqSet{rqFPBoundLoExePZeroSmartCTwoRustPct}{n/a}
\rqSet{rqFPBoundLoExePZeroKimiNum}{0}
\rqSet{rqFPBoundLoExePZeroKimiDen}{0}
\rqSet{rqFPBoundLoExePZeroKimiPct}{n/a}
\rqSet{rqFPBoundLoExePZeroGptFiveNum}{0}
\rqSet{rqFPBoundLoExePZeroGptFiveDen}{0}
\rqSet{rqFPBoundLoExePZeroGptFivePct}{n/a}
\rqSet{rqFPBoundLoExePZeroGeminiNum}{0}
\rqSet{rqFPBoundLoExePZeroGeminiDen}{0}
\rqSet{rqFPBoundLoExePZeroGeminiPct}{n/a}
\rqSet{rqFPBoundLoExePOneNum}{0}
\rqSet{rqFPBoundLoExePOneDen}{108}
\rqSet{rqFPBoundLoExePOnePct}{0.0}
\rqSet{rqFPBoundLoExePOneKiroNum}{0}
\rqSet{rqFPBoundLoExePOneKiroDen}{48}
\rqSet{rqFPBoundLoExePOneKiroPct}{0.0}
\rqSet{rqFPBoundLoExePOneClaudeNum}{0}
\rqSet{rqFPBoundLoExePOneClaudeDen}{48}
\rqSet{rqFPBoundLoExePOneClaudePct}{0.0}
\rqSet{rqFPBoundLoExePOneCodexNum}{0}
\rqSet{rqFPBoundLoExePOneCodexDen}{12}
\rqSet{rqFPBoundLoExePOneCodexPct}{0.0}
\rqSet{rqFPBoundLoExePOneCTwoRustNum}{0}
\rqSet{rqFPBoundLoExePOneCTwoRustDen}{0}
\rqSet{rqFPBoundLoExePOneCTwoRustPct}{n/a}
\rqSet{rqFPBoundLoExePOneLaertesNum}{0}
\rqSet{rqFPBoundLoExePOneLaertesDen}{0}
\rqSet{rqFPBoundLoExePOneLaertesPct}{n/a}
\rqSet{rqFPBoundLoExePOneCTwoSaferRustNum}{0}
\rqSet{rqFPBoundLoExePOneCTwoSaferRustDen}{0}
\rqSet{rqFPBoundLoExePOneCTwoSaferRustPct}{n/a}
\rqSet{rqFPBoundLoExePOneSmartCTwoRustNum}{0}
\rqSet{rqFPBoundLoExePOneSmartCTwoRustDen}{0}
\rqSet{rqFPBoundLoExePOneSmartCTwoRustPct}{n/a}
\rqSet{rqFPBoundLoExePOneKimiNum}{0}
\rqSet{rqFPBoundLoExePOneKimiDen}{0}
\rqSet{rqFPBoundLoExePOneKimiPct}{n/a}
\rqSet{rqFPBoundLoExePOneGptFiveNum}{0}
\rqSet{rqFPBoundLoExePOneGptFiveDen}{0}
\rqSet{rqFPBoundLoExePOneGptFivePct}{n/a}
\rqSet{rqFPBoundLoExePOneGeminiNum}{0}
\rqSet{rqFPBoundLoExePOneGeminiDen}{0}
\rqSet{rqFPBoundLoExePOneGeminiPct}{n/a}
\rqSet{rqFPBoundHiExeAllNum}{259}
\rqSet{rqFPBoundHiExeAllDen}{559}
\rqSet{rqFPBoundHiExeAllPct}{46.3}
\rqSet{rqFPBoundHiExeAllKiroNum}{37}
\rqSet{rqFPBoundHiExeAllKiroDen}{96}
\rqSet{rqFPBoundHiExeAllKiroPct}{38.5}
\rqSet{rqFPBoundHiExeAllClaudeNum}{32}
\rqSet{rqFPBoundHiExeAllClaudeDen}{101}
\rqSet{rqFPBoundHiExeAllClaudePct}{31.7}
\rqSet{rqFPBoundHiExeAllCodexNum}{18}
\rqSet{rqFPBoundHiExeAllCodexDen}{63}
\rqSet{rqFPBoundHiExeAllCodexPct}{28.6}
\rqSet{rqFPBoundHiExeAllCTwoRustNum}{13}
\rqSet{rqFPBoundHiExeAllCTwoRustDen}{54}
\rqSet{rqFPBoundHiExeAllCTwoRustPct}{24.1}
\rqSet{rqFPBoundHiExeAllLaertesNum}{15}
\rqSet{rqFPBoundHiExeAllLaertesDen}{56}
\rqSet{rqFPBoundHiExeAllLaertesPct}{26.8}
\rqSet{rqFPBoundHiExeAllCTwoSaferRustNum}{22}
\rqSet{rqFPBoundHiExeAllCTwoSaferRustDen}{42}
\rqSet{rqFPBoundHiExeAllCTwoSaferRustPct}{52.4}
\rqSet{rqFPBoundHiExeAllSmartCTwoRustNum}{34}
\rqSet{rqFPBoundHiExeAllSmartCTwoRustDen}{40}
\rqSet{rqFPBoundHiExeAllSmartCTwoRustPct}{85.0}
\rqSet{rqFPBoundHiExeAllKimiNum}{22}
\rqSet{rqFPBoundHiExeAllKimiDen}{28}
\rqSet{rqFPBoundHiExeAllKimiPct}{78.6}
\rqSet{rqFPBoundHiExeAllGptFiveNum}{32}
\rqSet{rqFPBoundHiExeAllGptFiveDen}{37}
\rqSet{rqFPBoundHiExeAllGptFivePct}{86.5}
\rqSet{rqFPBoundHiExeAllGeminiNum}{34}
\rqSet{rqFPBoundHiExeAllGeminiDen}{42}
\rqSet{rqFPBoundHiExeAllGeminiPct}{81.0}
\rqSet{rqFPBoundHiExeBOneNum}{216}
\rqSet{rqFPBoundHiExeBOneDen}{379}
\rqSet{rqFPBoundHiExeBOnePct}{57.0}
\rqSet{rqFPBoundHiExeBOneKiroNum}{33}
\rqSet{rqFPBoundHiExeBOneKiroDen}{43}
\rqSet{rqFPBoundHiExeBOneKiroPct}{76.7}
\rqSet{rqFPBoundHiExeBOneClaudeNum}{25}
\rqSet{rqFPBoundHiExeBOneClaudeDen}{43}
\rqSet{rqFPBoundHiExeBOneClaudePct}{58.1}
\rqSet{rqFPBoundHiExeBOneCodexNum}{14}
\rqSet{rqFPBoundHiExeBOneCodexDen}{42}
\rqSet{rqFPBoundHiExeBOneCodexPct}{33.3}
\rqSet{rqFPBoundHiExeBOneCTwoRustNum}{8}
\rqSet{rqFPBoundHiExeBOneCTwoRustDen}{43}
\rqSet{rqFPBoundHiExeBOneCTwoRustPct}{18.6}
\rqSet{rqFPBoundHiExeBOneLaertesNum}{8}
\rqSet{rqFPBoundHiExeBOneLaertesDen}{43}
\rqSet{rqFPBoundHiExeBOneLaertesPct}{18.6}
\rqSet{rqFPBoundHiExeBOneCTwoSaferRustNum}{15}
\rqSet{rqFPBoundHiExeBOneCTwoSaferRustDen}{32}
\rqSet{rqFPBoundHiExeBOneCTwoSaferRustPct}{46.9}
\rqSet{rqFPBoundHiExeBOneSmartCTwoRustNum}{33}
\rqSet{rqFPBoundHiExeBOneSmartCTwoRustDen}{38}
\rqSet{rqFPBoundHiExeBOneSmartCTwoRustPct}{86.8}
\rqSet{rqFPBoundHiExeBOneKimiNum}{21}
\rqSet{rqFPBoundHiExeBOneKimiDen}{26}
\rqSet{rqFPBoundHiExeBOneKimiPct}{80.8}
\rqSet{rqFPBoundHiExeBOneGptFiveNum}{29}
\rqSet{rqFPBoundHiExeBOneGptFiveDen}{33}
\rqSet{rqFPBoundHiExeBOneGptFivePct}{87.9}
\rqSet{rqFPBoundHiExeBOneGeminiNum}{30}
\rqSet{rqFPBoundHiExeBOneGeminiDen}{36}
\rqSet{rqFPBoundHiExeBOneGeminiPct}{83.3}
\rqSet{rqFPBoundHiExeBTwoNum}{37}
\rqSet{rqFPBoundHiExeBTwoDen}{66}
\rqSet{rqFPBoundHiExeBTwoPct}{56.1}
\rqSet{rqFPBoundHiExeBTwoKiroNum}{3}
\rqSet{rqFPBoundHiExeBTwoKiroDen}{4}
\rqSet{rqFPBoundHiExeBTwoKiroPct}{75.0}
\rqSet{rqFPBoundHiExeBTwoClaudeNum}{6}
\rqSet{rqFPBoundHiExeBTwoClaudeDen}{9}
\rqSet{rqFPBoundHiExeBTwoClaudePct}{66.7}
\rqSet{rqFPBoundHiExeBTwoCodexNum}{3}
\rqSet{rqFPBoundHiExeBTwoCodexDen}{8}
\rqSet{rqFPBoundHiExeBTwoCodexPct}{37.5}
\rqSet{rqFPBoundHiExeBTwoCTwoRustNum}{4}
\rqSet{rqFPBoundHiExeBTwoCTwoRustDen}{10}
\rqSet{rqFPBoundHiExeBTwoCTwoRustPct}{40.0}
\rqSet{rqFPBoundHiExeBTwoLaertesNum}{6}
\rqSet{rqFPBoundHiExeBTwoLaertesDen}{12}
\rqSet{rqFPBoundHiExeBTwoLaertesPct}{50.0}
\rqSet{rqFPBoundHiExeBTwoCTwoSaferRustNum}{6}
\rqSet{rqFPBoundHiExeBTwoCTwoSaferRustDen}{9}
\rqSet{rqFPBoundHiExeBTwoCTwoSaferRustPct}{66.7}
\rqSet{rqFPBoundHiExeBTwoSmartCTwoRustNum}{1}
\rqSet{rqFPBoundHiExeBTwoSmartCTwoRustDen}{2}
\rqSet{rqFPBoundHiExeBTwoSmartCTwoRustPct}{50.0}
\rqSet{rqFPBoundHiExeBTwoKimiNum}{1}
\rqSet{rqFPBoundHiExeBTwoKimiDen}{2}
\rqSet{rqFPBoundHiExeBTwoKimiPct}{50.0}
\rqSet{rqFPBoundHiExeBTwoGptFiveNum}{3}
\rqSet{rqFPBoundHiExeBTwoGptFiveDen}{4}
\rqSet{rqFPBoundHiExeBTwoGptFivePct}{75.0}
\rqSet{rqFPBoundHiExeBTwoGeminiNum}{4}
\rqSet{rqFPBoundHiExeBTwoGeminiDen}{6}
\rqSet{rqFPBoundHiExeBTwoGeminiPct}{66.7}
\rqSet{rqFPBoundHiExePZeroNum}{6}
\rqSet{rqFPBoundHiExePZeroDen}{6}
\rqSet{rqFPBoundHiExePZeroPct}{100.0}
\rqSet{rqFPBoundHiExePZeroKiroNum}{1}
\rqSet{rqFPBoundHiExePZeroKiroDen}{1}
\rqSet{rqFPBoundHiExePZeroKiroPct}{100.0}
\rqSet{rqFPBoundHiExePZeroClaudeNum}{1}
\rqSet{rqFPBoundHiExePZeroClaudeDen}{1}
\rqSet{rqFPBoundHiExePZeroClaudePct}{100.0}
\rqSet{rqFPBoundHiExePZeroCodexNum}{1}
\rqSet{rqFPBoundHiExePZeroCodexDen}{1}
\rqSet{rqFPBoundHiExePZeroCodexPct}{100.0}
\rqSet{rqFPBoundHiExePZeroCTwoRustNum}{1}
\rqSet{rqFPBoundHiExePZeroCTwoRustDen}{1}
\rqSet{rqFPBoundHiExePZeroCTwoRustPct}{100.0}
\rqSet{rqFPBoundHiExePZeroLaertesNum}{1}
\rqSet{rqFPBoundHiExePZeroLaertesDen}{1}
\rqSet{rqFPBoundHiExePZeroLaertesPct}{100.0}
\rqSet{rqFPBoundHiExePZeroCTwoSaferRustNum}{1}
\rqSet{rqFPBoundHiExePZeroCTwoSaferRustDen}{1}
\rqSet{rqFPBoundHiExePZeroCTwoSaferRustPct}{100.0}
\rqSet{rqFPBoundHiExePZeroSmartCTwoRustNum}{0}
\rqSet{rqFPBoundHiExePZeroSmartCTwoRustDen}{0}
\rqSet{rqFPBoundHiExePZeroSmartCTwoRustPct}{n/a}
\rqSet{rqFPBoundHiExePZeroKimiNum}{0}
\rqSet{rqFPBoundHiExePZeroKimiDen}{0}
\rqSet{rqFPBoundHiExePZeroKimiPct}{n/a}
\rqSet{rqFPBoundHiExePZeroGptFiveNum}{0}
\rqSet{rqFPBoundHiExePZeroGptFiveDen}{0}
\rqSet{rqFPBoundHiExePZeroGptFivePct}{n/a}
\rqSet{rqFPBoundHiExePZeroGeminiNum}{0}
\rqSet{rqFPBoundHiExePZeroGeminiDen}{0}
\rqSet{rqFPBoundHiExePZeroGeminiPct}{n/a}
\rqSet{rqFPBoundHiExePOneNum}{0}
\rqSet{rqFPBoundHiExePOneDen}{108}
\rqSet{rqFPBoundHiExePOnePct}{0.0}
\rqSet{rqFPBoundHiExePOneKiroNum}{0}
\rqSet{rqFPBoundHiExePOneKiroDen}{48}
\rqSet{rqFPBoundHiExePOneKiroPct}{0.0}
\rqSet{rqFPBoundHiExePOneClaudeNum}{0}
\rqSet{rqFPBoundHiExePOneClaudeDen}{48}
\rqSet{rqFPBoundHiExePOneClaudePct}{0.0}
\rqSet{rqFPBoundHiExePOneCodexNum}{0}
\rqSet{rqFPBoundHiExePOneCodexDen}{12}
\rqSet{rqFPBoundHiExePOneCodexPct}{0.0}
\rqSet{rqFPBoundHiExePOneCTwoRustNum}{0}
\rqSet{rqFPBoundHiExePOneCTwoRustDen}{0}
\rqSet{rqFPBoundHiExePOneCTwoRustPct}{n/a}
\rqSet{rqFPBoundHiExePOneLaertesNum}{0}
\rqSet{rqFPBoundHiExePOneLaertesDen}{0}
\rqSet{rqFPBoundHiExePOneLaertesPct}{n/a}
\rqSet{rqFPBoundHiExePOneCTwoSaferRustNum}{0}
\rqSet{rqFPBoundHiExePOneCTwoSaferRustDen}{0}
\rqSet{rqFPBoundHiExePOneCTwoSaferRustPct}{n/a}
\rqSet{rqFPBoundHiExePOneSmartCTwoRustNum}{0}
\rqSet{rqFPBoundHiExePOneSmartCTwoRustDen}{0}
\rqSet{rqFPBoundHiExePOneSmartCTwoRustPct}{n/a}
\rqSet{rqFPBoundHiExePOneKimiNum}{0}
\rqSet{rqFPBoundHiExePOneKimiDen}{0}
\rqSet{rqFPBoundHiExePOneKimiPct}{n/a}
\rqSet{rqFPBoundHiExePOneGptFiveNum}{0}
\rqSet{rqFPBoundHiExePOneGptFiveDen}{0}
\rqSet{rqFPBoundHiExePOneGptFivePct}{n/a}
\rqSet{rqFPBoundHiExePOneGeminiNum}{0}
\rqSet{rqFPBoundHiExePOneGeminiDen}{0}
\rqSet{rqFPBoundHiExePOneGeminiPct}{n/a}
\rqSet{rqSeedDetOneExeAllNum}{0}
\rqSet{rqSeedDetOneExeAllDen}{176}
\rqSet{rqSeedDetOneExeAllPct}{0.0}
\rqSet{rqSeedDetOneExeAllKiroNum}{0}
\rqSet{rqSeedDetOneExeAllKiroDen}{29}
\rqSet{rqSeedDetOneExeAllKiroPct}{0.0}
\rqSet{rqSeedDetOneExeAllClaudeNum}{0}
\rqSet{rqSeedDetOneExeAllClaudeDen}{24}
\rqSet{rqSeedDetOneExeAllClaudePct}{0.0}
\rqSet{rqSeedDetOneExeAllCodexNum}{0}
\rqSet{rqSeedDetOneExeAllCodexDen}{9}
\rqSet{rqSeedDetOneExeAllCodexPct}{0.0}
\rqSet{rqSeedDetOneExeAllCTwoRustNum}{0}
\rqSet{rqSeedDetOneExeAllCTwoRustDen}{3}
\rqSet{rqSeedDetOneExeAllCTwoRustPct}{0.0}
\rqSet{rqSeedDetOneExeAllLaertesNum}{0}
\rqSet{rqSeedDetOneExeAllLaertesDen}{3}
\rqSet{rqSeedDetOneExeAllLaertesPct}{0.0}
\rqSet{rqSeedDetOneExeAllCTwoSaferRustNum}{0}
\rqSet{rqSeedDetOneExeAllCTwoSaferRustDen}{12}
\rqSet{rqSeedDetOneExeAllCTwoSaferRustPct}{0.0}
\rqSet{rqSeedDetOneExeAllSmartCTwoRustNum}{0}
\rqSet{rqSeedDetOneExeAllSmartCTwoRustDen}{27}
\rqSet{rqSeedDetOneExeAllSmartCTwoRustPct}{0.0}
\rqSet{rqSeedDetOneExeAllKimiNum}{0}
\rqSet{rqSeedDetOneExeAllKimiDen}{18}
\rqSet{rqSeedDetOneExeAllKimiPct}{0.0}
\rqSet{rqSeedDetOneExeAllGptFiveNum}{0}
\rqSet{rqSeedDetOneExeAllGptFiveDen}{25}
\rqSet{rqSeedDetOneExeAllGptFivePct}{0.0}
\rqSet{rqSeedDetOneExeAllGeminiNum}{0}
\rqSet{rqSeedDetOneExeAllGeminiDen}{26}
\rqSet{rqSeedDetOneExeAllGeminiPct}{0.0}
\rqSet{rqSeedDetOneExeBOneNum}{0}
\rqSet{rqSeedDetOneExeBOneDen}{152}
\rqSet{rqSeedDetOneExeBOnePct}{0.0}
\rqSet{rqSeedDetOneExeBTwoNum}{0}
\rqSet{rqSeedDetOneExeBTwoDen}{24}
\rqSet{rqSeedDetOneExeBTwoPct}{0.0}
\rqSet{rqSeedDetOneExePZeroNum}{0}
\rqSet{rqSeedDetOneExePZeroDen}{0}
\rqSet{rqSeedDetOneExePZeroPct}{n/a}
\rqSet{rqSeedDetOneExePOneNum}{0}
\rqSet{rqSeedDetOneExePOneDen}{0}
\rqSet{rqSeedDetOneExePOnePct}{n/a}
\rqSet{rqSeedDetTwoExeAllNum}{4}
\rqSet{rqSeedDetTwoExeAllDen}{176}
\rqSet{rqSeedDetTwoExeAllPct}{2.3}
\rqSet{rqSeedDetTwoExeAllKiroNum}{1}
\rqSet{rqSeedDetTwoExeAllKiroDen}{29}
\rqSet{rqSeedDetTwoExeAllKiroPct}{3.4}
\rqSet{rqSeedDetTwoExeAllClaudeNum}{0}
\rqSet{rqSeedDetTwoExeAllClaudeDen}{24}
\rqSet{rqSeedDetTwoExeAllClaudePct}{0.0}
\rqSet{rqSeedDetTwoExeAllCodexNum}{0}
\rqSet{rqSeedDetTwoExeAllCodexDen}{9}
\rqSet{rqSeedDetTwoExeAllCodexPct}{0.0}
\rqSet{rqSeedDetTwoExeAllCTwoRustNum}{0}
\rqSet{rqSeedDetTwoExeAllCTwoRustDen}{3}
\rqSet{rqSeedDetTwoExeAllCTwoRustPct}{0.0}
\rqSet{rqSeedDetTwoExeAllLaertesNum}{0}
\rqSet{rqSeedDetTwoExeAllLaertesDen}{3}
\rqSet{rqSeedDetTwoExeAllLaertesPct}{0.0}
\rqSet{rqSeedDetTwoExeAllCTwoSaferRustNum}{1}
\rqSet{rqSeedDetTwoExeAllCTwoSaferRustDen}{12}
\rqSet{rqSeedDetTwoExeAllCTwoSaferRustPct}{8.3}
\rqSet{rqSeedDetTwoExeAllSmartCTwoRustNum}{0}
\rqSet{rqSeedDetTwoExeAllSmartCTwoRustDen}{27}
\rqSet{rqSeedDetTwoExeAllSmartCTwoRustPct}{0.0}
\rqSet{rqSeedDetTwoExeAllKimiNum}{1}
\rqSet{rqSeedDetTwoExeAllKimiDen}{18}
\rqSet{rqSeedDetTwoExeAllKimiPct}{5.6}
\rqSet{rqSeedDetTwoExeAllGptFiveNum}{1}
\rqSet{rqSeedDetTwoExeAllGptFiveDen}{25}
\rqSet{rqSeedDetTwoExeAllGptFivePct}{4.0}
\rqSet{rqSeedDetTwoExeAllGeminiNum}{0}
\rqSet{rqSeedDetTwoExeAllGeminiDen}{26}
\rqSet{rqSeedDetTwoExeAllGeminiPct}{0.0}
\rqSet{rqSeedDetTwoExeBOneNum}{4}
\rqSet{rqSeedDetTwoExeBOneDen}{152}
\rqSet{rqSeedDetTwoExeBOnePct}{2.6}
\rqSet{rqSeedDetTwoExeBTwoNum}{0}
\rqSet{rqSeedDetTwoExeBTwoDen}{24}
\rqSet{rqSeedDetTwoExeBTwoPct}{0.0}
\rqSet{rqSeedDetTwoExePZeroNum}{0}
\rqSet{rqSeedDetTwoExePZeroDen}{0}
\rqSet{rqSeedDetTwoExePZeroPct}{n/a}
\rqSet{rqSeedDetTwoExePOneNum}{0}
\rqSet{rqSeedDetTwoExePOneDen}{0}
\rqSet{rqSeedDetTwoExePOnePct}{n/a}
\rqSet{rqSeedDetThreeExeAllNum}{172}
\rqSet{rqSeedDetThreeExeAllDen}{176}
\rqSet{rqSeedDetThreeExeAllPct}{97.7}
\rqSet{rqSeedDetThreeExeAllKiroNum}{28}
\rqSet{rqSeedDetThreeExeAllKiroDen}{29}
\rqSet{rqSeedDetThreeExeAllKiroPct}{96.6}
\rqSet{rqSeedDetThreeExeAllClaudeNum}{24}
\rqSet{rqSeedDetThreeExeAllClaudeDen}{24}
\rqSet{rqSeedDetThreeExeAllClaudePct}{100.0}
\rqSet{rqSeedDetThreeExeAllCodexNum}{9}
\rqSet{rqSeedDetThreeExeAllCodexDen}{9}
\rqSet{rqSeedDetThreeExeAllCodexPct}{100.0}
\rqSet{rqSeedDetThreeExeAllCTwoRustNum}{3}
\rqSet{rqSeedDetThreeExeAllCTwoRustDen}{3}
\rqSet{rqSeedDetThreeExeAllCTwoRustPct}{100.0}
\rqSet{rqSeedDetThreeExeAllLaertesNum}{3}
\rqSet{rqSeedDetThreeExeAllLaertesDen}{3}
\rqSet{rqSeedDetThreeExeAllLaertesPct}{100.0}
\rqSet{rqSeedDetThreeExeAllCTwoSaferRustNum}{11}
\rqSet{rqSeedDetThreeExeAllCTwoSaferRustDen}{12}
\rqSet{rqSeedDetThreeExeAllCTwoSaferRustPct}{91.7}
\rqSet{rqSeedDetThreeExeAllSmartCTwoRustNum}{27}
\rqSet{rqSeedDetThreeExeAllSmartCTwoRustDen}{27}
\rqSet{rqSeedDetThreeExeAllSmartCTwoRustPct}{100.0}
\rqSet{rqSeedDetThreeExeAllKimiNum}{17}
\rqSet{rqSeedDetThreeExeAllKimiDen}{18}
\rqSet{rqSeedDetThreeExeAllKimiPct}{94.4}
\rqSet{rqSeedDetThreeExeAllGptFiveNum}{24}
\rqSet{rqSeedDetThreeExeAllGptFiveDen}{25}
\rqSet{rqSeedDetThreeExeAllGptFivePct}{96.0}
\rqSet{rqSeedDetThreeExeAllGeminiNum}{26}
\rqSet{rqSeedDetThreeExeAllGeminiDen}{26}
\rqSet{rqSeedDetThreeExeAllGeminiPct}{100.0}
\rqSet{rqSeedDetThreeExeBOneNum}{148}
\rqSet{rqSeedDetThreeExeBOneDen}{152}
\rqSet{rqSeedDetThreeExeBOnePct}{97.4}
\rqSet{rqSeedDetThreeExeBTwoNum}{24}
\rqSet{rqSeedDetThreeExeBTwoDen}{24}
\rqSet{rqSeedDetThreeExeBTwoPct}{100.0}
\rqSet{rqSeedDetThreeExePZeroNum}{0}
\rqSet{rqSeedDetThreeExePZeroDen}{0}
\rqSet{rqSeedDetThreeExePZeroPct}{n/a}
\rqSet{rqSeedDetThreeExePOneNum}{0}
\rqSet{rqSeedDetThreeExePOneDen}{0}
\rqSet{rqSeedDetThreeExePOnePct}{n/a}
\rqSet{rqFPSeedZeroExeAllNum}{188}
\rqSet{rqFPSeedZeroExeAllDen}{504}
\rqSet{rqFPSeedZeroExeAllPct}{37.3}
\rqSet{rqFPSeedZeroExeAllKiroNum}{30}
\rqSet{rqFPSeedZeroExeAllKiroDen}{91}
\rqSet{rqFPSeedZeroExeAllKiroPct}{33.0}
\rqSet{rqFPSeedZeroExeAllClaudeNum}{26}
\rqSet{rqFPSeedZeroExeAllClaudeDen}{96}
\rqSet{rqFPSeedZeroExeAllClaudePct}{27.1}
\rqSet{rqFPSeedZeroExeAllCodexNum}{9}
\rqSet{rqFPSeedZeroExeAllCodexDen}{57}
\rqSet{rqFPSeedZeroExeAllCodexPct}{15.8}
\rqSet{rqFPSeedZeroExeAllCTwoRustNum}{3}
\rqSet{rqFPSeedZeroExeAllCTwoRustDen}{48}
\rqSet{rqFPSeedZeroExeAllCTwoRustPct}{6.2}
\rqSet{rqFPSeedZeroExeAllLaertesNum}{3}
\rqSet{rqFPSeedZeroExeAllLaertesDen}{48}
\rqSet{rqFPSeedZeroExeAllLaertesPct}{6.2}
\rqSet{rqFPSeedZeroExeAllCTwoSaferRustNum}{14}
\rqSet{rqFPSeedZeroExeAllCTwoSaferRustDen}{36}
\rqSet{rqFPSeedZeroExeAllCTwoSaferRustPct}{38.9}
\rqSet{rqFPSeedZeroExeAllSmartCTwoRustNum}{29}
\rqSet{rqFPSeedZeroExeAllSmartCTwoRustDen}{35}
\rqSet{rqFPSeedZeroExeAllSmartCTwoRustPct}{82.9}
\rqSet{rqFPSeedZeroExeAllKimiNum}{19}
\rqSet{rqFPSeedZeroExeAllKimiDen}{25}
\rqSet{rqFPSeedZeroExeAllKimiPct}{76.0}
\rqSet{rqFPSeedZeroExeAllGptFiveNum}{27}
\rqSet{rqFPSeedZeroExeAllGptFiveDen}{32}
\rqSet{rqFPSeedZeroExeAllGptFivePct}{84.4}
\rqSet{rqFPSeedZeroExeAllGeminiNum}{28}
\rqSet{rqFPSeedZeroExeAllGeminiDen}{36}
\rqSet{rqFPSeedZeroExeAllGeminiPct}{77.8}
\rqSet{rqFPSeedZeroExeBOneNum}{162}
\rqSet{rqFPSeedZeroExeBOneDen}{334}
\rqSet{rqFPSeedZeroExeBOnePct}{48.5}
\rqSet{rqFPSeedZeroExeBTwoNum}{24}
\rqSet{rqFPSeedZeroExeBTwoDen}{56}
\rqSet{rqFPSeedZeroExeBTwoPct}{42.9}
\rqSet{rqFPSeedZeroExePZeroNum}{2}
\rqSet{rqFPSeedZeroExePZeroDen}{6}
\rqSet{rqFPSeedZeroExePZeroPct}{33.3}
\rqSet{rqFPSeedZeroExePOneNum}{0}
\rqSet{rqFPSeedZeroExePOneDen}{108}
\rqSet{rqFPSeedZeroExePOnePct}{0.0}
\rqSet{rqSeedMissingZeroExeAllNum}{0}
\rqSet{rqSeedMissingZeroExeAllDen}{548}
\rqSet{rqSeedMissingZeroExeAllPct}{0.0}
\rqSet{rqSeedMissingZeroExeBOneNum}{0}
\rqSet{rqSeedMissingZeroExeBOneDen}{378}
\rqSet{rqSeedMissingZeroExeBOnePct}{0.0}
\rqSet{rqSeedMissingZeroExeBTwoNum}{0}
\rqSet{rqSeedMissingZeroExeBTwoDen}{56}
\rqSet{rqSeedMissingZeroExeBTwoPct}{0.0}
\rqSet{rqSeedMissingZeroExePZeroNum}{0}
\rqSet{rqSeedMissingZeroExePZeroDen}{6}
\rqSet{rqSeedMissingZeroExePZeroPct}{0.0}
\rqSet{rqSeedMissingZeroExePOneNum}{0}
\rqSet{rqSeedMissingZeroExePOneDen}{108}
\rqSet{rqSeedMissingZeroExePOnePct}{0.0}
\rqSet{rqFPSeedOneExeAllNum}{182}
\rqSet{rqFPSeedOneExeAllDen}{363}
\rqSet{rqFPSeedOneExeAllPct}{50.1}
\rqSet{rqFPSeedOneExeAllKiroNum}{28}
\rqSet{rqFPSeedOneExeAllKiroDen}{40}
\rqSet{rqFPSeedOneExeAllKiroPct}{70.0}
\rqSet{rqFPSeedOneExeAllClaudeNum}{25}
\rqSet{rqFPSeedOneExeAllClaudeDen}{44}
\rqSet{rqFPSeedOneExeAllClaudePct}{56.8}
\rqSet{rqFPSeedOneExeAllCodexNum}{9}
\rqSet{rqFPSeedOneExeAllCodexDen}{41}
\rqSet{rqFPSeedOneExeAllCodexPct}{22.0}
\rqSet{rqFPSeedOneExeAllCTwoRustNum}{3}
\rqSet{rqFPSeedOneExeAllCTwoRustDen}{43}
\rqSet{rqFPSeedOneExeAllCTwoRustPct}{7.0}
\rqSet{rqFPSeedOneExeAllLaertesNum}{3}
\rqSet{rqFPSeedOneExeAllLaertesDen}{43}
\rqSet{rqFPSeedOneExeAllLaertesPct}{7.0}
\rqSet{rqFPSeedOneExeAllCTwoSaferRustNum}{13}
\rqSet{rqFPSeedOneExeAllCTwoSaferRustDen}{32}
\rqSet{rqFPSeedOneExeAllCTwoSaferRustPct}{40.6}
\rqSet{rqFPSeedOneExeAllSmartCTwoRustNum}{29}
\rqSet{rqFPSeedOneExeAllSmartCTwoRustDen}{34}
\rqSet{rqFPSeedOneExeAllSmartCTwoRustPct}{85.3}
\rqSet{rqFPSeedOneExeAllKimiNum}{18}
\rqSet{rqFPSeedOneExeAllKimiDen}{23}
\rqSet{rqFPSeedOneExeAllKimiPct}{78.3}
\rqSet{rqFPSeedOneExeAllGptFiveNum}{26}
\rqSet{rqFPSeedOneExeAllGptFiveDen}{30}
\rqSet{rqFPSeedOneExeAllGptFivePct}{86.7}
\rqSet{rqFPSeedOneExeAllGeminiNum}{28}
\rqSet{rqFPSeedOneExeAllGeminiDen}{33}
\rqSet{rqFPSeedOneExeAllGeminiPct}{84.8}
\rqSet{rqFPSeedOneExeBOneNum}{158}
\rqSet{rqFPSeedOneExeBOneDen}{315}
\rqSet{rqFPSeedOneExeBOnePct}{50.2}
\rqSet{rqFPSeedOneExeBTwoNum}{24}
\rqSet{rqFPSeedOneExeBTwoDen}{48}
\rqSet{rqFPSeedOneExeBTwoPct}{50.0}
\rqSet{rqFPSeedOneExePZeroNum}{0}
\rqSet{rqFPSeedOneExePZeroDen}{0}
\rqSet{rqFPSeedOneExePZeroPct}{n/a}
\rqSet{rqFPSeedOneExePOneNum}{0}
\rqSet{rqFPSeedOneExePOneDen}{0}
\rqSet{rqFPSeedOneExePOnePct}{n/a}
\rqSet{rqSeedMissingOneExeAllNum}{141}
\rqSet{rqSeedMissingOneExeAllDen}{548}
\rqSet{rqSeedMissingOneExeAllPct}{25.7}
\rqSet{rqSeedMissingOneExeBOneNum}{19}
\rqSet{rqSeedMissingOneExeBOneDen}{378}
\rqSet{rqSeedMissingOneExeBOnePct}{5.0}
\rqSet{rqSeedMissingOneExeBTwoNum}{8}
\rqSet{rqSeedMissingOneExeBTwoDen}{56}
\rqSet{rqSeedMissingOneExeBTwoPct}{14.3}
\rqSet{rqSeedMissingOneExePZeroNum}{6}
\rqSet{rqSeedMissingOneExePZeroDen}{6}
\rqSet{rqSeedMissingOneExePZeroPct}{100.0}
\rqSet{rqSeedMissingOneExePOneNum}{108}
\rqSet{rqSeedMissingOneExePOneDen}{108}
\rqSet{rqSeedMissingOneExePOnePct}{100.0}
\rqSet{rqFPSeedTwoExeAllNum}{176}
\rqSet{rqFPSeedTwoExeAllDen}{344}
\rqSet{rqFPSeedTwoExeAllPct}{51.2}
\rqSet{rqFPSeedTwoExeAllKiroNum}{29}
\rqSet{rqFPSeedTwoExeAllKiroDen}{38}
\rqSet{rqFPSeedTwoExeAllKiroPct}{76.3}
\rqSet{rqFPSeedTwoExeAllClaudeNum}{24}
\rqSet{rqFPSeedTwoExeAllClaudeDen}{42}
\rqSet{rqFPSeedTwoExeAllClaudePct}{57.1}
\rqSet{rqFPSeedTwoExeAllCodexNum}{9}
\rqSet{rqFPSeedTwoExeAllCodexDen}{39}
\rqSet{rqFPSeedTwoExeAllCodexPct}{23.1}
\rqSet{rqFPSeedTwoExeAllCTwoRustNum}{3}
\rqSet{rqFPSeedTwoExeAllCTwoRustDen}{41}
\rqSet{rqFPSeedTwoExeAllCTwoRustPct}{7.3}
\rqSet{rqFPSeedTwoExeAllLaertesNum}{3}
\rqSet{rqFPSeedTwoExeAllLaertesDen}{41}
\rqSet{rqFPSeedTwoExeAllLaertesPct}{7.3}
\rqSet{rqFPSeedTwoExeAllCTwoSaferRustNum}{12}
\rqSet{rqFPSeedTwoExeAllCTwoSaferRustDen}{30}
\rqSet{rqFPSeedTwoExeAllCTwoSaferRustPct}{40.0}
\rqSet{rqFPSeedTwoExeAllSmartCTwoRustNum}{27}
\rqSet{rqFPSeedTwoExeAllSmartCTwoRustDen}{32}
\rqSet{rqFPSeedTwoExeAllSmartCTwoRustPct}{84.4}
\rqSet{rqFPSeedTwoExeAllKimiNum}{18}
\rqSet{rqFPSeedTwoExeAllKimiDen}{22}
\rqSet{rqFPSeedTwoExeAllKimiPct}{81.8}
\rqSet{rqFPSeedTwoExeAllGptFiveNum}{25}
\rqSet{rqFPSeedTwoExeAllGptFiveDen}{28}
\rqSet{rqFPSeedTwoExeAllGptFivePct}{89.3}
\rqSet{rqFPSeedTwoExeAllGeminiNum}{26}
\rqSet{rqFPSeedTwoExeAllGeminiDen}{31}
\rqSet{rqFPSeedTwoExeAllGeminiPct}{83.9}
\rqSet{rqFPSeedTwoExeBOneNum}{152}
\rqSet{rqFPSeedTwoExeBOneDen}{296}
\rqSet{rqFPSeedTwoExeBOnePct}{51.4}
\rqSet{rqFPSeedTwoExeBTwoNum}{24}
\rqSet{rqFPSeedTwoExeBTwoDen}{48}
\rqSet{rqFPSeedTwoExeBTwoPct}{50.0}
\rqSet{rqFPSeedTwoExePZeroNum}{0}
\rqSet{rqFPSeedTwoExePZeroDen}{0}
\rqSet{rqFPSeedTwoExePZeroPct}{n/a}
\rqSet{rqFPSeedTwoExePOneNum}{0}
\rqSet{rqFPSeedTwoExePOneDen}{0}
\rqSet{rqFPSeedTwoExePOnePct}{n/a}
\rqSet{rqSeedMissingTwoExeAllNum}{160}
\rqSet{rqSeedMissingTwoExeAllDen}{548}
\rqSet{rqSeedMissingTwoExeAllPct}{29.2}
\rqSet{rqSeedMissingTwoExeBOneNum}{38}
\rqSet{rqSeedMissingTwoExeBOneDen}{378}
\rqSet{rqSeedMissingTwoExeBOnePct}{10.1}
\rqSet{rqSeedMissingTwoExeBTwoNum}{8}
\rqSet{rqSeedMissingTwoExeBTwoDen}{56}
\rqSet{rqSeedMissingTwoExeBTwoPct}{14.3}
\rqSet{rqSeedMissingTwoExePZeroNum}{6}
\rqSet{rqSeedMissingTwoExePZeroDen}{6}
\rqSet{rqSeedMissingTwoExePZeroPct}{100.0}
\rqSet{rqSeedMissingTwoExePOneNum}{108}
\rqSet{rqSeedMissingTwoExePOneDen}{108}
\rqSet{rqSeedMissingTwoExePOnePct}{100.0}
\rqSet{rqUnstableOnlyExeAllNum}{15}
\rqSet{rqUnstableOnlyExeAllDen}{476}
\rqSet{rqUnstableOnlyExeAllPct}{3.2}
\rqSet{rqUnstableOnlyExeAllKiroNum}{0}
\rqSet{rqUnstableOnlyExeAllKiroDen}{88}
\rqSet{rqUnstableOnlyExeAllKiroPct}{0.0}
\rqSet{rqUnstableOnlyExeAllClaudeNum}{2}
\rqSet{rqUnstableOnlyExeAllClaudeDen}{93}
\rqSet{rqUnstableOnlyExeAllClaudePct}{2.2}
\rqSet{rqUnstableOnlyExeAllCodexNum}{2}
\rqSet{rqUnstableOnlyExeAllCodexDen}{54}
\rqSet{rqUnstableOnlyExeAllCodexPct}{3.7}
\rqSet{rqUnstableOnlyExeAllCTwoRustNum}{4}
\rqSet{rqUnstableOnlyExeAllCTwoRustDen}{44}
\rqSet{rqUnstableOnlyExeAllCTwoRustPct}{9.1}
\rqSet{rqUnstableOnlyExeAllLaertesNum}{4}
\rqSet{rqUnstableOnlyExeAllLaertesDen}{44}
\rqSet{rqUnstableOnlyExeAllLaertesPct}{9.1}
\rqSet{rqUnstableOnlyExeAllCTwoSaferRustNum}{1}
\rqSet{rqUnstableOnlyExeAllCTwoSaferRustDen}{32}
\rqSet{rqUnstableOnlyExeAllCTwoSaferRustPct}{3.1}
\rqSet{rqUnstableOnlyExeAllSmartCTwoRustNum}{1}
\rqSet{rqUnstableOnlyExeAllSmartCTwoRustDen}{33}
\rqSet{rqUnstableOnlyExeAllSmartCTwoRustPct}{3.0}
\rqSet{rqUnstableOnlyExeAllKimiNum}{0}
\rqSet{rqUnstableOnlyExeAllKimiDen}{24}
\rqSet{rqUnstableOnlyExeAllKimiPct}{0.0}
\rqSet{rqUnstableOnlyExeAllGptFiveNum}{0}
\rqSet{rqUnstableOnlyExeAllGptFiveDen}{30}
\rqSet{rqUnstableOnlyExeAllGptFivePct}{0.0}
\rqSet{rqUnstableOnlyExeAllGeminiNum}{1}
\rqSet{rqUnstableOnlyExeAllGeminiDen}{34}
\rqSet{rqUnstableOnlyExeAllGeminiPct}{2.9}
\rqSet{rqUnstableOnlyExeBOneNum}{12}
\rqSet{rqUnstableOnlyExeBOneDen}{315}
\rqSet{rqUnstableOnlyExeBOnePct}{3.8}
\rqSet{rqUnstableOnlyExeBTwoNum}{3}
\rqSet{rqUnstableOnlyExeBTwoDen}{53}
\rqSet{rqUnstableOnlyExeBTwoPct}{5.7}
\rqSet{rqUnstableOnlyExePZeroNum}{0}
\rqSet{rqUnstableOnlyExePZeroDen}{0}
\rqSet{rqUnstableOnlyExePZeroPct}{n/a}
\rqSet{rqUnstableOnlyExePOneNum}{0}
\rqSet{rqUnstableOnlyExePOneDen}{108}
\rqSet{rqUnstableOnlyExePOnePct}{0.0}
\rqSet{rqInconclOnlyExeAllNum}{0}
\rqSet{rqInconclOnlyExeAllDen}{476}
\rqSet{rqInconclOnlyExeAllPct}{0.0}
\rqSet{rqInconclOnlyExeAllKiroNum}{0}
\rqSet{rqInconclOnlyExeAllKiroDen}{88}
\rqSet{rqInconclOnlyExeAllKiroPct}{0.0}
\rqSet{rqInconclOnlyExeAllClaudeNum}{0}
\rqSet{rqInconclOnlyExeAllClaudeDen}{93}
\rqSet{rqInconclOnlyExeAllClaudePct}{0.0}
\rqSet{rqInconclOnlyExeAllCodexNum}{0}
\rqSet{rqInconclOnlyExeAllCodexDen}{54}
\rqSet{rqInconclOnlyExeAllCodexPct}{0.0}
\rqSet{rqInconclOnlyExeAllCTwoRustNum}{0}
\rqSet{rqInconclOnlyExeAllCTwoRustDen}{44}
\rqSet{rqInconclOnlyExeAllCTwoRustPct}{0.0}
\rqSet{rqInconclOnlyExeAllLaertesNum}{0}
\rqSet{rqInconclOnlyExeAllLaertesDen}{44}
\rqSet{rqInconclOnlyExeAllLaertesPct}{0.0}
\rqSet{rqInconclOnlyExeAllCTwoSaferRustNum}{0}
\rqSet{rqInconclOnlyExeAllCTwoSaferRustDen}{32}
\rqSet{rqInconclOnlyExeAllCTwoSaferRustPct}{0.0}
\rqSet{rqInconclOnlyExeAllSmartCTwoRustNum}{0}
\rqSet{rqInconclOnlyExeAllSmartCTwoRustDen}{33}
\rqSet{rqInconclOnlyExeAllSmartCTwoRustPct}{0.0}
\rqSet{rqInconclOnlyExeAllKimiNum}{0}
\rqSet{rqInconclOnlyExeAllKimiDen}{24}
\rqSet{rqInconclOnlyExeAllKimiPct}{0.0}
\rqSet{rqInconclOnlyExeAllGptFiveNum}{0}
\rqSet{rqInconclOnlyExeAllGptFiveDen}{30}
\rqSet{rqInconclOnlyExeAllGptFivePct}{0.0}
\rqSet{rqInconclOnlyExeAllGeminiNum}{0}
\rqSet{rqInconclOnlyExeAllGeminiDen}{34}
\rqSet{rqInconclOnlyExeAllGeminiPct}{0.0}
\rqSet{rqInconclOnlyExeBOneNum}{0}
\rqSet{rqInconclOnlyExeBOneDen}{315}
\rqSet{rqInconclOnlyExeBOnePct}{0.0}
\rqSet{rqInconclOnlyExeBTwoNum}{0}
\rqSet{rqInconclOnlyExeBTwoDen}{53}
\rqSet{rqInconclOnlyExeBTwoPct}{0.0}
\rqSet{rqInconclOnlyExePZeroNum}{0}
\rqSet{rqInconclOnlyExePZeroDen}{0}
\rqSet{rqInconclOnlyExePZeroPct}{n/a}
\rqSet{rqInconclOnlyExePOneNum}{0}
\rqSet{rqInconclOnlyExePOneDen}{108}
\rqSet{rqInconclOnlyExePOnePct}{0.0}
\rqSet{rqInconclUnstableExeAllNum}{0}
\rqSet{rqInconclUnstableExeAllDen}{476}
\rqSet{rqInconclUnstableExeAllPct}{0.0}
\rqSet{rqInconclUnstableExeAllKiroNum}{0}
\rqSet{rqInconclUnstableExeAllKiroDen}{88}
\rqSet{rqInconclUnstableExeAllKiroPct}{0.0}
\rqSet{rqInconclUnstableExeAllClaudeNum}{0}
\rqSet{rqInconclUnstableExeAllClaudeDen}{93}
\rqSet{rqInconclUnstableExeAllClaudePct}{0.0}
\rqSet{rqInconclUnstableExeAllCodexNum}{0}
\rqSet{rqInconclUnstableExeAllCodexDen}{54}
\rqSet{rqInconclUnstableExeAllCodexPct}{0.0}
\rqSet{rqInconclUnstableExeAllCTwoRustNum}{0}
\rqSet{rqInconclUnstableExeAllCTwoRustDen}{44}
\rqSet{rqInconclUnstableExeAllCTwoRustPct}{0.0}
\rqSet{rqInconclUnstableExeAllLaertesNum}{0}
\rqSet{rqInconclUnstableExeAllLaertesDen}{44}
\rqSet{rqInconclUnstableExeAllLaertesPct}{0.0}
\rqSet{rqInconclUnstableExeAllCTwoSaferRustNum}{0}
\rqSet{rqInconclUnstableExeAllCTwoSaferRustDen}{32}
\rqSet{rqInconclUnstableExeAllCTwoSaferRustPct}{0.0}
\rqSet{rqInconclUnstableExeAllSmartCTwoRustNum}{0}
\rqSet{rqInconclUnstableExeAllSmartCTwoRustDen}{33}
\rqSet{rqInconclUnstableExeAllSmartCTwoRustPct}{0.0}
\rqSet{rqInconclUnstableExeAllKimiNum}{0}
\rqSet{rqInconclUnstableExeAllKimiDen}{24}
\rqSet{rqInconclUnstableExeAllKimiPct}{0.0}
\rqSet{rqInconclUnstableExeAllGptFiveNum}{0}
\rqSet{rqInconclUnstableExeAllGptFiveDen}{30}
\rqSet{rqInconclUnstableExeAllGptFivePct}{0.0}
\rqSet{rqInconclUnstableExeAllGeminiNum}{0}
\rqSet{rqInconclUnstableExeAllGeminiDen}{34}
\rqSet{rqInconclUnstableExeAllGeminiPct}{0.0}
\rqSet{rqInconclUnstableExeBOneNum}{0}
\rqSet{rqInconclUnstableExeBOneDen}{315}
\rqSet{rqInconclUnstableExeBOnePct}{0.0}
\rqSet{rqInconclUnstableExeBTwoNum}{0}
\rqSet{rqInconclUnstableExeBTwoDen}{53}
\rqSet{rqInconclUnstableExeBTwoPct}{0.0}
\rqSet{rqInconclUnstableExePZeroNum}{0}
\rqSet{rqInconclUnstableExePZeroDen}{0}
\rqSet{rqInconclUnstableExePZeroPct}{n/a}
\rqSet{rqInconclUnstableExePOneNum}{0}
\rqSet{rqInconclUnstableExePOneDen}{108}
\rqSet{rqInconclUnstableExePOnePct}{0.0}
\rqSet{rqInconclAnyExeAllNum}{0}
\rqSet{rqInconclAnyExeAllDen}{476}
\rqSet{rqInconclAnyExeAllPct}{0.0}
\rqSet{rqInconclAnyExeBOneNum}{0}
\rqSet{rqInconclAnyExeBOneDen}{315}
\rqSet{rqInconclAnyExeBOnePct}{0.0}
\rqSet{rqInconclAnyExeBTwoNum}{0}
\rqSet{rqInconclAnyExeBTwoDen}{53}
\rqSet{rqInconclAnyExeBTwoPct}{0.0}
\rqSet{rqInconclAnyExePZeroNum}{0}
\rqSet{rqInconclAnyExePZeroDen}{0}
\rqSet{rqInconclAnyExePZeroPct}{n/a}
\rqSet{rqInconclAnyExePOneNum}{0}
\rqSet{rqInconclAnyExePOneDen}{108}
\rqSet{rqInconclAnyExePOnePct}{0.0}
\rqSet{rqInsensExeAllNum}{0}
\rqSet{rqInsensExeAllDen}{132}
\rqSet{rqInsensExeAllPct}{0.0}
\rqSet{rqInsensExeAllKiroNum}{0}
\rqSet{rqInsensExeAllKiroDen}{50}
\rqSet{rqInsensExeAllKiroPct}{0.0}
\rqSet{rqInsensExeAllClaudeNum}{0}
\rqSet{rqInsensExeAllClaudeDen}{51}
\rqSet{rqInsensExeAllClaudePct}{0.0}
\rqSet{rqInsensExeAllCodexNum}{0}
\rqSet{rqInsensExeAllCodexDen}{15}
\rqSet{rqInsensExeAllCodexPct}{0.0}
\rqSet{rqInsensExeAllCTwoRustNum}{0}
\rqSet{rqInsensExeAllCTwoRustDen}{3}
\rqSet{rqInsensExeAllCTwoRustPct}{0.0}
\rqSet{rqInsensExeAllLaertesNum}{0}
\rqSet{rqInsensExeAllLaertesDen}{3}
\rqSet{rqInsensExeAllLaertesPct}{0.0}
\rqSet{rqInsensExeAllCTwoSaferRustNum}{0}
\rqSet{rqInsensExeAllCTwoSaferRustDen}{2}
\rqSet{rqInsensExeAllCTwoSaferRustPct}{0.0}
\rqSet{rqInsensExeAllSmartCTwoRustNum}{0}
\rqSet{rqInsensExeAllSmartCTwoRustDen}{1}
\rqSet{rqInsensExeAllSmartCTwoRustPct}{0.0}
\rqSet{rqInsensExeAllKimiNum}{0}
\rqSet{rqInsensExeAllKimiDen}{2}
\rqSet{rqInsensExeAllKimiPct}{0.0}
\rqSet{rqInsensExeAllGptFiveNum}{0}
\rqSet{rqInsensExeAllGptFiveDen}{2}
\rqSet{rqInsensExeAllGptFivePct}{0.0}
\rqSet{rqInsensExeAllGeminiNum}{0}
\rqSet{rqInsensExeAllGeminiDen}{3}
\rqSet{rqInsensExeAllGeminiPct}{0.0}
\rqSet{rqInsensExeBOneNum}{0}
\rqSet{rqInsensExeBOneDen}{19}
\rqSet{rqInsensExeBOnePct}{0.0}
\rqSet{rqInsensExeBTwoNum}{0}
\rqSet{rqInsensExeBTwoDen}{5}
\rqSet{rqInsensExeBTwoPct}{0.0}
\rqSet{rqInsensExePZeroNum}{0}
\rqSet{rqInsensExePZeroDen}{0}
\rqSet{rqInsensExePZeroPct}{n/a}
\rqSet{rqInsensExePOneNum}{0}
\rqSet{rqInsensExePOneDen}{108}
\rqSet{rqInsensExePOnePct}{0.0}
\rqSet{rqInsensRustWithInputSymsExeAllNum}{0}
\rqSet{rqInsensRustWithInputSymsExeAllDen}{132}
\rqSet{rqInsensRustWithInputSymsExeAllPct}{0.0}
\rqSet{rqInsensRustWithInputSymsExeBOneNum}{0}
\rqSet{rqInsensRustWithInputSymsExeBOneDen}{19}
\rqSet{rqInsensRustWithInputSymsExeBOnePct}{0.0}
\rqSet{rqInsensRustWithInputSymsExeBTwoNum}{0}
\rqSet{rqInsensRustWithInputSymsExeBTwoDen}{5}
\rqSet{rqInsensRustWithInputSymsExeBTwoPct}{0.0}
\rqSet{rqInsensRustWithInputSymsExePZeroNum}{0}
\rqSet{rqInsensRustWithInputSymsExePZeroDen}{0}
\rqSet{rqInsensRustWithInputSymsExePZeroPct}{n/a}
\rqSet{rqInsensRustWithInputSymsExePOneNum}{0}
\rqSet{rqInsensRustWithInputSymsExePOneDen}{108}
\rqSet{rqInsensRustWithInputSymsExePOnePct}{0.0}
\rqSet{rqInsensRustWithoutInputSymsExeAllNum}{132}
\rqSet{rqInsensRustWithoutInputSymsExeAllDen}{132}
\rqSet{rqInsensRustWithoutInputSymsExeAllPct}{100.0}
\rqSet{rqInsensRustWithoutInputSymsExeBOneNum}{19}
\rqSet{rqInsensRustWithoutInputSymsExeBOneDen}{19}
\rqSet{rqInsensRustWithoutInputSymsExeBOnePct}{100.0}
\rqSet{rqInsensRustWithoutInputSymsExeBTwoNum}{5}
\rqSet{rqInsensRustWithoutInputSymsExeBTwoDen}{5}
\rqSet{rqInsensRustWithoutInputSymsExeBTwoPct}{100.0}
\rqSet{rqInsensRustWithoutInputSymsExePZeroNum}{0}
\rqSet{rqInsensRustWithoutInputSymsExePZeroDen}{0}
\rqSet{rqInsensRustWithoutInputSymsExePZeroPct}{n/a}
\rqSet{rqInsensRustWithoutInputSymsExePOneNum}{108}
\rqSet{rqInsensRustWithoutInputSymsExePOneDen}{108}
\rqSet{rqInsensRustWithoutInputSymsExePOnePct}{100.0}
\rqSet{rqInsensRustUncheckedExeAllNum}{0}
\rqSet{rqInsensRustUncheckedExeAllDen}{132}
\rqSet{rqInsensRustUncheckedExeAllPct}{0.0}
\rqSet{rqInsensRustUncheckedExeBOneNum}{0}
\rqSet{rqInsensRustUncheckedExeBOneDen}{19}
\rqSet{rqInsensRustUncheckedExeBOnePct}{0.0}
\rqSet{rqInsensRustUncheckedExeBTwoNum}{0}
\rqSet{rqInsensRustUncheckedExeBTwoDen}{5}
\rqSet{rqInsensRustUncheckedExeBTwoPct}{0.0}
\rqSet{rqInsensRustUncheckedExePZeroNum}{0}
\rqSet{rqInsensRustUncheckedExePZeroDen}{0}
\rqSet{rqInsensRustUncheckedExePZeroPct}{n/a}
\rqSet{rqInsensRustUncheckedExePOneNum}{0}
\rqSet{rqInsensRustUncheckedExePOneDen}{108}
\rqSet{rqInsensRustUncheckedExePOnePct}{0.0}
\rqSet{rqInsensRustNotRecordedExeAllNum}{0}
\rqSet{rqInsensRustNotRecordedExeAllDen}{132}
\rqSet{rqInsensRustNotRecordedExeAllPct}{0.0}
\rqSet{rqInsensRustNotRecordedExeBOneNum}{0}
\rqSet{rqInsensRustNotRecordedExeBOneDen}{19}
\rqSet{rqInsensRustNotRecordedExeBOnePct}{0.0}
\rqSet{rqInsensRustNotRecordedExeBTwoNum}{0}
\rqSet{rqInsensRustNotRecordedExeBTwoDen}{5}
\rqSet{rqInsensRustNotRecordedExeBTwoPct}{0.0}
\rqSet{rqInsensRustNotRecordedExePZeroNum}{0}
\rqSet{rqInsensRustNotRecordedExePZeroDen}{0}
\rqSet{rqInsensRustNotRecordedExePZeroPct}{n/a}
\rqSet{rqInsensRustNotRecordedExePOneNum}{0}
\rqSet{rqInsensRustNotRecordedExePOneDen}{108}
\rqSet{rqInsensRustNotRecordedExePOnePct}{0.0}
\rqSet{rqNBLoExeAllNum}{158}
\rqSet{rqNBLoExeAllDen}{476}
\rqSet{rqNBLoExeAllPct}{33.2}
\rqSet{rqNBLoExeAllKiroNum}{26}
\rqSet{rqNBLoExeAllKiroDen}{88}
\rqSet{rqNBLoExeAllKiroPct}{29.5}
\rqSet{rqNBLoExeAllClaudeNum}{22}
\rqSet{rqNBLoExeAllClaudeDen}{93}
\rqSet{rqNBLoExeAllClaudePct}{23.7}
\rqSet{rqNBLoExeAllCodexNum}{7}
\rqSet{rqNBLoExeAllCodexDen}{54}
\rqSet{rqNBLoExeAllCodexPct}{13.0}
\rqSet{rqNBLoExeAllCTwoRustNum}{2}
\rqSet{rqNBLoExeAllCTwoRustDen}{44}
\rqSet{rqNBLoExeAllCTwoRustPct}{4.5}
\rqSet{rqNBLoExeAllLaertesNum}{2}
\rqSet{rqNBLoExeAllLaertesDen}{44}
\rqSet{rqNBLoExeAllLaertesPct}{4.5}
\rqSet{rqNBLoExeAllCTwoSaferRustNum}{11}
\rqSet{rqNBLoExeAllCTwoSaferRustDen}{32}
\rqSet{rqNBLoExeAllCTwoSaferRustPct}{34.4}
\rqSet{rqNBLoExeAllSmartCTwoRustNum}{25}
\rqSet{rqNBLoExeAllSmartCTwoRustDen}{33}
\rqSet{rqNBLoExeAllSmartCTwoRustPct}{75.8}
\rqSet{rqNBLoExeAllKimiNum}{16}
\rqSet{rqNBLoExeAllKimiDen}{24}
\rqSet{rqNBLoExeAllKimiPct}{66.7}
\rqSet{rqNBLoExeAllGptFiveNum}{23}
\rqSet{rqNBLoExeAllGptFiveDen}{30}
\rqSet{rqNBLoExeAllGptFivePct}{76.7}
\rqSet{rqNBLoExeAllGeminiNum}{24}
\rqSet{rqNBLoExeAllGeminiDen}{34}
\rqSet{rqNBLoExeAllGeminiPct}{70.6}
\rqSet{rqNBLoExeBOneNum}{135}
\rqSet{rqNBLoExeBOneDen}{315}
\rqSet{rqNBLoExeBOnePct}{42.9}
\rqSet{rqNBLoExeBOneKiroNum}{23}
\rqSet{rqNBLoExeBOneKiroDen}{36}
\rqSet{rqNBLoExeBOneKiroPct}{63.9}
\rqSet{rqNBLoExeBOneClaudeNum}{16}
\rqSet{rqNBLoExeBOneClaudeDen}{36}
\rqSet{rqNBLoExeBOneClaudePct}{44.4}
\rqSet{rqNBLoExeBOneCodexNum}{5}
\rqSet{rqNBLoExeBOneCodexDen}{34}
\rqSet{rqNBLoExeBOneCodexPct}{14.7}
\rqSet{rqNBLoExeBOneCTwoRustNum}{0}
\rqSet{rqNBLoExeBOneCTwoRustDen}{36}
\rqSet{rqNBLoExeBOneCTwoRustPct}{0.0}
\rqSet{rqNBLoExeBOneLaertesNum}{0}
\rqSet{rqNBLoExeBOneLaertesDen}{36}
\rqSet{rqNBLoExeBOneLaertesPct}{0.0}
\rqSet{rqNBLoExeBOneCTwoSaferRustNum}{9}
\rqSet{rqNBLoExeBOneCTwoSaferRustDen}{27}
\rqSet{rqNBLoExeBOneCTwoSaferRustPct}{33.3}
\rqSet{rqNBLoExeBOneSmartCTwoRustNum}{24}
\rqSet{rqNBLoExeBOneSmartCTwoRustDen}{31}
\rqSet{rqNBLoExeBOneSmartCTwoRustPct}{77.4}
\rqSet{rqNBLoExeBOneKimiNum}{16}
\rqSet{rqNBLoExeBOneKimiDen}{23}
\rqSet{rqNBLoExeBOneKimiPct}{69.6}
\rqSet{rqNBLoExeBOneGptFiveNum}{21}
\rqSet{rqNBLoExeBOneGptFiveDen}{27}
\rqSet{rqNBLoExeBOneGptFivePct}{77.8}
\rqSet{rqNBLoExeBOneGeminiNum}{21}
\rqSet{rqNBLoExeBOneGeminiDen}{29}
\rqSet{rqNBLoExeBOneGeminiPct}{72.4}
\rqSet{rqNBLoExeBTwoNum}{23}
\rqSet{rqNBLoExeBTwoDen}{53}
\rqSet{rqNBLoExeBTwoPct}{43.4}
\rqSet{rqNBLoExeBTwoKiroNum}{3}
\rqSet{rqNBLoExeBTwoKiroDen}{4}
\rqSet{rqNBLoExeBTwoKiroPct}{75.0}
\rqSet{rqNBLoExeBTwoClaudeNum}{6}
\rqSet{rqNBLoExeBTwoClaudeDen}{9}
\rqSet{rqNBLoExeBTwoClaudePct}{66.7}
\rqSet{rqNBLoExeBTwoCodexNum}{2}
\rqSet{rqNBLoExeBTwoCodexDen}{8}
\rqSet{rqNBLoExeBTwoCodexPct}{25.0}
\rqSet{rqNBLoExeBTwoCTwoRustNum}{2}
\rqSet{rqNBLoExeBTwoCTwoRustDen}{8}
\rqSet{rqNBLoExeBTwoCTwoRustPct}{25.0}
\rqSet{rqNBLoExeBTwoLaertesNum}{2}
\rqSet{rqNBLoExeBTwoLaertesDen}{8}
\rqSet{rqNBLoExeBTwoLaertesPct}{25.0}
\rqSet{rqNBLoExeBTwoCTwoSaferRustNum}{2}
\rqSet{rqNBLoExeBTwoCTwoSaferRustDen}{5}
\rqSet{rqNBLoExeBTwoCTwoSaferRustPct}{40.0}
\rqSet{rqNBLoExeBTwoSmartCTwoRustNum}{1}
\rqSet{rqNBLoExeBTwoSmartCTwoRustDen}{2}
\rqSet{rqNBLoExeBTwoSmartCTwoRustPct}{50.0}
\rqSet{rqNBLoExeBTwoKimiNum}{0}
\rqSet{rqNBLoExeBTwoKimiDen}{1}
\rqSet{rqNBLoExeBTwoKimiPct}{0.0}
\rqSet{rqNBLoExeBTwoGptFiveNum}{2}
\rqSet{rqNBLoExeBTwoGptFiveDen}{3}
\rqSet{rqNBLoExeBTwoGptFivePct}{66.7}
\rqSet{rqNBLoExeBTwoGeminiNum}{3}
\rqSet{rqNBLoExeBTwoGeminiDen}{5}
\rqSet{rqNBLoExeBTwoGeminiPct}{60.0}
\rqSet{rqNBLoExePZeroNum}{0}
\rqSet{rqNBLoExePZeroDen}{0}
\rqSet{rqNBLoExePZeroPct}{n/a}
\rqSet{rqNBLoExePZeroKiroNum}{0}
\rqSet{rqNBLoExePZeroKiroDen}{0}
\rqSet{rqNBLoExePZeroKiroPct}{n/a}
\rqSet{rqNBLoExePZeroClaudeNum}{0}
\rqSet{rqNBLoExePZeroClaudeDen}{0}
\rqSet{rqNBLoExePZeroClaudePct}{n/a}
\rqSet{rqNBLoExePZeroCodexNum}{0}
\rqSet{rqNBLoExePZeroCodexDen}{0}
\rqSet{rqNBLoExePZeroCodexPct}{n/a}
\rqSet{rqNBLoExePZeroCTwoRustNum}{0}
\rqSet{rqNBLoExePZeroCTwoRustDen}{0}
\rqSet{rqNBLoExePZeroCTwoRustPct}{n/a}
\rqSet{rqNBLoExePZeroLaertesNum}{0}
\rqSet{rqNBLoExePZeroLaertesDen}{0}
\rqSet{rqNBLoExePZeroLaertesPct}{n/a}
\rqSet{rqNBLoExePZeroCTwoSaferRustNum}{0}
\rqSet{rqNBLoExePZeroCTwoSaferRustDen}{0}
\rqSet{rqNBLoExePZeroCTwoSaferRustPct}{n/a}
\rqSet{rqNBLoExePZeroSmartCTwoRustNum}{0}
\rqSet{rqNBLoExePZeroSmartCTwoRustDen}{0}
\rqSet{rqNBLoExePZeroSmartCTwoRustPct}{n/a}
\rqSet{rqNBLoExePZeroKimiNum}{0}
\rqSet{rqNBLoExePZeroKimiDen}{0}
\rqSet{rqNBLoExePZeroKimiPct}{n/a}
\rqSet{rqNBLoExePZeroGptFiveNum}{0}
\rqSet{rqNBLoExePZeroGptFiveDen}{0}
\rqSet{rqNBLoExePZeroGptFivePct}{n/a}
\rqSet{rqNBLoExePZeroGeminiNum}{0}
\rqSet{rqNBLoExePZeroGeminiDen}{0}
\rqSet{rqNBLoExePZeroGeminiPct}{n/a}
\rqSet{rqNBLoExePOneNum}{0}
\rqSet{rqNBLoExePOneDen}{108}
\rqSet{rqNBLoExePOnePct}{0.0}
\rqSet{rqNBLoExePOneKiroNum}{0}
\rqSet{rqNBLoExePOneKiroDen}{48}
\rqSet{rqNBLoExePOneKiroPct}{0.0}
\rqSet{rqNBLoExePOneClaudeNum}{0}
\rqSet{rqNBLoExePOneClaudeDen}{48}
\rqSet{rqNBLoExePOneClaudePct}{0.0}
\rqSet{rqNBLoExePOneCodexNum}{0}
\rqSet{rqNBLoExePOneCodexDen}{12}
\rqSet{rqNBLoExePOneCodexPct}{0.0}
\rqSet{rqNBLoExePOneCTwoRustNum}{0}
\rqSet{rqNBLoExePOneCTwoRustDen}{0}
\rqSet{rqNBLoExePOneCTwoRustPct}{n/a}
\rqSet{rqNBLoExePOneLaertesNum}{0}
\rqSet{rqNBLoExePOneLaertesDen}{0}
\rqSet{rqNBLoExePOneLaertesPct}{n/a}
\rqSet{rqNBLoExePOneCTwoSaferRustNum}{0}
\rqSet{rqNBLoExePOneCTwoSaferRustDen}{0}
\rqSet{rqNBLoExePOneCTwoSaferRustPct}{n/a}
\rqSet{rqNBLoExePOneSmartCTwoRustNum}{0}
\rqSet{rqNBLoExePOneSmartCTwoRustDen}{0}
\rqSet{rqNBLoExePOneSmartCTwoRustPct}{n/a}
\rqSet{rqNBLoExePOneKimiNum}{0}
\rqSet{rqNBLoExePOneKimiDen}{0}
\rqSet{rqNBLoExePOneKimiPct}{n/a}
\rqSet{rqNBLoExePOneGptFiveNum}{0}
\rqSet{rqNBLoExePOneGptFiveDen}{0}
\rqSet{rqNBLoExePOneGptFivePct}{n/a}
\rqSet{rqNBLoExePOneGeminiNum}{0}
\rqSet{rqNBLoExePOneGeminiDen}{0}
\rqSet{rqNBLoExePOneGeminiPct}{n/a}
\rqSet{rqNBHiExeAllNum}{174}
\rqSet{rqNBHiExeAllDen}{476}
\rqSet{rqNBHiExeAllPct}{36.6}
\rqSet{rqNBHiExeAllKiroNum}{29}
\rqSet{rqNBHiExeAllKiroDen}{88}
\rqSet{rqNBHiExeAllKiroPct}{33.0}
\rqSet{rqNBHiExeAllClaudeNum}{24}
\rqSet{rqNBHiExeAllClaudeDen}{93}
\rqSet{rqNBHiExeAllClaudePct}{25.8}
\rqSet{rqNBHiExeAllCodexNum}{9}
\rqSet{rqNBHiExeAllCodexDen}{54}
\rqSet{rqNBHiExeAllCodexPct}{16.7}
\rqSet{rqNBHiExeAllCTwoRustNum}{2}
\rqSet{rqNBHiExeAllCTwoRustDen}{44}
\rqSet{rqNBHiExeAllCTwoRustPct}{4.5}
\rqSet{rqNBHiExeAllLaertesNum}{2}
\rqSet{rqNBHiExeAllLaertesDen}{44}
\rqSet{rqNBHiExeAllLaertesPct}{4.5}
\rqSet{rqNBHiExeAllCTwoSaferRustNum}{12}
\rqSet{rqNBHiExeAllCTwoSaferRustDen}{32}
\rqSet{rqNBHiExeAllCTwoSaferRustPct}{37.5}
\rqSet{rqNBHiExeAllSmartCTwoRustNum}{27}
\rqSet{rqNBHiExeAllSmartCTwoRustDen}{33}
\rqSet{rqNBHiExeAllSmartCTwoRustPct}{81.8}
\rqSet{rqNBHiExeAllKimiNum}{18}
\rqSet{rqNBHiExeAllKimiDen}{24}
\rqSet{rqNBHiExeAllKimiPct}{75.0}
\rqSet{rqNBHiExeAllGptFiveNum}{25}
\rqSet{rqNBHiExeAllGptFiveDen}{30}
\rqSet{rqNBHiExeAllGptFivePct}{83.3}
\rqSet{rqNBHiExeAllGeminiNum}{26}
\rqSet{rqNBHiExeAllGeminiDen}{34}
\rqSet{rqNBHiExeAllGeminiPct}{76.5}
\rqSet{rqNBHiExeBOneNum}{150}
\rqSet{rqNBHiExeBOneDen}{315}
\rqSet{rqNBHiExeBOnePct}{47.6}
\rqSet{rqNBHiExeBOneKiroNum}{26}
\rqSet{rqNBHiExeBOneKiroDen}{36}
\rqSet{rqNBHiExeBOneKiroPct}{72.2}
\rqSet{rqNBHiExeBOneClaudeNum}{18}
\rqSet{rqNBHiExeBOneClaudeDen}{36}
\rqSet{rqNBHiExeBOneClaudePct}{50.0}
\rqSet{rqNBHiExeBOneCodexNum}{6}
\rqSet{rqNBHiExeBOneCodexDen}{34}
\rqSet{rqNBHiExeBOneCodexPct}{17.6}
\rqSet{rqNBHiExeBOneCTwoRustNum}{0}
\rqSet{rqNBHiExeBOneCTwoRustDen}{36}
\rqSet{rqNBHiExeBOneCTwoRustPct}{0.0}
\rqSet{rqNBHiExeBOneLaertesNum}{0}
\rqSet{rqNBHiExeBOneLaertesDen}{36}
\rqSet{rqNBHiExeBOneLaertesPct}{0.0}
\rqSet{rqNBHiExeBOneCTwoSaferRustNum}{10}
\rqSet{rqNBHiExeBOneCTwoSaferRustDen}{27}
\rqSet{rqNBHiExeBOneCTwoSaferRustPct}{37.0}
\rqSet{rqNBHiExeBOneSmartCTwoRustNum}{26}
\rqSet{rqNBHiExeBOneSmartCTwoRustDen}{31}
\rqSet{rqNBHiExeBOneSmartCTwoRustPct}{83.9}
\rqSet{rqNBHiExeBOneKimiNum}{18}
\rqSet{rqNBHiExeBOneKimiDen}{23}
\rqSet{rqNBHiExeBOneKimiPct}{78.3}
\rqSet{rqNBHiExeBOneGptFiveNum}{23}
\rqSet{rqNBHiExeBOneGptFiveDen}{27}
\rqSet{rqNBHiExeBOneGptFivePct}{85.2}
\rqSet{rqNBHiExeBOneGeminiNum}{23}
\rqSet{rqNBHiExeBOneGeminiDen}{29}
\rqSet{rqNBHiExeBOneGeminiPct}{79.3}
\rqSet{rqNBHiExeBTwoNum}{24}
\rqSet{rqNBHiExeBTwoDen}{53}
\rqSet{rqNBHiExeBTwoPct}{45.3}
\rqSet{rqNBHiExeBTwoKiroNum}{3}
\rqSet{rqNBHiExeBTwoKiroDen}{4}
\rqSet{rqNBHiExeBTwoKiroPct}{75.0}
\rqSet{rqNBHiExeBTwoClaudeNum}{6}
\rqSet{rqNBHiExeBTwoClaudeDen}{9}
\rqSet{rqNBHiExeBTwoClaudePct}{66.7}
\rqSet{rqNBHiExeBTwoCodexNum}{3}
\rqSet{rqNBHiExeBTwoCodexDen}{8}
\rqSet{rqNBHiExeBTwoCodexPct}{37.5}
\rqSet{rqNBHiExeBTwoCTwoRustNum}{2}
\rqSet{rqNBHiExeBTwoCTwoRustDen}{8}
\rqSet{rqNBHiExeBTwoCTwoRustPct}{25.0}
\rqSet{rqNBHiExeBTwoLaertesNum}{2}
\rqSet{rqNBHiExeBTwoLaertesDen}{8}
\rqSet{rqNBHiExeBTwoLaertesPct}{25.0}
\rqSet{rqNBHiExeBTwoCTwoSaferRustNum}{2}
\rqSet{rqNBHiExeBTwoCTwoSaferRustDen}{5}
\rqSet{rqNBHiExeBTwoCTwoSaferRustPct}{40.0}
\rqSet{rqNBHiExeBTwoSmartCTwoRustNum}{1}
\rqSet{rqNBHiExeBTwoSmartCTwoRustDen}{2}
\rqSet{rqNBHiExeBTwoSmartCTwoRustPct}{50.0}
\rqSet{rqNBHiExeBTwoKimiNum}{0}
\rqSet{rqNBHiExeBTwoKimiDen}{1}
\rqSet{rqNBHiExeBTwoKimiPct}{0.0}
\rqSet{rqNBHiExeBTwoGptFiveNum}{2}
\rqSet{rqNBHiExeBTwoGptFiveDen}{3}
\rqSet{rqNBHiExeBTwoGptFivePct}{66.7}
\rqSet{rqNBHiExeBTwoGeminiNum}{3}
\rqSet{rqNBHiExeBTwoGeminiDen}{5}
\rqSet{rqNBHiExeBTwoGeminiPct}{60.0}
\rqSet{rqNBHiExePZeroNum}{0}
\rqSet{rqNBHiExePZeroDen}{0}
\rqSet{rqNBHiExePZeroPct}{n/a}
\rqSet{rqNBHiExePZeroKiroNum}{0}
\rqSet{rqNBHiExePZeroKiroDen}{0}
\rqSet{rqNBHiExePZeroKiroPct}{n/a}
\rqSet{rqNBHiExePZeroClaudeNum}{0}
\rqSet{rqNBHiExePZeroClaudeDen}{0}
\rqSet{rqNBHiExePZeroClaudePct}{n/a}
\rqSet{rqNBHiExePZeroCodexNum}{0}
\rqSet{rqNBHiExePZeroCodexDen}{0}
\rqSet{rqNBHiExePZeroCodexPct}{n/a}
\rqSet{rqNBHiExePZeroCTwoRustNum}{0}
\rqSet{rqNBHiExePZeroCTwoRustDen}{0}
\rqSet{rqNBHiExePZeroCTwoRustPct}{n/a}
\rqSet{rqNBHiExePZeroLaertesNum}{0}
\rqSet{rqNBHiExePZeroLaertesDen}{0}
\rqSet{rqNBHiExePZeroLaertesPct}{n/a}
\rqSet{rqNBHiExePZeroCTwoSaferRustNum}{0}
\rqSet{rqNBHiExePZeroCTwoSaferRustDen}{0}
\rqSet{rqNBHiExePZeroCTwoSaferRustPct}{n/a}
\rqSet{rqNBHiExePZeroSmartCTwoRustNum}{0}
\rqSet{rqNBHiExePZeroSmartCTwoRustDen}{0}
\rqSet{rqNBHiExePZeroSmartCTwoRustPct}{n/a}
\rqSet{rqNBHiExePZeroKimiNum}{0}
\rqSet{rqNBHiExePZeroKimiDen}{0}
\rqSet{rqNBHiExePZeroKimiPct}{n/a}
\rqSet{rqNBHiExePZeroGptFiveNum}{0}
\rqSet{rqNBHiExePZeroGptFiveDen}{0}
\rqSet{rqNBHiExePZeroGptFivePct}{n/a}
\rqSet{rqNBHiExePZeroGeminiNum}{0}
\rqSet{rqNBHiExePZeroGeminiDen}{0}
\rqSet{rqNBHiExePZeroGeminiPct}{n/a}
\rqSet{rqNBHiExePOneNum}{0}
\rqSet{rqNBHiExePOneDen}{108}
\rqSet{rqNBHiExePOnePct}{0.0}
\rqSet{rqNBHiExePOneKiroNum}{0}
\rqSet{rqNBHiExePOneKiroDen}{48}
\rqSet{rqNBHiExePOneKiroPct}{0.0}
\rqSet{rqNBHiExePOneClaudeNum}{0}
\rqSet{rqNBHiExePOneClaudeDen}{48}
\rqSet{rqNBHiExePOneClaudePct}{0.0}
\rqSet{rqNBHiExePOneCodexNum}{0}
\rqSet{rqNBHiExePOneCodexDen}{12}
\rqSet{rqNBHiExePOneCodexPct}{0.0}
\rqSet{rqNBHiExePOneCTwoRustNum}{0}
\rqSet{rqNBHiExePOneCTwoRustDen}{0}
\rqSet{rqNBHiExePOneCTwoRustPct}{n/a}
\rqSet{rqNBHiExePOneLaertesNum}{0}
\rqSet{rqNBHiExePOneLaertesDen}{0}
\rqSet{rqNBHiExePOneLaertesPct}{n/a}
\rqSet{rqNBHiExePOneCTwoSaferRustNum}{0}
\rqSet{rqNBHiExePOneCTwoSaferRustDen}{0}
\rqSet{rqNBHiExePOneCTwoSaferRustPct}{n/a}
\rqSet{rqNBHiExePOneSmartCTwoRustNum}{0}
\rqSet{rqNBHiExePOneSmartCTwoRustDen}{0}
\rqSet{rqNBHiExePOneSmartCTwoRustPct}{n/a}
\rqSet{rqNBHiExePOneKimiNum}{0}
\rqSet{rqNBHiExePOneKimiDen}{0}
\rqSet{rqNBHiExePOneKimiPct}{n/a}
\rqSet{rqNBHiExePOneGptFiveNum}{0}
\rqSet{rqNBHiExePOneGptFiveDen}{0}
\rqSet{rqNBHiExePOneGptFivePct}{n/a}
\rqSet{rqNBHiExePOneGeminiNum}{0}
\rqSet{rqNBHiExePOneGeminiDen}{0}
\rqSet{rqNBHiExePOneGeminiPct}{n/a}
\rqSet{rqBnuDivBOneWhitespaceExeAllNum}{204123}
\rqSet{rqBnuDivBOneWhitespaceExeAllDen}{6655127}
\rqSet{rqBnuDivBOneWhitespaceExeAllPct}{3.1}
\rqSet{rqBnuDivBOneWhitespaceExeBOneNum}{204123}
\rqSet{rqBnuDivBOneWhitespaceExeBOneDen}{6157997}
\rqSet{rqBnuDivBOneWhitespaceExeBOnePct}{3.3}
\rqSet{rqBnuDivBOneWhitespaceExeBTwoNum}{0}
\rqSet{rqBnuDivBOneWhitespaceExeBTwoDen}{497130}
\rqSet{rqBnuDivBOneWhitespaceExeBTwoPct}{0.0}
\rqSet{rqBnuDivBOneWhitespaceExePZeroNum}{0}
\rqSet{rqBnuDivBOneWhitespaceExePZeroDen}{0}
\rqSet{rqBnuDivBOneWhitespaceExePZeroPct}{n/a}
\rqSet{rqBnuDivBOneWhitespaceExePOneNum}{0}
\rqSet{rqBnuDivBOneWhitespaceExePOneDen}{0}
\rqSet{rqBnuDivBOneWhitespaceExePOnePct}{n/a}
\rqSet{rqBnuTransBOneWhitespaceExeAllNum}{9}
\rqSet{rqBnuTransBOneWhitespaceExeAllDen}{176}
\rqSet{rqBnuTransBOneWhitespaceExeAllPct}{5.1}
\rqSet{rqBnuTransBOneWhitespaceExeBOneNum}{9}
\rqSet{rqBnuTransBOneWhitespaceExeBOneDen}{152}
\rqSet{rqBnuTransBOneWhitespaceExeBOnePct}{5.9}
\rqSet{rqBnuTransBOneWhitespaceExeBTwoNum}{0}
\rqSet{rqBnuTransBOneWhitespaceExeBTwoDen}{24}
\rqSet{rqBnuTransBOneWhitespaceExeBTwoPct}{0.0}
\rqSet{rqBnuTransBOneWhitespaceExePZeroNum}{0}
\rqSet{rqBnuTransBOneWhitespaceExePZeroDen}{0}
\rqSet{rqBnuTransBOneWhitespaceExePZeroPct}{n/a}
\rqSet{rqBnuTransBOneWhitespaceExePOneNum}{0}
\rqSet{rqBnuTransBOneWhitespaceExePOneDen}{0}
\rqSet{rqBnuTransBOneWhitespaceExePOnePct}{n/a}
\rqSet{rqBnuDivBTwoNumericExeAllNum}{543}
\rqSet{rqBnuDivBTwoNumericExeAllDen}{6655127}
\rqSet{rqBnuDivBTwoNumericExeAllPct}{0.0}
\rqSet{rqBnuDivBTwoNumericExeBOneNum}{543}
\rqSet{rqBnuDivBTwoNumericExeBOneDen}{6157997}
\rqSet{rqBnuDivBTwoNumericExeBOnePct}{0.0}
\rqSet{rqBnuDivBTwoNumericExeBTwoNum}{0}
\rqSet{rqBnuDivBTwoNumericExeBTwoDen}{497130}
\rqSet{rqBnuDivBTwoNumericExeBTwoPct}{0.0}
\rqSet{rqBnuDivBTwoNumericExePZeroNum}{0}
\rqSet{rqBnuDivBTwoNumericExePZeroDen}{0}
\rqSet{rqBnuDivBTwoNumericExePZeroPct}{n/a}
\rqSet{rqBnuDivBTwoNumericExePOneNum}{0}
\rqSet{rqBnuDivBTwoNumericExePOneDen}{0}
\rqSet{rqBnuDivBTwoNumericExePOnePct}{n/a}
\rqSet{rqBnuTransBTwoNumericExeAllNum}{14}
\rqSet{rqBnuTransBTwoNumericExeAllDen}{176}
\rqSet{rqBnuTransBTwoNumericExeAllPct}{8.0}
\rqSet{rqBnuTransBTwoNumericExeBOneNum}{14}
\rqSet{rqBnuTransBTwoNumericExeBOneDen}{152}
\rqSet{rqBnuTransBTwoNumericExeBOnePct}{9.2}
\rqSet{rqBnuTransBTwoNumericExeBTwoNum}{0}
\rqSet{rqBnuTransBTwoNumericExeBTwoDen}{24}
\rqSet{rqBnuTransBTwoNumericExeBTwoPct}{0.0}
\rqSet{rqBnuTransBTwoNumericExePZeroNum}{0}
\rqSet{rqBnuTransBTwoNumericExePZeroDen}{0}
\rqSet{rqBnuTransBTwoNumericExePZeroPct}{n/a}
\rqSet{rqBnuTransBTwoNumericExePOneNum}{0}
\rqSet{rqBnuTransBTwoNumericExePOneDen}{0}
\rqSet{rqBnuTransBTwoNumericExePOnePct}{n/a}
\rqSet{rqBnuDivNOneStatusExeAllNum}{1355545}
\rqSet{rqBnuDivNOneStatusExeAllDen}{6655127}
\rqSet{rqBnuDivNOneStatusExeAllPct}{20.4}
\rqSet{rqBnuDivNOneStatusExeBOneNum}{1188374}
\rqSet{rqBnuDivNOneStatusExeBOneDen}{6157997}
\rqSet{rqBnuDivNOneStatusExeBOnePct}{19.3}
\rqSet{rqBnuDivNOneStatusExeBTwoNum}{167171}
\rqSet{rqBnuDivNOneStatusExeBTwoDen}{497130}
\rqSet{rqBnuDivNOneStatusExeBTwoPct}{33.6}
\rqSet{rqBnuDivNOneStatusExePZeroNum}{0}
\rqSet{rqBnuDivNOneStatusExePZeroDen}{0}
\rqSet{rqBnuDivNOneStatusExePZeroPct}{n/a}
\rqSet{rqBnuDivNOneStatusExePOneNum}{0}
\rqSet{rqBnuDivNOneStatusExePOneDen}{0}
\rqSet{rqBnuDivNOneStatusExePOnePct}{n/a}
\rqSet{rqBnuTransNOneStatusExeAllNum}{39}
\rqSet{rqBnuTransNOneStatusExeAllDen}{176}
\rqSet{rqBnuTransNOneStatusExeAllPct}{22.2}
\rqSet{rqBnuTransNOneStatusExeBOneNum}{24}
\rqSet{rqBnuTransNOneStatusExeBOneDen}{152}
\rqSet{rqBnuTransNOneStatusExeBOnePct}{15.8}
\rqSet{rqBnuTransNOneStatusExeBTwoNum}{15}
\rqSet{rqBnuTransNOneStatusExeBTwoDen}{24}
\rqSet{rqBnuTransNOneStatusExeBTwoPct}{62.5}
\rqSet{rqBnuTransNOneStatusExePZeroNum}{0}
\rqSet{rqBnuTransNOneStatusExePZeroDen}{0}
\rqSet{rqBnuTransNOneStatusExePZeroPct}{n/a}
\rqSet{rqBnuTransNOneStatusExePOneNum}{0}
\rqSet{rqBnuTransNOneStatusExePOneDen}{0}
\rqSet{rqBnuTransNOneStatusExePOnePct}{n/a}
\rqSet{rqBnuDivNTwoTimeoutExeAllNum}{0}
\rqSet{rqBnuDivNTwoTimeoutExeAllDen}{6655127}
\rqSet{rqBnuDivNTwoTimeoutExeAllPct}{0.0}
\rqSet{rqBnuDivNTwoTimeoutExeBOneNum}{0}
\rqSet{rqBnuDivNTwoTimeoutExeBOneDen}{6157997}
\rqSet{rqBnuDivNTwoTimeoutExeBOnePct}{0.0}
\rqSet{rqBnuDivNTwoTimeoutExeBTwoNum}{0}
\rqSet{rqBnuDivNTwoTimeoutExeBTwoDen}{497130}
\rqSet{rqBnuDivNTwoTimeoutExeBTwoPct}{0.0}
\rqSet{rqBnuDivNTwoTimeoutExePZeroNum}{0}
\rqSet{rqBnuDivNTwoTimeoutExePZeroDen}{0}
\rqSet{rqBnuDivNTwoTimeoutExePZeroPct}{n/a}
\rqSet{rqBnuDivNTwoTimeoutExePOneNum}{0}
\rqSet{rqBnuDivNTwoTimeoutExePOneDen}{0}
\rqSet{rqBnuDivNTwoTimeoutExePOnePct}{n/a}
\rqSet{rqBnuTransNTwoTimeoutExeAllNum}{0}
\rqSet{rqBnuTransNTwoTimeoutExeAllDen}{176}
\rqSet{rqBnuTransNTwoTimeoutExeAllPct}{0.0}
\rqSet{rqBnuTransNTwoTimeoutExeBOneNum}{0}
\rqSet{rqBnuTransNTwoTimeoutExeBOneDen}{152}
\rqSet{rqBnuTransNTwoTimeoutExeBOnePct}{0.0}
\rqSet{rqBnuTransNTwoTimeoutExeBTwoNum}{0}
\rqSet{rqBnuTransNTwoTimeoutExeBTwoDen}{24}
\rqSet{rqBnuTransNTwoTimeoutExeBTwoPct}{0.0}
\rqSet{rqBnuTransNTwoTimeoutExePZeroNum}{0}
\rqSet{rqBnuTransNTwoTimeoutExePZeroDen}{0}
\rqSet{rqBnuTransNTwoTimeoutExePZeroPct}{n/a}
\rqSet{rqBnuTransNTwoTimeoutExePOneNum}{0}
\rqSet{rqBnuTransNTwoTimeoutExePOneDen}{0}
\rqSet{rqBnuTransNTwoTimeoutExePOnePct}{n/a}
\rqSet{rqBnuDivNThreeLineExeAllNum}{4207962}
\rqSet{rqBnuDivNThreeLineExeAllDen}{6655127}
\rqSet{rqBnuDivNThreeLineExeAllPct}{63.2}
\rqSet{rqBnuDivNThreeLineExeBOneNum}{4076050}
\rqSet{rqBnuDivNThreeLineExeBOneDen}{6157997}
\rqSet{rqBnuDivNThreeLineExeBOnePct}{66.2}
\rqSet{rqBnuDivNThreeLineExeBTwoNum}{131912}
\rqSet{rqBnuDivNThreeLineExeBTwoDen}{497130}
\rqSet{rqBnuDivNThreeLineExeBTwoPct}{26.5}
\rqSet{rqBnuDivNThreeLineExePZeroNum}{0}
\rqSet{rqBnuDivNThreeLineExePZeroDen}{0}
\rqSet{rqBnuDivNThreeLineExePZeroPct}{n/a}
\rqSet{rqBnuDivNThreeLineExePOneNum}{0}
\rqSet{rqBnuDivNThreeLineExePOneDen}{0}
\rqSet{rqBnuDivNThreeLineExePOnePct}{n/a}
\rqSet{rqBnuTransNThreeLineExeAllNum}{142}
\rqSet{rqBnuTransNThreeLineExeAllDen}{176}
\rqSet{rqBnuTransNThreeLineExeAllPct}{80.7}
\rqSet{rqBnuTransNThreeLineExeBOneNum}{125}
\rqSet{rqBnuTransNThreeLineExeBOneDen}{152}
\rqSet{rqBnuTransNThreeLineExeBOnePct}{82.2}
\rqSet{rqBnuTransNThreeLineExeBTwoNum}{17}
\rqSet{rqBnuTransNThreeLineExeBTwoDen}{24}
\rqSet{rqBnuTransNThreeLineExeBTwoPct}{70.8}
\rqSet{rqBnuTransNThreeLineExePZeroNum}{0}
\rqSet{rqBnuTransNThreeLineExePZeroDen}{0}
\rqSet{rqBnuTransNThreeLineExePZeroPct}{n/a}
\rqSet{rqBnuTransNThreeLineExePOneNum}{0}
\rqSet{rqBnuTransNThreeLineExePOneDen}{0}
\rqSet{rqBnuTransNThreeLineExePOnePct}{n/a}
\rqSet{rqBnuDivNFourFieldExeAllNum}{0}
\rqSet{rqBnuDivNFourFieldExeAllDen}{6655127}
\rqSet{rqBnuDivNFourFieldExeAllPct}{0.0}
\rqSet{rqBnuDivNFourFieldExeBOneNum}{0}
\rqSet{rqBnuDivNFourFieldExeBOneDen}{6157997}
\rqSet{rqBnuDivNFourFieldExeBOnePct}{0.0}
\rqSet{rqBnuDivNFourFieldExeBTwoNum}{0}
\rqSet{rqBnuDivNFourFieldExeBTwoDen}{497130}
\rqSet{rqBnuDivNFourFieldExeBTwoPct}{0.0}
\rqSet{rqBnuDivNFourFieldExePZeroNum}{0}
\rqSet{rqBnuDivNFourFieldExePZeroDen}{0}
\rqSet{rqBnuDivNFourFieldExePZeroPct}{n/a}
\rqSet{rqBnuDivNFourFieldExePOneNum}{0}
\rqSet{rqBnuDivNFourFieldExePOneDen}{0}
\rqSet{rqBnuDivNFourFieldExePOnePct}{n/a}
\rqSet{rqBnuTransNFourFieldExeAllNum}{0}
\rqSet{rqBnuTransNFourFieldExeAllDen}{176}
\rqSet{rqBnuTransNFourFieldExeAllPct}{0.0}
\rqSet{rqBnuTransNFourFieldExeBOneNum}{0}
\rqSet{rqBnuTransNFourFieldExeBOneDen}{152}
\rqSet{rqBnuTransNFourFieldExeBOnePct}{0.0}
\rqSet{rqBnuTransNFourFieldExeBTwoNum}{0}
\rqSet{rqBnuTransNFourFieldExeBTwoDen}{24}
\rqSet{rqBnuTransNFourFieldExeBTwoPct}{0.0}
\rqSet{rqBnuTransNFourFieldExePZeroNum}{0}
\rqSet{rqBnuTransNFourFieldExePZeroDen}{0}
\rqSet{rqBnuTransNFourFieldExePZeroPct}{n/a}
\rqSet{rqBnuTransNFourFieldExePOneNum}{0}
\rqSet{rqBnuTransNFourFieldExePOneDen}{0}
\rqSet{rqBnuTransNFourFieldExePOnePct}{n/a}
\rqSet{rqBnuDivUResidualExeAllNum}{694656}
\rqSet{rqBnuDivUResidualExeAllDen}{6655127}
\rqSet{rqBnuDivUResidualExeAllPct}{10.4}
\rqSet{rqBnuDivUResidualExeBOneNum}{685671}
\rqSet{rqBnuDivUResidualExeBOneDen}{6157997}
\rqSet{rqBnuDivUResidualExeBOnePct}{11.1}
\rqSet{rqBnuDivUResidualExeBTwoNum}{8985}
\rqSet{rqBnuDivUResidualExeBTwoDen}{497130}
\rqSet{rqBnuDivUResidualExeBTwoPct}{1.8}
\rqSet{rqBnuDivUResidualExePZeroNum}{0}
\rqSet{rqBnuDivUResidualExePZeroDen}{0}
\rqSet{rqBnuDivUResidualExePZeroPct}{n/a}
\rqSet{rqBnuDivUResidualExePOneNum}{0}
\rqSet{rqBnuDivUResidualExePOneDen}{0}
\rqSet{rqBnuDivUResidualExePOnePct}{n/a}
\rqSet{rqBnuTransUResidualExeAllNum}{23}
\rqSet{rqBnuTransUResidualExeAllDen}{176}
\rqSet{rqBnuTransUResidualExeAllPct}{13.1}
\rqSet{rqBnuTransUResidualExeBOneNum}{21}
\rqSet{rqBnuTransUResidualExeBOneDen}{152}
\rqSet{rqBnuTransUResidualExeBOnePct}{13.8}
\rqSet{rqBnuTransUResidualExeBTwoNum}{2}
\rqSet{rqBnuTransUResidualExeBTwoDen}{24}
\rqSet{rqBnuTransUResidualExeBTwoPct}{8.3}
\rqSet{rqBnuTransUResidualExePZeroNum}{0}
\rqSet{rqBnuTransUResidualExePZeroDen}{0}
\rqSet{rqBnuTransUResidualExePZeroPct}{n/a}
\rqSet{rqBnuTransUResidualExePOneNum}{0}
\rqSet{rqBnuTransUResidualExePOneDen}{0}
\rqSet{rqBnuTransUResidualExePOnePct}{n/a}
\rqSet{rqBnuDivUStderrExeAllNum}{192298}
\rqSet{rqBnuDivUStderrExeAllDen}{6655127}
\rqSet{rqBnuDivUStderrExeAllPct}{2.9}
\rqSet{rqBnuDivUStderrExeBOneNum}{3236}
\rqSet{rqBnuDivUStderrExeBOneDen}{6157997}
\rqSet{rqBnuDivUStderrExeBOnePct}{0.1}
\rqSet{rqBnuDivUStderrExeBTwoNum}{189062}
\rqSet{rqBnuDivUStderrExeBTwoDen}{497130}
\rqSet{rqBnuDivUStderrExeBTwoPct}{38.0}
\rqSet{rqBnuDivUStderrExePZeroNum}{0}
\rqSet{rqBnuDivUStderrExePZeroDen}{0}
\rqSet{rqBnuDivUStderrExePZeroPct}{n/a}
\rqSet{rqBnuDivUStderrExePOneNum}{0}
\rqSet{rqBnuDivUStderrExePOneDen}{0}
\rqSet{rqBnuDivUStderrExePOnePct}{n/a}
\rqSet{rqBnuTransUStderrExeAllNum}{11}
\rqSet{rqBnuTransUStderrExeAllDen}{176}
\rqSet{rqBnuTransUStderrExeAllPct}{6.2}
\rqSet{rqBnuTransUStderrExeBOneNum}{1}
\rqSet{rqBnuTransUStderrExeBOneDen}{152}
\rqSet{rqBnuTransUStderrExeBOnePct}{0.7}
\rqSet{rqBnuTransUStderrExeBTwoNum}{10}
\rqSet{rqBnuTransUStderrExeBTwoDen}{24}
\rqSet{rqBnuTransUStderrExeBTwoPct}{41.7}
\rqSet{rqBnuTransUStderrExePZeroNum}{0}
\rqSet{rqBnuTransUStderrExePZeroDen}{0}
\rqSet{rqBnuTransUStderrExePZeroPct}{n/a}
\rqSet{rqBnuTransUStderrExePOneNum}{0}
\rqSet{rqBnuTransUStderrExePOneDen}{0}
\rqSet{rqBnuTransUStderrExePOnePct}{n/a}
\rqSet{rqRobCCrashExeAllNum}{176}
\rqSet{rqRobCCrashExeAllDen}{476}
\rqSet{rqRobCCrashExeAllPct}{37.0}
\rqSet{rqRobCCrashExeAllKiroNum}{29}
\rqSet{rqRobCCrashExeAllKiroDen}{88}
\rqSet{rqRobCCrashExeAllKiroPct}{33.0}
\rqSet{rqRobCCrashExeAllClaudeNum}{24}
\rqSet{rqRobCCrashExeAllClaudeDen}{93}
\rqSet{rqRobCCrashExeAllClaudePct}{25.8}
\rqSet{rqRobCCrashExeAllCodexNum}{9}
\rqSet{rqRobCCrashExeAllCodexDen}{54}
\rqSet{rqRobCCrashExeAllCodexPct}{16.7}
\rqSet{rqRobCCrashExeAllCTwoRustNum}{3}
\rqSet{rqRobCCrashExeAllCTwoRustDen}{44}
\rqSet{rqRobCCrashExeAllCTwoRustPct}{6.8}
\rqSet{rqRobCCrashExeAllLaertesNum}{3}
\rqSet{rqRobCCrashExeAllLaertesDen}{44}
\rqSet{rqRobCCrashExeAllLaertesPct}{6.8}
\rqSet{rqRobCCrashExeAllCTwoSaferRustNum}{12}
\rqSet{rqRobCCrashExeAllCTwoSaferRustDen}{32}
\rqSet{rqRobCCrashExeAllCTwoSaferRustPct}{37.5}
\rqSet{rqRobCCrashExeAllSmartCTwoRustNum}{27}
\rqSet{rqRobCCrashExeAllSmartCTwoRustDen}{33}
\rqSet{rqRobCCrashExeAllSmartCTwoRustPct}{81.8}
\rqSet{rqRobCCrashExeAllKimiNum}{18}
\rqSet{rqRobCCrashExeAllKimiDen}{24}
\rqSet{rqRobCCrashExeAllKimiPct}{75.0}
\rqSet{rqRobCCrashExeAllGptFiveNum}{25}
\rqSet{rqRobCCrashExeAllGptFiveDen}{30}
\rqSet{rqRobCCrashExeAllGptFivePct}{83.3}
\rqSet{rqRobCCrashExeAllGeminiNum}{26}
\rqSet{rqRobCCrashExeAllGeminiDen}{34}
\rqSet{rqRobCCrashExeAllGeminiPct}{76.5}
\rqSet{rqRobCCrashExeBOneNum}{152}
\rqSet{rqRobCCrashExeBOneDen}{315}
\rqSet{rqRobCCrashExeBOnePct}{48.3}
\rqSet{rqRobCCrashExeBTwoNum}{24}
\rqSet{rqRobCCrashExeBTwoDen}{53}
\rqSet{rqRobCCrashExeBTwoPct}{45.3}
\rqSet{rqRobCCrashExePZeroNum}{0}
\rqSet{rqRobCCrashExePZeroDen}{0}
\rqSet{rqRobCCrashExePZeroPct}{n/a}
\rqSet{rqRobCCrashExePOneNum}{0}
\rqSet{rqRobCCrashExePOneDen}{108}
\rqSet{rqRobCCrashExePOnePct}{0.0}
\rqSet{rqRobUnstableExeAllNum}{176}
\rqSet{rqRobUnstableExeAllDen}{476}
\rqSet{rqRobUnstableExeAllPct}{37.0}
\rqSet{rqRobUnstableExeAllKiroNum}{29}
\rqSet{rqRobUnstableExeAllKiroDen}{88}
\rqSet{rqRobUnstableExeAllKiroPct}{33.0}
\rqSet{rqRobUnstableExeAllClaudeNum}{24}
\rqSet{rqRobUnstableExeAllClaudeDen}{93}
\rqSet{rqRobUnstableExeAllClaudePct}{25.8}
\rqSet{rqRobUnstableExeAllCodexNum}{9}
\rqSet{rqRobUnstableExeAllCodexDen}{54}
\rqSet{rqRobUnstableExeAllCodexPct}{16.7}
\rqSet{rqRobUnstableExeAllCTwoRustNum}{3}
\rqSet{rqRobUnstableExeAllCTwoRustDen}{44}
\rqSet{rqRobUnstableExeAllCTwoRustPct}{6.8}
\rqSet{rqRobUnstableExeAllLaertesNum}{3}
\rqSet{rqRobUnstableExeAllLaertesDen}{44}
\rqSet{rqRobUnstableExeAllLaertesPct}{6.8}
\rqSet{rqRobUnstableExeAllCTwoSaferRustNum}{12}
\rqSet{rqRobUnstableExeAllCTwoSaferRustDen}{32}
\rqSet{rqRobUnstableExeAllCTwoSaferRustPct}{37.5}
\rqSet{rqRobUnstableExeAllSmartCTwoRustNum}{27}
\rqSet{rqRobUnstableExeAllSmartCTwoRustDen}{33}
\rqSet{rqRobUnstableExeAllSmartCTwoRustPct}{81.8}
\rqSet{rqRobUnstableExeAllKimiNum}{18}
\rqSet{rqRobUnstableExeAllKimiDen}{24}
\rqSet{rqRobUnstableExeAllKimiPct}{75.0}
\rqSet{rqRobUnstableExeAllGptFiveNum}{25}
\rqSet{rqRobUnstableExeAllGptFiveDen}{30}
\rqSet{rqRobUnstableExeAllGptFivePct}{83.3}
\rqSet{rqRobUnstableExeAllGeminiNum}{26}
\rqSet{rqRobUnstableExeAllGeminiDen}{34}
\rqSet{rqRobUnstableExeAllGeminiPct}{76.5}
\rqSet{rqRobUnstableExeBOneNum}{152}
\rqSet{rqRobUnstableExeBOneDen}{315}
\rqSet{rqRobUnstableExeBOnePct}{48.3}
\rqSet{rqRobUnstableExeBTwoNum}{24}
\rqSet{rqRobUnstableExeBTwoDen}{53}
\rqSet{rqRobUnstableExeBTwoPct}{45.3}
\rqSet{rqRobUnstableExePZeroNum}{0}
\rqSet{rqRobUnstableExePZeroDen}{0}
\rqSet{rqRobUnstableExePZeroPct}{n/a}
\rqSet{rqRobUnstableExePOneNum}{0}
\rqSet{rqRobUnstableExePOneDen}{108}
\rqSet{rqRobUnstableExePOnePct}{0.0}
\rqSet{rqRobCompleteExeAllNum}{170}
\rqSet{rqRobCompleteExeAllDen}{317}
\rqSet{rqRobCompleteExeAllPct}{53.6}
\rqSet{rqRobCompleteExeAllKiroNum}{27}
\rqSet{rqRobCompleteExeAllKiroDen}{34}
\rqSet{rqRobCompleteExeAllKiroPct}{79.4}
\rqSet{rqRobCompleteExeAllClaudeNum}{23}
\rqSet{rqRobCompleteExeAllClaudeDen}{38}
\rqSet{rqRobCompleteExeAllClaudePct}{60.5}
\rqSet{rqRobCompleteExeAllCodexNum}{9}
\rqSet{rqRobCompleteExeAllCodexDen}{35}
\rqSet{rqRobCompleteExeAllCodexPct}{25.7}
\rqSet{rqRobCompleteExeAllCTwoRustNum}{3}
\rqSet{rqRobCompleteExeAllCTwoRustDen}{39}
\rqSet{rqRobCompleteExeAllCTwoRustPct}{7.7}
\rqSet{rqRobCompleteExeAllLaertesNum}{3}
\rqSet{rqRobCompleteExeAllLaertesDen}{39}
\rqSet{rqRobCompleteExeAllLaertesPct}{7.7}
\rqSet{rqRobCompleteExeAllCTwoSaferRustNum}{12}
\rqSet{rqRobCompleteExeAllCTwoSaferRustDen}{28}
\rqSet{rqRobCompleteExeAllCTwoSaferRustPct}{42.9}
\rqSet{rqRobCompleteExeAllSmartCTwoRustNum}{26}
\rqSet{rqRobCompleteExeAllSmartCTwoRustDen}{30}
\rqSet{rqRobCompleteExeAllSmartCTwoRustPct}{86.7}
\rqSet{rqRobCompleteExeAllKimiNum}{18}
\rqSet{rqRobCompleteExeAllKimiDen}{20}
\rqSet{rqRobCompleteExeAllKimiPct}{90.0}
\rqSet{rqRobCompleteExeAllGptFiveNum}{24}
\rqSet{rqRobCompleteExeAllGptFiveDen}{26}
\rqSet{rqRobCompleteExeAllGptFivePct}{92.3}
\rqSet{rqRobCompleteExeAllGeminiNum}{25}
\rqSet{rqRobCompleteExeAllGeminiDen}{28}
\rqSet{rqRobCompleteExeAllGeminiPct}{89.3}
\rqSet{rqRobCompleteExeBOneNum}{150}
\rqSet{rqRobCompleteExeBOneDen}{281}
\rqSet{rqRobCompleteExeBOnePct}{53.4}
\rqSet{rqRobCompleteExeBTwoNum}{20}
\rqSet{rqRobCompleteExeBTwoDen}{36}
\rqSet{rqRobCompleteExeBTwoPct}{55.6}
\rqSet{rqRobCompleteExePZeroNum}{0}
\rqSet{rqRobCompleteExePZeroDen}{0}
\rqSet{rqRobCompleteExePZeroPct}{n/a}
\rqSet{rqRobCompleteExePOneNum}{0}
\rqSet{rqRobCompleteExePOneDen}{0}
\rqSet{rqRobCompleteExePOnePct}{n/a}
\rqSet{rqRobNoStderrExeAllNum}{176}
\rqSet{rqRobNoStderrExeAllDen}{476}
\rqSet{rqRobNoStderrExeAllPct}{37.0}
\rqSet{rqRobNoStderrExeAllKiroNum}{29}
\rqSet{rqRobNoStderrExeAllKiroDen}{88}
\rqSet{rqRobNoStderrExeAllKiroPct}{33.0}
\rqSet{rqRobNoStderrExeAllClaudeNum}{24}
\rqSet{rqRobNoStderrExeAllClaudeDen}{93}
\rqSet{rqRobNoStderrExeAllClaudePct}{25.8}
\rqSet{rqRobNoStderrExeAllCodexNum}{9}
\rqSet{rqRobNoStderrExeAllCodexDen}{54}
\rqSet{rqRobNoStderrExeAllCodexPct}{16.7}
\rqSet{rqRobNoStderrExeAllCTwoRustNum}{3}
\rqSet{rqRobNoStderrExeAllCTwoRustDen}{44}
\rqSet{rqRobNoStderrExeAllCTwoRustPct}{6.8}
\rqSet{rqRobNoStderrExeAllLaertesNum}{3}
\rqSet{rqRobNoStderrExeAllLaertesDen}{44}
\rqSet{rqRobNoStderrExeAllLaertesPct}{6.8}
\rqSet{rqRobNoStderrExeAllCTwoSaferRustNum}{12}
\rqSet{rqRobNoStderrExeAllCTwoSaferRustDen}{32}
\rqSet{rqRobNoStderrExeAllCTwoSaferRustPct}{37.5}
\rqSet{rqRobNoStderrExeAllSmartCTwoRustNum}{27}
\rqSet{rqRobNoStderrExeAllSmartCTwoRustDen}{33}
\rqSet{rqRobNoStderrExeAllSmartCTwoRustPct}{81.8}
\rqSet{rqRobNoStderrExeAllKimiNum}{18}
\rqSet{rqRobNoStderrExeAllKimiDen}{24}
\rqSet{rqRobNoStderrExeAllKimiPct}{75.0}
\rqSet{rqRobNoStderrExeAllGptFiveNum}{25}
\rqSet{rqRobNoStderrExeAllGptFiveDen}{30}
\rqSet{rqRobNoStderrExeAllGptFivePct}{83.3}
\rqSet{rqRobNoStderrExeAllGeminiNum}{26}
\rqSet{rqRobNoStderrExeAllGeminiDen}{34}
\rqSet{rqRobNoStderrExeAllGeminiPct}{76.5}
\rqSet{rqRobNoStderrExeBOneNum}{152}
\rqSet{rqRobNoStderrExeBOneDen}{315}
\rqSet{rqRobNoStderrExeBOnePct}{48.3}
\rqSet{rqRobNoStderrExeBTwoNum}{24}
\rqSet{rqRobNoStderrExeBTwoDen}{53}
\rqSet{rqRobNoStderrExeBTwoPct}{45.3}
\rqSet{rqRobNoStderrExePZeroNum}{0}
\rqSet{rqRobNoStderrExePZeroDen}{0}
\rqSet{rqRobNoStderrExePZeroPct}{n/a}
\rqSet{rqRobNoStderrExePOneNum}{0}
\rqSet{rqRobNoStderrExePOneDen}{108}
\rqSet{rqRobNoStderrExePOnePct}{0.0}
\rqSet{rqRobCapEightExeAllNum}{176}
\rqSet{rqRobCapEightExeAllDen}{476}
\rqSet{rqRobCapEightExeAllPct}{37.0}
\rqSet{rqRobCapEightExeAllKiroNum}{29}
\rqSet{rqRobCapEightExeAllKiroDen}{88}
\rqSet{rqRobCapEightExeAllKiroPct}{33.0}
\rqSet{rqRobCapEightExeAllClaudeNum}{24}
\rqSet{rqRobCapEightExeAllClaudeDen}{93}
\rqSet{rqRobCapEightExeAllClaudePct}{25.8}
\rqSet{rqRobCapEightExeAllCodexNum}{9}
\rqSet{rqRobCapEightExeAllCodexDen}{54}
\rqSet{rqRobCapEightExeAllCodexPct}{16.7}
\rqSet{rqRobCapEightExeAllCTwoRustNum}{3}
\rqSet{rqRobCapEightExeAllCTwoRustDen}{44}
\rqSet{rqRobCapEightExeAllCTwoRustPct}{6.8}
\rqSet{rqRobCapEightExeAllLaertesNum}{3}
\rqSet{rqRobCapEightExeAllLaertesDen}{44}
\rqSet{rqRobCapEightExeAllLaertesPct}{6.8}
\rqSet{rqRobCapEightExeAllCTwoSaferRustNum}{12}
\rqSet{rqRobCapEightExeAllCTwoSaferRustDen}{32}
\rqSet{rqRobCapEightExeAllCTwoSaferRustPct}{37.5}
\rqSet{rqRobCapEightExeAllSmartCTwoRustNum}{27}
\rqSet{rqRobCapEightExeAllSmartCTwoRustDen}{33}
\rqSet{rqRobCapEightExeAllSmartCTwoRustPct}{81.8}
\rqSet{rqRobCapEightExeAllKimiNum}{18}
\rqSet{rqRobCapEightExeAllKimiDen}{24}
\rqSet{rqRobCapEightExeAllKimiPct}{75.0}
\rqSet{rqRobCapEightExeAllGptFiveNum}{25}
\rqSet{rqRobCapEightExeAllGptFiveDen}{30}
\rqSet{rqRobCapEightExeAllGptFivePct}{83.3}
\rqSet{rqRobCapEightExeAllGeminiNum}{26}
\rqSet{rqRobCapEightExeAllGeminiDen}{34}
\rqSet{rqRobCapEightExeAllGeminiPct}{76.5}
\rqSet{rqRobCapEightExeBOneNum}{152}
\rqSet{rqRobCapEightExeBOneDen}{315}
\rqSet{rqRobCapEightExeBOnePct}{48.3}
\rqSet{rqRobCapEightExeBTwoNum}{24}
\rqSet{rqRobCapEightExeBTwoDen}{53}
\rqSet{rqRobCapEightExeBTwoPct}{45.3}
\rqSet{rqRobCapEightExePZeroNum}{0}
\rqSet{rqRobCapEightExePZeroDen}{0}
\rqSet{rqRobCapEightExePZeroPct}{n/a}
\rqSet{rqRobCapEightExePOneNum}{0}
\rqSet{rqRobCapEightExePOneDen}{108}
\rqSet{rqRobCapEightExePOnePct}{0.0}
\rqSet{rqRobInconclExeAllNum}{176}
\rqSet{rqRobInconclExeAllDen}{476}
\rqSet{rqRobInconclExeAllPct}{37.0}
\rqSet{rqRobInconclExeAllKiroNum}{29}
\rqSet{rqRobInconclExeAllKiroDen}{88}
\rqSet{rqRobInconclExeAllKiroPct}{33.0}
\rqSet{rqRobInconclExeAllClaudeNum}{24}
\rqSet{rqRobInconclExeAllClaudeDen}{93}
\rqSet{rqRobInconclExeAllClaudePct}{25.8}
\rqSet{rqRobInconclExeAllCodexNum}{9}
\rqSet{rqRobInconclExeAllCodexDen}{54}
\rqSet{rqRobInconclExeAllCodexPct}{16.7}
\rqSet{rqRobInconclExeAllCTwoRustNum}{3}
\rqSet{rqRobInconclExeAllCTwoRustDen}{44}
\rqSet{rqRobInconclExeAllCTwoRustPct}{6.8}
\rqSet{rqRobInconclExeAllLaertesNum}{3}
\rqSet{rqRobInconclExeAllLaertesDen}{44}
\rqSet{rqRobInconclExeAllLaertesPct}{6.8}
\rqSet{rqRobInconclExeAllCTwoSaferRustNum}{12}
\rqSet{rqRobInconclExeAllCTwoSaferRustDen}{32}
\rqSet{rqRobInconclExeAllCTwoSaferRustPct}{37.5}
\rqSet{rqRobInconclExeAllSmartCTwoRustNum}{27}
\rqSet{rqRobInconclExeAllSmartCTwoRustDen}{33}
\rqSet{rqRobInconclExeAllSmartCTwoRustPct}{81.8}
\rqSet{rqRobInconclExeAllKimiNum}{18}
\rqSet{rqRobInconclExeAllKimiDen}{24}
\rqSet{rqRobInconclExeAllKimiPct}{75.0}
\rqSet{rqRobInconclExeAllGptFiveNum}{25}
\rqSet{rqRobInconclExeAllGptFiveDen}{30}
\rqSet{rqRobInconclExeAllGptFivePct}{83.3}
\rqSet{rqRobInconclExeAllGeminiNum}{26}
\rqSet{rqRobInconclExeAllGeminiDen}{34}
\rqSet{rqRobInconclExeAllGeminiPct}{76.5}
\rqSet{rqRobInconclExeBOneNum}{152}
\rqSet{rqRobInconclExeBOneDen}{315}
\rqSet{rqRobInconclExeBOnePct}{48.3}
\rqSet{rqRobInconclExeBTwoNum}{24}
\rqSet{rqRobInconclExeBTwoDen}{53}
\rqSet{rqRobInconclExeBTwoPct}{45.3}
\rqSet{rqRobInconclExePZeroNum}{0}
\rqSet{rqRobInconclExePZeroDen}{0}
\rqSet{rqRobInconclExePZeroPct}{n/a}
\rqSet{rqRobInconclExePOneNum}{0}
\rqSet{rqRobInconclExePOneDen}{108}
\rqSet{rqRobInconclExePOnePct}{0.0}
\rqSet{rqCapHitExeAllNum}{0}
\rqSet{rqCapHitExeAllDen}{35894359}
\rqSet{rqCapHitExeAllPct}{0.0}
\rqSet{rqCapHitExeBOneNum}{0}
\rqSet{rqCapHitExeBOneDen}{29751287}
\rqSet{rqCapHitExeBOnePct}{0.0}
\rqSet{rqCapHitExeBTwoNum}{0}
\rqSet{rqCapHitExeBTwoDen}{6142964}
\rqSet{rqCapHitExeBTwoPct}{0.0}
\rqSet{rqCapHitExePZeroNum}{0}
\rqSet{rqCapHitExePZeroDen}{0}
\rqSet{rqCapHitExePZeroPct}{n/a}
\rqSet{rqCapHitExePOneNum}{0}
\rqSet{rqCapHitExePOneDen}{108}
\rqSet{rqCapHitExePOnePct}{0.0}
\rqSet{rqCapHitDivExeAllNum}{0}
\rqSet{rqCapHitDivExeAllDen}{0}
\rqSet{rqCapHitDivExeAllPct}{n/a}
\rqSet{rqCapHitDivExeBOneNum}{0}
\rqSet{rqCapHitDivExeBOneDen}{0}
\rqSet{rqCapHitDivExeBOnePct}{n/a}
\rqSet{rqCapHitDivExeBTwoNum}{0}
\rqSet{rqCapHitDivExeBTwoDen}{0}
\rqSet{rqCapHitDivExeBTwoPct}{n/a}
\rqSet{rqCapHitDivExePZeroNum}{0}
\rqSet{rqCapHitDivExePZeroDen}{0}
\rqSet{rqCapHitDivExePZeroPct}{n/a}
\rqSet{rqCapHitDivExePOneNum}{0}
\rqSet{rqCapHitDivExePOneDen}{0}
\rqSet{rqCapHitDivExePOnePct}{n/a}
\rqSet{rqCapSensInfraExeAllNum}{0}
\rqSet{rqCapSensInfraExeAllDen}{0}
\rqSet{rqCapSensInfraExeAllPct}{n/a}
\rqSet{rqCapSensInfraExeBOneNum}{0}
\rqSet{rqCapSensInfraExeBOneDen}{0}
\rqSet{rqCapSensInfraExeBOnePct}{n/a}
\rqSet{rqCapSensInfraExeBTwoNum}{0}
\rqSet{rqCapSensInfraExeBTwoDen}{0}
\rqSet{rqCapSensInfraExeBTwoPct}{n/a}
\rqSet{rqCapSensInfraExePZeroNum}{0}
\rqSet{rqCapSensInfraExePZeroDen}{0}
\rqSet{rqCapSensInfraExePZeroPct}{n/a}
\rqSet{rqCapSensInfraExePOneNum}{0}
\rqSet{rqCapSensInfraExePOneDen}{0}
\rqSet{rqCapSensInfraExePOnePct}{n/a}
\rqSet{rqCapSensNewDivExeAllNum}{0}
\rqSet{rqCapSensNewDivExeAllDen}{0}
\rqSet{rqCapSensNewDivExeBOneNum}{0}
\rqSet{rqCapSensNewDivExeBOneDen}{0}
\rqSet{rqCapSensNewDivExeBTwoNum}{0}
\rqSet{rqCapSensNewDivExeBTwoDen}{0}
\rqSet{rqCapSensNewDivExePZeroNum}{0}
\rqSet{rqCapSensNewDivExePZeroDen}{0}
\rqSet{rqCapSensNewDivExePOneNum}{0}
\rqSet{rqCapSensNewDivExePOneDen}{0}
\rqSet{rqCapSensTimeoutRerunExeAllNum}{0}
\rqSet{rqCapSensTimeoutRerunExeAllDen}{0}
\rqSet{rqCapSensTimeoutRerunExeAllPct}{n/a}
\rqSet{rqCapSensTimeoutRerunExeBOneNum}{0}
\rqSet{rqCapSensTimeoutRerunExeBOneDen}{0}
\rqSet{rqCapSensTimeoutRerunExeBOnePct}{n/a}
\rqSet{rqCapSensTimeoutRerunExeBTwoNum}{0}
\rqSet{rqCapSensTimeoutRerunExeBTwoDen}{0}
\rqSet{rqCapSensTimeoutRerunExeBTwoPct}{n/a}
\rqSet{rqCapSensTimeoutRerunExePZeroNum}{0}
\rqSet{rqCapSensTimeoutRerunExePZeroDen}{0}
\rqSet{rqCapSensTimeoutRerunExePZeroPct}{n/a}
\rqSet{rqCapSensTimeoutRerunExePOneNum}{0}
\rqSet{rqCapSensTimeoutRerunExePOneDen}{0}
\rqSet{rqCapSensTimeoutRerunExePOnePct}{n/a}
\rqSet{rqCapOneSideExeAllNum}{0}
\rqSet{rqCapOneSideExeAllDen}{476}
\rqSet{rqCapOneSideExeAllPct}{0.0}
\rqSet{rqCapOneSideExeBOneNum}{0}
\rqSet{rqCapOneSideExeBOneDen}{315}
\rqSet{rqCapOneSideExeBOnePct}{0.0}
\rqSet{rqCapOneSideExeBTwoNum}{0}
\rqSet{rqCapOneSideExeBTwoDen}{53}
\rqSet{rqCapOneSideExeBTwoPct}{0.0}
\rqSet{rqCapOneSideExePZeroNum}{0}
\rqSet{rqCapOneSideExePZeroDen}{0}
\rqSet{rqCapOneSideExePZeroPct}{n/a}
\rqSet{rqCapOneSideExePOneNum}{0}
\rqSet{rqCapOneSideExePOneDen}{108}
\rqSet{rqCapOneSideExePOnePct}{0.0}
\rqSet{rqUbDivExeAllNum}{845044}
\rqSet{rqUbDivExeAllDen}{7532662}
\rqSet{rqUbDivExeAllPct}{11.2}
\rqSet{rqUbDivExeAllKiroNum}{77759}
\rqSet{rqUbDivExeAllKiroDen}{1216931}
\rqSet{rqUbDivExeAllKiroPct}{6.4}
\rqSet{rqUbDivExeAllClaudeNum}{63238}
\rqSet{rqUbDivExeAllClaudeDen}{543412}
\rqSet{rqUbDivExeAllClaudePct}{11.6}
\rqSet{rqUbDivExeAllCodexNum}{37051}
\rqSet{rqUbDivExeAllCodexDen}{176657}
\rqSet{rqUbDivExeAllCodexPct}{21.0}
\rqSet{rqUbDivExeAllCTwoRustNum}{103572}
\rqSet{rqUbDivExeAllCTwoRustDen}{134882}
\rqSet{rqUbDivExeAllCTwoRustPct}{76.8}
\rqSet{rqUbDivExeAllLaertesNum}{103572}
\rqSet{rqUbDivExeAllLaertesDen}{134882}
\rqSet{rqUbDivExeAllLaertesPct}{76.8}
\rqSet{rqUbDivExeAllCTwoSaferRustNum}{77115}
\rqSet{rqUbDivExeAllCTwoSaferRustDen}{529321}
\rqSet{rqUbDivExeAllCTwoSaferRustPct}{14.6}
\rqSet{rqUbDivExeAllSmartCTwoRustNum}{101960}
\rqSet{rqUbDivExeAllSmartCTwoRustDen}{1223287}
\rqSet{rqUbDivExeAllSmartCTwoRustPct}{8.3}
\rqSet{rqUbDivExeAllKimiNum}{92453}
\rqSet{rqUbDivExeAllKimiDen}{1209030}
\rqSet{rqUbDivExeAllKimiPct}{7.6}
\rqSet{rqUbDivExeAllGptFiveNum}{93997}
\rqSet{rqUbDivExeAllGptFiveDen}{1189818}
\rqSet{rqUbDivExeAllGptFivePct}{7.9}
\rqSet{rqUbDivExeAllGeminiNum}{94327}
\rqSet{rqUbDivExeAllGeminiDen}{1174442}
\rqSet{rqUbDivExeAllGeminiPct}{8.0}
\rqSet{rqUbDivExeBOneNum}{843550}
\rqSet{rqUbDivExeBOneDen}{7001547}
\rqSet{rqUbDivExeBOnePct}{12.0}
\rqSet{rqUbDivExeBTwoNum}{1494}
\rqSet{rqUbDivExeBTwoDen}{531115}
\rqSet{rqUbDivExeBTwoPct}{0.3}
\rqSet{rqUbDivExePZeroNum}{0}
\rqSet{rqUbDivExePZeroDen}{0}
\rqSet{rqUbDivExePZeroPct}{n/a}
\rqSet{rqUbDivExePOneNum}{0}
\rqSet{rqUbDivExePOneDen}{0}
\rqSet{rqUbDivExePOnePct}{n/a}
\rqSet{rqUbDivStableExeAllNum}{506842}
\rqSet{rqUbDivStableExeAllDen}{7161969}
\rqSet{rqUbDivStableExeAllPct}{7.1}
\rqSet{rqUbDivStableExeBOneNum}{505392}
\rqSet{rqUbDivStableExeBOneDen}{6663389}
\rqSet{rqUbDivStableExeBOnePct}{7.6}
\rqSet{rqUbDivStableExeBTwoNum}{1450}
\rqSet{rqUbDivStableExeBTwoDen}{498580}
\rqSet{rqUbDivStableExeBTwoPct}{0.3}
\rqSet{rqUbDivStableExePZeroNum}{0}
\rqSet{rqUbDivStableExePZeroDen}{0}
\rqSet{rqUbDivStableExePZeroPct}{n/a}
\rqSet{rqUbDivStableExePOneNum}{0}
\rqSet{rqUbDivStableExePOneDen}{0}
\rqSet{rqUbDivStableExePOnePct}{n/a}
\rqSet{rqUbInputsExeAllNum}{140024}
\rqSet{rqUbInputsExeAllDen}{1926540}
\rqSet{rqUbInputsExeAllPct}{7.3}
\rqSet{rqUbInputsExeBOneNum}{139788}
\rqSet{rqUbInputsExeBOneDen}{1744834}
\rqSet{rqUbInputsExeBOnePct}{8.0}
\rqSet{rqUbInputsExeBTwoNum}{236}
\rqSet{rqUbInputsExeBTwoDen}{181706}
\rqSet{rqUbInputsExeBTwoPct}{0.1}
\rqSet{rqUbInputsExePZeroNum}{0}
\rqSet{rqUbInputsExePZeroDen}{0}
\rqSet{rqUbInputsExePZeroPct}{n/a}
\rqSet{rqUbInputsExePOneNum}{0}
\rqSet{rqUbInputsExePOneDen}{0}
\rqSet{rqUbInputsExePOnePct}{n/a}
\rqSet{rqUbInputsStableExeAllNum}{107066}
\rqSet{rqUbInputsStableExeAllDen}{1885509}
\rqSet{rqUbInputsStableExeAllPct}{5.7}
\rqSet{rqUbInputsStableExeBOneNum}{106840}
\rqSet{rqUbInputsStableExeBOneDen}{1711886}
\rqSet{rqUbInputsStableExeBOnePct}{6.2}
\rqSet{rqUbInputsStableExeBTwoNum}{226}
\rqSet{rqUbInputsStableExeBTwoDen}{173623}
\rqSet{rqUbInputsStableExeBTwoPct}{0.1}
\rqSet{rqUbInputsStableExePZeroNum}{0}
\rqSet{rqUbInputsStableExePZeroDen}{0}
\rqSet{rqUbInputsStableExePZeroPct}{n/a}
\rqSet{rqUbInputsStableExePOneNum}{0}
\rqSet{rqUbInputsStableExePOneDen}{0}
\rqSet{rqUbInputsStableExePOnePct}{n/a}
\rqSet{rqUbTransExeAllNum}{99}
\rqSet{rqUbTransExeAllDen}{206}
\rqSet{rqUbTransExeAllPct}{48.1}
\rqSet{rqUbTransExeBOneNum}{77}
\rqSet{rqUbTransExeBOneDen}{169}
\rqSet{rqUbTransExeBOnePct}{45.6}
\rqSet{rqUbTransExeBTwoNum}{22}
\rqSet{rqUbTransExeBTwoDen}{37}
\rqSet{rqUbTransExeBTwoPct}{59.5}
\rqSet{rqUbTransExePZeroNum}{0}
\rqSet{rqUbTransExePZeroDen}{0}
\rqSet{rqUbTransExePZeroPct}{n/a}
\rqSet{rqUbTransExePOneNum}{0}
\rqSet{rqUbTransExePOneDen}{0}
\rqSet{rqUbTransExePOnePct}{n/a}
\rqSet{rqStderrOnlyDivExeAllNum}{192298}
\rqSet{rqStderrOnlyDivExeAllDen}{6655127}
\rqSet{rqStderrOnlyDivExeAllPct}{2.9}
\rqSet{rqStderrOnlyDivExeAllKiroNum}{48900}
\rqSet{rqStderrOnlyDivExeAllKiroDen}{1139170}
\rqSet{rqStderrOnlyDivExeAllKiroPct}{4.3}
\rqSet{rqStderrOnlyDivExeAllClaudeNum}{30618}
\rqSet{rqStderrOnlyDivExeAllClaudeDen}{472054}
\rqSet{rqStderrOnlyDivExeAllClaudePct}{6.5}
\rqSet{rqStderrOnlyDivExeAllCodexNum}{14149}
\rqSet{rqStderrOnlyDivExeAllCodexDen}{139591}
\rqSet{rqStderrOnlyDivExeAllCodexPct}{10.1}
\rqSet{rqStderrOnlyDivExeAllCTwoRustNum}{0}
\rqSet{rqStderrOnlyDivExeAllCTwoRustDen}{23194}
\rqSet{rqStderrOnlyDivExeAllCTwoRustPct}{0.0}
\rqSet{rqStderrOnlyDivExeAllLaertesNum}{0}
\rqSet{rqStderrOnlyDivExeAllLaertesDen}{23192}
\rqSet{rqStderrOnlyDivExeAllLaertesPct}{0.0}
\rqSet{rqStderrOnlyDivExeAllCTwoSaferRustNum}{0}
\rqSet{rqStderrOnlyDivExeAllCTwoSaferRustDen}{444091}
\rqSet{rqStderrOnlyDivExeAllCTwoSaferRustPct}{0.0}
\rqSet{rqStderrOnlyDivExeAllSmartCTwoRustNum}{26507}
\rqSet{rqStderrOnlyDivExeAllSmartCTwoRustDen}{1121327}
\rqSet{rqStderrOnlyDivExeAllSmartCTwoRustPct}{2.4}
\rqSet{rqStderrOnlyDivExeAllKimiNum}{0}
\rqSet{rqStderrOnlyDivExeAllKimiDen}{1116577}
\rqSet{rqStderrOnlyDivExeAllKimiPct}{0.0}
\rqSet{rqStderrOnlyDivExeAllGptFiveNum}{23224}
\rqSet{rqStderrOnlyDivExeAllGptFiveDen}{1095821}
\rqSet{rqStderrOnlyDivExeAllGptFivePct}{2.1}
\rqSet{rqStderrOnlyDivExeAllGeminiNum}{48900}
\rqSet{rqStderrOnlyDivExeAllGeminiDen}{1080110}
\rqSet{rqStderrOnlyDivExeAllGeminiPct}{4.5}
\rqSet{rqStderrOnlyDivExeBOneNum}{3236}
\rqSet{rqStderrOnlyDivExeBOneDen}{6157997}
\rqSet{rqStderrOnlyDivExeBOnePct}{0.1}
\rqSet{rqStderrOnlyDivExeBTwoNum}{189062}
\rqSet{rqStderrOnlyDivExeBTwoDen}{497130}
\rqSet{rqStderrOnlyDivExeBTwoPct}{38.0}
\rqSet{rqStderrOnlyDivExePZeroNum}{0}
\rqSet{rqStderrOnlyDivExePZeroDen}{0}
\rqSet{rqStderrOnlyDivExePZeroPct}{n/a}
\rqSet{rqStderrOnlyDivExePOneNum}{0}
\rqSet{rqStderrOnlyDivExePOneDen}{0}
\rqSet{rqStderrOnlyDivExePOnePct}{n/a}
\rqSet{rqStderrOnlyTransExeAllNum}{0}
\rqSet{rqStderrOnlyTransExeAllDen}{176}
\rqSet{rqStderrOnlyTransExeAllPct}{0.0}
\rqSet{rqStderrOnlyTransExeBOneNum}{0}
\rqSet{rqStderrOnlyTransExeBOneDen}{152}
\rqSet{rqStderrOnlyTransExeBOnePct}{0.0}
\rqSet{rqStderrOnlyTransExeBTwoNum}{0}
\rqSet{rqStderrOnlyTransExeBTwoDen}{24}
\rqSet{rqStderrOnlyTransExeBTwoPct}{0.0}
\rqSet{rqStderrOnlyTransExePZeroNum}{0}
\rqSet{rqStderrOnlyTransExePZeroDen}{0}
\rqSet{rqStderrOnlyTransExePZeroPct}{n/a}
\rqSet{rqStderrOnlyTransExePOneNum}{0}
\rqSet{rqStderrOnlyTransExePOneDen}{0}
\rqSet{rqStderrOnlyTransExePOnePct}{n/a}
\rqSet{rqStderrOnlyInputsExeAllNum}{192325}
\rqSet{rqStderrOnlyInputsExeAllDen}{6655127}
\rqSet{rqStderrOnlyInputsExeBOneNum}{3236}
\rqSet{rqStderrOnlyInputsExeBOneDen}{6157997}
\rqSet{rqStderrOnlyInputsExeBTwoNum}{189089}
\rqSet{rqStderrOnlyInputsExeBTwoDen}{497130}
\rqSet{rqStderrOnlyInputsExePZeroNum}{0}
\rqSet{rqStderrOnlyInputsExePZeroDen}{0}
\rqSet{rqStderrOnlyInputsExePOneNum}{0}
\rqSet{rqStderrOnlyInputsExePOneDen}{0}
\rqSet{rqBothNtMatchExeAllNum}{0}
\rqSet{rqBothNtMatchExeAllDen}{35894359}
\rqSet{rqBothNtMatchExeAllPct}{0.0}
\rqSet{rqBothNtMatchExeBOneNum}{0}
\rqSet{rqBothNtMatchExeBOneDen}{29751287}
\rqSet{rqBothNtMatchExeBOnePct}{0.0}
\rqSet{rqBothNtMatchExeBTwoNum}{0}
\rqSet{rqBothNtMatchExeBTwoDen}{6142964}
\rqSet{rqBothNtMatchExeBTwoPct}{0.0}
\rqSet{rqBothNtMatchExePZeroNum}{0}
\rqSet{rqBothNtMatchExePZeroDen}{0}
\rqSet{rqBothNtMatchExePZeroPct}{n/a}
\rqSet{rqBothNtMatchExePOneNum}{0}
\rqSet{rqBothNtMatchExePOneDen}{108}
\rqSet{rqBothNtMatchExePOnePct}{0.0}
\rqSet{rqBothNtPrefixZeroExeAllNum}{0}
\rqSet{rqBothNtPrefixZeroExeAllDen}{0}
\rqSet{rqBothNtPrefixZeroExeAllPct}{n/a}
\rqSet{rqBothNtPrefixZeroExeBOneNum}{0}
\rqSet{rqBothNtPrefixZeroExeBOneDen}{0}
\rqSet{rqBothNtPrefixZeroExeBOnePct}{n/a}
\rqSet{rqBothNtPrefixZeroExeBTwoNum}{0}
\rqSet{rqBothNtPrefixZeroExeBTwoDen}{0}
\rqSet{rqBothNtPrefixZeroExeBTwoPct}{n/a}
\rqSet{rqBothNtPrefixZeroExePZeroNum}{0}
\rqSet{rqBothNtPrefixZeroExePZeroDen}{0}
\rqSet{rqBothNtPrefixZeroExePZeroPct}{n/a}
\rqSet{rqBothNtPrefixZeroExePOneNum}{0}
\rqSet{rqBothNtPrefixZeroExePOneDen}{0}
\rqSet{rqBothNtPrefixZeroExePOnePct}{n/a}
\rqSet{rqBothNtPrefixBytesExeAllNum}{0}
\rqSet{rqBothNtPrefixBytesExeAllDen}{0}
\rqSet{rqBothNtPrefixBytesExeBOneNum}{0}
\rqSet{rqBothNtPrefixBytesExeBOneDen}{0}
\rqSet{rqBothNtPrefixBytesExeBTwoNum}{0}
\rqSet{rqBothNtPrefixBytesExeBTwoDen}{0}
\rqSet{rqBothNtPrefixBytesExePZeroNum}{0}
\rqSet{rqBothNtPrefixBytesExePZeroDen}{0}
\rqSet{rqBothNtPrefixBytesExePOneNum}{0}
\rqSet{rqBothNtPrefixBytesExePOneDen}{0}
\rqSet{rqNtOneSideConfirmedExeAllNum}{0}
\rqSet{rqNtOneSideConfirmedExeAllDen}{6655127}
\rqSet{rqNtOneSideConfirmedExeAllPct}{0.0}
\rqSet{rqNtOneSideConfirmedExeBOneNum}{0}
\rqSet{rqNtOneSideConfirmedExeBOneDen}{6157997}
\rqSet{rqNtOneSideConfirmedExeBOnePct}{0.0}
\rqSet{rqNtOneSideConfirmedExeBTwoNum}{0}
\rqSet{rqNtOneSideConfirmedExeBTwoDen}{497130}
\rqSet{rqNtOneSideConfirmedExeBTwoPct}{0.0}
\rqSet{rqNtOneSideConfirmedExePZeroNum}{0}
\rqSet{rqNtOneSideConfirmedExePZeroDen}{0}
\rqSet{rqNtOneSideConfirmedExePZeroPct}{n/a}
\rqSet{rqNtOneSideConfirmedExePOneNum}{0}
\rqSet{rqNtOneSideConfirmedExePOneDen}{0}
\rqSet{rqNtOneSideConfirmedExePOnePct}{n/a}
\rqSet{rqNtOneSideConfirmedAllExeAllNum}{0}
\rqSet{rqNtOneSideConfirmedAllExeAllDen}{7532662}
\rqSet{rqNtOneSideConfirmedAllExeAllPct}{0.0}
\rqSet{rqNtOneSideConfirmedAllExeBOneNum}{0}
\rqSet{rqNtOneSideConfirmedAllExeBOneDen}{7001547}
\rqSet{rqNtOneSideConfirmedAllExeBOnePct}{0.0}
\rqSet{rqNtOneSideConfirmedAllExeBTwoNum}{0}
\rqSet{rqNtOneSideConfirmedAllExeBTwoDen}{531115}
\rqSet{rqNtOneSideConfirmedAllExeBTwoPct}{0.0}
\rqSet{rqNtOneSideConfirmedAllExePZeroNum}{0}
\rqSet{rqNtOneSideConfirmedAllExePZeroDen}{0}
\rqSet{rqNtOneSideConfirmedAllExePZeroPct}{n/a}
\rqSet{rqNtOneSideConfirmedAllExePOneNum}{0}
\rqSet{rqNtOneSideConfirmedAllExePOneDen}{0}
\rqSet{rqNtOneSideConfirmedAllExePOnePct}{n/a}
\rqSet{rqNtOneSideRerunAllExeAllNum}{0}
\rqSet{rqNtOneSideRerunAllExeAllDen}{7532662}
\rqSet{rqNtOneSideRerunAllExeAllPct}{0.0}
\rqSet{rqNtOneSideRerunAllExeBOneNum}{0}
\rqSet{rqNtOneSideRerunAllExeBOneDen}{7001547}
\rqSet{rqNtOneSideRerunAllExeBOnePct}{0.0}
\rqSet{rqNtOneSideRerunAllExeBTwoNum}{0}
\rqSet{rqNtOneSideRerunAllExeBTwoDen}{531115}
\rqSet{rqNtOneSideRerunAllExeBTwoPct}{0.0}
\rqSet{rqNtOneSideRerunAllExePZeroNum}{0}
\rqSet{rqNtOneSideRerunAllExePZeroDen}{0}
\rqSet{rqNtOneSideRerunAllExePZeroPct}{n/a}
\rqSet{rqNtOneSideRerunAllExePOneNum}{0}
\rqSet{rqNtOneSideRerunAllExePOneDen}{0}
\rqSet{rqNtOneSideRerunAllExePOnePct}{n/a}
\rqSet{rqNtOneSideTransAnyExeAllNum}{0}
\rqSet{rqNtOneSideTransAnyExeAllDen}{176}
\rqSet{rqNtOneSideTransAnyExeAllPct}{0.0}
\rqSet{rqNtOneSideTransAnyExeBOneNum}{0}
\rqSet{rqNtOneSideTransAnyExeBOneDen}{152}
\rqSet{rqNtOneSideTransAnyExeBOnePct}{0.0}
\rqSet{rqNtOneSideTransAnyExeBTwoNum}{0}
\rqSet{rqNtOneSideTransAnyExeBTwoDen}{24}
\rqSet{rqNtOneSideTransAnyExeBTwoPct}{0.0}
\rqSet{rqNtOneSideTransAnyExePZeroNum}{0}
\rqSet{rqNtOneSideTransAnyExePZeroDen}{0}
\rqSet{rqNtOneSideTransAnyExePZeroPct}{n/a}
\rqSet{rqNtOneSideTransAnyExePOneNum}{0}
\rqSet{rqNtOneSideTransAnyExePOneDen}{0}
\rqSet{rqNtOneSideTransAnyExePOnePct}{n/a}
\rqSet{rqNtOneSideTransOnlyExeAllNum}{0}
\rqSet{rqNtOneSideTransOnlyExeAllDen}{176}
\rqSet{rqNtOneSideTransOnlyExeAllPct}{0.0}
\rqSet{rqNtOneSideTransOnlyExeBOneNum}{0}
\rqSet{rqNtOneSideTransOnlyExeBOneDen}{152}
\rqSet{rqNtOneSideTransOnlyExeBOnePct}{0.0}
\rqSet{rqNtOneSideTransOnlyExeBTwoNum}{0}
\rqSet{rqNtOneSideTransOnlyExeBTwoDen}{24}
\rqSet{rqNtOneSideTransOnlyExeBTwoPct}{0.0}
\rqSet{rqNtOneSideTransOnlyExePZeroNum}{0}
\rqSet{rqNtOneSideTransOnlyExePZeroDen}{0}
\rqSet{rqNtOneSideTransOnlyExePZeroPct}{n/a}
\rqSet{rqNtOneSideTransOnlyExePOneNum}{0}
\rqSet{rqNtOneSideTransOnlyExePOneDen}{0}
\rqSet{rqNtOneSideTransOnlyExePOnePct}{n/a}
\rqSet{rqDivLayerSExeAllNum}{129}
\rqSet{rqDivLayerSExeAllDen}{11813}
\rqSet{rqDivLayerSExeAllPct}{1.1}
\rqSet{rqDivLayerSExeBOneNum}{129}
\rqSet{rqDivLayerSExeBOneDen}{3726}
\rqSet{rqDivLayerSExeBOnePct}{3.5}
\rqSet{rqDivLayerSExeBTwoNum}{0}
\rqSet{rqDivLayerSExeBTwoDen}{7979}
\rqSet{rqDivLayerSExeBTwoPct}{0.0}
\rqSet{rqDivLayerSExePZeroNum}{0}
\rqSet{rqDivLayerSExePZeroDen}{0}
\rqSet{rqDivLayerSExePZeroPct}{n/a}
\rqSet{rqDivLayerSExePOneNum}{0}
\rqSet{rqDivLayerSExePOneDen}{108}
\rqSet{rqDivLayerSExePOnePct}{0.0}
\rqSet{rqDivInputsLayerSExeAllNum}{63}
\rqSet{rqDivInputsLayerSExeAllDen}{2185}
\rqSet{rqDivInputsLayerSExeAllPct}{2.9}
\rqSet{rqDivInputsLayerSExeBOneNum}{63}
\rqSet{rqDivInputsLayerSExeBOneDen}{468}
\rqSet{rqDivInputsLayerSExeBOnePct}{13.5}
\rqSet{rqDivInputsLayerSExeBTwoNum}{0}
\rqSet{rqDivInputsLayerSExeBTwoDen}{1669}
\rqSet{rqDivInputsLayerSExeBTwoPct}{0.0}
\rqSet{rqDivInputsLayerSExePZeroNum}{0}
\rqSet{rqDivInputsLayerSExePZeroDen}{0}
\rqSet{rqDivInputsLayerSExePZeroPct}{n/a}
\rqSet{rqDivInputsLayerSExePOneNum}{0}
\rqSet{rqDivInputsLayerSExePOneDen}{48}
\rqSet{rqDivInputsLayerSExePOnePct}{0.0}
\rqSet{rqDivTransLayerSExeAllNum}{3}
\rqSet{rqDivTransLayerSExeAllDen}{176}
\rqSet{rqDivTransLayerSExeAllPct}{1.7}
\rqSet{rqDivTransLayerSExeBOneNum}{3}
\rqSet{rqDivTransLayerSExeBOneDen}{152}
\rqSet{rqDivTransLayerSExeBOnePct}{2.0}
\rqSet{rqDivTransLayerSExeBTwoNum}{0}
\rqSet{rqDivTransLayerSExeBTwoDen}{24}
\rqSet{rqDivTransLayerSExeBTwoPct}{0.0}
\rqSet{rqDivTransLayerSExePZeroNum}{0}
\rqSet{rqDivTransLayerSExePZeroDen}{0}
\rqSet{rqDivTransLayerSExePZeroPct}{n/a}
\rqSet{rqDivTransLayerSExePOneNum}{0}
\rqSet{rqDivTransLayerSExePOneDen}{0}
\rqSet{rqDivTransLayerSExePOnePct}{n/a}
\rqSet{rqOutOfDomainSExeAllNum}{0}
\rqSet{rqOutOfDomainSExeAllDen}{2328}
\rqSet{rqOutOfDomainSExeAllPct}{0.0}
\rqSet{rqOutOfDomainSExeBOneNum}{0}
\rqSet{rqOutOfDomainSExeBOneDen}{485}
\rqSet{rqOutOfDomainSExeBOnePct}{0.0}
\rqSet{rqOutOfDomainSExeBTwoNum}{0}
\rqSet{rqOutOfDomainSExeBTwoDen}{1765}
\rqSet{rqOutOfDomainSExeBTwoPct}{0.0}
\rqSet{rqOutOfDomainSExePZeroNum}{0}
\rqSet{rqOutOfDomainSExePZeroDen}{30}
\rqSet{rqOutOfDomainSExePZeroPct}{0.0}
\rqSet{rqOutOfDomainSExePOneNum}{0}
\rqSet{rqOutOfDomainSExePOneDen}{48}
\rqSet{rqOutOfDomainSExePOnePct}{0.0}
\rqSet{rqDivLayerMExeAllNum}{4080138}
\rqSet{rqDivLayerMExeAllDen}{20960510}
\rqSet{rqDivLayerMExeAllPct}{19.5}
\rqSet{rqDivLayerMExeBOneNum}{3807346}
\rqSet{rqDivLayerMExeBOneDen}{17531453}
\rqSet{rqDivLayerMExeBOnePct}{21.7}
\rqSet{rqDivLayerMExeBTwoNum}{272792}
\rqSet{rqDivLayerMExeBTwoDen}{3429057}
\rqSet{rqDivLayerMExeBTwoPct}{8.0}
\rqSet{rqDivLayerMExePZeroNum}{0}
\rqSet{rqDivLayerMExePZeroDen}{0}
\rqSet{rqDivLayerMExePZeroPct}{n/a}
\rqSet{rqDivLayerMExePOneNum}{0}
\rqSet{rqDivLayerMExePOneDen}{0}
\rqSet{rqDivLayerMExePOnePct}{n/a}
\rqSet{rqDivInputsLayerMExeAllNum}{1065638}
\rqSet{rqDivInputsLayerMExeAllDen}{2810308}
\rqSet{rqDivInputsLayerMExeAllPct}{37.9}
\rqSet{rqDivInputsLayerMExeBOneNum}{970199}
\rqSet{rqDivInputsLayerMExeBOneDen}{2057945}
\rqSet{rqDivInputsLayerMExeBOnePct}{47.1}
\rqSet{rqDivInputsLayerMExeBTwoNum}{95439}
\rqSet{rqDivInputsLayerMExeBTwoDen}{752363}
\rqSet{rqDivInputsLayerMExeBTwoPct}{12.7}
\rqSet{rqDivInputsLayerMExePZeroNum}{0}
\rqSet{rqDivInputsLayerMExePZeroDen}{0}
\rqSet{rqDivInputsLayerMExePZeroPct}{n/a}
\rqSet{rqDivInputsLayerMExePOneNum}{0}
\rqSet{rqDivInputsLayerMExePOneDen}{0}
\rqSet{rqDivInputsLayerMExePOnePct}{n/a}
\rqSet{rqDivTransLayerMExeAllNum}{172}
\rqSet{rqDivTransLayerMExeAllDen}{176}
\rqSet{rqDivTransLayerMExeAllPct}{97.7}
\rqSet{rqDivTransLayerMExeBOneNum}{148}
\rqSet{rqDivTransLayerMExeBOneDen}{152}
\rqSet{rqDivTransLayerMExeBOnePct}{97.4}
\rqSet{rqDivTransLayerMExeBTwoNum}{24}
\rqSet{rqDivTransLayerMExeBTwoDen}{24}
\rqSet{rqDivTransLayerMExeBTwoPct}{100.0}
\rqSet{rqDivTransLayerMExePZeroNum}{0}
\rqSet{rqDivTransLayerMExePZeroDen}{0}
\rqSet{rqDivTransLayerMExePZeroPct}{n/a}
\rqSet{rqDivTransLayerMExePOneNum}{0}
\rqSet{rqDivTransLayerMExePOneDen}{0}
\rqSet{rqDivTransLayerMExePOnePct}{n/a}
\rqSet{rqOutOfDomainMExeAllNum}{748271}
\rqSet{rqOutOfDomainMExeAllDen}{3688565}
\rqSet{rqOutOfDomainMExeAllPct}{20.3}
\rqSet{rqOutOfDomainMExeBOneNum}{687315}
\rqSet{rqOutOfDomainMExeBOneDen}{2818565}
\rqSet{rqOutOfDomainMExeBOnePct}{24.4}
\rqSet{rqOutOfDomainMExeBTwoNum}{58481}
\rqSet{rqOutOfDomainMExeBTwoDen}{840000}
\rqSet{rqOutOfDomainMExeBTwoPct}{7.0}
\rqSet{rqOutOfDomainMExePZeroNum}{2475}
\rqSet{rqOutOfDomainMExePZeroDen}{30000}
\rqSet{rqOutOfDomainMExePZeroPct}{8.2}
\rqSet{rqOutOfDomainMExePOneNum}{0}
\rqSet{rqOutOfDomainMExePOneDen}{0}
\rqSet{rqOutOfDomainMExePOnePct}{n/a}
\rqSet{rqDivLayerMSExeAllNum}{2573864}
\rqSet{rqDivLayerMSExeAllDen}{14908530}
\rqSet{rqDivLayerMSExeAllPct}{17.3}
\rqSet{rqDivLayerMSExeBOneNum}{2349756}
\rqSet{rqDivLayerMSExeBOneDen}{12209398}
\rqSet{rqDivLayerMSExeBOnePct}{19.2}
\rqSet{rqDivLayerMSExeBTwoNum}{224108}
\rqSet{rqDivLayerMSExeBTwoDen}{2699132}
\rqSet{rqDivLayerMSExeBTwoPct}{8.3}
\rqSet{rqDivLayerMSExePZeroNum}{0}
\rqSet{rqDivLayerMSExePZeroDen}{0}
\rqSet{rqDivLayerMSExePZeroPct}{n/a}
\rqSet{rqDivLayerMSExePOneNum}{0}
\rqSet{rqDivLayerMSExePOneDen}{0}
\rqSet{rqDivLayerMSExePOnePct}{n/a}
\rqSet{rqDivInputsLayerMSExeAllNum}{712306}
\rqSet{rqDivInputsLayerMSExeAllDen}{2028503}
\rqSet{rqDivInputsLayerMSExeAllPct}{35.1}
\rqSet{rqDivInputsLayerMSExeBOneNum}{634487}
\rqSet{rqDivInputsLayerMSExeBOneDen}{1441398}
\rqSet{rqDivInputsLayerMSExeBOnePct}{44.0}
\rqSet{rqDivInputsLayerMSExeBTwoNum}{77819}
\rqSet{rqDivInputsLayerMSExeBTwoDen}{587105}
\rqSet{rqDivInputsLayerMSExeBTwoPct}{13.3}
\rqSet{rqDivInputsLayerMSExePZeroNum}{0}
\rqSet{rqDivInputsLayerMSExePZeroDen}{0}
\rqSet{rqDivInputsLayerMSExePZeroPct}{n/a}
\rqSet{rqDivInputsLayerMSExePOneNum}{0}
\rqSet{rqDivInputsLayerMSExePOneDen}{0}
\rqSet{rqDivInputsLayerMSExePOnePct}{n/a}
\rqSet{rqDivTransLayerMSExeAllNum}{171}
\rqSet{rqDivTransLayerMSExeAllDen}{176}
\rqSet{rqDivTransLayerMSExeAllPct}{97.2}
\rqSet{rqDivTransLayerMSExeBOneNum}{147}
\rqSet{rqDivTransLayerMSExeBOneDen}{152}
\rqSet{rqDivTransLayerMSExeBOnePct}{96.7}
\rqSet{rqDivTransLayerMSExeBTwoNum}{24}
\rqSet{rqDivTransLayerMSExeBTwoDen}{24}
\rqSet{rqDivTransLayerMSExeBTwoPct}{100.0}
\rqSet{rqDivTransLayerMSExePZeroNum}{0}
\rqSet{rqDivTransLayerMSExePZeroDen}{0}
\rqSet{rqDivTransLayerMSExePZeroPct}{n/a}
\rqSet{rqDivTransLayerMSExePOneNum}{0}
\rqSet{rqDivTransLayerMSExePOneDen}{0}
\rqSet{rqDivTransLayerMSExePOnePct}{n/a}
\rqSet{rqOutOfDomainMSExeAllNum}{705759}
\rqSet{rqOutOfDomainMSExeAllDen}{3738007}
\rqSet{rqOutOfDomainMSExeAllPct}{18.9}
\rqSet{rqOutOfDomainMSExeBOneNum}{647849}
\rqSet{rqOutOfDomainMSExeBOneDen}{2868007}
\rqSet{rqOutOfDomainMSExeBOnePct}{22.6}
\rqSet{rqOutOfDomainMSExeBTwoNum}{55639}
\rqSet{rqOutOfDomainMSExeBTwoDen}{840000}
\rqSet{rqOutOfDomainMSExeBTwoPct}{6.6}
\rqSet{rqOutOfDomainMSExePZeroNum}{2271}
\rqSet{rqOutOfDomainMSExePZeroDen}{30000}
\rqSet{rqOutOfDomainMSExePZeroPct}{7.6}
\rqSet{rqOutOfDomainMSExePOneNum}{0}
\rqSet{rqOutOfDomainMSExePOneDen}{0}
\rqSet{rqOutOfDomainMSExePOnePct}{n/a}
\rqSet{rqDivLayerAExeAllNum}{996}
\rqSet{rqDivLayerAExeAllDen}{13506}
\rqSet{rqDivLayerAExeAllPct}{7.4}
\rqSet{rqDivLayerAExeBOneNum}{766}
\rqSet{rqDivLayerAExeBOneDen}{6710}
\rqSet{rqDivLayerAExeBOnePct}{11.4}
\rqSet{rqDivLayerAExeBTwoNum}{230}
\rqSet{rqDivLayerAExeBTwoDen}{6796}
\rqSet{rqDivLayerAExeBTwoPct}{3.4}
\rqSet{rqDivLayerAExePZeroNum}{0}
\rqSet{rqDivLayerAExePZeroDen}{0}
\rqSet{rqDivLayerAExePZeroPct}{n/a}
\rqSet{rqDivLayerAExePOneNum}{0}
\rqSet{rqDivLayerAExePOneDen}{0}
\rqSet{rqDivLayerAExePOnePct}{n/a}
\rqSet{rqDivInputsLayerAExeAllNum}{436}
\rqSet{rqDivInputsLayerAExeAllDen}{2312}
\rqSet{rqDivInputsLayerAExeAllPct}{18.9}
\rqSet{rqDivInputsLayerAExeBOneNum}{297}
\rqSet{rqDivInputsLayerAExeBOneDen}{789}
\rqSet{rqDivInputsLayerAExeBOnePct}{37.6}
\rqSet{rqDivInputsLayerAExeBTwoNum}{139}
\rqSet{rqDivInputsLayerAExeBTwoDen}{1523}
\rqSet{rqDivInputsLayerAExeBTwoPct}{9.1}
\rqSet{rqDivInputsLayerAExePZeroNum}{0}
\rqSet{rqDivInputsLayerAExePZeroDen}{0}
\rqSet{rqDivInputsLayerAExePZeroPct}{n/a}
\rqSet{rqDivInputsLayerAExePOneNum}{0}
\rqSet{rqDivInputsLayerAExePOneDen}{0}
\rqSet{rqDivInputsLayerAExePOnePct}{n/a}
\rqSet{rqDivTransLayerAExeAllNum}{128}
\rqSet{rqDivTransLayerAExeAllDen}{176}
\rqSet{rqDivTransLayerAExeAllPct}{72.7}
\rqSet{rqDivTransLayerAExeBOneNum}{110}
\rqSet{rqDivTransLayerAExeBOneDen}{152}
\rqSet{rqDivTransLayerAExeBOnePct}{72.4}
\rqSet{rqDivTransLayerAExeBTwoNum}{18}
\rqSet{rqDivTransLayerAExeBTwoDen}{24}
\rqSet{rqDivTransLayerAExeBTwoPct}{75.0}
\rqSet{rqDivTransLayerAExePZeroNum}{0}
\rqSet{rqDivTransLayerAExePZeroDen}{0}
\rqSet{rqDivTransLayerAExePZeroPct}{n/a}
\rqSet{rqDivTransLayerAExePOneNum}{0}
\rqSet{rqDivTransLayerAExePOneDen}{0}
\rqSet{rqDivTransLayerAExePOnePct}{n/a}
\rqSet{rqOutOfDomainAExeAllNum}{2343}
\rqSet{rqOutOfDomainAExeAllDen}{7171}
\rqSet{rqOutOfDomainAExeAllPct}{32.7}
\rqSet{rqOutOfDomainAExeBOneNum}{899}
\rqSet{rqOutOfDomainAExeBOneDen}{2277}
\rqSet{rqOutOfDomainAExeBOnePct}{39.5}
\rqSet{rqOutOfDomainAExeBTwoNum}{1432}
\rqSet{rqOutOfDomainAExeBTwoDen}{4817}
\rqSet{rqOutOfDomainAExeBTwoPct}{29.7}
\rqSet{rqOutOfDomainAExePZeroNum}{12}
\rqSet{rqOutOfDomainAExePZeroDen}{77}
\rqSet{rqOutOfDomainAExePZeroPct}{15.6}
\rqSet{rqOutOfDomainAExePOneNum}{0}
\rqSet{rqOutOfDomainAExePOneDen}{0}
\rqSet{rqOutOfDomainAExePOnePct}{n/a}
\rqSet{rqDivLayerGeneratedExeAllNum}{6654998}
\rqSet{rqDivLayerGeneratedExeAllDen}{35882546}
\rqSet{rqDivLayerGeneratedExeAllPct}{18.5}
\rqSet{rqDivLayerGeneratedExeBOneNum}{6157868}
\rqSet{rqDivLayerGeneratedExeBOneDen}{29747561}
\rqSet{rqDivLayerGeneratedExeBOnePct}{20.7}
\rqSet{rqDivLayerGeneratedExeBTwoNum}{497130}
\rqSet{rqDivLayerGeneratedExeBTwoDen}{6134985}
\rqSet{rqDivLayerGeneratedExeBTwoPct}{8.1}
\rqSet{rqDivLayerGeneratedExePZeroNum}{0}
\rqSet{rqDivLayerGeneratedExePZeroDen}{0}
\rqSet{rqDivLayerGeneratedExePZeroPct}{n/a}
\rqSet{rqDivLayerGeneratedExePOneNum}{0}
\rqSet{rqDivLayerGeneratedExePOneDen}{0}
\rqSet{rqDivLayerGeneratedExePOnePct}{n/a}
\rqSet{rqDivInputsLayerGeneratedExeAllNum}{1778380}
\rqSet{rqDivInputsLayerGeneratedExeAllDen}{4841123}
\rqSet{rqDivInputsLayerGeneratedExeAllPct}{36.7}
\rqSet{rqDivInputsLayerGeneratedExeBOneNum}{1604983}
\rqSet{rqDivInputsLayerGeneratedExeBOneDen}{3500132}
\rqSet{rqDivInputsLayerGeneratedExeBOnePct}{45.9}
\rqSet{rqDivInputsLayerGeneratedExeBTwoNum}{173397}
\rqSet{rqDivInputsLayerGeneratedExeBTwoDen}{1340991}
\rqSet{rqDivInputsLayerGeneratedExeBTwoPct}{12.9}
\rqSet{rqDivInputsLayerGeneratedExePZeroNum}{0}
\rqSet{rqDivInputsLayerGeneratedExePZeroDen}{0}
\rqSet{rqDivInputsLayerGeneratedExePZeroPct}{n/a}
\rqSet{rqDivInputsLayerGeneratedExePOneNum}{0}
\rqSet{rqDivInputsLayerGeneratedExePOneDen}{0}
\rqSet{rqDivInputsLayerGeneratedExePOnePct}{n/a}
\rqSet{rqDivTransLayerGeneratedExeAllNum}{176}
\rqSet{rqDivTransLayerGeneratedExeAllDen}{176}
\rqSet{rqDivTransLayerGeneratedExeAllPct}{100.0}
\rqSet{rqDivTransLayerGeneratedExeBOneNum}{152}
\rqSet{rqDivTransLayerGeneratedExeBOneDen}{152}
\rqSet{rqDivTransLayerGeneratedExeBOnePct}{100.0}
\rqSet{rqDivTransLayerGeneratedExeBTwoNum}{24}
\rqSet{rqDivTransLayerGeneratedExeBTwoDen}{24}
\rqSet{rqDivTransLayerGeneratedExeBTwoPct}{100.0}
\rqSet{rqDivTransLayerGeneratedExePZeroNum}{0}
\rqSet{rqDivTransLayerGeneratedExePZeroDen}{0}
\rqSet{rqDivTransLayerGeneratedExePZeroPct}{n/a}
\rqSet{rqDivTransLayerGeneratedExePOneNum}{0}
\rqSet{rqDivTransLayerGeneratedExePOneDen}{0}
\rqSet{rqDivTransLayerGeneratedExePOnePct}{n/a}
\rqSet{rqOutOfDomainGeneratedExeAllNum}{1456373}
\rqSet{rqOutOfDomainGeneratedExeAllDen}{7433743}
\rqSet{rqOutOfDomainGeneratedExeAllPct}{19.6}
\rqSet{rqOutOfDomainGeneratedExeBOneNum}{1336063}
\rqSet{rqOutOfDomainGeneratedExeBOneDen}{5688849}
\rqSet{rqOutOfDomainGeneratedExeBOnePct}{23.5}
\rqSet{rqOutOfDomainGeneratedExeBTwoNum}{115552}
\rqSet{rqOutOfDomainGeneratedExeBTwoDen}{1684817}
\rqSet{rqOutOfDomainGeneratedExeBTwoPct}{6.9}
\rqSet{rqOutOfDomainGeneratedExePZeroNum}{4758}
\rqSet{rqOutOfDomainGeneratedExePZeroDen}{60077}
\rqSet{rqOutOfDomainGeneratedExePZeroPct}{7.9}
\rqSet{rqOutOfDomainGeneratedExePOneNum}{0}
\rqSet{rqOutOfDomainGeneratedExePOneDen}{0}
\rqSet{rqOutOfDomainGeneratedExePOnePct}{n/a}
\rqSet{rqDivSourceSeedExeAllNum}{129}
\rqSet{rqDivSourceSeedExeAllDen}{11813}
\rqSet{rqDivSourceSeedExeAllPct}{1.1}
\rqSet{rqDivSourceSeedExeBOneNum}{129}
\rqSet{rqDivSourceSeedExeBOneDen}{3726}
\rqSet{rqDivSourceSeedExeBOnePct}{3.5}
\rqSet{rqDivSourceSeedExeBTwoNum}{0}
\rqSet{rqDivSourceSeedExeBTwoDen}{7979}
\rqSet{rqDivSourceSeedExeBTwoPct}{0.0}
\rqSet{rqDivSourceSeedExePZeroNum}{0}
\rqSet{rqDivSourceSeedExePZeroDen}{0}
\rqSet{rqDivSourceSeedExePZeroPct}{n/a}
\rqSet{rqDivSourceSeedExePOneNum}{0}
\rqSet{rqDivSourceSeedExePOneDen}{108}
\rqSet{rqDivSourceSeedExePOnePct}{0.0}
\rqSet{rqDivSourceRandomExeAllNum}{3282057}
\rqSet{rqDivSourceRandomExeAllDen}{22372673}
\rqSet{rqDivSourceRandomExeAllPct}{14.7}
\rqSet{rqDivSourceRandomExeBOneNum}{2992669}
\rqSet{rqDivSourceRandomExeBOneDen}{18228275}
\rqSet{rqDivSourceRandomExeBOnePct}{16.4}
\rqSet{rqDivSourceRandomExeBTwoNum}{289388}
\rqSet{rqDivSourceRandomExeBTwoDen}{4144398}
\rqSet{rqDivSourceRandomExeBTwoPct}{7.0}
\rqSet{rqDivSourceRandomExePZeroNum}{0}
\rqSet{rqDivSourceRandomExePZeroDen}{0}
\rqSet{rqDivSourceRandomExePZeroPct}{n/a}
\rqSet{rqDivSourceRandomExePOneNum}{0}
\rqSet{rqDivSourceRandomExePOneDen}{0}
\rqSet{rqDivSourceRandomExePOnePct}{n/a}
\rqSet{rqDivSourceDictExeAllNum}{21159}
\rqSet{rqDivSourceDictExeAllDen}{1338885}
\rqSet{rqDivSourceDictExeAllPct}{1.6}
\rqSet{rqDivSourceDictExeBOneNum}{12860}
\rqSet{rqDivSourceDictExeBOneDen}{297038}
\rqSet{rqDivSourceDictExeBOnePct}{4.3}
\rqSet{rqDivSourceDictExeBTwoNum}{8299}
\rqSet{rqDivSourceDictExeBTwoDen}{1041847}
\rqSet{rqDivSourceDictExeBTwoPct}{0.8}
\rqSet{rqDivSourceDictExePZeroNum}{0}
\rqSet{rqDivSourceDictExePZeroDen}{0}
\rqSet{rqDivSourceDictExePZeroPct}{n/a}
\rqSet{rqDivSourceDictExePOneNum}{0}
\rqSet{rqDivSourceDictExePOneDen}{0}
\rqSet{rqDivSourceDictExePOnePct}{n/a}
\rqSet{rqDivSourceSyntaxExeAllNum}{3350786}
\rqSet{rqDivSourceSyntaxExeAllDen}{12157482}
\rqSet{rqDivSourceSyntaxExeAllPct}{27.6}
\rqSet{rqDivSourceSyntaxExeBOneNum}{3151573}
\rqSet{rqDivSourceSyntaxExeBOneDen}{11215538}
\rqSet{rqDivSourceSyntaxExeBOnePct}{28.1}
\rqSet{rqDivSourceSyntaxExeBTwoNum}{199213}
\rqSet{rqDivSourceSyntaxExeBTwoDen}{941944}
\rqSet{rqDivSourceSyntaxExeBTwoPct}{21.1}
\rqSet{rqDivSourceSyntaxExePZeroNum}{0}
\rqSet{rqDivSourceSyntaxExePZeroDen}{0}
\rqSet{rqDivSourceSyntaxExePZeroPct}{n/a}
\rqSet{rqDivSourceSyntaxExePOneNum}{0}
\rqSet{rqDivSourceSyntaxExePOneDen}{0}
\rqSet{rqDivSourceSyntaxExePOnePct}{n/a}
\rqSet{rqDivSourceAflExeAllNum}{996}
\rqSet{rqDivSourceAflExeAllDen}{13506}
\rqSet{rqDivSourceAflExeAllPct}{7.4}
\rqSet{rqDivSourceAflExeBOneNum}{766}
\rqSet{rqDivSourceAflExeBOneDen}{6710}
\rqSet{rqDivSourceAflExeBOnePct}{11.4}
\rqSet{rqDivSourceAflExeBTwoNum}{230}
\rqSet{rqDivSourceAflExeBTwoDen}{6796}
\rqSet{rqDivSourceAflExeBTwoPct}{3.4}
\rqSet{rqDivSourceAflExePZeroNum}{0}
\rqSet{rqDivSourceAflExePZeroDen}{0}
\rqSet{rqDivSourceAflExePZeroPct}{n/a}
\rqSet{rqDivSourceAflExePOneNum}{0}
\rqSet{rqDivSourceAflExePOneDen}{0}
\rqSet{rqDivSourceAflExePOnePct}{n/a}
\rqSet{rqAttrInputsNoInputFieldsExeAllNum}{0}
\rqSet{rqAttrInputsNoInputFieldsExeAllDen}{35894359}
\rqSet{rqAttrInputsNoInputFieldsExeAllPct}{0.0}
\rqSet{rqAttrInputsNoInputFieldsExeAllKiroNum}{0}
\rqSet{rqAttrInputsNoInputFieldsExeAllKiroDen}{3947743}
\rqSet{rqAttrInputsNoInputFieldsExeAllKiroPct}{0.0}
\rqSet{rqAttrInputsNoInputFieldsExeAllClaudeNum}{0}
\rqSet{rqAttrInputsNoInputFieldsExeAllClaudeDen}{4546925}
\rqSet{rqAttrInputsNoInputFieldsExeAllClaudePct}{0.0}
\rqSet{rqAttrInputsNoInputFieldsExeAllCodexNum}{0}
\rqSet{rqAttrInputsNoInputFieldsExeAllCodexDen}{4138716}
\rqSet{rqAttrInputsNoInputFieldsExeAllCodexPct}{0.0}
\rqSet{rqAttrInputsNoInputFieldsExeAllCTwoRustNum}{0}
\rqSet{rqAttrInputsNoInputFieldsExeAllCTwoRustDen}{4552067}
\rqSet{rqAttrInputsNoInputFieldsExeAllCTwoRustPct}{0.0}
\rqSet{rqAttrInputsNoInputFieldsExeAllLaertesNum}{0}
\rqSet{rqAttrInputsNoInputFieldsExeAllLaertesDen}{4552067}
\rqSet{rqAttrInputsNoInputFieldsExeAllLaertesPct}{0.0}
\rqSet{rqAttrInputsNoInputFieldsExeAllCTwoSaferRustNum}{0}
\rqSet{rqAttrInputsNoInputFieldsExeAllCTwoSaferRustDen}{3214155}
\rqSet{rqAttrInputsNoInputFieldsExeAllCTwoSaferRustPct}{0.0}
\rqSet{rqAttrInputsNoInputFieldsExeAllSmartCTwoRustNum}{0}
\rqSet{rqAttrInputsNoInputFieldsExeAllSmartCTwoRustDen}{3232501}
\rqSet{rqAttrInputsNoInputFieldsExeAllSmartCTwoRustPct}{0.0}
\rqSet{rqAttrInputsNoInputFieldsExeAllKimiNum}{0}
\rqSet{rqAttrInputsNoInputFieldsExeAllKimiDen}{1880086}
\rqSet{rqAttrInputsNoInputFieldsExeAllKimiPct}{0.0}
\rqSet{rqAttrInputsNoInputFieldsExeAllGptFiveNum}{0}
\rqSet{rqAttrInputsNoInputFieldsExeAllGptFiveDen}{2782513}
\rqSet{rqAttrInputsNoInputFieldsExeAllGptFivePct}{0.0}
\rqSet{rqAttrInputsNoInputFieldsExeAllGeminiNum}{0}
\rqSet{rqAttrInputsNoInputFieldsExeAllGeminiDen}{3047586}
\rqSet{rqAttrInputsNoInputFieldsExeAllGeminiPct}{0.0}
\rqSet{rqAttrInputsNoInputFieldsExeBOneNum}{0}
\rqSet{rqAttrInputsNoInputFieldsExeBOneDen}{29751287}
\rqSet{rqAttrInputsNoInputFieldsExeBOnePct}{0.0}
\rqSet{rqAttrInputsNoInputFieldsExeBTwoNum}{0}
\rqSet{rqAttrInputsNoInputFieldsExeBTwoDen}{6142964}
\rqSet{rqAttrInputsNoInputFieldsExeBTwoPct}{0.0}
\rqSet{rqAttrInputsNoInputFieldsExePZeroNum}{0}
\rqSet{rqAttrInputsNoInputFieldsExePZeroDen}{0}
\rqSet{rqAttrInputsNoInputFieldsExePZeroPct}{n/a}
\rqSet{rqAttrInputsNoInputFieldsExePOneNum}{0}
\rqSet{rqAttrInputsNoInputFieldsExePOneDen}{108}
\rqSet{rqAttrInputsNoInputFieldsExePOnePct}{0.0}
\rqSet{rqAttrInputsCBuildFailedExeAllNum}{0}
\rqSet{rqAttrInputsCBuildFailedExeAllDen}{35894359}
\rqSet{rqAttrInputsCBuildFailedExeAllPct}{0.0}
\rqSet{rqAttrInputsCBuildFailedExeAllKiroNum}{0}
\rqSet{rqAttrInputsCBuildFailedExeAllKiroDen}{3947743}
\rqSet{rqAttrInputsCBuildFailedExeAllKiroPct}{0.0}
\rqSet{rqAttrInputsCBuildFailedExeAllClaudeNum}{0}
\rqSet{rqAttrInputsCBuildFailedExeAllClaudeDen}{4546925}
\rqSet{rqAttrInputsCBuildFailedExeAllClaudePct}{0.0}
\rqSet{rqAttrInputsCBuildFailedExeAllCodexNum}{0}
\rqSet{rqAttrInputsCBuildFailedExeAllCodexDen}{4138716}
\rqSet{rqAttrInputsCBuildFailedExeAllCodexPct}{0.0}
\rqSet{rqAttrInputsCBuildFailedExeAllCTwoRustNum}{0}
\rqSet{rqAttrInputsCBuildFailedExeAllCTwoRustDen}{4552067}
\rqSet{rqAttrInputsCBuildFailedExeAllCTwoRustPct}{0.0}
\rqSet{rqAttrInputsCBuildFailedExeAllLaertesNum}{0}
\rqSet{rqAttrInputsCBuildFailedExeAllLaertesDen}{4552067}
\rqSet{rqAttrInputsCBuildFailedExeAllLaertesPct}{0.0}
\rqSet{rqAttrInputsCBuildFailedExeAllCTwoSaferRustNum}{0}
\rqSet{rqAttrInputsCBuildFailedExeAllCTwoSaferRustDen}{3214155}
\rqSet{rqAttrInputsCBuildFailedExeAllCTwoSaferRustPct}{0.0}
\rqSet{rqAttrInputsCBuildFailedExeAllSmartCTwoRustNum}{0}
\rqSet{rqAttrInputsCBuildFailedExeAllSmartCTwoRustDen}{3232501}
\rqSet{rqAttrInputsCBuildFailedExeAllSmartCTwoRustPct}{0.0}
\rqSet{rqAttrInputsCBuildFailedExeAllKimiNum}{0}
\rqSet{rqAttrInputsCBuildFailedExeAllKimiDen}{1880086}
\rqSet{rqAttrInputsCBuildFailedExeAllKimiPct}{0.0}
\rqSet{rqAttrInputsCBuildFailedExeAllGptFiveNum}{0}
\rqSet{rqAttrInputsCBuildFailedExeAllGptFiveDen}{2782513}
\rqSet{rqAttrInputsCBuildFailedExeAllGptFivePct}{0.0}
\rqSet{rqAttrInputsCBuildFailedExeAllGeminiNum}{0}
\rqSet{rqAttrInputsCBuildFailedExeAllGeminiDen}{3047586}
\rqSet{rqAttrInputsCBuildFailedExeAllGeminiPct}{0.0}
\rqSet{rqAttrInputsCBuildFailedExeBOneNum}{0}
\rqSet{rqAttrInputsCBuildFailedExeBOneDen}{29751287}
\rqSet{rqAttrInputsCBuildFailedExeBOnePct}{0.0}
\rqSet{rqAttrInputsCBuildFailedExeBTwoNum}{0}
\rqSet{rqAttrInputsCBuildFailedExeBTwoDen}{6142964}
\rqSet{rqAttrInputsCBuildFailedExeBTwoPct}{0.0}
\rqSet{rqAttrInputsCBuildFailedExePZeroNum}{0}
\rqSet{rqAttrInputsCBuildFailedExePZeroDen}{0}
\rqSet{rqAttrInputsCBuildFailedExePZeroPct}{n/a}
\rqSet{rqAttrInputsCBuildFailedExePOneNum}{0}
\rqSet{rqAttrInputsCBuildFailedExePOneDen}{108}
\rqSet{rqAttrInputsCBuildFailedExePOnePct}{0.0}
\rqSet{rqAttrInputsRustBuildFailedExeAllNum}{0}
\rqSet{rqAttrInputsRustBuildFailedExeAllDen}{35894359}
\rqSet{rqAttrInputsRustBuildFailedExeAllPct}{0.0}
\rqSet{rqAttrInputsRustBuildFailedExeAllKiroNum}{0}
\rqSet{rqAttrInputsRustBuildFailedExeAllKiroDen}{3947743}
\rqSet{rqAttrInputsRustBuildFailedExeAllKiroPct}{0.0}
\rqSet{rqAttrInputsRustBuildFailedExeAllClaudeNum}{0}
\rqSet{rqAttrInputsRustBuildFailedExeAllClaudeDen}{4546925}
\rqSet{rqAttrInputsRustBuildFailedExeAllClaudePct}{0.0}
\rqSet{rqAttrInputsRustBuildFailedExeAllCodexNum}{0}
\rqSet{rqAttrInputsRustBuildFailedExeAllCodexDen}{4138716}
\rqSet{rqAttrInputsRustBuildFailedExeAllCodexPct}{0.0}
\rqSet{rqAttrInputsRustBuildFailedExeAllCTwoRustNum}{0}
\rqSet{rqAttrInputsRustBuildFailedExeAllCTwoRustDen}{4552067}
\rqSet{rqAttrInputsRustBuildFailedExeAllCTwoRustPct}{0.0}
\rqSet{rqAttrInputsRustBuildFailedExeAllLaertesNum}{0}
\rqSet{rqAttrInputsRustBuildFailedExeAllLaertesDen}{4552067}
\rqSet{rqAttrInputsRustBuildFailedExeAllLaertesPct}{0.0}
\rqSet{rqAttrInputsRustBuildFailedExeAllCTwoSaferRustNum}{0}
\rqSet{rqAttrInputsRustBuildFailedExeAllCTwoSaferRustDen}{3214155}
\rqSet{rqAttrInputsRustBuildFailedExeAllCTwoSaferRustPct}{0.0}
\rqSet{rqAttrInputsRustBuildFailedExeAllSmartCTwoRustNum}{0}
\rqSet{rqAttrInputsRustBuildFailedExeAllSmartCTwoRustDen}{3232501}
\rqSet{rqAttrInputsRustBuildFailedExeAllSmartCTwoRustPct}{0.0}
\rqSet{rqAttrInputsRustBuildFailedExeAllKimiNum}{0}
\rqSet{rqAttrInputsRustBuildFailedExeAllKimiDen}{1880086}
\rqSet{rqAttrInputsRustBuildFailedExeAllKimiPct}{0.0}
\rqSet{rqAttrInputsRustBuildFailedExeAllGptFiveNum}{0}
\rqSet{rqAttrInputsRustBuildFailedExeAllGptFiveDen}{2782513}
\rqSet{rqAttrInputsRustBuildFailedExeAllGptFivePct}{0.0}
\rqSet{rqAttrInputsRustBuildFailedExeAllGeminiNum}{0}
\rqSet{rqAttrInputsRustBuildFailedExeAllGeminiDen}{3047586}
\rqSet{rqAttrInputsRustBuildFailedExeAllGeminiPct}{0.0}
\rqSet{rqAttrInputsRustBuildFailedExeBOneNum}{0}
\rqSet{rqAttrInputsRustBuildFailedExeBOneDen}{29751287}
\rqSet{rqAttrInputsRustBuildFailedExeBOnePct}{0.0}
\rqSet{rqAttrInputsRustBuildFailedExeBTwoNum}{0}
\rqSet{rqAttrInputsRustBuildFailedExeBTwoDen}{6142964}
\rqSet{rqAttrInputsRustBuildFailedExeBTwoPct}{0.0}
\rqSet{rqAttrInputsRustBuildFailedExePZeroNum}{0}
\rqSet{rqAttrInputsRustBuildFailedExePZeroDen}{0}
\rqSet{rqAttrInputsRustBuildFailedExePZeroPct}{n/a}
\rqSet{rqAttrInputsRustBuildFailedExePOneNum}{0}
\rqSet{rqAttrInputsRustBuildFailedExePOneDen}{108}
\rqSet{rqAttrInputsRustBuildFailedExePOnePct}{0.0}
\rqSet{rqAttrInputsHarnessErrorExeAllNum}{0}
\rqSet{rqAttrInputsHarnessErrorExeAllDen}{35894359}
\rqSet{rqAttrInputsHarnessErrorExeAllPct}{0.0}
\rqSet{rqAttrInputsHarnessErrorExeAllKiroNum}{0}
\rqSet{rqAttrInputsHarnessErrorExeAllKiroDen}{3947743}
\rqSet{rqAttrInputsHarnessErrorExeAllKiroPct}{0.0}
\rqSet{rqAttrInputsHarnessErrorExeAllClaudeNum}{0}
\rqSet{rqAttrInputsHarnessErrorExeAllClaudeDen}{4546925}
\rqSet{rqAttrInputsHarnessErrorExeAllClaudePct}{0.0}
\rqSet{rqAttrInputsHarnessErrorExeAllCodexNum}{0}
\rqSet{rqAttrInputsHarnessErrorExeAllCodexDen}{4138716}
\rqSet{rqAttrInputsHarnessErrorExeAllCodexPct}{0.0}
\rqSet{rqAttrInputsHarnessErrorExeAllCTwoRustNum}{0}
\rqSet{rqAttrInputsHarnessErrorExeAllCTwoRustDen}{4552067}
\rqSet{rqAttrInputsHarnessErrorExeAllCTwoRustPct}{0.0}
\rqSet{rqAttrInputsHarnessErrorExeAllLaertesNum}{0}
\rqSet{rqAttrInputsHarnessErrorExeAllLaertesDen}{4552067}
\rqSet{rqAttrInputsHarnessErrorExeAllLaertesPct}{0.0}
\rqSet{rqAttrInputsHarnessErrorExeAllCTwoSaferRustNum}{0}
\rqSet{rqAttrInputsHarnessErrorExeAllCTwoSaferRustDen}{3214155}
\rqSet{rqAttrInputsHarnessErrorExeAllCTwoSaferRustPct}{0.0}
\rqSet{rqAttrInputsHarnessErrorExeAllSmartCTwoRustNum}{0}
\rqSet{rqAttrInputsHarnessErrorExeAllSmartCTwoRustDen}{3232501}
\rqSet{rqAttrInputsHarnessErrorExeAllSmartCTwoRustPct}{0.0}
\rqSet{rqAttrInputsHarnessErrorExeAllKimiNum}{0}
\rqSet{rqAttrInputsHarnessErrorExeAllKimiDen}{1880086}
\rqSet{rqAttrInputsHarnessErrorExeAllKimiPct}{0.0}
\rqSet{rqAttrInputsHarnessErrorExeAllGptFiveNum}{0}
\rqSet{rqAttrInputsHarnessErrorExeAllGptFiveDen}{2782513}
\rqSet{rqAttrInputsHarnessErrorExeAllGptFivePct}{0.0}
\rqSet{rqAttrInputsHarnessErrorExeAllGeminiNum}{0}
\rqSet{rqAttrInputsHarnessErrorExeAllGeminiDen}{3047586}
\rqSet{rqAttrInputsHarnessErrorExeAllGeminiPct}{0.0}
\rqSet{rqAttrInputsHarnessErrorExeBOneNum}{0}
\rqSet{rqAttrInputsHarnessErrorExeBOneDen}{29751287}
\rqSet{rqAttrInputsHarnessErrorExeBOnePct}{0.0}
\rqSet{rqAttrInputsHarnessErrorExeBTwoNum}{0}
\rqSet{rqAttrInputsHarnessErrorExeBTwoDen}{6142964}
\rqSet{rqAttrInputsHarnessErrorExeBTwoPct}{0.0}
\rqSet{rqAttrInputsHarnessErrorExePZeroNum}{0}
\rqSet{rqAttrInputsHarnessErrorExePZeroDen}{0}
\rqSet{rqAttrInputsHarnessErrorExePZeroPct}{n/a}
\rqSet{rqAttrInputsHarnessErrorExePOneNum}{0}
\rqSet{rqAttrInputsHarnessErrorExePOneDen}{108}
\rqSet{rqAttrInputsHarnessErrorExePOnePct}{0.0}
\rqSet{rqAttrInputsNoInputsExeAllNum}{0}
\rqSet{rqAttrInputsNoInputsExeAllDen}{35894359}
\rqSet{rqAttrInputsNoInputsExeAllPct}{0.0}
\rqSet{rqAttrInputsNoInputsExeAllKiroNum}{0}
\rqSet{rqAttrInputsNoInputsExeAllKiroDen}{3947743}
\rqSet{rqAttrInputsNoInputsExeAllKiroPct}{0.0}
\rqSet{rqAttrInputsNoInputsExeAllClaudeNum}{0}
\rqSet{rqAttrInputsNoInputsExeAllClaudeDen}{4546925}
\rqSet{rqAttrInputsNoInputsExeAllClaudePct}{0.0}
\rqSet{rqAttrInputsNoInputsExeAllCodexNum}{0}
\rqSet{rqAttrInputsNoInputsExeAllCodexDen}{4138716}
\rqSet{rqAttrInputsNoInputsExeAllCodexPct}{0.0}
\rqSet{rqAttrInputsNoInputsExeAllCTwoRustNum}{0}
\rqSet{rqAttrInputsNoInputsExeAllCTwoRustDen}{4552067}
\rqSet{rqAttrInputsNoInputsExeAllCTwoRustPct}{0.0}
\rqSet{rqAttrInputsNoInputsExeAllLaertesNum}{0}
\rqSet{rqAttrInputsNoInputsExeAllLaertesDen}{4552067}
\rqSet{rqAttrInputsNoInputsExeAllLaertesPct}{0.0}
\rqSet{rqAttrInputsNoInputsExeAllCTwoSaferRustNum}{0}
\rqSet{rqAttrInputsNoInputsExeAllCTwoSaferRustDen}{3214155}
\rqSet{rqAttrInputsNoInputsExeAllCTwoSaferRustPct}{0.0}
\rqSet{rqAttrInputsNoInputsExeAllSmartCTwoRustNum}{0}
\rqSet{rqAttrInputsNoInputsExeAllSmartCTwoRustDen}{3232501}
\rqSet{rqAttrInputsNoInputsExeAllSmartCTwoRustPct}{0.0}
\rqSet{rqAttrInputsNoInputsExeAllKimiNum}{0}
\rqSet{rqAttrInputsNoInputsExeAllKimiDen}{1880086}
\rqSet{rqAttrInputsNoInputsExeAllKimiPct}{0.0}
\rqSet{rqAttrInputsNoInputsExeAllGptFiveNum}{0}
\rqSet{rqAttrInputsNoInputsExeAllGptFiveDen}{2782513}
\rqSet{rqAttrInputsNoInputsExeAllGptFivePct}{0.0}
\rqSet{rqAttrInputsNoInputsExeAllGeminiNum}{0}
\rqSet{rqAttrInputsNoInputsExeAllGeminiDen}{3047586}
\rqSet{rqAttrInputsNoInputsExeAllGeminiPct}{0.0}
\rqSet{rqAttrInputsNoInputsExeBOneNum}{0}
\rqSet{rqAttrInputsNoInputsExeBOneDen}{29751287}
\rqSet{rqAttrInputsNoInputsExeBOnePct}{0.0}
\rqSet{rqAttrInputsNoInputsExeBTwoNum}{0}
\rqSet{rqAttrInputsNoInputsExeBTwoDen}{6142964}
\rqSet{rqAttrInputsNoInputsExeBTwoPct}{0.0}
\rqSet{rqAttrInputsNoInputsExePZeroNum}{0}
\rqSet{rqAttrInputsNoInputsExePZeroDen}{0}
\rqSet{rqAttrInputsNoInputsExePZeroPct}{n/a}
\rqSet{rqAttrInputsNoInputsExePOneNum}{0}
\rqSet{rqAttrInputsNoInputsExePOneDen}{108}
\rqSet{rqAttrInputsNoInputsExePOnePct}{0.0}
\rqSet{rqAttrInputsIncompleteExeAllNum}{0}
\rqSet{rqAttrInputsIncompleteExeAllDen}{35894359}
\rqSet{rqAttrInputsIncompleteExeAllPct}{0.0}
\rqSet{rqAttrInputsIncompleteExeAllKiroNum}{0}
\rqSet{rqAttrInputsIncompleteExeAllKiroDen}{3947743}
\rqSet{rqAttrInputsIncompleteExeAllKiroPct}{0.0}
\rqSet{rqAttrInputsIncompleteExeAllClaudeNum}{0}
\rqSet{rqAttrInputsIncompleteExeAllClaudeDen}{4546925}
\rqSet{rqAttrInputsIncompleteExeAllClaudePct}{0.0}
\rqSet{rqAttrInputsIncompleteExeAllCodexNum}{0}
\rqSet{rqAttrInputsIncompleteExeAllCodexDen}{4138716}
\rqSet{rqAttrInputsIncompleteExeAllCodexPct}{0.0}
\rqSet{rqAttrInputsIncompleteExeAllCTwoRustNum}{0}
\rqSet{rqAttrInputsIncompleteExeAllCTwoRustDen}{4552067}
\rqSet{rqAttrInputsIncompleteExeAllCTwoRustPct}{0.0}
\rqSet{rqAttrInputsIncompleteExeAllLaertesNum}{0}
\rqSet{rqAttrInputsIncompleteExeAllLaertesDen}{4552067}
\rqSet{rqAttrInputsIncompleteExeAllLaertesPct}{0.0}
\rqSet{rqAttrInputsIncompleteExeAllCTwoSaferRustNum}{0}
\rqSet{rqAttrInputsIncompleteExeAllCTwoSaferRustDen}{3214155}
\rqSet{rqAttrInputsIncompleteExeAllCTwoSaferRustPct}{0.0}
\rqSet{rqAttrInputsIncompleteExeAllSmartCTwoRustNum}{0}
\rqSet{rqAttrInputsIncompleteExeAllSmartCTwoRustDen}{3232501}
\rqSet{rqAttrInputsIncompleteExeAllSmartCTwoRustPct}{0.0}
\rqSet{rqAttrInputsIncompleteExeAllKimiNum}{0}
\rqSet{rqAttrInputsIncompleteExeAllKimiDen}{1880086}
\rqSet{rqAttrInputsIncompleteExeAllKimiPct}{0.0}
\rqSet{rqAttrInputsIncompleteExeAllGptFiveNum}{0}
\rqSet{rqAttrInputsIncompleteExeAllGptFiveDen}{2782513}
\rqSet{rqAttrInputsIncompleteExeAllGptFivePct}{0.0}
\rqSet{rqAttrInputsIncompleteExeAllGeminiNum}{0}
\rqSet{rqAttrInputsIncompleteExeAllGeminiDen}{3047586}
\rqSet{rqAttrInputsIncompleteExeAllGeminiPct}{0.0}
\rqSet{rqAttrInputsIncompleteExeBOneNum}{0}
\rqSet{rqAttrInputsIncompleteExeBOneDen}{29751287}
\rqSet{rqAttrInputsIncompleteExeBOnePct}{0.0}
\rqSet{rqAttrInputsIncompleteExeBTwoNum}{0}
\rqSet{rqAttrInputsIncompleteExeBTwoDen}{6142964}
\rqSet{rqAttrInputsIncompleteExeBTwoPct}{0.0}
\rqSet{rqAttrInputsIncompleteExePZeroNum}{0}
\rqSet{rqAttrInputsIncompleteExePZeroDen}{0}
\rqSet{rqAttrInputsIncompleteExePZeroPct}{n/a}
\rqSet{rqAttrInputsIncompleteExePOneNum}{0}
\rqSet{rqAttrInputsIncompleteExePOneDen}{108}
\rqSet{rqAttrInputsIncompleteExePOnePct}{0.0}
\rqSet{rqAttrInputsSanInfraExeAllNum}{0}
\rqSet{rqAttrInputsSanInfraExeAllDen}{35894359}
\rqSet{rqAttrInputsSanInfraExeAllPct}{0.0}
\rqSet{rqAttrInputsSanInfraExeAllKiroNum}{0}
\rqSet{rqAttrInputsSanInfraExeAllKiroDen}{3947743}
\rqSet{rqAttrInputsSanInfraExeAllKiroPct}{0.0}
\rqSet{rqAttrInputsSanInfraExeAllClaudeNum}{0}
\rqSet{rqAttrInputsSanInfraExeAllClaudeDen}{4546925}
\rqSet{rqAttrInputsSanInfraExeAllClaudePct}{0.0}
\rqSet{rqAttrInputsSanInfraExeAllCodexNum}{0}
\rqSet{rqAttrInputsSanInfraExeAllCodexDen}{4138716}
\rqSet{rqAttrInputsSanInfraExeAllCodexPct}{0.0}
\rqSet{rqAttrInputsSanInfraExeAllCTwoRustNum}{0}
\rqSet{rqAttrInputsSanInfraExeAllCTwoRustDen}{4552067}
\rqSet{rqAttrInputsSanInfraExeAllCTwoRustPct}{0.0}
\rqSet{rqAttrInputsSanInfraExeAllLaertesNum}{0}
\rqSet{rqAttrInputsSanInfraExeAllLaertesDen}{4552067}
\rqSet{rqAttrInputsSanInfraExeAllLaertesPct}{0.0}
\rqSet{rqAttrInputsSanInfraExeAllCTwoSaferRustNum}{0}
\rqSet{rqAttrInputsSanInfraExeAllCTwoSaferRustDen}{3214155}
\rqSet{rqAttrInputsSanInfraExeAllCTwoSaferRustPct}{0.0}
\rqSet{rqAttrInputsSanInfraExeAllSmartCTwoRustNum}{0}
\rqSet{rqAttrInputsSanInfraExeAllSmartCTwoRustDen}{3232501}
\rqSet{rqAttrInputsSanInfraExeAllSmartCTwoRustPct}{0.0}
\rqSet{rqAttrInputsSanInfraExeAllKimiNum}{0}
\rqSet{rqAttrInputsSanInfraExeAllKimiDen}{1880086}
\rqSet{rqAttrInputsSanInfraExeAllKimiPct}{0.0}
\rqSet{rqAttrInputsSanInfraExeAllGptFiveNum}{0}
\rqSet{rqAttrInputsSanInfraExeAllGptFiveDen}{2782513}
\rqSet{rqAttrInputsSanInfraExeAllGptFivePct}{0.0}
\rqSet{rqAttrInputsSanInfraExeAllGeminiNum}{0}
\rqSet{rqAttrInputsSanInfraExeAllGeminiDen}{3047586}
\rqSet{rqAttrInputsSanInfraExeAllGeminiPct}{0.0}
\rqSet{rqAttrInputsSanInfraExeBOneNum}{0}
\rqSet{rqAttrInputsSanInfraExeBOneDen}{29751287}
\rqSet{rqAttrInputsSanInfraExeBOnePct}{0.0}
\rqSet{rqAttrInputsSanInfraExeBTwoNum}{0}
\rqSet{rqAttrInputsSanInfraExeBTwoDen}{6142964}
\rqSet{rqAttrInputsSanInfraExeBTwoPct}{0.0}
\rqSet{rqAttrInputsSanInfraExePZeroNum}{0}
\rqSet{rqAttrInputsSanInfraExePZeroDen}{0}
\rqSet{rqAttrInputsSanInfraExePZeroPct}{n/a}
\rqSet{rqAttrInputsSanInfraExePOneNum}{0}
\rqSet{rqAttrInputsSanInfraExePOneDen}{108}
\rqSet{rqAttrInputsSanInfraExePOnePct}{0.0}
\rqSet{rqAttrInputsMissingSeedExeAllNum}{??}
\rqSet{rqAttrInputsMissingSeedExeAllDen}{35894359}
\rqSet{rqAttrInputsMissingSeedExeAllPct}{??}
\rqSet{rqAttrInputsMissingSeedExeAllKiroNum}{??}
\rqSet{rqAttrInputsMissingSeedExeAllKiroDen}{3947743}
\rqSet{rqAttrInputsMissingSeedExeAllKiroPct}{??}
\rqSet{rqAttrInputsMissingSeedExeAllClaudeNum}{??}
\rqSet{rqAttrInputsMissingSeedExeAllClaudeDen}{4546925}
\rqSet{rqAttrInputsMissingSeedExeAllClaudePct}{??}
\rqSet{rqAttrInputsMissingSeedExeAllCodexNum}{??}
\rqSet{rqAttrInputsMissingSeedExeAllCodexDen}{4138716}
\rqSet{rqAttrInputsMissingSeedExeAllCodexPct}{??}
\rqSet{rqAttrInputsMissingSeedExeAllCTwoRustNum}{??}
\rqSet{rqAttrInputsMissingSeedExeAllCTwoRustDen}{4552067}
\rqSet{rqAttrInputsMissingSeedExeAllCTwoRustPct}{??}
\rqSet{rqAttrInputsMissingSeedExeAllLaertesNum}{??}
\rqSet{rqAttrInputsMissingSeedExeAllLaertesDen}{4552067}
\rqSet{rqAttrInputsMissingSeedExeAllLaertesPct}{??}
\rqSet{rqAttrInputsMissingSeedExeAllCTwoSaferRustNum}{??}
\rqSet{rqAttrInputsMissingSeedExeAllCTwoSaferRustDen}{3214155}
\rqSet{rqAttrInputsMissingSeedExeAllCTwoSaferRustPct}{??}
\rqSet{rqAttrInputsMissingSeedExeAllSmartCTwoRustNum}{??}
\rqSet{rqAttrInputsMissingSeedExeAllSmartCTwoRustDen}{3232501}
\rqSet{rqAttrInputsMissingSeedExeAllSmartCTwoRustPct}{??}
\rqSet{rqAttrInputsMissingSeedExeAllKimiNum}{??}
\rqSet{rqAttrInputsMissingSeedExeAllKimiDen}{1880086}
\rqSet{rqAttrInputsMissingSeedExeAllKimiPct}{??}
\rqSet{rqAttrInputsMissingSeedExeAllGptFiveNum}{??}
\rqSet{rqAttrInputsMissingSeedExeAllGptFiveDen}{2782513}
\rqSet{rqAttrInputsMissingSeedExeAllGptFivePct}{??}
\rqSet{rqAttrInputsMissingSeedExeAllGeminiNum}{??}
\rqSet{rqAttrInputsMissingSeedExeAllGeminiDen}{3047586}
\rqSet{rqAttrInputsMissingSeedExeAllGeminiPct}{??}
\rqSet{rqAttrInputsMissingSeedExeBOneNum}{??}
\rqSet{rqAttrInputsMissingSeedExeBOneDen}{29751287}
\rqSet{rqAttrInputsMissingSeedExeBOnePct}{??}
\rqSet{rqAttrInputsMissingSeedExeBTwoNum}{??}
\rqSet{rqAttrInputsMissingSeedExeBTwoDen}{6142964}
\rqSet{rqAttrInputsMissingSeedExeBTwoPct}{??}
\rqSet{rqAttrInputsMissingSeedExePZeroNum}{??}
\rqSet{rqAttrInputsMissingSeedExePZeroDen}{0}
\rqSet{rqAttrInputsMissingSeedExePZeroPct}{??}
\rqSet{rqAttrInputsMissingSeedExePOneNum}{0}
\rqSet{rqAttrInputsMissingSeedExePOneDen}{108}
\rqSet{rqAttrInputsMissingSeedExePOnePct}{0.0}
\rqSet{rqAttrInputsMissingRecordExeAllNum}{??}
\rqSet{rqAttrInputsMissingRecordExeAllDen}{35894359}
\rqSet{rqAttrInputsMissingRecordExeAllPct}{??}
\rqSet{rqAttrInputsMissingRecordExeAllKiroNum}{??}
\rqSet{rqAttrInputsMissingRecordExeAllKiroDen}{3947743}
\rqSet{rqAttrInputsMissingRecordExeAllKiroPct}{??}
\rqSet{rqAttrInputsMissingRecordExeAllClaudeNum}{??}
\rqSet{rqAttrInputsMissingRecordExeAllClaudeDen}{4546925}
\rqSet{rqAttrInputsMissingRecordExeAllClaudePct}{??}
\rqSet{rqAttrInputsMissingRecordExeAllCodexNum}{??}
\rqSet{rqAttrInputsMissingRecordExeAllCodexDen}{4138716}
\rqSet{rqAttrInputsMissingRecordExeAllCodexPct}{??}
\rqSet{rqAttrInputsMissingRecordExeAllCTwoRustNum}{??}
\rqSet{rqAttrInputsMissingRecordExeAllCTwoRustDen}{4552067}
\rqSet{rqAttrInputsMissingRecordExeAllCTwoRustPct}{??}
\rqSet{rqAttrInputsMissingRecordExeAllLaertesNum}{??}
\rqSet{rqAttrInputsMissingRecordExeAllLaertesDen}{4552067}
\rqSet{rqAttrInputsMissingRecordExeAllLaertesPct}{??}
\rqSet{rqAttrInputsMissingRecordExeAllCTwoSaferRustNum}{??}
\rqSet{rqAttrInputsMissingRecordExeAllCTwoSaferRustDen}{3214155}
\rqSet{rqAttrInputsMissingRecordExeAllCTwoSaferRustPct}{??}
\rqSet{rqAttrInputsMissingRecordExeAllSmartCTwoRustNum}{??}
\rqSet{rqAttrInputsMissingRecordExeAllSmartCTwoRustDen}{3232501}
\rqSet{rqAttrInputsMissingRecordExeAllSmartCTwoRustPct}{??}
\rqSet{rqAttrInputsMissingRecordExeAllKimiNum}{??}
\rqSet{rqAttrInputsMissingRecordExeAllKimiDen}{1880086}
\rqSet{rqAttrInputsMissingRecordExeAllKimiPct}{??}
\rqSet{rqAttrInputsMissingRecordExeAllGptFiveNum}{??}
\rqSet{rqAttrInputsMissingRecordExeAllGptFiveDen}{2782513}
\rqSet{rqAttrInputsMissingRecordExeAllGptFivePct}{??}
\rqSet{rqAttrInputsMissingRecordExeAllGeminiNum}{??}
\rqSet{rqAttrInputsMissingRecordExeAllGeminiDen}{3047586}
\rqSet{rqAttrInputsMissingRecordExeAllGeminiPct}{??}
\rqSet{rqAttrInputsMissingRecordExeBOneNum}{??}
\rqSet{rqAttrInputsMissingRecordExeBOneDen}{29751287}
\rqSet{rqAttrInputsMissingRecordExeBOnePct}{??}
\rqSet{rqAttrInputsMissingRecordExeBTwoNum}{0}
\rqSet{rqAttrInputsMissingRecordExeBTwoDen}{6142964}
\rqSet{rqAttrInputsMissingRecordExeBTwoPct}{0.0}
\rqSet{rqAttrInputsMissingRecordExePZeroNum}{0}
\rqSet{rqAttrInputsMissingRecordExePZeroDen}{0}
\rqSet{rqAttrInputsMissingRecordExePZeroPct}{n/a}
\rqSet{rqAttrInputsMissingRecordExePOneNum}{0}
\rqSet{rqAttrInputsMissingRecordExePOneDen}{108}
\rqSet{rqAttrInputsMissingRecordExePOnePct}{0.0}
\rqSet{rqAttrInputsUnknownExeAllNum}{72}
\rqSet{rqAttrInputsUnknownExeAllDen}{548}
\rqSet{rqAttrInputsUnknownExeAllPct}{13.1}
\rqSet{rqAttrInputsUnknownExeAllKiroNum}{8}
\rqSet{rqAttrInputsUnknownExeAllKiroDen}{96}
\rqSet{rqAttrInputsUnknownExeAllKiroPct}{8.3}
\rqSet{rqAttrInputsUnknownExeAllClaudeNum}{8}
\rqSet{rqAttrInputsUnknownExeAllClaudeDen}{101}
\rqSet{rqAttrInputsUnknownExeAllClaudePct}{7.9}
\rqSet{rqAttrInputsUnknownExeAllCodexNum}{8}
\rqSet{rqAttrInputsUnknownExeAllCodexDen}{62}
\rqSet{rqAttrInputsUnknownExeAllCodexPct}{12.9}
\rqSet{rqAttrInputsUnknownExeAllCTwoRustNum}{9}
\rqSet{rqAttrInputsUnknownExeAllCTwoRustDen}{53}
\rqSet{rqAttrInputsUnknownExeAllCTwoRustPct}{17.0}
\rqSet{rqAttrInputsUnknownExeAllLaertesNum}{9}
\rqSet{rqAttrInputsUnknownExeAllLaertesDen}{53}
\rqSet{rqAttrInputsUnknownExeAllLaertesPct}{17.0}
\rqSet{rqAttrInputsUnknownExeAllCTwoSaferRustNum}{7}
\rqSet{rqAttrInputsUnknownExeAllCTwoSaferRustDen}{39}
\rqSet{rqAttrInputsUnknownExeAllCTwoSaferRustPct}{17.9}
\rqSet{rqAttrInputsUnknownExeAllSmartCTwoRustNum}{7}
\rqSet{rqAttrInputsUnknownExeAllSmartCTwoRustDen}{40}
\rqSet{rqAttrInputsUnknownExeAllSmartCTwoRustPct}{17.5}
\rqSet{rqAttrInputsUnknownExeAllKimiNum}{3}
\rqSet{rqAttrInputsUnknownExeAllKimiDen}{27}
\rqSet{rqAttrInputsUnknownExeAllKimiPct}{11.1}
\rqSet{rqAttrInputsUnknownExeAllGptFiveNum}{6}
\rqSet{rqAttrInputsUnknownExeAllGptFiveDen}{36}
\rqSet{rqAttrInputsUnknownExeAllGptFivePct}{16.7}
\rqSet{rqAttrInputsUnknownExeAllGeminiNum}{7}
\rqSet{rqAttrInputsUnknownExeAllGeminiDen}{41}
\rqSet{rqAttrInputsUnknownExeAllGeminiPct}{17.1}
\rqSet{rqAttrInputsUnknownExeBOneNum}{63}
\rqSet{rqAttrInputsUnknownExeBOneDen}{378}
\rqSet{rqAttrInputsUnknownExeBOnePct}{16.7}
\rqSet{rqAttrInputsUnknownExeBTwoNum}{3}
\rqSet{rqAttrInputsUnknownExeBTwoDen}{56}
\rqSet{rqAttrInputsUnknownExeBTwoPct}{5.4}
\rqSet{rqAttrInputsUnknownExePZeroNum}{6}
\rqSet{rqAttrInputsUnknownExePZeroDen}{6}
\rqSet{rqAttrInputsUnknownExePZeroPct}{100.0}
\rqSet{rqAttrInputsUnknownExePOneNum}{0}
\rqSet{rqAttrInputsUnknownExePOneDen}{108}
\rqSet{rqAttrInputsUnknownExePOnePct}{0.0}
\rqSet{rqAttrInputsSanUnresolvedExeAllNum}{0}
\rqSet{rqAttrInputsSanUnresolvedExeAllDen}{0}
\rqSet{rqAttrInputsSanUnresolvedExeBOneNum}{0}
\rqSet{rqAttrInputsSanUnresolvedExeBOneDen}{0}
\rqSet{rqAttrInputsSanUnresolvedExeBTwoNum}{0}
\rqSet{rqAttrInputsSanUnresolvedExeBTwoDen}{0}
\rqSet{rqAttrInputsSanUnresolvedExePZeroNum}{0}
\rqSet{rqAttrInputsSanUnresolvedExePZeroDen}{0}
\rqSet{rqAttrInputsSanUnresolvedExePOneNum}{0}
\rqSet{rqAttrInputsSanUnresolvedExePOneDen}{0}
\rqSet{rqNonRepAccOracleDriftExeAllNum}{0}
\rqSet{rqNonRepAccOracleDriftExeAllDen}{3}
\rqSet{rqNonRepAccOracleDriftExeAllPct}{0.0}
\rqSet{rqNonRepAccOracleDriftExeAllKiroNum}{0}
\rqSet{rqNonRepAccOracleDriftExeAllKiroDen}{0}
\rqSet{rqNonRepAccOracleDriftExeAllKiroPct}{n/a}
\rqSet{rqNonRepAccOracleDriftExeAllClaudeNum}{0}
\rqSet{rqNonRepAccOracleDriftExeAllClaudeDen}{0}
\rqSet{rqNonRepAccOracleDriftExeAllClaudePct}{n/a}
\rqSet{rqNonRepAccOracleDriftExeAllCodexNum}{0}
\rqSet{rqNonRepAccOracleDriftExeAllCodexDen}{0}
\rqSet{rqNonRepAccOracleDriftExeAllCodexPct}{n/a}
\rqSet{rqNonRepAccOracleDriftExeAllCTwoRustNum}{0}
\rqSet{rqNonRepAccOracleDriftExeAllCTwoRustDen}{0}
\rqSet{rqNonRepAccOracleDriftExeAllCTwoRustPct}{n/a}
\rqSet{rqNonRepAccOracleDriftExeAllLaertesNum}{0}
\rqSet{rqNonRepAccOracleDriftExeAllLaertesDen}{0}
\rqSet{rqNonRepAccOracleDriftExeAllLaertesPct}{n/a}
\rqSet{rqNonRepAccOracleDriftExeAllCTwoSaferRustNum}{0}
\rqSet{rqNonRepAccOracleDriftExeAllCTwoSaferRustDen}{0}
\rqSet{rqNonRepAccOracleDriftExeAllCTwoSaferRustPct}{n/a}
\rqSet{rqNonRepAccOracleDriftExeAllSmartCTwoRustNum}{0}
\rqSet{rqNonRepAccOracleDriftExeAllSmartCTwoRustDen}{0}
\rqSet{rqNonRepAccOracleDriftExeAllSmartCTwoRustPct}{n/a}
\rqSet{rqNonRepAccOracleDriftExeAllKimiNum}{0}
\rqSet{rqNonRepAccOracleDriftExeAllKimiDen}{1}
\rqSet{rqNonRepAccOracleDriftExeAllKimiPct}{0.0}
\rqSet{rqNonRepAccOracleDriftExeAllGptFiveNum}{0}
\rqSet{rqNonRepAccOracleDriftExeAllGptFiveDen}{1}
\rqSet{rqNonRepAccOracleDriftExeAllGptFivePct}{0.0}
\rqSet{rqNonRepAccOracleDriftExeAllGeminiNum}{0}
\rqSet{rqNonRepAccOracleDriftExeAllGeminiDen}{1}
\rqSet{rqNonRepAccOracleDriftExeAllGeminiPct}{0.0}
\rqSet{rqNonRepAccOracleDriftExeBOneNum}{0}
\rqSet{rqNonRepAccOracleDriftExeBOneDen}{0}
\rqSet{rqNonRepAccOracleDriftExeBOnePct}{n/a}
\rqSet{rqNonRepAccOracleDriftExeBTwoNum}{0}
\rqSet{rqNonRepAccOracleDriftExeBTwoDen}{3}
\rqSet{rqNonRepAccOracleDriftExeBTwoPct}{0.0}
\rqSet{rqNonRepAccOracleDriftExePZeroNum}{0}
\rqSet{rqNonRepAccOracleDriftExePZeroDen}{0}
\rqSet{rqNonRepAccOracleDriftExePZeroPct}{n/a}
\rqSet{rqNonRepAccOracleDriftExePOneNum}{0}
\rqSet{rqNonRepAccOracleDriftExePOneDen}{0}
\rqSet{rqNonRepAccOracleDriftExePOnePct}{n/a}
\rqSet{rqNonRepAccTranslationExeAllNum}{3}
\rqSet{rqNonRepAccTranslationExeAllDen}{3}
\rqSet{rqNonRepAccTranslationExeAllPct}{100.0}
\rqSet{rqNonRepAccTranslationExeAllKiroNum}{0}
\rqSet{rqNonRepAccTranslationExeAllKiroDen}{0}
\rqSet{rqNonRepAccTranslationExeAllKiroPct}{n/a}
\rqSet{rqNonRepAccTranslationExeAllClaudeNum}{0}
\rqSet{rqNonRepAccTranslationExeAllClaudeDen}{0}
\rqSet{rqNonRepAccTranslationExeAllClaudePct}{n/a}
\rqSet{rqNonRepAccTranslationExeAllCodexNum}{0}
\rqSet{rqNonRepAccTranslationExeAllCodexDen}{0}
\rqSet{rqNonRepAccTranslationExeAllCodexPct}{n/a}
\rqSet{rqNonRepAccTranslationExeAllCTwoRustNum}{0}
\rqSet{rqNonRepAccTranslationExeAllCTwoRustDen}{0}
\rqSet{rqNonRepAccTranslationExeAllCTwoRustPct}{n/a}
\rqSet{rqNonRepAccTranslationExeAllLaertesNum}{0}
\rqSet{rqNonRepAccTranslationExeAllLaertesDen}{0}
\rqSet{rqNonRepAccTranslationExeAllLaertesPct}{n/a}
\rqSet{rqNonRepAccTranslationExeAllCTwoSaferRustNum}{0}
\rqSet{rqNonRepAccTranslationExeAllCTwoSaferRustDen}{0}
\rqSet{rqNonRepAccTranslationExeAllCTwoSaferRustPct}{n/a}
\rqSet{rqNonRepAccTranslationExeAllSmartCTwoRustNum}{0}
\rqSet{rqNonRepAccTranslationExeAllSmartCTwoRustDen}{0}
\rqSet{rqNonRepAccTranslationExeAllSmartCTwoRustPct}{n/a}
\rqSet{rqNonRepAccTranslationExeAllKimiNum}{1}
\rqSet{rqNonRepAccTranslationExeAllKimiDen}{1}
\rqSet{rqNonRepAccTranslationExeAllKimiPct}{100.0}
\rqSet{rqNonRepAccTranslationExeAllGptFiveNum}{1}
\rqSet{rqNonRepAccTranslationExeAllGptFiveDen}{1}
\rqSet{rqNonRepAccTranslationExeAllGptFivePct}{100.0}
\rqSet{rqNonRepAccTranslationExeAllGeminiNum}{1}
\rqSet{rqNonRepAccTranslationExeAllGeminiDen}{1}
\rqSet{rqNonRepAccTranslationExeAllGeminiPct}{100.0}
\rqSet{rqNonRepAccTranslationExeBOneNum}{0}
\rqSet{rqNonRepAccTranslationExeBOneDen}{0}
\rqSet{rqNonRepAccTranslationExeBOnePct}{n/a}
\rqSet{rqNonRepAccTranslationExeBTwoNum}{3}
\rqSet{rqNonRepAccTranslationExeBTwoDen}{3}
\rqSet{rqNonRepAccTranslationExeBTwoPct}{100.0}
\rqSet{rqNonRepAccTranslationExePZeroNum}{0}
\rqSet{rqNonRepAccTranslationExePZeroDen}{0}
\rqSet{rqNonRepAccTranslationExePZeroPct}{n/a}
\rqSet{rqNonRepAccTranslationExePOneNum}{0}
\rqSet{rqNonRepAccTranslationExePOneDen}{0}
\rqSet{rqNonRepAccTranslationExePOnePct}{n/a}
\rqSet{rqNonRepAccUndetectedUbExeAllNum}{0}
\rqSet{rqNonRepAccUndetectedUbExeAllDen}{3}
\rqSet{rqNonRepAccUndetectedUbExeAllPct}{0.0}
\rqSet{rqNonRepAccUndetectedUbExeAllKiroNum}{0}
\rqSet{rqNonRepAccUndetectedUbExeAllKiroDen}{0}
\rqSet{rqNonRepAccUndetectedUbExeAllKiroPct}{n/a}
\rqSet{rqNonRepAccUndetectedUbExeAllClaudeNum}{0}
\rqSet{rqNonRepAccUndetectedUbExeAllClaudeDen}{0}
\rqSet{rqNonRepAccUndetectedUbExeAllClaudePct}{n/a}
\rqSet{rqNonRepAccUndetectedUbExeAllCodexNum}{0}
\rqSet{rqNonRepAccUndetectedUbExeAllCodexDen}{0}
\rqSet{rqNonRepAccUndetectedUbExeAllCodexPct}{n/a}
\rqSet{rqNonRepAccUndetectedUbExeAllCTwoRustNum}{0}
\rqSet{rqNonRepAccUndetectedUbExeAllCTwoRustDen}{0}
\rqSet{rqNonRepAccUndetectedUbExeAllCTwoRustPct}{n/a}
\rqSet{rqNonRepAccUndetectedUbExeAllLaertesNum}{0}
\rqSet{rqNonRepAccUndetectedUbExeAllLaertesDen}{0}
\rqSet{rqNonRepAccUndetectedUbExeAllLaertesPct}{n/a}
\rqSet{rqNonRepAccUndetectedUbExeAllCTwoSaferRustNum}{0}
\rqSet{rqNonRepAccUndetectedUbExeAllCTwoSaferRustDen}{0}
\rqSet{rqNonRepAccUndetectedUbExeAllCTwoSaferRustPct}{n/a}
\rqSet{rqNonRepAccUndetectedUbExeAllSmartCTwoRustNum}{0}
\rqSet{rqNonRepAccUndetectedUbExeAllSmartCTwoRustDen}{0}
\rqSet{rqNonRepAccUndetectedUbExeAllSmartCTwoRustPct}{n/a}
\rqSet{rqNonRepAccUndetectedUbExeAllKimiNum}{0}
\rqSet{rqNonRepAccUndetectedUbExeAllKimiDen}{1}
\rqSet{rqNonRepAccUndetectedUbExeAllKimiPct}{0.0}
\rqSet{rqNonRepAccUndetectedUbExeAllGptFiveNum}{0}
\rqSet{rqNonRepAccUndetectedUbExeAllGptFiveDen}{1}
\rqSet{rqNonRepAccUndetectedUbExeAllGptFivePct}{0.0}
\rqSet{rqNonRepAccUndetectedUbExeAllGeminiNum}{0}
\rqSet{rqNonRepAccUndetectedUbExeAllGeminiDen}{1}
\rqSet{rqNonRepAccUndetectedUbExeAllGeminiPct}{0.0}
\rqSet{rqNonRepAccUndetectedUbExeBOneNum}{0}
\rqSet{rqNonRepAccUndetectedUbExeBOneDen}{0}
\rqSet{rqNonRepAccUndetectedUbExeBOnePct}{n/a}
\rqSet{rqNonRepAccUndetectedUbExeBTwoNum}{0}
\rqSet{rqNonRepAccUndetectedUbExeBTwoDen}{3}
\rqSet{rqNonRepAccUndetectedUbExeBTwoPct}{0.0}
\rqSet{rqNonRepAccUndetectedUbExePZeroNum}{0}
\rqSet{rqNonRepAccUndetectedUbExePZeroDen}{0}
\rqSet{rqNonRepAccUndetectedUbExePZeroPct}{n/a}
\rqSet{rqNonRepAccUndetectedUbExePOneNum}{0}
\rqSet{rqNonRepAccUndetectedUbExePOneDen}{0}
\rqSet{rqNonRepAccUndetectedUbExePOnePct}{n/a}
\rqSet{rqNonRepAccUnknownExeAllNum}{0}
\rqSet{rqNonRepAccUnknownExeAllDen}{3}
\rqSet{rqNonRepAccUnknownExeAllPct}{0.0}
\rqSet{rqNonRepAccUnknownExeAllKiroNum}{0}
\rqSet{rqNonRepAccUnknownExeAllKiroDen}{0}
\rqSet{rqNonRepAccUnknownExeAllKiroPct}{n/a}
\rqSet{rqNonRepAccUnknownExeAllClaudeNum}{0}
\rqSet{rqNonRepAccUnknownExeAllClaudeDen}{0}
\rqSet{rqNonRepAccUnknownExeAllClaudePct}{n/a}
\rqSet{rqNonRepAccUnknownExeAllCodexNum}{0}
\rqSet{rqNonRepAccUnknownExeAllCodexDen}{0}
\rqSet{rqNonRepAccUnknownExeAllCodexPct}{n/a}
\rqSet{rqNonRepAccUnknownExeAllCTwoRustNum}{0}
\rqSet{rqNonRepAccUnknownExeAllCTwoRustDen}{0}
\rqSet{rqNonRepAccUnknownExeAllCTwoRustPct}{n/a}
\rqSet{rqNonRepAccUnknownExeAllLaertesNum}{0}
\rqSet{rqNonRepAccUnknownExeAllLaertesDen}{0}
\rqSet{rqNonRepAccUnknownExeAllLaertesPct}{n/a}
\rqSet{rqNonRepAccUnknownExeAllCTwoSaferRustNum}{0}
\rqSet{rqNonRepAccUnknownExeAllCTwoSaferRustDen}{0}
\rqSet{rqNonRepAccUnknownExeAllCTwoSaferRustPct}{n/a}
\rqSet{rqNonRepAccUnknownExeAllSmartCTwoRustNum}{0}
\rqSet{rqNonRepAccUnknownExeAllSmartCTwoRustDen}{0}
\rqSet{rqNonRepAccUnknownExeAllSmartCTwoRustPct}{n/a}
\rqSet{rqNonRepAccUnknownExeAllKimiNum}{0}
\rqSet{rqNonRepAccUnknownExeAllKimiDen}{1}
\rqSet{rqNonRepAccUnknownExeAllKimiPct}{0.0}
\rqSet{rqNonRepAccUnknownExeAllGptFiveNum}{0}
\rqSet{rqNonRepAccUnknownExeAllGptFiveDen}{1}
\rqSet{rqNonRepAccUnknownExeAllGptFivePct}{0.0}
\rqSet{rqNonRepAccUnknownExeAllGeminiNum}{0}
\rqSet{rqNonRepAccUnknownExeAllGeminiDen}{1}
\rqSet{rqNonRepAccUnknownExeAllGeminiPct}{0.0}
\rqSet{rqNonRepAccUnknownExeBOneNum}{0}
\rqSet{rqNonRepAccUnknownExeBOneDen}{0}
\rqSet{rqNonRepAccUnknownExeBOnePct}{n/a}
\rqSet{rqNonRepAccUnknownExeBTwoNum}{0}
\rqSet{rqNonRepAccUnknownExeBTwoDen}{3}
\rqSet{rqNonRepAccUnknownExeBTwoPct}{0.0}
\rqSet{rqNonRepAccUnknownExePZeroNum}{0}
\rqSet{rqNonRepAccUnknownExePZeroDen}{0}
\rqSet{rqNonRepAccUnknownExePZeroPct}{n/a}
\rqSet{rqNonRepAccUnknownExePOneNum}{0}
\rqSet{rqNonRepAccUnknownExePOneDen}{0}
\rqSet{rqNonRepAccUnknownExePOnePct}{n/a}
\rqSet{rqNonRepAccUnclassifiedExeAllNum}{0}
\rqSet{rqNonRepAccUnclassifiedExeAllDen}{3}
\rqSet{rqNonRepAccUnclassifiedExeAllPct}{0.0}
\rqSet{rqNonRepAccUnclassifiedExeAllKiroNum}{0}
\rqSet{rqNonRepAccUnclassifiedExeAllKiroDen}{0}
\rqSet{rqNonRepAccUnclassifiedExeAllKiroPct}{n/a}
\rqSet{rqNonRepAccUnclassifiedExeAllClaudeNum}{0}
\rqSet{rqNonRepAccUnclassifiedExeAllClaudeDen}{0}
\rqSet{rqNonRepAccUnclassifiedExeAllClaudePct}{n/a}
\rqSet{rqNonRepAccUnclassifiedExeAllCodexNum}{0}
\rqSet{rqNonRepAccUnclassifiedExeAllCodexDen}{0}
\rqSet{rqNonRepAccUnclassifiedExeAllCodexPct}{n/a}
\rqSet{rqNonRepAccUnclassifiedExeAllCTwoRustNum}{0}
\rqSet{rqNonRepAccUnclassifiedExeAllCTwoRustDen}{0}
\rqSet{rqNonRepAccUnclassifiedExeAllCTwoRustPct}{n/a}
\rqSet{rqNonRepAccUnclassifiedExeAllLaertesNum}{0}
\rqSet{rqNonRepAccUnclassifiedExeAllLaertesDen}{0}
\rqSet{rqNonRepAccUnclassifiedExeAllLaertesPct}{n/a}
\rqSet{rqNonRepAccUnclassifiedExeAllCTwoSaferRustNum}{0}
\rqSet{rqNonRepAccUnclassifiedExeAllCTwoSaferRustDen}{0}
\rqSet{rqNonRepAccUnclassifiedExeAllCTwoSaferRustPct}{n/a}
\rqSet{rqNonRepAccUnclassifiedExeAllSmartCTwoRustNum}{0}
\rqSet{rqNonRepAccUnclassifiedExeAllSmartCTwoRustDen}{0}
\rqSet{rqNonRepAccUnclassifiedExeAllSmartCTwoRustPct}{n/a}
\rqSet{rqNonRepAccUnclassifiedExeAllKimiNum}{0}
\rqSet{rqNonRepAccUnclassifiedExeAllKimiDen}{1}
\rqSet{rqNonRepAccUnclassifiedExeAllKimiPct}{0.0}
\rqSet{rqNonRepAccUnclassifiedExeAllGptFiveNum}{0}
\rqSet{rqNonRepAccUnclassifiedExeAllGptFiveDen}{1}
\rqSet{rqNonRepAccUnclassifiedExeAllGptFivePct}{0.0}
\rqSet{rqNonRepAccUnclassifiedExeAllGeminiNum}{0}
\rqSet{rqNonRepAccUnclassifiedExeAllGeminiDen}{1}
\rqSet{rqNonRepAccUnclassifiedExeAllGeminiPct}{0.0}
\rqSet{rqNonRepAccUnclassifiedExeBOneNum}{0}
\rqSet{rqNonRepAccUnclassifiedExeBOneDen}{0}
\rqSet{rqNonRepAccUnclassifiedExeBOnePct}{n/a}
\rqSet{rqNonRepAccUnclassifiedExeBTwoNum}{0}
\rqSet{rqNonRepAccUnclassifiedExeBTwoDen}{3}
\rqSet{rqNonRepAccUnclassifiedExeBTwoPct}{0.0}
\rqSet{rqNonRepAccUnclassifiedExePZeroNum}{0}
\rqSet{rqNonRepAccUnclassifiedExePZeroDen}{0}
\rqSet{rqNonRepAccUnclassifiedExePZeroPct}{n/a}
\rqSet{rqNonRepAccUnclassifiedExePOneNum}{0}
\rqSet{rqNonRepAccUnclassifiedExePOneDen}{0}
\rqSet{rqNonRepAccUnclassifiedExePOnePct}{n/a}
\rqSet{rqNonRepRejOracleDriftExeAllNum}{0}
\rqSet{rqNonRepRejOracleDriftExeAllDen}{0}
\rqSet{rqNonRepRejOracleDriftExeAllPct}{n/a}
\rqSet{rqNonRepRejOracleDriftExeAllKiroNum}{0}
\rqSet{rqNonRepRejOracleDriftExeAllKiroDen}{0}
\rqSet{rqNonRepRejOracleDriftExeAllKiroPct}{n/a}
\rqSet{rqNonRepRejOracleDriftExeAllClaudeNum}{0}
\rqSet{rqNonRepRejOracleDriftExeAllClaudeDen}{0}
\rqSet{rqNonRepRejOracleDriftExeAllClaudePct}{n/a}
\rqSet{rqNonRepRejOracleDriftExeAllCodexNum}{0}
\rqSet{rqNonRepRejOracleDriftExeAllCodexDen}{0}
\rqSet{rqNonRepRejOracleDriftExeAllCodexPct}{n/a}
\rqSet{rqNonRepRejOracleDriftExeAllCTwoRustNum}{0}
\rqSet{rqNonRepRejOracleDriftExeAllCTwoRustDen}{0}
\rqSet{rqNonRepRejOracleDriftExeAllCTwoRustPct}{n/a}
\rqSet{rqNonRepRejOracleDriftExeAllLaertesNum}{0}
\rqSet{rqNonRepRejOracleDriftExeAllLaertesDen}{0}
\rqSet{rqNonRepRejOracleDriftExeAllLaertesPct}{n/a}
\rqSet{rqNonRepRejOracleDriftExeAllCTwoSaferRustNum}{0}
\rqSet{rqNonRepRejOracleDriftExeAllCTwoSaferRustDen}{0}
\rqSet{rqNonRepRejOracleDriftExeAllCTwoSaferRustPct}{n/a}
\rqSet{rqNonRepRejOracleDriftExeAllSmartCTwoRustNum}{0}
\rqSet{rqNonRepRejOracleDriftExeAllSmartCTwoRustDen}{0}
\rqSet{rqNonRepRejOracleDriftExeAllSmartCTwoRustPct}{n/a}
\rqSet{rqNonRepRejOracleDriftExeAllKimiNum}{0}
\rqSet{rqNonRepRejOracleDriftExeAllKimiDen}{0}
\rqSet{rqNonRepRejOracleDriftExeAllKimiPct}{n/a}
\rqSet{rqNonRepRejOracleDriftExeAllGptFiveNum}{0}
\rqSet{rqNonRepRejOracleDriftExeAllGptFiveDen}{0}
\rqSet{rqNonRepRejOracleDriftExeAllGptFivePct}{n/a}
\rqSet{rqNonRepRejOracleDriftExeAllGeminiNum}{0}
\rqSet{rqNonRepRejOracleDriftExeAllGeminiDen}{0}
\rqSet{rqNonRepRejOracleDriftExeAllGeminiPct}{n/a}
\rqSet{rqNonRepRejOracleDriftExeBOneNum}{0}
\rqSet{rqNonRepRejOracleDriftExeBOneDen}{0}
\rqSet{rqNonRepRejOracleDriftExeBOnePct}{n/a}
\rqSet{rqNonRepRejOracleDriftExeBTwoNum}{0}
\rqSet{rqNonRepRejOracleDriftExeBTwoDen}{0}
\rqSet{rqNonRepRejOracleDriftExeBTwoPct}{n/a}
\rqSet{rqNonRepRejOracleDriftExePZeroNum}{0}
\rqSet{rqNonRepRejOracleDriftExePZeroDen}{0}
\rqSet{rqNonRepRejOracleDriftExePZeroPct}{n/a}
\rqSet{rqNonRepRejOracleDriftExePOneNum}{0}
\rqSet{rqNonRepRejOracleDriftExePOneDen}{0}
\rqSet{rqNonRepRejOracleDriftExePOnePct}{n/a}
\rqSet{rqNonRepRejUndetectedUbExeAllNum}{0}
\rqSet{rqNonRepRejUndetectedUbExeAllDen}{0}
\rqSet{rqNonRepRejUndetectedUbExeAllPct}{n/a}
\rqSet{rqNonRepRejUndetectedUbExeAllKiroNum}{0}
\rqSet{rqNonRepRejUndetectedUbExeAllKiroDen}{0}
\rqSet{rqNonRepRejUndetectedUbExeAllKiroPct}{n/a}
\rqSet{rqNonRepRejUndetectedUbExeAllClaudeNum}{0}
\rqSet{rqNonRepRejUndetectedUbExeAllClaudeDen}{0}
\rqSet{rqNonRepRejUndetectedUbExeAllClaudePct}{n/a}
\rqSet{rqNonRepRejUndetectedUbExeAllCodexNum}{0}
\rqSet{rqNonRepRejUndetectedUbExeAllCodexDen}{0}
\rqSet{rqNonRepRejUndetectedUbExeAllCodexPct}{n/a}
\rqSet{rqNonRepRejUndetectedUbExeAllCTwoRustNum}{0}
\rqSet{rqNonRepRejUndetectedUbExeAllCTwoRustDen}{0}
\rqSet{rqNonRepRejUndetectedUbExeAllCTwoRustPct}{n/a}
\rqSet{rqNonRepRejUndetectedUbExeAllLaertesNum}{0}
\rqSet{rqNonRepRejUndetectedUbExeAllLaertesDen}{0}
\rqSet{rqNonRepRejUndetectedUbExeAllLaertesPct}{n/a}
\rqSet{rqNonRepRejUndetectedUbExeAllCTwoSaferRustNum}{0}
\rqSet{rqNonRepRejUndetectedUbExeAllCTwoSaferRustDen}{0}
\rqSet{rqNonRepRejUndetectedUbExeAllCTwoSaferRustPct}{n/a}
\rqSet{rqNonRepRejUndetectedUbExeAllSmartCTwoRustNum}{0}
\rqSet{rqNonRepRejUndetectedUbExeAllSmartCTwoRustDen}{0}
\rqSet{rqNonRepRejUndetectedUbExeAllSmartCTwoRustPct}{n/a}
\rqSet{rqNonRepRejUndetectedUbExeAllKimiNum}{0}
\rqSet{rqNonRepRejUndetectedUbExeAllKimiDen}{0}
\rqSet{rqNonRepRejUndetectedUbExeAllKimiPct}{n/a}
\rqSet{rqNonRepRejUndetectedUbExeAllGptFiveNum}{0}
\rqSet{rqNonRepRejUndetectedUbExeAllGptFiveDen}{0}
\rqSet{rqNonRepRejUndetectedUbExeAllGptFivePct}{n/a}
\rqSet{rqNonRepRejUndetectedUbExeAllGeminiNum}{0}
\rqSet{rqNonRepRejUndetectedUbExeAllGeminiDen}{0}
\rqSet{rqNonRepRejUndetectedUbExeAllGeminiPct}{n/a}
\rqSet{rqNonRepRejUndetectedUbExeBOneNum}{0}
\rqSet{rqNonRepRejUndetectedUbExeBOneDen}{0}
\rqSet{rqNonRepRejUndetectedUbExeBOnePct}{n/a}
\rqSet{rqNonRepRejUndetectedUbExeBTwoNum}{0}
\rqSet{rqNonRepRejUndetectedUbExeBTwoDen}{0}
\rqSet{rqNonRepRejUndetectedUbExeBTwoPct}{n/a}
\rqSet{rqNonRepRejUndetectedUbExePZeroNum}{0}
\rqSet{rqNonRepRejUndetectedUbExePZeroDen}{0}
\rqSet{rqNonRepRejUndetectedUbExePZeroPct}{n/a}
\rqSet{rqNonRepRejUndetectedUbExePOneNum}{0}
\rqSet{rqNonRepRejUndetectedUbExePOneDen}{0}
\rqSet{rqNonRepRejUndetectedUbExePOnePct}{n/a}
\rqSet{rqNonRepRejUnknownExeAllNum}{0}
\rqSet{rqNonRepRejUnknownExeAllDen}{0}
\rqSet{rqNonRepRejUnknownExeAllPct}{n/a}
\rqSet{rqNonRepRejUnknownExeAllKiroNum}{0}
\rqSet{rqNonRepRejUnknownExeAllKiroDen}{0}
\rqSet{rqNonRepRejUnknownExeAllKiroPct}{n/a}
\rqSet{rqNonRepRejUnknownExeAllClaudeNum}{0}
\rqSet{rqNonRepRejUnknownExeAllClaudeDen}{0}
\rqSet{rqNonRepRejUnknownExeAllClaudePct}{n/a}
\rqSet{rqNonRepRejUnknownExeAllCodexNum}{0}
\rqSet{rqNonRepRejUnknownExeAllCodexDen}{0}
\rqSet{rqNonRepRejUnknownExeAllCodexPct}{n/a}
\rqSet{rqNonRepRejUnknownExeAllCTwoRustNum}{0}
\rqSet{rqNonRepRejUnknownExeAllCTwoRustDen}{0}
\rqSet{rqNonRepRejUnknownExeAllCTwoRustPct}{n/a}
\rqSet{rqNonRepRejUnknownExeAllLaertesNum}{0}
\rqSet{rqNonRepRejUnknownExeAllLaertesDen}{0}
\rqSet{rqNonRepRejUnknownExeAllLaertesPct}{n/a}
\rqSet{rqNonRepRejUnknownExeAllCTwoSaferRustNum}{0}
\rqSet{rqNonRepRejUnknownExeAllCTwoSaferRustDen}{0}
\rqSet{rqNonRepRejUnknownExeAllCTwoSaferRustPct}{n/a}
\rqSet{rqNonRepRejUnknownExeAllSmartCTwoRustNum}{0}
\rqSet{rqNonRepRejUnknownExeAllSmartCTwoRustDen}{0}
\rqSet{rqNonRepRejUnknownExeAllSmartCTwoRustPct}{n/a}
\rqSet{rqNonRepRejUnknownExeAllKimiNum}{0}
\rqSet{rqNonRepRejUnknownExeAllKimiDen}{0}
\rqSet{rqNonRepRejUnknownExeAllKimiPct}{n/a}
\rqSet{rqNonRepRejUnknownExeAllGptFiveNum}{0}
\rqSet{rqNonRepRejUnknownExeAllGptFiveDen}{0}
\rqSet{rqNonRepRejUnknownExeAllGptFivePct}{n/a}
\rqSet{rqNonRepRejUnknownExeAllGeminiNum}{0}
\rqSet{rqNonRepRejUnknownExeAllGeminiDen}{0}
\rqSet{rqNonRepRejUnknownExeAllGeminiPct}{n/a}
\rqSet{rqNonRepRejUnknownExeBOneNum}{0}
\rqSet{rqNonRepRejUnknownExeBOneDen}{0}
\rqSet{rqNonRepRejUnknownExeBOnePct}{n/a}
\rqSet{rqNonRepRejUnknownExeBTwoNum}{0}
\rqSet{rqNonRepRejUnknownExeBTwoDen}{0}
\rqSet{rqNonRepRejUnknownExeBTwoPct}{n/a}
\rqSet{rqNonRepRejUnknownExePZeroNum}{0}
\rqSet{rqNonRepRejUnknownExePZeroDen}{0}
\rqSet{rqNonRepRejUnknownExePZeroPct}{n/a}
\rqSet{rqNonRepRejUnknownExePOneNum}{0}
\rqSet{rqNonRepRejUnknownExePOneDen}{0}
\rqSet{rqNonRepRejUnknownExePOnePct}{n/a}
\rqSet{rqNonRepRejUnclassifiedExeAllNum}{0}
\rqSet{rqNonRepRejUnclassifiedExeAllDen}{0}
\rqSet{rqNonRepRejUnclassifiedExeAllPct}{n/a}
\rqSet{rqNonRepRejUnclassifiedExeAllKiroNum}{0}
\rqSet{rqNonRepRejUnclassifiedExeAllKiroDen}{0}
\rqSet{rqNonRepRejUnclassifiedExeAllKiroPct}{n/a}
\rqSet{rqNonRepRejUnclassifiedExeAllClaudeNum}{0}
\rqSet{rqNonRepRejUnclassifiedExeAllClaudeDen}{0}
\rqSet{rqNonRepRejUnclassifiedExeAllClaudePct}{n/a}
\rqSet{rqNonRepRejUnclassifiedExeAllCodexNum}{0}
\rqSet{rqNonRepRejUnclassifiedExeAllCodexDen}{0}
\rqSet{rqNonRepRejUnclassifiedExeAllCodexPct}{n/a}
\rqSet{rqNonRepRejUnclassifiedExeAllCTwoRustNum}{0}
\rqSet{rqNonRepRejUnclassifiedExeAllCTwoRustDen}{0}
\rqSet{rqNonRepRejUnclassifiedExeAllCTwoRustPct}{n/a}
\rqSet{rqNonRepRejUnclassifiedExeAllLaertesNum}{0}
\rqSet{rqNonRepRejUnclassifiedExeAllLaertesDen}{0}
\rqSet{rqNonRepRejUnclassifiedExeAllLaertesPct}{n/a}
\rqSet{rqNonRepRejUnclassifiedExeAllCTwoSaferRustNum}{0}
\rqSet{rqNonRepRejUnclassifiedExeAllCTwoSaferRustDen}{0}
\rqSet{rqNonRepRejUnclassifiedExeAllCTwoSaferRustPct}{n/a}
\rqSet{rqNonRepRejUnclassifiedExeAllSmartCTwoRustNum}{0}
\rqSet{rqNonRepRejUnclassifiedExeAllSmartCTwoRustDen}{0}
\rqSet{rqNonRepRejUnclassifiedExeAllSmartCTwoRustPct}{n/a}
\rqSet{rqNonRepRejUnclassifiedExeAllKimiNum}{0}
\rqSet{rqNonRepRejUnclassifiedExeAllKimiDen}{0}
\rqSet{rqNonRepRejUnclassifiedExeAllKimiPct}{n/a}
\rqSet{rqNonRepRejUnclassifiedExeAllGptFiveNum}{0}
\rqSet{rqNonRepRejUnclassifiedExeAllGptFiveDen}{0}
\rqSet{rqNonRepRejUnclassifiedExeAllGptFivePct}{n/a}
\rqSet{rqNonRepRejUnclassifiedExeAllGeminiNum}{0}
\rqSet{rqNonRepRejUnclassifiedExeAllGeminiDen}{0}
\rqSet{rqNonRepRejUnclassifiedExeAllGeminiPct}{n/a}
\rqSet{rqNonRepRejUnclassifiedExeBOneNum}{0}
\rqSet{rqNonRepRejUnclassifiedExeBOneDen}{0}
\rqSet{rqNonRepRejUnclassifiedExeBOnePct}{n/a}
\rqSet{rqNonRepRejUnclassifiedExeBTwoNum}{0}
\rqSet{rqNonRepRejUnclassifiedExeBTwoDen}{0}
\rqSet{rqNonRepRejUnclassifiedExeBTwoPct}{n/a}
\rqSet{rqNonRepRejUnclassifiedExePZeroNum}{0}
\rqSet{rqNonRepRejUnclassifiedExePZeroDen}{0}
\rqSet{rqNonRepRejUnclassifiedExePZeroPct}{n/a}
\rqSet{rqNonRepRejUnclassifiedExePOneNum}{0}
\rqSet{rqNonRepRejUnclassifiedExePOneDen}{0}
\rqSet{rqNonRepRejUnclassifiedExePOnePct}{n/a}
\rqSet{rqTermTotalCExeAllNum}{35894359}
\rqSet{rqTermTotalCExeAllDen}{35894359}
\rqSet{rqTermTotalCExeBOneNum}{29751287}
\rqSet{rqTermTotalCExeBOneDen}{29751287}
\rqSet{rqTermTotalCExeBTwoNum}{6142964}
\rqSet{rqTermTotalCExeBTwoDen}{6142964}
\rqSet{rqTermTotalCExePZeroNum}{0}
\rqSet{rqTermTotalCExePZeroDen}{0}
\rqSet{rqTermTotalCExePOneNum}{108}
\rqSet{rqTermTotalCExePOneDen}{108}
\rqSet{rqTermTotalRExeAllNum}{35894359}
\rqSet{rqTermTotalRExeAllDen}{35894359}
\rqSet{rqTermTotalRExeBOneNum}{29751287}
\rqSet{rqTermTotalRExeBOneDen}{29751287}
\rqSet{rqTermTotalRExeBTwoNum}{6142964}
\rqSet{rqTermTotalRExeBTwoDen}{6142964}
\rqSet{rqTermTotalRExePZeroNum}{0}
\rqSet{rqTermTotalRExePZeroDen}{0}
\rqSet{rqTermTotalRExePOneNum}{108}
\rqSet{rqTermTotalRExePOneDen}{108}
\rqSet{rqVSixSepExeAllNum}{0}
\rqSet{rqVSixSepExeAllDen}{0}
\rqSet{rqVSixSepExeAllPct}{n/a}
\rqSet{rqVSixSepExeAllKiroNum}{0}
\rqSet{rqVSixSepExeAllKiroDen}{0}
\rqSet{rqVSixSepExeAllKiroPct}{n/a}
\rqSet{rqVSixSepExeAllClaudeNum}{0}
\rqSet{rqVSixSepExeAllClaudeDen}{0}
\rqSet{rqVSixSepExeAllClaudePct}{n/a}
\rqSet{rqVSixSepExeAllCodexNum}{0}
\rqSet{rqVSixSepExeAllCodexDen}{0}
\rqSet{rqVSixSepExeAllCodexPct}{n/a}
\rqSet{rqVSixSepExeAllCTwoRustNum}{0}
\rqSet{rqVSixSepExeAllCTwoRustDen}{0}
\rqSet{rqVSixSepExeAllCTwoRustPct}{n/a}
\rqSet{rqVSixSepExeAllLaertesNum}{0}
\rqSet{rqVSixSepExeAllLaertesDen}{0}
\rqSet{rqVSixSepExeAllLaertesPct}{n/a}
\rqSet{rqVSixSepExeAllCTwoSaferRustNum}{0}
\rqSet{rqVSixSepExeAllCTwoSaferRustDen}{0}
\rqSet{rqVSixSepExeAllCTwoSaferRustPct}{n/a}
\rqSet{rqVSixSepExeAllSmartCTwoRustNum}{0}
\rqSet{rqVSixSepExeAllSmartCTwoRustDen}{0}
\rqSet{rqVSixSepExeAllSmartCTwoRustPct}{n/a}
\rqSet{rqVSixSepExeAllKimiNum}{0}
\rqSet{rqVSixSepExeAllKimiDen}{0}
\rqSet{rqVSixSepExeAllKimiPct}{n/a}
\rqSet{rqVSixSepExeAllGptFiveNum}{0}
\rqSet{rqVSixSepExeAllGptFiveDen}{0}
\rqSet{rqVSixSepExeAllGptFivePct}{n/a}
\rqSet{rqVSixSepExeAllGeminiNum}{0}
\rqSet{rqVSixSepExeAllGeminiDen}{0}
\rqSet{rqVSixSepExeAllGeminiPct}{n/a}
\rqSet{rqVSixSepExeBOneNum}{0}
\rqSet{rqVSixSepExeBOneDen}{0}
\rqSet{rqVSixSepExeBOnePct}{n/a}
\rqSet{rqVSixSepExeBTwoNum}{0}
\rqSet{rqVSixSepExeBTwoDen}{0}
\rqSet{rqVSixSepExeBTwoPct}{n/a}
\rqSet{rqVSixSepExePZeroNum}{0}
\rqSet{rqVSixSepExePZeroDen}{0}
\rqSet{rqVSixSepExePZeroPct}{n/a}
\rqSet{rqVSixSepExePOneNum}{0}
\rqSet{rqVSixSepExePOneDen}{0}
\rqSet{rqVSixSepExePOnePct}{n/a}
\rqSet{rqVSixRestExeAllNum}{176}
\rqSet{rqVSixRestExeAllDen}{476}
\rqSet{rqVSixRestExeAllPct}{37.0}
\rqSet{rqVSixRestExeAllKiroNum}{29}
\rqSet{rqVSixRestExeAllKiroDen}{88}
\rqSet{rqVSixRestExeAllKiroPct}{33.0}
\rqSet{rqVSixRestExeAllClaudeNum}{24}
\rqSet{rqVSixRestExeAllClaudeDen}{93}
\rqSet{rqVSixRestExeAllClaudePct}{25.8}
\rqSet{rqVSixRestExeAllCodexNum}{9}
\rqSet{rqVSixRestExeAllCodexDen}{54}
\rqSet{rqVSixRestExeAllCodexPct}{16.7}
\rqSet{rqVSixRestExeAllCTwoRustNum}{3}
\rqSet{rqVSixRestExeAllCTwoRustDen}{44}
\rqSet{rqVSixRestExeAllCTwoRustPct}{6.8}
\rqSet{rqVSixRestExeAllLaertesNum}{3}
\rqSet{rqVSixRestExeAllLaertesDen}{44}
\rqSet{rqVSixRestExeAllLaertesPct}{6.8}
\rqSet{rqVSixRestExeAllCTwoSaferRustNum}{12}
\rqSet{rqVSixRestExeAllCTwoSaferRustDen}{32}
\rqSet{rqVSixRestExeAllCTwoSaferRustPct}{37.5}
\rqSet{rqVSixRestExeAllSmartCTwoRustNum}{27}
\rqSet{rqVSixRestExeAllSmartCTwoRustDen}{33}
\rqSet{rqVSixRestExeAllSmartCTwoRustPct}{81.8}
\rqSet{rqVSixRestExeAllKimiNum}{18}
\rqSet{rqVSixRestExeAllKimiDen}{24}
\rqSet{rqVSixRestExeAllKimiPct}{75.0}
\rqSet{rqVSixRestExeAllGptFiveNum}{25}
\rqSet{rqVSixRestExeAllGptFiveDen}{30}
\rqSet{rqVSixRestExeAllGptFivePct}{83.3}
\rqSet{rqVSixRestExeAllGeminiNum}{26}
\rqSet{rqVSixRestExeAllGeminiDen}{34}
\rqSet{rqVSixRestExeAllGeminiPct}{76.5}
\rqSet{rqVSixRestExeBOneNum}{152}
\rqSet{rqVSixRestExeBOneDen}{315}
\rqSet{rqVSixRestExeBOnePct}{48.3}
\rqSet{rqVSixRestExeBTwoNum}{24}
\rqSet{rqVSixRestExeBTwoDen}{53}
\rqSet{rqVSixRestExeBTwoPct}{45.3}
\rqSet{rqVSixRestExePZeroNum}{0}
\rqSet{rqVSixRestExePZeroDen}{0}
\rqSet{rqVSixRestExePZeroPct}{n/a}
\rqSet{rqVSixRestExePOneNum}{0}
\rqSet{rqVSixRestExePOneDen}{108}
\rqSet{rqVSixRestExePOnePct}{0.0}
\rqSet{rqFFExeAllNum}{0}
\rqSet{rqFFExeAllDen}{92}
\rqSet{rqFFExeAllPct}{0.0}
\rqSet{rqFFExeAllKiroNum}{0}
\rqSet{rqFFExeAllKiroDen}{10}
\rqSet{rqFFExeAllKiroPct}{0.0}
\rqSet{rqFFExeAllClaudeNum}{0}
\rqSet{rqFFExeAllClaudeDen}{5}
\rqSet{rqFFExeAllClaudePct}{0.0}
\rqSet{rqFFExeAllCodexNum}{0}
\rqSet{rqFFExeAllCodexDen}{7}
\rqSet{rqFFExeAllCodexPct}{0.0}
\rqSet{rqFFExeAllCTwoRustNum}{0}
\rqSet{rqFFExeAllCTwoRustDen}{0}
\rqSet{rqFFExeAllCTwoRustPct}{n/a}
\rqSet{rqFFExeAllLaertesNum}{0}
\rqSet{rqFFExeAllLaertesDen}{0}
\rqSet{rqFFExeAllLaertesPct}{n/a}
\rqSet{rqFFExeAllCTwoSaferRustNum}{0}
\rqSet{rqFFExeAllCTwoSaferRustDen}{12}
\rqSet{rqFFExeAllCTwoSaferRustPct}{0.0}
\rqSet{rqFFExeAllSmartCTwoRustNum}{0}
\rqSet{rqFFExeAllSmartCTwoRustDen}{8}
\rqSet{rqFFExeAllSmartCTwoRustPct}{0.0}
\rqSet{rqFFExeAllKimiNum}{0}
\rqSet{rqFFExeAllKimiDen}{17}
\rqSet{rqFFExeAllKimiPct}{0.0}
\rqSet{rqFFExeAllGptFiveNum}{0}
\rqSet{rqFFExeAllGptFiveDen}{18}
\rqSet{rqFFExeAllGptFivePct}{0.0}
\rqSet{rqFFExeAllGeminiNum}{0}
\rqSet{rqFFExeAllGeminiDen}{15}
\rqSet{rqFFExeAllGeminiPct}{0.0}
\rqSet{rqFFExeBOneNum}{0}
\rqSet{rqFFExeBOneDen}{41}
\rqSet{rqFFExeBOnePct}{0.0}
\rqSet{rqFFExeBOneKiroNum}{0}
\rqSet{rqFFExeBOneKiroDen}{0}
\rqSet{rqFFExeBOneKiroPct}{n/a}
\rqSet{rqFFExeBOneClaudeNum}{0}
\rqSet{rqFFExeBOneClaudeDen}{0}
\rqSet{rqFFExeBOneClaudePct}{n/a}
\rqSet{rqFFExeBOneCodexNum}{0}
\rqSet{rqFFExeBOneCodexDen}{1}
\rqSet{rqFFExeBOneCodexPct}{0.0}
\rqSet{rqFFExeBOneCTwoRustNum}{0}
\rqSet{rqFFExeBOneCTwoRustDen}{0}
\rqSet{rqFFExeBOneCTwoRustPct}{n/a}
\rqSet{rqFFExeBOneLaertesNum}{0}
\rqSet{rqFFExeBOneLaertesDen}{0}
\rqSet{rqFFExeBOneLaertesPct}{n/a}
\rqSet{rqFFExeBOneCTwoSaferRustNum}{0}
\rqSet{rqFFExeBOneCTwoSaferRustDen}{11}
\rqSet{rqFFExeBOneCTwoSaferRustPct}{0.0}
\rqSet{rqFFExeBOneSmartCTwoRustNum}{0}
\rqSet{rqFFExeBOneSmartCTwoRustDen}{3}
\rqSet{rqFFExeBOneSmartCTwoRustPct}{0.0}
\rqSet{rqFFExeBOneKimiNum}{0}
\rqSet{rqFFExeBOneKimiDen}{12}
\rqSet{rqFFExeBOneKimiPct}{0.0}
\rqSet{rqFFExeBOneGptFiveNum}{0}
\rqSet{rqFFExeBOneGptFiveDen}{8}
\rqSet{rqFFExeBOneGptFivePct}{0.0}
\rqSet{rqFFExeBOneGeminiNum}{0}
\rqSet{rqFFExeBOneGeminiDen}{6}
\rqSet{rqFFExeBOneGeminiPct}{0.0}
\rqSet{rqFFExeBTwoNum}{0}
\rqSet{rqFFExeBTwoDen}{47}
\rqSet{rqFFExeBTwoPct}{0.0}
\rqSet{rqFFExeBTwoKiroNum}{0}
\rqSet{rqFFExeBTwoKiroDen}{10}
\rqSet{rqFFExeBTwoKiroPct}{0.0}
\rqSet{rqFFExeBTwoClaudeNum}{0}
\rqSet{rqFFExeBTwoClaudeDen}{5}
\rqSet{rqFFExeBTwoClaudePct}{0.0}
\rqSet{rqFFExeBTwoCodexNum}{0}
\rqSet{rqFFExeBTwoCodexDen}{6}
\rqSet{rqFFExeBTwoCodexPct}{0.0}
\rqSet{rqFFExeBTwoCTwoRustNum}{0}
\rqSet{rqFFExeBTwoCTwoRustDen}{0}
\rqSet{rqFFExeBTwoCTwoRustPct}{n/a}
\rqSet{rqFFExeBTwoLaertesNum}{0}
\rqSet{rqFFExeBTwoLaertesDen}{0}
\rqSet{rqFFExeBTwoLaertesPct}{n/a}
\rqSet{rqFFExeBTwoCTwoSaferRustNum}{0}
\rqSet{rqFFExeBTwoCTwoSaferRustDen}{1}
\rqSet{rqFFExeBTwoCTwoSaferRustPct}{0.0}
\rqSet{rqFFExeBTwoSmartCTwoRustNum}{0}
\rqSet{rqFFExeBTwoSmartCTwoRustDen}{4}
\rqSet{rqFFExeBTwoSmartCTwoRustPct}{0.0}
\rqSet{rqFFExeBTwoKimiNum}{0}
\rqSet{rqFFExeBTwoKimiDen}{4}
\rqSet{rqFFExeBTwoKimiPct}{0.0}
\rqSet{rqFFExeBTwoGptFiveNum}{0}
\rqSet{rqFFExeBTwoGptFiveDen}{9}
\rqSet{rqFFExeBTwoGptFivePct}{0.0}
\rqSet{rqFFExeBTwoGeminiNum}{0}
\rqSet{rqFFExeBTwoGeminiDen}{8}
\rqSet{rqFFExeBTwoGeminiPct}{0.0}
\rqSet{rqFFExePZeroNum}{0}
\rqSet{rqFFExePZeroDen}{4}
\rqSet{rqFFExePZeroPct}{0.0}
\rqSet{rqFFExePZeroKiroNum}{0}
\rqSet{rqFFExePZeroKiroDen}{0}
\rqSet{rqFFExePZeroKiroPct}{n/a}
\rqSet{rqFFExePZeroClaudeNum}{0}
\rqSet{rqFFExePZeroClaudeDen}{0}
\rqSet{rqFFExePZeroClaudePct}{n/a}
\rqSet{rqFFExePZeroCodexNum}{0}
\rqSet{rqFFExePZeroCodexDen}{0}
\rqSet{rqFFExePZeroCodexPct}{n/a}
\rqSet{rqFFExePZeroCTwoRustNum}{0}
\rqSet{rqFFExePZeroCTwoRustDen}{0}
\rqSet{rqFFExePZeroCTwoRustPct}{n/a}
\rqSet{rqFFExePZeroLaertesNum}{0}
\rqSet{rqFFExePZeroLaertesDen}{0}
\rqSet{rqFFExePZeroLaertesPct}{n/a}
\rqSet{rqFFExePZeroCTwoSaferRustNum}{0}
\rqSet{rqFFExePZeroCTwoSaferRustDen}{0}
\rqSet{rqFFExePZeroCTwoSaferRustPct}{n/a}
\rqSet{rqFFExePZeroSmartCTwoRustNum}{0}
\rqSet{rqFFExePZeroSmartCTwoRustDen}{1}
\rqSet{rqFFExePZeroSmartCTwoRustPct}{0.0}
\rqSet{rqFFExePZeroKimiNum}{0}
\rqSet{rqFFExePZeroKimiDen}{1}
\rqSet{rqFFExePZeroKimiPct}{0.0}
\rqSet{rqFFExePZeroGptFiveNum}{0}
\rqSet{rqFFExePZeroGptFiveDen}{1}
\rqSet{rqFFExePZeroGptFivePct}{0.0}
\rqSet{rqFFExePZeroGeminiNum}{0}
\rqSet{rqFFExePZeroGeminiDen}{1}
\rqSet{rqFFExePZeroGeminiPct}{0.0}
\rqSet{rqFFExePOneNum}{0}
\rqSet{rqFFExePOneDen}{0}
\rqSet{rqFFExePOnePct}{n/a}
\rqSet{rqFFExePOneKiroNum}{0}
\rqSet{rqFFExePOneKiroDen}{0}
\rqSet{rqFFExePOneKiroPct}{n/a}
\rqSet{rqFFExePOneClaudeNum}{0}
\rqSet{rqFFExePOneClaudeDen}{0}
\rqSet{rqFFExePOneClaudePct}{n/a}
\rqSet{rqFFExePOneCodexNum}{0}
\rqSet{rqFFExePOneCodexDen}{0}
\rqSet{rqFFExePOneCodexPct}{n/a}
\rqSet{rqFFExePOneCTwoRustNum}{0}
\rqSet{rqFFExePOneCTwoRustDen}{0}
\rqSet{rqFFExePOneCTwoRustPct}{n/a}
\rqSet{rqFFExePOneLaertesNum}{0}
\rqSet{rqFFExePOneLaertesDen}{0}
\rqSet{rqFFExePOneLaertesPct}{n/a}
\rqSet{rqFFExePOneCTwoSaferRustNum}{0}
\rqSet{rqFFExePOneCTwoSaferRustDen}{0}
\rqSet{rqFFExePOneCTwoSaferRustPct}{n/a}
\rqSet{rqFFExePOneSmartCTwoRustNum}{0}
\rqSet{rqFFExePOneSmartCTwoRustDen}{0}
\rqSet{rqFFExePOneSmartCTwoRustPct}{n/a}
\rqSet{rqFFExePOneKimiNum}{0}
\rqSet{rqFFExePOneKimiDen}{0}
\rqSet{rqFFExePOneKimiPct}{n/a}
\rqSet{rqFFExePOneGptFiveNum}{0}
\rqSet{rqFFExePOneGptFiveDen}{0}
\rqSet{rqFFExePOneGptFivePct}{n/a}
\rqSet{rqFFExePOneGeminiNum}{0}
\rqSet{rqFFExePOneGeminiDen}{0}
\rqSet{rqFFExePOneGeminiPct}{n/a}
\rqSet{rqFFDisagreeExeAllNum}{92}
\rqSet{rqFFDisagreeExeAllDen}{92}
\rqSet{rqFFDisagreeExeAllPct}{100.0}
\rqSet{rqFFDisagreeExeAllKiroNum}{10}
\rqSet{rqFFDisagreeExeAllKiroDen}{10}
\rqSet{rqFFDisagreeExeAllKiroPct}{100.0}
\rqSet{rqFFDisagreeExeAllClaudeNum}{5}
\rqSet{rqFFDisagreeExeAllClaudeDen}{5}
\rqSet{rqFFDisagreeExeAllClaudePct}{100.0}
\rqSet{rqFFDisagreeExeAllCodexNum}{7}
\rqSet{rqFFDisagreeExeAllCodexDen}{7}
\rqSet{rqFFDisagreeExeAllCodexPct}{100.0}
\rqSet{rqFFDisagreeExeAllCTwoRustNum}{0}
\rqSet{rqFFDisagreeExeAllCTwoRustDen}{0}
\rqSet{rqFFDisagreeExeAllCTwoRustPct}{n/a}
\rqSet{rqFFDisagreeExeAllLaertesNum}{0}
\rqSet{rqFFDisagreeExeAllLaertesDen}{0}
\rqSet{rqFFDisagreeExeAllLaertesPct}{n/a}
\rqSet{rqFFDisagreeExeAllCTwoSaferRustNum}{12}
\rqSet{rqFFDisagreeExeAllCTwoSaferRustDen}{12}
\rqSet{rqFFDisagreeExeAllCTwoSaferRustPct}{100.0}
\rqSet{rqFFDisagreeExeAllSmartCTwoRustNum}{8}
\rqSet{rqFFDisagreeExeAllSmartCTwoRustDen}{8}
\rqSet{rqFFDisagreeExeAllSmartCTwoRustPct}{100.0}
\rqSet{rqFFDisagreeExeAllKimiNum}{17}
\rqSet{rqFFDisagreeExeAllKimiDen}{17}
\rqSet{rqFFDisagreeExeAllKimiPct}{100.0}
\rqSet{rqFFDisagreeExeAllGptFiveNum}{18}
\rqSet{rqFFDisagreeExeAllGptFiveDen}{18}
\rqSet{rqFFDisagreeExeAllGptFivePct}{100.0}
\rqSet{rqFFDisagreeExeAllGeminiNum}{15}
\rqSet{rqFFDisagreeExeAllGeminiDen}{15}
\rqSet{rqFFDisagreeExeAllGeminiPct}{100.0}
\rqSet{rqFFDisagreeExeBOneNum}{41}
\rqSet{rqFFDisagreeExeBOneDen}{41}
\rqSet{rqFFDisagreeExeBOnePct}{100.0}
\rqSet{rqFFDisagreeExeBOneKiroNum}{0}
\rqSet{rqFFDisagreeExeBOneKiroDen}{0}
\rqSet{rqFFDisagreeExeBOneKiroPct}{n/a}
\rqSet{rqFFDisagreeExeBOneClaudeNum}{0}
\rqSet{rqFFDisagreeExeBOneClaudeDen}{0}
\rqSet{rqFFDisagreeExeBOneClaudePct}{n/a}
\rqSet{rqFFDisagreeExeBOneCodexNum}{1}
\rqSet{rqFFDisagreeExeBOneCodexDen}{1}
\rqSet{rqFFDisagreeExeBOneCodexPct}{100.0}
\rqSet{rqFFDisagreeExeBOneCTwoRustNum}{0}
\rqSet{rqFFDisagreeExeBOneCTwoRustDen}{0}
\rqSet{rqFFDisagreeExeBOneCTwoRustPct}{n/a}
\rqSet{rqFFDisagreeExeBOneLaertesNum}{0}
\rqSet{rqFFDisagreeExeBOneLaertesDen}{0}
\rqSet{rqFFDisagreeExeBOneLaertesPct}{n/a}
\rqSet{rqFFDisagreeExeBOneCTwoSaferRustNum}{11}
\rqSet{rqFFDisagreeExeBOneCTwoSaferRustDen}{11}
\rqSet{rqFFDisagreeExeBOneCTwoSaferRustPct}{100.0}
\rqSet{rqFFDisagreeExeBOneSmartCTwoRustNum}{3}
\rqSet{rqFFDisagreeExeBOneSmartCTwoRustDen}{3}
\rqSet{rqFFDisagreeExeBOneSmartCTwoRustPct}{100.0}
\rqSet{rqFFDisagreeExeBOneKimiNum}{12}
\rqSet{rqFFDisagreeExeBOneKimiDen}{12}
\rqSet{rqFFDisagreeExeBOneKimiPct}{100.0}
\rqSet{rqFFDisagreeExeBOneGptFiveNum}{8}
\rqSet{rqFFDisagreeExeBOneGptFiveDen}{8}
\rqSet{rqFFDisagreeExeBOneGptFivePct}{100.0}
\rqSet{rqFFDisagreeExeBOneGeminiNum}{6}
\rqSet{rqFFDisagreeExeBOneGeminiDen}{6}
\rqSet{rqFFDisagreeExeBOneGeminiPct}{100.0}
\rqSet{rqFFDisagreeExeBTwoNum}{47}
\rqSet{rqFFDisagreeExeBTwoDen}{47}
\rqSet{rqFFDisagreeExeBTwoPct}{100.0}
\rqSet{rqFFDisagreeExeBTwoKiroNum}{10}
\rqSet{rqFFDisagreeExeBTwoKiroDen}{10}
\rqSet{rqFFDisagreeExeBTwoKiroPct}{100.0}
\rqSet{rqFFDisagreeExeBTwoClaudeNum}{5}
\rqSet{rqFFDisagreeExeBTwoClaudeDen}{5}
\rqSet{rqFFDisagreeExeBTwoClaudePct}{100.0}
\rqSet{rqFFDisagreeExeBTwoCodexNum}{6}
\rqSet{rqFFDisagreeExeBTwoCodexDen}{6}
\rqSet{rqFFDisagreeExeBTwoCodexPct}{100.0}
\rqSet{rqFFDisagreeExeBTwoCTwoRustNum}{0}
\rqSet{rqFFDisagreeExeBTwoCTwoRustDen}{0}
\rqSet{rqFFDisagreeExeBTwoCTwoRustPct}{n/a}
\rqSet{rqFFDisagreeExeBTwoLaertesNum}{0}
\rqSet{rqFFDisagreeExeBTwoLaertesDen}{0}
\rqSet{rqFFDisagreeExeBTwoLaertesPct}{n/a}
\rqSet{rqFFDisagreeExeBTwoCTwoSaferRustNum}{1}
\rqSet{rqFFDisagreeExeBTwoCTwoSaferRustDen}{1}
\rqSet{rqFFDisagreeExeBTwoCTwoSaferRustPct}{100.0}
\rqSet{rqFFDisagreeExeBTwoSmartCTwoRustNum}{4}
\rqSet{rqFFDisagreeExeBTwoSmartCTwoRustDen}{4}
\rqSet{rqFFDisagreeExeBTwoSmartCTwoRustPct}{100.0}
\rqSet{rqFFDisagreeExeBTwoKimiNum}{4}
\rqSet{rqFFDisagreeExeBTwoKimiDen}{4}
\rqSet{rqFFDisagreeExeBTwoKimiPct}{100.0}
\rqSet{rqFFDisagreeExeBTwoGptFiveNum}{9}
\rqSet{rqFFDisagreeExeBTwoGptFiveDen}{9}
\rqSet{rqFFDisagreeExeBTwoGptFivePct}{100.0}
\rqSet{rqFFDisagreeExeBTwoGeminiNum}{8}
\rqSet{rqFFDisagreeExeBTwoGeminiDen}{8}
\rqSet{rqFFDisagreeExeBTwoGeminiPct}{100.0}
\rqSet{rqFFDisagreeExePZeroNum}{4}
\rqSet{rqFFDisagreeExePZeroDen}{4}
\rqSet{rqFFDisagreeExePZeroPct}{100.0}
\rqSet{rqFFDisagreeExePZeroKiroNum}{0}
\rqSet{rqFFDisagreeExePZeroKiroDen}{0}
\rqSet{rqFFDisagreeExePZeroKiroPct}{n/a}
\rqSet{rqFFDisagreeExePZeroClaudeNum}{0}
\rqSet{rqFFDisagreeExePZeroClaudeDen}{0}
\rqSet{rqFFDisagreeExePZeroClaudePct}{n/a}
\rqSet{rqFFDisagreeExePZeroCodexNum}{0}
\rqSet{rqFFDisagreeExePZeroCodexDen}{0}
\rqSet{rqFFDisagreeExePZeroCodexPct}{n/a}
\rqSet{rqFFDisagreeExePZeroCTwoRustNum}{0}
\rqSet{rqFFDisagreeExePZeroCTwoRustDen}{0}
\rqSet{rqFFDisagreeExePZeroCTwoRustPct}{n/a}
\rqSet{rqFFDisagreeExePZeroLaertesNum}{0}
\rqSet{rqFFDisagreeExePZeroLaertesDen}{0}
\rqSet{rqFFDisagreeExePZeroLaertesPct}{n/a}
\rqSet{rqFFDisagreeExePZeroCTwoSaferRustNum}{0}
\rqSet{rqFFDisagreeExePZeroCTwoSaferRustDen}{0}
\rqSet{rqFFDisagreeExePZeroCTwoSaferRustPct}{n/a}
\rqSet{rqFFDisagreeExePZeroSmartCTwoRustNum}{1}
\rqSet{rqFFDisagreeExePZeroSmartCTwoRustDen}{1}
\rqSet{rqFFDisagreeExePZeroSmartCTwoRustPct}{100.0}
\rqSet{rqFFDisagreeExePZeroKimiNum}{1}
\rqSet{rqFFDisagreeExePZeroKimiDen}{1}
\rqSet{rqFFDisagreeExePZeroKimiPct}{100.0}
\rqSet{rqFFDisagreeExePZeroGptFiveNum}{1}
\rqSet{rqFFDisagreeExePZeroGptFiveDen}{1}
\rqSet{rqFFDisagreeExePZeroGptFivePct}{100.0}
\rqSet{rqFFDisagreeExePZeroGeminiNum}{1}
\rqSet{rqFFDisagreeExePZeroGeminiDen}{1}
\rqSet{rqFFDisagreeExePZeroGeminiPct}{100.0}
\rqSet{rqFFDisagreeExePOneNum}{0}
\rqSet{rqFFDisagreeExePOneDen}{0}
\rqSet{rqFFDisagreeExePOnePct}{n/a}
\rqSet{rqFFDisagreeExePOneKiroNum}{0}
\rqSet{rqFFDisagreeExePOneKiroDen}{0}
\rqSet{rqFFDisagreeExePOneKiroPct}{n/a}
\rqSet{rqFFDisagreeExePOneClaudeNum}{0}
\rqSet{rqFFDisagreeExePOneClaudeDen}{0}
\rqSet{rqFFDisagreeExePOneClaudePct}{n/a}
\rqSet{rqFFDisagreeExePOneCodexNum}{0}
\rqSet{rqFFDisagreeExePOneCodexDen}{0}
\rqSet{rqFFDisagreeExePOneCodexPct}{n/a}
\rqSet{rqFFDisagreeExePOneCTwoRustNum}{0}
\rqSet{rqFFDisagreeExePOneCTwoRustDen}{0}
\rqSet{rqFFDisagreeExePOneCTwoRustPct}{n/a}
\rqSet{rqFFDisagreeExePOneLaertesNum}{0}
\rqSet{rqFFDisagreeExePOneLaertesDen}{0}
\rqSet{rqFFDisagreeExePOneLaertesPct}{n/a}
\rqSet{rqFFDisagreeExePOneCTwoSaferRustNum}{0}
\rqSet{rqFFDisagreeExePOneCTwoSaferRustDen}{0}
\rqSet{rqFFDisagreeExePOneCTwoSaferRustPct}{n/a}
\rqSet{rqFFDisagreeExePOneSmartCTwoRustNum}{0}
\rqSet{rqFFDisagreeExePOneSmartCTwoRustDen}{0}
\rqSet{rqFFDisagreeExePOneSmartCTwoRustPct}{n/a}
\rqSet{rqFFDisagreeExePOneKimiNum}{0}
\rqSet{rqFFDisagreeExePOneKimiDen}{0}
\rqSet{rqFFDisagreeExePOneKimiPct}{n/a}
\rqSet{rqFFDisagreeExePOneGptFiveNum}{0}
\rqSet{rqFFDisagreeExePOneGptFiveDen}{0}
\rqSet{rqFFDisagreeExePOneGptFivePct}{n/a}
\rqSet{rqFFDisagreeExePOneGeminiNum}{0}
\rqSet{rqFFDisagreeExePOneGeminiDen}{0}
\rqSet{rqFFDisagreeExePOneGeminiPct}{n/a}
\rqSet{rqFFNotAdjExeAllNum}{0}
\rqSet{rqFFNotAdjExeAllDen}{92}
\rqSet{rqFFNotAdjExeAllPct}{0.0}
\rqSet{rqFFNotAdjExeAllKiroNum}{0}
\rqSet{rqFFNotAdjExeAllKiroDen}{10}
\rqSet{rqFFNotAdjExeAllKiroPct}{0.0}
\rqSet{rqFFNotAdjExeAllClaudeNum}{0}
\rqSet{rqFFNotAdjExeAllClaudeDen}{5}
\rqSet{rqFFNotAdjExeAllClaudePct}{0.0}
\rqSet{rqFFNotAdjExeAllCodexNum}{0}
\rqSet{rqFFNotAdjExeAllCodexDen}{7}
\rqSet{rqFFNotAdjExeAllCodexPct}{0.0}
\rqSet{rqFFNotAdjExeAllCTwoRustNum}{0}
\rqSet{rqFFNotAdjExeAllCTwoRustDen}{0}
\rqSet{rqFFNotAdjExeAllCTwoRustPct}{n/a}
\rqSet{rqFFNotAdjExeAllLaertesNum}{0}
\rqSet{rqFFNotAdjExeAllLaertesDen}{0}
\rqSet{rqFFNotAdjExeAllLaertesPct}{n/a}
\rqSet{rqFFNotAdjExeAllCTwoSaferRustNum}{0}
\rqSet{rqFFNotAdjExeAllCTwoSaferRustDen}{12}
\rqSet{rqFFNotAdjExeAllCTwoSaferRustPct}{0.0}
\rqSet{rqFFNotAdjExeAllSmartCTwoRustNum}{0}
\rqSet{rqFFNotAdjExeAllSmartCTwoRustDen}{8}
\rqSet{rqFFNotAdjExeAllSmartCTwoRustPct}{0.0}
\rqSet{rqFFNotAdjExeAllKimiNum}{0}
\rqSet{rqFFNotAdjExeAllKimiDen}{17}
\rqSet{rqFFNotAdjExeAllKimiPct}{0.0}
\rqSet{rqFFNotAdjExeAllGptFiveNum}{0}
\rqSet{rqFFNotAdjExeAllGptFiveDen}{18}
\rqSet{rqFFNotAdjExeAllGptFivePct}{0.0}
\rqSet{rqFFNotAdjExeAllGeminiNum}{0}
\rqSet{rqFFNotAdjExeAllGeminiDen}{15}
\rqSet{rqFFNotAdjExeAllGeminiPct}{0.0}
\rqSet{rqFFNotAdjExeBOneNum}{0}
\rqSet{rqFFNotAdjExeBOneDen}{41}
\rqSet{rqFFNotAdjExeBOnePct}{0.0}
\rqSet{rqFFNotAdjExeBOneKiroNum}{0}
\rqSet{rqFFNotAdjExeBOneKiroDen}{0}
\rqSet{rqFFNotAdjExeBOneKiroPct}{n/a}
\rqSet{rqFFNotAdjExeBOneClaudeNum}{0}
\rqSet{rqFFNotAdjExeBOneClaudeDen}{0}
\rqSet{rqFFNotAdjExeBOneClaudePct}{n/a}
\rqSet{rqFFNotAdjExeBOneCodexNum}{0}
\rqSet{rqFFNotAdjExeBOneCodexDen}{1}
\rqSet{rqFFNotAdjExeBOneCodexPct}{0.0}
\rqSet{rqFFNotAdjExeBOneCTwoRustNum}{0}
\rqSet{rqFFNotAdjExeBOneCTwoRustDen}{0}
\rqSet{rqFFNotAdjExeBOneCTwoRustPct}{n/a}
\rqSet{rqFFNotAdjExeBOneLaertesNum}{0}
\rqSet{rqFFNotAdjExeBOneLaertesDen}{0}
\rqSet{rqFFNotAdjExeBOneLaertesPct}{n/a}
\rqSet{rqFFNotAdjExeBOneCTwoSaferRustNum}{0}
\rqSet{rqFFNotAdjExeBOneCTwoSaferRustDen}{11}
\rqSet{rqFFNotAdjExeBOneCTwoSaferRustPct}{0.0}
\rqSet{rqFFNotAdjExeBOneSmartCTwoRustNum}{0}
\rqSet{rqFFNotAdjExeBOneSmartCTwoRustDen}{3}
\rqSet{rqFFNotAdjExeBOneSmartCTwoRustPct}{0.0}
\rqSet{rqFFNotAdjExeBOneKimiNum}{0}
\rqSet{rqFFNotAdjExeBOneKimiDen}{12}
\rqSet{rqFFNotAdjExeBOneKimiPct}{0.0}
\rqSet{rqFFNotAdjExeBOneGptFiveNum}{0}
\rqSet{rqFFNotAdjExeBOneGptFiveDen}{8}
\rqSet{rqFFNotAdjExeBOneGptFivePct}{0.0}
\rqSet{rqFFNotAdjExeBOneGeminiNum}{0}
\rqSet{rqFFNotAdjExeBOneGeminiDen}{6}
\rqSet{rqFFNotAdjExeBOneGeminiPct}{0.0}
\rqSet{rqFFNotAdjExeBTwoNum}{0}
\rqSet{rqFFNotAdjExeBTwoDen}{47}
\rqSet{rqFFNotAdjExeBTwoPct}{0.0}
\rqSet{rqFFNotAdjExeBTwoKiroNum}{0}
\rqSet{rqFFNotAdjExeBTwoKiroDen}{10}
\rqSet{rqFFNotAdjExeBTwoKiroPct}{0.0}
\rqSet{rqFFNotAdjExeBTwoClaudeNum}{0}
\rqSet{rqFFNotAdjExeBTwoClaudeDen}{5}
\rqSet{rqFFNotAdjExeBTwoClaudePct}{0.0}
\rqSet{rqFFNotAdjExeBTwoCodexNum}{0}
\rqSet{rqFFNotAdjExeBTwoCodexDen}{6}
\rqSet{rqFFNotAdjExeBTwoCodexPct}{0.0}
\rqSet{rqFFNotAdjExeBTwoCTwoRustNum}{0}
\rqSet{rqFFNotAdjExeBTwoCTwoRustDen}{0}
\rqSet{rqFFNotAdjExeBTwoCTwoRustPct}{n/a}
\rqSet{rqFFNotAdjExeBTwoLaertesNum}{0}
\rqSet{rqFFNotAdjExeBTwoLaertesDen}{0}
\rqSet{rqFFNotAdjExeBTwoLaertesPct}{n/a}
\rqSet{rqFFNotAdjExeBTwoCTwoSaferRustNum}{0}
\rqSet{rqFFNotAdjExeBTwoCTwoSaferRustDen}{1}
\rqSet{rqFFNotAdjExeBTwoCTwoSaferRustPct}{0.0}
\rqSet{rqFFNotAdjExeBTwoSmartCTwoRustNum}{0}
\rqSet{rqFFNotAdjExeBTwoSmartCTwoRustDen}{4}
\rqSet{rqFFNotAdjExeBTwoSmartCTwoRustPct}{0.0}
\rqSet{rqFFNotAdjExeBTwoKimiNum}{0}
\rqSet{rqFFNotAdjExeBTwoKimiDen}{4}
\rqSet{rqFFNotAdjExeBTwoKimiPct}{0.0}
\rqSet{rqFFNotAdjExeBTwoGptFiveNum}{0}
\rqSet{rqFFNotAdjExeBTwoGptFiveDen}{9}
\rqSet{rqFFNotAdjExeBTwoGptFivePct}{0.0}
\rqSet{rqFFNotAdjExeBTwoGeminiNum}{0}
\rqSet{rqFFNotAdjExeBTwoGeminiDen}{8}
\rqSet{rqFFNotAdjExeBTwoGeminiPct}{0.0}
\rqSet{rqFFNotAdjExePZeroNum}{0}
\rqSet{rqFFNotAdjExePZeroDen}{4}
\rqSet{rqFFNotAdjExePZeroPct}{0.0}
\rqSet{rqFFNotAdjExePZeroKiroNum}{0}
\rqSet{rqFFNotAdjExePZeroKiroDen}{0}
\rqSet{rqFFNotAdjExePZeroKiroPct}{n/a}
\rqSet{rqFFNotAdjExePZeroClaudeNum}{0}
\rqSet{rqFFNotAdjExePZeroClaudeDen}{0}
\rqSet{rqFFNotAdjExePZeroClaudePct}{n/a}
\rqSet{rqFFNotAdjExePZeroCodexNum}{0}
\rqSet{rqFFNotAdjExePZeroCodexDen}{0}
\rqSet{rqFFNotAdjExePZeroCodexPct}{n/a}
\rqSet{rqFFNotAdjExePZeroCTwoRustNum}{0}
\rqSet{rqFFNotAdjExePZeroCTwoRustDen}{0}
\rqSet{rqFFNotAdjExePZeroCTwoRustPct}{n/a}
\rqSet{rqFFNotAdjExePZeroLaertesNum}{0}
\rqSet{rqFFNotAdjExePZeroLaertesDen}{0}
\rqSet{rqFFNotAdjExePZeroLaertesPct}{n/a}
\rqSet{rqFFNotAdjExePZeroCTwoSaferRustNum}{0}
\rqSet{rqFFNotAdjExePZeroCTwoSaferRustDen}{0}
\rqSet{rqFFNotAdjExePZeroCTwoSaferRustPct}{n/a}
\rqSet{rqFFNotAdjExePZeroSmartCTwoRustNum}{0}
\rqSet{rqFFNotAdjExePZeroSmartCTwoRustDen}{1}
\rqSet{rqFFNotAdjExePZeroSmartCTwoRustPct}{0.0}
\rqSet{rqFFNotAdjExePZeroKimiNum}{0}
\rqSet{rqFFNotAdjExePZeroKimiDen}{1}
\rqSet{rqFFNotAdjExePZeroKimiPct}{0.0}
\rqSet{rqFFNotAdjExePZeroGptFiveNum}{0}
\rqSet{rqFFNotAdjExePZeroGptFiveDen}{1}
\rqSet{rqFFNotAdjExePZeroGptFivePct}{0.0}
\rqSet{rqFFNotAdjExePZeroGeminiNum}{0}
\rqSet{rqFFNotAdjExePZeroGeminiDen}{1}
\rqSet{rqFFNotAdjExePZeroGeminiPct}{0.0}
\rqSet{rqFFNotAdjExePOneNum}{0}
\rqSet{rqFFNotAdjExePOneDen}{0}
\rqSet{rqFFNotAdjExePOnePct}{n/a}
\rqSet{rqFFNotAdjExePOneKiroNum}{0}
\rqSet{rqFFNotAdjExePOneKiroDen}{0}
\rqSet{rqFFNotAdjExePOneKiroPct}{n/a}
\rqSet{rqFFNotAdjExePOneClaudeNum}{0}
\rqSet{rqFFNotAdjExePOneClaudeDen}{0}
\rqSet{rqFFNotAdjExePOneClaudePct}{n/a}
\rqSet{rqFFNotAdjExePOneCodexNum}{0}
\rqSet{rqFFNotAdjExePOneCodexDen}{0}
\rqSet{rqFFNotAdjExePOneCodexPct}{n/a}
\rqSet{rqFFNotAdjExePOneCTwoRustNum}{0}
\rqSet{rqFFNotAdjExePOneCTwoRustDen}{0}
\rqSet{rqFFNotAdjExePOneCTwoRustPct}{n/a}
\rqSet{rqFFNotAdjExePOneLaertesNum}{0}
\rqSet{rqFFNotAdjExePOneLaertesDen}{0}
\rqSet{rqFFNotAdjExePOneLaertesPct}{n/a}
\rqSet{rqFFNotAdjExePOneCTwoSaferRustNum}{0}
\rqSet{rqFFNotAdjExePOneCTwoSaferRustDen}{0}
\rqSet{rqFFNotAdjExePOneCTwoSaferRustPct}{n/a}
\rqSet{rqFFNotAdjExePOneSmartCTwoRustNum}{0}
\rqSet{rqFFNotAdjExePOneSmartCTwoRustDen}{0}
\rqSet{rqFFNotAdjExePOneSmartCTwoRustPct}{n/a}
\rqSet{rqFFNotAdjExePOneKimiNum}{0}
\rqSet{rqFFNotAdjExePOneKimiDen}{0}
\rqSet{rqFFNotAdjExePOneKimiPct}{n/a}
\rqSet{rqFFNotAdjExePOneGptFiveNum}{0}
\rqSet{rqFFNotAdjExePOneGptFiveDen}{0}
\rqSet{rqFFNotAdjExePOneGptFivePct}{n/a}
\rqSet{rqFFNotAdjExePOneGeminiNum}{0}
\rqSet{rqFFNotAdjExePOneGeminiDen}{0}
\rqSet{rqFFNotAdjExePOneGeminiPct}{n/a}
\rqSet{rqFFNotReexExeAllNum}{36}
\rqSet{rqFFNotReexExeAllDen}{128}
\rqSet{rqFFNotReexExeAllPct}{28.1}
\rqSet{rqFFNotReexExeAllKiroNum}{0}
\rqSet{rqFFNotReexExeAllKiroDen}{10}
\rqSet{rqFFNotReexExeAllKiroPct}{0.0}
\rqSet{rqFFNotReexExeAllClaudeNum}{0}
\rqSet{rqFFNotReexExeAllClaudeDen}{5}
\rqSet{rqFFNotReexExeAllClaudePct}{0.0}
\rqSet{rqFFNotReexExeAllCodexNum}{36}
\rqSet{rqFFNotReexExeAllCodexDen}{43}
\rqSet{rqFFNotReexExeAllCodexPct}{83.7}
\rqSet{rqFFNotReexExeAllCTwoRustNum}{0}
\rqSet{rqFFNotReexExeAllCTwoRustDen}{0}
\rqSet{rqFFNotReexExeAllCTwoRustPct}{n/a}
\rqSet{rqFFNotReexExeAllLaertesNum}{0}
\rqSet{rqFFNotReexExeAllLaertesDen}{0}
\rqSet{rqFFNotReexExeAllLaertesPct}{n/a}
\rqSet{rqFFNotReexExeAllCTwoSaferRustNum}{0}
\rqSet{rqFFNotReexExeAllCTwoSaferRustDen}{12}
\rqSet{rqFFNotReexExeAllCTwoSaferRustPct}{0.0}
\rqSet{rqFFNotReexExeAllSmartCTwoRustNum}{0}
\rqSet{rqFFNotReexExeAllSmartCTwoRustDen}{8}
\rqSet{rqFFNotReexExeAllSmartCTwoRustPct}{0.0}
\rqSet{rqFFNotReexExeAllKimiNum}{0}
\rqSet{rqFFNotReexExeAllKimiDen}{17}
\rqSet{rqFFNotReexExeAllKimiPct}{0.0}
\rqSet{rqFFNotReexExeAllGptFiveNum}{0}
\rqSet{rqFFNotReexExeAllGptFiveDen}{18}
\rqSet{rqFFNotReexExeAllGptFivePct}{0.0}
\rqSet{rqFFNotReexExeAllGeminiNum}{0}
\rqSet{rqFFNotReexExeAllGeminiDen}{15}
\rqSet{rqFFNotReexExeAllGeminiPct}{0.0}
\rqSet{rqFFNotReexExeBOneNum}{0}
\rqSet{rqFFNotReexExeBOneDen}{41}
\rqSet{rqFFNotReexExeBOnePct}{0.0}
\rqSet{rqFFNotReexExeBTwoNum}{0}
\rqSet{rqFFNotReexExeBTwoDen}{47}
\rqSet{rqFFNotReexExeBTwoPct}{0.0}
\rqSet{rqFFNotReexExePZeroNum}{0}
\rqSet{rqFFNotReexExePZeroDen}{4}
\rqSet{rqFFNotReexExePZeroPct}{0.0}
\rqSet{rqFFNotReexExePOneNum}{36}
\rqSet{rqFFNotReexExePOneDen}{36}
\rqSet{rqFFNotReexExePOnePct}{100.0}
\rqSet{rqFFHarnessErrorExeAllNum}{0}
\rqSet{rqFFHarnessErrorExeAllDen}{128}
\rqSet{rqFFHarnessErrorExeAllPct}{0.0}
\rqSet{rqFFHarnessErrorExeAllKiroNum}{0}
\rqSet{rqFFHarnessErrorExeAllKiroDen}{10}
\rqSet{rqFFHarnessErrorExeAllKiroPct}{0.0}
\rqSet{rqFFHarnessErrorExeAllClaudeNum}{0}
\rqSet{rqFFHarnessErrorExeAllClaudeDen}{5}
\rqSet{rqFFHarnessErrorExeAllClaudePct}{0.0}
\rqSet{rqFFHarnessErrorExeAllCodexNum}{0}
\rqSet{rqFFHarnessErrorExeAllCodexDen}{43}
\rqSet{rqFFHarnessErrorExeAllCodexPct}{0.0}
\rqSet{rqFFHarnessErrorExeAllCTwoRustNum}{0}
\rqSet{rqFFHarnessErrorExeAllCTwoRustDen}{0}
\rqSet{rqFFHarnessErrorExeAllCTwoRustPct}{n/a}
\rqSet{rqFFHarnessErrorExeAllLaertesNum}{0}
\rqSet{rqFFHarnessErrorExeAllLaertesDen}{0}
\rqSet{rqFFHarnessErrorExeAllLaertesPct}{n/a}
\rqSet{rqFFHarnessErrorExeAllCTwoSaferRustNum}{0}
\rqSet{rqFFHarnessErrorExeAllCTwoSaferRustDen}{12}
\rqSet{rqFFHarnessErrorExeAllCTwoSaferRustPct}{0.0}
\rqSet{rqFFHarnessErrorExeAllSmartCTwoRustNum}{0}
\rqSet{rqFFHarnessErrorExeAllSmartCTwoRustDen}{8}
\rqSet{rqFFHarnessErrorExeAllSmartCTwoRustPct}{0.0}
\rqSet{rqFFHarnessErrorExeAllKimiNum}{0}
\rqSet{rqFFHarnessErrorExeAllKimiDen}{17}
\rqSet{rqFFHarnessErrorExeAllKimiPct}{0.0}
\rqSet{rqFFHarnessErrorExeAllGptFiveNum}{0}
\rqSet{rqFFHarnessErrorExeAllGptFiveDen}{18}
\rqSet{rqFFHarnessErrorExeAllGptFivePct}{0.0}
\rqSet{rqFFHarnessErrorExeAllGeminiNum}{0}
\rqSet{rqFFHarnessErrorExeAllGeminiDen}{15}
\rqSet{rqFFHarnessErrorExeAllGeminiPct}{0.0}
\rqSet{rqFFHarnessErrorExeBOneNum}{0}
\rqSet{rqFFHarnessErrorExeBOneDen}{41}
\rqSet{rqFFHarnessErrorExeBOnePct}{0.0}
\rqSet{rqFFHarnessErrorExeBTwoNum}{0}
\rqSet{rqFFHarnessErrorExeBTwoDen}{47}
\rqSet{rqFFHarnessErrorExeBTwoPct}{0.0}
\rqSet{rqFFHarnessErrorExePZeroNum}{0}
\rqSet{rqFFHarnessErrorExePZeroDen}{4}
\rqSet{rqFFHarnessErrorExePZeroPct}{0.0}
\rqSet{rqFFHarnessErrorExePOneNum}{0}
\rqSet{rqFFHarnessErrorExePOneDen}{36}
\rqSet{rqFFHarnessErrorExePOnePct}{0.0}
\rqSet{rqFFRrunRecExeAllNum}{128}
\rqSet{rqFFRrunRecExeAllDen}{128}
\rqSet{rqFFRrunRecExeAllPct}{100.0}
\rqSet{rqFFRrunRecExeAllKiroNum}{10}
\rqSet{rqFFRrunRecExeAllKiroDen}{10}
\rqSet{rqFFRrunRecExeAllKiroPct}{100.0}
\rqSet{rqFFRrunRecExeAllClaudeNum}{5}
\rqSet{rqFFRrunRecExeAllClaudeDen}{5}
\rqSet{rqFFRrunRecExeAllClaudePct}{100.0}
\rqSet{rqFFRrunRecExeAllCodexNum}{43}
\rqSet{rqFFRrunRecExeAllCodexDen}{43}
\rqSet{rqFFRrunRecExeAllCodexPct}{100.0}
\rqSet{rqFFRrunRecExeAllCTwoRustNum}{0}
\rqSet{rqFFRrunRecExeAllCTwoRustDen}{0}
\rqSet{rqFFRrunRecExeAllCTwoRustPct}{n/a}
\rqSet{rqFFRrunRecExeAllLaertesNum}{0}
\rqSet{rqFFRrunRecExeAllLaertesDen}{0}
\rqSet{rqFFRrunRecExeAllLaertesPct}{n/a}
\rqSet{rqFFRrunRecExeAllCTwoSaferRustNum}{12}
\rqSet{rqFFRrunRecExeAllCTwoSaferRustDen}{12}
\rqSet{rqFFRrunRecExeAllCTwoSaferRustPct}{100.0}
\rqSet{rqFFRrunRecExeAllSmartCTwoRustNum}{8}
\rqSet{rqFFRrunRecExeAllSmartCTwoRustDen}{8}
\rqSet{rqFFRrunRecExeAllSmartCTwoRustPct}{100.0}
\rqSet{rqFFRrunRecExeAllKimiNum}{17}
\rqSet{rqFFRrunRecExeAllKimiDen}{17}
\rqSet{rqFFRrunRecExeAllKimiPct}{100.0}
\rqSet{rqFFRrunRecExeAllGptFiveNum}{18}
\rqSet{rqFFRrunRecExeAllGptFiveDen}{18}
\rqSet{rqFFRrunRecExeAllGptFivePct}{100.0}
\rqSet{rqFFRrunRecExeAllGeminiNum}{15}
\rqSet{rqFFRrunRecExeAllGeminiDen}{15}
\rqSet{rqFFRrunRecExeAllGeminiPct}{100.0}
\rqSet{rqFFRrunRecExeBOneNum}{41}
\rqSet{rqFFRrunRecExeBOneDen}{41}
\rqSet{rqFFRrunRecExeBOnePct}{100.0}
\rqSet{rqFFRrunRecExeBTwoNum}{47}
\rqSet{rqFFRrunRecExeBTwoDen}{47}
\rqSet{rqFFRrunRecExeBTwoPct}{100.0}
\rqSet{rqFFRrunRecExePZeroNum}{4}
\rqSet{rqFFRrunRecExePZeroDen}{4}
\rqSet{rqFFRrunRecExePZeroPct}{100.0}
\rqSet{rqFFRrunRecExePOneNum}{36}
\rqSet{rqFFRrunRecExePOneDen}{36}
\rqSet{rqFFRrunRecExePOnePct}{100.0}
\rqSet{rqRZeroTestRRunExeAllNum}{128}
\rqSet{rqRZeroTestRRunExeAllDen}{128}
\rqSet{rqRZeroTestRRunExeAllPct}{100.0}
\rqSet{rqRZeroTestRRunExeAllKiroNum}{10}
\rqSet{rqRZeroTestRRunExeAllKiroDen}{10}
\rqSet{rqRZeroTestRRunExeAllKiroPct}{100.0}
\rqSet{rqRZeroTestRRunExeAllClaudeNum}{5}
\rqSet{rqRZeroTestRRunExeAllClaudeDen}{5}
\rqSet{rqRZeroTestRRunExeAllClaudePct}{100.0}
\rqSet{rqRZeroTestRRunExeAllCodexNum}{43}
\rqSet{rqRZeroTestRRunExeAllCodexDen}{43}
\rqSet{rqRZeroTestRRunExeAllCodexPct}{100.0}
\rqSet{rqRZeroTestRRunExeAllCTwoRustNum}{0}
\rqSet{rqRZeroTestRRunExeAllCTwoRustDen}{0}
\rqSet{rqRZeroTestRRunExeAllCTwoRustPct}{n/a}
\rqSet{rqRZeroTestRRunExeAllLaertesNum}{0}
\rqSet{rqRZeroTestRRunExeAllLaertesDen}{0}
\rqSet{rqRZeroTestRRunExeAllLaertesPct}{n/a}
\rqSet{rqRZeroTestRRunExeAllCTwoSaferRustNum}{12}
\rqSet{rqRZeroTestRRunExeAllCTwoSaferRustDen}{12}
\rqSet{rqRZeroTestRRunExeAllCTwoSaferRustPct}{100.0}
\rqSet{rqRZeroTestRRunExeAllSmartCTwoRustNum}{8}
\rqSet{rqRZeroTestRRunExeAllSmartCTwoRustDen}{8}
\rqSet{rqRZeroTestRRunExeAllSmartCTwoRustPct}{100.0}
\rqSet{rqRZeroTestRRunExeAllKimiNum}{17}
\rqSet{rqRZeroTestRRunExeAllKimiDen}{17}
\rqSet{rqRZeroTestRRunExeAllKimiPct}{100.0}
\rqSet{rqRZeroTestRRunExeAllGptFiveNum}{18}
\rqSet{rqRZeroTestRRunExeAllGptFiveDen}{18}
\rqSet{rqRZeroTestRRunExeAllGptFivePct}{100.0}
\rqSet{rqRZeroTestRRunExeAllGeminiNum}{15}
\rqSet{rqRZeroTestRRunExeAllGeminiDen}{15}
\rqSet{rqRZeroTestRRunExeAllGeminiPct}{100.0}
\rqSet{rqRZeroTestRRunExeBOneNum}{41}
\rqSet{rqRZeroTestRRunExeBOneDen}{41}
\rqSet{rqRZeroTestRRunExeBOnePct}{100.0}
\rqSet{rqRZeroTestRRunExeBTwoNum}{47}
\rqSet{rqRZeroTestRRunExeBTwoDen}{47}
\rqSet{rqRZeroTestRRunExeBTwoPct}{100.0}
\rqSet{rqRZeroTestRRunExePZeroNum}{4}
\rqSet{rqRZeroTestRRunExePZeroDen}{4}
\rqSet{rqRZeroTestRRunExePZeroPct}{100.0}
\rqSet{rqRZeroTestRRunExePOneNum}{36}
\rqSet{rqRZeroTestRRunExePOneDen}{36}
\rqSet{rqRZeroTestRRunExePOnePct}{100.0}
\rqSet{rqRZeroTestRNonrepExeAllNum}{0}
\rqSet{rqRZeroTestRNonrepExeAllDen}{128}
\rqSet{rqRZeroTestRNonrepExeAllPct}{0.0}
\rqSet{rqRZeroTestRNonrepExeAllKiroNum}{0}
\rqSet{rqRZeroTestRNonrepExeAllKiroDen}{10}
\rqSet{rqRZeroTestRNonrepExeAllKiroPct}{0.0}
\rqSet{rqRZeroTestRNonrepExeAllClaudeNum}{0}
\rqSet{rqRZeroTestRNonrepExeAllClaudeDen}{5}
\rqSet{rqRZeroTestRNonrepExeAllClaudePct}{0.0}
\rqSet{rqRZeroTestRNonrepExeAllCodexNum}{0}
\rqSet{rqRZeroTestRNonrepExeAllCodexDen}{43}
\rqSet{rqRZeroTestRNonrepExeAllCodexPct}{0.0}
\rqSet{rqRZeroTestRNonrepExeAllCTwoRustNum}{0}
\rqSet{rqRZeroTestRNonrepExeAllCTwoRustDen}{0}
\rqSet{rqRZeroTestRNonrepExeAllCTwoRustPct}{n/a}
\rqSet{rqRZeroTestRNonrepExeAllLaertesNum}{0}
\rqSet{rqRZeroTestRNonrepExeAllLaertesDen}{0}
\rqSet{rqRZeroTestRNonrepExeAllLaertesPct}{n/a}
\rqSet{rqRZeroTestRNonrepExeAllCTwoSaferRustNum}{0}
\rqSet{rqRZeroTestRNonrepExeAllCTwoSaferRustDen}{12}
\rqSet{rqRZeroTestRNonrepExeAllCTwoSaferRustPct}{0.0}
\rqSet{rqRZeroTestRNonrepExeAllSmartCTwoRustNum}{0}
\rqSet{rqRZeroTestRNonrepExeAllSmartCTwoRustDen}{8}
\rqSet{rqRZeroTestRNonrepExeAllSmartCTwoRustPct}{0.0}
\rqSet{rqRZeroTestRNonrepExeAllKimiNum}{0}
\rqSet{rqRZeroTestRNonrepExeAllKimiDen}{17}
\rqSet{rqRZeroTestRNonrepExeAllKimiPct}{0.0}
\rqSet{rqRZeroTestRNonrepExeAllGptFiveNum}{0}
\rqSet{rqRZeroTestRNonrepExeAllGptFiveDen}{18}
\rqSet{rqRZeroTestRNonrepExeAllGptFivePct}{0.0}
\rqSet{rqRZeroTestRNonrepExeAllGeminiNum}{0}
\rqSet{rqRZeroTestRNonrepExeAllGeminiDen}{15}
\rqSet{rqRZeroTestRNonrepExeAllGeminiPct}{0.0}
\rqSet{rqRZeroTestRNonrepExeBOneNum}{0}
\rqSet{rqRZeroTestRNonrepExeBOneDen}{41}
\rqSet{rqRZeroTestRNonrepExeBOnePct}{0.0}
\rqSet{rqRZeroTestRNonrepExeBTwoNum}{0}
\rqSet{rqRZeroTestRNonrepExeBTwoDen}{47}
\rqSet{rqRZeroTestRNonrepExeBTwoPct}{0.0}
\rqSet{rqRZeroTestRNonrepExePZeroNum}{0}
\rqSet{rqRZeroTestRNonrepExePZeroDen}{4}
\rqSet{rqRZeroTestRNonrepExePZeroPct}{0.0}
\rqSet{rqRZeroTestRNonrepExePOneNum}{0}
\rqSet{rqRZeroTestRNonrepExePOneDen}{36}
\rqSet{rqRZeroTestRNonrepExePOnePct}{0.0}
\rqSet{rqRZeroTestRNobuildExeAllNum}{0}
\rqSet{rqRZeroTestRNobuildExeAllDen}{128}
\rqSet{rqRZeroTestRNobuildExeAllPct}{0.0}
\rqSet{rqRZeroTestRNobuildExeAllKiroNum}{0}
\rqSet{rqRZeroTestRNobuildExeAllKiroDen}{10}
\rqSet{rqRZeroTestRNobuildExeAllKiroPct}{0.0}
\rqSet{rqRZeroTestRNobuildExeAllClaudeNum}{0}
\rqSet{rqRZeroTestRNobuildExeAllClaudeDen}{5}
\rqSet{rqRZeroTestRNobuildExeAllClaudePct}{0.0}
\rqSet{rqRZeroTestRNobuildExeAllCodexNum}{0}
\rqSet{rqRZeroTestRNobuildExeAllCodexDen}{43}
\rqSet{rqRZeroTestRNobuildExeAllCodexPct}{0.0}
\rqSet{rqRZeroTestRNobuildExeAllCTwoRustNum}{0}
\rqSet{rqRZeroTestRNobuildExeAllCTwoRustDen}{0}
\rqSet{rqRZeroTestRNobuildExeAllCTwoRustPct}{n/a}
\rqSet{rqRZeroTestRNobuildExeAllLaertesNum}{0}
\rqSet{rqRZeroTestRNobuildExeAllLaertesDen}{0}
\rqSet{rqRZeroTestRNobuildExeAllLaertesPct}{n/a}
\rqSet{rqRZeroTestRNobuildExeAllCTwoSaferRustNum}{0}
\rqSet{rqRZeroTestRNobuildExeAllCTwoSaferRustDen}{12}
\rqSet{rqRZeroTestRNobuildExeAllCTwoSaferRustPct}{0.0}
\rqSet{rqRZeroTestRNobuildExeAllSmartCTwoRustNum}{0}
\rqSet{rqRZeroTestRNobuildExeAllSmartCTwoRustDen}{8}
\rqSet{rqRZeroTestRNobuildExeAllSmartCTwoRustPct}{0.0}
\rqSet{rqRZeroTestRNobuildExeAllKimiNum}{0}
\rqSet{rqRZeroTestRNobuildExeAllKimiDen}{17}
\rqSet{rqRZeroTestRNobuildExeAllKimiPct}{0.0}
\rqSet{rqRZeroTestRNobuildExeAllGptFiveNum}{0}
\rqSet{rqRZeroTestRNobuildExeAllGptFiveDen}{18}
\rqSet{rqRZeroTestRNobuildExeAllGptFivePct}{0.0}
\rqSet{rqRZeroTestRNobuildExeAllGeminiNum}{0}
\rqSet{rqRZeroTestRNobuildExeAllGeminiDen}{15}
\rqSet{rqRZeroTestRNobuildExeAllGeminiPct}{0.0}
\rqSet{rqRZeroTestRNobuildExeBOneNum}{0}
\rqSet{rqRZeroTestRNobuildExeBOneDen}{41}
\rqSet{rqRZeroTestRNobuildExeBOnePct}{0.0}
\rqSet{rqRZeroTestRNobuildExeBTwoNum}{0}
\rqSet{rqRZeroTestRNobuildExeBTwoDen}{47}
\rqSet{rqRZeroTestRNobuildExeBTwoPct}{0.0}
\rqSet{rqRZeroTestRNobuildExePZeroNum}{0}
\rqSet{rqRZeroTestRNobuildExePZeroDen}{4}
\rqSet{rqRZeroTestRNobuildExePZeroPct}{0.0}
\rqSet{rqRZeroTestRNobuildExePOneNum}{0}
\rqSet{rqRZeroTestRNobuildExePOneDen}{36}
\rqSet{rqRZeroTestRNobuildExePOnePct}{0.0}
\rqSet{rqRZeroTestROtherExeAllNum}{0}
\rqSet{rqRZeroTestROtherExeAllDen}{128}
\rqSet{rqRZeroTestROtherExeAllPct}{0.0}
\rqSet{rqRZeroTestROtherExeAllKiroNum}{0}
\rqSet{rqRZeroTestROtherExeAllKiroDen}{10}
\rqSet{rqRZeroTestROtherExeAllKiroPct}{0.0}
\rqSet{rqRZeroTestROtherExeAllClaudeNum}{0}
\rqSet{rqRZeroTestROtherExeAllClaudeDen}{5}
\rqSet{rqRZeroTestROtherExeAllClaudePct}{0.0}
\rqSet{rqRZeroTestROtherExeAllCodexNum}{0}
\rqSet{rqRZeroTestROtherExeAllCodexDen}{43}
\rqSet{rqRZeroTestROtherExeAllCodexPct}{0.0}
\rqSet{rqRZeroTestROtherExeAllCTwoRustNum}{0}
\rqSet{rqRZeroTestROtherExeAllCTwoRustDen}{0}
\rqSet{rqRZeroTestROtherExeAllCTwoRustPct}{n/a}
\rqSet{rqRZeroTestROtherExeAllLaertesNum}{0}
\rqSet{rqRZeroTestROtherExeAllLaertesDen}{0}
\rqSet{rqRZeroTestROtherExeAllLaertesPct}{n/a}
\rqSet{rqRZeroTestROtherExeAllCTwoSaferRustNum}{0}
\rqSet{rqRZeroTestROtherExeAllCTwoSaferRustDen}{12}
\rqSet{rqRZeroTestROtherExeAllCTwoSaferRustPct}{0.0}
\rqSet{rqRZeroTestROtherExeAllSmartCTwoRustNum}{0}
\rqSet{rqRZeroTestROtherExeAllSmartCTwoRustDen}{8}
\rqSet{rqRZeroTestROtherExeAllSmartCTwoRustPct}{0.0}
\rqSet{rqRZeroTestROtherExeAllKimiNum}{0}
\rqSet{rqRZeroTestROtherExeAllKimiDen}{17}
\rqSet{rqRZeroTestROtherExeAllKimiPct}{0.0}
\rqSet{rqRZeroTestROtherExeAllGptFiveNum}{0}
\rqSet{rqRZeroTestROtherExeAllGptFiveDen}{18}
\rqSet{rqRZeroTestROtherExeAllGptFivePct}{0.0}
\rqSet{rqRZeroTestROtherExeAllGeminiNum}{0}
\rqSet{rqRZeroTestROtherExeAllGeminiDen}{15}
\rqSet{rqRZeroTestROtherExeAllGeminiPct}{0.0}
\rqSet{rqRZeroTestROtherExeBOneNum}{0}
\rqSet{rqRZeroTestROtherExeBOneDen}{41}
\rqSet{rqRZeroTestROtherExeBOnePct}{0.0}
\rqSet{rqRZeroTestROtherExeBTwoNum}{0}
\rqSet{rqRZeroTestROtherExeBTwoDen}{47}
\rqSet{rqRZeroTestROtherExeBTwoPct}{0.0}
\rqSet{rqRZeroTestROtherExePZeroNum}{0}
\rqSet{rqRZeroTestROtherExePZeroDen}{4}
\rqSet{rqRZeroTestROtherExePZeroPct}{0.0}
\rqSet{rqRZeroTestROtherExePOneNum}{0}
\rqSet{rqRZeroTestROtherExePOneDen}{36}
\rqSet{rqRZeroTestROtherExePOnePct}{0.0}
\rqSet{rqLrExeAllNum}{36}
\rqSet{rqLrExeAllDen}{128}
\rqSet{rqLrExeAllPct}{28.1}
\rqSet{rqLrExeAllKiroNum}{0}
\rqSet{rqLrExeAllKiroDen}{10}
\rqSet{rqLrExeAllKiroPct}{0.0}
\rqSet{rqLrExeAllClaudeNum}{0}
\rqSet{rqLrExeAllClaudeDen}{5}
\rqSet{rqLrExeAllClaudePct}{0.0}
\rqSet{rqLrExeAllCodexNum}{36}
\rqSet{rqLrExeAllCodexDen}{43}
\rqSet{rqLrExeAllCodexPct}{83.7}
\rqSet{rqLrExeAllCTwoRustNum}{0}
\rqSet{rqLrExeAllCTwoRustDen}{0}
\rqSet{rqLrExeAllCTwoRustPct}{n/a}
\rqSet{rqLrExeAllLaertesNum}{0}
\rqSet{rqLrExeAllLaertesDen}{0}
\rqSet{rqLrExeAllLaertesPct}{n/a}
\rqSet{rqLrExeAllCTwoSaferRustNum}{0}
\rqSet{rqLrExeAllCTwoSaferRustDen}{12}
\rqSet{rqLrExeAllCTwoSaferRustPct}{0.0}
\rqSet{rqLrExeAllSmartCTwoRustNum}{0}
\rqSet{rqLrExeAllSmartCTwoRustDen}{8}
\rqSet{rqLrExeAllSmartCTwoRustPct}{0.0}
\rqSet{rqLrExeAllKimiNum}{0}
\rqSet{rqLrExeAllKimiDen}{17}
\rqSet{rqLrExeAllKimiPct}{0.0}
\rqSet{rqLrExeAllGptFiveNum}{0}
\rqSet{rqLrExeAllGptFiveDen}{18}
\rqSet{rqLrExeAllGptFivePct}{0.0}
\rqSet{rqLrExeAllGeminiNum}{0}
\rqSet{rqLrExeAllGeminiDen}{15}
\rqSet{rqLrExeAllGeminiPct}{0.0}
\rqSet{rqLrExeBOneNum}{0}
\rqSet{rqLrExeBOneDen}{41}
\rqSet{rqLrExeBOnePct}{0.0}
\rqSet{rqLrExeBOneKiroNum}{0}
\rqSet{rqLrExeBOneKiroDen}{0}
\rqSet{rqLrExeBOneKiroPct}{n/a}
\rqSet{rqLrExeBOneClaudeNum}{0}
\rqSet{rqLrExeBOneClaudeDen}{0}
\rqSet{rqLrExeBOneClaudePct}{n/a}
\rqSet{rqLrExeBOneCodexNum}{0}
\rqSet{rqLrExeBOneCodexDen}{1}
\rqSet{rqLrExeBOneCodexPct}{0.0}
\rqSet{rqLrExeBOneCTwoRustNum}{0}
\rqSet{rqLrExeBOneCTwoRustDen}{0}
\rqSet{rqLrExeBOneCTwoRustPct}{n/a}
\rqSet{rqLrExeBOneLaertesNum}{0}
\rqSet{rqLrExeBOneLaertesDen}{0}
\rqSet{rqLrExeBOneLaertesPct}{n/a}
\rqSet{rqLrExeBOneCTwoSaferRustNum}{0}
\rqSet{rqLrExeBOneCTwoSaferRustDen}{11}
\rqSet{rqLrExeBOneCTwoSaferRustPct}{0.0}
\rqSet{rqLrExeBOneSmartCTwoRustNum}{0}
\rqSet{rqLrExeBOneSmartCTwoRustDen}{3}
\rqSet{rqLrExeBOneSmartCTwoRustPct}{0.0}
\rqSet{rqLrExeBOneKimiNum}{0}
\rqSet{rqLrExeBOneKimiDen}{12}
\rqSet{rqLrExeBOneKimiPct}{0.0}
\rqSet{rqLrExeBOneGptFiveNum}{0}
\rqSet{rqLrExeBOneGptFiveDen}{8}
\rqSet{rqLrExeBOneGptFivePct}{0.0}
\rqSet{rqLrExeBOneGeminiNum}{0}
\rqSet{rqLrExeBOneGeminiDen}{6}
\rqSet{rqLrExeBOneGeminiPct}{0.0}
\rqSet{rqLrExeBTwoNum}{0}
\rqSet{rqLrExeBTwoDen}{47}
\rqSet{rqLrExeBTwoPct}{0.0}
\rqSet{rqLrExeBTwoKiroNum}{0}
\rqSet{rqLrExeBTwoKiroDen}{10}
\rqSet{rqLrExeBTwoKiroPct}{0.0}
\rqSet{rqLrExeBTwoClaudeNum}{0}
\rqSet{rqLrExeBTwoClaudeDen}{5}
\rqSet{rqLrExeBTwoClaudePct}{0.0}
\rqSet{rqLrExeBTwoCodexNum}{0}
\rqSet{rqLrExeBTwoCodexDen}{6}
\rqSet{rqLrExeBTwoCodexPct}{0.0}
\rqSet{rqLrExeBTwoCTwoRustNum}{0}
\rqSet{rqLrExeBTwoCTwoRustDen}{0}
\rqSet{rqLrExeBTwoCTwoRustPct}{n/a}
\rqSet{rqLrExeBTwoLaertesNum}{0}
\rqSet{rqLrExeBTwoLaertesDen}{0}
\rqSet{rqLrExeBTwoLaertesPct}{n/a}
\rqSet{rqLrExeBTwoCTwoSaferRustNum}{0}
\rqSet{rqLrExeBTwoCTwoSaferRustDen}{1}
\rqSet{rqLrExeBTwoCTwoSaferRustPct}{0.0}
\rqSet{rqLrExeBTwoSmartCTwoRustNum}{0}
\rqSet{rqLrExeBTwoSmartCTwoRustDen}{4}
\rqSet{rqLrExeBTwoSmartCTwoRustPct}{0.0}
\rqSet{rqLrExeBTwoKimiNum}{0}
\rqSet{rqLrExeBTwoKimiDen}{4}
\rqSet{rqLrExeBTwoKimiPct}{0.0}
\rqSet{rqLrExeBTwoGptFiveNum}{0}
\rqSet{rqLrExeBTwoGptFiveDen}{9}
\rqSet{rqLrExeBTwoGptFivePct}{0.0}
\rqSet{rqLrExeBTwoGeminiNum}{0}
\rqSet{rqLrExeBTwoGeminiDen}{8}
\rqSet{rqLrExeBTwoGeminiPct}{0.0}
\rqSet{rqLrExePZeroNum}{0}
\rqSet{rqLrExePZeroDen}{4}
\rqSet{rqLrExePZeroPct}{0.0}
\rqSet{rqLrExePZeroKiroNum}{0}
\rqSet{rqLrExePZeroKiroDen}{0}
\rqSet{rqLrExePZeroKiroPct}{n/a}
\rqSet{rqLrExePZeroClaudeNum}{0}
\rqSet{rqLrExePZeroClaudeDen}{0}
\rqSet{rqLrExePZeroClaudePct}{n/a}
\rqSet{rqLrExePZeroCodexNum}{0}
\rqSet{rqLrExePZeroCodexDen}{0}
\rqSet{rqLrExePZeroCodexPct}{n/a}
\rqSet{rqLrExePZeroCTwoRustNum}{0}
\rqSet{rqLrExePZeroCTwoRustDen}{0}
\rqSet{rqLrExePZeroCTwoRustPct}{n/a}
\rqSet{rqLrExePZeroLaertesNum}{0}
\rqSet{rqLrExePZeroLaertesDen}{0}
\rqSet{rqLrExePZeroLaertesPct}{n/a}
\rqSet{rqLrExePZeroCTwoSaferRustNum}{0}
\rqSet{rqLrExePZeroCTwoSaferRustDen}{0}
\rqSet{rqLrExePZeroCTwoSaferRustPct}{n/a}
\rqSet{rqLrExePZeroSmartCTwoRustNum}{0}
\rqSet{rqLrExePZeroSmartCTwoRustDen}{1}
\rqSet{rqLrExePZeroSmartCTwoRustPct}{0.0}
\rqSet{rqLrExePZeroKimiNum}{0}
\rqSet{rqLrExePZeroKimiDen}{1}
\rqSet{rqLrExePZeroKimiPct}{0.0}
\rqSet{rqLrExePZeroGptFiveNum}{0}
\rqSet{rqLrExePZeroGptFiveDen}{1}
\rqSet{rqLrExePZeroGptFivePct}{0.0}
\rqSet{rqLrExePZeroGeminiNum}{0}
\rqSet{rqLrExePZeroGeminiDen}{1}
\rqSet{rqLrExePZeroGeminiPct}{0.0}
\rqSet{rqLrExePOneNum}{36}
\rqSet{rqLrExePOneDen}{36}
\rqSet{rqLrExePOnePct}{100.0}
\rqSet{rqLrExePOneKiroNum}{0}
\rqSet{rqLrExePOneKiroDen}{0}
\rqSet{rqLrExePOneKiroPct}{n/a}
\rqSet{rqLrExePOneClaudeNum}{0}
\rqSet{rqLrExePOneClaudeDen}{0}
\rqSet{rqLrExePOneClaudePct}{n/a}
\rqSet{rqLrExePOneCodexNum}{36}
\rqSet{rqLrExePOneCodexDen}{36}
\rqSet{rqLrExePOneCodexPct}{100.0}
\rqSet{rqLrExePOneCTwoRustNum}{0}
\rqSet{rqLrExePOneCTwoRustDen}{0}
\rqSet{rqLrExePOneCTwoRustPct}{n/a}
\rqSet{rqLrExePOneLaertesNum}{0}
\rqSet{rqLrExePOneLaertesDen}{0}
\rqSet{rqLrExePOneLaertesPct}{n/a}
\rqSet{rqLrExePOneCTwoSaferRustNum}{0}
\rqSet{rqLrExePOneCTwoSaferRustDen}{0}
\rqSet{rqLrExePOneCTwoSaferRustPct}{n/a}
\rqSet{rqLrExePOneSmartCTwoRustNum}{0}
\rqSet{rqLrExePOneSmartCTwoRustDen}{0}
\rqSet{rqLrExePOneSmartCTwoRustPct}{n/a}
\rqSet{rqLrExePOneKimiNum}{0}
\rqSet{rqLrExePOneKimiDen}{0}
\rqSet{rqLrExePOneKimiPct}{n/a}
\rqSet{rqLrExePOneGptFiveNum}{0}
\rqSet{rqLrExePOneGptFiveDen}{0}
\rqSet{rqLrExePOneGptFivePct}{n/a}
\rqSet{rqLrExePOneGeminiNum}{0}
\rqSet{rqLrExePOneGeminiDen}{0}
\rqSet{rqLrExePOneGeminiPct}{n/a}
\rqSet{rqFFReasonNotAdjudicableUbInCExeAllNum}{0}
\rqSet{rqFFReasonNotAdjudicableUbInCExeAllDen}{0}
\rqSet{rqFFReasonNotAdjudicableUbInCExeAllPct}{n/a}
\rqSet{rqFFReasonNotAdjudicableUbInCExeBOneNum}{0}
\rqSet{rqFFReasonNotAdjudicableUbInCExeBOneDen}{0}
\rqSet{rqFFReasonNotAdjudicableUbInCExeBOnePct}{n/a}
\rqSet{rqFFReasonNotAdjudicableUbInCExeBTwoNum}{0}
\rqSet{rqFFReasonNotAdjudicableUbInCExeBTwoDen}{0}
\rqSet{rqFFReasonNotAdjudicableUbInCExeBTwoPct}{n/a}
\rqSet{rqFFReasonNotAdjudicableUbInCExePZeroNum}{0}
\rqSet{rqFFReasonNotAdjudicableUbInCExePZeroDen}{0}
\rqSet{rqFFReasonNotAdjudicableUbInCExePZeroPct}{n/a}
\rqSet{rqFFReasonNotAdjudicableUbInCExePOneNum}{0}
\rqSet{rqFFReasonNotAdjudicableUbInCExePOneDen}{0}
\rqSet{rqFFReasonNotAdjudicableUbInCExePOnePct}{n/a}
\rqSet{rqFFReasonNotAdjudicableBothBudgetExeAllNum}{0}
\rqSet{rqFFReasonNotAdjudicableBothBudgetExeAllDen}{0}
\rqSet{rqFFReasonNotAdjudicableBothBudgetExeAllPct}{n/a}
\rqSet{rqFFReasonNotAdjudicableBothBudgetExeBOneNum}{0}
\rqSet{rqFFReasonNotAdjudicableBothBudgetExeBOneDen}{0}
\rqSet{rqFFReasonNotAdjudicableBothBudgetExeBOnePct}{n/a}
\rqSet{rqFFReasonNotAdjudicableBothBudgetExeBTwoNum}{0}
\rqSet{rqFFReasonNotAdjudicableBothBudgetExeBTwoDen}{0}
\rqSet{rqFFReasonNotAdjudicableBothBudgetExeBTwoPct}{n/a}
\rqSet{rqFFReasonNotAdjudicableBothBudgetExePZeroNum}{0}
\rqSet{rqFFReasonNotAdjudicableBothBudgetExePZeroDen}{0}
\rqSet{rqFFReasonNotAdjudicableBothBudgetExePZeroPct}{n/a}
\rqSet{rqFFReasonNotAdjudicableBothBudgetExePOneNum}{0}
\rqSet{rqFFReasonNotAdjudicableBothBudgetExePOneDen}{0}
\rqSet{rqFFReasonNotAdjudicableBothBudgetExePOnePct}{n/a}
\rqSet{rqFFReasonNotAdjudicableSanitizerInconclusiveExeAllNum}{0}
\rqSet{rqFFReasonNotAdjudicableSanitizerInconclusiveExeAllDen}{0}
\rqSet{rqFFReasonNotAdjudicableSanitizerInconclusiveExeAllPct}{n/a}
\rqSet{rqFFReasonNotAdjudicableSanitizerInconclusiveExeBOneNum}{0}
\rqSet{rqFFReasonNotAdjudicableSanitizerInconclusiveExeBOneDen}{0}
\rqSet{rqFFReasonNotAdjudicableSanitizerInconclusiveExeBOnePct}{n/a}
\rqSet{rqFFReasonNotAdjudicableSanitizerInconclusiveExeBTwoNum}{0}
\rqSet{rqFFReasonNotAdjudicableSanitizerInconclusiveExeBTwoDen}{0}
\rqSet{rqFFReasonNotAdjudicableSanitizerInconclusiveExeBTwoPct}{n/a}
\rqSet{rqFFReasonNotAdjudicableSanitizerInconclusiveExePZeroNum}{0}
\rqSet{rqFFReasonNotAdjudicableSanitizerInconclusiveExePZeroDen}{0}
\rqSet{rqFFReasonNotAdjudicableSanitizerInconclusiveExePZeroPct}{n/a}
\rqSet{rqFFReasonNotAdjudicableSanitizerInconclusiveExePOneNum}{0}
\rqSet{rqFFReasonNotAdjudicableSanitizerInconclusiveExePOneDen}{0}
\rqSet{rqFFReasonNotAdjudicableSanitizerInconclusiveExePOnePct}{n/a}
\rqSet{rqFFReasonNotAdjudicableCUnstableExeAllNum}{0}
\rqSet{rqFFReasonNotAdjudicableCUnstableExeAllDen}{0}
\rqSet{rqFFReasonNotAdjudicableCUnstableExeAllPct}{n/a}
\rqSet{rqFFReasonNotAdjudicableCUnstableExeBOneNum}{0}
\rqSet{rqFFReasonNotAdjudicableCUnstableExeBOneDen}{0}
\rqSet{rqFFReasonNotAdjudicableCUnstableExeBOnePct}{n/a}
\rqSet{rqFFReasonNotAdjudicableCUnstableExeBTwoNum}{0}
\rqSet{rqFFReasonNotAdjudicableCUnstableExeBTwoDen}{0}
\rqSet{rqFFReasonNotAdjudicableCUnstableExeBTwoPct}{n/a}
\rqSet{rqFFReasonNotAdjudicableCUnstableExePZeroNum}{0}
\rqSet{rqFFReasonNotAdjudicableCUnstableExePZeroDen}{0}
\rqSet{rqFFReasonNotAdjudicableCUnstableExePZeroPct}{n/a}
\rqSet{rqFFReasonNotAdjudicableCUnstableExePOneNum}{0}
\rqSet{rqFFReasonNotAdjudicableCUnstableExePOneDen}{0}
\rqSet{rqFFReasonNotAdjudicableCUnstableExePOnePct}{n/a}
\rqSet{rqFFReasonNotAdjudicableUnstableComparisonExeAllNum}{0}
\rqSet{rqFFReasonNotAdjudicableUnstableComparisonExeAllDen}{0}
\rqSet{rqFFReasonNotAdjudicableUnstableComparisonExeAllPct}{n/a}
\rqSet{rqFFReasonNotAdjudicableUnstableComparisonExeBOneNum}{0}
\rqSet{rqFFReasonNotAdjudicableUnstableComparisonExeBOneDen}{0}
\rqSet{rqFFReasonNotAdjudicableUnstableComparisonExeBOnePct}{n/a}
\rqSet{rqFFReasonNotAdjudicableUnstableComparisonExeBTwoNum}{0}
\rqSet{rqFFReasonNotAdjudicableUnstableComparisonExeBTwoDen}{0}
\rqSet{rqFFReasonNotAdjudicableUnstableComparisonExeBTwoPct}{n/a}
\rqSet{rqFFReasonNotAdjudicableUnstableComparisonExePZeroNum}{0}
\rqSet{rqFFReasonNotAdjudicableUnstableComparisonExePZeroDen}{0}
\rqSet{rqFFReasonNotAdjudicableUnstableComparisonExePZeroPct}{n/a}
\rqSet{rqFFReasonNotAdjudicableUnstableComparisonExePOneNum}{0}
\rqSet{rqFFReasonNotAdjudicableUnstableComparisonExePOneDen}{0}
\rqSet{rqFFReasonNotAdjudicableUnstableComparisonExePOnePct}{n/a}
\rqSet{rqFFReasonNotReExaminedNoFailingVectorExeAllNum}{0}
\rqSet{rqFFReasonNotReExaminedNoFailingVectorExeAllDen}{36}
\rqSet{rqFFReasonNotReExaminedNoFailingVectorExeAllPct}{0.0}
\rqSet{rqFFReasonNotReExaminedNoFailingVectorExeBOneNum}{0}
\rqSet{rqFFReasonNotReExaminedNoFailingVectorExeBOneDen}{0}
\rqSet{rqFFReasonNotReExaminedNoFailingVectorExeBOnePct}{n/a}
\rqSet{rqFFReasonNotReExaminedNoFailingVectorExeBTwoNum}{0}
\rqSet{rqFFReasonNotReExaminedNoFailingVectorExeBTwoDen}{0}
\rqSet{rqFFReasonNotReExaminedNoFailingVectorExeBTwoPct}{n/a}
\rqSet{rqFFReasonNotReExaminedNoFailingVectorExePZeroNum}{0}
\rqSet{rqFFReasonNotReExaminedNoFailingVectorExePZeroDen}{0}
\rqSet{rqFFReasonNotReExaminedNoFailingVectorExePZeroPct}{n/a}
\rqSet{rqFFReasonNotReExaminedNoFailingVectorExePOneNum}{0}
\rqSet{rqFFReasonNotReExaminedNoFailingVectorExePOneDen}{36}
\rqSet{rqFFReasonNotReExaminedNoFailingVectorExePOnePct}{0.0}
\rqSet{rqFFReasonNotReExaminedReportMissingExeAllNum}{0}
\rqSet{rqFFReasonNotReExaminedReportMissingExeAllDen}{36}
\rqSet{rqFFReasonNotReExaminedReportMissingExeAllPct}{0.0}
\rqSet{rqFFReasonNotReExaminedReportMissingExeBOneNum}{0}
\rqSet{rqFFReasonNotReExaminedReportMissingExeBOneDen}{0}
\rqSet{rqFFReasonNotReExaminedReportMissingExeBOnePct}{n/a}
\rqSet{rqFFReasonNotReExaminedReportMissingExeBTwoNum}{0}
\rqSet{rqFFReasonNotReExaminedReportMissingExeBTwoDen}{0}
\rqSet{rqFFReasonNotReExaminedReportMissingExeBTwoPct}{n/a}
\rqSet{rqFFReasonNotReExaminedReportMissingExePZeroNum}{0}
\rqSet{rqFFReasonNotReExaminedReportMissingExePZeroDen}{0}
\rqSet{rqFFReasonNotReExaminedReportMissingExePZeroPct}{n/a}
\rqSet{rqFFReasonNotReExaminedReportMissingExePOneNum}{0}
\rqSet{rqFFReasonNotReExaminedReportMissingExePOneDen}{36}
\rqSet{rqFFReasonNotReExaminedReportMissingExePOnePct}{0.0}
\rqSet{rqFFReasonNotReExaminedTimeoutOnlyRulePendingExeAllNum}{0}
\rqSet{rqFFReasonNotReExaminedTimeoutOnlyRulePendingExeAllDen}{36}
\rqSet{rqFFReasonNotReExaminedTimeoutOnlyRulePendingExeAllPct}{0.0}
\rqSet{rqFFReasonNotReExaminedTimeoutOnlyRulePendingExeBOneNum}{0}
\rqSet{rqFFReasonNotReExaminedTimeoutOnlyRulePendingExeBOneDen}{0}
\rqSet{rqFFReasonNotReExaminedTimeoutOnlyRulePendingExeBOnePct}{n/a}
\rqSet{rqFFReasonNotReExaminedTimeoutOnlyRulePendingExeBTwoNum}{0}
\rqSet{rqFFReasonNotReExaminedTimeoutOnlyRulePendingExeBTwoDen}{0}
\rqSet{rqFFReasonNotReExaminedTimeoutOnlyRulePendingExeBTwoPct}{n/a}
\rqSet{rqFFReasonNotReExaminedTimeoutOnlyRulePendingExePZeroNum}{0}
\rqSet{rqFFReasonNotReExaminedTimeoutOnlyRulePendingExePZeroDen}{0}
\rqSet{rqFFReasonNotReExaminedTimeoutOnlyRulePendingExePZeroPct}{n/a}
\rqSet{rqFFReasonNotReExaminedTimeoutOnlyRulePendingExePOneNum}{0}
\rqSet{rqFFReasonNotReExaminedTimeoutOnlyRulePendingExePOneDen}{36}
\rqSet{rqFFReasonNotReExaminedTimeoutOnlyRulePendingExePOnePct}{0.0}
\rqSet{rqFFReasonNotReExaminedBuildFailedInHarnessExeAllNum}{0}
\rqSet{rqFFReasonNotReExaminedBuildFailedInHarnessExeAllDen}{36}
\rqSet{rqFFReasonNotReExaminedBuildFailedInHarnessExeAllPct}{0.0}
\rqSet{rqFFReasonNotReExaminedBuildFailedInHarnessExeBOneNum}{0}
\rqSet{rqFFReasonNotReExaminedBuildFailedInHarnessExeBOneDen}{0}
\rqSet{rqFFReasonNotReExaminedBuildFailedInHarnessExeBOnePct}{n/a}
\rqSet{rqFFReasonNotReExaminedBuildFailedInHarnessExeBTwoNum}{0}
\rqSet{rqFFReasonNotReExaminedBuildFailedInHarnessExeBTwoDen}{0}
\rqSet{rqFFReasonNotReExaminedBuildFailedInHarnessExeBTwoPct}{n/a}
\rqSet{rqFFReasonNotReExaminedBuildFailedInHarnessExePZeroNum}{0}
\rqSet{rqFFReasonNotReExaminedBuildFailedInHarnessExePZeroDen}{0}
\rqSet{rqFFReasonNotReExaminedBuildFailedInHarnessExePZeroPct}{n/a}
\rqSet{rqFFReasonNotReExaminedBuildFailedInHarnessExePOneNum}{0}
\rqSet{rqFFReasonNotReExaminedBuildFailedInHarnessExePOneDen}{36}
\rqSet{rqFFReasonNotReExaminedBuildFailedInHarnessExePOnePct}{0.0}
\rqSet{rqFFReasonNotReExaminedIncompleteHarnessRunVectorUnreadableExeAllNum}{0}
\rqSet{rqFFReasonNotReExaminedIncompleteHarnessRunVectorUnreadableExeAllDen}{36}
\rqSet{rqFFReasonNotReExaminedIncompleteHarnessRunVectorUnreadableExeAllPct}{0.0}
\rqSet{rqFFReasonNotReExaminedIncompleteHarnessRunVectorUnreadableExeBOneNum}{0}
\rqSet{rqFFReasonNotReExaminedIncompleteHarnessRunVectorUnreadableExeBOneDen}{0}
\rqSet{rqFFReasonNotReExaminedIncompleteHarnessRunVectorUnreadableExeBOnePct}{n/a}
\rqSet{rqFFReasonNotReExaminedIncompleteHarnessRunVectorUnreadableExeBTwoNum}{0}
\rqSet{rqFFReasonNotReExaminedIncompleteHarnessRunVectorUnreadableExeBTwoDen}{0}
\rqSet{rqFFReasonNotReExaminedIncompleteHarnessRunVectorUnreadableExeBTwoPct}{n/a}
\rqSet{rqFFReasonNotReExaminedIncompleteHarnessRunVectorUnreadableExePZeroNum}{0}
\rqSet{rqFFReasonNotReExaminedIncompleteHarnessRunVectorUnreadableExePZeroDen}{0}
\rqSet{rqFFReasonNotReExaminedIncompleteHarnessRunVectorUnreadableExePZeroPct}{n/a}
\rqSet{rqFFReasonNotReExaminedIncompleteHarnessRunVectorUnreadableExePOneNum}{0}
\rqSet{rqFFReasonNotReExaminedIncompleteHarnessRunVectorUnreadableExePOneDen}{36}
\rqSet{rqFFReasonNotReExaminedIncompleteHarnessRunVectorUnreadableExePOnePct}{0.0}
\rqSet{rqFFReasonNotReExaminedIncompleteHarnessRunInvalidInputExeAllNum}{0}
\rqSet{rqFFReasonNotReExaminedIncompleteHarnessRunInvalidInputExeAllDen}{36}
\rqSet{rqFFReasonNotReExaminedIncompleteHarnessRunInvalidInputExeAllPct}{0.0}
\rqSet{rqFFReasonNotReExaminedIncompleteHarnessRunInvalidInputExeBOneNum}{0}
\rqSet{rqFFReasonNotReExaminedIncompleteHarnessRunInvalidInputExeBOneDen}{0}
\rqSet{rqFFReasonNotReExaminedIncompleteHarnessRunInvalidInputExeBOnePct}{n/a}
\rqSet{rqFFReasonNotReExaminedIncompleteHarnessRunInvalidInputExeBTwoNum}{0}
\rqSet{rqFFReasonNotReExaminedIncompleteHarnessRunInvalidInputExeBTwoDen}{0}
\rqSet{rqFFReasonNotReExaminedIncompleteHarnessRunInvalidInputExeBTwoPct}{n/a}
\rqSet{rqFFReasonNotReExaminedIncompleteHarnessRunInvalidInputExePZeroNum}{0}
\rqSet{rqFFReasonNotReExaminedIncompleteHarnessRunInvalidInputExePZeroDen}{0}
\rqSet{rqFFReasonNotReExaminedIncompleteHarnessRunInvalidInputExePZeroPct}{n/a}
\rqSet{rqFFReasonNotReExaminedIncompleteHarnessRunInvalidInputExePOneNum}{0}
\rqSet{rqFFReasonNotReExaminedIncompleteHarnessRunInvalidInputExePOneDen}{36}
\rqSet{rqFFReasonNotReExaminedIncompleteHarnessRunInvalidInputExePOnePct}{0.0}
\rqSet{rqFFReasonNotReExaminedIncompleteHarnessRunRunUnobservedExeAllNum}{0}
\rqSet{rqFFReasonNotReExaminedIncompleteHarnessRunRunUnobservedExeAllDen}{36}
\rqSet{rqFFReasonNotReExaminedIncompleteHarnessRunRunUnobservedExeAllPct}{0.0}
\rqSet{rqFFReasonNotReExaminedIncompleteHarnessRunRunUnobservedExeBOneNum}{0}
\rqSet{rqFFReasonNotReExaminedIncompleteHarnessRunRunUnobservedExeBOneDen}{0}
\rqSet{rqFFReasonNotReExaminedIncompleteHarnessRunRunUnobservedExeBOnePct}{n/a}
\rqSet{rqFFReasonNotReExaminedIncompleteHarnessRunRunUnobservedExeBTwoNum}{0}
\rqSet{rqFFReasonNotReExaminedIncompleteHarnessRunRunUnobservedExeBTwoDen}{0}
\rqSet{rqFFReasonNotReExaminedIncompleteHarnessRunRunUnobservedExeBTwoPct}{n/a}
\rqSet{rqFFReasonNotReExaminedIncompleteHarnessRunRunUnobservedExePZeroNum}{0}
\rqSet{rqFFReasonNotReExaminedIncompleteHarnessRunRunUnobservedExePZeroDen}{0}
\rqSet{rqFFReasonNotReExaminedIncompleteHarnessRunRunUnobservedExePZeroPct}{n/a}
\rqSet{rqFFReasonNotReExaminedIncompleteHarnessRunRunUnobservedExePOneNum}{0}
\rqSet{rqFFReasonNotReExaminedIncompleteHarnessRunRunUnobservedExePOneDen}{36}
\rqSet{rqFFReasonNotReExaminedIncompleteHarnessRunRunUnobservedExePOnePct}{0.0}
\rqSet{rqFFReasonNotReExaminedIncompleteHarnessRunSanBuildFailedExeAllNum}{36}
\rqSet{rqFFReasonNotReExaminedIncompleteHarnessRunSanBuildFailedExeAllDen}{36}
\rqSet{rqFFReasonNotReExaminedIncompleteHarnessRunSanBuildFailedExeAllPct}{100.0}
\rqSet{rqFFReasonNotReExaminedIncompleteHarnessRunSanBuildFailedExeBOneNum}{0}
\rqSet{rqFFReasonNotReExaminedIncompleteHarnessRunSanBuildFailedExeBOneDen}{0}
\rqSet{rqFFReasonNotReExaminedIncompleteHarnessRunSanBuildFailedExeBOnePct}{n/a}
\rqSet{rqFFReasonNotReExaminedIncompleteHarnessRunSanBuildFailedExeBTwoNum}{0}
\rqSet{rqFFReasonNotReExaminedIncompleteHarnessRunSanBuildFailedExeBTwoDen}{0}
\rqSet{rqFFReasonNotReExaminedIncompleteHarnessRunSanBuildFailedExeBTwoPct}{n/a}
\rqSet{rqFFReasonNotReExaminedIncompleteHarnessRunSanBuildFailedExePZeroNum}{0}
\rqSet{rqFFReasonNotReExaminedIncompleteHarnessRunSanBuildFailedExePZeroDen}{0}
\rqSet{rqFFReasonNotReExaminedIncompleteHarnessRunSanBuildFailedExePZeroPct}{n/a}
\rqSet{rqFFReasonNotReExaminedIncompleteHarnessRunSanBuildFailedExePOneNum}{36}
\rqSet{rqFFReasonNotReExaminedIncompleteHarnessRunSanBuildFailedExePOneDen}{36}
\rqSet{rqFFReasonNotReExaminedIncompleteHarnessRunSanBuildFailedExePOnePct}{100.0}
\rqSet{rqFFReasonNotReExaminedIncompleteHarnessRunSanInfraExeAllNum}{0}
\rqSet{rqFFReasonNotReExaminedIncompleteHarnessRunSanInfraExeAllDen}{36}
\rqSet{rqFFReasonNotReExaminedIncompleteHarnessRunSanInfraExeAllPct}{0.0}
\rqSet{rqFFReasonNotReExaminedIncompleteHarnessRunSanInfraExeBOneNum}{0}
\rqSet{rqFFReasonNotReExaminedIncompleteHarnessRunSanInfraExeBOneDen}{0}
\rqSet{rqFFReasonNotReExaminedIncompleteHarnessRunSanInfraExeBOnePct}{n/a}
\rqSet{rqFFReasonNotReExaminedIncompleteHarnessRunSanInfraExeBTwoNum}{0}
\rqSet{rqFFReasonNotReExaminedIncompleteHarnessRunSanInfraExeBTwoDen}{0}
\rqSet{rqFFReasonNotReExaminedIncompleteHarnessRunSanInfraExeBTwoPct}{n/a}
\rqSet{rqFFReasonNotReExaminedIncompleteHarnessRunSanInfraExePZeroNum}{0}
\rqSet{rqFFReasonNotReExaminedIncompleteHarnessRunSanInfraExePZeroDen}{0}
\rqSet{rqFFReasonNotReExaminedIncompleteHarnessRunSanInfraExePZeroPct}{n/a}
\rqSet{rqFFReasonNotReExaminedIncompleteHarnessRunSanInfraExePOneNum}{0}
\rqSet{rqFFReasonNotReExaminedIncompleteHarnessRunSanInfraExePOneDen}{36}
\rqSet{rqFFReasonNotReExaminedIncompleteHarnessRunSanInfraExePOnePct}{0.0}
\rqSet{rqFFChanStderrExeAllNum}{0}
\rqSet{rqFFChanStderrExeAllDen}{0}
\rqSet{rqFFChanStderrExeAllPct}{n/a}
\rqSet{rqFFChanStderrExeBOneNum}{0}
\rqSet{rqFFChanStderrExeBOneDen}{0}
\rqSet{rqFFChanStderrExeBOnePct}{n/a}
\rqSet{rqFFChanStderrExeBTwoNum}{0}
\rqSet{rqFFChanStderrExeBTwoDen}{0}
\rqSet{rqFFChanStderrExeBTwoPct}{n/a}
\rqSet{rqFFChanStderrExePZeroNum}{0}
\rqSet{rqFFChanStderrExePZeroDen}{0}
\rqSet{rqFFChanStderrExePZeroPct}{n/a}
\rqSet{rqFFChanStderrExePOneNum}{0}
\rqSet{rqFFChanStderrExePOneDen}{0}
\rqSet{rqFFChanStderrExePOnePct}{n/a}
\rqSet{rqFFChanStdoutExeAllNum}{0}
\rqSet{rqFFChanStdoutExeAllDen}{0}
\rqSet{rqFFChanStdoutExeAllPct}{n/a}
\rqSet{rqFFChanStdoutExeBOneNum}{0}
\rqSet{rqFFChanStdoutExeBOneDen}{0}
\rqSet{rqFFChanStdoutExeBOnePct}{n/a}
\rqSet{rqFFChanStdoutExeBTwoNum}{0}
\rqSet{rqFFChanStdoutExeBTwoDen}{0}
\rqSet{rqFFChanStdoutExeBTwoPct}{n/a}
\rqSet{rqFFChanStdoutExePZeroNum}{0}
\rqSet{rqFFChanStdoutExePZeroDen}{0}
\rqSet{rqFFChanStdoutExePZeroPct}{n/a}
\rqSet{rqFFChanStdoutExePOneNum}{0}
\rqSet{rqFFChanStdoutExePOneDen}{0}
\rqSet{rqFFChanStdoutExePOnePct}{n/a}
\rqSet{rqFFChanStateFloatExeAllNum}{0}
\rqSet{rqFFChanStateFloatExeAllDen}{0}
\rqSet{rqFFChanStateFloatExeAllPct}{n/a}
\rqSet{rqFFChanStateFloatExeBOneNum}{0}
\rqSet{rqFFChanStateFloatExeBOneDen}{0}
\rqSet{rqFFChanStateFloatExeBOnePct}{n/a}
\rqSet{rqFFChanStateFloatExeBTwoNum}{0}
\rqSet{rqFFChanStateFloatExeBTwoDen}{0}
\rqSet{rqFFChanStateFloatExeBTwoPct}{n/a}
\rqSet{rqFFChanStateFloatExePZeroNum}{0}
\rqSet{rqFFChanStateFloatExePZeroDen}{0}
\rqSet{rqFFChanStateFloatExePZeroPct}{n/a}
\rqSet{rqFFChanStateFloatExePOneNum}{0}
\rqSet{rqFFChanStateFloatExePOneDen}{0}
\rqSet{rqFFChanStateFloatExePOnePct}{n/a}
\rqSet{rqFFChanStateOtherExeAllNum}{0}
\rqSet{rqFFChanStateOtherExeAllDen}{0}
\rqSet{rqFFChanStateOtherExeAllPct}{n/a}
\rqSet{rqFFChanStateOtherExeBOneNum}{0}
\rqSet{rqFFChanStateOtherExeBOneDen}{0}
\rqSet{rqFFChanStateOtherExeBOnePct}{n/a}
\rqSet{rqFFChanStateOtherExeBTwoNum}{0}
\rqSet{rqFFChanStateOtherExeBTwoDen}{0}
\rqSet{rqFFChanStateOtherExeBTwoPct}{n/a}
\rqSet{rqFFChanStateOtherExePZeroNum}{0}
\rqSet{rqFFChanStateOtherExePZeroDen}{0}
\rqSet{rqFFChanStateOtherExePZeroPct}{n/a}
\rqSet{rqFFChanStateOtherExePOneNum}{0}
\rqSet{rqFFChanStateOtherExePOneDen}{0}
\rqSet{rqFFChanStateOtherExePOnePct}{n/a}
\rqSet{rqFFChanStateExeAllNum}{0}
\rqSet{rqFFChanStateExeAllDen}{0}
\rqSet{rqFFChanStateExeAllPct}{n/a}
\rqSet{rqFFChanStateExeBOneNum}{0}
\rqSet{rqFFChanStateExeBOneDen}{0}
\rqSet{rqFFChanStateExeBOnePct}{n/a}
\rqSet{rqFFChanStateExeBTwoNum}{0}
\rqSet{rqFFChanStateExeBTwoDen}{0}
\rqSet{rqFFChanStateExeBTwoPct}{n/a}
\rqSet{rqFFChanStateExePZeroNum}{0}
\rqSet{rqFFChanStateExePZeroDen}{0}
\rqSet{rqFFChanStateExePZeroPct}{n/a}
\rqSet{rqFFChanStateExePOneNum}{0}
\rqSet{rqFFChanStateExePOneDen}{0}
\rqSet{rqFFChanStateExePOnePct}{n/a}
\rqSet{rqFFChanResultPanicExeAllNum}{0}
\rqSet{rqFFChanResultPanicExeAllDen}{0}
\rqSet{rqFFChanResultPanicExeAllPct}{n/a}
\rqSet{rqFFChanResultPanicExeBOneNum}{0}
\rqSet{rqFFChanResultPanicExeBOneDen}{0}
\rqSet{rqFFChanResultPanicExeBOnePct}{n/a}
\rqSet{rqFFChanResultPanicExeBTwoNum}{0}
\rqSet{rqFFChanResultPanicExeBTwoDen}{0}
\rqSet{rqFFChanResultPanicExeBTwoPct}{n/a}
\rqSet{rqFFChanResultPanicExePZeroNum}{0}
\rqSet{rqFFChanResultPanicExePZeroDen}{0}
\rqSet{rqFFChanResultPanicExePZeroPct}{n/a}
\rqSet{rqFFChanResultPanicExePOneNum}{0}
\rqSet{rqFFChanResultPanicExePOneDen}{0}
\rqSet{rqFFChanResultPanicExePOnePct}{n/a}
\rqSet{rqFFChanResultSegmentationFaultExeAllNum}{0}
\rqSet{rqFFChanResultSegmentationFaultExeAllDen}{0}
\rqSet{rqFFChanResultSegmentationFaultExeAllPct}{n/a}
\rqSet{rqFFChanResultSegmentationFaultExeBOneNum}{0}
\rqSet{rqFFChanResultSegmentationFaultExeBOneDen}{0}
\rqSet{rqFFChanResultSegmentationFaultExeBOnePct}{n/a}
\rqSet{rqFFChanResultSegmentationFaultExeBTwoNum}{0}
\rqSet{rqFFChanResultSegmentationFaultExeBTwoDen}{0}
\rqSet{rqFFChanResultSegmentationFaultExeBTwoPct}{n/a}
\rqSet{rqFFChanResultSegmentationFaultExePZeroNum}{0}
\rqSet{rqFFChanResultSegmentationFaultExePZeroDen}{0}
\rqSet{rqFFChanResultSegmentationFaultExePZeroPct}{n/a}
\rqSet{rqFFChanResultSegmentationFaultExePOneNum}{0}
\rqSet{rqFFChanResultSegmentationFaultExePOneDen}{0}
\rqSet{rqFFChanResultSegmentationFaultExePOnePct}{n/a}
\rqSet{rqFFChanResultUnknownFailureExeAllNum}{0}
\rqSet{rqFFChanResultUnknownFailureExeAllDen}{0}
\rqSet{rqFFChanResultUnknownFailureExeAllPct}{n/a}
\rqSet{rqFFChanResultUnknownFailureExeBOneNum}{0}
\rqSet{rqFFChanResultUnknownFailureExeBOneDen}{0}
\rqSet{rqFFChanResultUnknownFailureExeBOnePct}{n/a}
\rqSet{rqFFChanResultUnknownFailureExeBTwoNum}{0}
\rqSet{rqFFChanResultUnknownFailureExeBTwoDen}{0}
\rqSet{rqFFChanResultUnknownFailureExeBTwoPct}{n/a}
\rqSet{rqFFChanResultUnknownFailureExePZeroNum}{0}
\rqSet{rqFFChanResultUnknownFailureExePZeroDen}{0}
\rqSet{rqFFChanResultUnknownFailureExePZeroPct}{n/a}
\rqSet{rqFFChanResultUnknownFailureExePOneNum}{0}
\rqSet{rqFFChanResultUnknownFailureExePOneDen}{0}
\rqSet{rqFFChanResultUnknownFailureExePOnePct}{n/a}
\rqSet{rqFFChanResultTimeoutExeAllNum}{0}
\rqSet{rqFFChanResultTimeoutExeAllDen}{0}
\rqSet{rqFFChanResultTimeoutExeAllPct}{n/a}
\rqSet{rqFFChanResultTimeoutExeBOneNum}{0}
\rqSet{rqFFChanResultTimeoutExeBOneDen}{0}
\rqSet{rqFFChanResultTimeoutExeBOnePct}{n/a}
\rqSet{rqFFChanResultTimeoutExeBTwoNum}{0}
\rqSet{rqFFChanResultTimeoutExeBTwoDen}{0}
\rqSet{rqFFChanResultTimeoutExeBTwoPct}{n/a}
\rqSet{rqFFChanResultTimeoutExePZeroNum}{0}
\rqSet{rqFFChanResultTimeoutExePZeroDen}{0}
\rqSet{rqFFChanResultTimeoutExePZeroPct}{n/a}
\rqSet{rqFFChanResultTimeoutExePOneNum}{0}
\rqSet{rqFFChanResultTimeoutExePOneDen}{0}
\rqSet{rqFFChanResultTimeoutExePOnePct}{n/a}
\rqSet{rqFFCTermAnyExeAllNum}{0}
\rqSet{rqFFCTermAnyExeAllDen}{810}
\rqSet{rqFFCTermAnyExeAllPct}{0.0}
\rqSet{rqFFCTermAnyExeBOneNum}{0}
\rqSet{rqFFCTermAnyExeBOneDen}{146}
\rqSet{rqFFCTermAnyExeBOnePct}{0.0}
\rqSet{rqFFCTermAnyExeBTwoNum}{0}
\rqSet{rqFFCTermAnyExeBTwoDen}{529}
\rqSet{rqFFCTermAnyExeBTwoPct}{0.0}
\rqSet{rqFFCTermAnyExePZeroNum}{0}
\rqSet{rqFFCTermAnyExePZeroDen}{99}
\rqSet{rqFFCTermAnyExePZeroPct}{0.0}
\rqSet{rqFFCTermAnyExePOneNum}{0}
\rqSet{rqFFCTermAnyExePOneDen}{36}
\rqSet{rqFFCTermAnyExePOnePct}{0.0}
\rqSet{rqFFCTermExitOneZeroOneExeAllNum}{0}
\rqSet{rqFFCTermExitOneZeroOneExeAllDen}{810}
\rqSet{rqFFCTermExitOneZeroOneExeAllPct}{0.0}
\rqSet{rqFFCTermExitOneZeroOneExeBOneNum}{0}
\rqSet{rqFFCTermExitOneZeroOneExeBOneDen}{146}
\rqSet{rqFFCTermExitOneZeroOneExeBOnePct}{0.0}
\rqSet{rqFFCTermExitOneZeroOneExeBTwoNum}{0}
\rqSet{rqFFCTermExitOneZeroOneExeBTwoDen}{529}
\rqSet{rqFFCTermExitOneZeroOneExeBTwoPct}{0.0}
\rqSet{rqFFCTermExitOneZeroOneExePZeroNum}{0}
\rqSet{rqFFCTermExitOneZeroOneExePZeroDen}{99}
\rqSet{rqFFCTermExitOneZeroOneExePZeroPct}{0.0}
\rqSet{rqFFCTermExitOneZeroOneExePOneNum}{0}
\rqSet{rqFFCTermExitOneZeroOneExePOneDen}{36}
\rqSet{rqFFCTermExitOneZeroOneExePOnePct}{0.0}
\rqSet{rqFFCTermSignalExeAllNum}{0}
\rqSet{rqFFCTermSignalExeAllDen}{810}
\rqSet{rqFFCTermSignalExeAllPct}{0.0}
\rqSet{rqFFCTermSignalExeBOneNum}{0}
\rqSet{rqFFCTermSignalExeBOneDen}{146}
\rqSet{rqFFCTermSignalExeBOnePct}{0.0}
\rqSet{rqFFCTermSignalExeBTwoNum}{0}
\rqSet{rqFFCTermSignalExeBTwoDen}{529}
\rqSet{rqFFCTermSignalExeBTwoPct}{0.0}
\rqSet{rqFFCTermSignalExePZeroNum}{0}
\rqSet{rqFFCTermSignalExePZeroDen}{99}
\rqSet{rqFFCTermSignalExePZeroPct}{0.0}
\rqSet{rqFFCTermSignalExePOneNum}{0}
\rqSet{rqFFCTermSignalExePOneDen}{36}
\rqSet{rqFFCTermSignalExePOnePct}{0.0}
\rqSet{rqFFCTermFFAnyExeAllNum}{0}
\rqSet{rqFFCTermFFAnyExeAllDen}{0}
\rqSet{rqFFCTermFFAnyExeAllPct}{n/a}
\rqSet{rqFFCTermFFAnyExeBOneNum}{0}
\rqSet{rqFFCTermFFAnyExeBOneDen}{0}
\rqSet{rqFFCTermFFAnyExeBOnePct}{n/a}
\rqSet{rqFFCTermFFAnyExeBTwoNum}{0}
\rqSet{rqFFCTermFFAnyExeBTwoDen}{0}
\rqSet{rqFFCTermFFAnyExeBTwoPct}{n/a}
\rqSet{rqFFCTermFFAnyExePZeroNum}{0}
\rqSet{rqFFCTermFFAnyExePZeroDen}{0}
\rqSet{rqFFCTermFFAnyExePZeroPct}{n/a}
\rqSet{rqFFCTermFFAnyExePOneNum}{0}
\rqSet{rqFFCTermFFAnyExePOneDen}{0}
\rqSet{rqFFCTermFFAnyExePOnePct}{n/a}
\rqSet{rqFFCTermFFExitOneZeroOneExeAllNum}{0}
\rqSet{rqFFCTermFFExitOneZeroOneExeAllDen}{0}
\rqSet{rqFFCTermFFExitOneZeroOneExeAllPct}{n/a}
\rqSet{rqFFCTermFFExitOneZeroOneExeBOneNum}{0}
\rqSet{rqFFCTermFFExitOneZeroOneExeBOneDen}{0}
\rqSet{rqFFCTermFFExitOneZeroOneExeBOnePct}{n/a}
\rqSet{rqFFCTermFFExitOneZeroOneExeBTwoNum}{0}
\rqSet{rqFFCTermFFExitOneZeroOneExeBTwoDen}{0}
\rqSet{rqFFCTermFFExitOneZeroOneExeBTwoPct}{n/a}
\rqSet{rqFFCTermFFExitOneZeroOneExePZeroNum}{0}
\rqSet{rqFFCTermFFExitOneZeroOneExePZeroDen}{0}
\rqSet{rqFFCTermFFExitOneZeroOneExePZeroPct}{n/a}
\rqSet{rqFFCTermFFExitOneZeroOneExePOneNum}{0}
\rqSet{rqFFCTermFFExitOneZeroOneExePOneDen}{0}
\rqSet{rqFFCTermFFExitOneZeroOneExePOnePct}{n/a}
\rqSet{rqFFCTermFFSignalExeAllNum}{0}
\rqSet{rqFFCTermFFSignalExeAllDen}{0}
\rqSet{rqFFCTermFFSignalExeAllPct}{n/a}
\rqSet{rqFFCTermFFSignalExeBOneNum}{0}
\rqSet{rqFFCTermFFSignalExeBOneDen}{0}
\rqSet{rqFFCTermFFSignalExeBOnePct}{n/a}
\rqSet{rqFFCTermFFSignalExeBTwoNum}{0}
\rqSet{rqFFCTermFFSignalExeBTwoDen}{0}
\rqSet{rqFFCTermFFSignalExeBTwoPct}{n/a}
\rqSet{rqFFCTermFFSignalExePZeroNum}{0}
\rqSet{rqFFCTermFFSignalExePZeroDen}{0}
\rqSet{rqFFCTermFFSignalExePZeroPct}{n/a}
\rqSet{rqFFCTermFFSignalExePOneNum}{0}
\rqSet{rqFFCTermFFSignalExePOneDen}{0}
\rqSet{rqFFCTermFFSignalExePOnePct}{n/a}
\rqSet{rqFFVecSupportsRejectionExeAllNum}{774}
\rqSet{rqFFVecSupportsRejectionExeAllDen}{810}
\rqSet{rqFFVecSupportsRejectionExeAllPct}{95.6}
\rqSet{rqFFVecSupportsRejectionExeBOneNum}{146}
\rqSet{rqFFVecSupportsRejectionExeBOneDen}{146}
\rqSet{rqFFVecSupportsRejectionExeBOnePct}{100.0}
\rqSet{rqFFVecSupportsRejectionExeBTwoNum}{529}
\rqSet{rqFFVecSupportsRejectionExeBTwoDen}{529}
\rqSet{rqFFVecSupportsRejectionExeBTwoPct}{100.0}
\rqSet{rqFFVecSupportsRejectionExePZeroNum}{99}
\rqSet{rqFFVecSupportsRejectionExePZeroDen}{99}
\rqSet{rqFFVecSupportsRejectionExePZeroPct}{100.0}
\rqSet{rqFFVecSupportsRejectionExePOneNum}{0}
\rqSet{rqFFVecSupportsRejectionExePOneDen}{36}
\rqSet{rqFFVecSupportsRejectionExePOnePct}{0.0}
\rqSet{rqFFVecAgreesExeAllNum}{0}
\rqSet{rqFFVecAgreesExeAllDen}{810}
\rqSet{rqFFVecAgreesExeAllPct}{0.0}
\rqSet{rqFFVecAgreesExeBOneNum}{0}
\rqSet{rqFFVecAgreesExeBOneDen}{146}
\rqSet{rqFFVecAgreesExeBOnePct}{0.0}
\rqSet{rqFFVecAgreesExeBTwoNum}{0}
\rqSet{rqFFVecAgreesExeBTwoDen}{529}
\rqSet{rqFFVecAgreesExeBTwoPct}{0.0}
\rqSet{rqFFVecAgreesExePZeroNum}{0}
\rqSet{rqFFVecAgreesExePZeroDen}{99}
\rqSet{rqFFVecAgreesExePZeroPct}{0.0}
\rqSet{rqFFVecAgreesExePOneNum}{0}
\rqSet{rqFFVecAgreesExePOneDen}{36}
\rqSet{rqFFVecAgreesExePOnePct}{0.0}
\rqSet{rqFFVecNotAdjudicableUbInCExeAllNum}{0}
\rqSet{rqFFVecNotAdjudicableUbInCExeAllDen}{810}
\rqSet{rqFFVecNotAdjudicableUbInCExeAllPct}{0.0}
\rqSet{rqFFVecNotAdjudicableUbInCExeBOneNum}{0}
\rqSet{rqFFVecNotAdjudicableUbInCExeBOneDen}{146}
\rqSet{rqFFVecNotAdjudicableUbInCExeBOnePct}{0.0}
\rqSet{rqFFVecNotAdjudicableUbInCExeBTwoNum}{0}
\rqSet{rqFFVecNotAdjudicableUbInCExeBTwoDen}{529}
\rqSet{rqFFVecNotAdjudicableUbInCExeBTwoPct}{0.0}
\rqSet{rqFFVecNotAdjudicableUbInCExePZeroNum}{0}
\rqSet{rqFFVecNotAdjudicableUbInCExePZeroDen}{99}
\rqSet{rqFFVecNotAdjudicableUbInCExePZeroPct}{0.0}
\rqSet{rqFFVecNotAdjudicableUbInCExePOneNum}{0}
\rqSet{rqFFVecNotAdjudicableUbInCExePOneDen}{36}
\rqSet{rqFFVecNotAdjudicableUbInCExePOnePct}{0.0}
\rqSet{rqFFVecNotAdjudicableBothBudgetExeAllNum}{0}
\rqSet{rqFFVecNotAdjudicableBothBudgetExeAllDen}{810}
\rqSet{rqFFVecNotAdjudicableBothBudgetExeAllPct}{0.0}
\rqSet{rqFFVecNotAdjudicableBothBudgetExeBOneNum}{0}
\rqSet{rqFFVecNotAdjudicableBothBudgetExeBOneDen}{146}
\rqSet{rqFFVecNotAdjudicableBothBudgetExeBOnePct}{0.0}
\rqSet{rqFFVecNotAdjudicableBothBudgetExeBTwoNum}{0}
\rqSet{rqFFVecNotAdjudicableBothBudgetExeBTwoDen}{529}
\rqSet{rqFFVecNotAdjudicableBothBudgetExeBTwoPct}{0.0}
\rqSet{rqFFVecNotAdjudicableBothBudgetExePZeroNum}{0}
\rqSet{rqFFVecNotAdjudicableBothBudgetExePZeroDen}{99}
\rqSet{rqFFVecNotAdjudicableBothBudgetExePZeroPct}{0.0}
\rqSet{rqFFVecNotAdjudicableBothBudgetExePOneNum}{0}
\rqSet{rqFFVecNotAdjudicableBothBudgetExePOneDen}{36}
\rqSet{rqFFVecNotAdjudicableBothBudgetExePOnePct}{0.0}
\rqSet{rqFFVecNotAdjudicableSanitizerInconclusiveExeAllNum}{0}
\rqSet{rqFFVecNotAdjudicableSanitizerInconclusiveExeAllDen}{810}
\rqSet{rqFFVecNotAdjudicableSanitizerInconclusiveExeAllPct}{0.0}
\rqSet{rqFFVecNotAdjudicableSanitizerInconclusiveExeBOneNum}{0}
\rqSet{rqFFVecNotAdjudicableSanitizerInconclusiveExeBOneDen}{146}
\rqSet{rqFFVecNotAdjudicableSanitizerInconclusiveExeBOnePct}{0.0}
\rqSet{rqFFVecNotAdjudicableSanitizerInconclusiveExeBTwoNum}{0}
\rqSet{rqFFVecNotAdjudicableSanitizerInconclusiveExeBTwoDen}{529}
\rqSet{rqFFVecNotAdjudicableSanitizerInconclusiveExeBTwoPct}{0.0}
\rqSet{rqFFVecNotAdjudicableSanitizerInconclusiveExePZeroNum}{0}
\rqSet{rqFFVecNotAdjudicableSanitizerInconclusiveExePZeroDen}{99}
\rqSet{rqFFVecNotAdjudicableSanitizerInconclusiveExePZeroPct}{0.0}
\rqSet{rqFFVecNotAdjudicableSanitizerInconclusiveExePOneNum}{0}
\rqSet{rqFFVecNotAdjudicableSanitizerInconclusiveExePOneDen}{36}
\rqSet{rqFFVecNotAdjudicableSanitizerInconclusiveExePOnePct}{0.0}
\rqSet{rqFFVecNotAdjudicableCUnstableExeAllNum}{0}
\rqSet{rqFFVecNotAdjudicableCUnstableExeAllDen}{810}
\rqSet{rqFFVecNotAdjudicableCUnstableExeAllPct}{0.0}
\rqSet{rqFFVecNotAdjudicableCUnstableExeBOneNum}{0}
\rqSet{rqFFVecNotAdjudicableCUnstableExeBOneDen}{146}
\rqSet{rqFFVecNotAdjudicableCUnstableExeBOnePct}{0.0}
\rqSet{rqFFVecNotAdjudicableCUnstableExeBTwoNum}{0}
\rqSet{rqFFVecNotAdjudicableCUnstableExeBTwoDen}{529}
\rqSet{rqFFVecNotAdjudicableCUnstableExeBTwoPct}{0.0}
\rqSet{rqFFVecNotAdjudicableCUnstableExePZeroNum}{0}
\rqSet{rqFFVecNotAdjudicableCUnstableExePZeroDen}{99}
\rqSet{rqFFVecNotAdjudicableCUnstableExePZeroPct}{0.0}
\rqSet{rqFFVecNotAdjudicableCUnstableExePOneNum}{0}
\rqSet{rqFFVecNotAdjudicableCUnstableExePOneDen}{36}
\rqSet{rqFFVecNotAdjudicableCUnstableExePOnePct}{0.0}
\rqSet{rqFFVecNotAdjudicableUnstableComparisonExeAllNum}{0}
\rqSet{rqFFVecNotAdjudicableUnstableComparisonExeAllDen}{810}
\rqSet{rqFFVecNotAdjudicableUnstableComparisonExeAllPct}{0.0}
\rqSet{rqFFVecNotAdjudicableUnstableComparisonExeBOneNum}{0}
\rqSet{rqFFVecNotAdjudicableUnstableComparisonExeBOneDen}{146}
\rqSet{rqFFVecNotAdjudicableUnstableComparisonExeBOnePct}{0.0}
\rqSet{rqFFVecNotAdjudicableUnstableComparisonExeBTwoNum}{0}
\rqSet{rqFFVecNotAdjudicableUnstableComparisonExeBTwoDen}{529}
\rqSet{rqFFVecNotAdjudicableUnstableComparisonExeBTwoPct}{0.0}
\rqSet{rqFFVecNotAdjudicableUnstableComparisonExePZeroNum}{0}
\rqSet{rqFFVecNotAdjudicableUnstableComparisonExePZeroDen}{99}
\rqSet{rqFFVecNotAdjudicableUnstableComparisonExePZeroPct}{0.0}
\rqSet{rqFFVecNotAdjudicableUnstableComparisonExePOneNum}{0}
\rqSet{rqFFVecNotAdjudicableUnstableComparisonExePOneDen}{36}
\rqSet{rqFFVecNotAdjudicableUnstableComparisonExePOnePct}{0.0}
\rqSet{rqFFVecIncompleteVectorUnreadableExeAllNum}{0}
\rqSet{rqFFVecIncompleteVectorUnreadableExeAllDen}{810}
\rqSet{rqFFVecIncompleteVectorUnreadableExeAllPct}{0.0}
\rqSet{rqFFVecIncompleteVectorUnreadableExeBOneNum}{0}
\rqSet{rqFFVecIncompleteVectorUnreadableExeBOneDen}{146}
\rqSet{rqFFVecIncompleteVectorUnreadableExeBOnePct}{0.0}
\rqSet{rqFFVecIncompleteVectorUnreadableExeBTwoNum}{0}
\rqSet{rqFFVecIncompleteVectorUnreadableExeBTwoDen}{529}
\rqSet{rqFFVecIncompleteVectorUnreadableExeBTwoPct}{0.0}
\rqSet{rqFFVecIncompleteVectorUnreadableExePZeroNum}{0}
\rqSet{rqFFVecIncompleteVectorUnreadableExePZeroDen}{99}
\rqSet{rqFFVecIncompleteVectorUnreadableExePZeroPct}{0.0}
\rqSet{rqFFVecIncompleteVectorUnreadableExePOneNum}{0}
\rqSet{rqFFVecIncompleteVectorUnreadableExePOneDen}{36}
\rqSet{rqFFVecIncompleteVectorUnreadableExePOnePct}{0.0}
\rqSet{rqFFVecIncompleteInvalidInputExeAllNum}{0}
\rqSet{rqFFVecIncompleteInvalidInputExeAllDen}{810}
\rqSet{rqFFVecIncompleteInvalidInputExeAllPct}{0.0}
\rqSet{rqFFVecIncompleteInvalidInputExeBOneNum}{0}
\rqSet{rqFFVecIncompleteInvalidInputExeBOneDen}{146}
\rqSet{rqFFVecIncompleteInvalidInputExeBOnePct}{0.0}
\rqSet{rqFFVecIncompleteInvalidInputExeBTwoNum}{0}
\rqSet{rqFFVecIncompleteInvalidInputExeBTwoDen}{529}
\rqSet{rqFFVecIncompleteInvalidInputExeBTwoPct}{0.0}
\rqSet{rqFFVecIncompleteInvalidInputExePZeroNum}{0}
\rqSet{rqFFVecIncompleteInvalidInputExePZeroDen}{99}
\rqSet{rqFFVecIncompleteInvalidInputExePZeroPct}{0.0}
\rqSet{rqFFVecIncompleteInvalidInputExePOneNum}{0}
\rqSet{rqFFVecIncompleteInvalidInputExePOneDen}{36}
\rqSet{rqFFVecIncompleteInvalidInputExePOnePct}{0.0}
\rqSet{rqFFVecIncompleteRunUnobservedExeAllNum}{0}
\rqSet{rqFFVecIncompleteRunUnobservedExeAllDen}{810}
\rqSet{rqFFVecIncompleteRunUnobservedExeAllPct}{0.0}
\rqSet{rqFFVecIncompleteRunUnobservedExeBOneNum}{0}
\rqSet{rqFFVecIncompleteRunUnobservedExeBOneDen}{146}
\rqSet{rqFFVecIncompleteRunUnobservedExeBOnePct}{0.0}
\rqSet{rqFFVecIncompleteRunUnobservedExeBTwoNum}{0}
\rqSet{rqFFVecIncompleteRunUnobservedExeBTwoDen}{529}
\rqSet{rqFFVecIncompleteRunUnobservedExeBTwoPct}{0.0}
\rqSet{rqFFVecIncompleteRunUnobservedExePZeroNum}{0}
\rqSet{rqFFVecIncompleteRunUnobservedExePZeroDen}{99}
\rqSet{rqFFVecIncompleteRunUnobservedExePZeroPct}{0.0}
\rqSet{rqFFVecIncompleteRunUnobservedExePOneNum}{0}
\rqSet{rqFFVecIncompleteRunUnobservedExePOneDen}{36}
\rqSet{rqFFVecIncompleteRunUnobservedExePOnePct}{0.0}
\rqSet{rqFFVecIncompleteSanBuildFailedExeAllNum}{36}
\rqSet{rqFFVecIncompleteSanBuildFailedExeAllDen}{810}
\rqSet{rqFFVecIncompleteSanBuildFailedExeAllPct}{4.4}
\rqSet{rqFFVecIncompleteSanBuildFailedExeBOneNum}{0}
\rqSet{rqFFVecIncompleteSanBuildFailedExeBOneDen}{146}
\rqSet{rqFFVecIncompleteSanBuildFailedExeBOnePct}{0.0}
\rqSet{rqFFVecIncompleteSanBuildFailedExeBTwoNum}{0}
\rqSet{rqFFVecIncompleteSanBuildFailedExeBTwoDen}{529}
\rqSet{rqFFVecIncompleteSanBuildFailedExeBTwoPct}{0.0}
\rqSet{rqFFVecIncompleteSanBuildFailedExePZeroNum}{0}
\rqSet{rqFFVecIncompleteSanBuildFailedExePZeroDen}{99}
\rqSet{rqFFVecIncompleteSanBuildFailedExePZeroPct}{0.0}
\rqSet{rqFFVecIncompleteSanBuildFailedExePOneNum}{36}
\rqSet{rqFFVecIncompleteSanBuildFailedExePOneDen}{36}
\rqSet{rqFFVecIncompleteSanBuildFailedExePOnePct}{100.0}
\rqSet{rqFFVecIncompleteSanInfraExeAllNum}{0}
\rqSet{rqFFVecIncompleteSanInfraExeAllDen}{810}
\rqSet{rqFFVecIncompleteSanInfraExeAllPct}{0.0}
\rqSet{rqFFVecIncompleteSanInfraExeBOneNum}{0}
\rqSet{rqFFVecIncompleteSanInfraExeBOneDen}{146}
\rqSet{rqFFVecIncompleteSanInfraExeBOnePct}{0.0}
\rqSet{rqFFVecIncompleteSanInfraExeBTwoNum}{0}
\rqSet{rqFFVecIncompleteSanInfraExeBTwoDen}{529}
\rqSet{rqFFVecIncompleteSanInfraExeBTwoPct}{0.0}
\rqSet{rqFFVecIncompleteSanInfraExePZeroNum}{0}
\rqSet{rqFFVecIncompleteSanInfraExePZeroDen}{99}
\rqSet{rqFFVecIncompleteSanInfraExePZeroPct}{0.0}
\rqSet{rqFFVecIncompleteSanInfraExePOneNum}{0}
\rqSet{rqFFVecIncompleteSanInfraExePOneDen}{36}
\rqSet{rqFFVecIncompleteSanInfraExePOnePct}{0.0}
\rqSet{rqRobFFNoStderrExeAllNum}{3}
\rqSet{rqRobFFNoStderrExeAllDen}{92}
\rqSet{rqRobFFNoStderrExeAllPct}{3.3}
\rqSet{rqRobFFNoStderrExeAllKiroNum}{0}
\rqSet{rqRobFFNoStderrExeAllKiroDen}{10}
\rqSet{rqRobFFNoStderrExeAllKiroPct}{0.0}
\rqSet{rqRobFFNoStderrExeAllClaudeNum}{0}
\rqSet{rqRobFFNoStderrExeAllClaudeDen}{5}
\rqSet{rqRobFFNoStderrExeAllClaudePct}{0.0}
\rqSet{rqRobFFNoStderrExeAllCodexNum}{0}
\rqSet{rqRobFFNoStderrExeAllCodexDen}{7}
\rqSet{rqRobFFNoStderrExeAllCodexPct}{0.0}
\rqSet{rqRobFFNoStderrExeAllCTwoRustNum}{0}
\rqSet{rqRobFFNoStderrExeAllCTwoRustDen}{0}
\rqSet{rqRobFFNoStderrExeAllCTwoRustPct}{n/a}
\rqSet{rqRobFFNoStderrExeAllLaertesNum}{0}
\rqSet{rqRobFFNoStderrExeAllLaertesDen}{0}
\rqSet{rqRobFFNoStderrExeAllLaertesPct}{n/a}
\rqSet{rqRobFFNoStderrExeAllCTwoSaferRustNum}{1}
\rqSet{rqRobFFNoStderrExeAllCTwoSaferRustDen}{12}
\rqSet{rqRobFFNoStderrExeAllCTwoSaferRustPct}{8.3}
\rqSet{rqRobFFNoStderrExeAllSmartCTwoRustNum}{0}
\rqSet{rqRobFFNoStderrExeAllSmartCTwoRustDen}{8}
\rqSet{rqRobFFNoStderrExeAllSmartCTwoRustPct}{0.0}
\rqSet{rqRobFFNoStderrExeAllKimiNum}{1}
\rqSet{rqRobFFNoStderrExeAllKimiDen}{17}
\rqSet{rqRobFFNoStderrExeAllKimiPct}{5.9}
\rqSet{rqRobFFNoStderrExeAllGptFiveNum}{1}
\rqSet{rqRobFFNoStderrExeAllGptFiveDen}{18}
\rqSet{rqRobFFNoStderrExeAllGptFivePct}{5.6}
\rqSet{rqRobFFNoStderrExeAllGeminiNum}{0}
\rqSet{rqRobFFNoStderrExeAllGeminiDen}{15}
\rqSet{rqRobFFNoStderrExeAllGeminiPct}{0.0}
\rqSet{rqRobFFNoStderrExeBOneNum}{3}
\rqSet{rqRobFFNoStderrExeBOneDen}{41}
\rqSet{rqRobFFNoStderrExeBOnePct}{7.3}
\rqSet{rqRobFFNoStderrExeBTwoNum}{0}
\rqSet{rqRobFFNoStderrExeBTwoDen}{47}
\rqSet{rqRobFFNoStderrExeBTwoPct}{0.0}
\rqSet{rqRobFFNoStderrExePZeroNum}{0}
\rqSet{rqRobFFNoStderrExePZeroDen}{4}
\rqSet{rqRobFFNoStderrExePZeroPct}{0.0}
\rqSet{rqRobFFNoStderrExePOneNum}{0}
\rqSet{rqRobFFNoStderrExePOneDen}{0}
\rqSet{rqRobFFNoStderrExePOnePct}{n/a}
\rqSet{rqRobFFNoStderrNotAdjExeAllNum}{0}
\rqSet{rqRobFFNoStderrNotAdjExeAllDen}{92}
\rqSet{rqRobFFNoStderrNotAdjExeAllPct}{0.0}
\rqSet{rqRobFFNoStderrNotAdjExeBOneNum}{0}
\rqSet{rqRobFFNoStderrNotAdjExeBOneDen}{41}
\rqSet{rqRobFFNoStderrNotAdjExeBOnePct}{0.0}
\rqSet{rqRobFFNoStderrNotAdjExeBTwoNum}{0}
\rqSet{rqRobFFNoStderrNotAdjExeBTwoDen}{47}
\rqSet{rqRobFFNoStderrNotAdjExeBTwoPct}{0.0}
\rqSet{rqRobFFNoStderrNotAdjExePZeroNum}{0}
\rqSet{rqRobFFNoStderrNotAdjExePZeroDen}{4}
\rqSet{rqRobFFNoStderrNotAdjExePZeroPct}{0.0}
\rqSet{rqRobFFNoStderrNotAdjExePOneNum}{0}
\rqSet{rqRobFFNoStderrNotAdjExePOneDen}{0}
\rqSet{rqRobFFNoStderrNotAdjExePOnePct}{n/a}
\rqSet{rqRobFFNoStderrNotReexExeAllNum}{0}
\rqSet{rqRobFFNoStderrNotReexExeAllDen}{92}
\rqSet{rqRobFFNoStderrNotReexExeAllPct}{0.0}
\rqSet{rqRobFFNoStderrNotReexExeBOneNum}{0}
\rqSet{rqRobFFNoStderrNotReexExeBOneDen}{41}
\rqSet{rqRobFFNoStderrNotReexExeBOnePct}{0.0}
\rqSet{rqRobFFNoStderrNotReexExeBTwoNum}{0}
\rqSet{rqRobFFNoStderrNotReexExeBTwoDen}{47}
\rqSet{rqRobFFNoStderrNotReexExeBTwoPct}{0.0}
\rqSet{rqRobFFNoStderrNotReexExePZeroNum}{0}
\rqSet{rqRobFFNoStderrNotReexExePZeroDen}{4}
\rqSet{rqRobFFNoStderrNotReexExePZeroPct}{0.0}
\rqSet{rqRobFFNoStderrNotReexExePOneNum}{0}
\rqSet{rqRobFFNoStderrNotReexExePOneDen}{0}
\rqSet{rqRobFFNoStderrNotReexExePOnePct}{n/a}
\rqSet{rqRobFFNoStderrChangedExeAllNum}{3}
\rqSet{rqRobFFNoStderrChangedExeAllDen}{92}
\rqSet{rqRobFFNoStderrChangedExeAllPct}{3.3}
\rqSet{rqRobFFNoStderrChangedExeBOneNum}{3}
\rqSet{rqRobFFNoStderrChangedExeBOneDen}{41}
\rqSet{rqRobFFNoStderrChangedExeBOnePct}{7.3}
\rqSet{rqRobFFNoStderrChangedExeBTwoNum}{0}
\rqSet{rqRobFFNoStderrChangedExeBTwoDen}{47}
\rqSet{rqRobFFNoStderrChangedExeBTwoPct}{0.0}
\rqSet{rqRobFFNoStderrChangedExePZeroNum}{0}
\rqSet{rqRobFFNoStderrChangedExePZeroDen}{4}
\rqSet{rqRobFFNoStderrChangedExePZeroPct}{0.0}
\rqSet{rqRobFFNoStderrChangedExePOneNum}{0}
\rqSet{rqRobFFNoStderrChangedExePOneDen}{0}
\rqSet{rqRobFFNoStderrChangedExePOnePct}{n/a}
\rqSet{rqSRExeAllNum}{559}
\rqSet{rqSRExeAllDen}{1060}
\rqSet{rqSRExeAllPct}{52.7}
\rqSet{rqSRExeAllKiroNum}{96}
\rqSet{rqSRExeAllKiroDen}{106}
\rqSet{rqSRExeAllKiroPct}{90.6}
\rqSet{rqSRExeAllClaudeNum}{101}
\rqSet{rqSRExeAllClaudeDen}{106}
\rqSet{rqSRExeAllClaudePct}{95.3}
\rqSet{rqSRExeAllCodexNum}{63}
\rqSet{rqSRExeAllCodexDen}{106}
\rqSet{rqSRExeAllCodexPct}{59.4}
\rqSet{rqSRExeAllCTwoRustNum}{54}
\rqSet{rqSRExeAllCTwoRustDen}{106}
\rqSet{rqSRExeAllCTwoRustPct}{50.9}
\rqSet{rqSRExeAllLaertesNum}{56}
\rqSet{rqSRExeAllLaertesDen}{106}
\rqSet{rqSRExeAllLaertesPct}{52.8}
\rqSet{rqSRExeAllCTwoSaferRustNum}{42}
\rqSet{rqSRExeAllCTwoSaferRustDen}{106}
\rqSet{rqSRExeAllCTwoSaferRustPct}{39.6}
\rqSet{rqSRExeAllSmartCTwoRustNum}{40}
\rqSet{rqSRExeAllSmartCTwoRustDen}{106}
\rqSet{rqSRExeAllSmartCTwoRustPct}{37.7}
\rqSet{rqSRExeAllKimiNum}{28}
\rqSet{rqSRExeAllKimiDen}{106}
\rqSet{rqSRExeAllKimiPct}{26.4}
\rqSet{rqSRExeAllGptFiveNum}{37}
\rqSet{rqSRExeAllGptFiveDen}{106}
\rqSet{rqSRExeAllGptFivePct}{34.9}
\rqSet{rqSRExeAllGeminiNum}{42}
\rqSet{rqSRExeAllGeminiDen}{106}
\rqSet{rqSRExeAllGeminiPct}{39.6}
\rqSet{rqSRExeBOneNum}{379}
\rqSet{rqSRExeBOneDen}{430}
\rqSet{rqSRExeBOnePct}{88.1}
\rqSet{rqSRExeBOneKiroNum}{43}
\rqSet{rqSRExeBOneKiroDen}{43}
\rqSet{rqSRExeBOneKiroPct}{100.0}
\rqSet{rqSRExeBOneClaudeNum}{43}
\rqSet{rqSRExeBOneClaudeDen}{43}
\rqSet{rqSRExeBOneClaudePct}{100.0}
\rqSet{rqSRExeBOneCodexNum}{42}
\rqSet{rqSRExeBOneCodexDen}{43}
\rqSet{rqSRExeBOneCodexPct}{97.7}
\rqSet{rqSRExeBOneCTwoRustNum}{43}
\rqSet{rqSRExeBOneCTwoRustDen}{43}
\rqSet{rqSRExeBOneCTwoRustPct}{100.0}
\rqSet{rqSRExeBOneLaertesNum}{43}
\rqSet{rqSRExeBOneLaertesDen}{43}
\rqSet{rqSRExeBOneLaertesPct}{100.0}
\rqSet{rqSRExeBOneCTwoSaferRustNum}{32}
\rqSet{rqSRExeBOneCTwoSaferRustDen}{43}
\rqSet{rqSRExeBOneCTwoSaferRustPct}{74.4}
\rqSet{rqSRExeBOneSmartCTwoRustNum}{38}
\rqSet{rqSRExeBOneSmartCTwoRustDen}{43}
\rqSet{rqSRExeBOneSmartCTwoRustPct}{88.4}
\rqSet{rqSRExeBOneKimiNum}{26}
\rqSet{rqSRExeBOneKimiDen}{43}
\rqSet{rqSRExeBOneKimiPct}{60.5}
\rqSet{rqSRExeBOneGptFiveNum}{33}
\rqSet{rqSRExeBOneGptFiveDen}{43}
\rqSet{rqSRExeBOneGptFivePct}{76.7}
\rqSet{rqSRExeBOneGeminiNum}{36}
\rqSet{rqSRExeBOneGeminiDen}{43}
\rqSet{rqSRExeBOneGeminiPct}{83.7}
\rqSet{rqSRExeBTwoNum}{66}
\rqSet{rqSRExeBTwoDen}{140}
\rqSet{rqSRExeBTwoPct}{47.1}
\rqSet{rqSRExeBTwoKiroNum}{4}
\rqSet{rqSRExeBTwoKiroDen}{14}
\rqSet{rqSRExeBTwoKiroPct}{28.6}
\rqSet{rqSRExeBTwoClaudeNum}{9}
\rqSet{rqSRExeBTwoClaudeDen}{14}
\rqSet{rqSRExeBTwoClaudePct}{64.3}
\rqSet{rqSRExeBTwoCodexNum}{8}
\rqSet{rqSRExeBTwoCodexDen}{14}
\rqSet{rqSRExeBTwoCodexPct}{57.1}
\rqSet{rqSRExeBTwoCTwoRustNum}{10}
\rqSet{rqSRExeBTwoCTwoRustDen}{14}
\rqSet{rqSRExeBTwoCTwoRustPct}{71.4}
\rqSet{rqSRExeBTwoLaertesNum}{12}
\rqSet{rqSRExeBTwoLaertesDen}{14}
\rqSet{rqSRExeBTwoLaertesPct}{85.7}
\rqSet{rqSRExeBTwoCTwoSaferRustNum}{9}
\rqSet{rqSRExeBTwoCTwoSaferRustDen}{14}
\rqSet{rqSRExeBTwoCTwoSaferRustPct}{64.3}
\rqSet{rqSRExeBTwoSmartCTwoRustNum}{2}
\rqSet{rqSRExeBTwoSmartCTwoRustDen}{14}
\rqSet{rqSRExeBTwoSmartCTwoRustPct}{14.3}
\rqSet{rqSRExeBTwoKimiNum}{2}
\rqSet{rqSRExeBTwoKimiDen}{14}
\rqSet{rqSRExeBTwoKimiPct}{14.3}
\rqSet{rqSRExeBTwoGptFiveNum}{4}
\rqSet{rqSRExeBTwoGptFiveDen}{14}
\rqSet{rqSRExeBTwoGptFivePct}{28.6}
\rqSet{rqSRExeBTwoGeminiNum}{6}
\rqSet{rqSRExeBTwoGeminiDen}{14}
\rqSet{rqSRExeBTwoGeminiPct}{42.9}
\rqSet{rqSRExePZeroNum}{6}
\rqSet{rqSRExePZeroDen}{10}
\rqSet{rqSRExePZeroPct}{60.0}
\rqSet{rqSRExePZeroKiroNum}{1}
\rqSet{rqSRExePZeroKiroDen}{1}
\rqSet{rqSRExePZeroKiroPct}{100.0}
\rqSet{rqSRExePZeroClaudeNum}{1}
\rqSet{rqSRExePZeroClaudeDen}{1}
\rqSet{rqSRExePZeroClaudePct}{100.0}
\rqSet{rqSRExePZeroCodexNum}{1}
\rqSet{rqSRExePZeroCodexDen}{1}
\rqSet{rqSRExePZeroCodexPct}{100.0}
\rqSet{rqSRExePZeroCTwoRustNum}{1}
\rqSet{rqSRExePZeroCTwoRustDen}{1}
\rqSet{rqSRExePZeroCTwoRustPct}{100.0}
\rqSet{rqSRExePZeroLaertesNum}{1}
\rqSet{rqSRExePZeroLaertesDen}{1}
\rqSet{rqSRExePZeroLaertesPct}{100.0}
\rqSet{rqSRExePZeroCTwoSaferRustNum}{1}
\rqSet{rqSRExePZeroCTwoSaferRustDen}{1}
\rqSet{rqSRExePZeroCTwoSaferRustPct}{100.0}
\rqSet{rqSRExePZeroSmartCTwoRustNum}{0}
\rqSet{rqSRExePZeroSmartCTwoRustDen}{1}
\rqSet{rqSRExePZeroSmartCTwoRustPct}{0.0}
\rqSet{rqSRExePZeroKimiNum}{0}
\rqSet{rqSRExePZeroKimiDen}{1}
\rqSet{rqSRExePZeroKimiPct}{0.0}
\rqSet{rqSRExePZeroGptFiveNum}{0}
\rqSet{rqSRExePZeroGptFiveDen}{1}
\rqSet{rqSRExePZeroGptFivePct}{0.0}
\rqSet{rqSRExePZeroGeminiNum}{0}
\rqSet{rqSRExePZeroGeminiDen}{1}
\rqSet{rqSRExePZeroGeminiPct}{0.0}
\rqSet{rqSRExePOneNum}{108}
\rqSet{rqSRExePOneDen}{480}
\rqSet{rqSRExePOnePct}{22.5}
\rqSet{rqSRExePOneKiroNum}{48}
\rqSet{rqSRExePOneKiroDen}{48}
\rqSet{rqSRExePOneKiroPct}{100.0}
\rqSet{rqSRExePOneClaudeNum}{48}
\rqSet{rqSRExePOneClaudeDen}{48}
\rqSet{rqSRExePOneClaudePct}{100.0}
\rqSet{rqSRExePOneCodexNum}{12}
\rqSet{rqSRExePOneCodexDen}{48}
\rqSet{rqSRExePOneCodexPct}{25.0}
\rqSet{rqSRExePOneCTwoRustNum}{0}
\rqSet{rqSRExePOneCTwoRustDen}{48}
\rqSet{rqSRExePOneCTwoRustPct}{0.0}
\rqSet{rqSRExePOneLaertesNum}{0}
\rqSet{rqSRExePOneLaertesDen}{48}
\rqSet{rqSRExePOneLaertesPct}{0.0}
\rqSet{rqSRExePOneCTwoSaferRustNum}{0}
\rqSet{rqSRExePOneCTwoSaferRustDen}{48}
\rqSet{rqSRExePOneCTwoSaferRustPct}{0.0}
\rqSet{rqSRExePOneSmartCTwoRustNum}{0}
\rqSet{rqSRExePOneSmartCTwoRustDen}{48}
\rqSet{rqSRExePOneSmartCTwoRustPct}{0.0}
\rqSet{rqSRExePOneKimiNum}{0}
\rqSet{rqSRExePOneKimiDen}{48}
\rqSet{rqSRExePOneKimiPct}{0.0}
\rqSet{rqSRExePOneGptFiveNum}{0}
\rqSet{rqSRExePOneGptFiveDen}{48}
\rqSet{rqSRExePOneGptFivePct}{0.0}
\rqSet{rqSRExePOneGeminiNum}{0}
\rqSet{rqSRExePOneGeminiDen}{48}
\rqSet{rqSRExePOneGeminiPct}{0.0}
\rqSet{rqSRLoExeAllNum}{300}
\rqSet{rqSRLoExeAllDen}{1060}
\rqSet{rqSRLoExeAllPct}{28.3}
\rqSet{rqSRLoExeAllKiroNum}{59}
\rqSet{rqSRLoExeAllKiroDen}{106}
\rqSet{rqSRLoExeAllKiroPct}{55.7}
\rqSet{rqSRLoExeAllClaudeNum}{69}
\rqSet{rqSRLoExeAllClaudeDen}{106}
\rqSet{rqSRLoExeAllClaudePct}{65.1}
\rqSet{rqSRLoExeAllCodexNum}{45}
\rqSet{rqSRLoExeAllCodexDen}{106}
\rqSet{rqSRLoExeAllCodexPct}{42.5}
\rqSet{rqSRLoExeAllCTwoRustNum}{41}
\rqSet{rqSRLoExeAllCTwoRustDen}{106}
\rqSet{rqSRLoExeAllCTwoRustPct}{38.7}
\rqSet{rqSRLoExeAllLaertesNum}{41}
\rqSet{rqSRLoExeAllLaertesDen}{106}
\rqSet{rqSRLoExeAllLaertesPct}{38.7}
\rqSet{rqSRLoExeAllCTwoSaferRustNum}{20}
\rqSet{rqSRLoExeAllCTwoSaferRustDen}{106}
\rqSet{rqSRLoExeAllCTwoSaferRustPct}{18.9}
\rqSet{rqSRLoExeAllSmartCTwoRustNum}{6}
\rqSet{rqSRLoExeAllSmartCTwoRustDen}{106}
\rqSet{rqSRLoExeAllSmartCTwoRustPct}{5.7}
\rqSet{rqSRLoExeAllKimiNum}{6}
\rqSet{rqSRLoExeAllKimiDen}{106}
\rqSet{rqSRLoExeAllKimiPct}{5.7}
\rqSet{rqSRLoExeAllGptFiveNum}{5}
\rqSet{rqSRLoExeAllGptFiveDen}{106}
\rqSet{rqSRLoExeAllGptFivePct}{4.7}
\rqSet{rqSRLoExeAllGeminiNum}{8}
\rqSet{rqSRLoExeAllGeminiDen}{106}
\rqSet{rqSRLoExeAllGeminiPct}{7.5}
\rqSet{rqSRLoExeBOneNum}{163}
\rqSet{rqSRLoExeBOneDen}{430}
\rqSet{rqSRLoExeBOnePct}{37.9}
\rqSet{rqSRLoExeBOneKiroNum}{10}
\rqSet{rqSRLoExeBOneKiroDen}{43}
\rqSet{rqSRLoExeBOneKiroPct}{23.3}
\rqSet{rqSRLoExeBOneClaudeNum}{18}
\rqSet{rqSRLoExeBOneClaudeDen}{43}
\rqSet{rqSRLoExeBOneClaudePct}{41.9}
\rqSet{rqSRLoExeBOneCodexNum}{28}
\rqSet{rqSRLoExeBOneCodexDen}{43}
\rqSet{rqSRLoExeBOneCodexPct}{65.1}
\rqSet{rqSRLoExeBOneCTwoRustNum}{35}
\rqSet{rqSRLoExeBOneCTwoRustDen}{43}
\rqSet{rqSRLoExeBOneCTwoRustPct}{81.4}
\rqSet{rqSRLoExeBOneLaertesNum}{35}
\rqSet{rqSRLoExeBOneLaertesDen}{43}
\rqSet{rqSRLoExeBOneLaertesPct}{81.4}
\rqSet{rqSRLoExeBOneCTwoSaferRustNum}{17}
\rqSet{rqSRLoExeBOneCTwoSaferRustDen}{43}
\rqSet{rqSRLoExeBOneCTwoSaferRustPct}{39.5}
\rqSet{rqSRLoExeBOneSmartCTwoRustNum}{5}
\rqSet{rqSRLoExeBOneSmartCTwoRustDen}{43}
\rqSet{rqSRLoExeBOneSmartCTwoRustPct}{11.6}
\rqSet{rqSRLoExeBOneKimiNum}{5}
\rqSet{rqSRLoExeBOneKimiDen}{43}
\rqSet{rqSRLoExeBOneKimiPct}{11.6}
\rqSet{rqSRLoExeBOneGptFiveNum}{4}
\rqSet{rqSRLoExeBOneGptFiveDen}{43}
\rqSet{rqSRLoExeBOneGptFivePct}{9.3}
\rqSet{rqSRLoExeBOneGeminiNum}{6}
\rqSet{rqSRLoExeBOneGeminiDen}{43}
\rqSet{rqSRLoExeBOneGeminiPct}{14.0}
\rqSet{rqSRLoExeBTwoNum}{29}
\rqSet{rqSRLoExeBTwoDen}{140}
\rqSet{rqSRLoExeBTwoPct}{20.7}
\rqSet{rqSRLoExeBTwoKiroNum}{1}
\rqSet{rqSRLoExeBTwoKiroDen}{14}
\rqSet{rqSRLoExeBTwoKiroPct}{7.1}
\rqSet{rqSRLoExeBTwoClaudeNum}{3}
\rqSet{rqSRLoExeBTwoClaudeDen}{14}
\rqSet{rqSRLoExeBTwoClaudePct}{21.4}
\rqSet{rqSRLoExeBTwoCodexNum}{5}
\rqSet{rqSRLoExeBTwoCodexDen}{14}
\rqSet{rqSRLoExeBTwoCodexPct}{35.7}
\rqSet{rqSRLoExeBTwoCTwoRustNum}{6}
\rqSet{rqSRLoExeBTwoCTwoRustDen}{14}
\rqSet{rqSRLoExeBTwoCTwoRustPct}{42.9}
\rqSet{rqSRLoExeBTwoLaertesNum}{6}
\rqSet{rqSRLoExeBTwoLaertesDen}{14}
\rqSet{rqSRLoExeBTwoLaertesPct}{42.9}
\rqSet{rqSRLoExeBTwoCTwoSaferRustNum}{3}
\rqSet{rqSRLoExeBTwoCTwoSaferRustDen}{14}
\rqSet{rqSRLoExeBTwoCTwoSaferRustPct}{21.4}
\rqSet{rqSRLoExeBTwoSmartCTwoRustNum}{1}
\rqSet{rqSRLoExeBTwoSmartCTwoRustDen}{14}
\rqSet{rqSRLoExeBTwoSmartCTwoRustPct}{7.1}
\rqSet{rqSRLoExeBTwoKimiNum}{1}
\rqSet{rqSRLoExeBTwoKimiDen}{14}
\rqSet{rqSRLoExeBTwoKimiPct}{7.1}
\rqSet{rqSRLoExeBTwoGptFiveNum}{1}
\rqSet{rqSRLoExeBTwoGptFiveDen}{14}
\rqSet{rqSRLoExeBTwoGptFivePct}{7.1}
\rqSet{rqSRLoExeBTwoGeminiNum}{2}
\rqSet{rqSRLoExeBTwoGeminiDen}{14}
\rqSet{rqSRLoExeBTwoGeminiPct}{14.3}
\rqSet{rqSRLoExePZeroNum}{0}
\rqSet{rqSRLoExePZeroDen}{10}
\rqSet{rqSRLoExePZeroPct}{0.0}
\rqSet{rqSRLoExePZeroKiroNum}{0}
\rqSet{rqSRLoExePZeroKiroDen}{1}
\rqSet{rqSRLoExePZeroKiroPct}{0.0}
\rqSet{rqSRLoExePZeroClaudeNum}{0}
\rqSet{rqSRLoExePZeroClaudeDen}{1}
\rqSet{rqSRLoExePZeroClaudePct}{0.0}
\rqSet{rqSRLoExePZeroCodexNum}{0}
\rqSet{rqSRLoExePZeroCodexDen}{1}
\rqSet{rqSRLoExePZeroCodexPct}{0.0}
\rqSet{rqSRLoExePZeroCTwoRustNum}{0}
\rqSet{rqSRLoExePZeroCTwoRustDen}{1}
\rqSet{rqSRLoExePZeroCTwoRustPct}{0.0}
\rqSet{rqSRLoExePZeroLaertesNum}{0}
\rqSet{rqSRLoExePZeroLaertesDen}{1}
\rqSet{rqSRLoExePZeroLaertesPct}{0.0}
\rqSet{rqSRLoExePZeroCTwoSaferRustNum}{0}
\rqSet{rqSRLoExePZeroCTwoSaferRustDen}{1}
\rqSet{rqSRLoExePZeroCTwoSaferRustPct}{0.0}
\rqSet{rqSRLoExePZeroSmartCTwoRustNum}{0}
\rqSet{rqSRLoExePZeroSmartCTwoRustDen}{1}
\rqSet{rqSRLoExePZeroSmartCTwoRustPct}{0.0}
\rqSet{rqSRLoExePZeroKimiNum}{0}
\rqSet{rqSRLoExePZeroKimiDen}{1}
\rqSet{rqSRLoExePZeroKimiPct}{0.0}
\rqSet{rqSRLoExePZeroGptFiveNum}{0}
\rqSet{rqSRLoExePZeroGptFiveDen}{1}
\rqSet{rqSRLoExePZeroGptFivePct}{0.0}
\rqSet{rqSRLoExePZeroGeminiNum}{0}
\rqSet{rqSRLoExePZeroGeminiDen}{1}
\rqSet{rqSRLoExePZeroGeminiPct}{0.0}
\rqSet{rqSRLoExePOneNum}{108}
\rqSet{rqSRLoExePOneDen}{480}
\rqSet{rqSRLoExePOnePct}{22.5}
\rqSet{rqSRLoExePOneKiroNum}{48}
\rqSet{rqSRLoExePOneKiroDen}{48}
\rqSet{rqSRLoExePOneKiroPct}{100.0}
\rqSet{rqSRLoExePOneClaudeNum}{48}
\rqSet{rqSRLoExePOneClaudeDen}{48}
\rqSet{rqSRLoExePOneClaudePct}{100.0}
\rqSet{rqSRLoExePOneCodexNum}{12}
\rqSet{rqSRLoExePOneCodexDen}{48}
\rqSet{rqSRLoExePOneCodexPct}{25.0}
\rqSet{rqSRLoExePOneCTwoRustNum}{0}
\rqSet{rqSRLoExePOneCTwoRustDen}{48}
\rqSet{rqSRLoExePOneCTwoRustPct}{0.0}
\rqSet{rqSRLoExePOneLaertesNum}{0}
\rqSet{rqSRLoExePOneLaertesDen}{48}
\rqSet{rqSRLoExePOneLaertesPct}{0.0}
\rqSet{rqSRLoExePOneCTwoSaferRustNum}{0}
\rqSet{rqSRLoExePOneCTwoSaferRustDen}{48}
\rqSet{rqSRLoExePOneCTwoSaferRustPct}{0.0}
\rqSet{rqSRLoExePOneSmartCTwoRustNum}{0}
\rqSet{rqSRLoExePOneSmartCTwoRustDen}{48}
\rqSet{rqSRLoExePOneSmartCTwoRustPct}{0.0}
\rqSet{rqSRLoExePOneKimiNum}{0}
\rqSet{rqSRLoExePOneKimiDen}{48}
\rqSet{rqSRLoExePOneKimiPct}{0.0}
\rqSet{rqSRLoExePOneGptFiveNum}{0}
\rqSet{rqSRLoExePOneGptFiveDen}{48}
\rqSet{rqSRLoExePOneGptFivePct}{0.0}
\rqSet{rqSRLoExePOneGeminiNum}{0}
\rqSet{rqSRLoExePOneGeminiDen}{48}
\rqSet{rqSRLoExePOneGeminiPct}{0.0}
\rqSet{rqSRHiExeAllNum}{419}
\rqSet{rqSRHiExeAllDen}{1060}
\rqSet{rqSRHiExeAllPct}{39.5}
\rqSet{rqSRHiExeAllKiroNum}{67}
\rqSet{rqSRHiExeAllKiroDen}{106}
\rqSet{rqSRHiExeAllKiroPct}{63.2}
\rqSet{rqSRHiExeAllClaudeNum}{77}
\rqSet{rqSRHiExeAllClaudeDen}{106}
\rqSet{rqSRHiExeAllClaudePct}{72.6}
\rqSet{rqSRHiExeAllCodexNum}{90}
\rqSet{rqSRHiExeAllCodexDen}{106}
\rqSet{rqSRHiExeAllCodexPct}{84.9}
\rqSet{rqSRHiExeAllCTwoRustNum}{51}
\rqSet{rqSRHiExeAllCTwoRustDen}{106}
\rqSet{rqSRHiExeAllCTwoRustPct}{48.1}
\rqSet{rqSRHiExeAllLaertesNum}{53}
\rqSet{rqSRHiExeAllLaertesDen}{106}
\rqSet{rqSRHiExeAllLaertesPct}{50.0}
\rqSet{rqSRHiExeAllCTwoSaferRustNum}{30}
\rqSet{rqSRHiExeAllCTwoSaferRustDen}{106}
\rqSet{rqSRHiExeAllCTwoSaferRustPct}{28.3}
\rqSet{rqSRHiExeAllSmartCTwoRustNum}{13}
\rqSet{rqSRHiExeAllSmartCTwoRustDen}{106}
\rqSet{rqSRHiExeAllSmartCTwoRustPct}{12.3}
\rqSet{rqSRHiExeAllKimiNum}{10}
\rqSet{rqSRHiExeAllKimiDen}{106}
\rqSet{rqSRHiExeAllKimiPct}{9.4}
\rqSet{rqSRHiExeAllGptFiveNum}{12}
\rqSet{rqSRHiExeAllGptFiveDen}{106}
\rqSet{rqSRHiExeAllGptFivePct}{11.3}
\rqSet{rqSRHiExeAllGeminiNum}{16}
\rqSet{rqSRHiExeAllGeminiDen}{106}
\rqSet{rqSRHiExeAllGeminiPct}{15.1}
\rqSet{rqSRHiExeBOneNum}{227}
\rqSet{rqSRHiExeBOneDen}{430}
\rqSet{rqSRHiExeBOnePct}{52.8}
\rqSet{rqSRHiExeBOneKiroNum}{17}
\rqSet{rqSRHiExeBOneKiroDen}{43}
\rqSet{rqSRHiExeBOneKiroPct}{39.5}
\rqSet{rqSRHiExeBOneClaudeNum}{25}
\rqSet{rqSRHiExeBOneClaudeDen}{43}
\rqSet{rqSRHiExeBOneClaudePct}{58.1}
\rqSet{rqSRHiExeBOneCodexNum}{36}
\rqSet{rqSRHiExeBOneCodexDen}{43}
\rqSet{rqSRHiExeBOneCodexPct}{83.7}
\rqSet{rqSRHiExeBOneCTwoRustNum}{42}
\rqSet{rqSRHiExeBOneCTwoRustDen}{43}
\rqSet{rqSRHiExeBOneCTwoRustPct}{97.7}
\rqSet{rqSRHiExeBOneLaertesNum}{42}
\rqSet{rqSRHiExeBOneLaertesDen}{43}
\rqSet{rqSRHiExeBOneLaertesPct}{97.7}
\rqSet{rqSRHiExeBOneCTwoSaferRustNum}{22}
\rqSet{rqSRHiExeBOneCTwoSaferRustDen}{43}
\rqSet{rqSRHiExeBOneCTwoSaferRustPct}{51.2}
\rqSet{rqSRHiExeBOneSmartCTwoRustNum}{12}
\rqSet{rqSRHiExeBOneSmartCTwoRustDen}{43}
\rqSet{rqSRHiExeBOneSmartCTwoRustPct}{27.9}
\rqSet{rqSRHiExeBOneKimiNum}{8}
\rqSet{rqSRHiExeBOneKimiDen}{43}
\rqSet{rqSRHiExeBOneKimiPct}{18.6}
\rqSet{rqSRHiExeBOneGptFiveNum}{10}
\rqSet{rqSRHiExeBOneGptFiveDen}{43}
\rqSet{rqSRHiExeBOneGptFivePct}{23.3}
\rqSet{rqSRHiExeBOneGeminiNum}{13}
\rqSet{rqSRHiExeBOneGeminiDen}{43}
\rqSet{rqSRHiExeBOneGeminiPct}{30.2}
\rqSet{rqSRHiExeBTwoNum}{42}
\rqSet{rqSRHiExeBTwoDen}{140}
\rqSet{rqSRHiExeBTwoPct}{30.0}
\rqSet{rqSRHiExeBTwoKiroNum}{1}
\rqSet{rqSRHiExeBTwoKiroDen}{14}
\rqSet{rqSRHiExeBTwoKiroPct}{7.1}
\rqSet{rqSRHiExeBTwoClaudeNum}{3}
\rqSet{rqSRHiExeBTwoClaudeDen}{14}
\rqSet{rqSRHiExeBTwoClaudePct}{21.4}
\rqSet{rqSRHiExeBTwoCodexNum}{5}
\rqSet{rqSRHiExeBTwoCodexDen}{14}
\rqSet{rqSRHiExeBTwoCodexPct}{35.7}
\rqSet{rqSRHiExeBTwoCTwoRustNum}{8}
\rqSet{rqSRHiExeBTwoCTwoRustDen}{14}
\rqSet{rqSRHiExeBTwoCTwoRustPct}{57.1}
\rqSet{rqSRHiExeBTwoLaertesNum}{10}
\rqSet{rqSRHiExeBTwoLaertesDen}{14}
\rqSet{rqSRHiExeBTwoLaertesPct}{71.4}
\rqSet{rqSRHiExeBTwoCTwoSaferRustNum}{7}
\rqSet{rqSRHiExeBTwoCTwoSaferRustDen}{14}
\rqSet{rqSRHiExeBTwoCTwoSaferRustPct}{50.0}
\rqSet{rqSRHiExeBTwoSmartCTwoRustNum}{1}
\rqSet{rqSRHiExeBTwoSmartCTwoRustDen}{14}
\rqSet{rqSRHiExeBTwoSmartCTwoRustPct}{7.1}
\rqSet{rqSRHiExeBTwoKimiNum}{2}
\rqSet{rqSRHiExeBTwoKimiDen}{14}
\rqSet{rqSRHiExeBTwoKimiPct}{14.3}
\rqSet{rqSRHiExeBTwoGptFiveNum}{2}
\rqSet{rqSRHiExeBTwoGptFiveDen}{14}
\rqSet{rqSRHiExeBTwoGptFivePct}{14.3}
\rqSet{rqSRHiExeBTwoGeminiNum}{3}
\rqSet{rqSRHiExeBTwoGeminiDen}{14}
\rqSet{rqSRHiExeBTwoGeminiPct}{21.4}
\rqSet{rqSRHiExePZeroNum}{6}
\rqSet{rqSRHiExePZeroDen}{10}
\rqSet{rqSRHiExePZeroPct}{60.0}
\rqSet{rqSRHiExePZeroKiroNum}{1}
\rqSet{rqSRHiExePZeroKiroDen}{1}
\rqSet{rqSRHiExePZeroKiroPct}{100.0}
\rqSet{rqSRHiExePZeroClaudeNum}{1}
\rqSet{rqSRHiExePZeroClaudeDen}{1}
\rqSet{rqSRHiExePZeroClaudePct}{100.0}
\rqSet{rqSRHiExePZeroCodexNum}{1}
\rqSet{rqSRHiExePZeroCodexDen}{1}
\rqSet{rqSRHiExePZeroCodexPct}{100.0}
\rqSet{rqSRHiExePZeroCTwoRustNum}{1}
\rqSet{rqSRHiExePZeroCTwoRustDen}{1}
\rqSet{rqSRHiExePZeroCTwoRustPct}{100.0}
\rqSet{rqSRHiExePZeroLaertesNum}{1}
\rqSet{rqSRHiExePZeroLaertesDen}{1}
\rqSet{rqSRHiExePZeroLaertesPct}{100.0}
\rqSet{rqSRHiExePZeroCTwoSaferRustNum}{1}
\rqSet{rqSRHiExePZeroCTwoSaferRustDen}{1}
\rqSet{rqSRHiExePZeroCTwoSaferRustPct}{100.0}
\rqSet{rqSRHiExePZeroSmartCTwoRustNum}{0}
\rqSet{rqSRHiExePZeroSmartCTwoRustDen}{1}
\rqSet{rqSRHiExePZeroSmartCTwoRustPct}{0.0}
\rqSet{rqSRHiExePZeroKimiNum}{0}
\rqSet{rqSRHiExePZeroKimiDen}{1}
\rqSet{rqSRHiExePZeroKimiPct}{0.0}
\rqSet{rqSRHiExePZeroGptFiveNum}{0}
\rqSet{rqSRHiExePZeroGptFiveDen}{1}
\rqSet{rqSRHiExePZeroGptFivePct}{0.0}
\rqSet{rqSRHiExePZeroGeminiNum}{0}
\rqSet{rqSRHiExePZeroGeminiDen}{1}
\rqSet{rqSRHiExePZeroGeminiPct}{0.0}
\rqSet{rqSRHiExePOneNum}{144}
\rqSet{rqSRHiExePOneDen}{480}
\rqSet{rqSRHiExePOnePct}{30.0}
\rqSet{rqSRHiExePOneKiroNum}{48}
\rqSet{rqSRHiExePOneKiroDen}{48}
\rqSet{rqSRHiExePOneKiroPct}{100.0}
\rqSet{rqSRHiExePOneClaudeNum}{48}
\rqSet{rqSRHiExePOneClaudeDen}{48}
\rqSet{rqSRHiExePOneClaudePct}{100.0}
\rqSet{rqSRHiExePOneCodexNum}{48}
\rqSet{rqSRHiExePOneCodexDen}{48}
\rqSet{rqSRHiExePOneCodexPct}{100.0}
\rqSet{rqSRHiExePOneCTwoRustNum}{0}
\rqSet{rqSRHiExePOneCTwoRustDen}{48}
\rqSet{rqSRHiExePOneCTwoRustPct}{0.0}
\rqSet{rqSRHiExePOneLaertesNum}{0}
\rqSet{rqSRHiExePOneLaertesDen}{48}
\rqSet{rqSRHiExePOneLaertesPct}{0.0}
\rqSet{rqSRHiExePOneCTwoSaferRustNum}{0}
\rqSet{rqSRHiExePOneCTwoSaferRustDen}{48}
\rqSet{rqSRHiExePOneCTwoSaferRustPct}{0.0}
\rqSet{rqSRHiExePOneSmartCTwoRustNum}{0}
\rqSet{rqSRHiExePOneSmartCTwoRustDen}{48}
\rqSet{rqSRHiExePOneSmartCTwoRustPct}{0.0}
\rqSet{rqSRHiExePOneKimiNum}{0}
\rqSet{rqSRHiExePOneKimiDen}{48}
\rqSet{rqSRHiExePOneKimiPct}{0.0}
\rqSet{rqSRHiExePOneGptFiveNum}{0}
\rqSet{rqSRHiExePOneGptFiveDen}{48}
\rqSet{rqSRHiExePOneGptFivePct}{0.0}
\rqSet{rqSRHiExePOneGeminiNum}{0}
\rqSet{rqSRHiExePOneGeminiDen}{48}
\rqSet{rqSRHiExePOneGeminiPct}{0.0}
\rqSet{rqSRWidthExeAllNum}{119}
\rqSet{rqSRWidthExeAllDen}{1060}
\rqSet{rqSRWidthExeAllPct}{11.2}
\rqSet{rqSRWidthExeAllKiroNum}{8}
\rqSet{rqSRWidthExeAllKiroDen}{106}
\rqSet{rqSRWidthExeAllKiroPct}{7.5}
\rqSet{rqSRWidthExeAllClaudeNum}{8}
\rqSet{rqSRWidthExeAllClaudeDen}{106}
\rqSet{rqSRWidthExeAllClaudePct}{7.5}
\rqSet{rqSRWidthExeAllCodexNum}{45}
\rqSet{rqSRWidthExeAllCodexDen}{106}
\rqSet{rqSRWidthExeAllCodexPct}{42.5}
\rqSet{rqSRWidthExeAllCTwoRustNum}{10}
\rqSet{rqSRWidthExeAllCTwoRustDen}{106}
\rqSet{rqSRWidthExeAllCTwoRustPct}{9.4}
\rqSet{rqSRWidthExeAllLaertesNum}{12}
\rqSet{rqSRWidthExeAllLaertesDen}{106}
\rqSet{rqSRWidthExeAllLaertesPct}{11.3}
\rqSet{rqSRWidthExeAllCTwoSaferRustNum}{10}
\rqSet{rqSRWidthExeAllCTwoSaferRustDen}{106}
\rqSet{rqSRWidthExeAllCTwoSaferRustPct}{9.4}
\rqSet{rqSRWidthExeAllSmartCTwoRustNum}{7}
\rqSet{rqSRWidthExeAllSmartCTwoRustDen}{106}
\rqSet{rqSRWidthExeAllSmartCTwoRustPct}{6.6}
\rqSet{rqSRWidthExeAllKimiNum}{4}
\rqSet{rqSRWidthExeAllKimiDen}{106}
\rqSet{rqSRWidthExeAllKimiPct}{3.8}
\rqSet{rqSRWidthExeAllGptFiveNum}{7}
\rqSet{rqSRWidthExeAllGptFiveDen}{106}
\rqSet{rqSRWidthExeAllGptFivePct}{6.6}
\rqSet{rqSRWidthExeAllGeminiNum}{8}
\rqSet{rqSRWidthExeAllGeminiDen}{106}
\rqSet{rqSRWidthExeAllGeminiPct}{7.5}
\rqSet{rqSRWidthExeBOneNum}{64}
\rqSet{rqSRWidthExeBOneDen}{430}
\rqSet{rqSRWidthExeBOnePct}{14.9}
\rqSet{rqSRWidthExeBOneKiroNum}{7}
\rqSet{rqSRWidthExeBOneKiroDen}{43}
\rqSet{rqSRWidthExeBOneKiroPct}{16.3}
\rqSet{rqSRWidthExeBOneClaudeNum}{7}
\rqSet{rqSRWidthExeBOneClaudeDen}{43}
\rqSet{rqSRWidthExeBOneClaudePct}{16.3}
\rqSet{rqSRWidthExeBOneCodexNum}{8}
\rqSet{rqSRWidthExeBOneCodexDen}{43}
\rqSet{rqSRWidthExeBOneCodexPct}{18.6}
\rqSet{rqSRWidthExeBOneCTwoRustNum}{7}
\rqSet{rqSRWidthExeBOneCTwoRustDen}{43}
\rqSet{rqSRWidthExeBOneCTwoRustPct}{16.3}
\rqSet{rqSRWidthExeBOneLaertesNum}{7}
\rqSet{rqSRWidthExeBOneLaertesDen}{43}
\rqSet{rqSRWidthExeBOneLaertesPct}{16.3}
\rqSet{rqSRWidthExeBOneCTwoSaferRustNum}{5}
\rqSet{rqSRWidthExeBOneCTwoSaferRustDen}{43}
\rqSet{rqSRWidthExeBOneCTwoSaferRustPct}{11.6}
\rqSet{rqSRWidthExeBOneSmartCTwoRustNum}{7}
\rqSet{rqSRWidthExeBOneSmartCTwoRustDen}{43}
\rqSet{rqSRWidthExeBOneSmartCTwoRustPct}{16.3}
\rqSet{rqSRWidthExeBOneKimiNum}{3}
\rqSet{rqSRWidthExeBOneKimiDen}{43}
\rqSet{rqSRWidthExeBOneKimiPct}{7.0}
\rqSet{rqSRWidthExeBOneGptFiveNum}{6}
\rqSet{rqSRWidthExeBOneGptFiveDen}{43}
\rqSet{rqSRWidthExeBOneGptFivePct}{14.0}
\rqSet{rqSRWidthExeBOneGeminiNum}{7}
\rqSet{rqSRWidthExeBOneGeminiDen}{43}
\rqSet{rqSRWidthExeBOneGeminiPct}{16.3}
\rqSet{rqSRWidthExeBTwoNum}{13}
\rqSet{rqSRWidthExeBTwoDen}{140}
\rqSet{rqSRWidthExeBTwoPct}{9.3}
\rqSet{rqSRWidthExeBTwoKiroNum}{0}
\rqSet{rqSRWidthExeBTwoKiroDen}{14}
\rqSet{rqSRWidthExeBTwoKiroPct}{0.0}
\rqSet{rqSRWidthExeBTwoClaudeNum}{0}
\rqSet{rqSRWidthExeBTwoClaudeDen}{14}
\rqSet{rqSRWidthExeBTwoClaudePct}{0.0}
\rqSet{rqSRWidthExeBTwoCodexNum}{0}
\rqSet{rqSRWidthExeBTwoCodexDen}{14}
\rqSet{rqSRWidthExeBTwoCodexPct}{0.0}
\rqSet{rqSRWidthExeBTwoCTwoRustNum}{2}
\rqSet{rqSRWidthExeBTwoCTwoRustDen}{14}
\rqSet{rqSRWidthExeBTwoCTwoRustPct}{14.3}
\rqSet{rqSRWidthExeBTwoLaertesNum}{4}
\rqSet{rqSRWidthExeBTwoLaertesDen}{14}
\rqSet{rqSRWidthExeBTwoLaertesPct}{28.6}
\rqSet{rqSRWidthExeBTwoCTwoSaferRustNum}{4}
\rqSet{rqSRWidthExeBTwoCTwoSaferRustDen}{14}
\rqSet{rqSRWidthExeBTwoCTwoSaferRustPct}{28.6}
\rqSet{rqSRWidthExeBTwoSmartCTwoRustNum}{0}
\rqSet{rqSRWidthExeBTwoSmartCTwoRustDen}{14}
\rqSet{rqSRWidthExeBTwoSmartCTwoRustPct}{0.0}
\rqSet{rqSRWidthExeBTwoKimiNum}{1}
\rqSet{rqSRWidthExeBTwoKimiDen}{14}
\rqSet{rqSRWidthExeBTwoKimiPct}{7.1}
\rqSet{rqSRWidthExeBTwoGptFiveNum}{1}
\rqSet{rqSRWidthExeBTwoGptFiveDen}{14}
\rqSet{rqSRWidthExeBTwoGptFivePct}{7.1}
\rqSet{rqSRWidthExeBTwoGeminiNum}{1}
\rqSet{rqSRWidthExeBTwoGeminiDen}{14}
\rqSet{rqSRWidthExeBTwoGeminiPct}{7.1}
\rqSet{rqSRWidthExePZeroNum}{6}
\rqSet{rqSRWidthExePZeroDen}{10}
\rqSet{rqSRWidthExePZeroPct}{60.0}
\rqSet{rqSRWidthExePZeroKiroNum}{1}
\rqSet{rqSRWidthExePZeroKiroDen}{1}
\rqSet{rqSRWidthExePZeroKiroPct}{100.0}
\rqSet{rqSRWidthExePZeroClaudeNum}{1}
\rqSet{rqSRWidthExePZeroClaudeDen}{1}
\rqSet{rqSRWidthExePZeroClaudePct}{100.0}
\rqSet{rqSRWidthExePZeroCodexNum}{1}
\rqSet{rqSRWidthExePZeroCodexDen}{1}
\rqSet{rqSRWidthExePZeroCodexPct}{100.0}
\rqSet{rqSRWidthExePZeroCTwoRustNum}{1}
\rqSet{rqSRWidthExePZeroCTwoRustDen}{1}
\rqSet{rqSRWidthExePZeroCTwoRustPct}{100.0}
\rqSet{rqSRWidthExePZeroLaertesNum}{1}
\rqSet{rqSRWidthExePZeroLaertesDen}{1}
\rqSet{rqSRWidthExePZeroLaertesPct}{100.0}
\rqSet{rqSRWidthExePZeroCTwoSaferRustNum}{1}
\rqSet{rqSRWidthExePZeroCTwoSaferRustDen}{1}
\rqSet{rqSRWidthExePZeroCTwoSaferRustPct}{100.0}
\rqSet{rqSRWidthExePZeroSmartCTwoRustNum}{0}
\rqSet{rqSRWidthExePZeroSmartCTwoRustDen}{1}
\rqSet{rqSRWidthExePZeroSmartCTwoRustPct}{0.0}
\rqSet{rqSRWidthExePZeroKimiNum}{0}
\rqSet{rqSRWidthExePZeroKimiDen}{1}
\rqSet{rqSRWidthExePZeroKimiPct}{0.0}
\rqSet{rqSRWidthExePZeroGptFiveNum}{0}
\rqSet{rqSRWidthExePZeroGptFiveDen}{1}
\rqSet{rqSRWidthExePZeroGptFivePct}{0.0}
\rqSet{rqSRWidthExePZeroGeminiNum}{0}
\rqSet{rqSRWidthExePZeroGeminiDen}{1}
\rqSet{rqSRWidthExePZeroGeminiPct}{0.0}
\rqSet{rqSRWidthExePOneNum}{36}
\rqSet{rqSRWidthExePOneDen}{480}
\rqSet{rqSRWidthExePOnePct}{7.5}
\rqSet{rqSRWidthExePOneKiroNum}{0}
\rqSet{rqSRWidthExePOneKiroDen}{48}
\rqSet{rqSRWidthExePOneKiroPct}{0.0}
\rqSet{rqSRWidthExePOneClaudeNum}{0}
\rqSet{rqSRWidthExePOneClaudeDen}{48}
\rqSet{rqSRWidthExePOneClaudePct}{0.0}
\rqSet{rqSRWidthExePOneCodexNum}{36}
\rqSet{rqSRWidthExePOneCodexDen}{48}
\rqSet{rqSRWidthExePOneCodexPct}{75.0}
\rqSet{rqSRWidthExePOneCTwoRustNum}{0}
\rqSet{rqSRWidthExePOneCTwoRustDen}{48}
\rqSet{rqSRWidthExePOneCTwoRustPct}{0.0}
\rqSet{rqSRWidthExePOneLaertesNum}{0}
\rqSet{rqSRWidthExePOneLaertesDen}{48}
\rqSet{rqSRWidthExePOneLaertesPct}{0.0}
\rqSet{rqSRWidthExePOneCTwoSaferRustNum}{0}
\rqSet{rqSRWidthExePOneCTwoSaferRustDen}{48}
\rqSet{rqSRWidthExePOneCTwoSaferRustPct}{0.0}
\rqSet{rqSRWidthExePOneSmartCTwoRustNum}{0}
\rqSet{rqSRWidthExePOneSmartCTwoRustDen}{48}
\rqSet{rqSRWidthExePOneSmartCTwoRustPct}{0.0}
\rqSet{rqSRWidthExePOneKimiNum}{0}
\rqSet{rqSRWidthExePOneKimiDen}{48}
\rqSet{rqSRWidthExePOneKimiPct}{0.0}
\rqSet{rqSRWidthExePOneGptFiveNum}{0}
\rqSet{rqSRWidthExePOneGptFiveDen}{48}
\rqSet{rqSRWidthExePOneGptFivePct}{0.0}
\rqSet{rqSRWidthExePOneGeminiNum}{0}
\rqSet{rqSRWidthExePOneGeminiDen}{48}
\rqSet{rqSRWidthExePOneGeminiPct}{0.0}
\rqSet{rqTermCExitExeAllNum}{35689763}
\rqSet{rqTermCExitExeAllDen}{35894359}
\rqSet{rqTermCExitExeAllPct}{99.4}
\rqSet{rqTermCExitExeBOneNum}{29548953}
\rqSet{rqTermCExitExeBOneDen}{29751287}
\rqSet{rqTermCExitExeBOnePct}{99.3}
\rqSet{rqTermCExitExeBTwoNum}{6140702}
\rqSet{rqTermCExitExeBTwoDen}{6142964}
\rqSet{rqTermCExitExeBTwoPct}{100.0}
\rqSet{rqTermCExitExePZeroNum}{0}
\rqSet{rqTermCExitExePZeroDen}{0}
\rqSet{rqTermCExitExePZeroPct}{n/a}
\rqSet{rqTermCExitExePOneNum}{108}
\rqSet{rqTermCExitExePOneDen}{108}
\rqSet{rqTermCExitExePOnePct}{100.0}
\rqSet{rqTermCSignalExeAllNum}{204596}
\rqSet{rqTermCSignalExeAllDen}{35894359}
\rqSet{rqTermCSignalExeAllPct}{0.6}
\rqSet{rqTermCSignalExeBOneNum}{202334}
\rqSet{rqTermCSignalExeBOneDen}{29751287}
\rqSet{rqTermCSignalExeBOnePct}{0.7}
\rqSet{rqTermCSignalExeBTwoNum}{2262}
\rqSet{rqTermCSignalExeBTwoDen}{6142964}
\rqSet{rqTermCSignalExeBTwoPct}{0.0}
\rqSet{rqTermCSignalExePZeroNum}{0}
\rqSet{rqTermCSignalExePZeroDen}{0}
\rqSet{rqTermCSignalExePZeroPct}{n/a}
\rqSet{rqTermCSignalExePOneNum}{0}
\rqSet{rqTermCSignalExePOneDen}{108}
\rqSet{rqTermCSignalExePOnePct}{0.0}
\rqSet{rqTermCTimeoutExeAllNum}{0}
\rqSet{rqTermCTimeoutExeAllDen}{35894359}
\rqSet{rqTermCTimeoutExeAllPct}{0.0}
\rqSet{rqTermCTimeoutExeBOneNum}{0}
\rqSet{rqTermCTimeoutExeBOneDen}{29751287}
\rqSet{rqTermCTimeoutExeBOnePct}{0.0}
\rqSet{rqTermCTimeoutExeBTwoNum}{0}
\rqSet{rqTermCTimeoutExeBTwoDen}{6142964}
\rqSet{rqTermCTimeoutExeBTwoPct}{0.0}
\rqSet{rqTermCTimeoutExePZeroNum}{0}
\rqSet{rqTermCTimeoutExePZeroDen}{0}
\rqSet{rqTermCTimeoutExePZeroPct}{n/a}
\rqSet{rqTermCTimeoutExePOneNum}{0}
\rqSet{rqTermCTimeoutExePOneDen}{108}
\rqSet{rqTermCTimeoutExePOnePct}{0.0}
\rqSet{rqTermCOutputLimitExeAllNum}{0}
\rqSet{rqTermCOutputLimitExeAllDen}{35894359}
\rqSet{rqTermCOutputLimitExeAllPct}{0.0}
\rqSet{rqTermCOutputLimitExeBOneNum}{0}
\rqSet{rqTermCOutputLimitExeBOneDen}{29751287}
\rqSet{rqTermCOutputLimitExeBOnePct}{0.0}
\rqSet{rqTermCOutputLimitExeBTwoNum}{0}
\rqSet{rqTermCOutputLimitExeBTwoDen}{6142964}
\rqSet{rqTermCOutputLimitExeBTwoPct}{0.0}
\rqSet{rqTermCOutputLimitExePZeroNum}{0}
\rqSet{rqTermCOutputLimitExePZeroDen}{0}
\rqSet{rqTermCOutputLimitExePZeroPct}{n/a}
\rqSet{rqTermCOutputLimitExePOneNum}{0}
\rqSet{rqTermCOutputLimitExePOneDen}{108}
\rqSet{rqTermCOutputLimitExePOnePct}{0.0}
\rqSet{rqTermCMemLimitExeAllNum}{0}
\rqSet{rqTermCMemLimitExeAllDen}{35894359}
\rqSet{rqTermCMemLimitExeAllPct}{0.0}
\rqSet{rqTermCMemLimitExeBOneNum}{0}
\rqSet{rqTermCMemLimitExeBOneDen}{29751287}
\rqSet{rqTermCMemLimitExeBOnePct}{0.0}
\rqSet{rqTermCMemLimitExeBTwoNum}{0}
\rqSet{rqTermCMemLimitExeBTwoDen}{6142964}
\rqSet{rqTermCMemLimitExeBTwoPct}{0.0}
\rqSet{rqTermCMemLimitExePZeroNum}{0}
\rqSet{rqTermCMemLimitExePZeroDen}{0}
\rqSet{rqTermCMemLimitExePZeroPct}{n/a}
\rqSet{rqTermCMemLimitExePOneNum}{0}
\rqSet{rqTermCMemLimitExePOneDen}{108}
\rqSet{rqTermCMemLimitExePOnePct}{0.0}
\rqSet{rqTermCInvalidInputExeAllNum}{0}
\rqSet{rqTermCInvalidInputExeAllDen}{35894359}
\rqSet{rqTermCInvalidInputExeAllPct}{0.0}
\rqSet{rqTermCInvalidInputExeBOneNum}{0}
\rqSet{rqTermCInvalidInputExeBOneDen}{29751287}
\rqSet{rqTermCInvalidInputExeBOnePct}{0.0}
\rqSet{rqTermCInvalidInputExeBTwoNum}{0}
\rqSet{rqTermCInvalidInputExeBTwoDen}{6142964}
\rqSet{rqTermCInvalidInputExeBTwoPct}{0.0}
\rqSet{rqTermCInvalidInputExePZeroNum}{0}
\rqSet{rqTermCInvalidInputExePZeroDen}{0}
\rqSet{rqTermCInvalidInputExePZeroPct}{n/a}
\rqSet{rqTermCInvalidInputExePOneNum}{0}
\rqSet{rqTermCInvalidInputExePOneDen}{108}
\rqSet{rqTermCInvalidInputExePOnePct}{0.0}
\rqSet{rqTermRExitExeAllNum}{35737782}
\rqSet{rqTermRExitExeAllDen}{35894359}
\rqSet{rqTermRExitExeAllPct}{99.6}
\rqSet{rqTermRExitExeBOneNum}{29594743}
\rqSet{rqTermRExitExeBOneDen}{29751287}
\rqSet{rqTermRExitExeBOnePct}{99.5}
\rqSet{rqTermRExitExeBTwoNum}{6142931}
\rqSet{rqTermRExitExeBTwoDen}{6142964}
\rqSet{rqTermRExitExeBTwoPct}{100.0}
\rqSet{rqTermRExitExePZeroNum}{0}
\rqSet{rqTermRExitExePZeroDen}{0}
\rqSet{rqTermRExitExePZeroPct}{n/a}
\rqSet{rqTermRExitExePOneNum}{108}
\rqSet{rqTermRExitExePOneDen}{108}
\rqSet{rqTermRExitExePOnePct}{100.0}
\rqSet{rqTermRSignalExeAllNum}{156577}
\rqSet{rqTermRSignalExeAllDen}{35894359}
\rqSet{rqTermRSignalExeAllPct}{0.4}
\rqSet{rqTermRSignalExeBOneNum}{156544}
\rqSet{rqTermRSignalExeBOneDen}{29751287}
\rqSet{rqTermRSignalExeBOnePct}{0.5}
\rqSet{rqTermRSignalExeBTwoNum}{33}
\rqSet{rqTermRSignalExeBTwoDen}{6142964}
\rqSet{rqTermRSignalExeBTwoPct}{0.0}
\rqSet{rqTermRSignalExePZeroNum}{0}
\rqSet{rqTermRSignalExePZeroDen}{0}
\rqSet{rqTermRSignalExePZeroPct}{n/a}
\rqSet{rqTermRSignalExePOneNum}{0}
\rqSet{rqTermRSignalExePOneDen}{108}
\rqSet{rqTermRSignalExePOnePct}{0.0}
\rqSet{rqTermRTimeoutExeAllNum}{0}
\rqSet{rqTermRTimeoutExeAllDen}{35894359}
\rqSet{rqTermRTimeoutExeAllPct}{0.0}
\rqSet{rqTermRTimeoutExeBOneNum}{0}
\rqSet{rqTermRTimeoutExeBOneDen}{29751287}
\rqSet{rqTermRTimeoutExeBOnePct}{0.0}
\rqSet{rqTermRTimeoutExeBTwoNum}{0}
\rqSet{rqTermRTimeoutExeBTwoDen}{6142964}
\rqSet{rqTermRTimeoutExeBTwoPct}{0.0}
\rqSet{rqTermRTimeoutExePZeroNum}{0}
\rqSet{rqTermRTimeoutExePZeroDen}{0}
\rqSet{rqTermRTimeoutExePZeroPct}{n/a}
\rqSet{rqTermRTimeoutExePOneNum}{0}
\rqSet{rqTermRTimeoutExePOneDen}{108}
\rqSet{rqTermRTimeoutExePOnePct}{0.0}
\rqSet{rqTermROutputLimitExeAllNum}{0}
\rqSet{rqTermROutputLimitExeAllDen}{35894359}
\rqSet{rqTermROutputLimitExeAllPct}{0.0}
\rqSet{rqTermROutputLimitExeBOneNum}{0}
\rqSet{rqTermROutputLimitExeBOneDen}{29751287}
\rqSet{rqTermROutputLimitExeBOnePct}{0.0}
\rqSet{rqTermROutputLimitExeBTwoNum}{0}
\rqSet{rqTermROutputLimitExeBTwoDen}{6142964}
\rqSet{rqTermROutputLimitExeBTwoPct}{0.0}
\rqSet{rqTermROutputLimitExePZeroNum}{0}
\rqSet{rqTermROutputLimitExePZeroDen}{0}
\rqSet{rqTermROutputLimitExePZeroPct}{n/a}
\rqSet{rqTermROutputLimitExePOneNum}{0}
\rqSet{rqTermROutputLimitExePOneDen}{108}
\rqSet{rqTermROutputLimitExePOnePct}{0.0}
\rqSet{rqTermRMemLimitExeAllNum}{0}
\rqSet{rqTermRMemLimitExeAllDen}{35894359}
\rqSet{rqTermRMemLimitExeAllPct}{0.0}
\rqSet{rqTermRMemLimitExeBOneNum}{0}
\rqSet{rqTermRMemLimitExeBOneDen}{29751287}
\rqSet{rqTermRMemLimitExeBOnePct}{0.0}
\rqSet{rqTermRMemLimitExeBTwoNum}{0}
\rqSet{rqTermRMemLimitExeBTwoDen}{6142964}
\rqSet{rqTermRMemLimitExeBTwoPct}{0.0}
\rqSet{rqTermRMemLimitExePZeroNum}{0}
\rqSet{rqTermRMemLimitExePZeroDen}{0}
\rqSet{rqTermRMemLimitExePZeroPct}{n/a}
\rqSet{rqTermRMemLimitExePOneNum}{0}
\rqSet{rqTermRMemLimitExePOneDen}{108}
\rqSet{rqTermRMemLimitExePOnePct}{0.0}
\rqSet{rqTermRInvalidInputExeAllNum}{0}
\rqSet{rqTermRInvalidInputExeAllDen}{35894359}
\rqSet{rqTermRInvalidInputExeAllPct}{0.0}
\rqSet{rqTermRInvalidInputExeBOneNum}{0}
\rqSet{rqTermRInvalidInputExeBOneDen}{29751287}
\rqSet{rqTermRInvalidInputExeBOnePct}{0.0}
\rqSet{rqTermRInvalidInputExeBTwoNum}{0}
\rqSet{rqTermRInvalidInputExeBTwoDen}{6142964}
\rqSet{rqTermRInvalidInputExeBTwoPct}{0.0}
\rqSet{rqTermRInvalidInputExePZeroNum}{0}
\rqSet{rqTermRInvalidInputExePZeroDen}{0}
\rqSet{rqTermRInvalidInputExePZeroPct}{n/a}
\rqSet{rqTermRInvalidInputExePOneNum}{0}
\rqSet{rqTermRInvalidInputExePOneDen}{108}
\rqSet{rqTermRInvalidInputExePOnePct}{0.0}
\rqSet{rqFFComboStateFloatExeAllNum}{0}
\rqSet{rqFFComboStateFloatExeAllDen}{0}
\rqSet{rqFFComboStateFloatExeAllPct}{n/a}
\rqSet{rqFFComboStateFloatExeBOneNum}{0}
\rqSet{rqFFComboStateFloatExeBOneDen}{0}
\rqSet{rqFFComboStateFloatExeBOnePct}{n/a}
\rqSet{rqFFComboStateFloatExeBTwoNum}{0}
\rqSet{rqFFComboStateFloatExeBTwoDen}{0}
\rqSet{rqFFComboStateFloatExeBTwoPct}{n/a}
\rqSet{rqFFComboStateFloatExePZeroNum}{0}
\rqSet{rqFFComboStateFloatExePZeroDen}{0}
\rqSet{rqFFComboStateFloatExePZeroPct}{n/a}
\rqSet{rqFFComboStateFloatExePOneNum}{0}
\rqSet{rqFFComboStateFloatExePOneDen}{0}
\rqSet{rqFFComboStateFloatExePOnePct}{n/a}
\rqSet{rqPzeroLibAllNum}{2190}
\rqSet{rqPzeroLibAllDen}{2190}
\rqSet{rqPzeroLibAllKiroNum}{219}
\rqSet{rqPzeroLibAllKiroDen}{219}
\rqSet{rqPzeroLibAllClaudeNum}{219}
\rqSet{rqPzeroLibAllClaudeDen}{219}
\rqSet{rqPzeroLibAllCodexNum}{219}
\rqSet{rqPzeroLibAllCodexDen}{219}
\rqSet{rqPzeroLibAllCTwoRustNum}{219}
\rqSet{rqPzeroLibAllCTwoRustDen}{219}
\rqSet{rqPzeroLibAllLaertesNum}{219}
\rqSet{rqPzeroLibAllLaertesDen}{219}
\rqSet{rqPzeroLibAllCTwoSaferRustNum}{219}
\rqSet{rqPzeroLibAllCTwoSaferRustDen}{219}
\rqSet{rqPzeroLibAllSmartCTwoRustNum}{219}
\rqSet{rqPzeroLibAllSmartCTwoRustDen}{219}
\rqSet{rqPzeroLibAllKimiNum}{219}
\rqSet{rqPzeroLibAllKimiDen}{219}
\rqSet{rqPzeroLibAllGptFiveNum}{219}
\rqSet{rqPzeroLibAllGptFiveDen}{219}
\rqSet{rqPzeroLibAllGeminiNum}{219}
\rqSet{rqPzeroLibAllGeminiDen}{219}
\rqSet{rqPzeroLibBOneNum}{770}
\rqSet{rqPzeroLibBOneDen}{770}
\rqSet{rqPzeroLibBOneKiroNum}{77}
\rqSet{rqPzeroLibBOneKiroDen}{77}
\rqSet{rqPzeroLibBOneClaudeNum}{77}
\rqSet{rqPzeroLibBOneClaudeDen}{77}
\rqSet{rqPzeroLibBOneCodexNum}{77}
\rqSet{rqPzeroLibBOneCodexDen}{77}
\rqSet{rqPzeroLibBOneCTwoRustNum}{77}
\rqSet{rqPzeroLibBOneCTwoRustDen}{77}
\rqSet{rqPzeroLibBOneLaertesNum}{77}
\rqSet{rqPzeroLibBOneLaertesDen}{77}
\rqSet{rqPzeroLibBOneCTwoSaferRustNum}{77}
\rqSet{rqPzeroLibBOneCTwoSaferRustDen}{77}
\rqSet{rqPzeroLibBOneSmartCTwoRustNum}{77}
\rqSet{rqPzeroLibBOneSmartCTwoRustDen}{77}
\rqSet{rqPzeroLibBOneKimiNum}{77}
\rqSet{rqPzeroLibBOneKimiDen}{77}
\rqSet{rqPzeroLibBOneGptFiveNum}{77}
\rqSet{rqPzeroLibBOneGptFiveDen}{77}
\rqSet{rqPzeroLibBOneGeminiNum}{77}
\rqSet{rqPzeroLibBOneGeminiDen}{77}
\rqSet{rqPzeroLibBTwoNum}{660}
\rqSet{rqPzeroLibBTwoDen}{660}
\rqSet{rqPzeroLibBTwoKiroNum}{66}
\rqSet{rqPzeroLibBTwoKiroDen}{66}
\rqSet{rqPzeroLibBTwoClaudeNum}{66}
\rqSet{rqPzeroLibBTwoClaudeDen}{66}
\rqSet{rqPzeroLibBTwoCodexNum}{66}
\rqSet{rqPzeroLibBTwoCodexDen}{66}
\rqSet{rqPzeroLibBTwoCTwoRustNum}{66}
\rqSet{rqPzeroLibBTwoCTwoRustDen}{66}
\rqSet{rqPzeroLibBTwoLaertesNum}{66}
\rqSet{rqPzeroLibBTwoLaertesDen}{66}
\rqSet{rqPzeroLibBTwoCTwoSaferRustNum}{66}
\rqSet{rqPzeroLibBTwoCTwoSaferRustDen}{66}
\rqSet{rqPzeroLibBTwoSmartCTwoRustNum}{66}
\rqSet{rqPzeroLibBTwoSmartCTwoRustDen}{66}
\rqSet{rqPzeroLibBTwoKimiNum}{66}
\rqSet{rqPzeroLibBTwoKimiDen}{66}
\rqSet{rqPzeroLibBTwoGptFiveNum}{66}
\rqSet{rqPzeroLibBTwoGptFiveDen}{66}
\rqSet{rqPzeroLibBTwoGeminiNum}{66}
\rqSet{rqPzeroLibBTwoGeminiDen}{66}
\rqSet{rqPzeroLibPZeroNum}{0}
\rqSet{rqPzeroLibPZeroDen}{0}
\rqSet{rqPzeroLibPZeroKiroNum}{0}
\rqSet{rqPzeroLibPZeroKiroDen}{0}
\rqSet{rqPzeroLibPZeroClaudeNum}{0}
\rqSet{rqPzeroLibPZeroClaudeDen}{0}
\rqSet{rqPzeroLibPZeroCodexNum}{0}
\rqSet{rqPzeroLibPZeroCodexDen}{0}
\rqSet{rqPzeroLibPZeroCTwoRustNum}{0}
\rqSet{rqPzeroLibPZeroCTwoRustDen}{0}
\rqSet{rqPzeroLibPZeroLaertesNum}{0}
\rqSet{rqPzeroLibPZeroLaertesDen}{0}
\rqSet{rqPzeroLibPZeroCTwoSaferRustNum}{0}
\rqSet{rqPzeroLibPZeroCTwoSaferRustDen}{0}
\rqSet{rqPzeroLibPZeroSmartCTwoRustNum}{0}
\rqSet{rqPzeroLibPZeroSmartCTwoRustDen}{0}
\rqSet{rqPzeroLibPZeroKimiNum}{0}
\rqSet{rqPzeroLibPZeroKimiDen}{0}
\rqSet{rqPzeroLibPZeroGptFiveNum}{0}
\rqSet{rqPzeroLibPZeroGptFiveDen}{0}
\rqSet{rqPzeroLibPZeroGeminiNum}{0}
\rqSet{rqPzeroLibPZeroGeminiDen}{0}
\rqSet{rqPzeroLibPOneNum}{760}
\rqSet{rqPzeroLibPOneDen}{760}
\rqSet{rqPzeroLibPOneKiroNum}{76}
\rqSet{rqPzeroLibPOneKiroDen}{76}
\rqSet{rqPzeroLibPOneClaudeNum}{76}
\rqSet{rqPzeroLibPOneClaudeDen}{76}
\rqSet{rqPzeroLibPOneCodexNum}{76}
\rqSet{rqPzeroLibPOneCodexDen}{76}
\rqSet{rqPzeroLibPOneCTwoRustNum}{76}
\rqSet{rqPzeroLibPOneCTwoRustDen}{76}
\rqSet{rqPzeroLibPOneLaertesNum}{76}
\rqSet{rqPzeroLibPOneLaertesDen}{76}
\rqSet{rqPzeroLibPOneCTwoSaferRustNum}{76}
\rqSet{rqPzeroLibPOneCTwoSaferRustDen}{76}
\rqSet{rqPzeroLibPOneSmartCTwoRustNum}{76}
\rqSet{rqPzeroLibPOneSmartCTwoRustDen}{76}
\rqSet{rqPzeroLibPOneKimiNum}{76}
\rqSet{rqPzeroLibPOneKimiDen}{76}
\rqSet{rqPzeroLibPOneGptFiveNum}{76}
\rqSet{rqPzeroLibPOneGptFiveDen}{76}
\rqSet{rqPzeroLibPOneGeminiNum}{76}
\rqSet{rqPzeroLibPOneGeminiDen}{76}
\rqSet{rqPaccLibAllNum}{1281}
\rqSet{rqPaccLibAllDen}{2190}
\rqSet{rqPaccLibAllPct}{58.5}
\rqSet{rqPaccLibAllKiroNum}{217}
\rqSet{rqPaccLibAllKiroDen}{219}
\rqSet{rqPaccLibAllKiroPct}{99.1}
\rqSet{rqPaccLibAllClaudeNum}{206}
\rqSet{rqPaccLibAllClaudeDen}{219}
\rqSet{rqPaccLibAllClaudePct}{94.1}
\rqSet{rqPaccLibAllCodexNum}{171}
\rqSet{rqPaccLibAllCodexDen}{219}
\rqSet{rqPaccLibAllCodexPct}{78.1}
\rqSet{rqPaccLibAllCTwoRustNum}{141}
\rqSet{rqPaccLibAllCTwoRustDen}{219}
\rqSet{rqPaccLibAllCTwoRustPct}{64.4}
\rqSet{rqPaccLibAllLaertesNum}{137}
\rqSet{rqPaccLibAllLaertesDen}{219}
\rqSet{rqPaccLibAllLaertesPct}{62.6}
\rqSet{rqPaccLibAllCTwoSaferRustNum}{105}
\rqSet{rqPaccLibAllCTwoSaferRustDen}{219}
\rqSet{rqPaccLibAllCTwoSaferRustPct}{47.9}
\rqSet{rqPaccLibAllSmartCTwoRustNum}{0}
\rqSet{rqPaccLibAllSmartCTwoRustDen}{219}
\rqSet{rqPaccLibAllSmartCTwoRustPct}{0.0}
\rqSet{rqPaccLibAllKimiNum}{85}
\rqSet{rqPaccLibAllKimiDen}{219}
\rqSet{rqPaccLibAllKimiPct}{38.8}
\rqSet{rqPaccLibAllGptFiveNum}{111}
\rqSet{rqPaccLibAllGptFiveDen}{219}
\rqSet{rqPaccLibAllGptFivePct}{50.7}
\rqSet{rqPaccLibAllGeminiNum}{108}
\rqSet{rqPaccLibAllGeminiDen}{219}
\rqSet{rqPaccLibAllGeminiPct}{49.3}
\rqSet{rqPaccLibBOneNum}{637}
\rqSet{rqPaccLibBOneDen}{770}
\rqSet{rqPaccLibBOnePct}{82.7}
\rqSet{rqPaccLibBOneKiroNum}{77}
\rqSet{rqPaccLibBOneKiroDen}{77}
\rqSet{rqPaccLibBOneKiroPct}{100.0}
\rqSet{rqPaccLibBOneClaudeNum}{77}
\rqSet{rqPaccLibBOneClaudeDen}{77}
\rqSet{rqPaccLibBOneClaudePct}{100.0}
\rqSet{rqPaccLibBOneCodexNum}{77}
\rqSet{rqPaccLibBOneCodexDen}{77}
\rqSet{rqPaccLibBOneCodexPct}{100.0}
\rqSet{rqPaccLibBOneCTwoRustNum}{77}
\rqSet{rqPaccLibBOneCTwoRustDen}{77}
\rqSet{rqPaccLibBOneCTwoRustPct}{100.0}
\rqSet{rqPaccLibBOneLaertesNum}{74}
\rqSet{rqPaccLibBOneLaertesDen}{77}
\rqSet{rqPaccLibBOneLaertesPct}{96.1}
\rqSet{rqPaccLibBOneCTwoSaferRustNum}{60}
\rqSet{rqPaccLibBOneCTwoSaferRustDen}{77}
\rqSet{rqPaccLibBOneCTwoSaferRustPct}{77.9}
\rqSet{rqPaccLibBOneSmartCTwoRustNum}{0}
\rqSet{rqPaccLibBOneSmartCTwoRustDen}{77}
\rqSet{rqPaccLibBOneSmartCTwoRustPct}{0.0}
\rqSet{rqPaccLibBOneKimiNum}{57}
\rqSet{rqPaccLibBOneKimiDen}{77}
\rqSet{rqPaccLibBOneKimiPct}{74.0}
\rqSet{rqPaccLibBOneGptFiveNum}{67}
\rqSet{rqPaccLibBOneGptFiveDen}{77}
\rqSet{rqPaccLibBOneGptFivePct}{87.0}
\rqSet{rqPaccLibBOneGeminiNum}{71}
\rqSet{rqPaccLibBOneGeminiDen}{77}
\rqSet{rqPaccLibBOneGeminiPct}{92.2}
\rqSet{rqPaccLibBTwoNum}{472}
\rqSet{rqPaccLibBTwoDen}{660}
\rqSet{rqPaccLibBTwoPct}{71.5}
\rqSet{rqPaccLibBTwoKiroNum}{64}
\rqSet{rqPaccLibBTwoKiroDen}{66}
\rqSet{rqPaccLibBTwoKiroPct}{97.0}
\rqSet{rqPaccLibBTwoClaudeNum}{64}
\rqSet{rqPaccLibBTwoClaudeDen}{66}
\rqSet{rqPaccLibBTwoClaudePct}{97.0}
\rqSet{rqPaccLibBTwoCodexNum}{63}
\rqSet{rqPaccLibBTwoCodexDen}{66}
\rqSet{rqPaccLibBTwoCodexPct}{95.5}
\rqSet{rqPaccLibBTwoCTwoRustNum}{64}
\rqSet{rqPaccLibBTwoCTwoRustDen}{66}
\rqSet{rqPaccLibBTwoCTwoRustPct}{97.0}
\rqSet{rqPaccLibBTwoLaertesNum}{63}
\rqSet{rqPaccLibBTwoLaertesDen}{66}
\rqSet{rqPaccLibBTwoLaertesPct}{95.5}
\rqSet{rqPaccLibBTwoCTwoSaferRustNum}{45}
\rqSet{rqPaccLibBTwoCTwoSaferRustDen}{66}
\rqSet{rqPaccLibBTwoCTwoSaferRustPct}{68.2}
\rqSet{rqPaccLibBTwoSmartCTwoRustNum}{0}
\rqSet{rqPaccLibBTwoSmartCTwoRustDen}{66}
\rqSet{rqPaccLibBTwoSmartCTwoRustPct}{0.0}
\rqSet{rqPaccLibBTwoKimiNum}{28}
\rqSet{rqPaccLibBTwoKimiDen}{66}
\rqSet{rqPaccLibBTwoKimiPct}{42.4}
\rqSet{rqPaccLibBTwoGptFiveNum}{44}
\rqSet{rqPaccLibBTwoGptFiveDen}{66}
\rqSet{rqPaccLibBTwoGptFivePct}{66.7}
\rqSet{rqPaccLibBTwoGeminiNum}{37}
\rqSet{rqPaccLibBTwoGeminiDen}{66}
\rqSet{rqPaccLibBTwoGeminiPct}{56.1}
\rqSet{rqPaccLibPZeroNum}{0}
\rqSet{rqPaccLibPZeroDen}{0}
\rqSet{rqPaccLibPZeroPct}{n/a}
\rqSet{rqPaccLibPZeroKiroNum}{0}
\rqSet{rqPaccLibPZeroKiroDen}{0}
\rqSet{rqPaccLibPZeroKiroPct}{n/a}
\rqSet{rqPaccLibPZeroClaudeNum}{0}
\rqSet{rqPaccLibPZeroClaudeDen}{0}
\rqSet{rqPaccLibPZeroClaudePct}{n/a}
\rqSet{rqPaccLibPZeroCodexNum}{0}
\rqSet{rqPaccLibPZeroCodexDen}{0}
\rqSet{rqPaccLibPZeroCodexPct}{n/a}
\rqSet{rqPaccLibPZeroCTwoRustNum}{0}
\rqSet{rqPaccLibPZeroCTwoRustDen}{0}
\rqSet{rqPaccLibPZeroCTwoRustPct}{n/a}
\rqSet{rqPaccLibPZeroLaertesNum}{0}
\rqSet{rqPaccLibPZeroLaertesDen}{0}
\rqSet{rqPaccLibPZeroLaertesPct}{n/a}
\rqSet{rqPaccLibPZeroCTwoSaferRustNum}{0}
\rqSet{rqPaccLibPZeroCTwoSaferRustDen}{0}
\rqSet{rqPaccLibPZeroCTwoSaferRustPct}{n/a}
\rqSet{rqPaccLibPZeroSmartCTwoRustNum}{0}
\rqSet{rqPaccLibPZeroSmartCTwoRustDen}{0}
\rqSet{rqPaccLibPZeroSmartCTwoRustPct}{n/a}
\rqSet{rqPaccLibPZeroKimiNum}{0}
\rqSet{rqPaccLibPZeroKimiDen}{0}
\rqSet{rqPaccLibPZeroKimiPct}{n/a}
\rqSet{rqPaccLibPZeroGptFiveNum}{0}
\rqSet{rqPaccLibPZeroGptFiveDen}{0}
\rqSet{rqPaccLibPZeroGptFivePct}{n/a}
\rqSet{rqPaccLibPZeroGeminiNum}{0}
\rqSet{rqPaccLibPZeroGeminiDen}{0}
\rqSet{rqPaccLibPZeroGeminiPct}{n/a}
\rqSet{rqPaccLibPOneNum}{172}
\rqSet{rqPaccLibPOneDen}{760}
\rqSet{rqPaccLibPOnePct}{22.6}
\rqSet{rqPaccLibPOneKiroNum}{76}
\rqSet{rqPaccLibPOneKiroDen}{76}
\rqSet{rqPaccLibPOneKiroPct}{100.0}
\rqSet{rqPaccLibPOneClaudeNum}{65}
\rqSet{rqPaccLibPOneClaudeDen}{76}
\rqSet{rqPaccLibPOneClaudePct}{85.5}
\rqSet{rqPaccLibPOneCodexNum}{31}
\rqSet{rqPaccLibPOneCodexDen}{76}
\rqSet{rqPaccLibPOneCodexPct}{40.8}
\rqSet{rqPaccLibPOneCTwoRustNum}{0}
\rqSet{rqPaccLibPOneCTwoRustDen}{76}
\rqSet{rqPaccLibPOneCTwoRustPct}{0.0}
\rqSet{rqPaccLibPOneLaertesNum}{0}
\rqSet{rqPaccLibPOneLaertesDen}{76}
\rqSet{rqPaccLibPOneLaertesPct}{0.0}
\rqSet{rqPaccLibPOneCTwoSaferRustNum}{0}
\rqSet{rqPaccLibPOneCTwoSaferRustDen}{76}
\rqSet{rqPaccLibPOneCTwoSaferRustPct}{0.0}
\rqSet{rqPaccLibPOneSmartCTwoRustNum}{0}
\rqSet{rqPaccLibPOneSmartCTwoRustDen}{76}
\rqSet{rqPaccLibPOneSmartCTwoRustPct}{0.0}
\rqSet{rqPaccLibPOneKimiNum}{0}
\rqSet{rqPaccLibPOneKimiDen}{76}
\rqSet{rqPaccLibPOneKimiPct}{0.0}
\rqSet{rqPaccLibPOneGptFiveNum}{0}
\rqSet{rqPaccLibPOneGptFiveDen}{76}
\rqSet{rqPaccLibPOneGptFivePct}{0.0}
\rqSet{rqPaccLibPOneGeminiNum}{0}
\rqSet{rqPaccLibPOneGeminiDen}{76}
\rqSet{rqPaccLibPOneGeminiPct}{0.0}
\rqSet{rqMissingLibAllNum}{76}
\rqSet{rqMissingLibAllDen}{2190}
\rqSet{rqMissingLibAllPct}{3.5}
\rqSet{rqMissingLibAllKiroNum}{0}
\rqSet{rqMissingLibAllKiroDen}{219}
\rqSet{rqMissingLibAllKiroPct}{0.0}
\rqSet{rqMissingLibAllClaudeNum}{0}
\rqSet{rqMissingLibAllClaudeDen}{219}
\rqSet{rqMissingLibAllClaudePct}{0.0}
\rqSet{rqMissingLibAllCodexNum}{0}
\rqSet{rqMissingLibAllCodexDen}{219}
\rqSet{rqMissingLibAllCodexPct}{0.0}
\rqSet{rqMissingLibAllCTwoRustNum}{0}
\rqSet{rqMissingLibAllCTwoRustDen}{219}
\rqSet{rqMissingLibAllCTwoRustPct}{0.0}
\rqSet{rqMissingLibAllLaertesNum}{0}
\rqSet{rqMissingLibAllLaertesDen}{219}
\rqSet{rqMissingLibAllLaertesPct}{0.0}
\rqSet{rqMissingLibAllCTwoSaferRustNum}{0}
\rqSet{rqMissingLibAllCTwoSaferRustDen}{219}
\rqSet{rqMissingLibAllCTwoSaferRustPct}{0.0}
\rqSet{rqMissingLibAllSmartCTwoRustNum}{0}
\rqSet{rqMissingLibAllSmartCTwoRustDen}{219}
\rqSet{rqMissingLibAllSmartCTwoRustPct}{0.0}
\rqSet{rqMissingLibAllKimiNum}{76}
\rqSet{rqMissingLibAllKimiDen}{219}
\rqSet{rqMissingLibAllKimiPct}{34.7}
\rqSet{rqMissingLibAllGptFiveNum}{0}
\rqSet{rqMissingLibAllGptFiveDen}{219}
\rqSet{rqMissingLibAllGptFivePct}{0.0}
\rqSet{rqMissingLibAllGeminiNum}{0}
\rqSet{rqMissingLibAllGeminiDen}{219}
\rqSet{rqMissingLibAllGeminiPct}{0.0}
\rqSet{rqMissingLibBOneNum}{0}
\rqSet{rqMissingLibBOneDen}{770}
\rqSet{rqMissingLibBOnePct}{0.0}
\rqSet{rqMissingLibBOneKiroNum}{0}
\rqSet{rqMissingLibBOneKiroDen}{77}
\rqSet{rqMissingLibBOneKiroPct}{0.0}
\rqSet{rqMissingLibBOneClaudeNum}{0}
\rqSet{rqMissingLibBOneClaudeDen}{77}
\rqSet{rqMissingLibBOneClaudePct}{0.0}
\rqSet{rqMissingLibBOneCodexNum}{0}
\rqSet{rqMissingLibBOneCodexDen}{77}
\rqSet{rqMissingLibBOneCodexPct}{0.0}
\rqSet{rqMissingLibBOneCTwoRustNum}{0}
\rqSet{rqMissingLibBOneCTwoRustDen}{77}
\rqSet{rqMissingLibBOneCTwoRustPct}{0.0}
\rqSet{rqMissingLibBOneLaertesNum}{0}
\rqSet{rqMissingLibBOneLaertesDen}{77}
\rqSet{rqMissingLibBOneLaertesPct}{0.0}
\rqSet{rqMissingLibBOneCTwoSaferRustNum}{0}
\rqSet{rqMissingLibBOneCTwoSaferRustDen}{77}
\rqSet{rqMissingLibBOneCTwoSaferRustPct}{0.0}
\rqSet{rqMissingLibBOneSmartCTwoRustNum}{0}
\rqSet{rqMissingLibBOneSmartCTwoRustDen}{77}
\rqSet{rqMissingLibBOneSmartCTwoRustPct}{0.0}
\rqSet{rqMissingLibBOneKimiNum}{0}
\rqSet{rqMissingLibBOneKimiDen}{77}
\rqSet{rqMissingLibBOneKimiPct}{0.0}
\rqSet{rqMissingLibBOneGptFiveNum}{0}
\rqSet{rqMissingLibBOneGptFiveDen}{77}
\rqSet{rqMissingLibBOneGptFivePct}{0.0}
\rqSet{rqMissingLibBOneGeminiNum}{0}
\rqSet{rqMissingLibBOneGeminiDen}{77}
\rqSet{rqMissingLibBOneGeminiPct}{0.0}
\rqSet{rqMissingLibBTwoNum}{0}
\rqSet{rqMissingLibBTwoDen}{660}
\rqSet{rqMissingLibBTwoPct}{0.0}
\rqSet{rqMissingLibBTwoKiroNum}{0}
\rqSet{rqMissingLibBTwoKiroDen}{66}
\rqSet{rqMissingLibBTwoKiroPct}{0.0}
\rqSet{rqMissingLibBTwoClaudeNum}{0}
\rqSet{rqMissingLibBTwoClaudeDen}{66}
\rqSet{rqMissingLibBTwoClaudePct}{0.0}
\rqSet{rqMissingLibBTwoCodexNum}{0}
\rqSet{rqMissingLibBTwoCodexDen}{66}
\rqSet{rqMissingLibBTwoCodexPct}{0.0}
\rqSet{rqMissingLibBTwoCTwoRustNum}{0}
\rqSet{rqMissingLibBTwoCTwoRustDen}{66}
\rqSet{rqMissingLibBTwoCTwoRustPct}{0.0}
\rqSet{rqMissingLibBTwoLaertesNum}{0}
\rqSet{rqMissingLibBTwoLaertesDen}{66}
\rqSet{rqMissingLibBTwoLaertesPct}{0.0}
\rqSet{rqMissingLibBTwoCTwoSaferRustNum}{0}
\rqSet{rqMissingLibBTwoCTwoSaferRustDen}{66}
\rqSet{rqMissingLibBTwoCTwoSaferRustPct}{0.0}
\rqSet{rqMissingLibBTwoSmartCTwoRustNum}{0}
\rqSet{rqMissingLibBTwoSmartCTwoRustDen}{66}
\rqSet{rqMissingLibBTwoSmartCTwoRustPct}{0.0}
\rqSet{rqMissingLibBTwoKimiNum}{0}
\rqSet{rqMissingLibBTwoKimiDen}{66}
\rqSet{rqMissingLibBTwoKimiPct}{0.0}
\rqSet{rqMissingLibBTwoGptFiveNum}{0}
\rqSet{rqMissingLibBTwoGptFiveDen}{66}
\rqSet{rqMissingLibBTwoGptFivePct}{0.0}
\rqSet{rqMissingLibBTwoGeminiNum}{0}
\rqSet{rqMissingLibBTwoGeminiDen}{66}
\rqSet{rqMissingLibBTwoGeminiPct}{0.0}
\rqSet{rqMissingLibPZeroNum}{0}
\rqSet{rqMissingLibPZeroDen}{0}
\rqSet{rqMissingLibPZeroPct}{n/a}
\rqSet{rqMissingLibPZeroKiroNum}{0}
\rqSet{rqMissingLibPZeroKiroDen}{0}
\rqSet{rqMissingLibPZeroKiroPct}{n/a}
\rqSet{rqMissingLibPZeroClaudeNum}{0}
\rqSet{rqMissingLibPZeroClaudeDen}{0}
\rqSet{rqMissingLibPZeroClaudePct}{n/a}
\rqSet{rqMissingLibPZeroCodexNum}{0}
\rqSet{rqMissingLibPZeroCodexDen}{0}
\rqSet{rqMissingLibPZeroCodexPct}{n/a}
\rqSet{rqMissingLibPZeroCTwoRustNum}{0}
\rqSet{rqMissingLibPZeroCTwoRustDen}{0}
\rqSet{rqMissingLibPZeroCTwoRustPct}{n/a}
\rqSet{rqMissingLibPZeroLaertesNum}{0}
\rqSet{rqMissingLibPZeroLaertesDen}{0}
\rqSet{rqMissingLibPZeroLaertesPct}{n/a}
\rqSet{rqMissingLibPZeroCTwoSaferRustNum}{0}
\rqSet{rqMissingLibPZeroCTwoSaferRustDen}{0}
\rqSet{rqMissingLibPZeroCTwoSaferRustPct}{n/a}
\rqSet{rqMissingLibPZeroSmartCTwoRustNum}{0}
\rqSet{rqMissingLibPZeroSmartCTwoRustDen}{0}
\rqSet{rqMissingLibPZeroSmartCTwoRustPct}{n/a}
\rqSet{rqMissingLibPZeroKimiNum}{0}
\rqSet{rqMissingLibPZeroKimiDen}{0}
\rqSet{rqMissingLibPZeroKimiPct}{n/a}
\rqSet{rqMissingLibPZeroGptFiveNum}{0}
\rqSet{rqMissingLibPZeroGptFiveDen}{0}
\rqSet{rqMissingLibPZeroGptFivePct}{n/a}
\rqSet{rqMissingLibPZeroGeminiNum}{0}
\rqSet{rqMissingLibPZeroGeminiDen}{0}
\rqSet{rqMissingLibPZeroGeminiPct}{n/a}
\rqSet{rqMissingLibPOneNum}{76}
\rqSet{rqMissingLibPOneDen}{760}
\rqSet{rqMissingLibPOnePct}{10.0}
\rqSet{rqMissingLibPOneKiroNum}{0}
\rqSet{rqMissingLibPOneKiroDen}{76}
\rqSet{rqMissingLibPOneKiroPct}{0.0}
\rqSet{rqMissingLibPOneClaudeNum}{0}
\rqSet{rqMissingLibPOneClaudeDen}{76}
\rqSet{rqMissingLibPOneClaudePct}{0.0}
\rqSet{rqMissingLibPOneCodexNum}{0}
\rqSet{rqMissingLibPOneCodexDen}{76}
\rqSet{rqMissingLibPOneCodexPct}{0.0}
\rqSet{rqMissingLibPOneCTwoRustNum}{0}
\rqSet{rqMissingLibPOneCTwoRustDen}{76}
\rqSet{rqMissingLibPOneCTwoRustPct}{0.0}
\rqSet{rqMissingLibPOneLaertesNum}{0}
\rqSet{rqMissingLibPOneLaertesDen}{76}
\rqSet{rqMissingLibPOneLaertesPct}{0.0}
\rqSet{rqMissingLibPOneCTwoSaferRustNum}{0}
\rqSet{rqMissingLibPOneCTwoSaferRustDen}{76}
\rqSet{rqMissingLibPOneCTwoSaferRustPct}{0.0}
\rqSet{rqMissingLibPOneSmartCTwoRustNum}{0}
\rqSet{rqMissingLibPOneSmartCTwoRustDen}{76}
\rqSet{rqMissingLibPOneSmartCTwoRustPct}{0.0}
\rqSet{rqMissingLibPOneKimiNum}{76}
\rqSet{rqMissingLibPOneKimiDen}{76}
\rqSet{rqMissingLibPOneKimiPct}{100.0}
\rqSet{rqMissingLibPOneGptFiveNum}{0}
\rqSet{rqMissingLibPOneGptFiveDen}{76}
\rqSet{rqMissingLibPOneGptFivePct}{0.0}
\rqSet{rqMissingLibPOneGeminiNum}{0}
\rqSet{rqMissingLibPOneGeminiDen}{76}
\rqSet{rqMissingLibPOneGeminiPct}{0.0}
\rqSet{rqRZeroBuildLibAllNum}{695}
\rqSet{rqRZeroBuildLibAllDen}{2190}
\rqSet{rqRZeroBuildLibAllPct}{31.7}
\rqSet{rqRZeroBuildLibAllKiroNum}{0}
\rqSet{rqRZeroBuildLibAllKiroDen}{219}
\rqSet{rqRZeroBuildLibAllKiroPct}{0.0}
\rqSet{rqRZeroBuildLibAllClaudeNum}{0}
\rqSet{rqRZeroBuildLibAllClaudeDen}{219}
\rqSet{rqRZeroBuildLibAllClaudePct}{0.0}
\rqSet{rqRZeroBuildLibAllCodexNum}{1}
\rqSet{rqRZeroBuildLibAllCodexDen}{219}
\rqSet{rqRZeroBuildLibAllCodexPct}{0.5}
\rqSet{rqRZeroBuildLibAllCTwoRustNum}{77}
\rqSet{rqRZeroBuildLibAllCTwoRustDen}{219}
\rqSet{rqRZeroBuildLibAllCTwoRustPct}{35.2}
\rqSet{rqRZeroBuildLibAllLaertesNum}{81}
\rqSet{rqRZeroBuildLibAllLaertesDen}{219}
\rqSet{rqRZeroBuildLibAllLaertesPct}{37.0}
\rqSet{rqRZeroBuildLibAllCTwoSaferRustNum}{89}
\rqSet{rqRZeroBuildLibAllCTwoSaferRustDen}{219}
\rqSet{rqRZeroBuildLibAllCTwoSaferRustPct}{40.6}
\rqSet{rqRZeroBuildLibAllSmartCTwoRustNum}{219}
\rqSet{rqRZeroBuildLibAllSmartCTwoRustDen}{219}
\rqSet{rqRZeroBuildLibAllSmartCTwoRustPct}{100.0}
\rqSet{rqRZeroBuildLibAllKimiNum}{37}
\rqSet{rqRZeroBuildLibAllKimiDen}{219}
\rqSet{rqRZeroBuildLibAllKimiPct}{16.9}
\rqSet{rqRZeroBuildLibAllGptFiveNum}{93}
\rqSet{rqRZeroBuildLibAllGptFiveDen}{219}
\rqSet{rqRZeroBuildLibAllGptFivePct}{42.5}
\rqSet{rqRZeroBuildLibAllGeminiNum}{98}
\rqSet{rqRZeroBuildLibAllGeminiDen}{219}
\rqSet{rqRZeroBuildLibAllGeminiPct}{44.7}
\rqSet{rqRZeroBuildLibBOneNum}{95}
\rqSet{rqRZeroBuildLibBOneDen}{770}
\rqSet{rqRZeroBuildLibBOnePct}{12.3}
\rqSet{rqRZeroBuildLibBOneKiroNum}{0}
\rqSet{rqRZeroBuildLibBOneKiroDen}{77}
\rqSet{rqRZeroBuildLibBOneKiroPct}{0.0}
\rqSet{rqRZeroBuildLibBOneClaudeNum}{0}
\rqSet{rqRZeroBuildLibBOneClaudeDen}{77}
\rqSet{rqRZeroBuildLibBOneClaudePct}{0.0}
\rqSet{rqRZeroBuildLibBOneCodexNum}{0}
\rqSet{rqRZeroBuildLibBOneCodexDen}{77}
\rqSet{rqRZeroBuildLibBOneCodexPct}{0.0}
\rqSet{rqRZeroBuildLibBOneCTwoRustNum}{0}
\rqSet{rqRZeroBuildLibBOneCTwoRustDen}{77}
\rqSet{rqRZeroBuildLibBOneCTwoRustPct}{0.0}
\rqSet{rqRZeroBuildLibBOneLaertesNum}{3}
\rqSet{rqRZeroBuildLibBOneLaertesDen}{77}
\rqSet{rqRZeroBuildLibBOneLaertesPct}{3.9}
\rqSet{rqRZeroBuildLibBOneCTwoSaferRustNum}{0}
\rqSet{rqRZeroBuildLibBOneCTwoSaferRustDen}{77}
\rqSet{rqRZeroBuildLibBOneCTwoSaferRustPct}{0.0}
\rqSet{rqRZeroBuildLibBOneSmartCTwoRustNum}{77}
\rqSet{rqRZeroBuildLibBOneSmartCTwoRustDen}{77}
\rqSet{rqRZeroBuildLibBOneSmartCTwoRustPct}{100.0}
\rqSet{rqRZeroBuildLibBOneKimiNum}{9}
\rqSet{rqRZeroBuildLibBOneKimiDen}{77}
\rqSet{rqRZeroBuildLibBOneKimiPct}{11.7}
\rqSet{rqRZeroBuildLibBOneGptFiveNum}{3}
\rqSet{rqRZeroBuildLibBOneGptFiveDen}{77}
\rqSet{rqRZeroBuildLibBOneGptFivePct}{3.9}
\rqSet{rqRZeroBuildLibBOneGeminiNum}{3}
\rqSet{rqRZeroBuildLibBOneGeminiDen}{77}
\rqSet{rqRZeroBuildLibBOneGeminiPct}{3.9}
\rqSet{rqRZeroBuildLibBTwoNum}{144}
\rqSet{rqRZeroBuildLibBTwoDen}{660}
\rqSet{rqRZeroBuildLibBTwoPct}{21.8}
\rqSet{rqRZeroBuildLibBTwoKiroNum}{0}
\rqSet{rqRZeroBuildLibBTwoKiroDen}{66}
\rqSet{rqRZeroBuildLibBTwoKiroPct}{0.0}
\rqSet{rqRZeroBuildLibBTwoClaudeNum}{0}
\rqSet{rqRZeroBuildLibBTwoClaudeDen}{66}
\rqSet{rqRZeroBuildLibBTwoClaudePct}{0.0}
\rqSet{rqRZeroBuildLibBTwoCodexNum}{1}
\rqSet{rqRZeroBuildLibBTwoCodexDen}{66}
\rqSet{rqRZeroBuildLibBTwoCodexPct}{1.5}
\rqSet{rqRZeroBuildLibBTwoCTwoRustNum}{1}
\rqSet{rqRZeroBuildLibBTwoCTwoRustDen}{66}
\rqSet{rqRZeroBuildLibBTwoCTwoRustPct}{1.5}
\rqSet{rqRZeroBuildLibBTwoLaertesNum}{2}
\rqSet{rqRZeroBuildLibBTwoLaertesDen}{66}
\rqSet{rqRZeroBuildLibBTwoLaertesPct}{3.0}
\rqSet{rqRZeroBuildLibBTwoCTwoSaferRustNum}{13}
\rqSet{rqRZeroBuildLibBTwoCTwoSaferRustDen}{66}
\rqSet{rqRZeroBuildLibBTwoCTwoSaferRustPct}{19.7}
\rqSet{rqRZeroBuildLibBTwoSmartCTwoRustNum}{66}
\rqSet{rqRZeroBuildLibBTwoSmartCTwoRustDen}{66}
\rqSet{rqRZeroBuildLibBTwoSmartCTwoRustPct}{100.0}
\rqSet{rqRZeroBuildLibBTwoKimiNum}{28}
\rqSet{rqRZeroBuildLibBTwoKimiDen}{66}
\rqSet{rqRZeroBuildLibBTwoKimiPct}{42.4}
\rqSet{rqRZeroBuildLibBTwoGptFiveNum}{14}
\rqSet{rqRZeroBuildLibBTwoGptFiveDen}{66}
\rqSet{rqRZeroBuildLibBTwoGptFivePct}{21.2}
\rqSet{rqRZeroBuildLibBTwoGeminiNum}{19}
\rqSet{rqRZeroBuildLibBTwoGeminiDen}{66}
\rqSet{rqRZeroBuildLibBTwoGeminiPct}{28.8}
\rqSet{rqRZeroBuildLibPZeroNum}{0}
\rqSet{rqRZeroBuildLibPZeroDen}{0}
\rqSet{rqRZeroBuildLibPZeroPct}{n/a}
\rqSet{rqRZeroBuildLibPZeroKiroNum}{0}
\rqSet{rqRZeroBuildLibPZeroKiroDen}{0}
\rqSet{rqRZeroBuildLibPZeroKiroPct}{n/a}
\rqSet{rqRZeroBuildLibPZeroClaudeNum}{0}
\rqSet{rqRZeroBuildLibPZeroClaudeDen}{0}
\rqSet{rqRZeroBuildLibPZeroClaudePct}{n/a}
\rqSet{rqRZeroBuildLibPZeroCodexNum}{0}
\rqSet{rqRZeroBuildLibPZeroCodexDen}{0}
\rqSet{rqRZeroBuildLibPZeroCodexPct}{n/a}
\rqSet{rqRZeroBuildLibPZeroCTwoRustNum}{0}
\rqSet{rqRZeroBuildLibPZeroCTwoRustDen}{0}
\rqSet{rqRZeroBuildLibPZeroCTwoRustPct}{n/a}
\rqSet{rqRZeroBuildLibPZeroLaertesNum}{0}
\rqSet{rqRZeroBuildLibPZeroLaertesDen}{0}
\rqSet{rqRZeroBuildLibPZeroLaertesPct}{n/a}
\rqSet{rqRZeroBuildLibPZeroCTwoSaferRustNum}{0}
\rqSet{rqRZeroBuildLibPZeroCTwoSaferRustDen}{0}
\rqSet{rqRZeroBuildLibPZeroCTwoSaferRustPct}{n/a}
\rqSet{rqRZeroBuildLibPZeroSmartCTwoRustNum}{0}
\rqSet{rqRZeroBuildLibPZeroSmartCTwoRustDen}{0}
\rqSet{rqRZeroBuildLibPZeroSmartCTwoRustPct}{n/a}
\rqSet{rqRZeroBuildLibPZeroKimiNum}{0}
\rqSet{rqRZeroBuildLibPZeroKimiDen}{0}
\rqSet{rqRZeroBuildLibPZeroKimiPct}{n/a}
\rqSet{rqRZeroBuildLibPZeroGptFiveNum}{0}
\rqSet{rqRZeroBuildLibPZeroGptFiveDen}{0}
\rqSet{rqRZeroBuildLibPZeroGptFivePct}{n/a}
\rqSet{rqRZeroBuildLibPZeroGeminiNum}{0}
\rqSet{rqRZeroBuildLibPZeroGeminiDen}{0}
\rqSet{rqRZeroBuildLibPZeroGeminiPct}{n/a}
\rqSet{rqRZeroBuildLibPOneNum}{456}
\rqSet{rqRZeroBuildLibPOneDen}{760}
\rqSet{rqRZeroBuildLibPOnePct}{60.0}
\rqSet{rqRZeroBuildLibPOneKiroNum}{0}
\rqSet{rqRZeroBuildLibPOneKiroDen}{76}
\rqSet{rqRZeroBuildLibPOneKiroPct}{0.0}
\rqSet{rqRZeroBuildLibPOneClaudeNum}{0}
\rqSet{rqRZeroBuildLibPOneClaudeDen}{76}
\rqSet{rqRZeroBuildLibPOneClaudePct}{0.0}
\rqSet{rqRZeroBuildLibPOneCodexNum}{0}
\rqSet{rqRZeroBuildLibPOneCodexDen}{76}
\rqSet{rqRZeroBuildLibPOneCodexPct}{0.0}
\rqSet{rqRZeroBuildLibPOneCTwoRustNum}{76}
\rqSet{rqRZeroBuildLibPOneCTwoRustDen}{76}
\rqSet{rqRZeroBuildLibPOneCTwoRustPct}{100.0}
\rqSet{rqRZeroBuildLibPOneLaertesNum}{76}
\rqSet{rqRZeroBuildLibPOneLaertesDen}{76}
\rqSet{rqRZeroBuildLibPOneLaertesPct}{100.0}
\rqSet{rqRZeroBuildLibPOneCTwoSaferRustNum}{76}
\rqSet{rqRZeroBuildLibPOneCTwoSaferRustDen}{76}
\rqSet{rqRZeroBuildLibPOneCTwoSaferRustPct}{100.0}
\rqSet{rqRZeroBuildLibPOneSmartCTwoRustNum}{76}
\rqSet{rqRZeroBuildLibPOneSmartCTwoRustDen}{76}
\rqSet{rqRZeroBuildLibPOneSmartCTwoRustPct}{100.0}
\rqSet{rqRZeroBuildLibPOneKimiNum}{0}
\rqSet{rqRZeroBuildLibPOneKimiDen}{76}
\rqSet{rqRZeroBuildLibPOneKimiPct}{0.0}
\rqSet{rqRZeroBuildLibPOneGptFiveNum}{76}
\rqSet{rqRZeroBuildLibPOneGptFiveDen}{76}
\rqSet{rqRZeroBuildLibPOneGptFivePct}{100.0}
\rqSet{rqRZeroBuildLibPOneGeminiNum}{76}
\rqSet{rqRZeroBuildLibPOneGeminiDen}{76}
\rqSet{rqRZeroBuildLibPOneGeminiPct}{100.0}
\rqSet{rqRZeroTestLibAllNum}{138}
\rqSet{rqRZeroTestLibAllDen}{2190}
\rqSet{rqRZeroTestLibAllPct}{6.3}
\rqSet{rqRZeroTestLibAllKiroNum}{2}
\rqSet{rqRZeroTestLibAllKiroDen}{219}
\rqSet{rqRZeroTestLibAllKiroPct}{0.9}
\rqSet{rqRZeroTestLibAllClaudeNum}{13}
\rqSet{rqRZeroTestLibAllClaudeDen}{219}
\rqSet{rqRZeroTestLibAllClaudePct}{5.9}
\rqSet{rqRZeroTestLibAllCodexNum}{47}
\rqSet{rqRZeroTestLibAllCodexDen}{219}
\rqSet{rqRZeroTestLibAllCodexPct}{21.5}
\rqSet{rqRZeroTestLibAllCTwoRustNum}{1}
\rqSet{rqRZeroTestLibAllCTwoRustDen}{219}
\rqSet{rqRZeroTestLibAllCTwoRustPct}{0.5}
\rqSet{rqRZeroTestLibAllLaertesNum}{1}
\rqSet{rqRZeroTestLibAllLaertesDen}{219}
\rqSet{rqRZeroTestLibAllLaertesPct}{0.5}
\rqSet{rqRZeroTestLibAllCTwoSaferRustNum}{25}
\rqSet{rqRZeroTestLibAllCTwoSaferRustDen}{219}
\rqSet{rqRZeroTestLibAllCTwoSaferRustPct}{11.4}
\rqSet{rqRZeroTestLibAllSmartCTwoRustNum}{0}
\rqSet{rqRZeroTestLibAllSmartCTwoRustDen}{219}
\rqSet{rqRZeroTestLibAllSmartCTwoRustPct}{0.0}
\rqSet{rqRZeroTestLibAllKimiNum}{21}
\rqSet{rqRZeroTestLibAllKimiDen}{219}
\rqSet{rqRZeroTestLibAllKimiPct}{9.6}
\rqSet{rqRZeroTestLibAllGptFiveNum}{15}
\rqSet{rqRZeroTestLibAllGptFiveDen}{219}
\rqSet{rqRZeroTestLibAllGptFivePct}{6.8}
\rqSet{rqRZeroTestLibAllGeminiNum}{13}
\rqSet{rqRZeroTestLibAllGeminiDen}{219}
\rqSet{rqRZeroTestLibAllGeminiPct}{5.9}
\rqSet{rqRZeroTestLibBOneNum}{38}
\rqSet{rqRZeroTestLibBOneDen}{770}
\rqSet{rqRZeroTestLibBOnePct}{4.9}
\rqSet{rqRZeroTestLibBOneKiroNum}{0}
\rqSet{rqRZeroTestLibBOneKiroDen}{77}
\rqSet{rqRZeroTestLibBOneKiroPct}{0.0}
\rqSet{rqRZeroTestLibBOneClaudeNum}{0}
\rqSet{rqRZeroTestLibBOneClaudeDen}{77}
\rqSet{rqRZeroTestLibBOneClaudePct}{0.0}
\rqSet{rqRZeroTestLibBOneCodexNum}{0}
\rqSet{rqRZeroTestLibBOneCodexDen}{77}
\rqSet{rqRZeroTestLibBOneCodexPct}{0.0}
\rqSet{rqRZeroTestLibBOneCTwoRustNum}{0}
\rqSet{rqRZeroTestLibBOneCTwoRustDen}{77}
\rqSet{rqRZeroTestLibBOneCTwoRustPct}{0.0}
\rqSet{rqRZeroTestLibBOneLaertesNum}{0}
\rqSet{rqRZeroTestLibBOneLaertesDen}{77}
\rqSet{rqRZeroTestLibBOneLaertesPct}{0.0}
\rqSet{rqRZeroTestLibBOneCTwoSaferRustNum}{17}
\rqSet{rqRZeroTestLibBOneCTwoSaferRustDen}{77}
\rqSet{rqRZeroTestLibBOneCTwoSaferRustPct}{22.1}
\rqSet{rqRZeroTestLibBOneSmartCTwoRustNum}{0}
\rqSet{rqRZeroTestLibBOneSmartCTwoRustDen}{77}
\rqSet{rqRZeroTestLibBOneSmartCTwoRustPct}{0.0}
\rqSet{rqRZeroTestLibBOneKimiNum}{11}
\rqSet{rqRZeroTestLibBOneKimiDen}{77}
\rqSet{rqRZeroTestLibBOneKimiPct}{14.3}
\rqSet{rqRZeroTestLibBOneGptFiveNum}{7}
\rqSet{rqRZeroTestLibBOneGptFiveDen}{77}
\rqSet{rqRZeroTestLibBOneGptFivePct}{9.1}
\rqSet{rqRZeroTestLibBOneGeminiNum}{3}
\rqSet{rqRZeroTestLibBOneGeminiDen}{77}
\rqSet{rqRZeroTestLibBOneGeminiPct}{3.9}
\rqSet{rqRZeroTestLibBTwoNum}{44}
\rqSet{rqRZeroTestLibBTwoDen}{660}
\rqSet{rqRZeroTestLibBTwoPct}{6.7}
\rqSet{rqRZeroTestLibBTwoKiroNum}{2}
\rqSet{rqRZeroTestLibBTwoKiroDen}{66}
\rqSet{rqRZeroTestLibBTwoKiroPct}{3.0}
\rqSet{rqRZeroTestLibBTwoClaudeNum}{2}
\rqSet{rqRZeroTestLibBTwoClaudeDen}{66}
\rqSet{rqRZeroTestLibBTwoClaudePct}{3.0}
\rqSet{rqRZeroTestLibBTwoCodexNum}{2}
\rqSet{rqRZeroTestLibBTwoCodexDen}{66}
\rqSet{rqRZeroTestLibBTwoCodexPct}{3.0}
\rqSet{rqRZeroTestLibBTwoCTwoRustNum}{1}
\rqSet{rqRZeroTestLibBTwoCTwoRustDen}{66}
\rqSet{rqRZeroTestLibBTwoCTwoRustPct}{1.5}
\rqSet{rqRZeroTestLibBTwoLaertesNum}{1}
\rqSet{rqRZeroTestLibBTwoLaertesDen}{66}
\rqSet{rqRZeroTestLibBTwoLaertesPct}{1.5}
\rqSet{rqRZeroTestLibBTwoCTwoSaferRustNum}{8}
\rqSet{rqRZeroTestLibBTwoCTwoSaferRustDen}{66}
\rqSet{rqRZeroTestLibBTwoCTwoSaferRustPct}{12.1}
\rqSet{rqRZeroTestLibBTwoSmartCTwoRustNum}{0}
\rqSet{rqRZeroTestLibBTwoSmartCTwoRustDen}{66}
\rqSet{rqRZeroTestLibBTwoSmartCTwoRustPct}{0.0}
\rqSet{rqRZeroTestLibBTwoKimiNum}{10}
\rqSet{rqRZeroTestLibBTwoKimiDen}{66}
\rqSet{rqRZeroTestLibBTwoKimiPct}{15.2}
\rqSet{rqRZeroTestLibBTwoGptFiveNum}{8}
\rqSet{rqRZeroTestLibBTwoGptFiveDen}{66}
\rqSet{rqRZeroTestLibBTwoGptFivePct}{12.1}
\rqSet{rqRZeroTestLibBTwoGeminiNum}{10}
\rqSet{rqRZeroTestLibBTwoGeminiDen}{66}
\rqSet{rqRZeroTestLibBTwoGeminiPct}{15.2}
\rqSet{rqRZeroTestLibPZeroNum}{0}
\rqSet{rqRZeroTestLibPZeroDen}{0}
\rqSet{rqRZeroTestLibPZeroPct}{n/a}
\rqSet{rqRZeroTestLibPZeroKiroNum}{0}
\rqSet{rqRZeroTestLibPZeroKiroDen}{0}
\rqSet{rqRZeroTestLibPZeroKiroPct}{n/a}
\rqSet{rqRZeroTestLibPZeroClaudeNum}{0}
\rqSet{rqRZeroTestLibPZeroClaudeDen}{0}
\rqSet{rqRZeroTestLibPZeroClaudePct}{n/a}
\rqSet{rqRZeroTestLibPZeroCodexNum}{0}
\rqSet{rqRZeroTestLibPZeroCodexDen}{0}
\rqSet{rqRZeroTestLibPZeroCodexPct}{n/a}
\rqSet{rqRZeroTestLibPZeroCTwoRustNum}{0}
\rqSet{rqRZeroTestLibPZeroCTwoRustDen}{0}
\rqSet{rqRZeroTestLibPZeroCTwoRustPct}{n/a}
\rqSet{rqRZeroTestLibPZeroLaertesNum}{0}
\rqSet{rqRZeroTestLibPZeroLaertesDen}{0}
\rqSet{rqRZeroTestLibPZeroLaertesPct}{n/a}
\rqSet{rqRZeroTestLibPZeroCTwoSaferRustNum}{0}
\rqSet{rqRZeroTestLibPZeroCTwoSaferRustDen}{0}
\rqSet{rqRZeroTestLibPZeroCTwoSaferRustPct}{n/a}
\rqSet{rqRZeroTestLibPZeroSmartCTwoRustNum}{0}
\rqSet{rqRZeroTestLibPZeroSmartCTwoRustDen}{0}
\rqSet{rqRZeroTestLibPZeroSmartCTwoRustPct}{n/a}
\rqSet{rqRZeroTestLibPZeroKimiNum}{0}
\rqSet{rqRZeroTestLibPZeroKimiDen}{0}
\rqSet{rqRZeroTestLibPZeroKimiPct}{n/a}
\rqSet{rqRZeroTestLibPZeroGptFiveNum}{0}
\rqSet{rqRZeroTestLibPZeroGptFiveDen}{0}
\rqSet{rqRZeroTestLibPZeroGptFivePct}{n/a}
\rqSet{rqRZeroTestLibPZeroGeminiNum}{0}
\rqSet{rqRZeroTestLibPZeroGeminiDen}{0}
\rqSet{rqRZeroTestLibPZeroGeminiPct}{n/a}
\rqSet{rqRZeroTestLibPOneNum}{56}
\rqSet{rqRZeroTestLibPOneDen}{760}
\rqSet{rqRZeroTestLibPOnePct}{7.4}
\rqSet{rqRZeroTestLibPOneKiroNum}{0}
\rqSet{rqRZeroTestLibPOneKiroDen}{76}
\rqSet{rqRZeroTestLibPOneKiroPct}{0.0}
\rqSet{rqRZeroTestLibPOneClaudeNum}{11}
\rqSet{rqRZeroTestLibPOneClaudeDen}{76}
\rqSet{rqRZeroTestLibPOneClaudePct}{14.5}
\rqSet{rqRZeroTestLibPOneCodexNum}{45}
\rqSet{rqRZeroTestLibPOneCodexDen}{76}
\rqSet{rqRZeroTestLibPOneCodexPct}{59.2}
\rqSet{rqRZeroTestLibPOneCTwoRustNum}{0}
\rqSet{rqRZeroTestLibPOneCTwoRustDen}{76}
\rqSet{rqRZeroTestLibPOneCTwoRustPct}{0.0}
\rqSet{rqRZeroTestLibPOneLaertesNum}{0}
\rqSet{rqRZeroTestLibPOneLaertesDen}{76}
\rqSet{rqRZeroTestLibPOneLaertesPct}{0.0}
\rqSet{rqRZeroTestLibPOneCTwoSaferRustNum}{0}
\rqSet{rqRZeroTestLibPOneCTwoSaferRustDen}{76}
\rqSet{rqRZeroTestLibPOneCTwoSaferRustPct}{0.0}
\rqSet{rqRZeroTestLibPOneSmartCTwoRustNum}{0}
\rqSet{rqRZeroTestLibPOneSmartCTwoRustDen}{76}
\rqSet{rqRZeroTestLibPOneSmartCTwoRustPct}{0.0}
\rqSet{rqRZeroTestLibPOneKimiNum}{0}
\rqSet{rqRZeroTestLibPOneKimiDen}{76}
\rqSet{rqRZeroTestLibPOneKimiPct}{0.0}
\rqSet{rqRZeroTestLibPOneGptFiveNum}{0}
\rqSet{rqRZeroTestLibPOneGptFiveDen}{76}
\rqSet{rqRZeroTestLibPOneGptFivePct}{0.0}
\rqSet{rqRZeroTestLibPOneGeminiNum}{0}
\rqSet{rqRZeroTestLibPOneGeminiDen}{76}
\rqSet{rqRZeroTestLibPOneGeminiPct}{0.0}
\rqSet{rqPrepLibAllNum}{1264}
\rqSet{rqPrepLibAllDen}{1281}
\rqSet{rqPrepLibAllPct}{98.7}
\rqSet{rqPrepLibAllKiroNum}{216}
\rqSet{rqPrepLibAllKiroDen}{217}
\rqSet{rqPrepLibAllKiroPct}{99.5}
\rqSet{rqPrepLibAllClaudeNum}{205}
\rqSet{rqPrepLibAllClaudeDen}{206}
\rqSet{rqPrepLibAllClaudePct}{99.5}
\rqSet{rqPrepLibAllCodexNum}{167}
\rqSet{rqPrepLibAllCodexDen}{171}
\rqSet{rqPrepLibAllCodexPct}{97.7}
\rqSet{rqPrepLibAllCTwoRustNum}{140}
\rqSet{rqPrepLibAllCTwoRustDen}{141}
\rqSet{rqPrepLibAllCTwoRustPct}{99.3}
\rqSet{rqPrepLibAllLaertesNum}{136}
\rqSet{rqPrepLibAllLaertesDen}{137}
\rqSet{rqPrepLibAllLaertesPct}{99.3}
\rqSet{rqPrepLibAllCTwoSaferRustNum}{103}
\rqSet{rqPrepLibAllCTwoSaferRustDen}{105}
\rqSet{rqPrepLibAllCTwoSaferRustPct}{98.1}
\rqSet{rqPrepLibAllSmartCTwoRustNum}{0}
\rqSet{rqPrepLibAllSmartCTwoRustDen}{0}
\rqSet{rqPrepLibAllSmartCTwoRustPct}{n/a}
\rqSet{rqPrepLibAllKimiNum}{83}
\rqSet{rqPrepLibAllKimiDen}{85}
\rqSet{rqPrepLibAllKimiPct}{97.6}
\rqSet{rqPrepLibAllGptFiveNum}{108}
\rqSet{rqPrepLibAllGptFiveDen}{111}
\rqSet{rqPrepLibAllGptFivePct}{97.3}
\rqSet{rqPrepLibAllGeminiNum}{106}
\rqSet{rqPrepLibAllGeminiDen}{108}
\rqSet{rqPrepLibAllGeminiPct}{98.1}
\rqSet{rqPrepLibBOneNum}{636}
\rqSet{rqPrepLibBOneDen}{637}
\rqSet{rqPrepLibBOnePct}{99.8}
\rqSet{rqPrepLibBOneKiroNum}{77}
\rqSet{rqPrepLibBOneKiroDen}{77}
\rqSet{rqPrepLibBOneKiroPct}{100.0}
\rqSet{rqPrepLibBOneClaudeNum}{77}
\rqSet{rqPrepLibBOneClaudeDen}{77}
\rqSet{rqPrepLibBOneClaudePct}{100.0}
\rqSet{rqPrepLibBOneCodexNum}{77}
\rqSet{rqPrepLibBOneCodexDen}{77}
\rqSet{rqPrepLibBOneCodexPct}{100.0}
\rqSet{rqPrepLibBOneCTwoRustNum}{77}
\rqSet{rqPrepLibBOneCTwoRustDen}{77}
\rqSet{rqPrepLibBOneCTwoRustPct}{100.0}
\rqSet{rqPrepLibBOneLaertesNum}{74}
\rqSet{rqPrepLibBOneLaertesDen}{74}
\rqSet{rqPrepLibBOneLaertesPct}{100.0}
\rqSet{rqPrepLibBOneCTwoSaferRustNum}{59}
\rqSet{rqPrepLibBOneCTwoSaferRustDen}{60}
\rqSet{rqPrepLibBOneCTwoSaferRustPct}{98.3}
\rqSet{rqPrepLibBOneSmartCTwoRustNum}{0}
\rqSet{rqPrepLibBOneSmartCTwoRustDen}{0}
\rqSet{rqPrepLibBOneSmartCTwoRustPct}{n/a}
\rqSet{rqPrepLibBOneKimiNum}{57}
\rqSet{rqPrepLibBOneKimiDen}{57}
\rqSet{rqPrepLibBOneKimiPct}{100.0}
\rqSet{rqPrepLibBOneGptFiveNum}{67}
\rqSet{rqPrepLibBOneGptFiveDen}{67}
\rqSet{rqPrepLibBOneGptFivePct}{100.0}
\rqSet{rqPrepLibBOneGeminiNum}{71}
\rqSet{rqPrepLibBOneGeminiDen}{71}
\rqSet{rqPrepLibBOneGeminiPct}{100.0}
\rqSet{rqPrepLibBTwoNum}{459}
\rqSet{rqPrepLibBTwoDen}{472}
\rqSet{rqPrepLibBTwoPct}{97.2}
\rqSet{rqPrepLibBTwoKiroNum}{63}
\rqSet{rqPrepLibBTwoKiroDen}{64}
\rqSet{rqPrepLibBTwoKiroPct}{98.4}
\rqSet{rqPrepLibBTwoClaudeNum}{63}
\rqSet{rqPrepLibBTwoClaudeDen}{64}
\rqSet{rqPrepLibBTwoClaudePct}{98.4}
\rqSet{rqPrepLibBTwoCodexNum}{62}
\rqSet{rqPrepLibBTwoCodexDen}{63}
\rqSet{rqPrepLibBTwoCodexPct}{98.4}
\rqSet{rqPrepLibBTwoCTwoRustNum}{63}
\rqSet{rqPrepLibBTwoCTwoRustDen}{64}
\rqSet{rqPrepLibBTwoCTwoRustPct}{98.4}
\rqSet{rqPrepLibBTwoLaertesNum}{62}
\rqSet{rqPrepLibBTwoLaertesDen}{63}
\rqSet{rqPrepLibBTwoLaertesPct}{98.4}
\rqSet{rqPrepLibBTwoCTwoSaferRustNum}{44}
\rqSet{rqPrepLibBTwoCTwoSaferRustDen}{45}
\rqSet{rqPrepLibBTwoCTwoSaferRustPct}{97.8}
\rqSet{rqPrepLibBTwoSmartCTwoRustNum}{0}
\rqSet{rqPrepLibBTwoSmartCTwoRustDen}{0}
\rqSet{rqPrepLibBTwoSmartCTwoRustPct}{n/a}
\rqSet{rqPrepLibBTwoKimiNum}{26}
\rqSet{rqPrepLibBTwoKimiDen}{28}
\rqSet{rqPrepLibBTwoKimiPct}{92.9}
\rqSet{rqPrepLibBTwoGptFiveNum}{41}
\rqSet{rqPrepLibBTwoGptFiveDen}{44}
\rqSet{rqPrepLibBTwoGptFivePct}{93.2}
\rqSet{rqPrepLibBTwoGeminiNum}{35}
\rqSet{rqPrepLibBTwoGeminiDen}{37}
\rqSet{rqPrepLibBTwoGeminiPct}{94.6}
\rqSet{rqPrepLibPZeroNum}{0}
\rqSet{rqPrepLibPZeroDen}{0}
\rqSet{rqPrepLibPZeroPct}{n/a}
\rqSet{rqPrepLibPZeroKiroNum}{0}
\rqSet{rqPrepLibPZeroKiroDen}{0}
\rqSet{rqPrepLibPZeroKiroPct}{n/a}
\rqSet{rqPrepLibPZeroClaudeNum}{0}
\rqSet{rqPrepLibPZeroClaudeDen}{0}
\rqSet{rqPrepLibPZeroClaudePct}{n/a}
\rqSet{rqPrepLibPZeroCodexNum}{0}
\rqSet{rqPrepLibPZeroCodexDen}{0}
\rqSet{rqPrepLibPZeroCodexPct}{n/a}
\rqSet{rqPrepLibPZeroCTwoRustNum}{0}
\rqSet{rqPrepLibPZeroCTwoRustDen}{0}
\rqSet{rqPrepLibPZeroCTwoRustPct}{n/a}
\rqSet{rqPrepLibPZeroLaertesNum}{0}
\rqSet{rqPrepLibPZeroLaertesDen}{0}
\rqSet{rqPrepLibPZeroLaertesPct}{n/a}
\rqSet{rqPrepLibPZeroCTwoSaferRustNum}{0}
\rqSet{rqPrepLibPZeroCTwoSaferRustDen}{0}
\rqSet{rqPrepLibPZeroCTwoSaferRustPct}{n/a}
\rqSet{rqPrepLibPZeroSmartCTwoRustNum}{0}
\rqSet{rqPrepLibPZeroSmartCTwoRustDen}{0}
\rqSet{rqPrepLibPZeroSmartCTwoRustPct}{n/a}
\rqSet{rqPrepLibPZeroKimiNum}{0}
\rqSet{rqPrepLibPZeroKimiDen}{0}
\rqSet{rqPrepLibPZeroKimiPct}{n/a}
\rqSet{rqPrepLibPZeroGptFiveNum}{0}
\rqSet{rqPrepLibPZeroGptFiveDen}{0}
\rqSet{rqPrepLibPZeroGptFivePct}{n/a}
\rqSet{rqPrepLibPZeroGeminiNum}{0}
\rqSet{rqPrepLibPZeroGeminiDen}{0}
\rqSet{rqPrepLibPZeroGeminiPct}{n/a}
\rqSet{rqPrepLibPOneNum}{169}
\rqSet{rqPrepLibPOneDen}{172}
\rqSet{rqPrepLibPOnePct}{98.3}
\rqSet{rqPrepLibPOneKiroNum}{76}
\rqSet{rqPrepLibPOneKiroDen}{76}
\rqSet{rqPrepLibPOneKiroPct}{100.0}
\rqSet{rqPrepLibPOneClaudeNum}{65}
\rqSet{rqPrepLibPOneClaudeDen}{65}
\rqSet{rqPrepLibPOneClaudePct}{100.0}
\rqSet{rqPrepLibPOneCodexNum}{28}
\rqSet{rqPrepLibPOneCodexDen}{31}
\rqSet{rqPrepLibPOneCodexPct}{90.3}
\rqSet{rqPrepLibPOneCTwoRustNum}{0}
\rqSet{rqPrepLibPOneCTwoRustDen}{0}
\rqSet{rqPrepLibPOneCTwoRustPct}{n/a}
\rqSet{rqPrepLibPOneLaertesNum}{0}
\rqSet{rqPrepLibPOneLaertesDen}{0}
\rqSet{rqPrepLibPOneLaertesPct}{n/a}
\rqSet{rqPrepLibPOneCTwoSaferRustNum}{0}
\rqSet{rqPrepLibPOneCTwoSaferRustDen}{0}
\rqSet{rqPrepLibPOneCTwoSaferRustPct}{n/a}
\rqSet{rqPrepLibPOneSmartCTwoRustNum}{0}
\rqSet{rqPrepLibPOneSmartCTwoRustDen}{0}
\rqSet{rqPrepLibPOneSmartCTwoRustPct}{n/a}
\rqSet{rqPrepLibPOneKimiNum}{0}
\rqSet{rqPrepLibPOneKimiDen}{0}
\rqSet{rqPrepLibPOneKimiPct}{n/a}
\rqSet{rqPrepLibPOneGptFiveNum}{0}
\rqSet{rqPrepLibPOneGptFiveDen}{0}
\rqSet{rqPrepLibPOneGptFivePct}{n/a}
\rqSet{rqPrepLibPOneGeminiNum}{0}
\rqSet{rqPrepLibPOneGeminiDen}{0}
\rqSet{rqPrepLibPOneGeminiPct}{n/a}
\rqSet{rqPtestLibAllNum}{1088}
\rqSet{rqPtestLibAllDen}{1281}
\rqSet{rqPtestLibAllPct}{84.9}
\rqSet{rqPtestLibAllKiroNum}{192}
\rqSet{rqPtestLibAllKiroDen}{217}
\rqSet{rqPtestLibAllKiroPct}{88.5}
\rqSet{rqPtestLibAllClaudeNum}{181}
\rqSet{rqPtestLibAllClaudeDen}{206}
\rqSet{rqPtestLibAllClaudePct}{87.9}
\rqSet{rqPtestLibAllCodexNum}{143}
\rqSet{rqPtestLibAllCodexDen}{171}
\rqSet{rqPtestLibAllCodexPct}{83.6}
\rqSet{rqPtestLibAllCTwoRustNum}{116}
\rqSet{rqPtestLibAllCTwoRustDen}{141}
\rqSet{rqPtestLibAllCTwoRustPct}{82.3}
\rqSet{rqPtestLibAllLaertesNum}{114}
\rqSet{rqPtestLibAllLaertesDen}{137}
\rqSet{rqPtestLibAllLaertesPct}{83.2}
\rqSet{rqPtestLibAllCTwoSaferRustNum}{87}
\rqSet{rqPtestLibAllCTwoSaferRustDen}{105}
\rqSet{rqPtestLibAllCTwoSaferRustPct}{82.9}
\rqSet{rqPtestLibAllSmartCTwoRustNum}{0}
\rqSet{rqPtestLibAllSmartCTwoRustDen}{0}
\rqSet{rqPtestLibAllSmartCTwoRustPct}{n/a}
\rqSet{rqPtestLibAllKimiNum}{73}
\rqSet{rqPtestLibAllKimiDen}{85}
\rqSet{rqPtestLibAllKimiPct}{85.9}
\rqSet{rqPtestLibAllGptFiveNum}{92}
\rqSet{rqPtestLibAllGptFiveDen}{111}
\rqSet{rqPtestLibAllGptFivePct}{82.9}
\rqSet{rqPtestLibAllGeminiNum}{90}
\rqSet{rqPtestLibAllGeminiDen}{108}
\rqSet{rqPtestLibAllGeminiPct}{83.3}
\rqSet{rqPtestLibBOneNum}{540}
\rqSet{rqPtestLibBOneDen}{637}
\rqSet{rqPtestLibBOnePct}{84.8}
\rqSet{rqPtestLibBOneKiroNum}{65}
\rqSet{rqPtestLibBOneKiroDen}{77}
\rqSet{rqPtestLibBOneKiroPct}{84.4}
\rqSet{rqPtestLibBOneClaudeNum}{65}
\rqSet{rqPtestLibBOneClaudeDen}{77}
\rqSet{rqPtestLibBOneClaudePct}{84.4}
\rqSet{rqPtestLibBOneCodexNum}{65}
\rqSet{rqPtestLibBOneCodexDen}{77}
\rqSet{rqPtestLibBOneCodexPct}{84.4}
\rqSet{rqPtestLibBOneCTwoRustNum}{65}
\rqSet{rqPtestLibBOneCTwoRustDen}{77}
\rqSet{rqPtestLibBOneCTwoRustPct}{84.4}
\rqSet{rqPtestLibBOneLaertesNum}{63}
\rqSet{rqPtestLibBOneLaertesDen}{74}
\rqSet{rqPtestLibBOneLaertesPct}{85.1}
\rqSet{rqPtestLibBOneCTwoSaferRustNum}{50}
\rqSet{rqPtestLibBOneCTwoSaferRustDen}{60}
\rqSet{rqPtestLibBOneCTwoSaferRustPct}{83.3}
\rqSet{rqPtestLibBOneSmartCTwoRustNum}{0}
\rqSet{rqPtestLibBOneSmartCTwoRustDen}{0}
\rqSet{rqPtestLibBOneSmartCTwoRustPct}{n/a}
\rqSet{rqPtestLibBOneKimiNum}{49}
\rqSet{rqPtestLibBOneKimiDen}{57}
\rqSet{rqPtestLibBOneKimiPct}{86.0}
\rqSet{rqPtestLibBOneGptFiveNum}{57}
\rqSet{rqPtestLibBOneGptFiveDen}{67}
\rqSet{rqPtestLibBOneGptFivePct}{85.1}
\rqSet{rqPtestLibBOneGeminiNum}{61}
\rqSet{rqPtestLibBOneGeminiDen}{71}
\rqSet{rqPtestLibBOneGeminiPct}{85.9}
\rqSet{rqPtestLibBTwoNum}{379}
\rqSet{rqPtestLibBTwoDen}{472}
\rqSet{rqPtestLibBTwoPct}{80.3}
\rqSet{rqPtestLibBTwoKiroNum}{51}
\rqSet{rqPtestLibBTwoKiroDen}{64}
\rqSet{rqPtestLibBTwoKiroPct}{79.7}
\rqSet{rqPtestLibBTwoClaudeNum}{51}
\rqSet{rqPtestLibBTwoClaudeDen}{64}
\rqSet{rqPtestLibBTwoClaudePct}{79.7}
\rqSet{rqPtestLibBTwoCodexNum}{50}
\rqSet{rqPtestLibBTwoCodexDen}{63}
\rqSet{rqPtestLibBTwoCodexPct}{79.4}
\rqSet{rqPtestLibBTwoCTwoRustNum}{51}
\rqSet{rqPtestLibBTwoCTwoRustDen}{64}
\rqSet{rqPtestLibBTwoCTwoRustPct}{79.7}
\rqSet{rqPtestLibBTwoLaertesNum}{51}
\rqSet{rqPtestLibBTwoLaertesDen}{63}
\rqSet{rqPtestLibBTwoLaertesPct}{81.0}
\rqSet{rqPtestLibBTwoCTwoSaferRustNum}{37}
\rqSet{rqPtestLibBTwoCTwoSaferRustDen}{45}
\rqSet{rqPtestLibBTwoCTwoSaferRustPct}{82.2}
\rqSet{rqPtestLibBTwoSmartCTwoRustNum}{0}
\rqSet{rqPtestLibBTwoSmartCTwoRustDen}{0}
\rqSet{rqPtestLibBTwoSmartCTwoRustPct}{n/a}
\rqSet{rqPtestLibBTwoKimiNum}{24}
\rqSet{rqPtestLibBTwoKimiDen}{28}
\rqSet{rqPtestLibBTwoKimiPct}{85.7}
\rqSet{rqPtestLibBTwoGptFiveNum}{35}
\rqSet{rqPtestLibBTwoGptFiveDen}{44}
\rqSet{rqPtestLibBTwoGptFivePct}{79.5}
\rqSet{rqPtestLibBTwoGeminiNum}{29}
\rqSet{rqPtestLibBTwoGeminiDen}{37}
\rqSet{rqPtestLibBTwoGeminiPct}{78.4}
\rqSet{rqPtestLibPZeroNum}{0}
\rqSet{rqPtestLibPZeroDen}{0}
\rqSet{rqPtestLibPZeroPct}{n/a}
\rqSet{rqPtestLibPZeroKiroNum}{0}
\rqSet{rqPtestLibPZeroKiroDen}{0}
\rqSet{rqPtestLibPZeroKiroPct}{n/a}
\rqSet{rqPtestLibPZeroClaudeNum}{0}
\rqSet{rqPtestLibPZeroClaudeDen}{0}
\rqSet{rqPtestLibPZeroClaudePct}{n/a}
\rqSet{rqPtestLibPZeroCodexNum}{0}
\rqSet{rqPtestLibPZeroCodexDen}{0}
\rqSet{rqPtestLibPZeroCodexPct}{n/a}
\rqSet{rqPtestLibPZeroCTwoRustNum}{0}
\rqSet{rqPtestLibPZeroCTwoRustDen}{0}
\rqSet{rqPtestLibPZeroCTwoRustPct}{n/a}
\rqSet{rqPtestLibPZeroLaertesNum}{0}
\rqSet{rqPtestLibPZeroLaertesDen}{0}
\rqSet{rqPtestLibPZeroLaertesPct}{n/a}
\rqSet{rqPtestLibPZeroCTwoSaferRustNum}{0}
\rqSet{rqPtestLibPZeroCTwoSaferRustDen}{0}
\rqSet{rqPtestLibPZeroCTwoSaferRustPct}{n/a}
\rqSet{rqPtestLibPZeroSmartCTwoRustNum}{0}
\rqSet{rqPtestLibPZeroSmartCTwoRustDen}{0}
\rqSet{rqPtestLibPZeroSmartCTwoRustPct}{n/a}
\rqSet{rqPtestLibPZeroKimiNum}{0}
\rqSet{rqPtestLibPZeroKimiDen}{0}
\rqSet{rqPtestLibPZeroKimiPct}{n/a}
\rqSet{rqPtestLibPZeroGptFiveNum}{0}
\rqSet{rqPtestLibPZeroGptFiveDen}{0}
\rqSet{rqPtestLibPZeroGptFivePct}{n/a}
\rqSet{rqPtestLibPZeroGeminiNum}{0}
\rqSet{rqPtestLibPZeroGeminiDen}{0}
\rqSet{rqPtestLibPZeroGeminiPct}{n/a}
\rqSet{rqPtestLibPOneNum}{169}
\rqSet{rqPtestLibPOneDen}{172}
\rqSet{rqPtestLibPOnePct}{98.3}
\rqSet{rqPtestLibPOneKiroNum}{76}
\rqSet{rqPtestLibPOneKiroDen}{76}
\rqSet{rqPtestLibPOneKiroPct}{100.0}
\rqSet{rqPtestLibPOneClaudeNum}{65}
\rqSet{rqPtestLibPOneClaudeDen}{65}
\rqSet{rqPtestLibPOneClaudePct}{100.0}
\rqSet{rqPtestLibPOneCodexNum}{28}
\rqSet{rqPtestLibPOneCodexDen}{31}
\rqSet{rqPtestLibPOneCodexPct}{90.3}
\rqSet{rqPtestLibPOneCTwoRustNum}{0}
\rqSet{rqPtestLibPOneCTwoRustDen}{0}
\rqSet{rqPtestLibPOneCTwoRustPct}{n/a}
\rqSet{rqPtestLibPOneLaertesNum}{0}
\rqSet{rqPtestLibPOneLaertesDen}{0}
\rqSet{rqPtestLibPOneLaertesPct}{n/a}
\rqSet{rqPtestLibPOneCTwoSaferRustNum}{0}
\rqSet{rqPtestLibPOneCTwoSaferRustDen}{0}
\rqSet{rqPtestLibPOneCTwoSaferRustPct}{n/a}
\rqSet{rqPtestLibPOneSmartCTwoRustNum}{0}
\rqSet{rqPtestLibPOneSmartCTwoRustDen}{0}
\rqSet{rqPtestLibPOneSmartCTwoRustPct}{n/a}
\rqSet{rqPtestLibPOneKimiNum}{0}
\rqSet{rqPtestLibPOneKimiDen}{0}
\rqSet{rqPtestLibPOneKimiPct}{n/a}
\rqSet{rqPtestLibPOneGptFiveNum}{0}
\rqSet{rqPtestLibPOneGptFiveDen}{0}
\rqSet{rqPtestLibPOneGptFivePct}{n/a}
\rqSet{rqPtestLibPOneGeminiNum}{0}
\rqSet{rqPtestLibPOneGeminiDen}{0}
\rqSet{rqPtestLibPOneGeminiPct}{n/a}
\rqSet{rqLpLibAllNum}{193}
\rqSet{rqLpLibAllDen}{1281}
\rqSet{rqLpLibAllPct}{15.1}
\rqSet{rqLpLibAllKiroNum}{25}
\rqSet{rqLpLibAllKiroDen}{217}
\rqSet{rqLpLibAllKiroPct}{11.5}
\rqSet{rqLpLibAllClaudeNum}{25}
\rqSet{rqLpLibAllClaudeDen}{206}
\rqSet{rqLpLibAllClaudePct}{12.1}
\rqSet{rqLpLibAllCodexNum}{28}
\rqSet{rqLpLibAllCodexDen}{171}
\rqSet{rqLpLibAllCodexPct}{16.4}
\rqSet{rqLpLibAllCTwoRustNum}{25}
\rqSet{rqLpLibAllCTwoRustDen}{141}
\rqSet{rqLpLibAllCTwoRustPct}{17.7}
\rqSet{rqLpLibAllLaertesNum}{23}
\rqSet{rqLpLibAllLaertesDen}{137}
\rqSet{rqLpLibAllLaertesPct}{16.8}
\rqSet{rqLpLibAllCTwoSaferRustNum}{18}
\rqSet{rqLpLibAllCTwoSaferRustDen}{105}
\rqSet{rqLpLibAllCTwoSaferRustPct}{17.1}
\rqSet{rqLpLibAllSmartCTwoRustNum}{0}
\rqSet{rqLpLibAllSmartCTwoRustDen}{0}
\rqSet{rqLpLibAllSmartCTwoRustPct}{n/a}
\rqSet{rqLpLibAllKimiNum}{12}
\rqSet{rqLpLibAllKimiDen}{85}
\rqSet{rqLpLibAllKimiPct}{14.1}
\rqSet{rqLpLibAllGptFiveNum}{19}
\rqSet{rqLpLibAllGptFiveDen}{111}
\rqSet{rqLpLibAllGptFivePct}{17.1}
\rqSet{rqLpLibAllGeminiNum}{18}
\rqSet{rqLpLibAllGeminiDen}{108}
\rqSet{rqLpLibAllGeminiPct}{16.7}
\rqSet{rqLpLibBOneNum}{97}
\rqSet{rqLpLibBOneDen}{637}
\rqSet{rqLpLibBOnePct}{15.2}
\rqSet{rqLpLibBOneKiroNum}{12}
\rqSet{rqLpLibBOneKiroDen}{77}
\rqSet{rqLpLibBOneKiroPct}{15.6}
\rqSet{rqLpLibBOneClaudeNum}{12}
\rqSet{rqLpLibBOneClaudeDen}{77}
\rqSet{rqLpLibBOneClaudePct}{15.6}
\rqSet{rqLpLibBOneCodexNum}{12}
\rqSet{rqLpLibBOneCodexDen}{77}
\rqSet{rqLpLibBOneCodexPct}{15.6}
\rqSet{rqLpLibBOneCTwoRustNum}{12}
\rqSet{rqLpLibBOneCTwoRustDen}{77}
\rqSet{rqLpLibBOneCTwoRustPct}{15.6}
\rqSet{rqLpLibBOneLaertesNum}{11}
\rqSet{rqLpLibBOneLaertesDen}{74}
\rqSet{rqLpLibBOneLaertesPct}{14.9}
\rqSet{rqLpLibBOneCTwoSaferRustNum}{10}
\rqSet{rqLpLibBOneCTwoSaferRustDen}{60}
\rqSet{rqLpLibBOneCTwoSaferRustPct}{16.7}
\rqSet{rqLpLibBOneSmartCTwoRustNum}{0}
\rqSet{rqLpLibBOneSmartCTwoRustDen}{0}
\rqSet{rqLpLibBOneSmartCTwoRustPct}{n/a}
\rqSet{rqLpLibBOneKimiNum}{8}
\rqSet{rqLpLibBOneKimiDen}{57}
\rqSet{rqLpLibBOneKimiPct}{14.0}
\rqSet{rqLpLibBOneGptFiveNum}{10}
\rqSet{rqLpLibBOneGptFiveDen}{67}
\rqSet{rqLpLibBOneGptFivePct}{14.9}
\rqSet{rqLpLibBOneGeminiNum}{10}
\rqSet{rqLpLibBOneGeminiDen}{71}
\rqSet{rqLpLibBOneGeminiPct}{14.1}
\rqSet{rqLpLibBTwoNum}{93}
\rqSet{rqLpLibBTwoDen}{472}
\rqSet{rqLpLibBTwoPct}{19.7}
\rqSet{rqLpLibBTwoKiroNum}{13}
\rqSet{rqLpLibBTwoKiroDen}{64}
\rqSet{rqLpLibBTwoKiroPct}{20.3}
\rqSet{rqLpLibBTwoClaudeNum}{13}
\rqSet{rqLpLibBTwoClaudeDen}{64}
\rqSet{rqLpLibBTwoClaudePct}{20.3}
\rqSet{rqLpLibBTwoCodexNum}{13}
\rqSet{rqLpLibBTwoCodexDen}{63}
\rqSet{rqLpLibBTwoCodexPct}{20.6}
\rqSet{rqLpLibBTwoCTwoRustNum}{13}
\rqSet{rqLpLibBTwoCTwoRustDen}{64}
\rqSet{rqLpLibBTwoCTwoRustPct}{20.3}
\rqSet{rqLpLibBTwoLaertesNum}{12}
\rqSet{rqLpLibBTwoLaertesDen}{63}
\rqSet{rqLpLibBTwoLaertesPct}{19.0}
\rqSet{rqLpLibBTwoCTwoSaferRustNum}{8}
\rqSet{rqLpLibBTwoCTwoSaferRustDen}{45}
\rqSet{rqLpLibBTwoCTwoSaferRustPct}{17.8}
\rqSet{rqLpLibBTwoSmartCTwoRustNum}{0}
\rqSet{rqLpLibBTwoSmartCTwoRustDen}{0}
\rqSet{rqLpLibBTwoSmartCTwoRustPct}{n/a}
\rqSet{rqLpLibBTwoKimiNum}{4}
\rqSet{rqLpLibBTwoKimiDen}{28}
\rqSet{rqLpLibBTwoKimiPct}{14.3}
\rqSet{rqLpLibBTwoGptFiveNum}{9}
\rqSet{rqLpLibBTwoGptFiveDen}{44}
\rqSet{rqLpLibBTwoGptFivePct}{20.5}
\rqSet{rqLpLibBTwoGeminiNum}{8}
\rqSet{rqLpLibBTwoGeminiDen}{37}
\rqSet{rqLpLibBTwoGeminiPct}{21.6}
\rqSet{rqLpLibPZeroNum}{0}
\rqSet{rqLpLibPZeroDen}{0}
\rqSet{rqLpLibPZeroPct}{n/a}
\rqSet{rqLpLibPZeroKiroNum}{0}
\rqSet{rqLpLibPZeroKiroDen}{0}
\rqSet{rqLpLibPZeroKiroPct}{n/a}
\rqSet{rqLpLibPZeroClaudeNum}{0}
\rqSet{rqLpLibPZeroClaudeDen}{0}
\rqSet{rqLpLibPZeroClaudePct}{n/a}
\rqSet{rqLpLibPZeroCodexNum}{0}
\rqSet{rqLpLibPZeroCodexDen}{0}
\rqSet{rqLpLibPZeroCodexPct}{n/a}
\rqSet{rqLpLibPZeroCTwoRustNum}{0}
\rqSet{rqLpLibPZeroCTwoRustDen}{0}
\rqSet{rqLpLibPZeroCTwoRustPct}{n/a}
\rqSet{rqLpLibPZeroLaertesNum}{0}
\rqSet{rqLpLibPZeroLaertesDen}{0}
\rqSet{rqLpLibPZeroLaertesPct}{n/a}
\rqSet{rqLpLibPZeroCTwoSaferRustNum}{0}
\rqSet{rqLpLibPZeroCTwoSaferRustDen}{0}
\rqSet{rqLpLibPZeroCTwoSaferRustPct}{n/a}
\rqSet{rqLpLibPZeroSmartCTwoRustNum}{0}
\rqSet{rqLpLibPZeroSmartCTwoRustDen}{0}
\rqSet{rqLpLibPZeroSmartCTwoRustPct}{n/a}
\rqSet{rqLpLibPZeroKimiNum}{0}
\rqSet{rqLpLibPZeroKimiDen}{0}
\rqSet{rqLpLibPZeroKimiPct}{n/a}
\rqSet{rqLpLibPZeroGptFiveNum}{0}
\rqSet{rqLpLibPZeroGptFiveDen}{0}
\rqSet{rqLpLibPZeroGptFivePct}{n/a}
\rqSet{rqLpLibPZeroGeminiNum}{0}
\rqSet{rqLpLibPZeroGeminiDen}{0}
\rqSet{rqLpLibPZeroGeminiPct}{n/a}
\rqSet{rqLpLibPOneNum}{3}
\rqSet{rqLpLibPOneDen}{172}
\rqSet{rqLpLibPOnePct}{1.7}
\rqSet{rqLpLibPOneKiroNum}{0}
\rqSet{rqLpLibPOneKiroDen}{76}
\rqSet{rqLpLibPOneKiroPct}{0.0}
\rqSet{rqLpLibPOneClaudeNum}{0}
\rqSet{rqLpLibPOneClaudeDen}{65}
\rqSet{rqLpLibPOneClaudePct}{0.0}
\rqSet{rqLpLibPOneCodexNum}{3}
\rqSet{rqLpLibPOneCodexDen}{31}
\rqSet{rqLpLibPOneCodexPct}{9.7}
\rqSet{rqLpLibPOneCTwoRustNum}{0}
\rqSet{rqLpLibPOneCTwoRustDen}{0}
\rqSet{rqLpLibPOneCTwoRustPct}{n/a}
\rqSet{rqLpLibPOneLaertesNum}{0}
\rqSet{rqLpLibPOneLaertesDen}{0}
\rqSet{rqLpLibPOneLaertesPct}{n/a}
\rqSet{rqLpLibPOneCTwoSaferRustNum}{0}
\rqSet{rqLpLibPOneCTwoSaferRustDen}{0}
\rqSet{rqLpLibPOneCTwoSaferRustPct}{n/a}
\rqSet{rqLpLibPOneSmartCTwoRustNum}{0}
\rqSet{rqLpLibPOneSmartCTwoRustDen}{0}
\rqSet{rqLpLibPOneSmartCTwoRustPct}{n/a}
\rqSet{rqLpLibPOneKimiNum}{0}
\rqSet{rqLpLibPOneKimiDen}{0}
\rqSet{rqLpLibPOneKimiPct}{n/a}
\rqSet{rqLpLibPOneGptFiveNum}{0}
\rqSet{rqLpLibPOneGptFiveDen}{0}
\rqSet{rqLpLibPOneGptFivePct}{n/a}
\rqSet{rqLpLibPOneGeminiNum}{0}
\rqSet{rqLpLibPOneGeminiDen}{0}
\rqSet{rqLpLibPOneGeminiPct}{n/a}
\rqSet{rqPreRepBuildLibAllNum}{8}
\rqSet{rqPreRepBuildLibAllDen}{1281}
\rqSet{rqPreRepBuildLibAllPct}{0.6}
\rqSet{rqPreRepBuildLibAllKiroNum}{1}
\rqSet{rqPreRepBuildLibAllKiroDen}{217}
\rqSet{rqPreRepBuildLibAllKiroPct}{0.5}
\rqSet{rqPreRepBuildLibAllClaudeNum}{1}
\rqSet{rqPreRepBuildLibAllClaudeDen}{206}
\rqSet{rqPreRepBuildLibAllClaudePct}{0.5}
\rqSet{rqPreRepBuildLibAllCodexNum}{1}
\rqSet{rqPreRepBuildLibAllCodexDen}{171}
\rqSet{rqPreRepBuildLibAllCodexPct}{0.6}
\rqSet{rqPreRepBuildLibAllCTwoRustNum}{1}
\rqSet{rqPreRepBuildLibAllCTwoRustDen}{141}
\rqSet{rqPreRepBuildLibAllCTwoRustPct}{0.7}
\rqSet{rqPreRepBuildLibAllLaertesNum}{1}
\rqSet{rqPreRepBuildLibAllLaertesDen}{137}
\rqSet{rqPreRepBuildLibAllLaertesPct}{0.7}
\rqSet{rqPreRepBuildLibAllCTwoSaferRustNum}{1}
\rqSet{rqPreRepBuildLibAllCTwoSaferRustDen}{105}
\rqSet{rqPreRepBuildLibAllCTwoSaferRustPct}{1.0}
\rqSet{rqPreRepBuildLibAllSmartCTwoRustNum}{0}
\rqSet{rqPreRepBuildLibAllSmartCTwoRustDen}{0}
\rqSet{rqPreRepBuildLibAllSmartCTwoRustPct}{n/a}
\rqSet{rqPreRepBuildLibAllKimiNum}{0}
\rqSet{rqPreRepBuildLibAllKimiDen}{85}
\rqSet{rqPreRepBuildLibAllKimiPct}{0.0}
\rqSet{rqPreRepBuildLibAllGptFiveNum}{1}
\rqSet{rqPreRepBuildLibAllGptFiveDen}{111}
\rqSet{rqPreRepBuildLibAllGptFivePct}{0.9}
\rqSet{rqPreRepBuildLibAllGeminiNum}{1}
\rqSet{rqPreRepBuildLibAllGeminiDen}{108}
\rqSet{rqPreRepBuildLibAllGeminiPct}{0.9}
\rqSet{rqPreRepBuildLibBOneNum}{0}
\rqSet{rqPreRepBuildLibBOneDen}{637}
\rqSet{rqPreRepBuildLibBOnePct}{0.0}
\rqSet{rqPreRepBuildLibBTwoNum}{8}
\rqSet{rqPreRepBuildLibBTwoDen}{472}
\rqSet{rqPreRepBuildLibBTwoPct}{1.7}
\rqSet{rqPreRepBuildLibPZeroNum}{0}
\rqSet{rqPreRepBuildLibPZeroDen}{0}
\rqSet{rqPreRepBuildLibPZeroPct}{n/a}
\rqSet{rqPreRepBuildLibPOneNum}{0}
\rqSet{rqPreRepBuildLibPOneDen}{172}
\rqSet{rqPreRepBuildLibPOnePct}{0.0}
\rqSet{rqPreRepTimeoutLibAllNum}{3}
\rqSet{rqPreRepTimeoutLibAllDen}{1281}
\rqSet{rqPreRepTimeoutLibAllPct}{0.2}
\rqSet{rqPreRepTimeoutLibAllKiroNum}{0}
\rqSet{rqPreRepTimeoutLibAllKiroDen}{217}
\rqSet{rqPreRepTimeoutLibAllKiroPct}{0.0}
\rqSet{rqPreRepTimeoutLibAllClaudeNum}{0}
\rqSet{rqPreRepTimeoutLibAllClaudeDen}{206}
\rqSet{rqPreRepTimeoutLibAllClaudePct}{0.0}
\rqSet{rqPreRepTimeoutLibAllCodexNum}{0}
\rqSet{rqPreRepTimeoutLibAllCodexDen}{171}
\rqSet{rqPreRepTimeoutLibAllCodexPct}{0.0}
\rqSet{rqPreRepTimeoutLibAllCTwoRustNum}{0}
\rqSet{rqPreRepTimeoutLibAllCTwoRustDen}{141}
\rqSet{rqPreRepTimeoutLibAllCTwoRustPct}{0.0}
\rqSet{rqPreRepTimeoutLibAllLaertesNum}{0}
\rqSet{rqPreRepTimeoutLibAllLaertesDen}{137}
\rqSet{rqPreRepTimeoutLibAllLaertesPct}{0.0}
\rqSet{rqPreRepTimeoutLibAllCTwoSaferRustNum}{0}
\rqSet{rqPreRepTimeoutLibAllCTwoSaferRustDen}{105}
\rqSet{rqPreRepTimeoutLibAllCTwoSaferRustPct}{0.0}
\rqSet{rqPreRepTimeoutLibAllSmartCTwoRustNum}{0}
\rqSet{rqPreRepTimeoutLibAllSmartCTwoRustDen}{0}
\rqSet{rqPreRepTimeoutLibAllSmartCTwoRustPct}{n/a}
\rqSet{rqPreRepTimeoutLibAllKimiNum}{2}
\rqSet{rqPreRepTimeoutLibAllKimiDen}{85}
\rqSet{rqPreRepTimeoutLibAllKimiPct}{2.4}
\rqSet{rqPreRepTimeoutLibAllGptFiveNum}{1}
\rqSet{rqPreRepTimeoutLibAllGptFiveDen}{111}
\rqSet{rqPreRepTimeoutLibAllGptFivePct}{0.9}
\rqSet{rqPreRepTimeoutLibAllGeminiNum}{0}
\rqSet{rqPreRepTimeoutLibAllGeminiDen}{108}
\rqSet{rqPreRepTimeoutLibAllGeminiPct}{0.0}
\rqSet{rqPreRepTimeoutLibBOneNum}{0}
\rqSet{rqPreRepTimeoutLibBOneDen}{637}
\rqSet{rqPreRepTimeoutLibBOnePct}{0.0}
\rqSet{rqPreRepTimeoutLibBTwoNum}{3}
\rqSet{rqPreRepTimeoutLibBTwoDen}{472}
\rqSet{rqPreRepTimeoutLibBTwoPct}{0.6}
\rqSet{rqPreRepTimeoutLibPZeroNum}{0}
\rqSet{rqPreRepTimeoutLibPZeroDen}{0}
\rqSet{rqPreRepTimeoutLibPZeroPct}{n/a}
\rqSet{rqPreRepTimeoutLibPOneNum}{0}
\rqSet{rqPreRepTimeoutLibPOneDen}{172}
\rqSet{rqPreRepTimeoutLibPOnePct}{0.0}
\rqSet{rqPreRepUnstableLibAllNum}{0}
\rqSet{rqPreRepUnstableLibAllDen}{1281}
\rqSet{rqPreRepUnstableLibAllPct}{0.0}
\rqSet{rqPreRepUnstableLibAllKiroNum}{0}
\rqSet{rqPreRepUnstableLibAllKiroDen}{217}
\rqSet{rqPreRepUnstableLibAllKiroPct}{0.0}
\rqSet{rqPreRepUnstableLibAllClaudeNum}{0}
\rqSet{rqPreRepUnstableLibAllClaudeDen}{206}
\rqSet{rqPreRepUnstableLibAllClaudePct}{0.0}
\rqSet{rqPreRepUnstableLibAllCodexNum}{0}
\rqSet{rqPreRepUnstableLibAllCodexDen}{171}
\rqSet{rqPreRepUnstableLibAllCodexPct}{0.0}
\rqSet{rqPreRepUnstableLibAllCTwoRustNum}{0}
\rqSet{rqPreRepUnstableLibAllCTwoRustDen}{141}
\rqSet{rqPreRepUnstableLibAllCTwoRustPct}{0.0}
\rqSet{rqPreRepUnstableLibAllLaertesNum}{0}
\rqSet{rqPreRepUnstableLibAllLaertesDen}{137}
\rqSet{rqPreRepUnstableLibAllLaertesPct}{0.0}
\rqSet{rqPreRepUnstableLibAllCTwoSaferRustNum}{0}
\rqSet{rqPreRepUnstableLibAllCTwoSaferRustDen}{105}
\rqSet{rqPreRepUnstableLibAllCTwoSaferRustPct}{0.0}
\rqSet{rqPreRepUnstableLibAllSmartCTwoRustNum}{0}
\rqSet{rqPreRepUnstableLibAllSmartCTwoRustDen}{0}
\rqSet{rqPreRepUnstableLibAllSmartCTwoRustPct}{n/a}
\rqSet{rqPreRepUnstableLibAllKimiNum}{0}
\rqSet{rqPreRepUnstableLibAllKimiDen}{85}
\rqSet{rqPreRepUnstableLibAllKimiPct}{0.0}
\rqSet{rqPreRepUnstableLibAllGptFiveNum}{0}
\rqSet{rqPreRepUnstableLibAllGptFiveDen}{111}
\rqSet{rqPreRepUnstableLibAllGptFivePct}{0.0}
\rqSet{rqPreRepUnstableLibAllGeminiNum}{0}
\rqSet{rqPreRepUnstableLibAllGeminiDen}{108}
\rqSet{rqPreRepUnstableLibAllGeminiPct}{0.0}
\rqSet{rqPreRepUnstableLibBOneNum}{0}
\rqSet{rqPreRepUnstableLibBOneDen}{637}
\rqSet{rqPreRepUnstableLibBOnePct}{0.0}
\rqSet{rqPreRepUnstableLibBTwoNum}{0}
\rqSet{rqPreRepUnstableLibBTwoDen}{472}
\rqSet{rqPreRepUnstableLibBTwoPct}{0.0}
\rqSet{rqPreRepUnstableLibPZeroNum}{0}
\rqSet{rqPreRepUnstableLibPZeroDen}{0}
\rqSet{rqPreRepUnstableLibPZeroPct}{n/a}
\rqSet{rqPreRepUnstableLibPOneNum}{0}
\rqSet{rqPreRepUnstableLibPOneDen}{172}
\rqSet{rqPreRepUnstableLibPOnePct}{0.0}
\rqSet{rqPreRepNonReproducedLibAllNum}{6}
\rqSet{rqPreRepNonReproducedLibAllDen}{1281}
\rqSet{rqPreRepNonReproducedLibAllPct}{0.5}
\rqSet{rqPreRepNonReproducedLibAllKiroNum}{0}
\rqSet{rqPreRepNonReproducedLibAllKiroDen}{217}
\rqSet{rqPreRepNonReproducedLibAllKiroPct}{0.0}
\rqSet{rqPreRepNonReproducedLibAllClaudeNum}{0}
\rqSet{rqPreRepNonReproducedLibAllClaudeDen}{206}
\rqSet{rqPreRepNonReproducedLibAllClaudePct}{0.0}
\rqSet{rqPreRepNonReproducedLibAllCodexNum}{3}
\rqSet{rqPreRepNonReproducedLibAllCodexDen}{171}
\rqSet{rqPreRepNonReproducedLibAllCodexPct}{1.8}
\rqSet{rqPreRepNonReproducedLibAllCTwoRustNum}{0}
\rqSet{rqPreRepNonReproducedLibAllCTwoRustDen}{141}
\rqSet{rqPreRepNonReproducedLibAllCTwoRustPct}{0.0}
\rqSet{rqPreRepNonReproducedLibAllLaertesNum}{0}
\rqSet{rqPreRepNonReproducedLibAllLaertesDen}{137}
\rqSet{rqPreRepNonReproducedLibAllLaertesPct}{0.0}
\rqSet{rqPreRepNonReproducedLibAllCTwoSaferRustNum}{1}
\rqSet{rqPreRepNonReproducedLibAllCTwoSaferRustDen}{105}
\rqSet{rqPreRepNonReproducedLibAllCTwoSaferRustPct}{1.0}
\rqSet{rqPreRepNonReproducedLibAllSmartCTwoRustNum}{0}
\rqSet{rqPreRepNonReproducedLibAllSmartCTwoRustDen}{0}
\rqSet{rqPreRepNonReproducedLibAllSmartCTwoRustPct}{n/a}
\rqSet{rqPreRepNonReproducedLibAllKimiNum}{0}
\rqSet{rqPreRepNonReproducedLibAllKimiDen}{85}
\rqSet{rqPreRepNonReproducedLibAllKimiPct}{0.0}
\rqSet{rqPreRepNonReproducedLibAllGptFiveNum}{1}
\rqSet{rqPreRepNonReproducedLibAllGptFiveDen}{111}
\rqSet{rqPreRepNonReproducedLibAllGptFivePct}{0.9}
\rqSet{rqPreRepNonReproducedLibAllGeminiNum}{1}
\rqSet{rqPreRepNonReproducedLibAllGeminiDen}{108}
\rqSet{rqPreRepNonReproducedLibAllGeminiPct}{0.9}
\rqSet{rqPreRepNonReproducedLibBOneNum}{1}
\rqSet{rqPreRepNonReproducedLibBOneDen}{637}
\rqSet{rqPreRepNonReproducedLibBOnePct}{0.2}
\rqSet{rqPreRepNonReproducedLibBTwoNum}{2}
\rqSet{rqPreRepNonReproducedLibBTwoDen}{472}
\rqSet{rqPreRepNonReproducedLibBTwoPct}{0.4}
\rqSet{rqPreRepNonReproducedLibPZeroNum}{0}
\rqSet{rqPreRepNonReproducedLibPZeroDen}{0}
\rqSet{rqPreRepNonReproducedLibPZeroPct}{n/a}
\rqSet{rqPreRepNonReproducedLibPOneNum}{3}
\rqSet{rqPreRepNonReproducedLibPOneDen}{172}
\rqSet{rqPreRepNonReproducedLibPOnePct}{1.7}
\rqSet{rqLeaveTotalLibAllNum}{176}
\rqSet{rqLeaveTotalLibAllDen}{1264}
\rqSet{rqLeaveTotalLibAllPct}{13.9}
\rqSet{rqLeaveTotalLibAllKiroNum}{24}
\rqSet{rqLeaveTotalLibAllKiroDen}{216}
\rqSet{rqLeaveTotalLibAllKiroPct}{11.1}
\rqSet{rqLeaveTotalLibAllClaudeNum}{24}
\rqSet{rqLeaveTotalLibAllClaudeDen}{205}
\rqSet{rqLeaveTotalLibAllClaudePct}{11.7}
\rqSet{rqLeaveTotalLibAllCodexNum}{24}
\rqSet{rqLeaveTotalLibAllCodexDen}{167}
\rqSet{rqLeaveTotalLibAllCodexPct}{14.4}
\rqSet{rqLeaveTotalLibAllCTwoRustNum}{24}
\rqSet{rqLeaveTotalLibAllCTwoRustDen}{140}
\rqSet{rqLeaveTotalLibAllCTwoRustPct}{17.1}
\rqSet{rqLeaveTotalLibAllLaertesNum}{22}
\rqSet{rqLeaveTotalLibAllLaertesDen}{136}
\rqSet{rqLeaveTotalLibAllLaertesPct}{16.2}
\rqSet{rqLeaveTotalLibAllCTwoSaferRustNum}{16}
\rqSet{rqLeaveTotalLibAllCTwoSaferRustDen}{103}
\rqSet{rqLeaveTotalLibAllCTwoSaferRustPct}{15.5}
\rqSet{rqLeaveTotalLibAllSmartCTwoRustNum}{0}
\rqSet{rqLeaveTotalLibAllSmartCTwoRustDen}{0}
\rqSet{rqLeaveTotalLibAllSmartCTwoRustPct}{n/a}
\rqSet{rqLeaveTotalLibAllKimiNum}{10}
\rqSet{rqLeaveTotalLibAllKimiDen}{83}
\rqSet{rqLeaveTotalLibAllKimiPct}{12.0}
\rqSet{rqLeaveTotalLibAllGptFiveNum}{16}
\rqSet{rqLeaveTotalLibAllGptFiveDen}{108}
\rqSet{rqLeaveTotalLibAllGptFivePct}{14.8}
\rqSet{rqLeaveTotalLibAllGeminiNum}{16}
\rqSet{rqLeaveTotalLibAllGeminiDen}{106}
\rqSet{rqLeaveTotalLibAllGeminiPct}{15.1}
\rqSet{rqLeaveTotalLibBOneNum}{96}
\rqSet{rqLeaveTotalLibBOneDen}{636}
\rqSet{rqLeaveTotalLibBOnePct}{15.1}
\rqSet{rqLeaveTotalLibBTwoNum}{80}
\rqSet{rqLeaveTotalLibBTwoDen}{459}
\rqSet{rqLeaveTotalLibBTwoPct}{17.4}
\rqSet{rqLeaveTotalLibPZeroNum}{0}
\rqSet{rqLeaveTotalLibPZeroDen}{0}
\rqSet{rqLeaveTotalLibPZeroPct}{n/a}
\rqSet{rqLeaveTotalLibPOneNum}{0}
\rqSet{rqLeaveTotalLibPOneDen}{169}
\rqSet{rqLeaveTotalLibPOnePct}{0.0}
\rqSet{rqLeaveNoInputFieldsLibAllNum}{21}
\rqSet{rqLeaveNoInputFieldsLibAllDen}{1264}
\rqSet{rqLeaveNoInputFieldsLibAllPct}{1.7}
\rqSet{rqLeaveNoInputFieldsLibAllKiroNum}{3}
\rqSet{rqLeaveNoInputFieldsLibAllKiroDen}{216}
\rqSet{rqLeaveNoInputFieldsLibAllKiroPct}{1.4}
\rqSet{rqLeaveNoInputFieldsLibAllClaudeNum}{3}
\rqSet{rqLeaveNoInputFieldsLibAllClaudeDen}{205}
\rqSet{rqLeaveNoInputFieldsLibAllClaudePct}{1.5}
\rqSet{rqLeaveNoInputFieldsLibAllCodexNum}{3}
\rqSet{rqLeaveNoInputFieldsLibAllCodexDen}{167}
\rqSet{rqLeaveNoInputFieldsLibAllCodexPct}{1.8}
\rqSet{rqLeaveNoInputFieldsLibAllCTwoRustNum}{3}
\rqSet{rqLeaveNoInputFieldsLibAllCTwoRustDen}{140}
\rqSet{rqLeaveNoInputFieldsLibAllCTwoRustPct}{2.1}
\rqSet{rqLeaveNoInputFieldsLibAllLaertesNum}{2}
\rqSet{rqLeaveNoInputFieldsLibAllLaertesDen}{136}
\rqSet{rqLeaveNoInputFieldsLibAllLaertesPct}{1.5}
\rqSet{rqLeaveNoInputFieldsLibAllCTwoSaferRustNum}{2}
\rqSet{rqLeaveNoInputFieldsLibAllCTwoSaferRustDen}{103}
\rqSet{rqLeaveNoInputFieldsLibAllCTwoSaferRustPct}{1.9}
\rqSet{rqLeaveNoInputFieldsLibAllSmartCTwoRustNum}{0}
\rqSet{rqLeaveNoInputFieldsLibAllSmartCTwoRustDen}{0}
\rqSet{rqLeaveNoInputFieldsLibAllSmartCTwoRustPct}{n/a}
\rqSet{rqLeaveNoInputFieldsLibAllKimiNum}{1}
\rqSet{rqLeaveNoInputFieldsLibAllKimiDen}{83}
\rqSet{rqLeaveNoInputFieldsLibAllKimiPct}{1.2}
\rqSet{rqLeaveNoInputFieldsLibAllGptFiveNum}{2}
\rqSet{rqLeaveNoInputFieldsLibAllGptFiveDen}{108}
\rqSet{rqLeaveNoInputFieldsLibAllGptFivePct}{1.9}
\rqSet{rqLeaveNoInputFieldsLibAllGeminiNum}{2}
\rqSet{rqLeaveNoInputFieldsLibAllGeminiDen}{106}
\rqSet{rqLeaveNoInputFieldsLibAllGeminiPct}{1.9}
\rqSet{rqLeaveNoInputFieldsLibBOneNum}{16}
\rqSet{rqLeaveNoInputFieldsLibBOneDen}{636}
\rqSet{rqLeaveNoInputFieldsLibBOnePct}{2.5}
\rqSet{rqLeaveNoInputFieldsLibBTwoNum}{5}
\rqSet{rqLeaveNoInputFieldsLibBTwoDen}{459}
\rqSet{rqLeaveNoInputFieldsLibBTwoPct}{1.1}
\rqSet{rqLeaveNoInputFieldsLibPZeroNum}{0}
\rqSet{rqLeaveNoInputFieldsLibPZeroDen}{0}
\rqSet{rqLeaveNoInputFieldsLibPZeroPct}{n/a}
\rqSet{rqLeaveNoInputFieldsLibPOneNum}{0}
\rqSet{rqLeaveNoInputFieldsLibPOneDen}{169}
\rqSet{rqLeaveNoInputFieldsLibPOnePct}{0.0}
\rqSet{rqLeaveCBuildFailedLibAllNum}{0}
\rqSet{rqLeaveCBuildFailedLibAllDen}{1264}
\rqSet{rqLeaveCBuildFailedLibAllPct}{0.0}
\rqSet{rqLeaveCBuildFailedLibAllKiroNum}{0}
\rqSet{rqLeaveCBuildFailedLibAllKiroDen}{216}
\rqSet{rqLeaveCBuildFailedLibAllKiroPct}{0.0}
\rqSet{rqLeaveCBuildFailedLibAllClaudeNum}{0}
\rqSet{rqLeaveCBuildFailedLibAllClaudeDen}{205}
\rqSet{rqLeaveCBuildFailedLibAllClaudePct}{0.0}
\rqSet{rqLeaveCBuildFailedLibAllCodexNum}{0}
\rqSet{rqLeaveCBuildFailedLibAllCodexDen}{167}
\rqSet{rqLeaveCBuildFailedLibAllCodexPct}{0.0}
\rqSet{rqLeaveCBuildFailedLibAllCTwoRustNum}{0}
\rqSet{rqLeaveCBuildFailedLibAllCTwoRustDen}{140}
\rqSet{rqLeaveCBuildFailedLibAllCTwoRustPct}{0.0}
\rqSet{rqLeaveCBuildFailedLibAllLaertesNum}{0}
\rqSet{rqLeaveCBuildFailedLibAllLaertesDen}{136}
\rqSet{rqLeaveCBuildFailedLibAllLaertesPct}{0.0}
\rqSet{rqLeaveCBuildFailedLibAllCTwoSaferRustNum}{0}
\rqSet{rqLeaveCBuildFailedLibAllCTwoSaferRustDen}{103}
\rqSet{rqLeaveCBuildFailedLibAllCTwoSaferRustPct}{0.0}
\rqSet{rqLeaveCBuildFailedLibAllSmartCTwoRustNum}{0}
\rqSet{rqLeaveCBuildFailedLibAllSmartCTwoRustDen}{0}
\rqSet{rqLeaveCBuildFailedLibAllSmartCTwoRustPct}{n/a}
\rqSet{rqLeaveCBuildFailedLibAllKimiNum}{0}
\rqSet{rqLeaveCBuildFailedLibAllKimiDen}{83}
\rqSet{rqLeaveCBuildFailedLibAllKimiPct}{0.0}
\rqSet{rqLeaveCBuildFailedLibAllGptFiveNum}{0}
\rqSet{rqLeaveCBuildFailedLibAllGptFiveDen}{108}
\rqSet{rqLeaveCBuildFailedLibAllGptFivePct}{0.0}
\rqSet{rqLeaveCBuildFailedLibAllGeminiNum}{0}
\rqSet{rqLeaveCBuildFailedLibAllGeminiDen}{106}
\rqSet{rqLeaveCBuildFailedLibAllGeminiPct}{0.0}
\rqSet{rqLeaveCBuildFailedLibBOneNum}{0}
\rqSet{rqLeaveCBuildFailedLibBOneDen}{636}
\rqSet{rqLeaveCBuildFailedLibBOnePct}{0.0}
\rqSet{rqLeaveCBuildFailedLibBTwoNum}{0}
\rqSet{rqLeaveCBuildFailedLibBTwoDen}{459}
\rqSet{rqLeaveCBuildFailedLibBTwoPct}{0.0}
\rqSet{rqLeaveCBuildFailedLibPZeroNum}{0}
\rqSet{rqLeaveCBuildFailedLibPZeroDen}{0}
\rqSet{rqLeaveCBuildFailedLibPZeroPct}{n/a}
\rqSet{rqLeaveCBuildFailedLibPOneNum}{0}
\rqSet{rqLeaveCBuildFailedLibPOneDen}{169}
\rqSet{rqLeaveCBuildFailedLibPOnePct}{0.0}
\rqSet{rqLeaveRustBuildFailedLibAllNum}{0}
\rqSet{rqLeaveRustBuildFailedLibAllDen}{1264}
\rqSet{rqLeaveRustBuildFailedLibAllPct}{0.0}
\rqSet{rqLeaveRustBuildFailedLibAllKiroNum}{0}
\rqSet{rqLeaveRustBuildFailedLibAllKiroDen}{216}
\rqSet{rqLeaveRustBuildFailedLibAllKiroPct}{0.0}
\rqSet{rqLeaveRustBuildFailedLibAllClaudeNum}{0}
\rqSet{rqLeaveRustBuildFailedLibAllClaudeDen}{205}
\rqSet{rqLeaveRustBuildFailedLibAllClaudePct}{0.0}
\rqSet{rqLeaveRustBuildFailedLibAllCodexNum}{0}
\rqSet{rqLeaveRustBuildFailedLibAllCodexDen}{167}
\rqSet{rqLeaveRustBuildFailedLibAllCodexPct}{0.0}
\rqSet{rqLeaveRustBuildFailedLibAllCTwoRustNum}{0}
\rqSet{rqLeaveRustBuildFailedLibAllCTwoRustDen}{140}
\rqSet{rqLeaveRustBuildFailedLibAllCTwoRustPct}{0.0}
\rqSet{rqLeaveRustBuildFailedLibAllLaertesNum}{0}
\rqSet{rqLeaveRustBuildFailedLibAllLaertesDen}{136}
\rqSet{rqLeaveRustBuildFailedLibAllLaertesPct}{0.0}
\rqSet{rqLeaveRustBuildFailedLibAllCTwoSaferRustNum}{0}
\rqSet{rqLeaveRustBuildFailedLibAllCTwoSaferRustDen}{103}
\rqSet{rqLeaveRustBuildFailedLibAllCTwoSaferRustPct}{0.0}
\rqSet{rqLeaveRustBuildFailedLibAllSmartCTwoRustNum}{0}
\rqSet{rqLeaveRustBuildFailedLibAllSmartCTwoRustDen}{0}
\rqSet{rqLeaveRustBuildFailedLibAllSmartCTwoRustPct}{n/a}
\rqSet{rqLeaveRustBuildFailedLibAllKimiNum}{0}
\rqSet{rqLeaveRustBuildFailedLibAllKimiDen}{83}
\rqSet{rqLeaveRustBuildFailedLibAllKimiPct}{0.0}
\rqSet{rqLeaveRustBuildFailedLibAllGptFiveNum}{0}
\rqSet{rqLeaveRustBuildFailedLibAllGptFiveDen}{108}
\rqSet{rqLeaveRustBuildFailedLibAllGptFivePct}{0.0}
\rqSet{rqLeaveRustBuildFailedLibAllGeminiNum}{0}
\rqSet{rqLeaveRustBuildFailedLibAllGeminiDen}{106}
\rqSet{rqLeaveRustBuildFailedLibAllGeminiPct}{0.0}
\rqSet{rqLeaveRustBuildFailedLibBOneNum}{0}
\rqSet{rqLeaveRustBuildFailedLibBOneDen}{636}
\rqSet{rqLeaveRustBuildFailedLibBOnePct}{0.0}
\rqSet{rqLeaveRustBuildFailedLibBTwoNum}{0}
\rqSet{rqLeaveRustBuildFailedLibBTwoDen}{459}
\rqSet{rqLeaveRustBuildFailedLibBTwoPct}{0.0}
\rqSet{rqLeaveRustBuildFailedLibPZeroNum}{0}
\rqSet{rqLeaveRustBuildFailedLibPZeroDen}{0}
\rqSet{rqLeaveRustBuildFailedLibPZeroPct}{n/a}
\rqSet{rqLeaveRustBuildFailedLibPOneNum}{0}
\rqSet{rqLeaveRustBuildFailedLibPOneDen}{169}
\rqSet{rqLeaveRustBuildFailedLibPOnePct}{0.0}
\rqSet{rqLeaveHarnessErrorLibAllNum}{0}
\rqSet{rqLeaveHarnessErrorLibAllDen}{1264}
\rqSet{rqLeaveHarnessErrorLibAllPct}{0.0}
\rqSet{rqLeaveHarnessErrorLibAllKiroNum}{0}
\rqSet{rqLeaveHarnessErrorLibAllKiroDen}{216}
\rqSet{rqLeaveHarnessErrorLibAllKiroPct}{0.0}
\rqSet{rqLeaveHarnessErrorLibAllClaudeNum}{0}
\rqSet{rqLeaveHarnessErrorLibAllClaudeDen}{205}
\rqSet{rqLeaveHarnessErrorLibAllClaudePct}{0.0}
\rqSet{rqLeaveHarnessErrorLibAllCodexNum}{0}
\rqSet{rqLeaveHarnessErrorLibAllCodexDen}{167}
\rqSet{rqLeaveHarnessErrorLibAllCodexPct}{0.0}
\rqSet{rqLeaveHarnessErrorLibAllCTwoRustNum}{0}
\rqSet{rqLeaveHarnessErrorLibAllCTwoRustDen}{140}
\rqSet{rqLeaveHarnessErrorLibAllCTwoRustPct}{0.0}
\rqSet{rqLeaveHarnessErrorLibAllLaertesNum}{0}
\rqSet{rqLeaveHarnessErrorLibAllLaertesDen}{136}
\rqSet{rqLeaveHarnessErrorLibAllLaertesPct}{0.0}
\rqSet{rqLeaveHarnessErrorLibAllCTwoSaferRustNum}{0}
\rqSet{rqLeaveHarnessErrorLibAllCTwoSaferRustDen}{103}
\rqSet{rqLeaveHarnessErrorLibAllCTwoSaferRustPct}{0.0}
\rqSet{rqLeaveHarnessErrorLibAllSmartCTwoRustNum}{0}
\rqSet{rqLeaveHarnessErrorLibAllSmartCTwoRustDen}{0}
\rqSet{rqLeaveHarnessErrorLibAllSmartCTwoRustPct}{n/a}
\rqSet{rqLeaveHarnessErrorLibAllKimiNum}{0}
\rqSet{rqLeaveHarnessErrorLibAllKimiDen}{83}
\rqSet{rqLeaveHarnessErrorLibAllKimiPct}{0.0}
\rqSet{rqLeaveHarnessErrorLibAllGptFiveNum}{0}
\rqSet{rqLeaveHarnessErrorLibAllGptFiveDen}{108}
\rqSet{rqLeaveHarnessErrorLibAllGptFivePct}{0.0}
\rqSet{rqLeaveHarnessErrorLibAllGeminiNum}{0}
\rqSet{rqLeaveHarnessErrorLibAllGeminiDen}{106}
\rqSet{rqLeaveHarnessErrorLibAllGeminiPct}{0.0}
\rqSet{rqLeaveHarnessErrorLibBOneNum}{0}
\rqSet{rqLeaveHarnessErrorLibBOneDen}{636}
\rqSet{rqLeaveHarnessErrorLibBOnePct}{0.0}
\rqSet{rqLeaveHarnessErrorLibBTwoNum}{0}
\rqSet{rqLeaveHarnessErrorLibBTwoDen}{459}
\rqSet{rqLeaveHarnessErrorLibBTwoPct}{0.0}
\rqSet{rqLeaveHarnessErrorLibPZeroNum}{0}
\rqSet{rqLeaveHarnessErrorLibPZeroDen}{0}
\rqSet{rqLeaveHarnessErrorLibPZeroPct}{n/a}
\rqSet{rqLeaveHarnessErrorLibPOneNum}{0}
\rqSet{rqLeaveHarnessErrorLibPOneDen}{169}
\rqSet{rqLeaveHarnessErrorLibPOnePct}{0.0}
\rqSet{rqLeaveNoInputsLibAllNum}{0}
\rqSet{rqLeaveNoInputsLibAllDen}{1264}
\rqSet{rqLeaveNoInputsLibAllPct}{0.0}
\rqSet{rqLeaveNoInputsLibAllKiroNum}{0}
\rqSet{rqLeaveNoInputsLibAllKiroDen}{216}
\rqSet{rqLeaveNoInputsLibAllKiroPct}{0.0}
\rqSet{rqLeaveNoInputsLibAllClaudeNum}{0}
\rqSet{rqLeaveNoInputsLibAllClaudeDen}{205}
\rqSet{rqLeaveNoInputsLibAllClaudePct}{0.0}
\rqSet{rqLeaveNoInputsLibAllCodexNum}{0}
\rqSet{rqLeaveNoInputsLibAllCodexDen}{167}
\rqSet{rqLeaveNoInputsLibAllCodexPct}{0.0}
\rqSet{rqLeaveNoInputsLibAllCTwoRustNum}{0}
\rqSet{rqLeaveNoInputsLibAllCTwoRustDen}{140}
\rqSet{rqLeaveNoInputsLibAllCTwoRustPct}{0.0}
\rqSet{rqLeaveNoInputsLibAllLaertesNum}{0}
\rqSet{rqLeaveNoInputsLibAllLaertesDen}{136}
\rqSet{rqLeaveNoInputsLibAllLaertesPct}{0.0}
\rqSet{rqLeaveNoInputsLibAllCTwoSaferRustNum}{0}
\rqSet{rqLeaveNoInputsLibAllCTwoSaferRustDen}{103}
\rqSet{rqLeaveNoInputsLibAllCTwoSaferRustPct}{0.0}
\rqSet{rqLeaveNoInputsLibAllSmartCTwoRustNum}{0}
\rqSet{rqLeaveNoInputsLibAllSmartCTwoRustDen}{0}
\rqSet{rqLeaveNoInputsLibAllSmartCTwoRustPct}{n/a}
\rqSet{rqLeaveNoInputsLibAllKimiNum}{0}
\rqSet{rqLeaveNoInputsLibAllKimiDen}{83}
\rqSet{rqLeaveNoInputsLibAllKimiPct}{0.0}
\rqSet{rqLeaveNoInputsLibAllGptFiveNum}{0}
\rqSet{rqLeaveNoInputsLibAllGptFiveDen}{108}
\rqSet{rqLeaveNoInputsLibAllGptFivePct}{0.0}
\rqSet{rqLeaveNoInputsLibAllGeminiNum}{0}
\rqSet{rqLeaveNoInputsLibAllGeminiDen}{106}
\rqSet{rqLeaveNoInputsLibAllGeminiPct}{0.0}
\rqSet{rqLeaveNoInputsLibBOneNum}{0}
\rqSet{rqLeaveNoInputsLibBOneDen}{636}
\rqSet{rqLeaveNoInputsLibBOnePct}{0.0}
\rqSet{rqLeaveNoInputsLibBTwoNum}{0}
\rqSet{rqLeaveNoInputsLibBTwoDen}{459}
\rqSet{rqLeaveNoInputsLibBTwoPct}{0.0}
\rqSet{rqLeaveNoInputsLibPZeroNum}{0}
\rqSet{rqLeaveNoInputsLibPZeroDen}{0}
\rqSet{rqLeaveNoInputsLibPZeroPct}{n/a}
\rqSet{rqLeaveNoInputsLibPOneNum}{0}
\rqSet{rqLeaveNoInputsLibPOneDen}{169}
\rqSet{rqLeaveNoInputsLibPOnePct}{0.0}
\rqSet{rqLeaveIncompleteLibAllNum}{64}
\rqSet{rqLeaveIncompleteLibAllDen}{1264}
\rqSet{rqLeaveIncompleteLibAllPct}{5.1}
\rqSet{rqLeaveIncompleteLibAllKiroNum}{9}
\rqSet{rqLeaveIncompleteLibAllKiroDen}{216}
\rqSet{rqLeaveIncompleteLibAllKiroPct}{4.2}
\rqSet{rqLeaveIncompleteLibAllClaudeNum}{9}
\rqSet{rqLeaveIncompleteLibAllClaudeDen}{205}
\rqSet{rqLeaveIncompleteLibAllClaudePct}{4.4}
\rqSet{rqLeaveIncompleteLibAllCodexNum}{9}
\rqSet{rqLeaveIncompleteLibAllCodexDen}{167}
\rqSet{rqLeaveIncompleteLibAllCodexPct}{5.4}
\rqSet{rqLeaveIncompleteLibAllCTwoRustNum}{9}
\rqSet{rqLeaveIncompleteLibAllCTwoRustDen}{140}
\rqSet{rqLeaveIncompleteLibAllCTwoRustPct}{6.4}
\rqSet{rqLeaveIncompleteLibAllLaertesNum}{9}
\rqSet{rqLeaveIncompleteLibAllLaertesDen}{136}
\rqSet{rqLeaveIncompleteLibAllLaertesPct}{6.6}
\rqSet{rqLeaveIncompleteLibAllCTwoSaferRustNum}{5}
\rqSet{rqLeaveIncompleteLibAllCTwoSaferRustDen}{103}
\rqSet{rqLeaveIncompleteLibAllCTwoSaferRustPct}{4.9}
\rqSet{rqLeaveIncompleteLibAllSmartCTwoRustNum}{0}
\rqSet{rqLeaveIncompleteLibAllSmartCTwoRustDen}{0}
\rqSet{rqLeaveIncompleteLibAllSmartCTwoRustPct}{n/a}
\rqSet{rqLeaveIncompleteLibAllKimiNum}{4}
\rqSet{rqLeaveIncompleteLibAllKimiDen}{83}
\rqSet{rqLeaveIncompleteLibAllKimiPct}{4.8}
\rqSet{rqLeaveIncompleteLibAllGptFiveNum}{5}
\rqSet{rqLeaveIncompleteLibAllGptFiveDen}{108}
\rqSet{rqLeaveIncompleteLibAllGptFivePct}{4.6}
\rqSet{rqLeaveIncompleteLibAllGeminiNum}{5}
\rqSet{rqLeaveIncompleteLibAllGeminiDen}{106}
\rqSet{rqLeaveIncompleteLibAllGeminiPct}{4.7}
\rqSet{rqLeaveIncompleteLibBOneNum}{32}
\rqSet{rqLeaveIncompleteLibBOneDen}{636}
\rqSet{rqLeaveIncompleteLibBOnePct}{5.0}
\rqSet{rqLeaveIncompleteLibBTwoNum}{32}
\rqSet{rqLeaveIncompleteLibBTwoDen}{459}
\rqSet{rqLeaveIncompleteLibBTwoPct}{7.0}
\rqSet{rqLeaveIncompleteLibPZeroNum}{0}
\rqSet{rqLeaveIncompleteLibPZeroDen}{0}
\rqSet{rqLeaveIncompleteLibPZeroPct}{n/a}
\rqSet{rqLeaveIncompleteLibPOneNum}{0}
\rqSet{rqLeaveIncompleteLibPOneDen}{169}
\rqSet{rqLeaveIncompleteLibPOnePct}{0.0}
\rqSet{rqLeaveSanInfraLibAllNum}{0}
\rqSet{rqLeaveSanInfraLibAllDen}{1264}
\rqSet{rqLeaveSanInfraLibAllPct}{0.0}
\rqSet{rqLeaveSanInfraLibAllKiroNum}{0}
\rqSet{rqLeaveSanInfraLibAllKiroDen}{216}
\rqSet{rqLeaveSanInfraLibAllKiroPct}{0.0}
\rqSet{rqLeaveSanInfraLibAllClaudeNum}{0}
\rqSet{rqLeaveSanInfraLibAllClaudeDen}{205}
\rqSet{rqLeaveSanInfraLibAllClaudePct}{0.0}
\rqSet{rqLeaveSanInfraLibAllCodexNum}{0}
\rqSet{rqLeaveSanInfraLibAllCodexDen}{167}
\rqSet{rqLeaveSanInfraLibAllCodexPct}{0.0}
\rqSet{rqLeaveSanInfraLibAllCTwoRustNum}{0}
\rqSet{rqLeaveSanInfraLibAllCTwoRustDen}{140}
\rqSet{rqLeaveSanInfraLibAllCTwoRustPct}{0.0}
\rqSet{rqLeaveSanInfraLibAllLaertesNum}{0}
\rqSet{rqLeaveSanInfraLibAllLaertesDen}{136}
\rqSet{rqLeaveSanInfraLibAllLaertesPct}{0.0}
\rqSet{rqLeaveSanInfraLibAllCTwoSaferRustNum}{0}
\rqSet{rqLeaveSanInfraLibAllCTwoSaferRustDen}{103}
\rqSet{rqLeaveSanInfraLibAllCTwoSaferRustPct}{0.0}
\rqSet{rqLeaveSanInfraLibAllSmartCTwoRustNum}{0}
\rqSet{rqLeaveSanInfraLibAllSmartCTwoRustDen}{0}
\rqSet{rqLeaveSanInfraLibAllSmartCTwoRustPct}{n/a}
\rqSet{rqLeaveSanInfraLibAllKimiNum}{0}
\rqSet{rqLeaveSanInfraLibAllKimiDen}{83}
\rqSet{rqLeaveSanInfraLibAllKimiPct}{0.0}
\rqSet{rqLeaveSanInfraLibAllGptFiveNum}{0}
\rqSet{rqLeaveSanInfraLibAllGptFiveDen}{108}
\rqSet{rqLeaveSanInfraLibAllGptFivePct}{0.0}
\rqSet{rqLeaveSanInfraLibAllGeminiNum}{0}
\rqSet{rqLeaveSanInfraLibAllGeminiDen}{106}
\rqSet{rqLeaveSanInfraLibAllGeminiPct}{0.0}
\rqSet{rqLeaveSanInfraLibBOneNum}{0}
\rqSet{rqLeaveSanInfraLibBOneDen}{636}
\rqSet{rqLeaveSanInfraLibBOnePct}{0.0}
\rqSet{rqLeaveSanInfraLibBTwoNum}{0}
\rqSet{rqLeaveSanInfraLibBTwoDen}{459}
\rqSet{rqLeaveSanInfraLibBTwoPct}{0.0}
\rqSet{rqLeaveSanInfraLibPZeroNum}{0}
\rqSet{rqLeaveSanInfraLibPZeroDen}{0}
\rqSet{rqLeaveSanInfraLibPZeroPct}{n/a}
\rqSet{rqLeaveSanInfraLibPOneNum}{0}
\rqSet{rqLeaveSanInfraLibPOneDen}{169}
\rqSet{rqLeaveSanInfraLibPOnePct}{0.0}
\rqSet{rqLeaveMissingSeedLibAllNum}{68}
\rqSet{rqLeaveMissingSeedLibAllDen}{1264}
\rqSet{rqLeaveMissingSeedLibAllPct}{5.4}
\rqSet{rqLeaveMissingSeedLibAllKiroNum}{9}
\rqSet{rqLeaveMissingSeedLibAllKiroDen}{216}
\rqSet{rqLeaveMissingSeedLibAllKiroPct}{4.2}
\rqSet{rqLeaveMissingSeedLibAllClaudeNum}{9}
\rqSet{rqLeaveMissingSeedLibAllClaudeDen}{205}
\rqSet{rqLeaveMissingSeedLibAllClaudePct}{4.4}
\rqSet{rqLeaveMissingSeedLibAllCodexNum}{9}
\rqSet{rqLeaveMissingSeedLibAllCodexDen}{167}
\rqSet{rqLeaveMissingSeedLibAllCodexPct}{5.4}
\rqSet{rqLeaveMissingSeedLibAllCTwoRustNum}{9}
\rqSet{rqLeaveMissingSeedLibAllCTwoRustDen}{140}
\rqSet{rqLeaveMissingSeedLibAllCTwoRustPct}{6.4}
\rqSet{rqLeaveMissingSeedLibAllLaertesNum}{8}
\rqSet{rqLeaveMissingSeedLibAllLaertesDen}{136}
\rqSet{rqLeaveMissingSeedLibAllLaertesPct}{5.9}
\rqSet{rqLeaveMissingSeedLibAllCTwoSaferRustNum}{6}
\rqSet{rqLeaveMissingSeedLibAllCTwoSaferRustDen}{103}
\rqSet{rqLeaveMissingSeedLibAllCTwoSaferRustPct}{5.8}
\rqSet{rqLeaveMissingSeedLibAllSmartCTwoRustNum}{0}
\rqSet{rqLeaveMissingSeedLibAllSmartCTwoRustDen}{0}
\rqSet{rqLeaveMissingSeedLibAllSmartCTwoRustPct}{n/a}
\rqSet{rqLeaveMissingSeedLibAllKimiNum}{5}
\rqSet{rqLeaveMissingSeedLibAllKimiDen}{83}
\rqSet{rqLeaveMissingSeedLibAllKimiPct}{6.0}
\rqSet{rqLeaveMissingSeedLibAllGptFiveNum}{7}
\rqSet{rqLeaveMissingSeedLibAllGptFiveDen}{108}
\rqSet{rqLeaveMissingSeedLibAllGptFivePct}{6.5}
\rqSet{rqLeaveMissingSeedLibAllGeminiNum}{6}
\rqSet{rqLeaveMissingSeedLibAllGeminiDen}{106}
\rqSet{rqLeaveMissingSeedLibAllGeminiPct}{5.7}
\rqSet{rqLeaveMissingSeedLibBOneNum}{48}
\rqSet{rqLeaveMissingSeedLibBOneDen}{636}
\rqSet{rqLeaveMissingSeedLibBOnePct}{7.5}
\rqSet{rqLeaveMissingSeedLibBTwoNum}{20}
\rqSet{rqLeaveMissingSeedLibBTwoDen}{459}
\rqSet{rqLeaveMissingSeedLibBTwoPct}{4.4}
\rqSet{rqLeaveMissingSeedLibPZeroNum}{0}
\rqSet{rqLeaveMissingSeedLibPZeroDen}{0}
\rqSet{rqLeaveMissingSeedLibPZeroPct}{n/a}
\rqSet{rqLeaveMissingSeedLibPOneNum}{0}
\rqSet{rqLeaveMissingSeedLibPOneDen}{169}
\rqSet{rqLeaveMissingSeedLibPOnePct}{0.0}
\rqSet{rqLeaveMissingRecordLibAllNum}{23}
\rqSet{rqLeaveMissingRecordLibAllDen}{1264}
\rqSet{rqLeaveMissingRecordLibAllPct}{1.8}
\rqSet{rqLeaveMissingRecordLibAllKiroNum}{3}
\rqSet{rqLeaveMissingRecordLibAllKiroDen}{216}
\rqSet{rqLeaveMissingRecordLibAllKiroPct}{1.4}
\rqSet{rqLeaveMissingRecordLibAllClaudeNum}{3}
\rqSet{rqLeaveMissingRecordLibAllClaudeDen}{205}
\rqSet{rqLeaveMissingRecordLibAllClaudePct}{1.5}
\rqSet{rqLeaveMissingRecordLibAllCodexNum}{3}
\rqSet{rqLeaveMissingRecordLibAllCodexDen}{167}
\rqSet{rqLeaveMissingRecordLibAllCodexPct}{1.8}
\rqSet{rqLeaveMissingRecordLibAllCTwoRustNum}{3}
\rqSet{rqLeaveMissingRecordLibAllCTwoRustDen}{140}
\rqSet{rqLeaveMissingRecordLibAllCTwoRustPct}{2.1}
\rqSet{rqLeaveMissingRecordLibAllLaertesNum}{3}
\rqSet{rqLeaveMissingRecordLibAllLaertesDen}{136}
\rqSet{rqLeaveMissingRecordLibAllLaertesPct}{2.2}
\rqSet{rqLeaveMissingRecordLibAllCTwoSaferRustNum}{3}
\rqSet{rqLeaveMissingRecordLibAllCTwoSaferRustDen}{103}
\rqSet{rqLeaveMissingRecordLibAllCTwoSaferRustPct}{2.9}
\rqSet{rqLeaveMissingRecordLibAllSmartCTwoRustNum}{0}
\rqSet{rqLeaveMissingRecordLibAllSmartCTwoRustDen}{0}
\rqSet{rqLeaveMissingRecordLibAllSmartCTwoRustPct}{n/a}
\rqSet{rqLeaveMissingRecordLibAllKimiNum}{0}
\rqSet{rqLeaveMissingRecordLibAllKimiDen}{83}
\rqSet{rqLeaveMissingRecordLibAllKimiPct}{0.0}
\rqSet{rqLeaveMissingRecordLibAllGptFiveNum}{2}
\rqSet{rqLeaveMissingRecordLibAllGptFiveDen}{108}
\rqSet{rqLeaveMissingRecordLibAllGptFivePct}{1.9}
\rqSet{rqLeaveMissingRecordLibAllGeminiNum}{3}
\rqSet{rqLeaveMissingRecordLibAllGeminiDen}{106}
\rqSet{rqLeaveMissingRecordLibAllGeminiPct}{2.8}
\rqSet{rqLeaveMissingRecordLibBOneNum}{0}
\rqSet{rqLeaveMissingRecordLibBOneDen}{636}
\rqSet{rqLeaveMissingRecordLibBOnePct}{0.0}
\rqSet{rqLeaveMissingRecordLibBTwoNum}{23}
\rqSet{rqLeaveMissingRecordLibBTwoDen}{459}
\rqSet{rqLeaveMissingRecordLibBTwoPct}{5.0}
\rqSet{rqLeaveMissingRecordLibPZeroNum}{0}
\rqSet{rqLeaveMissingRecordLibPZeroDen}{0}
\rqSet{rqLeaveMissingRecordLibPZeroPct}{n/a}
\rqSet{rqLeaveMissingRecordLibPOneNum}{0}
\rqSet{rqLeaveMissingRecordLibPOneDen}{169}
\rqSet{rqLeaveMissingRecordLibPOnePct}{0.0}
\rqSet{rqLeaveOtherLibAllNum}{0}
\rqSet{rqLeaveOtherLibAllDen}{1264}
\rqSet{rqLeaveOtherLibAllPct}{0.0}
\rqSet{rqLeaveOtherLibAllKiroNum}{0}
\rqSet{rqLeaveOtherLibAllKiroDen}{216}
\rqSet{rqLeaveOtherLibAllKiroPct}{0.0}
\rqSet{rqLeaveOtherLibAllClaudeNum}{0}
\rqSet{rqLeaveOtherLibAllClaudeDen}{205}
\rqSet{rqLeaveOtherLibAllClaudePct}{0.0}
\rqSet{rqLeaveOtherLibAllCodexNum}{0}
\rqSet{rqLeaveOtherLibAllCodexDen}{167}
\rqSet{rqLeaveOtherLibAllCodexPct}{0.0}
\rqSet{rqLeaveOtherLibAllCTwoRustNum}{0}
\rqSet{rqLeaveOtherLibAllCTwoRustDen}{140}
\rqSet{rqLeaveOtherLibAllCTwoRustPct}{0.0}
\rqSet{rqLeaveOtherLibAllLaertesNum}{0}
\rqSet{rqLeaveOtherLibAllLaertesDen}{136}
\rqSet{rqLeaveOtherLibAllLaertesPct}{0.0}
\rqSet{rqLeaveOtherLibAllCTwoSaferRustNum}{0}
\rqSet{rqLeaveOtherLibAllCTwoSaferRustDen}{103}
\rqSet{rqLeaveOtherLibAllCTwoSaferRustPct}{0.0}
\rqSet{rqLeaveOtherLibAllSmartCTwoRustNum}{0}
\rqSet{rqLeaveOtherLibAllSmartCTwoRustDen}{0}
\rqSet{rqLeaveOtherLibAllSmartCTwoRustPct}{n/a}
\rqSet{rqLeaveOtherLibAllKimiNum}{0}
\rqSet{rqLeaveOtherLibAllKimiDen}{83}
\rqSet{rqLeaveOtherLibAllKimiPct}{0.0}
\rqSet{rqLeaveOtherLibAllGptFiveNum}{0}
\rqSet{rqLeaveOtherLibAllGptFiveDen}{108}
\rqSet{rqLeaveOtherLibAllGptFivePct}{0.0}
\rqSet{rqLeaveOtherLibAllGeminiNum}{0}
\rqSet{rqLeaveOtherLibAllGeminiDen}{106}
\rqSet{rqLeaveOtherLibAllGeminiPct}{0.0}
\rqSet{rqLeaveOtherLibBOneNum}{0}
\rqSet{rqLeaveOtherLibBOneDen}{636}
\rqSet{rqLeaveOtherLibBOnePct}{0.0}
\rqSet{rqLeaveOtherLibBTwoNum}{0}
\rqSet{rqLeaveOtherLibBTwoDen}{459}
\rqSet{rqLeaveOtherLibBTwoPct}{0.0}
\rqSet{rqLeaveOtherLibPZeroNum}{0}
\rqSet{rqLeaveOtherLibPZeroDen}{0}
\rqSet{rqLeaveOtherLibPZeroPct}{n/a}
\rqSet{rqLeaveOtherLibPOneNum}{0}
\rqSet{rqLeaveOtherLibPOneDen}{169}
\rqSet{rqLeaveOtherLibPOnePct}{0.0}
\rqSet{rqUnexpectedLibAllNum}{0}
\rqSet{rqUnexpectedLibAllDen}{1264}
\rqSet{rqUnexpectedLibBOneNum}{0}
\rqSet{rqUnexpectedLibBOneDen}{636}
\rqSet{rqUnexpectedLibBTwoNum}{0}
\rqSet{rqUnexpectedLibBTwoDen}{459}
\rqSet{rqUnexpectedLibPZeroNum}{0}
\rqSet{rqUnexpectedLibPZeroDen}{0}
\rqSet{rqUnexpectedLibPOneNum}{0}
\rqSet{rqUnexpectedLibPOneDen}{169}
\rqSet{rqFPLibAllNum}{65}
\rqSet{rqFPLibAllDen}{1088}
\rqSet{rqFPLibAllPct}{6.0}
\rqSet{rqFPLibAllKiroNum}{7}
\rqSet{rqFPLibAllKiroDen}{192}
\rqSet{rqFPLibAllKiroPct}{3.6}
\rqSet{rqFPLibAllClaudeNum}{6}
\rqSet{rqFPLibAllClaudeDen}{181}
\rqSet{rqFPLibAllClaudePct}{3.3}
\rqSet{rqFPLibAllCodexNum}{2}
\rqSet{rqFPLibAllCodexDen}{143}
\rqSet{rqFPLibAllCodexPct}{1.4}
\rqSet{rqFPLibAllCTwoRustNum}{3}
\rqSet{rqFPLibAllCTwoRustDen}{116}
\rqSet{rqFPLibAllCTwoRustPct}{2.6}
\rqSet{rqFPLibAllLaertesNum}{3}
\rqSet{rqFPLibAllLaertesDen}{114}
\rqSet{rqFPLibAllLaertesPct}{2.6}
\rqSet{rqFPLibAllCTwoSaferRustNum}{7}
\rqSet{rqFPLibAllCTwoSaferRustDen}{87}
\rqSet{rqFPLibAllCTwoSaferRustPct}{8.0}
\rqSet{rqFPLibAllSmartCTwoRustNum}{0}
\rqSet{rqFPLibAllSmartCTwoRustDen}{0}
\rqSet{rqFPLibAllSmartCTwoRustPct}{n/a}
\rqSet{rqFPLibAllKimiNum}{16}
\rqSet{rqFPLibAllKimiDen}{73}
\rqSet{rqFPLibAllKimiPct}{21.9}
\rqSet{rqFPLibAllGptFiveNum}{15}
\rqSet{rqFPLibAllGptFiveDen}{92}
\rqSet{rqFPLibAllGptFivePct}{16.3}
\rqSet{rqFPLibAllGeminiNum}{6}
\rqSet{rqFPLibAllGeminiDen}{90}
\rqSet{rqFPLibAllGeminiPct}{6.7}
\rqSet{rqFPLibBOneNum}{38}
\rqSet{rqFPLibBOneDen}{540}
\rqSet{rqFPLibBOnePct}{7.0}
\rqSet{rqFPLibBOneKiroNum}{2}
\rqSet{rqFPLibBOneKiroDen}{65}
\rqSet{rqFPLibBOneKiroPct}{3.1}
\rqSet{rqFPLibBOneClaudeNum}{2}
\rqSet{rqFPLibBOneClaudeDen}{65}
\rqSet{rqFPLibBOneClaudePct}{3.1}
\rqSet{rqFPLibBOneCodexNum}{2}
\rqSet{rqFPLibBOneCodexDen}{65}
\rqSet{rqFPLibBOneCodexPct}{3.1}
\rqSet{rqFPLibBOneCTwoRustNum}{3}
\rqSet{rqFPLibBOneCTwoRustDen}{65}
\rqSet{rqFPLibBOneCTwoRustPct}{4.6}
\rqSet{rqFPLibBOneLaertesNum}{3}
\rqSet{rqFPLibBOneLaertesDen}{63}
\rqSet{rqFPLibBOneLaertesPct}{4.8}
\rqSet{rqFPLibBOneCTwoSaferRustNum}{5}
\rqSet{rqFPLibBOneCTwoSaferRustDen}{50}
\rqSet{rqFPLibBOneCTwoSaferRustPct}{10.0}
\rqSet{rqFPLibBOneSmartCTwoRustNum}{0}
\rqSet{rqFPLibBOneSmartCTwoRustDen}{0}
\rqSet{rqFPLibBOneSmartCTwoRustPct}{n/a}
\rqSet{rqFPLibBOneKimiNum}{9}
\rqSet{rqFPLibBOneKimiDen}{49}
\rqSet{rqFPLibBOneKimiPct}{18.4}
\rqSet{rqFPLibBOneGptFiveNum}{9}
\rqSet{rqFPLibBOneGptFiveDen}{57}
\rqSet{rqFPLibBOneGptFivePct}{15.8}
\rqSet{rqFPLibBOneGeminiNum}{3}
\rqSet{rqFPLibBOneGeminiDen}{61}
\rqSet{rqFPLibBOneGeminiPct}{4.9}
\rqSet{rqFPLibBTwoNum}{27}
\rqSet{rqFPLibBTwoDen}{379}
\rqSet{rqFPLibBTwoPct}{7.1}
\rqSet{rqFPLibBTwoKiroNum}{5}
\rqSet{rqFPLibBTwoKiroDen}{51}
\rqSet{rqFPLibBTwoKiroPct}{9.8}
\rqSet{rqFPLibBTwoClaudeNum}{4}
\rqSet{rqFPLibBTwoClaudeDen}{51}
\rqSet{rqFPLibBTwoClaudePct}{7.8}
\rqSet{rqFPLibBTwoCodexNum}{0}
\rqSet{rqFPLibBTwoCodexDen}{50}
\rqSet{rqFPLibBTwoCodexPct}{0.0}
\rqSet{rqFPLibBTwoCTwoRustNum}{0}
\rqSet{rqFPLibBTwoCTwoRustDen}{51}
\rqSet{rqFPLibBTwoCTwoRustPct}{0.0}
\rqSet{rqFPLibBTwoLaertesNum}{0}
\rqSet{rqFPLibBTwoLaertesDen}{51}
\rqSet{rqFPLibBTwoLaertesPct}{0.0}
\rqSet{rqFPLibBTwoCTwoSaferRustNum}{2}
\rqSet{rqFPLibBTwoCTwoSaferRustDen}{37}
\rqSet{rqFPLibBTwoCTwoSaferRustPct}{5.4}
\rqSet{rqFPLibBTwoSmartCTwoRustNum}{0}
\rqSet{rqFPLibBTwoSmartCTwoRustDen}{0}
\rqSet{rqFPLibBTwoSmartCTwoRustPct}{n/a}
\rqSet{rqFPLibBTwoKimiNum}{7}
\rqSet{rqFPLibBTwoKimiDen}{24}
\rqSet{rqFPLibBTwoKimiPct}{29.2}
\rqSet{rqFPLibBTwoGptFiveNum}{6}
\rqSet{rqFPLibBTwoGptFiveDen}{35}
\rqSet{rqFPLibBTwoGptFivePct}{17.1}
\rqSet{rqFPLibBTwoGeminiNum}{3}
\rqSet{rqFPLibBTwoGeminiDen}{29}
\rqSet{rqFPLibBTwoGeminiPct}{10.3}
\rqSet{rqFPLibPZeroNum}{0}
\rqSet{rqFPLibPZeroDen}{0}
\rqSet{rqFPLibPZeroPct}{n/a}
\rqSet{rqFPLibPZeroKiroNum}{0}
\rqSet{rqFPLibPZeroKiroDen}{0}
\rqSet{rqFPLibPZeroKiroPct}{n/a}
\rqSet{rqFPLibPZeroClaudeNum}{0}
\rqSet{rqFPLibPZeroClaudeDen}{0}
\rqSet{rqFPLibPZeroClaudePct}{n/a}
\rqSet{rqFPLibPZeroCodexNum}{0}
\rqSet{rqFPLibPZeroCodexDen}{0}
\rqSet{rqFPLibPZeroCodexPct}{n/a}
\rqSet{rqFPLibPZeroCTwoRustNum}{0}
\rqSet{rqFPLibPZeroCTwoRustDen}{0}
\rqSet{rqFPLibPZeroCTwoRustPct}{n/a}
\rqSet{rqFPLibPZeroLaertesNum}{0}
\rqSet{rqFPLibPZeroLaertesDen}{0}
\rqSet{rqFPLibPZeroLaertesPct}{n/a}
\rqSet{rqFPLibPZeroCTwoSaferRustNum}{0}
\rqSet{rqFPLibPZeroCTwoSaferRustDen}{0}
\rqSet{rqFPLibPZeroCTwoSaferRustPct}{n/a}
\rqSet{rqFPLibPZeroSmartCTwoRustNum}{0}
\rqSet{rqFPLibPZeroSmartCTwoRustDen}{0}
\rqSet{rqFPLibPZeroSmartCTwoRustPct}{n/a}
\rqSet{rqFPLibPZeroKimiNum}{0}
\rqSet{rqFPLibPZeroKimiDen}{0}
\rqSet{rqFPLibPZeroKimiPct}{n/a}
\rqSet{rqFPLibPZeroGptFiveNum}{0}
\rqSet{rqFPLibPZeroGptFiveDen}{0}
\rqSet{rqFPLibPZeroGptFivePct}{n/a}
\rqSet{rqFPLibPZeroGeminiNum}{0}
\rqSet{rqFPLibPZeroGeminiDen}{0}
\rqSet{rqFPLibPZeroGeminiPct}{n/a}
\rqSet{rqFPLibPOneNum}{0}
\rqSet{rqFPLibPOneDen}{169}
\rqSet{rqFPLibPOnePct}{0.0}
\rqSet{rqFPLibPOneKiroNum}{0}
\rqSet{rqFPLibPOneKiroDen}{76}
\rqSet{rqFPLibPOneKiroPct}{0.0}
\rqSet{rqFPLibPOneClaudeNum}{0}
\rqSet{rqFPLibPOneClaudeDen}{65}
\rqSet{rqFPLibPOneClaudePct}{0.0}
\rqSet{rqFPLibPOneCodexNum}{0}
\rqSet{rqFPLibPOneCodexDen}{28}
\rqSet{rqFPLibPOneCodexPct}{0.0}
\rqSet{rqFPLibPOneCTwoRustNum}{0}
\rqSet{rqFPLibPOneCTwoRustDen}{0}
\rqSet{rqFPLibPOneCTwoRustPct}{n/a}
\rqSet{rqFPLibPOneLaertesNum}{0}
\rqSet{rqFPLibPOneLaertesDen}{0}
\rqSet{rqFPLibPOneLaertesPct}{n/a}
\rqSet{rqFPLibPOneCTwoSaferRustNum}{0}
\rqSet{rqFPLibPOneCTwoSaferRustDen}{0}
\rqSet{rqFPLibPOneCTwoSaferRustPct}{n/a}
\rqSet{rqFPLibPOneSmartCTwoRustNum}{0}
\rqSet{rqFPLibPOneSmartCTwoRustDen}{0}
\rqSet{rqFPLibPOneSmartCTwoRustPct}{n/a}
\rqSet{rqFPLibPOneKimiNum}{0}
\rqSet{rqFPLibPOneKimiDen}{0}
\rqSet{rqFPLibPOneKimiPct}{n/a}
\rqSet{rqFPLibPOneGptFiveNum}{0}
\rqSet{rqFPLibPOneGptFiveDen}{0}
\rqSet{rqFPLibPOneGptFivePct}{n/a}
\rqSet{rqFPLibPOneGeminiNum}{0}
\rqSet{rqFPLibPOneGeminiDen}{0}
\rqSet{rqFPLibPOneGeminiPct}{n/a}
\rqSet{rqFPBoundLoLibAllNum}{65}
\rqSet{rqFPBoundLoLibAllDen}{1281}
\rqSet{rqFPBoundLoLibAllPct}{5.1}
\rqSet{rqFPBoundLoLibAllKiroNum}{7}
\rqSet{rqFPBoundLoLibAllKiroDen}{217}
\rqSet{rqFPBoundLoLibAllKiroPct}{3.2}
\rqSet{rqFPBoundLoLibAllClaudeNum}{6}
\rqSet{rqFPBoundLoLibAllClaudeDen}{206}
\rqSet{rqFPBoundLoLibAllClaudePct}{2.9}
\rqSet{rqFPBoundLoLibAllCodexNum}{2}
\rqSet{rqFPBoundLoLibAllCodexDen}{171}
\rqSet{rqFPBoundLoLibAllCodexPct}{1.2}
\rqSet{rqFPBoundLoLibAllCTwoRustNum}{3}
\rqSet{rqFPBoundLoLibAllCTwoRustDen}{141}
\rqSet{rqFPBoundLoLibAllCTwoRustPct}{2.1}
\rqSet{rqFPBoundLoLibAllLaertesNum}{3}
\rqSet{rqFPBoundLoLibAllLaertesDen}{137}
\rqSet{rqFPBoundLoLibAllLaertesPct}{2.2}
\rqSet{rqFPBoundLoLibAllCTwoSaferRustNum}{7}
\rqSet{rqFPBoundLoLibAllCTwoSaferRustDen}{105}
\rqSet{rqFPBoundLoLibAllCTwoSaferRustPct}{6.7}
\rqSet{rqFPBoundLoLibAllSmartCTwoRustNum}{0}
\rqSet{rqFPBoundLoLibAllSmartCTwoRustDen}{0}
\rqSet{rqFPBoundLoLibAllSmartCTwoRustPct}{n/a}
\rqSet{rqFPBoundLoLibAllKimiNum}{16}
\rqSet{rqFPBoundLoLibAllKimiDen}{85}
\rqSet{rqFPBoundLoLibAllKimiPct}{18.8}
\rqSet{rqFPBoundLoLibAllGptFiveNum}{15}
\rqSet{rqFPBoundLoLibAllGptFiveDen}{111}
\rqSet{rqFPBoundLoLibAllGptFivePct}{13.5}
\rqSet{rqFPBoundLoLibAllGeminiNum}{6}
\rqSet{rqFPBoundLoLibAllGeminiDen}{108}
\rqSet{rqFPBoundLoLibAllGeminiPct}{5.6}
\rqSet{rqFPBoundLoLibBOneNum}{38}
\rqSet{rqFPBoundLoLibBOneDen}{637}
\rqSet{rqFPBoundLoLibBOnePct}{6.0}
\rqSet{rqFPBoundLoLibBOneKiroNum}{2}
\rqSet{rqFPBoundLoLibBOneKiroDen}{77}
\rqSet{rqFPBoundLoLibBOneKiroPct}{2.6}
\rqSet{rqFPBoundLoLibBOneClaudeNum}{2}
\rqSet{rqFPBoundLoLibBOneClaudeDen}{77}
\rqSet{rqFPBoundLoLibBOneClaudePct}{2.6}
\rqSet{rqFPBoundLoLibBOneCodexNum}{2}
\rqSet{rqFPBoundLoLibBOneCodexDen}{77}
\rqSet{rqFPBoundLoLibBOneCodexPct}{2.6}
\rqSet{rqFPBoundLoLibBOneCTwoRustNum}{3}
\rqSet{rqFPBoundLoLibBOneCTwoRustDen}{77}
\rqSet{rqFPBoundLoLibBOneCTwoRustPct}{3.9}
\rqSet{rqFPBoundLoLibBOneLaertesNum}{3}
\rqSet{rqFPBoundLoLibBOneLaertesDen}{74}
\rqSet{rqFPBoundLoLibBOneLaertesPct}{4.1}
\rqSet{rqFPBoundLoLibBOneCTwoSaferRustNum}{5}
\rqSet{rqFPBoundLoLibBOneCTwoSaferRustDen}{60}
\rqSet{rqFPBoundLoLibBOneCTwoSaferRustPct}{8.3}
\rqSet{rqFPBoundLoLibBOneSmartCTwoRustNum}{0}
\rqSet{rqFPBoundLoLibBOneSmartCTwoRustDen}{0}
\rqSet{rqFPBoundLoLibBOneSmartCTwoRustPct}{n/a}
\rqSet{rqFPBoundLoLibBOneKimiNum}{9}
\rqSet{rqFPBoundLoLibBOneKimiDen}{57}
\rqSet{rqFPBoundLoLibBOneKimiPct}{15.8}
\rqSet{rqFPBoundLoLibBOneGptFiveNum}{9}
\rqSet{rqFPBoundLoLibBOneGptFiveDen}{67}
\rqSet{rqFPBoundLoLibBOneGptFivePct}{13.4}
\rqSet{rqFPBoundLoLibBOneGeminiNum}{3}
\rqSet{rqFPBoundLoLibBOneGeminiDen}{71}
\rqSet{rqFPBoundLoLibBOneGeminiPct}{4.2}
\rqSet{rqFPBoundLoLibBTwoNum}{27}
\rqSet{rqFPBoundLoLibBTwoDen}{472}
\rqSet{rqFPBoundLoLibBTwoPct}{5.7}
\rqSet{rqFPBoundLoLibBTwoKiroNum}{5}
\rqSet{rqFPBoundLoLibBTwoKiroDen}{64}
\rqSet{rqFPBoundLoLibBTwoKiroPct}{7.8}
\rqSet{rqFPBoundLoLibBTwoClaudeNum}{4}
\rqSet{rqFPBoundLoLibBTwoClaudeDen}{64}
\rqSet{rqFPBoundLoLibBTwoClaudePct}{6.2}
\rqSet{rqFPBoundLoLibBTwoCodexNum}{0}
\rqSet{rqFPBoundLoLibBTwoCodexDen}{63}
\rqSet{rqFPBoundLoLibBTwoCodexPct}{0.0}
\rqSet{rqFPBoundLoLibBTwoCTwoRustNum}{0}
\rqSet{rqFPBoundLoLibBTwoCTwoRustDen}{64}
\rqSet{rqFPBoundLoLibBTwoCTwoRustPct}{0.0}
\rqSet{rqFPBoundLoLibBTwoLaertesNum}{0}
\rqSet{rqFPBoundLoLibBTwoLaertesDen}{63}
\rqSet{rqFPBoundLoLibBTwoLaertesPct}{0.0}
\rqSet{rqFPBoundLoLibBTwoCTwoSaferRustNum}{2}
\rqSet{rqFPBoundLoLibBTwoCTwoSaferRustDen}{45}
\rqSet{rqFPBoundLoLibBTwoCTwoSaferRustPct}{4.4}
\rqSet{rqFPBoundLoLibBTwoSmartCTwoRustNum}{0}
\rqSet{rqFPBoundLoLibBTwoSmartCTwoRustDen}{0}
\rqSet{rqFPBoundLoLibBTwoSmartCTwoRustPct}{n/a}
\rqSet{rqFPBoundLoLibBTwoKimiNum}{7}
\rqSet{rqFPBoundLoLibBTwoKimiDen}{28}
\rqSet{rqFPBoundLoLibBTwoKimiPct}{25.0}
\rqSet{rqFPBoundLoLibBTwoGptFiveNum}{6}
\rqSet{rqFPBoundLoLibBTwoGptFiveDen}{44}
\rqSet{rqFPBoundLoLibBTwoGptFivePct}{13.6}
\rqSet{rqFPBoundLoLibBTwoGeminiNum}{3}
\rqSet{rqFPBoundLoLibBTwoGeminiDen}{37}
\rqSet{rqFPBoundLoLibBTwoGeminiPct}{8.1}
\rqSet{rqFPBoundLoLibPZeroNum}{0}
\rqSet{rqFPBoundLoLibPZeroDen}{0}
\rqSet{rqFPBoundLoLibPZeroPct}{n/a}
\rqSet{rqFPBoundLoLibPZeroKiroNum}{0}
\rqSet{rqFPBoundLoLibPZeroKiroDen}{0}
\rqSet{rqFPBoundLoLibPZeroKiroPct}{n/a}
\rqSet{rqFPBoundLoLibPZeroClaudeNum}{0}
\rqSet{rqFPBoundLoLibPZeroClaudeDen}{0}
\rqSet{rqFPBoundLoLibPZeroClaudePct}{n/a}
\rqSet{rqFPBoundLoLibPZeroCodexNum}{0}
\rqSet{rqFPBoundLoLibPZeroCodexDen}{0}
\rqSet{rqFPBoundLoLibPZeroCodexPct}{n/a}
\rqSet{rqFPBoundLoLibPZeroCTwoRustNum}{0}
\rqSet{rqFPBoundLoLibPZeroCTwoRustDen}{0}
\rqSet{rqFPBoundLoLibPZeroCTwoRustPct}{n/a}
\rqSet{rqFPBoundLoLibPZeroLaertesNum}{0}
\rqSet{rqFPBoundLoLibPZeroLaertesDen}{0}
\rqSet{rqFPBoundLoLibPZeroLaertesPct}{n/a}
\rqSet{rqFPBoundLoLibPZeroCTwoSaferRustNum}{0}
\rqSet{rqFPBoundLoLibPZeroCTwoSaferRustDen}{0}
\rqSet{rqFPBoundLoLibPZeroCTwoSaferRustPct}{n/a}
\rqSet{rqFPBoundLoLibPZeroSmartCTwoRustNum}{0}
\rqSet{rqFPBoundLoLibPZeroSmartCTwoRustDen}{0}
\rqSet{rqFPBoundLoLibPZeroSmartCTwoRustPct}{n/a}
\rqSet{rqFPBoundLoLibPZeroKimiNum}{0}
\rqSet{rqFPBoundLoLibPZeroKimiDen}{0}
\rqSet{rqFPBoundLoLibPZeroKimiPct}{n/a}
\rqSet{rqFPBoundLoLibPZeroGptFiveNum}{0}
\rqSet{rqFPBoundLoLibPZeroGptFiveDen}{0}
\rqSet{rqFPBoundLoLibPZeroGptFivePct}{n/a}
\rqSet{rqFPBoundLoLibPZeroGeminiNum}{0}
\rqSet{rqFPBoundLoLibPZeroGeminiDen}{0}
\rqSet{rqFPBoundLoLibPZeroGeminiPct}{n/a}
\rqSet{rqFPBoundLoLibPOneNum}{0}
\rqSet{rqFPBoundLoLibPOneDen}{172}
\rqSet{rqFPBoundLoLibPOnePct}{0.0}
\rqSet{rqFPBoundLoLibPOneKiroNum}{0}
\rqSet{rqFPBoundLoLibPOneKiroDen}{76}
\rqSet{rqFPBoundLoLibPOneKiroPct}{0.0}
\rqSet{rqFPBoundLoLibPOneClaudeNum}{0}
\rqSet{rqFPBoundLoLibPOneClaudeDen}{65}
\rqSet{rqFPBoundLoLibPOneClaudePct}{0.0}
\rqSet{rqFPBoundLoLibPOneCodexNum}{0}
\rqSet{rqFPBoundLoLibPOneCodexDen}{31}
\rqSet{rqFPBoundLoLibPOneCodexPct}{0.0}
\rqSet{rqFPBoundLoLibPOneCTwoRustNum}{0}
\rqSet{rqFPBoundLoLibPOneCTwoRustDen}{0}
\rqSet{rqFPBoundLoLibPOneCTwoRustPct}{n/a}
\rqSet{rqFPBoundLoLibPOneLaertesNum}{0}
\rqSet{rqFPBoundLoLibPOneLaertesDen}{0}
\rqSet{rqFPBoundLoLibPOneLaertesPct}{n/a}
\rqSet{rqFPBoundLoLibPOneCTwoSaferRustNum}{0}
\rqSet{rqFPBoundLoLibPOneCTwoSaferRustDen}{0}
\rqSet{rqFPBoundLoLibPOneCTwoSaferRustPct}{n/a}
\rqSet{rqFPBoundLoLibPOneSmartCTwoRustNum}{0}
\rqSet{rqFPBoundLoLibPOneSmartCTwoRustDen}{0}
\rqSet{rqFPBoundLoLibPOneSmartCTwoRustPct}{n/a}
\rqSet{rqFPBoundLoLibPOneKimiNum}{0}
\rqSet{rqFPBoundLoLibPOneKimiDen}{0}
\rqSet{rqFPBoundLoLibPOneKimiPct}{n/a}
\rqSet{rqFPBoundLoLibPOneGptFiveNum}{0}
\rqSet{rqFPBoundLoLibPOneGptFiveDen}{0}
\rqSet{rqFPBoundLoLibPOneGptFivePct}{n/a}
\rqSet{rqFPBoundLoLibPOneGeminiNum}{0}
\rqSet{rqFPBoundLoLibPOneGeminiDen}{0}
\rqSet{rqFPBoundLoLibPOneGeminiPct}{n/a}
\rqSet{rqFPBoundHiLibAllNum}{258}
\rqSet{rqFPBoundHiLibAllDen}{1281}
\rqSet{rqFPBoundHiLibAllPct}{20.1}
\rqSet{rqFPBoundHiLibAllKiroNum}{32}
\rqSet{rqFPBoundHiLibAllKiroDen}{217}
\rqSet{rqFPBoundHiLibAllKiroPct}{14.7}
\rqSet{rqFPBoundHiLibAllClaudeNum}{31}
\rqSet{rqFPBoundHiLibAllClaudeDen}{206}
\rqSet{rqFPBoundHiLibAllClaudePct}{15.0}
\rqSet{rqFPBoundHiLibAllCodexNum}{30}
\rqSet{rqFPBoundHiLibAllCodexDen}{171}
\rqSet{rqFPBoundHiLibAllCodexPct}{17.5}
\rqSet{rqFPBoundHiLibAllCTwoRustNum}{28}
\rqSet{rqFPBoundHiLibAllCTwoRustDen}{141}
\rqSet{rqFPBoundHiLibAllCTwoRustPct}{19.9}
\rqSet{rqFPBoundHiLibAllLaertesNum}{26}
\rqSet{rqFPBoundHiLibAllLaertesDen}{137}
\rqSet{rqFPBoundHiLibAllLaertesPct}{19.0}
\rqSet{rqFPBoundHiLibAllCTwoSaferRustNum}{25}
\rqSet{rqFPBoundHiLibAllCTwoSaferRustDen}{105}
\rqSet{rqFPBoundHiLibAllCTwoSaferRustPct}{23.8}
\rqSet{rqFPBoundHiLibAllSmartCTwoRustNum}{0}
\rqSet{rqFPBoundHiLibAllSmartCTwoRustDen}{0}
\rqSet{rqFPBoundHiLibAllSmartCTwoRustPct}{n/a}
\rqSet{rqFPBoundHiLibAllKimiNum}{28}
\rqSet{rqFPBoundHiLibAllKimiDen}{85}
\rqSet{rqFPBoundHiLibAllKimiPct}{32.9}
\rqSet{rqFPBoundHiLibAllGptFiveNum}{34}
\rqSet{rqFPBoundHiLibAllGptFiveDen}{111}
\rqSet{rqFPBoundHiLibAllGptFivePct}{30.6}
\rqSet{rqFPBoundHiLibAllGeminiNum}{24}
\rqSet{rqFPBoundHiLibAllGeminiDen}{108}
\rqSet{rqFPBoundHiLibAllGeminiPct}{22.2}
\rqSet{rqFPBoundHiLibBOneNum}{135}
\rqSet{rqFPBoundHiLibBOneDen}{637}
\rqSet{rqFPBoundHiLibBOnePct}{21.2}
\rqSet{rqFPBoundHiLibBOneKiroNum}{14}
\rqSet{rqFPBoundHiLibBOneKiroDen}{77}
\rqSet{rqFPBoundHiLibBOneKiroPct}{18.2}
\rqSet{rqFPBoundHiLibBOneClaudeNum}{14}
\rqSet{rqFPBoundHiLibBOneClaudeDen}{77}
\rqSet{rqFPBoundHiLibBOneClaudePct}{18.2}
\rqSet{rqFPBoundHiLibBOneCodexNum}{14}
\rqSet{rqFPBoundHiLibBOneCodexDen}{77}
\rqSet{rqFPBoundHiLibBOneCodexPct}{18.2}
\rqSet{rqFPBoundHiLibBOneCTwoRustNum}{15}
\rqSet{rqFPBoundHiLibBOneCTwoRustDen}{77}
\rqSet{rqFPBoundHiLibBOneCTwoRustPct}{19.5}
\rqSet{rqFPBoundHiLibBOneLaertesNum}{14}
\rqSet{rqFPBoundHiLibBOneLaertesDen}{74}
\rqSet{rqFPBoundHiLibBOneLaertesPct}{18.9}
\rqSet{rqFPBoundHiLibBOneCTwoSaferRustNum}{15}
\rqSet{rqFPBoundHiLibBOneCTwoSaferRustDen}{60}
\rqSet{rqFPBoundHiLibBOneCTwoSaferRustPct}{25.0}
\rqSet{rqFPBoundHiLibBOneSmartCTwoRustNum}{0}
\rqSet{rqFPBoundHiLibBOneSmartCTwoRustDen}{0}
\rqSet{rqFPBoundHiLibBOneSmartCTwoRustPct}{n/a}
\rqSet{rqFPBoundHiLibBOneKimiNum}{17}
\rqSet{rqFPBoundHiLibBOneKimiDen}{57}
\rqSet{rqFPBoundHiLibBOneKimiPct}{29.8}
\rqSet{rqFPBoundHiLibBOneGptFiveNum}{19}
\rqSet{rqFPBoundHiLibBOneGptFiveDen}{67}
\rqSet{rqFPBoundHiLibBOneGptFivePct}{28.4}
\rqSet{rqFPBoundHiLibBOneGeminiNum}{13}
\rqSet{rqFPBoundHiLibBOneGeminiDen}{71}
\rqSet{rqFPBoundHiLibBOneGeminiPct}{18.3}
\rqSet{rqFPBoundHiLibBTwoNum}{120}
\rqSet{rqFPBoundHiLibBTwoDen}{472}
\rqSet{rqFPBoundHiLibBTwoPct}{25.4}
\rqSet{rqFPBoundHiLibBTwoKiroNum}{18}
\rqSet{rqFPBoundHiLibBTwoKiroDen}{64}
\rqSet{rqFPBoundHiLibBTwoKiroPct}{28.1}
\rqSet{rqFPBoundHiLibBTwoClaudeNum}{17}
\rqSet{rqFPBoundHiLibBTwoClaudeDen}{64}
\rqSet{rqFPBoundHiLibBTwoClaudePct}{26.6}
\rqSet{rqFPBoundHiLibBTwoCodexNum}{13}
\rqSet{rqFPBoundHiLibBTwoCodexDen}{63}
\rqSet{rqFPBoundHiLibBTwoCodexPct}{20.6}
\rqSet{rqFPBoundHiLibBTwoCTwoRustNum}{13}
\rqSet{rqFPBoundHiLibBTwoCTwoRustDen}{64}
\rqSet{rqFPBoundHiLibBTwoCTwoRustPct}{20.3}
\rqSet{rqFPBoundHiLibBTwoLaertesNum}{12}
\rqSet{rqFPBoundHiLibBTwoLaertesDen}{63}
\rqSet{rqFPBoundHiLibBTwoLaertesPct}{19.0}
\rqSet{rqFPBoundHiLibBTwoCTwoSaferRustNum}{10}
\rqSet{rqFPBoundHiLibBTwoCTwoSaferRustDen}{45}
\rqSet{rqFPBoundHiLibBTwoCTwoSaferRustPct}{22.2}
\rqSet{rqFPBoundHiLibBTwoSmartCTwoRustNum}{0}
\rqSet{rqFPBoundHiLibBTwoSmartCTwoRustDen}{0}
\rqSet{rqFPBoundHiLibBTwoSmartCTwoRustPct}{n/a}
\rqSet{rqFPBoundHiLibBTwoKimiNum}{11}
\rqSet{rqFPBoundHiLibBTwoKimiDen}{28}
\rqSet{rqFPBoundHiLibBTwoKimiPct}{39.3}
\rqSet{rqFPBoundHiLibBTwoGptFiveNum}{15}
\rqSet{rqFPBoundHiLibBTwoGptFiveDen}{44}
\rqSet{rqFPBoundHiLibBTwoGptFivePct}{34.1}
\rqSet{rqFPBoundHiLibBTwoGeminiNum}{11}
\rqSet{rqFPBoundHiLibBTwoGeminiDen}{37}
\rqSet{rqFPBoundHiLibBTwoGeminiPct}{29.7}
\rqSet{rqFPBoundHiLibPZeroNum}{0}
\rqSet{rqFPBoundHiLibPZeroDen}{0}
\rqSet{rqFPBoundHiLibPZeroPct}{n/a}
\rqSet{rqFPBoundHiLibPZeroKiroNum}{0}
\rqSet{rqFPBoundHiLibPZeroKiroDen}{0}
\rqSet{rqFPBoundHiLibPZeroKiroPct}{n/a}
\rqSet{rqFPBoundHiLibPZeroClaudeNum}{0}
\rqSet{rqFPBoundHiLibPZeroClaudeDen}{0}
\rqSet{rqFPBoundHiLibPZeroClaudePct}{n/a}
\rqSet{rqFPBoundHiLibPZeroCodexNum}{0}
\rqSet{rqFPBoundHiLibPZeroCodexDen}{0}
\rqSet{rqFPBoundHiLibPZeroCodexPct}{n/a}
\rqSet{rqFPBoundHiLibPZeroCTwoRustNum}{0}
\rqSet{rqFPBoundHiLibPZeroCTwoRustDen}{0}
\rqSet{rqFPBoundHiLibPZeroCTwoRustPct}{n/a}
\rqSet{rqFPBoundHiLibPZeroLaertesNum}{0}
\rqSet{rqFPBoundHiLibPZeroLaertesDen}{0}
\rqSet{rqFPBoundHiLibPZeroLaertesPct}{n/a}
\rqSet{rqFPBoundHiLibPZeroCTwoSaferRustNum}{0}
\rqSet{rqFPBoundHiLibPZeroCTwoSaferRustDen}{0}
\rqSet{rqFPBoundHiLibPZeroCTwoSaferRustPct}{n/a}
\rqSet{rqFPBoundHiLibPZeroSmartCTwoRustNum}{0}
\rqSet{rqFPBoundHiLibPZeroSmartCTwoRustDen}{0}
\rqSet{rqFPBoundHiLibPZeroSmartCTwoRustPct}{n/a}
\rqSet{rqFPBoundHiLibPZeroKimiNum}{0}
\rqSet{rqFPBoundHiLibPZeroKimiDen}{0}
\rqSet{rqFPBoundHiLibPZeroKimiPct}{n/a}
\rqSet{rqFPBoundHiLibPZeroGptFiveNum}{0}
\rqSet{rqFPBoundHiLibPZeroGptFiveDen}{0}
\rqSet{rqFPBoundHiLibPZeroGptFivePct}{n/a}
\rqSet{rqFPBoundHiLibPZeroGeminiNum}{0}
\rqSet{rqFPBoundHiLibPZeroGeminiDen}{0}
\rqSet{rqFPBoundHiLibPZeroGeminiPct}{n/a}
\rqSet{rqFPBoundHiLibPOneNum}{3}
\rqSet{rqFPBoundHiLibPOneDen}{172}
\rqSet{rqFPBoundHiLibPOnePct}{1.7}
\rqSet{rqFPBoundHiLibPOneKiroNum}{0}
\rqSet{rqFPBoundHiLibPOneKiroDen}{76}
\rqSet{rqFPBoundHiLibPOneKiroPct}{0.0}
\rqSet{rqFPBoundHiLibPOneClaudeNum}{0}
\rqSet{rqFPBoundHiLibPOneClaudeDen}{65}
\rqSet{rqFPBoundHiLibPOneClaudePct}{0.0}
\rqSet{rqFPBoundHiLibPOneCodexNum}{3}
\rqSet{rqFPBoundHiLibPOneCodexDen}{31}
\rqSet{rqFPBoundHiLibPOneCodexPct}{9.7}
\rqSet{rqFPBoundHiLibPOneCTwoRustNum}{0}
\rqSet{rqFPBoundHiLibPOneCTwoRustDen}{0}
\rqSet{rqFPBoundHiLibPOneCTwoRustPct}{n/a}
\rqSet{rqFPBoundHiLibPOneLaertesNum}{0}
\rqSet{rqFPBoundHiLibPOneLaertesDen}{0}
\rqSet{rqFPBoundHiLibPOneLaertesPct}{n/a}
\rqSet{rqFPBoundHiLibPOneCTwoSaferRustNum}{0}
\rqSet{rqFPBoundHiLibPOneCTwoSaferRustDen}{0}
\rqSet{rqFPBoundHiLibPOneCTwoSaferRustPct}{n/a}
\rqSet{rqFPBoundHiLibPOneSmartCTwoRustNum}{0}
\rqSet{rqFPBoundHiLibPOneSmartCTwoRustDen}{0}
\rqSet{rqFPBoundHiLibPOneSmartCTwoRustPct}{n/a}
\rqSet{rqFPBoundHiLibPOneKimiNum}{0}
\rqSet{rqFPBoundHiLibPOneKimiDen}{0}
\rqSet{rqFPBoundHiLibPOneKimiPct}{n/a}
\rqSet{rqFPBoundHiLibPOneGptFiveNum}{0}
\rqSet{rqFPBoundHiLibPOneGptFiveDen}{0}
\rqSet{rqFPBoundHiLibPOneGptFivePct}{n/a}
\rqSet{rqFPBoundHiLibPOneGeminiNum}{0}
\rqSet{rqFPBoundHiLibPOneGeminiDen}{0}
\rqSet{rqFPBoundHiLibPOneGeminiPct}{n/a}
\rqSet{rqSeedDetOneLibAllNum}{4}
\rqSet{rqSeedDetOneLibAllDen}{65}
\rqSet{rqSeedDetOneLibAllPct}{6.2}
\rqSet{rqSeedDetOneLibAllKiroNum}{0}
\rqSet{rqSeedDetOneLibAllKiroDen}{7}
\rqSet{rqSeedDetOneLibAllKiroPct}{0.0}
\rqSet{rqSeedDetOneLibAllClaudeNum}{0}
\rqSet{rqSeedDetOneLibAllClaudeDen}{6}
\rqSet{rqSeedDetOneLibAllClaudePct}{0.0}
\rqSet{rqSeedDetOneLibAllCodexNum}{0}
\rqSet{rqSeedDetOneLibAllCodexDen}{2}
\rqSet{rqSeedDetOneLibAllCodexPct}{0.0}
\rqSet{rqSeedDetOneLibAllCTwoRustNum}{0}
\rqSet{rqSeedDetOneLibAllCTwoRustDen}{3}
\rqSet{rqSeedDetOneLibAllCTwoRustPct}{0.0}
\rqSet{rqSeedDetOneLibAllLaertesNum}{0}
\rqSet{rqSeedDetOneLibAllLaertesDen}{3}
\rqSet{rqSeedDetOneLibAllLaertesPct}{0.0}
\rqSet{rqSeedDetOneLibAllCTwoSaferRustNum}{0}
\rqSet{rqSeedDetOneLibAllCTwoSaferRustDen}{7}
\rqSet{rqSeedDetOneLibAllCTwoSaferRustPct}{0.0}
\rqSet{rqSeedDetOneLibAllSmartCTwoRustNum}{0}
\rqSet{rqSeedDetOneLibAllSmartCTwoRustDen}{0}
\rqSet{rqSeedDetOneLibAllSmartCTwoRustPct}{n/a}
\rqSet{rqSeedDetOneLibAllKimiNum}{2}
\rqSet{rqSeedDetOneLibAllKimiDen}{16}
\rqSet{rqSeedDetOneLibAllKimiPct}{12.5}
\rqSet{rqSeedDetOneLibAllGptFiveNum}{1}
\rqSet{rqSeedDetOneLibAllGptFiveDen}{15}
\rqSet{rqSeedDetOneLibAllGptFivePct}{6.7}
\rqSet{rqSeedDetOneLibAllGeminiNum}{1}
\rqSet{rqSeedDetOneLibAllGeminiDen}{6}
\rqSet{rqSeedDetOneLibAllGeminiPct}{16.7}
\rqSet{rqSeedDetOneLibBOneNum}{3}
\rqSet{rqSeedDetOneLibBOneDen}{38}
\rqSet{rqSeedDetOneLibBOnePct}{7.9}
\rqSet{rqSeedDetOneLibBTwoNum}{1}
\rqSet{rqSeedDetOneLibBTwoDen}{27}
\rqSet{rqSeedDetOneLibBTwoPct}{3.7}
\rqSet{rqSeedDetOneLibPZeroNum}{0}
\rqSet{rqSeedDetOneLibPZeroDen}{0}
\rqSet{rqSeedDetOneLibPZeroPct}{n/a}
\rqSet{rqSeedDetOneLibPOneNum}{0}
\rqSet{rqSeedDetOneLibPOneDen}{0}
\rqSet{rqSeedDetOneLibPOnePct}{n/a}
\rqSet{rqSeedDetTwoLibAllNum}{1}
\rqSet{rqSeedDetTwoLibAllDen}{65}
\rqSet{rqSeedDetTwoLibAllPct}{1.5}
\rqSet{rqSeedDetTwoLibAllKiroNum}{0}
\rqSet{rqSeedDetTwoLibAllKiroDen}{7}
\rqSet{rqSeedDetTwoLibAllKiroPct}{0.0}
\rqSet{rqSeedDetTwoLibAllClaudeNum}{0}
\rqSet{rqSeedDetTwoLibAllClaudeDen}{6}
\rqSet{rqSeedDetTwoLibAllClaudePct}{0.0}
\rqSet{rqSeedDetTwoLibAllCodexNum}{0}
\rqSet{rqSeedDetTwoLibAllCodexDen}{2}
\rqSet{rqSeedDetTwoLibAllCodexPct}{0.0}
\rqSet{rqSeedDetTwoLibAllCTwoRustNum}{0}
\rqSet{rqSeedDetTwoLibAllCTwoRustDen}{3}
\rqSet{rqSeedDetTwoLibAllCTwoRustPct}{0.0}
\rqSet{rqSeedDetTwoLibAllLaertesNum}{0}
\rqSet{rqSeedDetTwoLibAllLaertesDen}{3}
\rqSet{rqSeedDetTwoLibAllLaertesPct}{0.0}
\rqSet{rqSeedDetTwoLibAllCTwoSaferRustNum}{0}
\rqSet{rqSeedDetTwoLibAllCTwoSaferRustDen}{7}
\rqSet{rqSeedDetTwoLibAllCTwoSaferRustPct}{0.0}
\rqSet{rqSeedDetTwoLibAllSmartCTwoRustNum}{0}
\rqSet{rqSeedDetTwoLibAllSmartCTwoRustDen}{0}
\rqSet{rqSeedDetTwoLibAllSmartCTwoRustPct}{n/a}
\rqSet{rqSeedDetTwoLibAllKimiNum}{1}
\rqSet{rqSeedDetTwoLibAllKimiDen}{16}
\rqSet{rqSeedDetTwoLibAllKimiPct}{6.2}
\rqSet{rqSeedDetTwoLibAllGptFiveNum}{0}
\rqSet{rqSeedDetTwoLibAllGptFiveDen}{15}
\rqSet{rqSeedDetTwoLibAllGptFivePct}{0.0}
\rqSet{rqSeedDetTwoLibAllGeminiNum}{0}
\rqSet{rqSeedDetTwoLibAllGeminiDen}{6}
\rqSet{rqSeedDetTwoLibAllGeminiPct}{0.0}
\rqSet{rqSeedDetTwoLibBOneNum}{0}
\rqSet{rqSeedDetTwoLibBOneDen}{38}
\rqSet{rqSeedDetTwoLibBOnePct}{0.0}
\rqSet{rqSeedDetTwoLibBTwoNum}{1}
\rqSet{rqSeedDetTwoLibBTwoDen}{27}
\rqSet{rqSeedDetTwoLibBTwoPct}{3.7}
\rqSet{rqSeedDetTwoLibPZeroNum}{0}
\rqSet{rqSeedDetTwoLibPZeroDen}{0}
\rqSet{rqSeedDetTwoLibPZeroPct}{n/a}
\rqSet{rqSeedDetTwoLibPOneNum}{0}
\rqSet{rqSeedDetTwoLibPOneDen}{0}
\rqSet{rqSeedDetTwoLibPOnePct}{n/a}
\rqSet{rqSeedDetThreeLibAllNum}{60}
\rqSet{rqSeedDetThreeLibAllDen}{65}
\rqSet{rqSeedDetThreeLibAllPct}{92.3}
\rqSet{rqSeedDetThreeLibAllKiroNum}{7}
\rqSet{rqSeedDetThreeLibAllKiroDen}{7}
\rqSet{rqSeedDetThreeLibAllKiroPct}{100.0}
\rqSet{rqSeedDetThreeLibAllClaudeNum}{6}
\rqSet{rqSeedDetThreeLibAllClaudeDen}{6}
\rqSet{rqSeedDetThreeLibAllClaudePct}{100.0}
\rqSet{rqSeedDetThreeLibAllCodexNum}{2}
\rqSet{rqSeedDetThreeLibAllCodexDen}{2}
\rqSet{rqSeedDetThreeLibAllCodexPct}{100.0}
\rqSet{rqSeedDetThreeLibAllCTwoRustNum}{3}
\rqSet{rqSeedDetThreeLibAllCTwoRustDen}{3}
\rqSet{rqSeedDetThreeLibAllCTwoRustPct}{100.0}
\rqSet{rqSeedDetThreeLibAllLaertesNum}{3}
\rqSet{rqSeedDetThreeLibAllLaertesDen}{3}
\rqSet{rqSeedDetThreeLibAllLaertesPct}{100.0}
\rqSet{rqSeedDetThreeLibAllCTwoSaferRustNum}{7}
\rqSet{rqSeedDetThreeLibAllCTwoSaferRustDen}{7}
\rqSet{rqSeedDetThreeLibAllCTwoSaferRustPct}{100.0}
\rqSet{rqSeedDetThreeLibAllSmartCTwoRustNum}{0}
\rqSet{rqSeedDetThreeLibAllSmartCTwoRustDen}{0}
\rqSet{rqSeedDetThreeLibAllSmartCTwoRustPct}{n/a}
\rqSet{rqSeedDetThreeLibAllKimiNum}{13}
\rqSet{rqSeedDetThreeLibAllKimiDen}{16}
\rqSet{rqSeedDetThreeLibAllKimiPct}{81.2}
\rqSet{rqSeedDetThreeLibAllGptFiveNum}{14}
\rqSet{rqSeedDetThreeLibAllGptFiveDen}{15}
\rqSet{rqSeedDetThreeLibAllGptFivePct}{93.3}
\rqSet{rqSeedDetThreeLibAllGeminiNum}{5}
\rqSet{rqSeedDetThreeLibAllGeminiDen}{6}
\rqSet{rqSeedDetThreeLibAllGeminiPct}{83.3}
\rqSet{rqSeedDetThreeLibBOneNum}{35}
\rqSet{rqSeedDetThreeLibBOneDen}{38}
\rqSet{rqSeedDetThreeLibBOnePct}{92.1}
\rqSet{rqSeedDetThreeLibBTwoNum}{25}
\rqSet{rqSeedDetThreeLibBTwoDen}{27}
\rqSet{rqSeedDetThreeLibBTwoPct}{92.6}
\rqSet{rqSeedDetThreeLibPZeroNum}{0}
\rqSet{rqSeedDetThreeLibPZeroDen}{0}
\rqSet{rqSeedDetThreeLibPZeroPct}{n/a}
\rqSet{rqSeedDetThreeLibPOneNum}{0}
\rqSet{rqSeedDetThreeLibPOneDen}{0}
\rqSet{rqSeedDetThreeLibPOnePct}{n/a}
\rqSet{rqFPSeedZeroLibAllNum}{62}
\rqSet{rqFPSeedZeroLibAllDen}{1156}
\rqSet{rqFPSeedZeroLibAllPct}{5.4}
\rqSet{rqFPSeedZeroLibAllKiroNum}{7}
\rqSet{rqFPSeedZeroLibAllKiroDen}{201}
\rqSet{rqFPSeedZeroLibAllKiroPct}{3.5}
\rqSet{rqFPSeedZeroLibAllClaudeNum}{6}
\rqSet{rqFPSeedZeroLibAllClaudeDen}{190}
\rqSet{rqFPSeedZeroLibAllClaudePct}{3.2}
\rqSet{rqFPSeedZeroLibAllCodexNum}{3}
\rqSet{rqFPSeedZeroLibAllCodexDen}{152}
\rqSet{rqFPSeedZeroLibAllCodexPct}{2.0}
\rqSet{rqFPSeedZeroLibAllCTwoRustNum}{3}
\rqSet{rqFPSeedZeroLibAllCTwoRustDen}{125}
\rqSet{rqFPSeedZeroLibAllCTwoRustPct}{2.4}
\rqSet{rqFPSeedZeroLibAllLaertesNum}{3}
\rqSet{rqFPSeedZeroLibAllLaertesDen}{122}
\rqSet{rqFPSeedZeroLibAllLaertesPct}{2.5}
\rqSet{rqFPSeedZeroLibAllCTwoSaferRustNum}{7}
\rqSet{rqFPSeedZeroLibAllCTwoSaferRustDen}{93}
\rqSet{rqFPSeedZeroLibAllCTwoSaferRustPct}{7.5}
\rqSet{rqFPSeedZeroLibAllSmartCTwoRustNum}{0}
\rqSet{rqFPSeedZeroLibAllSmartCTwoRustDen}{0}
\rqSet{rqFPSeedZeroLibAllSmartCTwoRustPct}{n/a}
\rqSet{rqFPSeedZeroLibAllKimiNum}{14}
\rqSet{rqFPSeedZeroLibAllKimiDen}{78}
\rqSet{rqFPSeedZeroLibAllKimiPct}{17.9}
\rqSet{rqFPSeedZeroLibAllGptFiveNum}{14}
\rqSet{rqFPSeedZeroLibAllGptFiveDen}{99}
\rqSet{rqFPSeedZeroLibAllGptFivePct}{14.1}
\rqSet{rqFPSeedZeroLibAllGeminiNum}{5}
\rqSet{rqFPSeedZeroLibAllGeminiDen}{96}
\rqSet{rqFPSeedZeroLibAllGeminiPct}{5.2}
\rqSet{rqFPSeedZeroLibBOneNum}{35}
\rqSet{rqFPSeedZeroLibBOneDen}{588}
\rqSet{rqFPSeedZeroLibBOnePct}{6.0}
\rqSet{rqFPSeedZeroLibBTwoNum}{27}
\rqSet{rqFPSeedZeroLibBTwoDen}{399}
\rqSet{rqFPSeedZeroLibBTwoPct}{6.8}
\rqSet{rqFPSeedZeroLibPZeroNum}{0}
\rqSet{rqFPSeedZeroLibPZeroDen}{0}
\rqSet{rqFPSeedZeroLibPZeroPct}{n/a}
\rqSet{rqFPSeedZeroLibPOneNum}{0}
\rqSet{rqFPSeedZeroLibPOneDen}{169}
\rqSet{rqFPSeedZeroLibPOnePct}{0.0}
\rqSet{rqSeedMissingZeroLibAllNum}{0}
\rqSet{rqSeedMissingZeroLibAllDen}{1264}
\rqSet{rqSeedMissingZeroLibAllPct}{0.0}
\rqSet{rqSeedMissingZeroLibBOneNum}{0}
\rqSet{rqSeedMissingZeroLibBOneDen}{636}
\rqSet{rqSeedMissingZeroLibBOnePct}{0.0}
\rqSet{rqSeedMissingZeroLibBTwoNum}{0}
\rqSet{rqSeedMissingZeroLibBTwoDen}{459}
\rqSet{rqSeedMissingZeroLibBTwoPct}{0.0}
\rqSet{rqSeedMissingZeroLibPZeroNum}{0}
\rqSet{rqSeedMissingZeroLibPZeroDen}{0}
\rqSet{rqSeedMissingZeroLibPZeroPct}{n/a}
\rqSet{rqSeedMissingZeroLibPOneNum}{0}
\rqSet{rqSeedMissingZeroLibPOneDen}{169}
\rqSet{rqSeedMissingZeroLibPOnePct}{0.0}
\rqSet{rqFPSeedOneLibAllNum}{63}
\rqSet{rqFPSeedOneLibAllDen}{1111}
\rqSet{rqFPSeedOneLibAllPct}{5.7}
\rqSet{rqFPSeedOneLibAllKiroNum}{7}
\rqSet{rqFPSeedOneLibAllKiroDen}{195}
\rqSet{rqFPSeedOneLibAllKiroPct}{3.6}
\rqSet{rqFPSeedOneLibAllClaudeNum}{6}
\rqSet{rqFPSeedOneLibAllClaudeDen}{184}
\rqSet{rqFPSeedOneLibAllClaudePct}{3.3}
\rqSet{rqFPSeedOneLibAllCodexNum}{2}
\rqSet{rqFPSeedOneLibAllCodexDen}{146}
\rqSet{rqFPSeedOneLibAllCodexPct}{1.4}
\rqSet{rqFPSeedOneLibAllCTwoRustNum}{3}
\rqSet{rqFPSeedOneLibAllCTwoRustDen}{119}
\rqSet{rqFPSeedOneLibAllCTwoRustPct}{2.5}
\rqSet{rqFPSeedOneLibAllLaertesNum}{3}
\rqSet{rqFPSeedOneLibAllLaertesDen}{116}
\rqSet{rqFPSeedOneLibAllLaertesPct}{2.6}
\rqSet{rqFPSeedOneLibAllCTwoSaferRustNum}{7}
\rqSet{rqFPSeedOneLibAllCTwoSaferRustDen}{89}
\rqSet{rqFPSeedOneLibAllCTwoSaferRustPct}{7.9}
\rqSet{rqFPSeedOneLibAllSmartCTwoRustNum}{0}
\rqSet{rqFPSeedOneLibAllSmartCTwoRustDen}{0}
\rqSet{rqFPSeedOneLibAllSmartCTwoRustPct}{n/a}
\rqSet{rqFPSeedOneLibAllKimiNum}{15}
\rqSet{rqFPSeedOneLibAllKimiDen}{75}
\rqSet{rqFPSeedOneLibAllKimiPct}{20.0}
\rqSet{rqFPSeedOneLibAllGptFiveNum}{15}
\rqSet{rqFPSeedOneLibAllGptFiveDen}{94}
\rqSet{rqFPSeedOneLibAllGptFivePct}{16.0}
\rqSet{rqFPSeedOneLibAllGeminiNum}{5}
\rqSet{rqFPSeedOneLibAllGeminiDen}{93}
\rqSet{rqFPSeedOneLibAllGeminiPct}{5.4}
\rqSet{rqFPSeedOneLibBOneNum}{38}
\rqSet{rqFPSeedOneLibBOneDen}{563}
\rqSet{rqFPSeedOneLibBOnePct}{6.7}
\rqSet{rqFPSeedOneLibBTwoNum}{25}
\rqSet{rqFPSeedOneLibBTwoDen}{379}
\rqSet{rqFPSeedOneLibBTwoPct}{6.6}
\rqSet{rqFPSeedOneLibPZeroNum}{0}
\rqSet{rqFPSeedOneLibPZeroDen}{0}
\rqSet{rqFPSeedOneLibPZeroPct}{n/a}
\rqSet{rqFPSeedOneLibPOneNum}{0}
\rqSet{rqFPSeedOneLibPOneDen}{169}
\rqSet{rqFPSeedOneLibPOnePct}{0.0}
\rqSet{rqSeedMissingOneLibAllNum}{45}
\rqSet{rqSeedMissingOneLibAllDen}{1264}
\rqSet{rqSeedMissingOneLibAllPct}{3.6}
\rqSet{rqSeedMissingOneLibBOneNum}{25}
\rqSet{rqSeedMissingOneLibBOneDen}{636}
\rqSet{rqSeedMissingOneLibBOnePct}{3.9}
\rqSet{rqSeedMissingOneLibBTwoNum}{20}
\rqSet{rqSeedMissingOneLibBTwoDen}{459}
\rqSet{rqSeedMissingOneLibBTwoPct}{4.4}
\rqSet{rqSeedMissingOneLibPZeroNum}{0}
\rqSet{rqSeedMissingOneLibPZeroDen}{0}
\rqSet{rqSeedMissingOneLibPZeroPct}{n/a}
\rqSet{rqSeedMissingOneLibPOneNum}{0}
\rqSet{rqSeedMissingOneLibPOneDen}{169}
\rqSet{rqSeedMissingOneLibPOnePct}{0.0}
\rqSet{rqFPSeedTwoLibAllNum}{62}
\rqSet{rqFPSeedTwoLibAllDen}{1095}
\rqSet{rqFPSeedTwoLibAllPct}{5.7}
\rqSet{rqFPSeedTwoLibAllKiroNum}{7}
\rqSet{rqFPSeedTwoLibAllKiroDen}{193}
\rqSet{rqFPSeedTwoLibAllKiroPct}{3.6}
\rqSet{rqFPSeedTwoLibAllClaudeNum}{6}
\rqSet{rqFPSeedTwoLibAllClaudeDen}{182}
\rqSet{rqFPSeedTwoLibAllClaudePct}{3.3}
\rqSet{rqFPSeedTwoLibAllCodexNum}{2}
\rqSet{rqFPSeedTwoLibAllCodexDen}{144}
\rqSet{rqFPSeedTwoLibAllCodexPct}{1.4}
\rqSet{rqFPSeedTwoLibAllCTwoRustNum}{3}
\rqSet{rqFPSeedTwoLibAllCTwoRustDen}{117}
\rqSet{rqFPSeedTwoLibAllCTwoRustPct}{2.6}
\rqSet{rqFPSeedTwoLibAllLaertesNum}{3}
\rqSet{rqFPSeedTwoLibAllLaertesDen}{115}
\rqSet{rqFPSeedTwoLibAllLaertesPct}{2.6}
\rqSet{rqFPSeedTwoLibAllCTwoSaferRustNum}{7}
\rqSet{rqFPSeedTwoLibAllCTwoSaferRustDen}{88}
\rqSet{rqFPSeedTwoLibAllCTwoSaferRustPct}{8.0}
\rqSet{rqFPSeedTwoLibAllSmartCTwoRustNum}{0}
\rqSet{rqFPSeedTwoLibAllSmartCTwoRustDen}{0}
\rqSet{rqFPSeedTwoLibAllSmartCTwoRustPct}{n/a}
\rqSet{rqFPSeedTwoLibAllKimiNum}{14}
\rqSet{rqFPSeedTwoLibAllKimiDen}{73}
\rqSet{rqFPSeedTwoLibAllKimiPct}{19.2}
\rqSet{rqFPSeedTwoLibAllGptFiveNum}{14}
\rqSet{rqFPSeedTwoLibAllGptFiveDen}{92}
\rqSet{rqFPSeedTwoLibAllGptFivePct}{15.2}
\rqSet{rqFPSeedTwoLibAllGeminiNum}{6}
\rqSet{rqFPSeedTwoLibAllGeminiDen}{91}
\rqSet{rqFPSeedTwoLibAllGeminiPct}{6.6}
\rqSet{rqFPSeedTwoLibBOneNum}{35}
\rqSet{rqFPSeedTwoLibBOneDen}{540}
\rqSet{rqFPSeedTwoLibBOnePct}{6.5}
\rqSet{rqFPSeedTwoLibBTwoNum}{27}
\rqSet{rqFPSeedTwoLibBTwoDen}{386}
\rqSet{rqFPSeedTwoLibBTwoPct}{7.0}
\rqSet{rqFPSeedTwoLibPZeroNum}{0}
\rqSet{rqFPSeedTwoLibPZeroDen}{0}
\rqSet{rqFPSeedTwoLibPZeroPct}{n/a}
\rqSet{rqFPSeedTwoLibPOneNum}{0}
\rqSet{rqFPSeedTwoLibPOneDen}{169}
\rqSet{rqFPSeedTwoLibPOnePct}{0.0}
\rqSet{rqSeedMissingTwoLibAllNum}{61}
\rqSet{rqSeedMissingTwoLibAllDen}{1264}
\rqSet{rqSeedMissingTwoLibAllPct}{4.8}
\rqSet{rqSeedMissingTwoLibBOneNum}{48}
\rqSet{rqSeedMissingTwoLibBOneDen}{636}
\rqSet{rqSeedMissingTwoLibBOnePct}{7.5}
\rqSet{rqSeedMissingTwoLibBTwoNum}{13}
\rqSet{rqSeedMissingTwoLibBTwoDen}{459}
\rqSet{rqSeedMissingTwoLibBTwoPct}{2.8}
\rqSet{rqSeedMissingTwoLibPZeroNum}{0}
\rqSet{rqSeedMissingTwoLibPZeroDen}{0}
\rqSet{rqSeedMissingTwoLibPZeroPct}{n/a}
\rqSet{rqSeedMissingTwoLibPOneNum}{0}
\rqSet{rqSeedMissingTwoLibPOneDen}{169}
\rqSet{rqSeedMissingTwoLibPOnePct}{0.0}
\rqSet{rqUnstableOnlyLibAllNum}{91}
\rqSet{rqUnstableOnlyLibAllDen}{1088}
\rqSet{rqUnstableOnlyLibAllPct}{8.4}
\rqSet{rqUnstableOnlyLibAllKiroNum}{13}
\rqSet{rqUnstableOnlyLibAllKiroDen}{192}
\rqSet{rqUnstableOnlyLibAllKiroPct}{6.8}
\rqSet{rqUnstableOnlyLibAllClaudeNum}{12}
\rqSet{rqUnstableOnlyLibAllClaudeDen}{181}
\rqSet{rqUnstableOnlyLibAllClaudePct}{6.6}
\rqSet{rqUnstableOnlyLibAllCodexNum}{13}
\rqSet{rqUnstableOnlyLibAllCodexDen}{143}
\rqSet{rqUnstableOnlyLibAllCodexPct}{9.1}
\rqSet{rqUnstableOnlyLibAllCTwoRustNum}{13}
\rqSet{rqUnstableOnlyLibAllCTwoRustDen}{116}
\rqSet{rqUnstableOnlyLibAllCTwoRustPct}{11.2}
\rqSet{rqUnstableOnlyLibAllLaertesNum}{10}
\rqSet{rqUnstableOnlyLibAllLaertesDen}{114}
\rqSet{rqUnstableOnlyLibAllLaertesPct}{8.8}
\rqSet{rqUnstableOnlyLibAllCTwoSaferRustNum}{8}
\rqSet{rqUnstableOnlyLibAllCTwoSaferRustDen}{87}
\rqSet{rqUnstableOnlyLibAllCTwoSaferRustPct}{9.2}
\rqSet{rqUnstableOnlyLibAllSmartCTwoRustNum}{0}
\rqSet{rqUnstableOnlyLibAllSmartCTwoRustDen}{0}
\rqSet{rqUnstableOnlyLibAllSmartCTwoRustPct}{n/a}
\rqSet{rqUnstableOnlyLibAllKimiNum}{5}
\rqSet{rqUnstableOnlyLibAllKimiDen}{73}
\rqSet{rqUnstableOnlyLibAllKimiPct}{6.8}
\rqSet{rqUnstableOnlyLibAllGptFiveNum}{8}
\rqSet{rqUnstableOnlyLibAllGptFiveDen}{92}
\rqSet{rqUnstableOnlyLibAllGptFivePct}{8.7}
\rqSet{rqUnstableOnlyLibAllGeminiNum}{9}
\rqSet{rqUnstableOnlyLibAllGeminiDen}{90}
\rqSet{rqUnstableOnlyLibAllGeminiPct}{10.0}
\rqSet{rqUnstableOnlyLibBOneNum}{59}
\rqSet{rqUnstableOnlyLibBOneDen}{540}
\rqSet{rqUnstableOnlyLibBOnePct}{10.9}
\rqSet{rqUnstableOnlyLibBTwoNum}{32}
\rqSet{rqUnstableOnlyLibBTwoDen}{379}
\rqSet{rqUnstableOnlyLibBTwoPct}{8.4}
\rqSet{rqUnstableOnlyLibPZeroNum}{0}
\rqSet{rqUnstableOnlyLibPZeroDen}{0}
\rqSet{rqUnstableOnlyLibPZeroPct}{n/a}
\rqSet{rqUnstableOnlyLibPOneNum}{0}
\rqSet{rqUnstableOnlyLibPOneDen}{169}
\rqSet{rqUnstableOnlyLibPOnePct}{0.0}
\rqSet{rqInconclOnlyLibAllNum}{8}
\rqSet{rqInconclOnlyLibAllDen}{1088}
\rqSet{rqInconclOnlyLibAllPct}{0.7}
\rqSet{rqInconclOnlyLibAllKiroNum}{1}
\rqSet{rqInconclOnlyLibAllKiroDen}{192}
\rqSet{rqInconclOnlyLibAllKiroPct}{0.5}
\rqSet{rqInconclOnlyLibAllClaudeNum}{1}
\rqSet{rqInconclOnlyLibAllClaudeDen}{181}
\rqSet{rqInconclOnlyLibAllClaudePct}{0.6}
\rqSet{rqInconclOnlyLibAllCodexNum}{0}
\rqSet{rqInconclOnlyLibAllCodexDen}{143}
\rqSet{rqInconclOnlyLibAllCodexPct}{0.0}
\rqSet{rqInconclOnlyLibAllCTwoRustNum}{0}
\rqSet{rqInconclOnlyLibAllCTwoRustDen}{116}
\rqSet{rqInconclOnlyLibAllCTwoRustPct}{0.0}
\rqSet{rqInconclOnlyLibAllLaertesNum}{0}
\rqSet{rqInconclOnlyLibAllLaertesDen}{114}
\rqSet{rqInconclOnlyLibAllLaertesPct}{0.0}
\rqSet{rqInconclOnlyLibAllCTwoSaferRustNum}{2}
\rqSet{rqInconclOnlyLibAllCTwoSaferRustDen}{87}
\rqSet{rqInconclOnlyLibAllCTwoSaferRustPct}{2.3}
\rqSet{rqInconclOnlyLibAllSmartCTwoRustNum}{0}
\rqSet{rqInconclOnlyLibAllSmartCTwoRustDen}{0}
\rqSet{rqInconclOnlyLibAllSmartCTwoRustPct}{n/a}
\rqSet{rqInconclOnlyLibAllKimiNum}{1}
\rqSet{rqInconclOnlyLibAllKimiDen}{73}
\rqSet{rqInconclOnlyLibAllKimiPct}{1.4}
\rqSet{rqInconclOnlyLibAllGptFiveNum}{2}
\rqSet{rqInconclOnlyLibAllGptFiveDen}{92}
\rqSet{rqInconclOnlyLibAllGptFivePct}{2.2}
\rqSet{rqInconclOnlyLibAllGeminiNum}{1}
\rqSet{rqInconclOnlyLibAllGeminiDen}{90}
\rqSet{rqInconclOnlyLibAllGeminiPct}{1.1}
\rqSet{rqInconclOnlyLibBOneNum}{0}
\rqSet{rqInconclOnlyLibBOneDen}{540}
\rqSet{rqInconclOnlyLibBOnePct}{0.0}
\rqSet{rqInconclOnlyLibBTwoNum}{8}
\rqSet{rqInconclOnlyLibBTwoDen}{379}
\rqSet{rqInconclOnlyLibBTwoPct}{2.1}
\rqSet{rqInconclOnlyLibPZeroNum}{0}
\rqSet{rqInconclOnlyLibPZeroDen}{0}
\rqSet{rqInconclOnlyLibPZeroPct}{n/a}
\rqSet{rqInconclOnlyLibPOneNum}{0}
\rqSet{rqInconclOnlyLibPOneDen}{169}
\rqSet{rqInconclOnlyLibPOnePct}{0.0}
\rqSet{rqInconclUnstableLibAllNum}{0}
\rqSet{rqInconclUnstableLibAllDen}{1088}
\rqSet{rqInconclUnstableLibAllPct}{0.0}
\rqSet{rqInconclUnstableLibAllKiroNum}{0}
\rqSet{rqInconclUnstableLibAllKiroDen}{192}
\rqSet{rqInconclUnstableLibAllKiroPct}{0.0}
\rqSet{rqInconclUnstableLibAllClaudeNum}{0}
\rqSet{rqInconclUnstableLibAllClaudeDen}{181}
\rqSet{rqInconclUnstableLibAllClaudePct}{0.0}
\rqSet{rqInconclUnstableLibAllCodexNum}{0}
\rqSet{rqInconclUnstableLibAllCodexDen}{143}
\rqSet{rqInconclUnstableLibAllCodexPct}{0.0}
\rqSet{rqInconclUnstableLibAllCTwoRustNum}{0}
\rqSet{rqInconclUnstableLibAllCTwoRustDen}{116}
\rqSet{rqInconclUnstableLibAllCTwoRustPct}{0.0}
\rqSet{rqInconclUnstableLibAllLaertesNum}{0}
\rqSet{rqInconclUnstableLibAllLaertesDen}{114}
\rqSet{rqInconclUnstableLibAllLaertesPct}{0.0}
\rqSet{rqInconclUnstableLibAllCTwoSaferRustNum}{0}
\rqSet{rqInconclUnstableLibAllCTwoSaferRustDen}{87}
\rqSet{rqInconclUnstableLibAllCTwoSaferRustPct}{0.0}
\rqSet{rqInconclUnstableLibAllSmartCTwoRustNum}{0}
\rqSet{rqInconclUnstableLibAllSmartCTwoRustDen}{0}
\rqSet{rqInconclUnstableLibAllSmartCTwoRustPct}{n/a}
\rqSet{rqInconclUnstableLibAllKimiNum}{0}
\rqSet{rqInconclUnstableLibAllKimiDen}{73}
\rqSet{rqInconclUnstableLibAllKimiPct}{0.0}
\rqSet{rqInconclUnstableLibAllGptFiveNum}{0}
\rqSet{rqInconclUnstableLibAllGptFiveDen}{92}
\rqSet{rqInconclUnstableLibAllGptFivePct}{0.0}
\rqSet{rqInconclUnstableLibAllGeminiNum}{0}
\rqSet{rqInconclUnstableLibAllGeminiDen}{90}
\rqSet{rqInconclUnstableLibAllGeminiPct}{0.0}
\rqSet{rqInconclUnstableLibBOneNum}{0}
\rqSet{rqInconclUnstableLibBOneDen}{540}
\rqSet{rqInconclUnstableLibBOnePct}{0.0}
\rqSet{rqInconclUnstableLibBTwoNum}{0}
\rqSet{rqInconclUnstableLibBTwoDen}{379}
\rqSet{rqInconclUnstableLibBTwoPct}{0.0}
\rqSet{rqInconclUnstableLibPZeroNum}{0}
\rqSet{rqInconclUnstableLibPZeroDen}{0}
\rqSet{rqInconclUnstableLibPZeroPct}{n/a}
\rqSet{rqInconclUnstableLibPOneNum}{0}
\rqSet{rqInconclUnstableLibPOneDen}{169}
\rqSet{rqInconclUnstableLibPOnePct}{0.0}
\rqSet{rqInconclAnyLibAllNum}{12}
\rqSet{rqInconclAnyLibAllDen}{1088}
\rqSet{rqInconclAnyLibAllPct}{1.1}
\rqSet{rqInconclAnyLibBOneNum}{0}
\rqSet{rqInconclAnyLibBOneDen}{540}
\rqSet{rqInconclAnyLibBOnePct}{0.0}
\rqSet{rqInconclAnyLibBTwoNum}{12}
\rqSet{rqInconclAnyLibBTwoDen}{379}
\rqSet{rqInconclAnyLibBTwoPct}{3.2}
\rqSet{rqInconclAnyLibPZeroNum}{0}
\rqSet{rqInconclAnyLibPZeroDen}{0}
\rqSet{rqInconclAnyLibPZeroPct}{n/a}
\rqSet{rqInconclAnyLibPOneNum}{0}
\rqSet{rqInconclAnyLibPOneDen}{169}
\rqSet{rqInconclAnyLibPOnePct}{0.0}
\rqSet{rqNBLoLibAllNum}{56}
\rqSet{rqNBLoLibAllDen}{1088}
\rqSet{rqNBLoLibAllPct}{5.1}
\rqSet{rqNBLoLibAllKiroNum}{6}
\rqSet{rqNBLoLibAllKiroDen}{192}
\rqSet{rqNBLoLibAllKiroPct}{3.1}
\rqSet{rqNBLoLibAllClaudeNum}{5}
\rqSet{rqNBLoLibAllClaudeDen}{181}
\rqSet{rqNBLoLibAllClaudePct}{2.8}
\rqSet{rqNBLoLibAllCodexNum}{1}
\rqSet{rqNBLoLibAllCodexDen}{143}
\rqSet{rqNBLoLibAllCodexPct}{0.7}
\rqSet{rqNBLoLibAllCTwoRustNum}{2}
\rqSet{rqNBLoLibAllCTwoRustDen}{116}
\rqSet{rqNBLoLibAllCTwoRustPct}{1.7}
\rqSet{rqNBLoLibAllLaertesNum}{2}
\rqSet{rqNBLoLibAllLaertesDen}{114}
\rqSet{rqNBLoLibAllLaertesPct}{1.8}
\rqSet{rqNBLoLibAllCTwoSaferRustNum}{6}
\rqSet{rqNBLoLibAllCTwoSaferRustDen}{87}
\rqSet{rqNBLoLibAllCTwoSaferRustPct}{6.9}
\rqSet{rqNBLoLibAllSmartCTwoRustNum}{0}
\rqSet{rqNBLoLibAllSmartCTwoRustDen}{0}
\rqSet{rqNBLoLibAllSmartCTwoRustPct}{n/a}
\rqSet{rqNBLoLibAllKimiNum}{15}
\rqSet{rqNBLoLibAllKimiDen}{73}
\rqSet{rqNBLoLibAllKimiPct}{20.5}
\rqSet{rqNBLoLibAllGptFiveNum}{14}
\rqSet{rqNBLoLibAllGptFiveDen}{92}
\rqSet{rqNBLoLibAllGptFivePct}{15.2}
\rqSet{rqNBLoLibAllGeminiNum}{5}
\rqSet{rqNBLoLibAllGeminiDen}{90}
\rqSet{rqNBLoLibAllGeminiPct}{5.6}
\rqSet{rqNBLoLibBOneNum}{30}
\rqSet{rqNBLoLibBOneDen}{540}
\rqSet{rqNBLoLibBOnePct}{5.6}
\rqSet{rqNBLoLibBOneKiroNum}{1}
\rqSet{rqNBLoLibBOneKiroDen}{65}
\rqSet{rqNBLoLibBOneKiroPct}{1.5}
\rqSet{rqNBLoLibBOneClaudeNum}{1}
\rqSet{rqNBLoLibBOneClaudeDen}{65}
\rqSet{rqNBLoLibBOneClaudePct}{1.5}
\rqSet{rqNBLoLibBOneCodexNum}{1}
\rqSet{rqNBLoLibBOneCodexDen}{65}
\rqSet{rqNBLoLibBOneCodexPct}{1.5}
\rqSet{rqNBLoLibBOneCTwoRustNum}{2}
\rqSet{rqNBLoLibBOneCTwoRustDen}{65}
\rqSet{rqNBLoLibBOneCTwoRustPct}{3.1}
\rqSet{rqNBLoLibBOneLaertesNum}{2}
\rqSet{rqNBLoLibBOneLaertesDen}{63}
\rqSet{rqNBLoLibBOneLaertesPct}{3.2}
\rqSet{rqNBLoLibBOneCTwoSaferRustNum}{5}
\rqSet{rqNBLoLibBOneCTwoSaferRustDen}{50}
\rqSet{rqNBLoLibBOneCTwoSaferRustPct}{10.0}
\rqSet{rqNBLoLibBOneSmartCTwoRustNum}{0}
\rqSet{rqNBLoLibBOneSmartCTwoRustDen}{0}
\rqSet{rqNBLoLibBOneSmartCTwoRustPct}{n/a}
\rqSet{rqNBLoLibBOneKimiNum}{8}
\rqSet{rqNBLoLibBOneKimiDen}{49}
\rqSet{rqNBLoLibBOneKimiPct}{16.3}
\rqSet{rqNBLoLibBOneGptFiveNum}{8}
\rqSet{rqNBLoLibBOneGptFiveDen}{57}
\rqSet{rqNBLoLibBOneGptFivePct}{14.0}
\rqSet{rqNBLoLibBOneGeminiNum}{2}
\rqSet{rqNBLoLibBOneGeminiDen}{61}
\rqSet{rqNBLoLibBOneGeminiPct}{3.3}
\rqSet{rqNBLoLibBTwoNum}{26}
\rqSet{rqNBLoLibBTwoDen}{379}
\rqSet{rqNBLoLibBTwoPct}{6.9}
\rqSet{rqNBLoLibBTwoKiroNum}{5}
\rqSet{rqNBLoLibBTwoKiroDen}{51}
\rqSet{rqNBLoLibBTwoKiroPct}{9.8}
\rqSet{rqNBLoLibBTwoClaudeNum}{4}
\rqSet{rqNBLoLibBTwoClaudeDen}{51}
\rqSet{rqNBLoLibBTwoClaudePct}{7.8}
\rqSet{rqNBLoLibBTwoCodexNum}{0}
\rqSet{rqNBLoLibBTwoCodexDen}{50}
\rqSet{rqNBLoLibBTwoCodexPct}{0.0}
\rqSet{rqNBLoLibBTwoCTwoRustNum}{0}
\rqSet{rqNBLoLibBTwoCTwoRustDen}{51}
\rqSet{rqNBLoLibBTwoCTwoRustPct}{0.0}
\rqSet{rqNBLoLibBTwoLaertesNum}{0}
\rqSet{rqNBLoLibBTwoLaertesDen}{51}
\rqSet{rqNBLoLibBTwoLaertesPct}{0.0}
\rqSet{rqNBLoLibBTwoCTwoSaferRustNum}{1}
\rqSet{rqNBLoLibBTwoCTwoSaferRustDen}{37}
\rqSet{rqNBLoLibBTwoCTwoSaferRustPct}{2.7}
\rqSet{rqNBLoLibBTwoSmartCTwoRustNum}{0}
\rqSet{rqNBLoLibBTwoSmartCTwoRustDen}{0}
\rqSet{rqNBLoLibBTwoSmartCTwoRustPct}{n/a}
\rqSet{rqNBLoLibBTwoKimiNum}{7}
\rqSet{rqNBLoLibBTwoKimiDen}{24}
\rqSet{rqNBLoLibBTwoKimiPct}{29.2}
\rqSet{rqNBLoLibBTwoGptFiveNum}{6}
\rqSet{rqNBLoLibBTwoGptFiveDen}{35}
\rqSet{rqNBLoLibBTwoGptFivePct}{17.1}
\rqSet{rqNBLoLibBTwoGeminiNum}{3}
\rqSet{rqNBLoLibBTwoGeminiDen}{29}
\rqSet{rqNBLoLibBTwoGeminiPct}{10.3}
\rqSet{rqNBLoLibPZeroNum}{0}
\rqSet{rqNBLoLibPZeroDen}{0}
\rqSet{rqNBLoLibPZeroPct}{n/a}
\rqSet{rqNBLoLibPZeroKiroNum}{0}
\rqSet{rqNBLoLibPZeroKiroDen}{0}
\rqSet{rqNBLoLibPZeroKiroPct}{n/a}
\rqSet{rqNBLoLibPZeroClaudeNum}{0}
\rqSet{rqNBLoLibPZeroClaudeDen}{0}
\rqSet{rqNBLoLibPZeroClaudePct}{n/a}
\rqSet{rqNBLoLibPZeroCodexNum}{0}
\rqSet{rqNBLoLibPZeroCodexDen}{0}
\rqSet{rqNBLoLibPZeroCodexPct}{n/a}
\rqSet{rqNBLoLibPZeroCTwoRustNum}{0}
\rqSet{rqNBLoLibPZeroCTwoRustDen}{0}
\rqSet{rqNBLoLibPZeroCTwoRustPct}{n/a}
\rqSet{rqNBLoLibPZeroLaertesNum}{0}
\rqSet{rqNBLoLibPZeroLaertesDen}{0}
\rqSet{rqNBLoLibPZeroLaertesPct}{n/a}
\rqSet{rqNBLoLibPZeroCTwoSaferRustNum}{0}
\rqSet{rqNBLoLibPZeroCTwoSaferRustDen}{0}
\rqSet{rqNBLoLibPZeroCTwoSaferRustPct}{n/a}
\rqSet{rqNBLoLibPZeroSmartCTwoRustNum}{0}
\rqSet{rqNBLoLibPZeroSmartCTwoRustDen}{0}
\rqSet{rqNBLoLibPZeroSmartCTwoRustPct}{n/a}
\rqSet{rqNBLoLibPZeroKimiNum}{0}
\rqSet{rqNBLoLibPZeroKimiDen}{0}
\rqSet{rqNBLoLibPZeroKimiPct}{n/a}
\rqSet{rqNBLoLibPZeroGptFiveNum}{0}
\rqSet{rqNBLoLibPZeroGptFiveDen}{0}
\rqSet{rqNBLoLibPZeroGptFivePct}{n/a}
\rqSet{rqNBLoLibPZeroGeminiNum}{0}
\rqSet{rqNBLoLibPZeroGeminiDen}{0}
\rqSet{rqNBLoLibPZeroGeminiPct}{n/a}
\rqSet{rqNBLoLibPOneNum}{0}
\rqSet{rqNBLoLibPOneDen}{169}
\rqSet{rqNBLoLibPOnePct}{0.0}
\rqSet{rqNBLoLibPOneKiroNum}{0}
\rqSet{rqNBLoLibPOneKiroDen}{76}
\rqSet{rqNBLoLibPOneKiroPct}{0.0}
\rqSet{rqNBLoLibPOneClaudeNum}{0}
\rqSet{rqNBLoLibPOneClaudeDen}{65}
\rqSet{rqNBLoLibPOneClaudePct}{0.0}
\rqSet{rqNBLoLibPOneCodexNum}{0}
\rqSet{rqNBLoLibPOneCodexDen}{28}
\rqSet{rqNBLoLibPOneCodexPct}{0.0}
\rqSet{rqNBLoLibPOneCTwoRustNum}{0}
\rqSet{rqNBLoLibPOneCTwoRustDen}{0}
\rqSet{rqNBLoLibPOneCTwoRustPct}{n/a}
\rqSet{rqNBLoLibPOneLaertesNum}{0}
\rqSet{rqNBLoLibPOneLaertesDen}{0}
\rqSet{rqNBLoLibPOneLaertesPct}{n/a}
\rqSet{rqNBLoLibPOneCTwoSaferRustNum}{0}
\rqSet{rqNBLoLibPOneCTwoSaferRustDen}{0}
\rqSet{rqNBLoLibPOneCTwoSaferRustPct}{n/a}
\rqSet{rqNBLoLibPOneSmartCTwoRustNum}{0}
\rqSet{rqNBLoLibPOneSmartCTwoRustDen}{0}
\rqSet{rqNBLoLibPOneSmartCTwoRustPct}{n/a}
\rqSet{rqNBLoLibPOneKimiNum}{0}
\rqSet{rqNBLoLibPOneKimiDen}{0}
\rqSet{rqNBLoLibPOneKimiPct}{n/a}
\rqSet{rqNBLoLibPOneGptFiveNum}{0}
\rqSet{rqNBLoLibPOneGptFiveDen}{0}
\rqSet{rqNBLoLibPOneGptFivePct}{n/a}
\rqSet{rqNBLoLibPOneGeminiNum}{0}
\rqSet{rqNBLoLibPOneGeminiDen}{0}
\rqSet{rqNBLoLibPOneGeminiPct}{n/a}
\rqSet{rqNBHiLibAllNum}{59}
\rqSet{rqNBHiLibAllDen}{1088}
\rqSet{rqNBHiLibAllPct}{5.4}
\rqSet{rqNBHiLibAllKiroNum}{7}
\rqSet{rqNBHiLibAllKiroDen}{192}
\rqSet{rqNBHiLibAllKiroPct}{3.6}
\rqSet{rqNBHiLibAllClaudeNum}{5}
\rqSet{rqNBHiLibAllClaudeDen}{181}
\rqSet{rqNBHiLibAllClaudePct}{2.8}
\rqSet{rqNBHiLibAllCodexNum}{2}
\rqSet{rqNBHiLibAllCodexDen}{143}
\rqSet{rqNBHiLibAllCodexPct}{1.4}
\rqSet{rqNBHiLibAllCTwoRustNum}{2}
\rqSet{rqNBHiLibAllCTwoRustDen}{116}
\rqSet{rqNBHiLibAllCTwoRustPct}{1.7}
\rqSet{rqNBHiLibAllLaertesNum}{2}
\rqSet{rqNBHiLibAllLaertesDen}{114}
\rqSet{rqNBHiLibAllLaertesPct}{1.8}
\rqSet{rqNBHiLibAllCTwoSaferRustNum}{7}
\rqSet{rqNBHiLibAllCTwoSaferRustDen}{87}
\rqSet{rqNBHiLibAllCTwoSaferRustPct}{8.0}
\rqSet{rqNBHiLibAllSmartCTwoRustNum}{0}
\rqSet{rqNBHiLibAllSmartCTwoRustDen}{0}
\rqSet{rqNBHiLibAllSmartCTwoRustPct}{n/a}
\rqSet{rqNBHiLibAllKimiNum}{15}
\rqSet{rqNBHiLibAllKimiDen}{73}
\rqSet{rqNBHiLibAllKimiPct}{20.5}
\rqSet{rqNBHiLibAllGptFiveNum}{14}
\rqSet{rqNBHiLibAllGptFiveDen}{92}
\rqSet{rqNBHiLibAllGptFivePct}{15.2}
\rqSet{rqNBHiLibAllGeminiNum}{5}
\rqSet{rqNBHiLibAllGeminiDen}{90}
\rqSet{rqNBHiLibAllGeminiPct}{5.6}
\rqSet{rqNBHiLibBOneNum}{32}
\rqSet{rqNBHiLibBOneDen}{540}
\rqSet{rqNBHiLibBOnePct}{5.9}
\rqSet{rqNBHiLibBOneKiroNum}{2}
\rqSet{rqNBHiLibBOneKiroDen}{65}
\rqSet{rqNBHiLibBOneKiroPct}{3.1}
\rqSet{rqNBHiLibBOneClaudeNum}{1}
\rqSet{rqNBHiLibBOneClaudeDen}{65}
\rqSet{rqNBHiLibBOneClaudePct}{1.5}
\rqSet{rqNBHiLibBOneCodexNum}{2}
\rqSet{rqNBHiLibBOneCodexDen}{65}
\rqSet{rqNBHiLibBOneCodexPct}{3.1}
\rqSet{rqNBHiLibBOneCTwoRustNum}{2}
\rqSet{rqNBHiLibBOneCTwoRustDen}{65}
\rqSet{rqNBHiLibBOneCTwoRustPct}{3.1}
\rqSet{rqNBHiLibBOneLaertesNum}{2}
\rqSet{rqNBHiLibBOneLaertesDen}{63}
\rqSet{rqNBHiLibBOneLaertesPct}{3.2}
\rqSet{rqNBHiLibBOneCTwoSaferRustNum}{5}
\rqSet{rqNBHiLibBOneCTwoSaferRustDen}{50}
\rqSet{rqNBHiLibBOneCTwoSaferRustPct}{10.0}
\rqSet{rqNBHiLibBOneSmartCTwoRustNum}{0}
\rqSet{rqNBHiLibBOneSmartCTwoRustDen}{0}
\rqSet{rqNBHiLibBOneSmartCTwoRustPct}{n/a}
\rqSet{rqNBHiLibBOneKimiNum}{8}
\rqSet{rqNBHiLibBOneKimiDen}{49}
\rqSet{rqNBHiLibBOneKimiPct}{16.3}
\rqSet{rqNBHiLibBOneGptFiveNum}{8}
\rqSet{rqNBHiLibBOneGptFiveDen}{57}
\rqSet{rqNBHiLibBOneGptFivePct}{14.0}
\rqSet{rqNBHiLibBOneGeminiNum}{2}
\rqSet{rqNBHiLibBOneGeminiDen}{61}
\rqSet{rqNBHiLibBOneGeminiPct}{3.3}
\rqSet{rqNBHiLibBTwoNum}{27}
\rqSet{rqNBHiLibBTwoDen}{379}
\rqSet{rqNBHiLibBTwoPct}{7.1}
\rqSet{rqNBHiLibBTwoKiroNum}{5}
\rqSet{rqNBHiLibBTwoKiroDen}{51}
\rqSet{rqNBHiLibBTwoKiroPct}{9.8}
\rqSet{rqNBHiLibBTwoClaudeNum}{4}
\rqSet{rqNBHiLibBTwoClaudeDen}{51}
\rqSet{rqNBHiLibBTwoClaudePct}{7.8}
\rqSet{rqNBHiLibBTwoCodexNum}{0}
\rqSet{rqNBHiLibBTwoCodexDen}{50}
\rqSet{rqNBHiLibBTwoCodexPct}{0.0}
\rqSet{rqNBHiLibBTwoCTwoRustNum}{0}
\rqSet{rqNBHiLibBTwoCTwoRustDen}{51}
\rqSet{rqNBHiLibBTwoCTwoRustPct}{0.0}
\rqSet{rqNBHiLibBTwoLaertesNum}{0}
\rqSet{rqNBHiLibBTwoLaertesDen}{51}
\rqSet{rqNBHiLibBTwoLaertesPct}{0.0}
\rqSet{rqNBHiLibBTwoCTwoSaferRustNum}{2}
\rqSet{rqNBHiLibBTwoCTwoSaferRustDen}{37}
\rqSet{rqNBHiLibBTwoCTwoSaferRustPct}{5.4}
\rqSet{rqNBHiLibBTwoSmartCTwoRustNum}{0}
\rqSet{rqNBHiLibBTwoSmartCTwoRustDen}{0}
\rqSet{rqNBHiLibBTwoSmartCTwoRustPct}{n/a}
\rqSet{rqNBHiLibBTwoKimiNum}{7}
\rqSet{rqNBHiLibBTwoKimiDen}{24}
\rqSet{rqNBHiLibBTwoKimiPct}{29.2}
\rqSet{rqNBHiLibBTwoGptFiveNum}{6}
\rqSet{rqNBHiLibBTwoGptFiveDen}{35}
\rqSet{rqNBHiLibBTwoGptFivePct}{17.1}
\rqSet{rqNBHiLibBTwoGeminiNum}{3}
\rqSet{rqNBHiLibBTwoGeminiDen}{29}
\rqSet{rqNBHiLibBTwoGeminiPct}{10.3}
\rqSet{rqNBHiLibPZeroNum}{0}
\rqSet{rqNBHiLibPZeroDen}{0}
\rqSet{rqNBHiLibPZeroPct}{n/a}
\rqSet{rqNBHiLibPZeroKiroNum}{0}
\rqSet{rqNBHiLibPZeroKiroDen}{0}
\rqSet{rqNBHiLibPZeroKiroPct}{n/a}
\rqSet{rqNBHiLibPZeroClaudeNum}{0}
\rqSet{rqNBHiLibPZeroClaudeDen}{0}
\rqSet{rqNBHiLibPZeroClaudePct}{n/a}
\rqSet{rqNBHiLibPZeroCodexNum}{0}
\rqSet{rqNBHiLibPZeroCodexDen}{0}
\rqSet{rqNBHiLibPZeroCodexPct}{n/a}
\rqSet{rqNBHiLibPZeroCTwoRustNum}{0}
\rqSet{rqNBHiLibPZeroCTwoRustDen}{0}
\rqSet{rqNBHiLibPZeroCTwoRustPct}{n/a}
\rqSet{rqNBHiLibPZeroLaertesNum}{0}
\rqSet{rqNBHiLibPZeroLaertesDen}{0}
\rqSet{rqNBHiLibPZeroLaertesPct}{n/a}
\rqSet{rqNBHiLibPZeroCTwoSaferRustNum}{0}
\rqSet{rqNBHiLibPZeroCTwoSaferRustDen}{0}
\rqSet{rqNBHiLibPZeroCTwoSaferRustPct}{n/a}
\rqSet{rqNBHiLibPZeroSmartCTwoRustNum}{0}
\rqSet{rqNBHiLibPZeroSmartCTwoRustDen}{0}
\rqSet{rqNBHiLibPZeroSmartCTwoRustPct}{n/a}
\rqSet{rqNBHiLibPZeroKimiNum}{0}
\rqSet{rqNBHiLibPZeroKimiDen}{0}
\rqSet{rqNBHiLibPZeroKimiPct}{n/a}
\rqSet{rqNBHiLibPZeroGptFiveNum}{0}
\rqSet{rqNBHiLibPZeroGptFiveDen}{0}
\rqSet{rqNBHiLibPZeroGptFivePct}{n/a}
\rqSet{rqNBHiLibPZeroGeminiNum}{0}
\rqSet{rqNBHiLibPZeroGeminiDen}{0}
\rqSet{rqNBHiLibPZeroGeminiPct}{n/a}
\rqSet{rqNBHiLibPOneNum}{0}
\rqSet{rqNBHiLibPOneDen}{169}
\rqSet{rqNBHiLibPOnePct}{0.0}
\rqSet{rqNBHiLibPOneKiroNum}{0}
\rqSet{rqNBHiLibPOneKiroDen}{76}
\rqSet{rqNBHiLibPOneKiroPct}{0.0}
\rqSet{rqNBHiLibPOneClaudeNum}{0}
\rqSet{rqNBHiLibPOneClaudeDen}{65}
\rqSet{rqNBHiLibPOneClaudePct}{0.0}
\rqSet{rqNBHiLibPOneCodexNum}{0}
\rqSet{rqNBHiLibPOneCodexDen}{28}
\rqSet{rqNBHiLibPOneCodexPct}{0.0}
\rqSet{rqNBHiLibPOneCTwoRustNum}{0}
\rqSet{rqNBHiLibPOneCTwoRustDen}{0}
\rqSet{rqNBHiLibPOneCTwoRustPct}{n/a}
\rqSet{rqNBHiLibPOneLaertesNum}{0}
\rqSet{rqNBHiLibPOneLaertesDen}{0}
\rqSet{rqNBHiLibPOneLaertesPct}{n/a}
\rqSet{rqNBHiLibPOneCTwoSaferRustNum}{0}
\rqSet{rqNBHiLibPOneCTwoSaferRustDen}{0}
\rqSet{rqNBHiLibPOneCTwoSaferRustPct}{n/a}
\rqSet{rqNBHiLibPOneSmartCTwoRustNum}{0}
\rqSet{rqNBHiLibPOneSmartCTwoRustDen}{0}
\rqSet{rqNBHiLibPOneSmartCTwoRustPct}{n/a}
\rqSet{rqNBHiLibPOneKimiNum}{0}
\rqSet{rqNBHiLibPOneKimiDen}{0}
\rqSet{rqNBHiLibPOneKimiPct}{n/a}
\rqSet{rqNBHiLibPOneGptFiveNum}{0}
\rqSet{rqNBHiLibPOneGptFiveDen}{0}
\rqSet{rqNBHiLibPOneGptFivePct}{n/a}
\rqSet{rqNBHiLibPOneGeminiNum}{0}
\rqSet{rqNBHiLibPOneGeminiDen}{0}
\rqSet{rqNBHiLibPOneGeminiPct}{n/a}
\rqSet{rqBnuDivBOneWhitespaceLibAllNum}{18}
\rqSet{rqBnuDivBOneWhitespaceLibAllDen}{18194}
\rqSet{rqBnuDivBOneWhitespaceLibAllPct}{0.1}
\rqSet{rqBnuDivBOneWhitespaceLibBOneNum}{18}
\rqSet{rqBnuDivBOneWhitespaceLibBOneDen}{10507}
\rqSet{rqBnuDivBOneWhitespaceLibBOnePct}{0.2}
\rqSet{rqBnuDivBOneWhitespaceLibBTwoNum}{0}
\rqSet{rqBnuDivBOneWhitespaceLibBTwoDen}{7687}
\rqSet{rqBnuDivBOneWhitespaceLibBTwoPct}{0.0}
\rqSet{rqBnuDivBOneWhitespaceLibPZeroNum}{0}
\rqSet{rqBnuDivBOneWhitespaceLibPZeroDen}{0}
\rqSet{rqBnuDivBOneWhitespaceLibPZeroPct}{n/a}
\rqSet{rqBnuDivBOneWhitespaceLibPOneNum}{0}
\rqSet{rqBnuDivBOneWhitespaceLibPOneDen}{0}
\rqSet{rqBnuDivBOneWhitespaceLibPOnePct}{n/a}
\rqSet{rqBnuTransBOneWhitespaceLibAllNum}{6}
\rqSet{rqBnuTransBOneWhitespaceLibAllDen}{65}
\rqSet{rqBnuTransBOneWhitespaceLibAllPct}{9.2}
\rqSet{rqBnuTransBOneWhitespaceLibBOneNum}{6}
\rqSet{rqBnuTransBOneWhitespaceLibBOneDen}{38}
\rqSet{rqBnuTransBOneWhitespaceLibBOnePct}{15.8}
\rqSet{rqBnuTransBOneWhitespaceLibBTwoNum}{0}
\rqSet{rqBnuTransBOneWhitespaceLibBTwoDen}{27}
\rqSet{rqBnuTransBOneWhitespaceLibBTwoPct}{0.0}
\rqSet{rqBnuTransBOneWhitespaceLibPZeroNum}{0}
\rqSet{rqBnuTransBOneWhitespaceLibPZeroDen}{0}
\rqSet{rqBnuTransBOneWhitespaceLibPZeroPct}{n/a}
\rqSet{rqBnuTransBOneWhitespaceLibPOneNum}{0}
\rqSet{rqBnuTransBOneWhitespaceLibPOneDen}{0}
\rqSet{rqBnuTransBOneWhitespaceLibPOnePct}{n/a}
\rqSet{rqBnuDivBTwoNumericLibAllNum}{0}
\rqSet{rqBnuDivBTwoNumericLibAllDen}{18194}
\rqSet{rqBnuDivBTwoNumericLibAllPct}{0.0}
\rqSet{rqBnuDivBTwoNumericLibBOneNum}{0}
\rqSet{rqBnuDivBTwoNumericLibBOneDen}{10507}
\rqSet{rqBnuDivBTwoNumericLibBOnePct}{0.0}
\rqSet{rqBnuDivBTwoNumericLibBTwoNum}{0}
\rqSet{rqBnuDivBTwoNumericLibBTwoDen}{7687}
\rqSet{rqBnuDivBTwoNumericLibBTwoPct}{0.0}
\rqSet{rqBnuDivBTwoNumericLibPZeroNum}{0}
\rqSet{rqBnuDivBTwoNumericLibPZeroDen}{0}
\rqSet{rqBnuDivBTwoNumericLibPZeroPct}{n/a}
\rqSet{rqBnuDivBTwoNumericLibPOneNum}{0}
\rqSet{rqBnuDivBTwoNumericLibPOneDen}{0}
\rqSet{rqBnuDivBTwoNumericLibPOnePct}{n/a}
\rqSet{rqBnuTransBTwoNumericLibAllNum}{0}
\rqSet{rqBnuTransBTwoNumericLibAllDen}{65}
\rqSet{rqBnuTransBTwoNumericLibAllPct}{0.0}
\rqSet{rqBnuTransBTwoNumericLibBOneNum}{0}
\rqSet{rqBnuTransBTwoNumericLibBOneDen}{38}
\rqSet{rqBnuTransBTwoNumericLibBOnePct}{0.0}
\rqSet{rqBnuTransBTwoNumericLibBTwoNum}{0}
\rqSet{rqBnuTransBTwoNumericLibBTwoDen}{27}
\rqSet{rqBnuTransBTwoNumericLibBTwoPct}{0.0}
\rqSet{rqBnuTransBTwoNumericLibPZeroNum}{0}
\rqSet{rqBnuTransBTwoNumericLibPZeroDen}{0}
\rqSet{rqBnuTransBTwoNumericLibPZeroPct}{n/a}
\rqSet{rqBnuTransBTwoNumericLibPOneNum}{0}
\rqSet{rqBnuTransBTwoNumericLibPOneDen}{0}
\rqSet{rqBnuTransBTwoNumericLibPOnePct}{n/a}
\rqSet{rqBnuDivNOneStatusLibAllNum}{5798}
\rqSet{rqBnuDivNOneStatusLibAllDen}{18194}
\rqSet{rqBnuDivNOneStatusLibAllPct}{31.9}
\rqSet{rqBnuDivNOneStatusLibBOneNum}{5004}
\rqSet{rqBnuDivNOneStatusLibBOneDen}{10507}
\rqSet{rqBnuDivNOneStatusLibBOnePct}{47.6}
\rqSet{rqBnuDivNOneStatusLibBTwoNum}{794}
\rqSet{rqBnuDivNOneStatusLibBTwoDen}{7687}
\rqSet{rqBnuDivNOneStatusLibBTwoPct}{10.3}
\rqSet{rqBnuDivNOneStatusLibPZeroNum}{0}
\rqSet{rqBnuDivNOneStatusLibPZeroDen}{0}
\rqSet{rqBnuDivNOneStatusLibPZeroPct}{n/a}
\rqSet{rqBnuDivNOneStatusLibPOneNum}{0}
\rqSet{rqBnuDivNOneStatusLibPOneDen}{0}
\rqSet{rqBnuDivNOneStatusLibPOnePct}{n/a}
\rqSet{rqBnuTransNOneStatusLibAllNum}{25}
\rqSet{rqBnuTransNOneStatusLibAllDen}{65}
\rqSet{rqBnuTransNOneStatusLibAllPct}{38.5}
\rqSet{rqBnuTransNOneStatusLibBOneNum}{11}
\rqSet{rqBnuTransNOneStatusLibBOneDen}{38}
\rqSet{rqBnuTransNOneStatusLibBOnePct}{28.9}
\rqSet{rqBnuTransNOneStatusLibBTwoNum}{14}
\rqSet{rqBnuTransNOneStatusLibBTwoDen}{27}
\rqSet{rqBnuTransNOneStatusLibBTwoPct}{51.9}
\rqSet{rqBnuTransNOneStatusLibPZeroNum}{0}
\rqSet{rqBnuTransNOneStatusLibPZeroDen}{0}
\rqSet{rqBnuTransNOneStatusLibPZeroPct}{n/a}
\rqSet{rqBnuTransNOneStatusLibPOneNum}{0}
\rqSet{rqBnuTransNOneStatusLibPOneDen}{0}
\rqSet{rqBnuTransNOneStatusLibPOnePct}{n/a}
\rqSet{rqBnuDivNTwoTimeoutLibAllNum}{0}
\rqSet{rqBnuDivNTwoTimeoutLibAllDen}{18194}
\rqSet{rqBnuDivNTwoTimeoutLibAllPct}{0.0}
\rqSet{rqBnuDivNTwoTimeoutLibBOneNum}{0}
\rqSet{rqBnuDivNTwoTimeoutLibBOneDen}{10507}
\rqSet{rqBnuDivNTwoTimeoutLibBOnePct}{0.0}
\rqSet{rqBnuDivNTwoTimeoutLibBTwoNum}{0}
\rqSet{rqBnuDivNTwoTimeoutLibBTwoDen}{7687}
\rqSet{rqBnuDivNTwoTimeoutLibBTwoPct}{0.0}
\rqSet{rqBnuDivNTwoTimeoutLibPZeroNum}{0}
\rqSet{rqBnuDivNTwoTimeoutLibPZeroDen}{0}
\rqSet{rqBnuDivNTwoTimeoutLibPZeroPct}{n/a}
\rqSet{rqBnuDivNTwoTimeoutLibPOneNum}{0}
\rqSet{rqBnuDivNTwoTimeoutLibPOneDen}{0}
\rqSet{rqBnuDivNTwoTimeoutLibPOnePct}{n/a}
\rqSet{rqBnuTransNTwoTimeoutLibAllNum}{0}
\rqSet{rqBnuTransNTwoTimeoutLibAllDen}{65}
\rqSet{rqBnuTransNTwoTimeoutLibAllPct}{0.0}
\rqSet{rqBnuTransNTwoTimeoutLibBOneNum}{0}
\rqSet{rqBnuTransNTwoTimeoutLibBOneDen}{38}
\rqSet{rqBnuTransNTwoTimeoutLibBOnePct}{0.0}
\rqSet{rqBnuTransNTwoTimeoutLibBTwoNum}{0}
\rqSet{rqBnuTransNTwoTimeoutLibBTwoDen}{27}
\rqSet{rqBnuTransNTwoTimeoutLibBTwoPct}{0.0}
\rqSet{rqBnuTransNTwoTimeoutLibPZeroNum}{0}
\rqSet{rqBnuTransNTwoTimeoutLibPZeroDen}{0}
\rqSet{rqBnuTransNTwoTimeoutLibPZeroPct}{n/a}
\rqSet{rqBnuTransNTwoTimeoutLibPOneNum}{0}
\rqSet{rqBnuTransNTwoTimeoutLibPOneDen}{0}
\rqSet{rqBnuTransNTwoTimeoutLibPOnePct}{n/a}
\rqSet{rqBnuDivNThreeLineLibAllNum}{1688}
\rqSet{rqBnuDivNThreeLineLibAllDen}{18194}
\rqSet{rqBnuDivNThreeLineLibAllPct}{9.3}
\rqSet{rqBnuDivNThreeLineLibBOneNum}{1688}
\rqSet{rqBnuDivNThreeLineLibBOneDen}{10507}
\rqSet{rqBnuDivNThreeLineLibBOnePct}{16.1}
\rqSet{rqBnuDivNThreeLineLibBTwoNum}{0}
\rqSet{rqBnuDivNThreeLineLibBTwoDen}{7687}
\rqSet{rqBnuDivNThreeLineLibBTwoPct}{0.0}
\rqSet{rqBnuDivNThreeLineLibPZeroNum}{0}
\rqSet{rqBnuDivNThreeLineLibPZeroDen}{0}
\rqSet{rqBnuDivNThreeLineLibPZeroPct}{n/a}
\rqSet{rqBnuDivNThreeLineLibPOneNum}{0}
\rqSet{rqBnuDivNThreeLineLibPOneDen}{0}
\rqSet{rqBnuDivNThreeLineLibPOnePct}{n/a}
\rqSet{rqBnuTransNThreeLineLibAllNum}{6}
\rqSet{rqBnuTransNThreeLineLibAllDen}{65}
\rqSet{rqBnuTransNThreeLineLibAllPct}{9.2}
\rqSet{rqBnuTransNThreeLineLibBOneNum}{6}
\rqSet{rqBnuTransNThreeLineLibBOneDen}{38}
\rqSet{rqBnuTransNThreeLineLibBOnePct}{15.8}
\rqSet{rqBnuTransNThreeLineLibBTwoNum}{0}
\rqSet{rqBnuTransNThreeLineLibBTwoDen}{27}
\rqSet{rqBnuTransNThreeLineLibBTwoPct}{0.0}
\rqSet{rqBnuTransNThreeLineLibPZeroNum}{0}
\rqSet{rqBnuTransNThreeLineLibPZeroDen}{0}
\rqSet{rqBnuTransNThreeLineLibPZeroPct}{n/a}
\rqSet{rqBnuTransNThreeLineLibPOneNum}{0}
\rqSet{rqBnuTransNThreeLineLibPOneDen}{0}
\rqSet{rqBnuTransNThreeLineLibPOnePct}{n/a}
\rqSet{rqBnuDivNFourFieldLibAllNum}{8031}
\rqSet{rqBnuDivNFourFieldLibAllDen}{18194}
\rqSet{rqBnuDivNFourFieldLibAllPct}{44.1}
\rqSet{rqBnuDivNFourFieldLibBOneNum}{3791}
\rqSet{rqBnuDivNFourFieldLibBOneDen}{10507}
\rqSet{rqBnuDivNFourFieldLibBOnePct}{36.1}
\rqSet{rqBnuDivNFourFieldLibBTwoNum}{4240}
\rqSet{rqBnuDivNFourFieldLibBTwoDen}{7687}
\rqSet{rqBnuDivNFourFieldLibBTwoPct}{55.2}
\rqSet{rqBnuDivNFourFieldLibPZeroNum}{0}
\rqSet{rqBnuDivNFourFieldLibPZeroDen}{0}
\rqSet{rqBnuDivNFourFieldLibPZeroPct}{n/a}
\rqSet{rqBnuDivNFourFieldLibPOneNum}{0}
\rqSet{rqBnuDivNFourFieldLibPOneDen}{0}
\rqSet{rqBnuDivNFourFieldLibPOnePct}{n/a}
\rqSet{rqBnuTransNFourFieldLibAllNum}{30}
\rqSet{rqBnuTransNFourFieldLibAllDen}{65}
\rqSet{rqBnuTransNFourFieldLibAllPct}{46.2}
\rqSet{rqBnuTransNFourFieldLibBOneNum}{16}
\rqSet{rqBnuTransNFourFieldLibBOneDen}{38}
\rqSet{rqBnuTransNFourFieldLibBOnePct}{42.1}
\rqSet{rqBnuTransNFourFieldLibBTwoNum}{14}
\rqSet{rqBnuTransNFourFieldLibBTwoDen}{27}
\rqSet{rqBnuTransNFourFieldLibBTwoPct}{51.9}
\rqSet{rqBnuTransNFourFieldLibPZeroNum}{0}
\rqSet{rqBnuTransNFourFieldLibPZeroDen}{0}
\rqSet{rqBnuTransNFourFieldLibPZeroPct}{n/a}
\rqSet{rqBnuTransNFourFieldLibPOneNum}{0}
\rqSet{rqBnuTransNFourFieldLibPOneDen}{0}
\rqSet{rqBnuTransNFourFieldLibPOnePct}{n/a}
\rqSet{rqBnuDivUResidualLibAllNum}{1861}
\rqSet{rqBnuDivUResidualLibAllDen}{18194}
\rqSet{rqBnuDivUResidualLibAllPct}{10.2}
\rqSet{rqBnuDivUResidualLibBOneNum}{6}
\rqSet{rqBnuDivUResidualLibBOneDen}{10507}
\rqSet{rqBnuDivUResidualLibBOnePct}{0.1}
\rqSet{rqBnuDivUResidualLibBTwoNum}{1855}
\rqSet{rqBnuDivUResidualLibBTwoDen}{7687}
\rqSet{rqBnuDivUResidualLibBTwoPct}{24.1}
\rqSet{rqBnuDivUResidualLibPZeroNum}{0}
\rqSet{rqBnuDivUResidualLibPZeroDen}{0}
\rqSet{rqBnuDivUResidualLibPZeroPct}{n/a}
\rqSet{rqBnuDivUResidualLibPOneNum}{0}
\rqSet{rqBnuDivUResidualLibPOneDen}{0}
\rqSet{rqBnuDivUResidualLibPOnePct}{n/a}
\rqSet{rqBnuTransUResidualLibAllNum}{3}
\rqSet{rqBnuTransUResidualLibAllDen}{65}
\rqSet{rqBnuTransUResidualLibAllPct}{4.6}
\rqSet{rqBnuTransUResidualLibBOneNum}{2}
\rqSet{rqBnuTransUResidualLibBOneDen}{38}
\rqSet{rqBnuTransUResidualLibBOnePct}{5.3}
\rqSet{rqBnuTransUResidualLibBTwoNum}{1}
\rqSet{rqBnuTransUResidualLibBTwoDen}{27}
\rqSet{rqBnuTransUResidualLibBTwoPct}{3.7}
\rqSet{rqBnuTransUResidualLibPZeroNum}{0}
\rqSet{rqBnuTransUResidualLibPZeroDen}{0}
\rqSet{rqBnuTransUResidualLibPZeroPct}{n/a}
\rqSet{rqBnuTransUResidualLibPOneNum}{0}
\rqSet{rqBnuTransUResidualLibPOneDen}{0}
\rqSet{rqBnuTransUResidualLibPOnePct}{n/a}
\rqSet{rqBnuDivUStderrLibAllNum}{798}
\rqSet{rqBnuDivUStderrLibAllDen}{18194}
\rqSet{rqBnuDivUStderrLibAllPct}{4.4}
\rqSet{rqBnuDivUStderrLibBOneNum}{0}
\rqSet{rqBnuDivUStderrLibBOneDen}{10507}
\rqSet{rqBnuDivUStderrLibBOnePct}{0.0}
\rqSet{rqBnuDivUStderrLibBTwoNum}{798}
\rqSet{rqBnuDivUStderrLibBTwoDen}{7687}
\rqSet{rqBnuDivUStderrLibBTwoPct}{10.4}
\rqSet{rqBnuDivUStderrLibPZeroNum}{0}
\rqSet{rqBnuDivUStderrLibPZeroDen}{0}
\rqSet{rqBnuDivUStderrLibPZeroPct}{n/a}
\rqSet{rqBnuDivUStderrLibPOneNum}{0}
\rqSet{rqBnuDivUStderrLibPOneDen}{0}
\rqSet{rqBnuDivUStderrLibPOnePct}{n/a}
\rqSet{rqBnuTransUStderrLibAllNum}{3}
\rqSet{rqBnuTransUStderrLibAllDen}{65}
\rqSet{rqBnuTransUStderrLibAllPct}{4.6}
\rqSet{rqBnuTransUStderrLibBOneNum}{0}
\rqSet{rqBnuTransUStderrLibBOneDen}{38}
\rqSet{rqBnuTransUStderrLibBOnePct}{0.0}
\rqSet{rqBnuTransUStderrLibBTwoNum}{3}
\rqSet{rqBnuTransUStderrLibBTwoDen}{27}
\rqSet{rqBnuTransUStderrLibBTwoPct}{11.1}
\rqSet{rqBnuTransUStderrLibPZeroNum}{0}
\rqSet{rqBnuTransUStderrLibPZeroDen}{0}
\rqSet{rqBnuTransUStderrLibPZeroPct}{n/a}
\rqSet{rqBnuTransUStderrLibPOneNum}{0}
\rqSet{rqBnuTransUStderrLibPOneDen}{0}
\rqSet{rqBnuTransUStderrLibPOnePct}{n/a}
\rqSet{rqRobCCrashLibAllNum}{55}
\rqSet{rqRobCCrashLibAllDen}{1088}
\rqSet{rqRobCCrashLibAllPct}{5.1}
\rqSet{rqRobCCrashLibAllKiroNum}{4}
\rqSet{rqRobCCrashLibAllKiroDen}{192}
\rqSet{rqRobCCrashLibAllKiroPct}{2.1}
\rqSet{rqRobCCrashLibAllClaudeNum}{3}
\rqSet{rqRobCCrashLibAllClaudeDen}{181}
\rqSet{rqRobCCrashLibAllClaudePct}{1.7}
\rqSet{rqRobCCrashLibAllCodexNum}{2}
\rqSet{rqRobCCrashLibAllCodexDen}{143}
\rqSet{rqRobCCrashLibAllCodexPct}{1.4}
\rqSet{rqRobCCrashLibAllCTwoRustNum}{3}
\rqSet{rqRobCCrashLibAllCTwoRustDen}{116}
\rqSet{rqRobCCrashLibAllCTwoRustPct}{2.6}
\rqSet{rqRobCCrashLibAllLaertesNum}{3}
\rqSet{rqRobCCrashLibAllLaertesDen}{114}
\rqSet{rqRobCCrashLibAllLaertesPct}{2.6}
\rqSet{rqRobCCrashLibAllCTwoSaferRustNum}{7}
\rqSet{rqRobCCrashLibAllCTwoSaferRustDen}{87}
\rqSet{rqRobCCrashLibAllCTwoSaferRustPct}{8.0}
\rqSet{rqRobCCrashLibAllSmartCTwoRustNum}{0}
\rqSet{rqRobCCrashLibAllSmartCTwoRustDen}{0}
\rqSet{rqRobCCrashLibAllSmartCTwoRustPct}{n/a}
\rqSet{rqRobCCrashLibAllKimiNum}{14}
\rqSet{rqRobCCrashLibAllKimiDen}{73}
\rqSet{rqRobCCrashLibAllKimiPct}{19.2}
\rqSet{rqRobCCrashLibAllGptFiveNum}{13}
\rqSet{rqRobCCrashLibAllGptFiveDen}{92}
\rqSet{rqRobCCrashLibAllGptFivePct}{14.1}
\rqSet{rqRobCCrashLibAllGeminiNum}{6}
\rqSet{rqRobCCrashLibAllGeminiDen}{90}
\rqSet{rqRobCCrashLibAllGeminiPct}{6.7}
\rqSet{rqRobCCrashLibBOneNum}{38}
\rqSet{rqRobCCrashLibBOneDen}{540}
\rqSet{rqRobCCrashLibBOnePct}{7.0}
\rqSet{rqRobCCrashLibBTwoNum}{17}
\rqSet{rqRobCCrashLibBTwoDen}{379}
\rqSet{rqRobCCrashLibBTwoPct}{4.5}
\rqSet{rqRobCCrashLibPZeroNum}{0}
\rqSet{rqRobCCrashLibPZeroDen}{0}
\rqSet{rqRobCCrashLibPZeroPct}{n/a}
\rqSet{rqRobCCrashLibPOneNum}{0}
\rqSet{rqRobCCrashLibPOneDen}{169}
\rqSet{rqRobCCrashLibPOnePct}{0.0}
\rqSet{rqRobUnstableLibAllNum}{90}
\rqSet{rqRobUnstableLibAllDen}{1088}
\rqSet{rqRobUnstableLibAllPct}{8.3}
\rqSet{rqRobUnstableLibAllKiroNum}{11}
\rqSet{rqRobUnstableLibAllKiroDen}{192}
\rqSet{rqRobUnstableLibAllKiroPct}{5.7}
\rqSet{rqRobUnstableLibAllClaudeNum}{9}
\rqSet{rqRobUnstableLibAllClaudeDen}{181}
\rqSet{rqRobUnstableLibAllClaudePct}{5.0}
\rqSet{rqRobUnstableLibAllCodexNum}{5}
\rqSet{rqRobUnstableLibAllCodexDen}{143}
\rqSet{rqRobUnstableLibAllCodexPct}{3.5}
\rqSet{rqRobUnstableLibAllCTwoRustNum}{6}
\rqSet{rqRobUnstableLibAllCTwoRustDen}{116}
\rqSet{rqRobUnstableLibAllCTwoRustPct}{5.2}
\rqSet{rqRobUnstableLibAllLaertesNum}{4}
\rqSet{rqRobUnstableLibAllLaertesDen}{114}
\rqSet{rqRobUnstableLibAllLaertesPct}{3.5}
\rqSet{rqRobUnstableLibAllCTwoSaferRustNum}{10}
\rqSet{rqRobUnstableLibAllCTwoSaferRustDen}{87}
\rqSet{rqRobUnstableLibAllCTwoSaferRustPct}{11.5}
\rqSet{rqRobUnstableLibAllSmartCTwoRustNum}{0}
\rqSet{rqRobUnstableLibAllSmartCTwoRustDen}{0}
\rqSet{rqRobUnstableLibAllSmartCTwoRustPct}{n/a}
\rqSet{rqRobUnstableLibAllKimiNum}{18}
\rqSet{rqRobUnstableLibAllKimiDen}{73}
\rqSet{rqRobUnstableLibAllKimiPct}{24.7}
\rqSet{rqRobUnstableLibAllGptFiveNum}{18}
\rqSet{rqRobUnstableLibAllGptFiveDen}{92}
\rqSet{rqRobUnstableLibAllGptFivePct}{19.6}
\rqSet{rqRobUnstableLibAllGeminiNum}{9}
\rqSet{rqRobUnstableLibAllGeminiDen}{90}
\rqSet{rqRobUnstableLibAllGeminiPct}{10.0}
\rqSet{rqRobUnstableLibBOneNum}{62}
\rqSet{rqRobUnstableLibBOneDen}{540}
\rqSet{rqRobUnstableLibBOnePct}{11.5}
\rqSet{rqRobUnstableLibBTwoNum}{28}
\rqSet{rqRobUnstableLibBTwoDen}{379}
\rqSet{rqRobUnstableLibBTwoPct}{7.4}
\rqSet{rqRobUnstableLibPZeroNum}{0}
\rqSet{rqRobUnstableLibPZeroDen}{0}
\rqSet{rqRobUnstableLibPZeroPct}{n/a}
\rqSet{rqRobUnstableLibPOneNum}{0}
\rqSet{rqRobUnstableLibPOneDen}{169}
\rqSet{rqRobUnstableLibPOnePct}{0.0}
\rqSet{rqRobExactLibAllNum}{??}
\rqSet{rqRobExactLibAllDen}{1088}
\rqSet{rqRobExactLibAllPct}{??}
\rqSet{rqRobExactLibAllKiroNum}{??}
\rqSet{rqRobExactLibAllKiroDen}{192}
\rqSet{rqRobExactLibAllKiroPct}{??}
\rqSet{rqRobExactLibAllClaudeNum}{??}
\rqSet{rqRobExactLibAllClaudeDen}{181}
\rqSet{rqRobExactLibAllClaudePct}{??}
\rqSet{rqRobExactLibAllCodexNum}{??}
\rqSet{rqRobExactLibAllCodexDen}{143}
\rqSet{rqRobExactLibAllCodexPct}{??}
\rqSet{rqRobExactLibAllCTwoRustNum}{??}
\rqSet{rqRobExactLibAllCTwoRustDen}{116}
\rqSet{rqRobExactLibAllCTwoRustPct}{??}
\rqSet{rqRobExactLibAllLaertesNum}{??}
\rqSet{rqRobExactLibAllLaertesDen}{114}
\rqSet{rqRobExactLibAllLaertesPct}{??}
\rqSet{rqRobExactLibAllCTwoSaferRustNum}{??}
\rqSet{rqRobExactLibAllCTwoSaferRustDen}{87}
\rqSet{rqRobExactLibAllCTwoSaferRustPct}{??}
\rqSet{rqRobExactLibAllSmartCTwoRustNum}{0}
\rqSet{rqRobExactLibAllSmartCTwoRustDen}{0}
\rqSet{rqRobExactLibAllSmartCTwoRustPct}{n/a}
\rqSet{rqRobExactLibAllKimiNum}{??}
\rqSet{rqRobExactLibAllKimiDen}{73}
\rqSet{rqRobExactLibAllKimiPct}{??}
\rqSet{rqRobExactLibAllGptFiveNum}{??}
\rqSet{rqRobExactLibAllGptFiveDen}{92}
\rqSet{rqRobExactLibAllGptFivePct}{??}
\rqSet{rqRobExactLibAllGeminiNum}{??}
\rqSet{rqRobExactLibAllGeminiDen}{90}
\rqSet{rqRobExactLibAllGeminiPct}{??}
\rqSet{rqRobExactLibBOneNum}{??}
\rqSet{rqRobExactLibBOneDen}{540}
\rqSet{rqRobExactLibBOnePct}{??}
\rqSet{rqRobExactLibBTwoNum}{??}
\rqSet{rqRobExactLibBTwoDen}{379}
\rqSet{rqRobExactLibBTwoPct}{??}
\rqSet{rqRobExactLibPZeroNum}{0}
\rqSet{rqRobExactLibPZeroDen}{0}
\rqSet{rqRobExactLibPZeroPct}{n/a}
\rqSet{rqRobExactLibPOneNum}{??}
\rqSet{rqRobExactLibPOneDen}{169}
\rqSet{rqRobExactLibPOnePct}{??}
\rqSet{rqRobOrigFloatLibAllNum}{??}
\rqSet{rqRobOrigFloatLibAllDen}{1088}
\rqSet{rqRobOrigFloatLibAllPct}{??}
\rqSet{rqRobOrigFloatLibAllKiroNum}{??}
\rqSet{rqRobOrigFloatLibAllKiroDen}{192}
\rqSet{rqRobOrigFloatLibAllKiroPct}{??}
\rqSet{rqRobOrigFloatLibAllClaudeNum}{??}
\rqSet{rqRobOrigFloatLibAllClaudeDen}{181}
\rqSet{rqRobOrigFloatLibAllClaudePct}{??}
\rqSet{rqRobOrigFloatLibAllCodexNum}{??}
\rqSet{rqRobOrigFloatLibAllCodexDen}{143}
\rqSet{rqRobOrigFloatLibAllCodexPct}{??}
\rqSet{rqRobOrigFloatLibAllCTwoRustNum}{??}
\rqSet{rqRobOrigFloatLibAllCTwoRustDen}{116}
\rqSet{rqRobOrigFloatLibAllCTwoRustPct}{??}
\rqSet{rqRobOrigFloatLibAllLaertesNum}{??}
\rqSet{rqRobOrigFloatLibAllLaertesDen}{114}
\rqSet{rqRobOrigFloatLibAllLaertesPct}{??}
\rqSet{rqRobOrigFloatLibAllCTwoSaferRustNum}{??}
\rqSet{rqRobOrigFloatLibAllCTwoSaferRustDen}{87}
\rqSet{rqRobOrigFloatLibAllCTwoSaferRustPct}{??}
\rqSet{rqRobOrigFloatLibAllSmartCTwoRustNum}{0}
\rqSet{rqRobOrigFloatLibAllSmartCTwoRustDen}{0}
\rqSet{rqRobOrigFloatLibAllSmartCTwoRustPct}{n/a}
\rqSet{rqRobOrigFloatLibAllKimiNum}{??}
\rqSet{rqRobOrigFloatLibAllKimiDen}{73}
\rqSet{rqRobOrigFloatLibAllKimiPct}{??}
\rqSet{rqRobOrigFloatLibAllGptFiveNum}{??}
\rqSet{rqRobOrigFloatLibAllGptFiveDen}{92}
\rqSet{rqRobOrigFloatLibAllGptFivePct}{??}
\rqSet{rqRobOrigFloatLibAllGeminiNum}{??}
\rqSet{rqRobOrigFloatLibAllGeminiDen}{90}
\rqSet{rqRobOrigFloatLibAllGeminiPct}{??}
\rqSet{rqRobOrigFloatLibBOneNum}{??}
\rqSet{rqRobOrigFloatLibBOneDen}{540}
\rqSet{rqRobOrigFloatLibBOnePct}{??}
\rqSet{rqRobOrigFloatLibBTwoNum}{??}
\rqSet{rqRobOrigFloatLibBTwoDen}{379}
\rqSet{rqRobOrigFloatLibBTwoPct}{??}
\rqSet{rqRobOrigFloatLibPZeroNum}{0}
\rqSet{rqRobOrigFloatLibPZeroDen}{0}
\rqSet{rqRobOrigFloatLibPZeroPct}{n/a}
\rqSet{rqRobOrigFloatLibPOneNum}{??}
\rqSet{rqRobOrigFloatLibPOneDen}{169}
\rqSet{rqRobOrigFloatLibPOnePct}{??}
\rqSet{rqRobOrigFloatRemovedDivLibAllNum}{0}
\rqSet{rqRobOrigFloatRemovedDivLibAllDen}{18194}
\rqSet{rqRobOrigFloatRemovedDivLibAllPct}{0.0}
\rqSet{rqRobOrigFloatRemovedDivLibBOneNum}{0}
\rqSet{rqRobOrigFloatRemovedDivLibBOneDen}{10507}
\rqSet{rqRobOrigFloatRemovedDivLibBOnePct}{0.0}
\rqSet{rqRobOrigFloatRemovedDivLibBTwoNum}{0}
\rqSet{rqRobOrigFloatRemovedDivLibBTwoDen}{7687}
\rqSet{rqRobOrigFloatRemovedDivLibBTwoPct}{0.0}
\rqSet{rqRobOrigFloatRemovedDivLibPZeroNum}{0}
\rqSet{rqRobOrigFloatRemovedDivLibPZeroDen}{0}
\rqSet{rqRobOrigFloatRemovedDivLibPZeroPct}{n/a}
\rqSet{rqRobOrigFloatRemovedDivLibPOneNum}{0}
\rqSet{rqRobOrigFloatRemovedDivLibPOneDen}{0}
\rqSet{rqRobOrigFloatRemovedDivLibPOnePct}{n/a}
\rqSet{rqRobOrigFloatAllRemovedLibAllNum}{0}
\rqSet{rqRobOrigFloatAllRemovedLibAllDen}{65}
\rqSet{rqRobOrigFloatAllRemovedLibAllPct}{0.0}
\rqSet{rqRobOrigFloatAllRemovedLibBOneNum}{0}
\rqSet{rqRobOrigFloatAllRemovedLibBOneDen}{38}
\rqSet{rqRobOrigFloatAllRemovedLibBOnePct}{0.0}
\rqSet{rqRobOrigFloatAllRemovedLibBTwoNum}{0}
\rqSet{rqRobOrigFloatAllRemovedLibBTwoDen}{27}
\rqSet{rqRobOrigFloatAllRemovedLibBTwoPct}{0.0}
\rqSet{rqRobOrigFloatAllRemovedLibPZeroNum}{0}
\rqSet{rqRobOrigFloatAllRemovedLibPZeroDen}{0}
\rqSet{rqRobOrigFloatAllRemovedLibPZeroPct}{n/a}
\rqSet{rqRobOrigFloatAllRemovedLibPOneNum}{0}
\rqSet{rqRobOrigFloatAllRemovedLibPOneDen}{0}
\rqSet{rqRobOrigFloatAllRemovedLibPOnePct}{n/a}
\rqSet{rqRobNoStderrLibAllNum}{65}
\rqSet{rqRobNoStderrLibAllDen}{1088}
\rqSet{rqRobNoStderrLibAllPct}{6.0}
\rqSet{rqRobNoStderrLibAllKiroNum}{7}
\rqSet{rqRobNoStderrLibAllKiroDen}{192}
\rqSet{rqRobNoStderrLibAllKiroPct}{3.6}
\rqSet{rqRobNoStderrLibAllClaudeNum}{6}
\rqSet{rqRobNoStderrLibAllClaudeDen}{181}
\rqSet{rqRobNoStderrLibAllClaudePct}{3.3}
\rqSet{rqRobNoStderrLibAllCodexNum}{2}
\rqSet{rqRobNoStderrLibAllCodexDen}{143}
\rqSet{rqRobNoStderrLibAllCodexPct}{1.4}
\rqSet{rqRobNoStderrLibAllCTwoRustNum}{3}
\rqSet{rqRobNoStderrLibAllCTwoRustDen}{116}
\rqSet{rqRobNoStderrLibAllCTwoRustPct}{2.6}
\rqSet{rqRobNoStderrLibAllLaertesNum}{3}
\rqSet{rqRobNoStderrLibAllLaertesDen}{114}
\rqSet{rqRobNoStderrLibAllLaertesPct}{2.6}
\rqSet{rqRobNoStderrLibAllCTwoSaferRustNum}{7}
\rqSet{rqRobNoStderrLibAllCTwoSaferRustDen}{87}
\rqSet{rqRobNoStderrLibAllCTwoSaferRustPct}{8.0}
\rqSet{rqRobNoStderrLibAllSmartCTwoRustNum}{0}
\rqSet{rqRobNoStderrLibAllSmartCTwoRustDen}{0}
\rqSet{rqRobNoStderrLibAllSmartCTwoRustPct}{n/a}
\rqSet{rqRobNoStderrLibAllKimiNum}{16}
\rqSet{rqRobNoStderrLibAllKimiDen}{73}
\rqSet{rqRobNoStderrLibAllKimiPct}{21.9}
\rqSet{rqRobNoStderrLibAllGptFiveNum}{15}
\rqSet{rqRobNoStderrLibAllGptFiveDen}{92}
\rqSet{rqRobNoStderrLibAllGptFivePct}{16.3}
\rqSet{rqRobNoStderrLibAllGeminiNum}{6}
\rqSet{rqRobNoStderrLibAllGeminiDen}{90}
\rqSet{rqRobNoStderrLibAllGeminiPct}{6.7}
\rqSet{rqRobNoStderrLibBOneNum}{38}
\rqSet{rqRobNoStderrLibBOneDen}{540}
\rqSet{rqRobNoStderrLibBOnePct}{7.0}
\rqSet{rqRobNoStderrLibBTwoNum}{27}
\rqSet{rqRobNoStderrLibBTwoDen}{379}
\rqSet{rqRobNoStderrLibBTwoPct}{7.1}
\rqSet{rqRobNoStderrLibPZeroNum}{0}
\rqSet{rqRobNoStderrLibPZeroDen}{0}
\rqSet{rqRobNoStderrLibPZeroPct}{n/a}
\rqSet{rqRobNoStderrLibPOneNum}{0}
\rqSet{rqRobNoStderrLibPOneDen}{169}
\rqSet{rqRobNoStderrLibPOnePct}{0.0}
\rqSet{rqRobCapEightLibAllNum}{65}
\rqSet{rqRobCapEightLibAllDen}{1088}
\rqSet{rqRobCapEightLibAllPct}{6.0}
\rqSet{rqRobCapEightLibAllKiroNum}{7}
\rqSet{rqRobCapEightLibAllKiroDen}{192}
\rqSet{rqRobCapEightLibAllKiroPct}{3.6}
\rqSet{rqRobCapEightLibAllClaudeNum}{6}
\rqSet{rqRobCapEightLibAllClaudeDen}{181}
\rqSet{rqRobCapEightLibAllClaudePct}{3.3}
\rqSet{rqRobCapEightLibAllCodexNum}{2}
\rqSet{rqRobCapEightLibAllCodexDen}{143}
\rqSet{rqRobCapEightLibAllCodexPct}{1.4}
\rqSet{rqRobCapEightLibAllCTwoRustNum}{3}
\rqSet{rqRobCapEightLibAllCTwoRustDen}{116}
\rqSet{rqRobCapEightLibAllCTwoRustPct}{2.6}
\rqSet{rqRobCapEightLibAllLaertesNum}{3}
\rqSet{rqRobCapEightLibAllLaertesDen}{114}
\rqSet{rqRobCapEightLibAllLaertesPct}{2.6}
\rqSet{rqRobCapEightLibAllCTwoSaferRustNum}{7}
\rqSet{rqRobCapEightLibAllCTwoSaferRustDen}{87}
\rqSet{rqRobCapEightLibAllCTwoSaferRustPct}{8.0}
\rqSet{rqRobCapEightLibAllSmartCTwoRustNum}{0}
\rqSet{rqRobCapEightLibAllSmartCTwoRustDen}{0}
\rqSet{rqRobCapEightLibAllSmartCTwoRustPct}{n/a}
\rqSet{rqRobCapEightLibAllKimiNum}{16}
\rqSet{rqRobCapEightLibAllKimiDen}{73}
\rqSet{rqRobCapEightLibAllKimiPct}{21.9}
\rqSet{rqRobCapEightLibAllGptFiveNum}{15}
\rqSet{rqRobCapEightLibAllGptFiveDen}{92}
\rqSet{rqRobCapEightLibAllGptFivePct}{16.3}
\rqSet{rqRobCapEightLibAllGeminiNum}{6}
\rqSet{rqRobCapEightLibAllGeminiDen}{90}
\rqSet{rqRobCapEightLibAllGeminiPct}{6.7}
\rqSet{rqRobCapEightLibBOneNum}{38}
\rqSet{rqRobCapEightLibBOneDen}{540}
\rqSet{rqRobCapEightLibBOnePct}{7.0}
\rqSet{rqRobCapEightLibBTwoNum}{27}
\rqSet{rqRobCapEightLibBTwoDen}{379}
\rqSet{rqRobCapEightLibBTwoPct}{7.1}
\rqSet{rqRobCapEightLibPZeroNum}{0}
\rqSet{rqRobCapEightLibPZeroDen}{0}
\rqSet{rqRobCapEightLibPZeroPct}{n/a}
\rqSet{rqRobCapEightLibPOneNum}{0}
\rqSet{rqRobCapEightLibPOneDen}{169}
\rqSet{rqRobCapEightLibPOnePct}{0.0}
\rqSet{rqRobInconclLibAllNum}{73}
\rqSet{rqRobInconclLibAllDen}{1088}
\rqSet{rqRobInconclLibAllPct}{6.7}
\rqSet{rqRobInconclLibAllKiroNum}{8}
\rqSet{rqRobInconclLibAllKiroDen}{192}
\rqSet{rqRobInconclLibAllKiroPct}{4.2}
\rqSet{rqRobInconclLibAllClaudeNum}{7}
\rqSet{rqRobInconclLibAllClaudeDen}{181}
\rqSet{rqRobInconclLibAllClaudePct}{3.9}
\rqSet{rqRobInconclLibAllCodexNum}{2}
\rqSet{rqRobInconclLibAllCodexDen}{143}
\rqSet{rqRobInconclLibAllCodexPct}{1.4}
\rqSet{rqRobInconclLibAllCTwoRustNum}{3}
\rqSet{rqRobInconclLibAllCTwoRustDen}{116}
\rqSet{rqRobInconclLibAllCTwoRustPct}{2.6}
\rqSet{rqRobInconclLibAllLaertesNum}{3}
\rqSet{rqRobInconclLibAllLaertesDen}{114}
\rqSet{rqRobInconclLibAllLaertesPct}{2.6}
\rqSet{rqRobInconclLibAllCTwoSaferRustNum}{9}
\rqSet{rqRobInconclLibAllCTwoSaferRustDen}{87}
\rqSet{rqRobInconclLibAllCTwoSaferRustPct}{10.3}
\rqSet{rqRobInconclLibAllSmartCTwoRustNum}{0}
\rqSet{rqRobInconclLibAllSmartCTwoRustDen}{0}
\rqSet{rqRobInconclLibAllSmartCTwoRustPct}{n/a}
\rqSet{rqRobInconclLibAllKimiNum}{17}
\rqSet{rqRobInconclLibAllKimiDen}{73}
\rqSet{rqRobInconclLibAllKimiPct}{23.3}
\rqSet{rqRobInconclLibAllGptFiveNum}{17}
\rqSet{rqRobInconclLibAllGptFiveDen}{92}
\rqSet{rqRobInconclLibAllGptFivePct}{18.5}
\rqSet{rqRobInconclLibAllGeminiNum}{7}
\rqSet{rqRobInconclLibAllGeminiDen}{90}
\rqSet{rqRobInconclLibAllGeminiPct}{7.8}
\rqSet{rqRobInconclLibBOneNum}{38}
\rqSet{rqRobInconclLibBOneDen}{540}
\rqSet{rqRobInconclLibBOnePct}{7.0}
\rqSet{rqRobInconclLibBTwoNum}{35}
\rqSet{rqRobInconclLibBTwoDen}{379}
\rqSet{rqRobInconclLibBTwoPct}{9.2}
\rqSet{rqRobInconclLibPZeroNum}{0}
\rqSet{rqRobInconclLibPZeroDen}{0}
\rqSet{rqRobInconclLibPZeroPct}{n/a}
\rqSet{rqRobInconclLibPOneNum}{0}
\rqSet{rqRobInconclLibPOneDen}{169}
\rqSet{rqRobInconclLibPOnePct}{0.0}
\rqSet{rqExactTransLibAllNum}{1}
\rqSet{rqExactTransLibAllDen}{1088}
\rqSet{rqExactTransLibAllPct}{0.1}
\rqSet{rqExactTransLibBOneNum}{0}
\rqSet{rqExactTransLibBOneDen}{540}
\rqSet{rqExactTransLibBOnePct}{0.0}
\rqSet{rqExactTransLibBTwoNum}{1}
\rqSet{rqExactTransLibBTwoDen}{379}
\rqSet{rqExactTransLibBTwoPct}{0.3}
\rqSet{rqExactTransLibPZeroNum}{0}
\rqSet{rqExactTransLibPZeroDen}{0}
\rqSet{rqExactTransLibPZeroPct}{n/a}
\rqSet{rqExactTransLibPOneNum}{0}
\rqSet{rqExactTransLibPOneDen}{169}
\rqSet{rqExactTransLibPOnePct}{0.0}
\rqSet{rqMatchTolNotExactLibAllNum}{1}
\rqSet{rqMatchTolNotExactLibAllDen}{6554751}
\rqSet{rqMatchTolNotExactLibAllPct}{0.0}
\rqSet{rqMatchTolNotExactLibBOneNum}{0}
\rqSet{rqMatchTolNotExactLibBOneDen}{3259995}
\rqSet{rqMatchTolNotExactLibBOnePct}{0.0}
\rqSet{rqMatchTolNotExactLibBTwoNum}{1}
\rqSet{rqMatchTolNotExactLibBTwoDen}{2280249}
\rqSet{rqMatchTolNotExactLibBTwoPct}{0.0}
\rqSet{rqMatchTolNotExactLibPZeroNum}{0}
\rqSet{rqMatchTolNotExactLibPZeroDen}{0}
\rqSet{rqMatchTolNotExactLibPZeroPct}{n/a}
\rqSet{rqMatchTolNotExactLibPOneNum}{0}
\rqSet{rqMatchTolNotExactLibPOneDen}{1014507}
\rqSet{rqMatchTolNotExactLibPOnePct}{0.0}
\rqSet{rqCapHitLibAllNum}{0}
\rqSet{rqCapHitLibAllDen}{6554751}
\rqSet{rqCapHitLibAllPct}{0.0}
\rqSet{rqCapHitLibBOneNum}{0}
\rqSet{rqCapHitLibBOneDen}{3259995}
\rqSet{rqCapHitLibBOnePct}{0.0}
\rqSet{rqCapHitLibBTwoNum}{0}
\rqSet{rqCapHitLibBTwoDen}{2280249}
\rqSet{rqCapHitLibBTwoPct}{0.0}
\rqSet{rqCapHitLibPZeroNum}{0}
\rqSet{rqCapHitLibPZeroDen}{0}
\rqSet{rqCapHitLibPZeroPct}{n/a}
\rqSet{rqCapHitLibPOneNum}{0}
\rqSet{rqCapHitLibPOneDen}{1014507}
\rqSet{rqCapHitLibPOnePct}{0.0}
\rqSet{rqCapHitDivLibAllNum}{0}
\rqSet{rqCapHitDivLibAllDen}{0}
\rqSet{rqCapHitDivLibAllPct}{n/a}
\rqSet{rqCapHitDivLibBOneNum}{0}
\rqSet{rqCapHitDivLibBOneDen}{0}
\rqSet{rqCapHitDivLibBOnePct}{n/a}
\rqSet{rqCapHitDivLibBTwoNum}{0}
\rqSet{rqCapHitDivLibBTwoDen}{0}
\rqSet{rqCapHitDivLibBTwoPct}{n/a}
\rqSet{rqCapHitDivLibPZeroNum}{0}
\rqSet{rqCapHitDivLibPZeroDen}{0}
\rqSet{rqCapHitDivLibPZeroPct}{n/a}
\rqSet{rqCapHitDivLibPOneNum}{0}
\rqSet{rqCapHitDivLibPOneDen}{0}
\rqSet{rqCapHitDivLibPOnePct}{n/a}
\rqSet{rqCapSensInfraLibAllNum}{0}
\rqSet{rqCapSensInfraLibAllDen}{0}
\rqSet{rqCapSensInfraLibAllPct}{n/a}
\rqSet{rqCapSensInfraLibBOneNum}{0}
\rqSet{rqCapSensInfraLibBOneDen}{0}
\rqSet{rqCapSensInfraLibBOnePct}{n/a}
\rqSet{rqCapSensInfraLibBTwoNum}{0}
\rqSet{rqCapSensInfraLibBTwoDen}{0}
\rqSet{rqCapSensInfraLibBTwoPct}{n/a}
\rqSet{rqCapSensInfraLibPZeroNum}{0}
\rqSet{rqCapSensInfraLibPZeroDen}{0}
\rqSet{rqCapSensInfraLibPZeroPct}{n/a}
\rqSet{rqCapSensInfraLibPOneNum}{0}
\rqSet{rqCapSensInfraLibPOneDen}{0}
\rqSet{rqCapSensInfraLibPOnePct}{n/a}
\rqSet{rqCapSensNewDivLibAllNum}{0}
\rqSet{rqCapSensNewDivLibAllDen}{0}
\rqSet{rqCapSensNewDivLibBOneNum}{0}
\rqSet{rqCapSensNewDivLibBOneDen}{0}
\rqSet{rqCapSensNewDivLibBTwoNum}{0}
\rqSet{rqCapSensNewDivLibBTwoDen}{0}
\rqSet{rqCapSensNewDivLibPZeroNum}{0}
\rqSet{rqCapSensNewDivLibPZeroDen}{0}
\rqSet{rqCapSensNewDivLibPOneNum}{0}
\rqSet{rqCapSensNewDivLibPOneDen}{0}
\rqSet{rqCapSensTimeoutRerunLibAllNum}{0}
\rqSet{rqCapSensTimeoutRerunLibAllDen}{0}
\rqSet{rqCapSensTimeoutRerunLibAllPct}{n/a}
\rqSet{rqCapSensTimeoutRerunLibBOneNum}{0}
\rqSet{rqCapSensTimeoutRerunLibBOneDen}{0}
\rqSet{rqCapSensTimeoutRerunLibBOnePct}{n/a}
\rqSet{rqCapSensTimeoutRerunLibBTwoNum}{0}
\rqSet{rqCapSensTimeoutRerunLibBTwoDen}{0}
\rqSet{rqCapSensTimeoutRerunLibBTwoPct}{n/a}
\rqSet{rqCapSensTimeoutRerunLibPZeroNum}{0}
\rqSet{rqCapSensTimeoutRerunLibPZeroDen}{0}
\rqSet{rqCapSensTimeoutRerunLibPZeroPct}{n/a}
\rqSet{rqCapSensTimeoutRerunLibPOneNum}{0}
\rqSet{rqCapSensTimeoutRerunLibPOneDen}{0}
\rqSet{rqCapSensTimeoutRerunLibPOnePct}{n/a}
\rqSet{rqCapOneSideLibAllNum}{0}
\rqSet{rqCapOneSideLibAllDen}{1088}
\rqSet{rqCapOneSideLibAllPct}{0.0}
\rqSet{rqCapOneSideLibBOneNum}{0}
\rqSet{rqCapOneSideLibBOneDen}{540}
\rqSet{rqCapOneSideLibBOnePct}{0.0}
\rqSet{rqCapOneSideLibBTwoNum}{0}
\rqSet{rqCapOneSideLibBTwoDen}{379}
\rqSet{rqCapOneSideLibBTwoPct}{0.0}
\rqSet{rqCapOneSideLibPZeroNum}{0}
\rqSet{rqCapOneSideLibPZeroDen}{0}
\rqSet{rqCapOneSideLibPZeroPct}{n/a}
\rqSet{rqCapOneSideLibPOneNum}{0}
\rqSet{rqCapOneSideLibPOneDen}{169}
\rqSet{rqCapOneSideLibPOnePct}{0.0}
\rqSet{rqUbDivLibAllNum}{158667}
\rqSet{rqUbDivLibAllDen}{199794}
\rqSet{rqUbDivLibAllPct}{79.4}
\rqSet{rqUbDivLibAllKiroNum}{14321}
\rqSet{rqUbDivLibAllKiroDen}{18442}
\rqSet{rqUbDivLibAllKiroPct}{77.7}
\rqSet{rqUbDivLibAllClaudeNum}{11252}
\rqSet{rqUbDivLibAllClaudeDen}{14690}
\rqSet{rqUbDivLibAllClaudePct}{76.6}
\rqSet{rqUbDivLibAllCodexNum}{12603}
\rqSet{rqUbDivLibAllCodexDen}{14880}
\rqSet{rqUbDivLibAllCodexPct}{84.7}
\rqSet{rqUbDivLibAllCTwoRustNum}{17142}
\rqSet{rqUbDivLibAllCTwoRustDen}{19432}
\rqSet{rqUbDivLibAllCTwoRustPct}{88.2}
\rqSet{rqUbDivLibAllLaertesNum}{17139}
\rqSet{rqUbDivLibAllLaertesDen}{17426}
\rqSet{rqUbDivLibAllLaertesPct}{98.4}
\rqSet{rqUbDivLibAllCTwoSaferRustNum}{18613}
\rqSet{rqUbDivLibAllCTwoSaferRustDen}{26627}
\rqSet{rqUbDivLibAllCTwoSaferRustPct}{69.9}
\rqSet{rqUbDivLibAllSmartCTwoRustNum}{0}
\rqSet{rqUbDivLibAllSmartCTwoRustDen}{0}
\rqSet{rqUbDivLibAllSmartCTwoRustPct}{n/a}
\rqSet{rqUbDivLibAllKimiNum}{14828}
\rqSet{rqUbDivLibAllKimiDen}{23227}
\rqSet{rqUbDivLibAllKimiPct}{63.8}
\rqSet{rqUbDivLibAllGptFiveNum}{26076}
\rqSet{rqUbDivLibAllGptFiveDen}{34848}
\rqSet{rqUbDivLibAllGptFivePct}{74.8}
\rqSet{rqUbDivLibAllGeminiNum}{26693}
\rqSet{rqUbDivLibAllGeminiDen}{30222}
\rqSet{rqUbDivLibAllGeminiPct}{88.3}
\rqSet{rqUbDivLibBOneNum}{79803}
\rqSet{rqUbDivLibBOneDen}{106695}
\rqSet{rqUbDivLibBOnePct}{74.8}
\rqSet{rqUbDivLibBTwoNum}{78864}
\rqSet{rqUbDivLibBTwoDen}{93099}
\rqSet{rqUbDivLibBTwoPct}{84.7}
\rqSet{rqUbDivLibPZeroNum}{0}
\rqSet{rqUbDivLibPZeroDen}{0}
\rqSet{rqUbDivLibPZeroPct}{n/a}
\rqSet{rqUbDivLibPOneNum}{0}
\rqSet{rqUbDivLibPOneDen}{0}
\rqSet{rqUbDivLibPOnePct}{n/a}
\rqSet{rqUbDivStableLibAllNum}{127673}
\rqSet{rqUbDivStableLibAllDen}{152414}
\rqSet{rqUbDivStableLibAllPct}{83.8}
\rqSet{rqUbDivStableLibBOneNum}{60854}
\rqSet{rqUbDivStableLibBOneDen}{71361}
\rqSet{rqUbDivStableLibBOnePct}{85.3}
\rqSet{rqUbDivStableLibBTwoNum}{66819}
\rqSet{rqUbDivStableLibBTwoDen}{81053}
\rqSet{rqUbDivStableLibBTwoPct}{82.4}
\rqSet{rqUbDivStableLibPZeroNum}{0}
\rqSet{rqUbDivStableLibPZeroDen}{0}
\rqSet{rqUbDivStableLibPZeroPct}{n/a}
\rqSet{rqUbDivStableLibPOneNum}{0}
\rqSet{rqUbDivStableLibPOneDen}{0}
\rqSet{rqUbDivStableLibPOnePct}{n/a}
\rqSet{rqUbInputsLibAllNum}{34208}
\rqSet{rqUbInputsLibAllDen}{51636}
\rqSet{rqUbInputsLibAllPct}{66.2}
\rqSet{rqUbInputsLibBOneNum}{16099}
\rqSet{rqUbInputsLibBOneDen}{24253}
\rqSet{rqUbInputsLibBOnePct}{66.4}
\rqSet{rqUbInputsLibBTwoNum}{18109}
\rqSet{rqUbInputsLibBTwoDen}{27383}
\rqSet{rqUbInputsLibBTwoPct}{66.1}
\rqSet{rqUbInputsLibPZeroNum}{0}
\rqSet{rqUbInputsLibPZeroDen}{0}
\rqSet{rqUbInputsLibPZeroPct}{n/a}
\rqSet{rqUbInputsLibPOneNum}{0}
\rqSet{rqUbInputsLibPOneDen}{0}
\rqSet{rqUbInputsLibPOnePct}{n/a}
\rqSet{rqUbInputsStableLibAllNum}{30611}
\rqSet{rqUbInputsStableLibAllDen}{45994}
\rqSet{rqUbInputsStableLibAllPct}{66.6}
\rqSet{rqUbInputsStableLibBOneNum}{14253}
\rqSet{rqUbInputsStableLibBOneDen}{20363}
\rqSet{rqUbInputsStableLibBOnePct}{70.0}
\rqSet{rqUbInputsStableLibBTwoNum}{16358}
\rqSet{rqUbInputsStableLibBTwoDen}{25631}
\rqSet{rqUbInputsStableLibBTwoPct}{63.8}
\rqSet{rqUbInputsStableLibPZeroNum}{0}
\rqSet{rqUbInputsStableLibPZeroDen}{0}
\rqSet{rqUbInputsStableLibPZeroPct}{n/a}
\rqSet{rqUbInputsStableLibPOneNum}{0}
\rqSet{rqUbInputsStableLibPOneDen}{0}
\rqSet{rqUbInputsStableLibPOnePct}{n/a}
\rqSet{rqUbTransLibAllNum}{197}
\rqSet{rqUbTransLibAllDen}{254}
\rqSet{rqUbTransLibAllPct}{77.6}
\rqSet{rqUbTransLibBOneNum}{97}
\rqSet{rqUbTransLibBOneDen}{138}
\rqSet{rqUbTransLibBOnePct}{70.3}
\rqSet{rqUbTransLibBTwoNum}{100}
\rqSet{rqUbTransLibBTwoDen}{116}
\rqSet{rqUbTransLibBTwoPct}{86.2}
\rqSet{rqUbTransLibPZeroNum}{0}
\rqSet{rqUbTransLibPZeroDen}{0}
\rqSet{rqUbTransLibPZeroPct}{n/a}
\rqSet{rqUbTransLibPOneNum}{0}
\rqSet{rqUbTransLibPOneDen}{0}
\rqSet{rqUbTransLibPOnePct}{n/a}
\rqSet{rqStderrOnlyDivLibAllNum}{798}
\rqSet{rqStderrOnlyDivLibAllDen}{18194}
\rqSet{rqStderrOnlyDivLibAllPct}{4.4}
\rqSet{rqStderrOnlyDivLibAllKiroNum}{7}
\rqSet{rqStderrOnlyDivLibAllKiroDen}{825}
\rqSet{rqStderrOnlyDivLibAllKiroPct}{0.8}
\rqSet{rqStderrOnlyDivLibAllClaudeNum}{0}
\rqSet{rqStderrOnlyDivLibAllClaudeDen}{1239}
\rqSet{rqStderrOnlyDivLibAllClaudePct}{0.0}
\rqSet{rqStderrOnlyDivLibAllCodexNum}{0}
\rqSet{rqStderrOnlyDivLibAllCodexDen}{234}
\rqSet{rqStderrOnlyDivLibAllCodexPct}{0.0}
\rqSet{rqStderrOnlyDivLibAllCTwoRustNum}{0}
\rqSet{rqStderrOnlyDivLibAllCTwoRustDen}{247}
\rqSet{rqStderrOnlyDivLibAllCTwoRustPct}{0.0}
\rqSet{rqStderrOnlyDivLibAllLaertesNum}{0}
\rqSet{rqStderrOnlyDivLibAllLaertesDen}{247}
\rqSet{rqStderrOnlyDivLibAllLaertesPct}{0.0}
\rqSet{rqStderrOnlyDivLibAllCTwoSaferRustNum}{18}
\rqSet{rqStderrOnlyDivLibAllCTwoSaferRustDen}{3731}
\rqSet{rqStderrOnlyDivLibAllCTwoSaferRustPct}{0.5}
\rqSet{rqStderrOnlyDivLibAllSmartCTwoRustNum}{0}
\rqSet{rqStderrOnlyDivLibAllSmartCTwoRustDen}{0}
\rqSet{rqStderrOnlyDivLibAllSmartCTwoRustPct}{n/a}
\rqSet{rqStderrOnlyDivLibAllKimiNum}{0}
\rqSet{rqStderrOnlyDivLibAllKimiDen}{5739}
\rqSet{rqStderrOnlyDivLibAllKimiPct}{0.0}
\rqSet{rqStderrOnlyDivLibAllGptFiveNum}{773}
\rqSet{rqStderrOnlyDivLibAllGptFiveDen}{4602}
\rqSet{rqStderrOnlyDivLibAllGptFivePct}{16.8}
\rqSet{rqStderrOnlyDivLibAllGeminiNum}{0}
\rqSet{rqStderrOnlyDivLibAllGeminiDen}{1330}
\rqSet{rqStderrOnlyDivLibAllGeminiPct}{0.0}
\rqSet{rqStderrOnlyDivLibBOneNum}{0}
\rqSet{rqStderrOnlyDivLibBOneDen}{10507}
\rqSet{rqStderrOnlyDivLibBOnePct}{0.0}
\rqSet{rqStderrOnlyDivLibBTwoNum}{798}
\rqSet{rqStderrOnlyDivLibBTwoDen}{7687}
\rqSet{rqStderrOnlyDivLibBTwoPct}{10.4}
\rqSet{rqStderrOnlyDivLibPZeroNum}{0}
\rqSet{rqStderrOnlyDivLibPZeroDen}{0}
\rqSet{rqStderrOnlyDivLibPZeroPct}{n/a}
\rqSet{rqStderrOnlyDivLibPOneNum}{0}
\rqSet{rqStderrOnlyDivLibPOneDen}{0}
\rqSet{rqStderrOnlyDivLibPOnePct}{n/a}
\rqSet{rqStderrOnlyTransLibAllNum}{0}
\rqSet{rqStderrOnlyTransLibAllDen}{65}
\rqSet{rqStderrOnlyTransLibAllPct}{0.0}
\rqSet{rqStderrOnlyTransLibBOneNum}{0}
\rqSet{rqStderrOnlyTransLibBOneDen}{38}
\rqSet{rqStderrOnlyTransLibBOnePct}{0.0}
\rqSet{rqStderrOnlyTransLibBTwoNum}{0}
\rqSet{rqStderrOnlyTransLibBTwoDen}{27}
\rqSet{rqStderrOnlyTransLibBTwoPct}{0.0}
\rqSet{rqStderrOnlyTransLibPZeroNum}{0}
\rqSet{rqStderrOnlyTransLibPZeroDen}{0}
\rqSet{rqStderrOnlyTransLibPZeroPct}{n/a}
\rqSet{rqStderrOnlyTransLibPOneNum}{0}
\rqSet{rqStderrOnlyTransLibPOneDen}{0}
\rqSet{rqStderrOnlyTransLibPOnePct}{n/a}
\rqSet{rqStderrOnlyInputsLibAllNum}{3774}
\rqSet{rqStderrOnlyInputsLibAllDen}{18194}
\rqSet{rqStderrOnlyInputsLibBOneNum}{0}
\rqSet{rqStderrOnlyInputsLibBOneDen}{10507}
\rqSet{rqStderrOnlyInputsLibBTwoNum}{3774}
\rqSet{rqStderrOnlyInputsLibBTwoDen}{7687}
\rqSet{rqStderrOnlyInputsLibPZeroNum}{0}
\rqSet{rqStderrOnlyInputsLibPZeroDen}{0}
\rqSet{rqStderrOnlyInputsLibPOneNum}{0}
\rqSet{rqStderrOnlyInputsLibPOneDen}{0}
\rqSet{rqBothNtMatchLibAllNum}{12422}
\rqSet{rqBothNtMatchLibAllDen}{6554751}
\rqSet{rqBothNtMatchLibAllPct}{0.2}
\rqSet{rqBothNtMatchLibBOneNum}{0}
\rqSet{rqBothNtMatchLibBOneDen}{3259995}
\rqSet{rqBothNtMatchLibBOnePct}{0.0}
\rqSet{rqBothNtMatchLibBTwoNum}{12422}
\rqSet{rqBothNtMatchLibBTwoDen}{2280249}
\rqSet{rqBothNtMatchLibBTwoPct}{0.5}
\rqSet{rqBothNtMatchLibPZeroNum}{0}
\rqSet{rqBothNtMatchLibPZeroDen}{0}
\rqSet{rqBothNtMatchLibPZeroPct}{n/a}
\rqSet{rqBothNtMatchLibPOneNum}{0}
\rqSet{rqBothNtMatchLibPOneDen}{1014507}
\rqSet{rqBothNtMatchLibPOnePct}{0.0}
\rqSet{rqBothNtPrefixZeroLibAllNum}{0}
\rqSet{rqBothNtPrefixZeroLibAllDen}{12422}
\rqSet{rqBothNtPrefixZeroLibAllPct}{0.0}
\rqSet{rqBothNtPrefixZeroLibBOneNum}{0}
\rqSet{rqBothNtPrefixZeroLibBOneDen}{0}
\rqSet{rqBothNtPrefixZeroLibBOnePct}{n/a}
\rqSet{rqBothNtPrefixZeroLibBTwoNum}{0}
\rqSet{rqBothNtPrefixZeroLibBTwoDen}{12422}
\rqSet{rqBothNtPrefixZeroLibBTwoPct}{0.0}
\rqSet{rqBothNtPrefixZeroLibPZeroNum}{0}
\rqSet{rqBothNtPrefixZeroLibPZeroDen}{0}
\rqSet{rqBothNtPrefixZeroLibPZeroPct}{n/a}
\rqSet{rqBothNtPrefixZeroLibPOneNum}{0}
\rqSet{rqBothNtPrefixZeroLibPOneDen}{0}
\rqSet{rqBothNtPrefixZeroLibPOnePct}{n/a}
\rqSet{rqBothNtPrefixBytesLibAllNum}{0}
\rqSet{rqBothNtPrefixBytesLibAllDen}{12422}
\rqSet{rqBothNtPrefixBytesLibBOneNum}{0}
\rqSet{rqBothNtPrefixBytesLibBOneDen}{0}
\rqSet{rqBothNtPrefixBytesLibBTwoNum}{0}
\rqSet{rqBothNtPrefixBytesLibBTwoDen}{12422}
\rqSet{rqBothNtPrefixBytesLibPZeroNum}{0}
\rqSet{rqBothNtPrefixBytesLibPZeroDen}{0}
\rqSet{rqBothNtPrefixBytesLibPOneNum}{0}
\rqSet{rqBothNtPrefixBytesLibPOneDen}{0}
\rqSet{rqNtOneSideConfirmedLibAllNum}{22}
\rqSet{rqNtOneSideConfirmedLibAllDen}{18194}
\rqSet{rqNtOneSideConfirmedLibAllPct}{0.1}
\rqSet{rqNtOneSideConfirmedLibBOneNum}{0}
\rqSet{rqNtOneSideConfirmedLibBOneDen}{10507}
\rqSet{rqNtOneSideConfirmedLibBOnePct}{0.0}
\rqSet{rqNtOneSideConfirmedLibBTwoNum}{22}
\rqSet{rqNtOneSideConfirmedLibBTwoDen}{7687}
\rqSet{rqNtOneSideConfirmedLibBTwoPct}{0.3}
\rqSet{rqNtOneSideConfirmedLibPZeroNum}{0}
\rqSet{rqNtOneSideConfirmedLibPZeroDen}{0}
\rqSet{rqNtOneSideConfirmedLibPZeroPct}{n/a}
\rqSet{rqNtOneSideConfirmedLibPOneNum}{0}
\rqSet{rqNtOneSideConfirmedLibPOneDen}{0}
\rqSet{rqNtOneSideConfirmedLibPOnePct}{n/a}
\rqSet{rqNtOneSideConfirmedAllLibAllNum}{6335}
\rqSet{rqNtOneSideConfirmedAllLibAllDen}{199794}
\rqSet{rqNtOneSideConfirmedAllLibAllPct}{3.2}
\rqSet{rqNtOneSideConfirmedAllLibBOneNum}{3835}
\rqSet{rqNtOneSideConfirmedAllLibBOneDen}{106695}
\rqSet{rqNtOneSideConfirmedAllLibBOnePct}{3.6}
\rqSet{rqNtOneSideConfirmedAllLibBTwoNum}{2500}
\rqSet{rqNtOneSideConfirmedAllLibBTwoDen}{93099}
\rqSet{rqNtOneSideConfirmedAllLibBTwoPct}{2.7}
\rqSet{rqNtOneSideConfirmedAllLibPZeroNum}{0}
\rqSet{rqNtOneSideConfirmedAllLibPZeroDen}{0}
\rqSet{rqNtOneSideConfirmedAllLibPZeroPct}{n/a}
\rqSet{rqNtOneSideConfirmedAllLibPOneNum}{0}
\rqSet{rqNtOneSideConfirmedAllLibPOneDen}{0}
\rqSet{rqNtOneSideConfirmedAllLibPOnePct}{n/a}
\rqSet{rqNtOneSideRerunAllLibAllNum}{6389}
\rqSet{rqNtOneSideRerunAllLibAllDen}{199794}
\rqSet{rqNtOneSideRerunAllLibAllPct}{3.2}
\rqSet{rqNtOneSideRerunAllLibBOneNum}{3835}
\rqSet{rqNtOneSideRerunAllLibBOneDen}{106695}
\rqSet{rqNtOneSideRerunAllLibBOnePct}{3.6}
\rqSet{rqNtOneSideRerunAllLibBTwoNum}{2554}
\rqSet{rqNtOneSideRerunAllLibBTwoDen}{93099}
\rqSet{rqNtOneSideRerunAllLibBTwoPct}{2.7}
\rqSet{rqNtOneSideRerunAllLibPZeroNum}{0}
\rqSet{rqNtOneSideRerunAllLibPZeroDen}{0}
\rqSet{rqNtOneSideRerunAllLibPZeroPct}{n/a}
\rqSet{rqNtOneSideRerunAllLibPOneNum}{0}
\rqSet{rqNtOneSideRerunAllLibPOneDen}{0}
\rqSet{rqNtOneSideRerunAllLibPOnePct}{n/a}
\rqSet{rqNtOneSideTransAnyLibAllNum}{3}
\rqSet{rqNtOneSideTransAnyLibAllDen}{65}
\rqSet{rqNtOneSideTransAnyLibAllPct}{4.6}
\rqSet{rqNtOneSideTransAnyLibBOneNum}{0}
\rqSet{rqNtOneSideTransAnyLibBOneDen}{38}
\rqSet{rqNtOneSideTransAnyLibBOnePct}{0.0}
\rqSet{rqNtOneSideTransAnyLibBTwoNum}{3}
\rqSet{rqNtOneSideTransAnyLibBTwoDen}{27}
\rqSet{rqNtOneSideTransAnyLibBTwoPct}{11.1}
\rqSet{rqNtOneSideTransAnyLibPZeroNum}{0}
\rqSet{rqNtOneSideTransAnyLibPZeroDen}{0}
\rqSet{rqNtOneSideTransAnyLibPZeroPct}{n/a}
\rqSet{rqNtOneSideTransAnyLibPOneNum}{0}
\rqSet{rqNtOneSideTransAnyLibPOneDen}{0}
\rqSet{rqNtOneSideTransAnyLibPOnePct}{n/a}
\rqSet{rqNtOneSideTransOnlyLibAllNum}{0}
\rqSet{rqNtOneSideTransOnlyLibAllDen}{65}
\rqSet{rqNtOneSideTransOnlyLibAllPct}{0.0}
\rqSet{rqNtOneSideTransOnlyLibBOneNum}{0}
\rqSet{rqNtOneSideTransOnlyLibBOneDen}{38}
\rqSet{rqNtOneSideTransOnlyLibBOnePct}{0.0}
\rqSet{rqNtOneSideTransOnlyLibBTwoNum}{0}
\rqSet{rqNtOneSideTransOnlyLibBTwoDen}{27}
\rqSet{rqNtOneSideTransOnlyLibBTwoPct}{0.0}
\rqSet{rqNtOneSideTransOnlyLibPZeroNum}{0}
\rqSet{rqNtOneSideTransOnlyLibPZeroDen}{0}
\rqSet{rqNtOneSideTransOnlyLibPZeroPct}{n/a}
\rqSet{rqNtOneSideTransOnlyLibPOneNum}{0}
\rqSet{rqNtOneSideTransOnlyLibPOneDen}{0}
\rqSet{rqNtOneSideTransOnlyLibPOnePct}{n/a}
\rqSet{rqDivLayerVectorLibAllNum}{0}
\rqSet{rqDivLayerVectorLibAllDen}{26751}
\rqSet{rqDivLayerVectorLibAllPct}{0.0}
\rqSet{rqDivLayerVectorLibBOneNum}{0}
\rqSet{rqDivLayerVectorLibBOneDen}{19995}
\rqSet{rqDivLayerVectorLibBOnePct}{0.0}
\rqSet{rqDivLayerVectorLibBTwoNum}{0}
\rqSet{rqDivLayerVectorLibBTwoDen}{6249}
\rqSet{rqDivLayerVectorLibBTwoPct}{0.0}
\rqSet{rqDivLayerVectorLibPZeroNum}{0}
\rqSet{rqDivLayerVectorLibPZeroDen}{0}
\rqSet{rqDivLayerVectorLibPZeroPct}{n/a}
\rqSet{rqDivLayerVectorLibPOneNum}{0}
\rqSet{rqDivLayerVectorLibPOneDen}{507}
\rqSet{rqDivLayerVectorLibPOnePct}{0.0}
\rqSet{rqDivInputsLayerVectorLibAllNum}{0}
\rqSet{rqDivInputsLayerVectorLibAllDen}{3459}
\rqSet{rqDivInputsLayerVectorLibAllPct}{0.0}
\rqSet{rqDivInputsLayerVectorLibBOneNum}{0}
\rqSet{rqDivInputsLayerVectorLibBOneDen}{2388}
\rqSet{rqDivInputsLayerVectorLibBOnePct}{0.0}
\rqSet{rqDivInputsLayerVectorLibBTwoNum}{0}
\rqSet{rqDivInputsLayerVectorLibBTwoDen}{843}
\rqSet{rqDivInputsLayerVectorLibBTwoPct}{0.0}
\rqSet{rqDivInputsLayerVectorLibPZeroNum}{0}
\rqSet{rqDivInputsLayerVectorLibPZeroDen}{0}
\rqSet{rqDivInputsLayerVectorLibPZeroPct}{n/a}
\rqSet{rqDivInputsLayerVectorLibPOneNum}{0}
\rqSet{rqDivInputsLayerVectorLibPOneDen}{228}
\rqSet{rqDivInputsLayerVectorLibPOnePct}{0.0}
\rqSet{rqDivTransLayerVectorLibAllNum}{0}
\rqSet{rqDivTransLayerVectorLibAllDen}{65}
\rqSet{rqDivTransLayerVectorLibAllPct}{0.0}
\rqSet{rqDivTransLayerVectorLibBOneNum}{0}
\rqSet{rqDivTransLayerVectorLibBOneDen}{38}
\rqSet{rqDivTransLayerVectorLibBOnePct}{0.0}
\rqSet{rqDivTransLayerVectorLibBTwoNum}{0}
\rqSet{rqDivTransLayerVectorLibBTwoDen}{27}
\rqSet{rqDivTransLayerVectorLibBTwoPct}{0.0}
\rqSet{rqDivTransLayerVectorLibPZeroNum}{0}
\rqSet{rqDivTransLayerVectorLibPZeroDen}{0}
\rqSet{rqDivTransLayerVectorLibPZeroPct}{n/a}
\rqSet{rqDivTransLayerVectorLibPOneNum}{0}
\rqSet{rqDivTransLayerVectorLibPOneDen}{0}
\rqSet{rqDivTransLayerVectorLibPOnePct}{n/a}
\rqSet{rqOutOfDomainVectorLibAllNum}{0}
\rqSet{rqOutOfDomainVectorLibAllDen}{3841}
\rqSet{rqOutOfDomainVectorLibAllPct}{0.0}
\rqSet{rqOutOfDomainVectorLibBOneNum}{0}
\rqSet{rqOutOfDomainVectorLibBOneDen}{2680}
\rqSet{rqOutOfDomainVectorLibBOnePct}{0.0}
\rqSet{rqOutOfDomainVectorLibBTwoNum}{0}
\rqSet{rqOutOfDomainVectorLibBTwoDen}{933}
\rqSet{rqOutOfDomainVectorLibBTwoPct}{0.0}
\rqSet{rqOutOfDomainVectorLibPZeroNum}{0}
\rqSet{rqOutOfDomainVectorLibPZeroDen}{0}
\rqSet{rqOutOfDomainVectorLibPZeroPct}{n/a}
\rqSet{rqOutOfDomainVectorLibPOneNum}{0}
\rqSet{rqOutOfDomainVectorLibPOneDen}{228}
\rqSet{rqOutOfDomainVectorLibPOnePct}{0.0}
\rqSet{rqDivLayerRandomLibAllNum}{13403}
\rqSet{rqDivLayerRandomLibAllDen}{4933776}
\rqSet{rqDivLayerRandomLibAllPct}{0.3}
\rqSet{rqDivLayerRandomLibBOneNum}{8363}
\rqSet{rqDivLayerRandomLibBOneDen}{2650673}
\rqSet{rqDivLayerRandomLibBOnePct}{0.3}
\rqSet{rqDivLayerRandomLibBTwoNum}{5040}
\rqSet{rqDivLayerRandomLibBTwoDen}{1670503}
\rqSet{rqDivLayerRandomLibBTwoPct}{0.3}
\rqSet{rqDivLayerRandomLibPZeroNum}{0}
\rqSet{rqDivLayerRandomLibPZeroDen}{0}
\rqSet{rqDivLayerRandomLibPZeroPct}{n/a}
\rqSet{rqDivLayerRandomLibPOneNum}{0}
\rqSet{rqDivLayerRandomLibPOneDen}{612600}
\rqSet{rqDivLayerRandomLibPOnePct}{0.0}
\rqSet{rqDivInputsLayerRandomLibAllNum}{8786}
\rqSet{rqDivInputsLayerRandomLibAllDen}{815767}
\rqSet{rqDivInputsLayerRandomLibAllPct}{1.1}
\rqSet{rqDivInputsLayerRandomLibBOneNum}{4780}
\rqSet{rqDivInputsLayerRandomLibBOneDen}{318157}
\rqSet{rqDivInputsLayerRandomLibBOnePct}{1.5}
\rqSet{rqDivInputsLayerRandomLibBTwoNum}{4006}
\rqSet{rqDivInputsLayerRandomLibBTwoDen}{224010}
\rqSet{rqDivInputsLayerRandomLibBTwoPct}{1.8}
\rqSet{rqDivInputsLayerRandomLibPZeroNum}{0}
\rqSet{rqDivInputsLayerRandomLibPZeroDen}{0}
\rqSet{rqDivInputsLayerRandomLibPZeroPct}{n/a}
\rqSet{rqDivInputsLayerRandomLibPOneNum}{0}
\rqSet{rqDivInputsLayerRandomLibPOneDen}{273600}
\rqSet{rqDivInputsLayerRandomLibPOnePct}{0.0}
\rqSet{rqDivTransLayerRandomLibAllNum}{59}
\rqSet{rqDivTransLayerRandomLibAllDen}{65}
\rqSet{rqDivTransLayerRandomLibAllPct}{90.8}
\rqSet{rqDivTransLayerRandomLibBOneNum}{37}
\rqSet{rqDivTransLayerRandomLibBOneDen}{38}
\rqSet{rqDivTransLayerRandomLibBOnePct}{97.4}
\rqSet{rqDivTransLayerRandomLibBTwoNum}{22}
\rqSet{rqDivTransLayerRandomLibBTwoDen}{27}
\rqSet{rqDivTransLayerRandomLibBTwoPct}{81.5}
\rqSet{rqDivTransLayerRandomLibPZeroNum}{0}
\rqSet{rqDivTransLayerRandomLibPZeroDen}{0}
\rqSet{rqDivTransLayerRandomLibPZeroPct}{n/a}
\rqSet{rqDivTransLayerRandomLibPOneNum}{0}
\rqSet{rqDivTransLayerRandomLibPOneDen}{0}
\rqSet{rqDivTransLayerRandomLibPOnePct}{n/a}
\rqSet{rqOutOfDomainRandomLibAllNum}{0}
\rqSet{rqOutOfDomainRandomLibAllDen}{885352}
\rqSet{rqOutOfDomainRandomLibAllPct}{0.0}
\rqSet{rqOutOfDomainRandomLibBOneNum}{0}
\rqSet{rqOutOfDomainRandomLibBOneDen}{357982}
\rqSet{rqOutOfDomainRandomLibBOnePct}{0.0}
\rqSet{rqOutOfDomainRandomLibBTwoNum}{0}
\rqSet{rqOutOfDomainRandomLibBTwoDen}{253770}
\rqSet{rqOutOfDomainRandomLibBTwoPct}{0.0}
\rqSet{rqOutOfDomainRandomLibPZeroNum}{0}
\rqSet{rqOutOfDomainRandomLibPZeroDen}{0}
\rqSet{rqOutOfDomainRandomLibPZeroPct}{n/a}
\rqSet{rqOutOfDomainRandomLibPOneNum}{0}
\rqSet{rqOutOfDomainRandomLibPOneDen}{273600}
\rqSet{rqOutOfDomainRandomLibPOnePct}{0.0}
\rqSet{rqDivLayerDictLibAllNum}{4791}
\rqSet{rqDivLayerDictLibAllDen}{1496973}
\rqSet{rqDivLayerDictLibAllPct}{0.3}
\rqSet{rqDivLayerDictLibBOneNum}{2144}
\rqSet{rqDivLayerDictLibBOneDen}{589303}
\rqSet{rqDivLayerDictLibBOnePct}{0.4}
\rqSet{rqDivLayerDictLibBTwoNum}{2647}
\rqSet{rqDivLayerDictLibBTwoDen}{603470}
\rqSet{rqDivLayerDictLibBTwoPct}{0.4}
\rqSet{rqDivLayerDictLibPZeroNum}{0}
\rqSet{rqDivLayerDictLibPZeroDen}{0}
\rqSet{rqDivLayerDictLibPZeroPct}{n/a}
\rqSet{rqDivLayerDictLibPOneNum}{0}
\rqSet{rqDivLayerDictLibPOneDen}{304200}
\rqSet{rqDivLayerDictLibPOnePct}{0.0}
\rqSet{rqDivInputsLayerDictLibAllNum}{3571}
\rqSet{rqDivInputsLayerDictLibAllDen}{290627}
\rqSet{rqDivInputsLayerDictLibAllPct}{1.2}
\rqSet{rqDivInputsLayerDictLibBOneNum}{1330}
\rqSet{rqDivInputsLayerDictLibBOneDen}{71840}
\rqSet{rqDivInputsLayerDictLibBOnePct}{1.9}
\rqSet{rqDivInputsLayerDictLibBTwoNum}{2241}
\rqSet{rqDivInputsLayerDictLibBTwoDen}{81987}
\rqSet{rqDivInputsLayerDictLibBTwoPct}{2.7}
\rqSet{rqDivInputsLayerDictLibPZeroNum}{0}
\rqSet{rqDivInputsLayerDictLibPZeroDen}{0}
\rqSet{rqDivInputsLayerDictLibPZeroPct}{n/a}
\rqSet{rqDivInputsLayerDictLibPOneNum}{0}
\rqSet{rqDivInputsLayerDictLibPOneDen}{136800}
\rqSet{rqDivInputsLayerDictLibPOnePct}{0.0}
\rqSet{rqDivTransLayerDictLibAllNum}{37}
\rqSet{rqDivTransLayerDictLibAllDen}{65}
\rqSet{rqDivTransLayerDictLibAllPct}{56.9}
\rqSet{rqDivTransLayerDictLibBOneNum}{16}
\rqSet{rqDivTransLayerDictLibBOneDen}{38}
\rqSet{rqDivTransLayerDictLibBOnePct}{42.1}
\rqSet{rqDivTransLayerDictLibBTwoNum}{21}
\rqSet{rqDivTransLayerDictLibBTwoDen}{27}
\rqSet{rqDivTransLayerDictLibBTwoPct}{77.8}
\rqSet{rqDivTransLayerDictLibPZeroNum}{0}
\rqSet{rqDivTransLayerDictLibPZeroDen}{0}
\rqSet{rqDivTransLayerDictLibPZeroPct}{n/a}
\rqSet{rqDivTransLayerDictLibPOneNum}{0}
\rqSet{rqDivTransLayerDictLibPOneDen}{0}
\rqSet{rqDivTransLayerDictLibPOnePct}{n/a}
\rqSet{rqOutOfDomainDictLibAllNum}{0}
\rqSet{rqOutOfDomainDictLibAllDen}{301038}
\rqSet{rqOutOfDomainDictLibAllPct}{0.0}
\rqSet{rqOutOfDomainDictLibBOneNum}{0}
\rqSet{rqOutOfDomainDictLibBOneDen}{74014}
\rqSet{rqOutOfDomainDictLibBOnePct}{0.0}
\rqSet{rqOutOfDomainDictLibBTwoNum}{0}
\rqSet{rqOutOfDomainDictLibBTwoDen}{90224}
\rqSet{rqOutOfDomainDictLibBTwoPct}{0.0}
\rqSet{rqOutOfDomainDictLibPZeroNum}{0}
\rqSet{rqOutOfDomainDictLibPZeroDen}{0}
\rqSet{rqOutOfDomainDictLibPZeroPct}{n/a}
\rqSet{rqOutOfDomainDictLibPOneNum}{0}
\rqSet{rqOutOfDomainDictLibPOneDen}{136800}
\rqSet{rqOutOfDomainDictLibPOnePct}{0.0}
\rqSet{rqDivLayerSweepLibAllNum}{0}
\rqSet{rqDivLayerSweepLibAllDen}{97200}
\rqSet{rqDivLayerSweepLibAllPct}{0.0}
\rqSet{rqDivLayerSweepLibBOneNum}{0}
\rqSet{rqDivLayerSweepLibBOneDen}{0}
\rqSet{rqDivLayerSweepLibBOnePct}{n/a}
\rqSet{rqDivLayerSweepLibBTwoNum}{0}
\rqSet{rqDivLayerSweepLibBTwoDen}{0}
\rqSet{rqDivLayerSweepLibBTwoPct}{n/a}
\rqSet{rqDivLayerSweepLibPZeroNum}{0}
\rqSet{rqDivLayerSweepLibPZeroDen}{0}
\rqSet{rqDivLayerSweepLibPZeroPct}{n/a}
\rqSet{rqDivLayerSweepLibPOneNum}{0}
\rqSet{rqDivLayerSweepLibPOneDen}{97200}
\rqSet{rqDivLayerSweepLibPOnePct}{0.0}
\rqSet{rqDivInputsLayerSweepLibAllNum}{0}
\rqSet{rqDivInputsLayerSweepLibAllDen}{45600}
\rqSet{rqDivInputsLayerSweepLibAllPct}{0.0}
\rqSet{rqDivInputsLayerSweepLibBOneNum}{0}
\rqSet{rqDivInputsLayerSweepLibBOneDen}{0}
\rqSet{rqDivInputsLayerSweepLibBOnePct}{n/a}
\rqSet{rqDivInputsLayerSweepLibBTwoNum}{0}
\rqSet{rqDivInputsLayerSweepLibBTwoDen}{0}
\rqSet{rqDivInputsLayerSweepLibBTwoPct}{n/a}
\rqSet{rqDivInputsLayerSweepLibPZeroNum}{0}
\rqSet{rqDivInputsLayerSweepLibPZeroDen}{0}
\rqSet{rqDivInputsLayerSweepLibPZeroPct}{n/a}
\rqSet{rqDivInputsLayerSweepLibPOneNum}{0}
\rqSet{rqDivInputsLayerSweepLibPOneDen}{45600}
\rqSet{rqDivInputsLayerSweepLibPOnePct}{0.0}
\rqSet{rqDivTransLayerSweepLibAllNum}{0}
\rqSet{rqDivTransLayerSweepLibAllDen}{65}
\rqSet{rqDivTransLayerSweepLibAllPct}{0.0}
\rqSet{rqDivTransLayerSweepLibBOneNum}{0}
\rqSet{rqDivTransLayerSweepLibBOneDen}{38}
\rqSet{rqDivTransLayerSweepLibBOnePct}{0.0}
\rqSet{rqDivTransLayerSweepLibBTwoNum}{0}
\rqSet{rqDivTransLayerSweepLibBTwoDen}{27}
\rqSet{rqDivTransLayerSweepLibBTwoPct}{0.0}
\rqSet{rqDivTransLayerSweepLibPZeroNum}{0}
\rqSet{rqDivTransLayerSweepLibPZeroDen}{0}
\rqSet{rqDivTransLayerSweepLibPZeroPct}{n/a}
\rqSet{rqDivTransLayerSweepLibPOneNum}{0}
\rqSet{rqDivTransLayerSweepLibPOneDen}{0}
\rqSet{rqDivTransLayerSweepLibPOnePct}{n/a}
\rqSet{rqOutOfDomainSweepLibAllNum}{0}
\rqSet{rqOutOfDomainSweepLibAllDen}{45600}
\rqSet{rqOutOfDomainSweepLibAllPct}{0.0}
\rqSet{rqOutOfDomainSweepLibBOneNum}{0}
\rqSet{rqOutOfDomainSweepLibBOneDen}{0}
\rqSet{rqOutOfDomainSweepLibBOnePct}{n/a}
\rqSet{rqOutOfDomainSweepLibBTwoNum}{0}
\rqSet{rqOutOfDomainSweepLibBTwoDen}{0}
\rqSet{rqOutOfDomainSweepLibBTwoPct}{n/a}
\rqSet{rqOutOfDomainSweepLibPZeroNum}{0}
\rqSet{rqOutOfDomainSweepLibPZeroDen}{0}
\rqSet{rqOutOfDomainSweepLibPZeroPct}{n/a}
\rqSet{rqOutOfDomainSweepLibPOneNum}{0}
\rqSet{rqOutOfDomainSweepLibPOneDen}{45600}
\rqSet{rqOutOfDomainSweepLibPOnePct}{0.0}
\rqSet{rqDivLayerGeneratedLibAllNum}{18194}
\rqSet{rqDivLayerGeneratedLibAllDen}{6528000}
\rqSet{rqDivLayerGeneratedLibAllPct}{0.3}
\rqSet{rqDivLayerGeneratedLibBOneNum}{10507}
\rqSet{rqDivLayerGeneratedLibBOneDen}{3240000}
\rqSet{rqDivLayerGeneratedLibBOnePct}{0.3}
\rqSet{rqDivLayerGeneratedLibBTwoNum}{7687}
\rqSet{rqDivLayerGeneratedLibBTwoDen}{2274000}
\rqSet{rqDivLayerGeneratedLibBTwoPct}{0.3}
\rqSet{rqDivLayerGeneratedLibPZeroNum}{0}
\rqSet{rqDivLayerGeneratedLibPZeroDen}{0}
\rqSet{rqDivLayerGeneratedLibPZeroPct}{n/a}
\rqSet{rqDivLayerGeneratedLibPOneNum}{0}
\rqSet{rqDivLayerGeneratedLibPOneDen}{1014000}
\rqSet{rqDivLayerGeneratedLibPOnePct}{0.0}
\rqSet{rqDivInputsLayerGeneratedLibAllNum}{12357}
\rqSet{rqDivInputsLayerGeneratedLibAllDen}{1152000}
\rqSet{rqDivInputsLayerGeneratedLibAllPct}{1.1}
\rqSet{rqDivInputsLayerGeneratedLibBOneNum}{6110}
\rqSet{rqDivInputsLayerGeneratedLibBOneDen}{390000}
\rqSet{rqDivInputsLayerGeneratedLibBOnePct}{1.6}
\rqSet{rqDivInputsLayerGeneratedLibBTwoNum}{6247}
\rqSet{rqDivInputsLayerGeneratedLibBTwoDen}{306000}
\rqSet{rqDivInputsLayerGeneratedLibBTwoPct}{2.0}
\rqSet{rqDivInputsLayerGeneratedLibPZeroNum}{0}
\rqSet{rqDivInputsLayerGeneratedLibPZeroDen}{0}
\rqSet{rqDivInputsLayerGeneratedLibPZeroPct}{n/a}
\rqSet{rqDivInputsLayerGeneratedLibPOneNum}{0}
\rqSet{rqDivInputsLayerGeneratedLibPOneDen}{456000}
\rqSet{rqDivInputsLayerGeneratedLibPOnePct}{0.0}
\rqSet{rqDivTransLayerGeneratedLibAllNum}{65}
\rqSet{rqDivTransLayerGeneratedLibAllDen}{65}
\rqSet{rqDivTransLayerGeneratedLibAllPct}{100.0}
\rqSet{rqDivTransLayerGeneratedLibBOneNum}{38}
\rqSet{rqDivTransLayerGeneratedLibBOneDen}{38}
\rqSet{rqDivTransLayerGeneratedLibBOnePct}{100.0}
\rqSet{rqDivTransLayerGeneratedLibBTwoNum}{27}
\rqSet{rqDivTransLayerGeneratedLibBTwoDen}{27}
\rqSet{rqDivTransLayerGeneratedLibBTwoPct}{100.0}
\rqSet{rqDivTransLayerGeneratedLibPZeroNum}{0}
\rqSet{rqDivTransLayerGeneratedLibPZeroDen}{0}
\rqSet{rqDivTransLayerGeneratedLibPZeroPct}{n/a}
\rqSet{rqDivTransLayerGeneratedLibPOneNum}{0}
\rqSet{rqDivTransLayerGeneratedLibPOneDen}{0}
\rqSet{rqDivTransLayerGeneratedLibPOnePct}{n/a}
\rqSet{rqOutOfDomainGeneratedLibAllNum}{0}
\rqSet{rqOutOfDomainGeneratedLibAllDen}{1232000}
\rqSet{rqOutOfDomainGeneratedLibAllPct}{0.0}
\rqSet{rqOutOfDomainGeneratedLibBOneNum}{0}
\rqSet{rqOutOfDomainGeneratedLibBOneDen}{432000}
\rqSet{rqOutOfDomainGeneratedLibBOnePct}{0.0}
\rqSet{rqOutOfDomainGeneratedLibBTwoNum}{0}
\rqSet{rqOutOfDomainGeneratedLibBTwoDen}{344000}
\rqSet{rqOutOfDomainGeneratedLibBTwoPct}{0.0}
\rqSet{rqOutOfDomainGeneratedLibPZeroNum}{0}
\rqSet{rqOutOfDomainGeneratedLibPZeroDen}{0}
\rqSet{rqOutOfDomainGeneratedLibPZeroPct}{n/a}
\rqSet{rqOutOfDomainGeneratedLibPOneNum}{0}
\rqSet{rqOutOfDomainGeneratedLibPOneDen}{456000}
\rqSet{rqOutOfDomainGeneratedLibPOnePct}{0.0}
\rqSet{rqAttrInputsNoInputFieldsLibAllNum}{0}
\rqSet{rqAttrInputsNoInputFieldsLibAllDen}{6939729}
\rqSet{rqAttrInputsNoInputFieldsLibAllPct}{0.0}
\rqSet{rqAttrInputsNoInputFieldsLibAllKiroNum}{0}
\rqSet{rqAttrInputsNoInputFieldsLibAllKiroDen}{1209621}
\rqSet{rqAttrInputsNoInputFieldsLibAllKiroPct}{0.0}
\rqSet{rqAttrInputsNoInputFieldsLibAllClaudeNum}{0}
\rqSet{rqAttrInputsNoInputFieldsLibAllClaudeDen}{1143588}
\rqSet{rqAttrInputsNoInputFieldsLibAllClaudePct}{0.0}
\rqSet{rqAttrInputsNoInputFieldsLibAllCodexNum}{0}
\rqSet{rqAttrInputsNoInputFieldsLibAllCodexDen}{915456}
\rqSet{rqAttrInputsNoInputFieldsLibAllCodexPct}{0.0}
\rqSet{rqAttrInputsNoInputFieldsLibAllCTwoRustNum}{0}
\rqSet{rqAttrInputsNoInputFieldsLibAllCTwoRustDen}{753393}
\rqSet{rqAttrInputsNoInputFieldsLibAllCTwoRustPct}{0.0}
\rqSet{rqAttrInputsNoInputFieldsLibAllLaertesNum}{0}
\rqSet{rqAttrInputsNoInputFieldsLibAllLaertesDen}{741390}
\rqSet{rqAttrInputsNoInputFieldsLibAllLaertesPct}{0.0}
\rqSet{rqAttrInputsNoInputFieldsLibAllCTwoSaferRustNum}{0}
\rqSet{rqAttrInputsNoInputFieldsLibAllCTwoSaferRustDen}{554526}
\rqSet{rqAttrInputsNoInputFieldsLibAllCTwoSaferRustPct}{0.0}
\rqSet{rqAttrInputsNoInputFieldsLibAllSmartCTwoRustNum}{0}
\rqSet{rqAttrInputsNoInputFieldsLibAllSmartCTwoRustDen}{0}
\rqSet{rqAttrInputsNoInputFieldsLibAllSmartCTwoRustPct}{n/a}
\rqSet{rqAttrInputsNoInputFieldsLibAllKimiNum}{0}
\rqSet{rqAttrInputsNoInputFieldsLibAllKimiDen}{464415}
\rqSet{rqAttrInputsNoInputFieldsLibAllKimiPct}{0.0}
\rqSet{rqAttrInputsNoInputFieldsLibAllGptFiveNum}{0}
\rqSet{rqAttrInputsNoInputFieldsLibAllGptFiveDen}{584547}
\rqSet{rqAttrInputsNoInputFieldsLibAllGptFivePct}{0.0}
\rqSet{rqAttrInputsNoInputFieldsLibAllGeminiNum}{0}
\rqSet{rqAttrInputsNoInputFieldsLibAllGeminiDen}{572793}
\rqSet{rqAttrInputsNoInputFieldsLibAllGeminiPct}{0.0}
\rqSet{rqAttrInputsNoInputFieldsLibBOneNum}{0}
\rqSet{rqAttrInputsNoInputFieldsLibBOneDen}{3452556}
\rqSet{rqAttrInputsNoInputFieldsLibBOnePct}{0.0}
\rqSet{rqAttrInputsNoInputFieldsLibBTwoNum}{0}
\rqSet{rqAttrInputsNoInputFieldsLibBTwoDen}{2472666}
\rqSet{rqAttrInputsNoInputFieldsLibBTwoPct}{0.0}
\rqSet{rqAttrInputsNoInputFieldsLibPZeroNum}{0}
\rqSet{rqAttrInputsNoInputFieldsLibPZeroDen}{0}
\rqSet{rqAttrInputsNoInputFieldsLibPZeroPct}{n/a}
\rqSet{rqAttrInputsNoInputFieldsLibPOneNum}{0}
\rqSet{rqAttrInputsNoInputFieldsLibPOneDen}{1014507}
\rqSet{rqAttrInputsNoInputFieldsLibPOnePct}{0.0}
\rqSet{rqAttrInputsCBuildFailedLibAllNum}{0}
\rqSet{rqAttrInputsCBuildFailedLibAllDen}{6939729}
\rqSet{rqAttrInputsCBuildFailedLibAllPct}{0.0}
\rqSet{rqAttrInputsCBuildFailedLibAllKiroNum}{0}
\rqSet{rqAttrInputsCBuildFailedLibAllKiroDen}{1209621}
\rqSet{rqAttrInputsCBuildFailedLibAllKiroPct}{0.0}
\rqSet{rqAttrInputsCBuildFailedLibAllClaudeNum}{0}
\rqSet{rqAttrInputsCBuildFailedLibAllClaudeDen}{1143588}
\rqSet{rqAttrInputsCBuildFailedLibAllClaudePct}{0.0}
\rqSet{rqAttrInputsCBuildFailedLibAllCodexNum}{0}
\rqSet{rqAttrInputsCBuildFailedLibAllCodexDen}{915456}
\rqSet{rqAttrInputsCBuildFailedLibAllCodexPct}{0.0}
\rqSet{rqAttrInputsCBuildFailedLibAllCTwoRustNum}{0}
\rqSet{rqAttrInputsCBuildFailedLibAllCTwoRustDen}{753393}
\rqSet{rqAttrInputsCBuildFailedLibAllCTwoRustPct}{0.0}
\rqSet{rqAttrInputsCBuildFailedLibAllLaertesNum}{0}
\rqSet{rqAttrInputsCBuildFailedLibAllLaertesDen}{741390}
\rqSet{rqAttrInputsCBuildFailedLibAllLaertesPct}{0.0}
\rqSet{rqAttrInputsCBuildFailedLibAllCTwoSaferRustNum}{0}
\rqSet{rqAttrInputsCBuildFailedLibAllCTwoSaferRustDen}{554526}
\rqSet{rqAttrInputsCBuildFailedLibAllCTwoSaferRustPct}{0.0}
\rqSet{rqAttrInputsCBuildFailedLibAllSmartCTwoRustNum}{0}
\rqSet{rqAttrInputsCBuildFailedLibAllSmartCTwoRustDen}{0}
\rqSet{rqAttrInputsCBuildFailedLibAllSmartCTwoRustPct}{n/a}
\rqSet{rqAttrInputsCBuildFailedLibAllKimiNum}{0}
\rqSet{rqAttrInputsCBuildFailedLibAllKimiDen}{464415}
\rqSet{rqAttrInputsCBuildFailedLibAllKimiPct}{0.0}
\rqSet{rqAttrInputsCBuildFailedLibAllGptFiveNum}{0}
\rqSet{rqAttrInputsCBuildFailedLibAllGptFiveDen}{584547}
\rqSet{rqAttrInputsCBuildFailedLibAllGptFivePct}{0.0}
\rqSet{rqAttrInputsCBuildFailedLibAllGeminiNum}{0}
\rqSet{rqAttrInputsCBuildFailedLibAllGeminiDen}{572793}
\rqSet{rqAttrInputsCBuildFailedLibAllGeminiPct}{0.0}
\rqSet{rqAttrInputsCBuildFailedLibBOneNum}{0}
\rqSet{rqAttrInputsCBuildFailedLibBOneDen}{3452556}
\rqSet{rqAttrInputsCBuildFailedLibBOnePct}{0.0}
\rqSet{rqAttrInputsCBuildFailedLibBTwoNum}{0}
\rqSet{rqAttrInputsCBuildFailedLibBTwoDen}{2472666}
\rqSet{rqAttrInputsCBuildFailedLibBTwoPct}{0.0}
\rqSet{rqAttrInputsCBuildFailedLibPZeroNum}{0}
\rqSet{rqAttrInputsCBuildFailedLibPZeroDen}{0}
\rqSet{rqAttrInputsCBuildFailedLibPZeroPct}{n/a}
\rqSet{rqAttrInputsCBuildFailedLibPOneNum}{0}
\rqSet{rqAttrInputsCBuildFailedLibPOneDen}{1014507}
\rqSet{rqAttrInputsCBuildFailedLibPOnePct}{0.0}
\rqSet{rqAttrInputsRustBuildFailedLibAllNum}{0}
\rqSet{rqAttrInputsRustBuildFailedLibAllDen}{6939729}
\rqSet{rqAttrInputsRustBuildFailedLibAllPct}{0.0}
\rqSet{rqAttrInputsRustBuildFailedLibAllKiroNum}{0}
\rqSet{rqAttrInputsRustBuildFailedLibAllKiroDen}{1209621}
\rqSet{rqAttrInputsRustBuildFailedLibAllKiroPct}{0.0}
\rqSet{rqAttrInputsRustBuildFailedLibAllClaudeNum}{0}
\rqSet{rqAttrInputsRustBuildFailedLibAllClaudeDen}{1143588}
\rqSet{rqAttrInputsRustBuildFailedLibAllClaudePct}{0.0}
\rqSet{rqAttrInputsRustBuildFailedLibAllCodexNum}{0}
\rqSet{rqAttrInputsRustBuildFailedLibAllCodexDen}{915456}
\rqSet{rqAttrInputsRustBuildFailedLibAllCodexPct}{0.0}
\rqSet{rqAttrInputsRustBuildFailedLibAllCTwoRustNum}{0}
\rqSet{rqAttrInputsRustBuildFailedLibAllCTwoRustDen}{753393}
\rqSet{rqAttrInputsRustBuildFailedLibAllCTwoRustPct}{0.0}
\rqSet{rqAttrInputsRustBuildFailedLibAllLaertesNum}{0}
\rqSet{rqAttrInputsRustBuildFailedLibAllLaertesDen}{741390}
\rqSet{rqAttrInputsRustBuildFailedLibAllLaertesPct}{0.0}
\rqSet{rqAttrInputsRustBuildFailedLibAllCTwoSaferRustNum}{0}
\rqSet{rqAttrInputsRustBuildFailedLibAllCTwoSaferRustDen}{554526}
\rqSet{rqAttrInputsRustBuildFailedLibAllCTwoSaferRustPct}{0.0}
\rqSet{rqAttrInputsRustBuildFailedLibAllSmartCTwoRustNum}{0}
\rqSet{rqAttrInputsRustBuildFailedLibAllSmartCTwoRustDen}{0}
\rqSet{rqAttrInputsRustBuildFailedLibAllSmartCTwoRustPct}{n/a}
\rqSet{rqAttrInputsRustBuildFailedLibAllKimiNum}{0}
\rqSet{rqAttrInputsRustBuildFailedLibAllKimiDen}{464415}
\rqSet{rqAttrInputsRustBuildFailedLibAllKimiPct}{0.0}
\rqSet{rqAttrInputsRustBuildFailedLibAllGptFiveNum}{0}
\rqSet{rqAttrInputsRustBuildFailedLibAllGptFiveDen}{584547}
\rqSet{rqAttrInputsRustBuildFailedLibAllGptFivePct}{0.0}
\rqSet{rqAttrInputsRustBuildFailedLibAllGeminiNum}{0}
\rqSet{rqAttrInputsRustBuildFailedLibAllGeminiDen}{572793}
\rqSet{rqAttrInputsRustBuildFailedLibAllGeminiPct}{0.0}
\rqSet{rqAttrInputsRustBuildFailedLibBOneNum}{0}
\rqSet{rqAttrInputsRustBuildFailedLibBOneDen}{3452556}
\rqSet{rqAttrInputsRustBuildFailedLibBOnePct}{0.0}
\rqSet{rqAttrInputsRustBuildFailedLibBTwoNum}{0}
\rqSet{rqAttrInputsRustBuildFailedLibBTwoDen}{2472666}
\rqSet{rqAttrInputsRustBuildFailedLibBTwoPct}{0.0}
\rqSet{rqAttrInputsRustBuildFailedLibPZeroNum}{0}
\rqSet{rqAttrInputsRustBuildFailedLibPZeroDen}{0}
\rqSet{rqAttrInputsRustBuildFailedLibPZeroPct}{n/a}
\rqSet{rqAttrInputsRustBuildFailedLibPOneNum}{0}
\rqSet{rqAttrInputsRustBuildFailedLibPOneDen}{1014507}
\rqSet{rqAttrInputsRustBuildFailedLibPOnePct}{0.0}
\rqSet{rqAttrInputsHarnessErrorLibAllNum}{0}
\rqSet{rqAttrInputsHarnessErrorLibAllDen}{6939729}
\rqSet{rqAttrInputsHarnessErrorLibAllPct}{0.0}
\rqSet{rqAttrInputsHarnessErrorLibAllKiroNum}{0}
\rqSet{rqAttrInputsHarnessErrorLibAllKiroDen}{1209621}
\rqSet{rqAttrInputsHarnessErrorLibAllKiroPct}{0.0}
\rqSet{rqAttrInputsHarnessErrorLibAllClaudeNum}{0}
\rqSet{rqAttrInputsHarnessErrorLibAllClaudeDen}{1143588}
\rqSet{rqAttrInputsHarnessErrorLibAllClaudePct}{0.0}
\rqSet{rqAttrInputsHarnessErrorLibAllCodexNum}{0}
\rqSet{rqAttrInputsHarnessErrorLibAllCodexDen}{915456}
\rqSet{rqAttrInputsHarnessErrorLibAllCodexPct}{0.0}
\rqSet{rqAttrInputsHarnessErrorLibAllCTwoRustNum}{0}
\rqSet{rqAttrInputsHarnessErrorLibAllCTwoRustDen}{753393}
\rqSet{rqAttrInputsHarnessErrorLibAllCTwoRustPct}{0.0}
\rqSet{rqAttrInputsHarnessErrorLibAllLaertesNum}{0}
\rqSet{rqAttrInputsHarnessErrorLibAllLaertesDen}{741390}
\rqSet{rqAttrInputsHarnessErrorLibAllLaertesPct}{0.0}
\rqSet{rqAttrInputsHarnessErrorLibAllCTwoSaferRustNum}{0}
\rqSet{rqAttrInputsHarnessErrorLibAllCTwoSaferRustDen}{554526}
\rqSet{rqAttrInputsHarnessErrorLibAllCTwoSaferRustPct}{0.0}
\rqSet{rqAttrInputsHarnessErrorLibAllSmartCTwoRustNum}{0}
\rqSet{rqAttrInputsHarnessErrorLibAllSmartCTwoRustDen}{0}
\rqSet{rqAttrInputsHarnessErrorLibAllSmartCTwoRustPct}{n/a}
\rqSet{rqAttrInputsHarnessErrorLibAllKimiNum}{0}
\rqSet{rqAttrInputsHarnessErrorLibAllKimiDen}{464415}
\rqSet{rqAttrInputsHarnessErrorLibAllKimiPct}{0.0}
\rqSet{rqAttrInputsHarnessErrorLibAllGptFiveNum}{0}
\rqSet{rqAttrInputsHarnessErrorLibAllGptFiveDen}{584547}
\rqSet{rqAttrInputsHarnessErrorLibAllGptFivePct}{0.0}
\rqSet{rqAttrInputsHarnessErrorLibAllGeminiNum}{0}
\rqSet{rqAttrInputsHarnessErrorLibAllGeminiDen}{572793}
\rqSet{rqAttrInputsHarnessErrorLibAllGeminiPct}{0.0}
\rqSet{rqAttrInputsHarnessErrorLibBOneNum}{0}
\rqSet{rqAttrInputsHarnessErrorLibBOneDen}{3452556}
\rqSet{rqAttrInputsHarnessErrorLibBOnePct}{0.0}
\rqSet{rqAttrInputsHarnessErrorLibBTwoNum}{0}
\rqSet{rqAttrInputsHarnessErrorLibBTwoDen}{2472666}
\rqSet{rqAttrInputsHarnessErrorLibBTwoPct}{0.0}
\rqSet{rqAttrInputsHarnessErrorLibPZeroNum}{0}
\rqSet{rqAttrInputsHarnessErrorLibPZeroDen}{0}
\rqSet{rqAttrInputsHarnessErrorLibPZeroPct}{n/a}
\rqSet{rqAttrInputsHarnessErrorLibPOneNum}{0}
\rqSet{rqAttrInputsHarnessErrorLibPOneDen}{1014507}
\rqSet{rqAttrInputsHarnessErrorLibPOnePct}{0.0}
\rqSet{rqAttrInputsNoInputsLibAllNum}{0}
\rqSet{rqAttrInputsNoInputsLibAllDen}{6939729}
\rqSet{rqAttrInputsNoInputsLibAllPct}{0.0}
\rqSet{rqAttrInputsNoInputsLibAllKiroNum}{0}
\rqSet{rqAttrInputsNoInputsLibAllKiroDen}{1209621}
\rqSet{rqAttrInputsNoInputsLibAllKiroPct}{0.0}
\rqSet{rqAttrInputsNoInputsLibAllClaudeNum}{0}
\rqSet{rqAttrInputsNoInputsLibAllClaudeDen}{1143588}
\rqSet{rqAttrInputsNoInputsLibAllClaudePct}{0.0}
\rqSet{rqAttrInputsNoInputsLibAllCodexNum}{0}
\rqSet{rqAttrInputsNoInputsLibAllCodexDen}{915456}
\rqSet{rqAttrInputsNoInputsLibAllCodexPct}{0.0}
\rqSet{rqAttrInputsNoInputsLibAllCTwoRustNum}{0}
\rqSet{rqAttrInputsNoInputsLibAllCTwoRustDen}{753393}
\rqSet{rqAttrInputsNoInputsLibAllCTwoRustPct}{0.0}
\rqSet{rqAttrInputsNoInputsLibAllLaertesNum}{0}
\rqSet{rqAttrInputsNoInputsLibAllLaertesDen}{741390}
\rqSet{rqAttrInputsNoInputsLibAllLaertesPct}{0.0}
\rqSet{rqAttrInputsNoInputsLibAllCTwoSaferRustNum}{0}
\rqSet{rqAttrInputsNoInputsLibAllCTwoSaferRustDen}{554526}
\rqSet{rqAttrInputsNoInputsLibAllCTwoSaferRustPct}{0.0}
\rqSet{rqAttrInputsNoInputsLibAllSmartCTwoRustNum}{0}
\rqSet{rqAttrInputsNoInputsLibAllSmartCTwoRustDen}{0}
\rqSet{rqAttrInputsNoInputsLibAllSmartCTwoRustPct}{n/a}
\rqSet{rqAttrInputsNoInputsLibAllKimiNum}{0}
\rqSet{rqAttrInputsNoInputsLibAllKimiDen}{464415}
\rqSet{rqAttrInputsNoInputsLibAllKimiPct}{0.0}
\rqSet{rqAttrInputsNoInputsLibAllGptFiveNum}{0}
\rqSet{rqAttrInputsNoInputsLibAllGptFiveDen}{584547}
\rqSet{rqAttrInputsNoInputsLibAllGptFivePct}{0.0}
\rqSet{rqAttrInputsNoInputsLibAllGeminiNum}{0}
\rqSet{rqAttrInputsNoInputsLibAllGeminiDen}{572793}
\rqSet{rqAttrInputsNoInputsLibAllGeminiPct}{0.0}
\rqSet{rqAttrInputsNoInputsLibBOneNum}{0}
\rqSet{rqAttrInputsNoInputsLibBOneDen}{3452556}
\rqSet{rqAttrInputsNoInputsLibBOnePct}{0.0}
\rqSet{rqAttrInputsNoInputsLibBTwoNum}{0}
\rqSet{rqAttrInputsNoInputsLibBTwoDen}{2472666}
\rqSet{rqAttrInputsNoInputsLibBTwoPct}{0.0}
\rqSet{rqAttrInputsNoInputsLibPZeroNum}{0}
\rqSet{rqAttrInputsNoInputsLibPZeroDen}{0}
\rqSet{rqAttrInputsNoInputsLibPZeroPct}{n/a}
\rqSet{rqAttrInputsNoInputsLibPOneNum}{0}
\rqSet{rqAttrInputsNoInputsLibPOneDen}{1014507}
\rqSet{rqAttrInputsNoInputsLibPOnePct}{0.0}
\rqSet{rqAttrInputsIncompleteLibAllNum}{32696}
\rqSet{rqAttrInputsIncompleteLibAllDen}{6939729}
\rqSet{rqAttrInputsIncompleteLibAllPct}{0.5}
\rqSet{rqAttrInputsIncompleteLibAllKiroNum}{4488}
\rqSet{rqAttrInputsIncompleteLibAllKiroDen}{1209621}
\rqSet{rqAttrInputsIncompleteLibAllKiroPct}{0.4}
\rqSet{rqAttrInputsIncompleteLibAllClaudeNum}{4495}
\rqSet{rqAttrInputsIncompleteLibAllClaudeDen}{1143588}
\rqSet{rqAttrInputsIncompleteLibAllClaudePct}{0.4}
\rqSet{rqAttrInputsIncompleteLibAllCodexNum}{4568}
\rqSet{rqAttrInputsIncompleteLibAllCodexDen}{915456}
\rqSet{rqAttrInputsIncompleteLibAllCodexPct}{0.5}
\rqSet{rqAttrInputsIncompleteLibAllCTwoRustNum}{4504}
\rqSet{rqAttrInputsIncompleteLibAllCTwoRustDen}{753393}
\rqSet{rqAttrInputsIncompleteLibAllCTwoRustPct}{0.6}
\rqSet{rqAttrInputsIncompleteLibAllLaertesNum}{4502}
\rqSet{rqAttrInputsIncompleteLibAllLaertesDen}{741390}
\rqSet{rqAttrInputsIncompleteLibAllLaertesPct}{0.6}
\rqSet{rqAttrInputsIncompleteLibAllCTwoSaferRustNum}{3080}
\rqSet{rqAttrInputsIncompleteLibAllCTwoSaferRustDen}{554526}
\rqSet{rqAttrInputsIncompleteLibAllCTwoSaferRustPct}{0.6}
\rqSet{rqAttrInputsIncompleteLibAllSmartCTwoRustNum}{0}
\rqSet{rqAttrInputsIncompleteLibAllSmartCTwoRustDen}{0}
\rqSet{rqAttrInputsIncompleteLibAllSmartCTwoRustPct}{n/a}
\rqSet{rqAttrInputsIncompleteLibAllKimiNum}{380}
\rqSet{rqAttrInputsIncompleteLibAllKimiDen}{464415}
\rqSet{rqAttrInputsIncompleteLibAllKimiPct}{0.1}
\rqSet{rqAttrInputsIncompleteLibAllGptFiveNum}{3338}
\rqSet{rqAttrInputsIncompleteLibAllGptFiveDen}{584547}
\rqSet{rqAttrInputsIncompleteLibAllGptFivePct}{0.6}
\rqSet{rqAttrInputsIncompleteLibAllGeminiNum}{3341}
\rqSet{rqAttrInputsIncompleteLibAllGeminiDen}{572793}
\rqSet{rqAttrInputsIncompleteLibAllGeminiPct}{0.6}
\rqSet{rqAttrInputsIncompleteLibBOneNum}{2216}
\rqSet{rqAttrInputsIncompleteLibBOneDen}{3452556}
\rqSet{rqAttrInputsIncompleteLibBOnePct}{0.1}
\rqSet{rqAttrInputsIncompleteLibBTwoNum}{30480}
\rqSet{rqAttrInputsIncompleteLibBTwoDen}{2472666}
\rqSet{rqAttrInputsIncompleteLibBTwoPct}{1.2}
\rqSet{rqAttrInputsIncompleteLibPZeroNum}{0}
\rqSet{rqAttrInputsIncompleteLibPZeroDen}{0}
\rqSet{rqAttrInputsIncompleteLibPZeroPct}{n/a}
\rqSet{rqAttrInputsIncompleteLibPOneNum}{0}
\rqSet{rqAttrInputsIncompleteLibPOneDen}{1014507}
\rqSet{rqAttrInputsIncompleteLibPOnePct}{0.0}
\rqSet{rqAttrInputsSanInfraLibAllNum}{0}
\rqSet{rqAttrInputsSanInfraLibAllDen}{6939729}
\rqSet{rqAttrInputsSanInfraLibAllPct}{0.0}
\rqSet{rqAttrInputsSanInfraLibAllKiroNum}{0}
\rqSet{rqAttrInputsSanInfraLibAllKiroDen}{1209621}
\rqSet{rqAttrInputsSanInfraLibAllKiroPct}{0.0}
\rqSet{rqAttrInputsSanInfraLibAllClaudeNum}{0}
\rqSet{rqAttrInputsSanInfraLibAllClaudeDen}{1143588}
\rqSet{rqAttrInputsSanInfraLibAllClaudePct}{0.0}
\rqSet{rqAttrInputsSanInfraLibAllCodexNum}{0}
\rqSet{rqAttrInputsSanInfraLibAllCodexDen}{915456}
\rqSet{rqAttrInputsSanInfraLibAllCodexPct}{0.0}
\rqSet{rqAttrInputsSanInfraLibAllCTwoRustNum}{0}
\rqSet{rqAttrInputsSanInfraLibAllCTwoRustDen}{753393}
\rqSet{rqAttrInputsSanInfraLibAllCTwoRustPct}{0.0}
\rqSet{rqAttrInputsSanInfraLibAllLaertesNum}{0}
\rqSet{rqAttrInputsSanInfraLibAllLaertesDen}{741390}
\rqSet{rqAttrInputsSanInfraLibAllLaertesPct}{0.0}
\rqSet{rqAttrInputsSanInfraLibAllCTwoSaferRustNum}{0}
\rqSet{rqAttrInputsSanInfraLibAllCTwoSaferRustDen}{554526}
\rqSet{rqAttrInputsSanInfraLibAllCTwoSaferRustPct}{0.0}
\rqSet{rqAttrInputsSanInfraLibAllSmartCTwoRustNum}{0}
\rqSet{rqAttrInputsSanInfraLibAllSmartCTwoRustDen}{0}
\rqSet{rqAttrInputsSanInfraLibAllSmartCTwoRustPct}{n/a}
\rqSet{rqAttrInputsSanInfraLibAllKimiNum}{0}
\rqSet{rqAttrInputsSanInfraLibAllKimiDen}{464415}
\rqSet{rqAttrInputsSanInfraLibAllKimiPct}{0.0}
\rqSet{rqAttrInputsSanInfraLibAllGptFiveNum}{0}
\rqSet{rqAttrInputsSanInfraLibAllGptFiveDen}{584547}
\rqSet{rqAttrInputsSanInfraLibAllGptFivePct}{0.0}
\rqSet{rqAttrInputsSanInfraLibAllGeminiNum}{0}
\rqSet{rqAttrInputsSanInfraLibAllGeminiDen}{572793}
\rqSet{rqAttrInputsSanInfraLibAllGeminiPct}{0.0}
\rqSet{rqAttrInputsSanInfraLibBOneNum}{0}
\rqSet{rqAttrInputsSanInfraLibBOneDen}{3452556}
\rqSet{rqAttrInputsSanInfraLibBOnePct}{0.0}
\rqSet{rqAttrInputsSanInfraLibBTwoNum}{0}
\rqSet{rqAttrInputsSanInfraLibBTwoDen}{2472666}
\rqSet{rqAttrInputsSanInfraLibBTwoPct}{0.0}
\rqSet{rqAttrInputsSanInfraLibPZeroNum}{0}
\rqSet{rqAttrInputsSanInfraLibPZeroDen}{0}
\rqSet{rqAttrInputsSanInfraLibPZeroPct}{n/a}
\rqSet{rqAttrInputsSanInfraLibPOneNum}{0}
\rqSet{rqAttrInputsSanInfraLibPOneDen}{1014507}
\rqSet{rqAttrInputsSanInfraLibPOnePct}{0.0}
\rqSet{rqAttrInputsMissingSeedLibAllNum}{??}
\rqSet{rqAttrInputsMissingSeedLibAllDen}{6939729}
\rqSet{rqAttrInputsMissingSeedLibAllPct}{??}
\rqSet{rqAttrInputsMissingSeedLibAllKiroNum}{??}
\rqSet{rqAttrInputsMissingSeedLibAllKiroDen}{1209621}
\rqSet{rqAttrInputsMissingSeedLibAllKiroPct}{??}
\rqSet{rqAttrInputsMissingSeedLibAllClaudeNum}{??}
\rqSet{rqAttrInputsMissingSeedLibAllClaudeDen}{1143588}
\rqSet{rqAttrInputsMissingSeedLibAllClaudePct}{??}
\rqSet{rqAttrInputsMissingSeedLibAllCodexNum}{??}
\rqSet{rqAttrInputsMissingSeedLibAllCodexDen}{915456}
\rqSet{rqAttrInputsMissingSeedLibAllCodexPct}{??}
\rqSet{rqAttrInputsMissingSeedLibAllCTwoRustNum}{??}
\rqSet{rqAttrInputsMissingSeedLibAllCTwoRustDen}{753393}
\rqSet{rqAttrInputsMissingSeedLibAllCTwoRustPct}{??}
\rqSet{rqAttrInputsMissingSeedLibAllLaertesNum}{??}
\rqSet{rqAttrInputsMissingSeedLibAllLaertesDen}{741390}
\rqSet{rqAttrInputsMissingSeedLibAllLaertesPct}{??}
\rqSet{rqAttrInputsMissingSeedLibAllCTwoSaferRustNum}{??}
\rqSet{rqAttrInputsMissingSeedLibAllCTwoSaferRustDen}{554526}
\rqSet{rqAttrInputsMissingSeedLibAllCTwoSaferRustPct}{??}
\rqSet{rqAttrInputsMissingSeedLibAllSmartCTwoRustNum}{0}
\rqSet{rqAttrInputsMissingSeedLibAllSmartCTwoRustDen}{0}
\rqSet{rqAttrInputsMissingSeedLibAllSmartCTwoRustPct}{n/a}
\rqSet{rqAttrInputsMissingSeedLibAllKimiNum}{??}
\rqSet{rqAttrInputsMissingSeedLibAllKimiDen}{464415}
\rqSet{rqAttrInputsMissingSeedLibAllKimiPct}{??}
\rqSet{rqAttrInputsMissingSeedLibAllGptFiveNum}{??}
\rqSet{rqAttrInputsMissingSeedLibAllGptFiveDen}{584547}
\rqSet{rqAttrInputsMissingSeedLibAllGptFivePct}{??}
\rqSet{rqAttrInputsMissingSeedLibAllGeminiNum}{??}
\rqSet{rqAttrInputsMissingSeedLibAllGeminiDen}{572793}
\rqSet{rqAttrInputsMissingSeedLibAllGeminiPct}{??}
\rqSet{rqAttrInputsMissingSeedLibBOneNum}{??}
\rqSet{rqAttrInputsMissingSeedLibBOneDen}{3452556}
\rqSet{rqAttrInputsMissingSeedLibBOnePct}{??}
\rqSet{rqAttrInputsMissingSeedLibBTwoNum}{??}
\rqSet{rqAttrInputsMissingSeedLibBTwoDen}{2472666}
\rqSet{rqAttrInputsMissingSeedLibBTwoPct}{??}
\rqSet{rqAttrInputsMissingSeedLibPZeroNum}{0}
\rqSet{rqAttrInputsMissingSeedLibPZeroDen}{0}
\rqSet{rqAttrInputsMissingSeedLibPZeroPct}{n/a}
\rqSet{rqAttrInputsMissingSeedLibPOneNum}{0}
\rqSet{rqAttrInputsMissingSeedLibPOneDen}{1014507}
\rqSet{rqAttrInputsMissingSeedLibPOnePct}{0.0}
\rqSet{rqAttrInputsMissingRecordLibAllNum}{??}
\rqSet{rqAttrInputsMissingRecordLibAllDen}{6939729}
\rqSet{rqAttrInputsMissingRecordLibAllPct}{??}
\rqSet{rqAttrInputsMissingRecordLibAllKiroNum}{??}
\rqSet{rqAttrInputsMissingRecordLibAllKiroDen}{1209621}
\rqSet{rqAttrInputsMissingRecordLibAllKiroPct}{??}
\rqSet{rqAttrInputsMissingRecordLibAllClaudeNum}{??}
\rqSet{rqAttrInputsMissingRecordLibAllClaudeDen}{1143588}
\rqSet{rqAttrInputsMissingRecordLibAllClaudePct}{??}
\rqSet{rqAttrInputsMissingRecordLibAllCodexNum}{??}
\rqSet{rqAttrInputsMissingRecordLibAllCodexDen}{915456}
\rqSet{rqAttrInputsMissingRecordLibAllCodexPct}{??}
\rqSet{rqAttrInputsMissingRecordLibAllCTwoRustNum}{??}
\rqSet{rqAttrInputsMissingRecordLibAllCTwoRustDen}{753393}
\rqSet{rqAttrInputsMissingRecordLibAllCTwoRustPct}{??}
\rqSet{rqAttrInputsMissingRecordLibAllLaertesNum}{??}
\rqSet{rqAttrInputsMissingRecordLibAllLaertesDen}{741390}
\rqSet{rqAttrInputsMissingRecordLibAllLaertesPct}{??}
\rqSet{rqAttrInputsMissingRecordLibAllCTwoSaferRustNum}{??}
\rqSet{rqAttrInputsMissingRecordLibAllCTwoSaferRustDen}{554526}
\rqSet{rqAttrInputsMissingRecordLibAllCTwoSaferRustPct}{??}
\rqSet{rqAttrInputsMissingRecordLibAllSmartCTwoRustNum}{0}
\rqSet{rqAttrInputsMissingRecordLibAllSmartCTwoRustDen}{0}
\rqSet{rqAttrInputsMissingRecordLibAllSmartCTwoRustPct}{n/a}
\rqSet{rqAttrInputsMissingRecordLibAllKimiNum}{0}
\rqSet{rqAttrInputsMissingRecordLibAllKimiDen}{464415}
\rqSet{rqAttrInputsMissingRecordLibAllKimiPct}{0.0}
\rqSet{rqAttrInputsMissingRecordLibAllGptFiveNum}{??}
\rqSet{rqAttrInputsMissingRecordLibAllGptFiveDen}{584547}
\rqSet{rqAttrInputsMissingRecordLibAllGptFivePct}{??}
\rqSet{rqAttrInputsMissingRecordLibAllGeminiNum}{??}
\rqSet{rqAttrInputsMissingRecordLibAllGeminiDen}{572793}
\rqSet{rqAttrInputsMissingRecordLibAllGeminiPct}{??}
\rqSet{rqAttrInputsMissingRecordLibBOneNum}{0}
\rqSet{rqAttrInputsMissingRecordLibBOneDen}{3452556}
\rqSet{rqAttrInputsMissingRecordLibBOnePct}{0.0}
\rqSet{rqAttrInputsMissingRecordLibBTwoNum}{??}
\rqSet{rqAttrInputsMissingRecordLibBTwoDen}{2472666}
\rqSet{rqAttrInputsMissingRecordLibBTwoPct}{??}
\rqSet{rqAttrInputsMissingRecordLibPZeroNum}{0}
\rqSet{rqAttrInputsMissingRecordLibPZeroDen}{0}
\rqSet{rqAttrInputsMissingRecordLibPZeroPct}{n/a}
\rqSet{rqAttrInputsMissingRecordLibPOneNum}{0}
\rqSet{rqAttrInputsMissingRecordLibPOneDen}{1014507}
\rqSet{rqAttrInputsMissingRecordLibPOnePct}{0.0}
\rqSet{rqAttrInputsUnknownLibAllNum}{91}
\rqSet{rqAttrInputsUnknownLibAllDen}{1264}
\rqSet{rqAttrInputsUnknownLibAllPct}{7.2}
\rqSet{rqAttrInputsUnknownLibAllKiroNum}{12}
\rqSet{rqAttrInputsUnknownLibAllKiroDen}{216}
\rqSet{rqAttrInputsUnknownLibAllKiroPct}{5.6}
\rqSet{rqAttrInputsUnknownLibAllClaudeNum}{12}
\rqSet{rqAttrInputsUnknownLibAllClaudeDen}{205}
\rqSet{rqAttrInputsUnknownLibAllClaudePct}{5.9}
\rqSet{rqAttrInputsUnknownLibAllCodexNum}{12}
\rqSet{rqAttrInputsUnknownLibAllCodexDen}{167}
\rqSet{rqAttrInputsUnknownLibAllCodexPct}{7.2}
\rqSet{rqAttrInputsUnknownLibAllCTwoRustNum}{12}
\rqSet{rqAttrInputsUnknownLibAllCTwoRustDen}{140}
\rqSet{rqAttrInputsUnknownLibAllCTwoRustPct}{8.6}
\rqSet{rqAttrInputsUnknownLibAllLaertesNum}{11}
\rqSet{rqAttrInputsUnknownLibAllLaertesDen}{136}
\rqSet{rqAttrInputsUnknownLibAllLaertesPct}{8.1}
\rqSet{rqAttrInputsUnknownLibAllCTwoSaferRustNum}{9}
\rqSet{rqAttrInputsUnknownLibAllCTwoSaferRustDen}{103}
\rqSet{rqAttrInputsUnknownLibAllCTwoSaferRustPct}{8.7}
\rqSet{rqAttrInputsUnknownLibAllSmartCTwoRustNum}{0}
\rqSet{rqAttrInputsUnknownLibAllSmartCTwoRustDen}{0}
\rqSet{rqAttrInputsUnknownLibAllSmartCTwoRustPct}{n/a}
\rqSet{rqAttrInputsUnknownLibAllKimiNum}{5}
\rqSet{rqAttrInputsUnknownLibAllKimiDen}{83}
\rqSet{rqAttrInputsUnknownLibAllKimiPct}{6.0}
\rqSet{rqAttrInputsUnknownLibAllGptFiveNum}{9}
\rqSet{rqAttrInputsUnknownLibAllGptFiveDen}{108}
\rqSet{rqAttrInputsUnknownLibAllGptFivePct}{8.3}
\rqSet{rqAttrInputsUnknownLibAllGeminiNum}{9}
\rqSet{rqAttrInputsUnknownLibAllGeminiDen}{106}
\rqSet{rqAttrInputsUnknownLibAllGeminiPct}{8.5}
\rqSet{rqAttrInputsUnknownLibBOneNum}{48}
\rqSet{rqAttrInputsUnknownLibBOneDen}{636}
\rqSet{rqAttrInputsUnknownLibBOnePct}{7.5}
\rqSet{rqAttrInputsUnknownLibBTwoNum}{43}
\rqSet{rqAttrInputsUnknownLibBTwoDen}{459}
\rqSet{rqAttrInputsUnknownLibBTwoPct}{9.4}
\rqSet{rqAttrInputsUnknownLibPZeroNum}{0}
\rqSet{rqAttrInputsUnknownLibPZeroDen}{0}
\rqSet{rqAttrInputsUnknownLibPZeroPct}{n/a}
\rqSet{rqAttrInputsUnknownLibPOneNum}{0}
\rqSet{rqAttrInputsUnknownLibPOneDen}{169}
\rqSet{rqAttrInputsUnknownLibPOnePct}{0.0}
\rqSet{rqAttrInputsSanUnresolvedLibAllNum}{0}
\rqSet{rqAttrInputsSanUnresolvedLibAllDen}{0}
\rqSet{rqAttrInputsSanUnresolvedLibBOneNum}{0}
\rqSet{rqAttrInputsSanUnresolvedLibBOneDen}{0}
\rqSet{rqAttrInputsSanUnresolvedLibBTwoNum}{0}
\rqSet{rqAttrInputsSanUnresolvedLibBTwoDen}{0}
\rqSet{rqAttrInputsSanUnresolvedLibPZeroNum}{0}
\rqSet{rqAttrInputsSanUnresolvedLibPZeroDen}{0}
\rqSet{rqAttrInputsSanUnresolvedLibPOneNum}{0}
\rqSet{rqAttrInputsSanUnresolvedLibPOneDen}{0}
\rqSet{rqNonRepAccOracleDriftLibAllNum}{0}
\rqSet{rqNonRepAccOracleDriftLibAllDen}{6}
\rqSet{rqNonRepAccOracleDriftLibAllPct}{0.0}
\rqSet{rqNonRepAccOracleDriftLibAllKiroNum}{0}
\rqSet{rqNonRepAccOracleDriftLibAllKiroDen}{0}
\rqSet{rqNonRepAccOracleDriftLibAllKiroPct}{n/a}
\rqSet{rqNonRepAccOracleDriftLibAllClaudeNum}{0}
\rqSet{rqNonRepAccOracleDriftLibAllClaudeDen}{0}
\rqSet{rqNonRepAccOracleDriftLibAllClaudePct}{n/a}
\rqSet{rqNonRepAccOracleDriftLibAllCodexNum}{0}
\rqSet{rqNonRepAccOracleDriftLibAllCodexDen}{3}
\rqSet{rqNonRepAccOracleDriftLibAllCodexPct}{0.0}
\rqSet{rqNonRepAccOracleDriftLibAllCTwoRustNum}{0}
\rqSet{rqNonRepAccOracleDriftLibAllCTwoRustDen}{0}
\rqSet{rqNonRepAccOracleDriftLibAllCTwoRustPct}{n/a}
\rqSet{rqNonRepAccOracleDriftLibAllLaertesNum}{0}
\rqSet{rqNonRepAccOracleDriftLibAllLaertesDen}{0}
\rqSet{rqNonRepAccOracleDriftLibAllLaertesPct}{n/a}
\rqSet{rqNonRepAccOracleDriftLibAllCTwoSaferRustNum}{0}
\rqSet{rqNonRepAccOracleDriftLibAllCTwoSaferRustDen}{1}
\rqSet{rqNonRepAccOracleDriftLibAllCTwoSaferRustPct}{0.0}
\rqSet{rqNonRepAccOracleDriftLibAllSmartCTwoRustNum}{0}
\rqSet{rqNonRepAccOracleDriftLibAllSmartCTwoRustDen}{0}
\rqSet{rqNonRepAccOracleDriftLibAllSmartCTwoRustPct}{n/a}
\rqSet{rqNonRepAccOracleDriftLibAllKimiNum}{0}
\rqSet{rqNonRepAccOracleDriftLibAllKimiDen}{0}
\rqSet{rqNonRepAccOracleDriftLibAllKimiPct}{n/a}
\rqSet{rqNonRepAccOracleDriftLibAllGptFiveNum}{0}
\rqSet{rqNonRepAccOracleDriftLibAllGptFiveDen}{1}
\rqSet{rqNonRepAccOracleDriftLibAllGptFivePct}{0.0}
\rqSet{rqNonRepAccOracleDriftLibAllGeminiNum}{0}
\rqSet{rqNonRepAccOracleDriftLibAllGeminiDen}{1}
\rqSet{rqNonRepAccOracleDriftLibAllGeminiPct}{0.0}
\rqSet{rqNonRepAccOracleDriftLibBOneNum}{0}
\rqSet{rqNonRepAccOracleDriftLibBOneDen}{1}
\rqSet{rqNonRepAccOracleDriftLibBOnePct}{0.0}
\rqSet{rqNonRepAccOracleDriftLibBTwoNum}{0}
\rqSet{rqNonRepAccOracleDriftLibBTwoDen}{2}
\rqSet{rqNonRepAccOracleDriftLibBTwoPct}{0.0}
\rqSet{rqNonRepAccOracleDriftLibPZeroNum}{0}
\rqSet{rqNonRepAccOracleDriftLibPZeroDen}{0}
\rqSet{rqNonRepAccOracleDriftLibPZeroPct}{n/a}
\rqSet{rqNonRepAccOracleDriftLibPOneNum}{0}
\rqSet{rqNonRepAccOracleDriftLibPOneDen}{3}
\rqSet{rqNonRepAccOracleDriftLibPOnePct}{0.0}
\rqSet{rqNonRepAccTranslationLibAllNum}{4}
\rqSet{rqNonRepAccTranslationLibAllDen}{6}
\rqSet{rqNonRepAccTranslationLibAllPct}{66.7}
\rqSet{rqNonRepAccTranslationLibAllKiroNum}{0}
\rqSet{rqNonRepAccTranslationLibAllKiroDen}{0}
\rqSet{rqNonRepAccTranslationLibAllKiroPct}{n/a}
\rqSet{rqNonRepAccTranslationLibAllClaudeNum}{0}
\rqSet{rqNonRepAccTranslationLibAllClaudeDen}{0}
\rqSet{rqNonRepAccTranslationLibAllClaudePct}{n/a}
\rqSet{rqNonRepAccTranslationLibAllCodexNum}{3}
\rqSet{rqNonRepAccTranslationLibAllCodexDen}{3}
\rqSet{rqNonRepAccTranslationLibAllCodexPct}{100.0}
\rqSet{rqNonRepAccTranslationLibAllCTwoRustNum}{0}
\rqSet{rqNonRepAccTranslationLibAllCTwoRustDen}{0}
\rqSet{rqNonRepAccTranslationLibAllCTwoRustPct}{n/a}
\rqSet{rqNonRepAccTranslationLibAllLaertesNum}{0}
\rqSet{rqNonRepAccTranslationLibAllLaertesDen}{0}
\rqSet{rqNonRepAccTranslationLibAllLaertesPct}{n/a}
\rqSet{rqNonRepAccTranslationLibAllCTwoSaferRustNum}{1}
\rqSet{rqNonRepAccTranslationLibAllCTwoSaferRustDen}{1}
\rqSet{rqNonRepAccTranslationLibAllCTwoSaferRustPct}{100.0}
\rqSet{rqNonRepAccTranslationLibAllSmartCTwoRustNum}{0}
\rqSet{rqNonRepAccTranslationLibAllSmartCTwoRustDen}{0}
\rqSet{rqNonRepAccTranslationLibAllSmartCTwoRustPct}{n/a}
\rqSet{rqNonRepAccTranslationLibAllKimiNum}{0}
\rqSet{rqNonRepAccTranslationLibAllKimiDen}{0}
\rqSet{rqNonRepAccTranslationLibAllKimiPct}{n/a}
\rqSet{rqNonRepAccTranslationLibAllGptFiveNum}{0}
\rqSet{rqNonRepAccTranslationLibAllGptFiveDen}{1}
\rqSet{rqNonRepAccTranslationLibAllGptFivePct}{0.0}
\rqSet{rqNonRepAccTranslationLibAllGeminiNum}{0}
\rqSet{rqNonRepAccTranslationLibAllGeminiDen}{1}
\rqSet{rqNonRepAccTranslationLibAllGeminiPct}{0.0}
\rqSet{rqNonRepAccTranslationLibBOneNum}{1}
\rqSet{rqNonRepAccTranslationLibBOneDen}{1}
\rqSet{rqNonRepAccTranslationLibBOnePct}{100.0}
\rqSet{rqNonRepAccTranslationLibBTwoNum}{0}
\rqSet{rqNonRepAccTranslationLibBTwoDen}{2}
\rqSet{rqNonRepAccTranslationLibBTwoPct}{0.0}
\rqSet{rqNonRepAccTranslationLibPZeroNum}{0}
\rqSet{rqNonRepAccTranslationLibPZeroDen}{0}
\rqSet{rqNonRepAccTranslationLibPZeroPct}{n/a}
\rqSet{rqNonRepAccTranslationLibPOneNum}{3}
\rqSet{rqNonRepAccTranslationLibPOneDen}{3}
\rqSet{rqNonRepAccTranslationLibPOnePct}{100.0}
\rqSet{rqNonRepAccUndetectedUbLibAllNum}{0}
\rqSet{rqNonRepAccUndetectedUbLibAllDen}{6}
\rqSet{rqNonRepAccUndetectedUbLibAllPct}{0.0}
\rqSet{rqNonRepAccUndetectedUbLibAllKiroNum}{0}
\rqSet{rqNonRepAccUndetectedUbLibAllKiroDen}{0}
\rqSet{rqNonRepAccUndetectedUbLibAllKiroPct}{n/a}
\rqSet{rqNonRepAccUndetectedUbLibAllClaudeNum}{0}
\rqSet{rqNonRepAccUndetectedUbLibAllClaudeDen}{0}
\rqSet{rqNonRepAccUndetectedUbLibAllClaudePct}{n/a}
\rqSet{rqNonRepAccUndetectedUbLibAllCodexNum}{0}
\rqSet{rqNonRepAccUndetectedUbLibAllCodexDen}{3}
\rqSet{rqNonRepAccUndetectedUbLibAllCodexPct}{0.0}
\rqSet{rqNonRepAccUndetectedUbLibAllCTwoRustNum}{0}
\rqSet{rqNonRepAccUndetectedUbLibAllCTwoRustDen}{0}
\rqSet{rqNonRepAccUndetectedUbLibAllCTwoRustPct}{n/a}
\rqSet{rqNonRepAccUndetectedUbLibAllLaertesNum}{0}
\rqSet{rqNonRepAccUndetectedUbLibAllLaertesDen}{0}
\rqSet{rqNonRepAccUndetectedUbLibAllLaertesPct}{n/a}
\rqSet{rqNonRepAccUndetectedUbLibAllCTwoSaferRustNum}{0}
\rqSet{rqNonRepAccUndetectedUbLibAllCTwoSaferRustDen}{1}
\rqSet{rqNonRepAccUndetectedUbLibAllCTwoSaferRustPct}{0.0}
\rqSet{rqNonRepAccUndetectedUbLibAllSmartCTwoRustNum}{0}
\rqSet{rqNonRepAccUndetectedUbLibAllSmartCTwoRustDen}{0}
\rqSet{rqNonRepAccUndetectedUbLibAllSmartCTwoRustPct}{n/a}
\rqSet{rqNonRepAccUndetectedUbLibAllKimiNum}{0}
\rqSet{rqNonRepAccUndetectedUbLibAllKimiDen}{0}
\rqSet{rqNonRepAccUndetectedUbLibAllKimiPct}{n/a}
\rqSet{rqNonRepAccUndetectedUbLibAllGptFiveNum}{0}
\rqSet{rqNonRepAccUndetectedUbLibAllGptFiveDen}{1}
\rqSet{rqNonRepAccUndetectedUbLibAllGptFivePct}{0.0}
\rqSet{rqNonRepAccUndetectedUbLibAllGeminiNum}{0}
\rqSet{rqNonRepAccUndetectedUbLibAllGeminiDen}{1}
\rqSet{rqNonRepAccUndetectedUbLibAllGeminiPct}{0.0}
\rqSet{rqNonRepAccUndetectedUbLibBOneNum}{0}
\rqSet{rqNonRepAccUndetectedUbLibBOneDen}{1}
\rqSet{rqNonRepAccUndetectedUbLibBOnePct}{0.0}
\rqSet{rqNonRepAccUndetectedUbLibBTwoNum}{0}
\rqSet{rqNonRepAccUndetectedUbLibBTwoDen}{2}
\rqSet{rqNonRepAccUndetectedUbLibBTwoPct}{0.0}
\rqSet{rqNonRepAccUndetectedUbLibPZeroNum}{0}
\rqSet{rqNonRepAccUndetectedUbLibPZeroDen}{0}
\rqSet{rqNonRepAccUndetectedUbLibPZeroPct}{n/a}
\rqSet{rqNonRepAccUndetectedUbLibPOneNum}{0}
\rqSet{rqNonRepAccUndetectedUbLibPOneDen}{3}
\rqSet{rqNonRepAccUndetectedUbLibPOnePct}{0.0}
\rqSet{rqNonRepAccUnknownLibAllNum}{2}
\rqSet{rqNonRepAccUnknownLibAllDen}{6}
\rqSet{rqNonRepAccUnknownLibAllPct}{33.3}
\rqSet{rqNonRepAccUnknownLibAllKiroNum}{0}
\rqSet{rqNonRepAccUnknownLibAllKiroDen}{0}
\rqSet{rqNonRepAccUnknownLibAllKiroPct}{n/a}
\rqSet{rqNonRepAccUnknownLibAllClaudeNum}{0}
\rqSet{rqNonRepAccUnknownLibAllClaudeDen}{0}
\rqSet{rqNonRepAccUnknownLibAllClaudePct}{n/a}
\rqSet{rqNonRepAccUnknownLibAllCodexNum}{0}
\rqSet{rqNonRepAccUnknownLibAllCodexDen}{3}
\rqSet{rqNonRepAccUnknownLibAllCodexPct}{0.0}
\rqSet{rqNonRepAccUnknownLibAllCTwoRustNum}{0}
\rqSet{rqNonRepAccUnknownLibAllCTwoRustDen}{0}
\rqSet{rqNonRepAccUnknownLibAllCTwoRustPct}{n/a}
\rqSet{rqNonRepAccUnknownLibAllLaertesNum}{0}
\rqSet{rqNonRepAccUnknownLibAllLaertesDen}{0}
\rqSet{rqNonRepAccUnknownLibAllLaertesPct}{n/a}
\rqSet{rqNonRepAccUnknownLibAllCTwoSaferRustNum}{0}
\rqSet{rqNonRepAccUnknownLibAllCTwoSaferRustDen}{1}
\rqSet{rqNonRepAccUnknownLibAllCTwoSaferRustPct}{0.0}
\rqSet{rqNonRepAccUnknownLibAllSmartCTwoRustNum}{0}
\rqSet{rqNonRepAccUnknownLibAllSmartCTwoRustDen}{0}
\rqSet{rqNonRepAccUnknownLibAllSmartCTwoRustPct}{n/a}
\rqSet{rqNonRepAccUnknownLibAllKimiNum}{0}
\rqSet{rqNonRepAccUnknownLibAllKimiDen}{0}
\rqSet{rqNonRepAccUnknownLibAllKimiPct}{n/a}
\rqSet{rqNonRepAccUnknownLibAllGptFiveNum}{1}
\rqSet{rqNonRepAccUnknownLibAllGptFiveDen}{1}
\rqSet{rqNonRepAccUnknownLibAllGptFivePct}{100.0}
\rqSet{rqNonRepAccUnknownLibAllGeminiNum}{1}
\rqSet{rqNonRepAccUnknownLibAllGeminiDen}{1}
\rqSet{rqNonRepAccUnknownLibAllGeminiPct}{100.0}
\rqSet{rqNonRepAccUnknownLibBOneNum}{0}
\rqSet{rqNonRepAccUnknownLibBOneDen}{1}
\rqSet{rqNonRepAccUnknownLibBOnePct}{0.0}
\rqSet{rqNonRepAccUnknownLibBTwoNum}{2}
\rqSet{rqNonRepAccUnknownLibBTwoDen}{2}
\rqSet{rqNonRepAccUnknownLibBTwoPct}{100.0}
\rqSet{rqNonRepAccUnknownLibPZeroNum}{0}
\rqSet{rqNonRepAccUnknownLibPZeroDen}{0}
\rqSet{rqNonRepAccUnknownLibPZeroPct}{n/a}
\rqSet{rqNonRepAccUnknownLibPOneNum}{0}
\rqSet{rqNonRepAccUnknownLibPOneDen}{3}
\rqSet{rqNonRepAccUnknownLibPOnePct}{0.0}
\rqSet{rqNonRepAccUnclassifiedLibAllNum}{0}
\rqSet{rqNonRepAccUnclassifiedLibAllDen}{6}
\rqSet{rqNonRepAccUnclassifiedLibAllPct}{0.0}
\rqSet{rqNonRepAccUnclassifiedLibAllKiroNum}{0}
\rqSet{rqNonRepAccUnclassifiedLibAllKiroDen}{0}
\rqSet{rqNonRepAccUnclassifiedLibAllKiroPct}{n/a}
\rqSet{rqNonRepAccUnclassifiedLibAllClaudeNum}{0}
\rqSet{rqNonRepAccUnclassifiedLibAllClaudeDen}{0}
\rqSet{rqNonRepAccUnclassifiedLibAllClaudePct}{n/a}
\rqSet{rqNonRepAccUnclassifiedLibAllCodexNum}{0}
\rqSet{rqNonRepAccUnclassifiedLibAllCodexDen}{3}
\rqSet{rqNonRepAccUnclassifiedLibAllCodexPct}{0.0}
\rqSet{rqNonRepAccUnclassifiedLibAllCTwoRustNum}{0}
\rqSet{rqNonRepAccUnclassifiedLibAllCTwoRustDen}{0}
\rqSet{rqNonRepAccUnclassifiedLibAllCTwoRustPct}{n/a}
\rqSet{rqNonRepAccUnclassifiedLibAllLaertesNum}{0}
\rqSet{rqNonRepAccUnclassifiedLibAllLaertesDen}{0}
\rqSet{rqNonRepAccUnclassifiedLibAllLaertesPct}{n/a}
\rqSet{rqNonRepAccUnclassifiedLibAllCTwoSaferRustNum}{0}
\rqSet{rqNonRepAccUnclassifiedLibAllCTwoSaferRustDen}{1}
\rqSet{rqNonRepAccUnclassifiedLibAllCTwoSaferRustPct}{0.0}
\rqSet{rqNonRepAccUnclassifiedLibAllSmartCTwoRustNum}{0}
\rqSet{rqNonRepAccUnclassifiedLibAllSmartCTwoRustDen}{0}
\rqSet{rqNonRepAccUnclassifiedLibAllSmartCTwoRustPct}{n/a}
\rqSet{rqNonRepAccUnclassifiedLibAllKimiNum}{0}
\rqSet{rqNonRepAccUnclassifiedLibAllKimiDen}{0}
\rqSet{rqNonRepAccUnclassifiedLibAllKimiPct}{n/a}
\rqSet{rqNonRepAccUnclassifiedLibAllGptFiveNum}{0}
\rqSet{rqNonRepAccUnclassifiedLibAllGptFiveDen}{1}
\rqSet{rqNonRepAccUnclassifiedLibAllGptFivePct}{0.0}
\rqSet{rqNonRepAccUnclassifiedLibAllGeminiNum}{0}
\rqSet{rqNonRepAccUnclassifiedLibAllGeminiDen}{1}
\rqSet{rqNonRepAccUnclassifiedLibAllGeminiPct}{0.0}
\rqSet{rqNonRepAccUnclassifiedLibBOneNum}{0}
\rqSet{rqNonRepAccUnclassifiedLibBOneDen}{1}
\rqSet{rqNonRepAccUnclassifiedLibBOnePct}{0.0}
\rqSet{rqNonRepAccUnclassifiedLibBTwoNum}{0}
\rqSet{rqNonRepAccUnclassifiedLibBTwoDen}{2}
\rqSet{rqNonRepAccUnclassifiedLibBTwoPct}{0.0}
\rqSet{rqNonRepAccUnclassifiedLibPZeroNum}{0}
\rqSet{rqNonRepAccUnclassifiedLibPZeroDen}{0}
\rqSet{rqNonRepAccUnclassifiedLibPZeroPct}{n/a}
\rqSet{rqNonRepAccUnclassifiedLibPOneNum}{0}
\rqSet{rqNonRepAccUnclassifiedLibPOneDen}{3}
\rqSet{rqNonRepAccUnclassifiedLibPOnePct}{0.0}
\rqSet{rqNonRepRejOracleDriftLibAllNum}{2}
\rqSet{rqNonRepRejOracleDriftLibAllDen}{15}
\rqSet{rqNonRepRejOracleDriftLibAllPct}{13.3}
\rqSet{rqNonRepRejOracleDriftLibAllKiroNum}{1}
\rqSet{rqNonRepRejOracleDriftLibAllKiroDen}{2}
\rqSet{rqNonRepRejOracleDriftLibAllKiroPct}{50.0}
\rqSet{rqNonRepRejOracleDriftLibAllClaudeNum}{1}
\rqSet{rqNonRepRejOracleDriftLibAllClaudeDen}{2}
\rqSet{rqNonRepRejOracleDriftLibAllClaudePct}{50.0}
\rqSet{rqNonRepRejOracleDriftLibAllCodexNum}{0}
\rqSet{rqNonRepRejOracleDriftLibAllCodexDen}{1}
\rqSet{rqNonRepRejOracleDriftLibAllCodexPct}{0.0}
\rqSet{rqNonRepRejOracleDriftLibAllCTwoRustNum}{0}
\rqSet{rqNonRepRejOracleDriftLibAllCTwoRustDen}{1}
\rqSet{rqNonRepRejOracleDriftLibAllCTwoRustPct}{0.0}
\rqSet{rqNonRepRejOracleDriftLibAllLaertesNum}{0}
\rqSet{rqNonRepRejOracleDriftLibAllLaertesDen}{1}
\rqSet{rqNonRepRejOracleDriftLibAllLaertesPct}{0.0}
\rqSet{rqNonRepRejOracleDriftLibAllCTwoSaferRustNum}{0}
\rqSet{rqNonRepRejOracleDriftLibAllCTwoSaferRustDen}{2}
\rqSet{rqNonRepRejOracleDriftLibAllCTwoSaferRustPct}{0.0}
\rqSet{rqNonRepRejOracleDriftLibAllSmartCTwoRustNum}{0}
\rqSet{rqNonRepRejOracleDriftLibAllSmartCTwoRustDen}{0}
\rqSet{rqNonRepRejOracleDriftLibAllSmartCTwoRustPct}{n/a}
\rqSet{rqNonRepRejOracleDriftLibAllKimiNum}{0}
\rqSet{rqNonRepRejOracleDriftLibAllKimiDen}{2}
\rqSet{rqNonRepRejOracleDriftLibAllKimiPct}{0.0}
\rqSet{rqNonRepRejOracleDriftLibAllGptFiveNum}{0}
\rqSet{rqNonRepRejOracleDriftLibAllGptFiveDen}{2}
\rqSet{rqNonRepRejOracleDriftLibAllGptFivePct}{0.0}
\rqSet{rqNonRepRejOracleDriftLibAllGeminiNum}{0}
\rqSet{rqNonRepRejOracleDriftLibAllGeminiDen}{2}
\rqSet{rqNonRepRejOracleDriftLibAllGeminiPct}{0.0}
\rqSet{rqNonRepRejOracleDriftLibBOneNum}{0}
\rqSet{rqNonRepRejOracleDriftLibBOneDen}{5}
\rqSet{rqNonRepRejOracleDriftLibBOnePct}{0.0}
\rqSet{rqNonRepRejOracleDriftLibBTwoNum}{2}
\rqSet{rqNonRepRejOracleDriftLibBTwoDen}{10}
\rqSet{rqNonRepRejOracleDriftLibBTwoPct}{20.0}
\rqSet{rqNonRepRejOracleDriftLibPZeroNum}{0}
\rqSet{rqNonRepRejOracleDriftLibPZeroDen}{0}
\rqSet{rqNonRepRejOracleDriftLibPZeroPct}{n/a}
\rqSet{rqNonRepRejOracleDriftLibPOneNum}{0}
\rqSet{rqNonRepRejOracleDriftLibPOneDen}{0}
\rqSet{rqNonRepRejOracleDriftLibPOnePct}{n/a}
\rqSet{rqNonRepRejUndetectedUbLibAllNum}{0}
\rqSet{rqNonRepRejUndetectedUbLibAllDen}{15}
\rqSet{rqNonRepRejUndetectedUbLibAllPct}{0.0}
\rqSet{rqNonRepRejUndetectedUbLibAllKiroNum}{0}
\rqSet{rqNonRepRejUndetectedUbLibAllKiroDen}{2}
\rqSet{rqNonRepRejUndetectedUbLibAllKiroPct}{0.0}
\rqSet{rqNonRepRejUndetectedUbLibAllClaudeNum}{0}
\rqSet{rqNonRepRejUndetectedUbLibAllClaudeDen}{2}
\rqSet{rqNonRepRejUndetectedUbLibAllClaudePct}{0.0}
\rqSet{rqNonRepRejUndetectedUbLibAllCodexNum}{0}
\rqSet{rqNonRepRejUndetectedUbLibAllCodexDen}{1}
\rqSet{rqNonRepRejUndetectedUbLibAllCodexPct}{0.0}
\rqSet{rqNonRepRejUndetectedUbLibAllCTwoRustNum}{0}
\rqSet{rqNonRepRejUndetectedUbLibAllCTwoRustDen}{1}
\rqSet{rqNonRepRejUndetectedUbLibAllCTwoRustPct}{0.0}
\rqSet{rqNonRepRejUndetectedUbLibAllLaertesNum}{0}
\rqSet{rqNonRepRejUndetectedUbLibAllLaertesDen}{1}
\rqSet{rqNonRepRejUndetectedUbLibAllLaertesPct}{0.0}
\rqSet{rqNonRepRejUndetectedUbLibAllCTwoSaferRustNum}{0}
\rqSet{rqNonRepRejUndetectedUbLibAllCTwoSaferRustDen}{2}
\rqSet{rqNonRepRejUndetectedUbLibAllCTwoSaferRustPct}{0.0}
\rqSet{rqNonRepRejUndetectedUbLibAllSmartCTwoRustNum}{0}
\rqSet{rqNonRepRejUndetectedUbLibAllSmartCTwoRustDen}{0}
\rqSet{rqNonRepRejUndetectedUbLibAllSmartCTwoRustPct}{n/a}
\rqSet{rqNonRepRejUndetectedUbLibAllKimiNum}{0}
\rqSet{rqNonRepRejUndetectedUbLibAllKimiDen}{2}
\rqSet{rqNonRepRejUndetectedUbLibAllKimiPct}{0.0}
\rqSet{rqNonRepRejUndetectedUbLibAllGptFiveNum}{0}
\rqSet{rqNonRepRejUndetectedUbLibAllGptFiveDen}{2}
\rqSet{rqNonRepRejUndetectedUbLibAllGptFivePct}{0.0}
\rqSet{rqNonRepRejUndetectedUbLibAllGeminiNum}{0}
\rqSet{rqNonRepRejUndetectedUbLibAllGeminiDen}{2}
\rqSet{rqNonRepRejUndetectedUbLibAllGeminiPct}{0.0}
\rqSet{rqNonRepRejUndetectedUbLibBOneNum}{0}
\rqSet{rqNonRepRejUndetectedUbLibBOneDen}{5}
\rqSet{rqNonRepRejUndetectedUbLibBOnePct}{0.0}
\rqSet{rqNonRepRejUndetectedUbLibBTwoNum}{0}
\rqSet{rqNonRepRejUndetectedUbLibBTwoDen}{10}
\rqSet{rqNonRepRejUndetectedUbLibBTwoPct}{0.0}
\rqSet{rqNonRepRejUndetectedUbLibPZeroNum}{0}
\rqSet{rqNonRepRejUndetectedUbLibPZeroDen}{0}
\rqSet{rqNonRepRejUndetectedUbLibPZeroPct}{n/a}
\rqSet{rqNonRepRejUndetectedUbLibPOneNum}{0}
\rqSet{rqNonRepRejUndetectedUbLibPOneDen}{0}
\rqSet{rqNonRepRejUndetectedUbLibPOnePct}{n/a}
\rqSet{rqNonRepRejUnknownLibAllNum}{13}
\rqSet{rqNonRepRejUnknownLibAllDen}{15}
\rqSet{rqNonRepRejUnknownLibAllPct}{86.7}
\rqSet{rqNonRepRejUnknownLibAllKiroNum}{1}
\rqSet{rqNonRepRejUnknownLibAllKiroDen}{2}
\rqSet{rqNonRepRejUnknownLibAllKiroPct}{50.0}
\rqSet{rqNonRepRejUnknownLibAllClaudeNum}{1}
\rqSet{rqNonRepRejUnknownLibAllClaudeDen}{2}
\rqSet{rqNonRepRejUnknownLibAllClaudePct}{50.0}
\rqSet{rqNonRepRejUnknownLibAllCodexNum}{1}
\rqSet{rqNonRepRejUnknownLibAllCodexDen}{1}
\rqSet{rqNonRepRejUnknownLibAllCodexPct}{100.0}
\rqSet{rqNonRepRejUnknownLibAllCTwoRustNum}{1}
\rqSet{rqNonRepRejUnknownLibAllCTwoRustDen}{1}
\rqSet{rqNonRepRejUnknownLibAllCTwoRustPct}{100.0}
\rqSet{rqNonRepRejUnknownLibAllLaertesNum}{1}
\rqSet{rqNonRepRejUnknownLibAllLaertesDen}{1}
\rqSet{rqNonRepRejUnknownLibAllLaertesPct}{100.0}
\rqSet{rqNonRepRejUnknownLibAllCTwoSaferRustNum}{2}
\rqSet{rqNonRepRejUnknownLibAllCTwoSaferRustDen}{2}
\rqSet{rqNonRepRejUnknownLibAllCTwoSaferRustPct}{100.0}
\rqSet{rqNonRepRejUnknownLibAllSmartCTwoRustNum}{0}
\rqSet{rqNonRepRejUnknownLibAllSmartCTwoRustDen}{0}
\rqSet{rqNonRepRejUnknownLibAllSmartCTwoRustPct}{n/a}
\rqSet{rqNonRepRejUnknownLibAllKimiNum}{2}
\rqSet{rqNonRepRejUnknownLibAllKimiDen}{2}
\rqSet{rqNonRepRejUnknownLibAllKimiPct}{100.0}
\rqSet{rqNonRepRejUnknownLibAllGptFiveNum}{2}
\rqSet{rqNonRepRejUnknownLibAllGptFiveDen}{2}
\rqSet{rqNonRepRejUnknownLibAllGptFivePct}{100.0}
\rqSet{rqNonRepRejUnknownLibAllGeminiNum}{2}
\rqSet{rqNonRepRejUnknownLibAllGeminiDen}{2}
\rqSet{rqNonRepRejUnknownLibAllGeminiPct}{100.0}
\rqSet{rqNonRepRejUnknownLibBOneNum}{5}
\rqSet{rqNonRepRejUnknownLibBOneDen}{5}
\rqSet{rqNonRepRejUnknownLibBOnePct}{100.0}
\rqSet{rqNonRepRejUnknownLibBTwoNum}{8}
\rqSet{rqNonRepRejUnknownLibBTwoDen}{10}
\rqSet{rqNonRepRejUnknownLibBTwoPct}{80.0}
\rqSet{rqNonRepRejUnknownLibPZeroNum}{0}
\rqSet{rqNonRepRejUnknownLibPZeroDen}{0}
\rqSet{rqNonRepRejUnknownLibPZeroPct}{n/a}
\rqSet{rqNonRepRejUnknownLibPOneNum}{0}
\rqSet{rqNonRepRejUnknownLibPOneDen}{0}
\rqSet{rqNonRepRejUnknownLibPOnePct}{n/a}
\rqSet{rqNonRepRejUnclassifiedLibAllNum}{0}
\rqSet{rqNonRepRejUnclassifiedLibAllDen}{15}
\rqSet{rqNonRepRejUnclassifiedLibAllPct}{0.0}
\rqSet{rqNonRepRejUnclassifiedLibAllKiroNum}{0}
\rqSet{rqNonRepRejUnclassifiedLibAllKiroDen}{2}
\rqSet{rqNonRepRejUnclassifiedLibAllKiroPct}{0.0}
\rqSet{rqNonRepRejUnclassifiedLibAllClaudeNum}{0}
\rqSet{rqNonRepRejUnclassifiedLibAllClaudeDen}{2}
\rqSet{rqNonRepRejUnclassifiedLibAllClaudePct}{0.0}
\rqSet{rqNonRepRejUnclassifiedLibAllCodexNum}{0}
\rqSet{rqNonRepRejUnclassifiedLibAllCodexDen}{1}
\rqSet{rqNonRepRejUnclassifiedLibAllCodexPct}{0.0}
\rqSet{rqNonRepRejUnclassifiedLibAllCTwoRustNum}{0}
\rqSet{rqNonRepRejUnclassifiedLibAllCTwoRustDen}{1}
\rqSet{rqNonRepRejUnclassifiedLibAllCTwoRustPct}{0.0}
\rqSet{rqNonRepRejUnclassifiedLibAllLaertesNum}{0}
\rqSet{rqNonRepRejUnclassifiedLibAllLaertesDen}{1}
\rqSet{rqNonRepRejUnclassifiedLibAllLaertesPct}{0.0}
\rqSet{rqNonRepRejUnclassifiedLibAllCTwoSaferRustNum}{0}
\rqSet{rqNonRepRejUnclassifiedLibAllCTwoSaferRustDen}{2}
\rqSet{rqNonRepRejUnclassifiedLibAllCTwoSaferRustPct}{0.0}
\rqSet{rqNonRepRejUnclassifiedLibAllSmartCTwoRustNum}{0}
\rqSet{rqNonRepRejUnclassifiedLibAllSmartCTwoRustDen}{0}
\rqSet{rqNonRepRejUnclassifiedLibAllSmartCTwoRustPct}{n/a}
\rqSet{rqNonRepRejUnclassifiedLibAllKimiNum}{0}
\rqSet{rqNonRepRejUnclassifiedLibAllKimiDen}{2}
\rqSet{rqNonRepRejUnclassifiedLibAllKimiPct}{0.0}
\rqSet{rqNonRepRejUnclassifiedLibAllGptFiveNum}{0}
\rqSet{rqNonRepRejUnclassifiedLibAllGptFiveDen}{2}
\rqSet{rqNonRepRejUnclassifiedLibAllGptFivePct}{0.0}
\rqSet{rqNonRepRejUnclassifiedLibAllGeminiNum}{0}
\rqSet{rqNonRepRejUnclassifiedLibAllGeminiDen}{2}
\rqSet{rqNonRepRejUnclassifiedLibAllGeminiPct}{0.0}
\rqSet{rqNonRepRejUnclassifiedLibBOneNum}{0}
\rqSet{rqNonRepRejUnclassifiedLibBOneDen}{5}
\rqSet{rqNonRepRejUnclassifiedLibBOnePct}{0.0}
\rqSet{rqNonRepRejUnclassifiedLibBTwoNum}{0}
\rqSet{rqNonRepRejUnclassifiedLibBTwoDen}{10}
\rqSet{rqNonRepRejUnclassifiedLibBTwoPct}{0.0}
\rqSet{rqNonRepRejUnclassifiedLibPZeroNum}{0}
\rqSet{rqNonRepRejUnclassifiedLibPZeroDen}{0}
\rqSet{rqNonRepRejUnclassifiedLibPZeroPct}{n/a}
\rqSet{rqNonRepRejUnclassifiedLibPOneNum}{0}
\rqSet{rqNonRepRejUnclassifiedLibPOneDen}{0}
\rqSet{rqNonRepRejUnclassifiedLibPOnePct}{n/a}
\rqSet{rqTermTotalCLibAllNum}{6554751}
\rqSet{rqTermTotalCLibAllDen}{6554751}
\rqSet{rqTermTotalCLibBOneNum}{3259995}
\rqSet{rqTermTotalCLibBOneDen}{3259995}
\rqSet{rqTermTotalCLibBTwoNum}{2280249}
\rqSet{rqTermTotalCLibBTwoDen}{2280249}
\rqSet{rqTermTotalCLibPZeroNum}{0}
\rqSet{rqTermTotalCLibPZeroDen}{0}
\rqSet{rqTermTotalCLibPOneNum}{1014507}
\rqSet{rqTermTotalCLibPOneDen}{1014507}
\rqSet{rqTermTotalRLibAllNum}{6554751}
\rqSet{rqTermTotalRLibAllDen}{6554751}
\rqSet{rqTermTotalRLibBOneNum}{3259995}
\rqSet{rqTermTotalRLibBOneDen}{3259995}
\rqSet{rqTermTotalRLibBTwoNum}{2280249}
\rqSet{rqTermTotalRLibBTwoDen}{2280249}
\rqSet{rqTermTotalRLibPZeroNum}{0}
\rqSet{rqTermTotalRLibPZeroDen}{0}
\rqSet{rqTermTotalRLibPOneNum}{1014507}
\rqSet{rqTermTotalRLibPOneDen}{1014507}
\rqSet{rqVSixSepLibAllNum}{0}
\rqSet{rqVSixSepLibAllDen}{16}
\rqSet{rqVSixSepLibAllPct}{0.0}
\rqSet{rqVSixSepLibAllKiroNum}{0}
\rqSet{rqVSixSepLibAllKiroDen}{2}
\rqSet{rqVSixSepLibAllKiroPct}{0.0}
\rqSet{rqVSixSepLibAllClaudeNum}{0}
\rqSet{rqVSixSepLibAllClaudeDen}{2}
\rqSet{rqVSixSepLibAllClaudePct}{0.0}
\rqSet{rqVSixSepLibAllCodexNum}{0}
\rqSet{rqVSixSepLibAllCodexDen}{2}
\rqSet{rqVSixSepLibAllCodexPct}{0.0}
\rqSet{rqVSixSepLibAllCTwoRustNum}{0}
\rqSet{rqVSixSepLibAllCTwoRustDen}{2}
\rqSet{rqVSixSepLibAllCTwoRustPct}{0.0}
\rqSet{rqVSixSepLibAllLaertesNum}{0}
\rqSet{rqVSixSepLibAllLaertesDen}{0}
\rqSet{rqVSixSepLibAllLaertesPct}{n/a}
\rqSet{rqVSixSepLibAllCTwoSaferRustNum}{0}
\rqSet{rqVSixSepLibAllCTwoSaferRustDen}{2}
\rqSet{rqVSixSepLibAllCTwoSaferRustPct}{0.0}
\rqSet{rqVSixSepLibAllSmartCTwoRustNum}{0}
\rqSet{rqVSixSepLibAllSmartCTwoRustDen}{0}
\rqSet{rqVSixSepLibAllSmartCTwoRustPct}{n/a}
\rqSet{rqVSixSepLibAllKimiNum}{0}
\rqSet{rqVSixSepLibAllKimiDen}{2}
\rqSet{rqVSixSepLibAllKimiPct}{0.0}
\rqSet{rqVSixSepLibAllGptFiveNum}{0}
\rqSet{rqVSixSepLibAllGptFiveDen}{2}
\rqSet{rqVSixSepLibAllGptFivePct}{0.0}
\rqSet{rqVSixSepLibAllGeminiNum}{0}
\rqSet{rqVSixSepLibAllGeminiDen}{2}
\rqSet{rqVSixSepLibAllGeminiPct}{0.0}
\rqSet{rqVSixSepLibBOneNum}{0}
\rqSet{rqVSixSepLibBOneDen}{16}
\rqSet{rqVSixSepLibBOnePct}{0.0}
\rqSet{rqVSixSepLibBTwoNum}{0}
\rqSet{rqVSixSepLibBTwoDen}{0}
\rqSet{rqVSixSepLibBTwoPct}{n/a}
\rqSet{rqVSixSepLibPZeroNum}{0}
\rqSet{rqVSixSepLibPZeroDen}{0}
\rqSet{rqVSixSepLibPZeroPct}{n/a}
\rqSet{rqVSixSepLibPOneNum}{0}
\rqSet{rqVSixSepLibPOneDen}{0}
\rqSet{rqVSixSepLibPOnePct}{n/a}
\rqSet{rqVSixRestLibAllNum}{65}
\rqSet{rqVSixRestLibAllDen}{1072}
\rqSet{rqVSixRestLibAllPct}{6.1}
\rqSet{rqVSixRestLibAllKiroNum}{7}
\rqSet{rqVSixRestLibAllKiroDen}{190}
\rqSet{rqVSixRestLibAllKiroPct}{3.7}
\rqSet{rqVSixRestLibAllClaudeNum}{6}
\rqSet{rqVSixRestLibAllClaudeDen}{179}
\rqSet{rqVSixRestLibAllClaudePct}{3.4}
\rqSet{rqVSixRestLibAllCodexNum}{2}
\rqSet{rqVSixRestLibAllCodexDen}{141}
\rqSet{rqVSixRestLibAllCodexPct}{1.4}
\rqSet{rqVSixRestLibAllCTwoRustNum}{3}
\rqSet{rqVSixRestLibAllCTwoRustDen}{114}
\rqSet{rqVSixRestLibAllCTwoRustPct}{2.6}
\rqSet{rqVSixRestLibAllLaertesNum}{3}
\rqSet{rqVSixRestLibAllLaertesDen}{114}
\rqSet{rqVSixRestLibAllLaertesPct}{2.6}
\rqSet{rqVSixRestLibAllCTwoSaferRustNum}{7}
\rqSet{rqVSixRestLibAllCTwoSaferRustDen}{85}
\rqSet{rqVSixRestLibAllCTwoSaferRustPct}{8.2}
\rqSet{rqVSixRestLibAllSmartCTwoRustNum}{0}
\rqSet{rqVSixRestLibAllSmartCTwoRustDen}{0}
\rqSet{rqVSixRestLibAllSmartCTwoRustPct}{n/a}
\rqSet{rqVSixRestLibAllKimiNum}{16}
\rqSet{rqVSixRestLibAllKimiDen}{71}
\rqSet{rqVSixRestLibAllKimiPct}{22.5}
\rqSet{rqVSixRestLibAllGptFiveNum}{15}
\rqSet{rqVSixRestLibAllGptFiveDen}{90}
\rqSet{rqVSixRestLibAllGptFivePct}{16.7}
\rqSet{rqVSixRestLibAllGeminiNum}{6}
\rqSet{rqVSixRestLibAllGeminiDen}{88}
\rqSet{rqVSixRestLibAllGeminiPct}{6.8}
\rqSet{rqVSixRestLibBOneNum}{38}
\rqSet{rqVSixRestLibBOneDen}{524}
\rqSet{rqVSixRestLibBOnePct}{7.3}
\rqSet{rqVSixRestLibBTwoNum}{27}
\rqSet{rqVSixRestLibBTwoDen}{379}
\rqSet{rqVSixRestLibBTwoPct}{7.1}
\rqSet{rqVSixRestLibPZeroNum}{0}
\rqSet{rqVSixRestLibPZeroDen}{0}
\rqSet{rqVSixRestLibPZeroPct}{n/a}
\rqSet{rqVSixRestLibPOneNum}{0}
\rqSet{rqVSixRestLibPOneDen}{169}
\rqSet{rqVSixRestLibPOnePct}{0.0}
\rqSet{rqFFLibAllNum}{1}
\rqSet{rqFFLibAllDen}{111}
\rqSet{rqFFLibAllPct}{0.9}
\rqSet{rqFFLibAllKiroNum}{0}
\rqSet{rqFFLibAllKiroDen}{0}
\rqSet{rqFFLibAllKiroPct}{n/a}
\rqSet{rqFFLibAllClaudeNum}{0}
\rqSet{rqFFLibAllClaudeDen}{11}
\rqSet{rqFFLibAllClaudePct}{0.0}
\rqSet{rqFFLibAllCodexNum}{0}
\rqSet{rqFFLibAllCodexDen}{45}
\rqSet{rqFFLibAllCodexPct}{0.0}
\rqSet{rqFFLibAllCTwoRustNum}{0}
\rqSet{rqFFLibAllCTwoRustDen}{0}
\rqSet{rqFFLibAllCTwoRustPct}{n/a}
\rqSet{rqFFLibAllLaertesNum}{0}
\rqSet{rqFFLibAllLaertesDen}{0}
\rqSet{rqFFLibAllLaertesPct}{n/a}
\rqSet{rqFFLibAllCTwoSaferRustNum}{0}
\rqSet{rqFFLibAllCTwoSaferRustDen}{23}
\rqSet{rqFFLibAllCTwoSaferRustPct}{0.0}
\rqSet{rqFFLibAllSmartCTwoRustNum}{0}
\rqSet{rqFFLibAllSmartCTwoRustDen}{0}
\rqSet{rqFFLibAllSmartCTwoRustPct}{n/a}
\rqSet{rqFFLibAllKimiNum}{0}
\rqSet{rqFFLibAllKimiDen}{15}
\rqSet{rqFFLibAllKimiPct}{0.0}
\rqSet{rqFFLibAllGptFiveNum}{1}
\rqSet{rqFFLibAllGptFiveDen}{13}
\rqSet{rqFFLibAllGptFivePct}{7.7}
\rqSet{rqFFLibAllGeminiNum}{0}
\rqSet{rqFFLibAllGeminiDen}{4}
\rqSet{rqFFLibAllGeminiPct}{0.0}
\rqSet{rqFFLibBOneNum}{0}
\rqSet{rqFFLibBOneDen}{31}
\rqSet{rqFFLibBOnePct}{0.0}
\rqSet{rqFFLibBOneKiroNum}{0}
\rqSet{rqFFLibBOneKiroDen}{0}
\rqSet{rqFFLibBOneKiroPct}{n/a}
\rqSet{rqFFLibBOneClaudeNum}{0}
\rqSet{rqFFLibBOneClaudeDen}{0}
\rqSet{rqFFLibBOneClaudePct}{n/a}
\rqSet{rqFFLibBOneCodexNum}{0}
\rqSet{rqFFLibBOneCodexDen}{0}
\rqSet{rqFFLibBOneCodexPct}{n/a}
\rqSet{rqFFLibBOneCTwoRustNum}{0}
\rqSet{rqFFLibBOneCTwoRustDen}{0}
\rqSet{rqFFLibBOneCTwoRustPct}{n/a}
\rqSet{rqFFLibBOneLaertesNum}{0}
\rqSet{rqFFLibBOneLaertesDen}{0}
\rqSet{rqFFLibBOneLaertesPct}{n/a}
\rqSet{rqFFLibBOneCTwoSaferRustNum}{0}
\rqSet{rqFFLibBOneCTwoSaferRustDen}{16}
\rqSet{rqFFLibBOneCTwoSaferRustPct}{0.0}
\rqSet{rqFFLibBOneSmartCTwoRustNum}{0}
\rqSet{rqFFLibBOneSmartCTwoRustDen}{0}
\rqSet{rqFFLibBOneSmartCTwoRustPct}{n/a}
\rqSet{rqFFLibBOneKimiNum}{0}
\rqSet{rqFFLibBOneKimiDen}{9}
\rqSet{rqFFLibBOneKimiPct}{0.0}
\rqSet{rqFFLibBOneGptFiveNum}{0}
\rqSet{rqFFLibBOneGptFiveDen}{5}
\rqSet{rqFFLibBOneGptFivePct}{0.0}
\rqSet{rqFFLibBOneGeminiNum}{0}
\rqSet{rqFFLibBOneGeminiDen}{1}
\rqSet{rqFFLibBOneGeminiPct}{0.0}
\rqSet{rqFFLibBTwoNum}{1}
\rqSet{rqFFLibBTwoDen}{24}
\rqSet{rqFFLibBTwoPct}{4.2}
\rqSet{rqFFLibBTwoKiroNum}{0}
\rqSet{rqFFLibBTwoKiroDen}{0}
\rqSet{rqFFLibBTwoKiroPct}{n/a}
\rqSet{rqFFLibBTwoClaudeNum}{0}
\rqSet{rqFFLibBTwoClaudeDen}{0}
\rqSet{rqFFLibBTwoClaudePct}{n/a}
\rqSet{rqFFLibBTwoCodexNum}{0}
\rqSet{rqFFLibBTwoCodexDen}{0}
\rqSet{rqFFLibBTwoCodexPct}{n/a}
\rqSet{rqFFLibBTwoCTwoRustNum}{0}
\rqSet{rqFFLibBTwoCTwoRustDen}{0}
\rqSet{rqFFLibBTwoCTwoRustPct}{n/a}
\rqSet{rqFFLibBTwoLaertesNum}{0}
\rqSet{rqFFLibBTwoLaertesDen}{0}
\rqSet{rqFFLibBTwoLaertesPct}{n/a}
\rqSet{rqFFLibBTwoCTwoSaferRustNum}{0}
\rqSet{rqFFLibBTwoCTwoSaferRustDen}{7}
\rqSet{rqFFLibBTwoCTwoSaferRustPct}{0.0}
\rqSet{rqFFLibBTwoSmartCTwoRustNum}{0}
\rqSet{rqFFLibBTwoSmartCTwoRustDen}{0}
\rqSet{rqFFLibBTwoSmartCTwoRustPct}{n/a}
\rqSet{rqFFLibBTwoKimiNum}{0}
\rqSet{rqFFLibBTwoKimiDen}{6}
\rqSet{rqFFLibBTwoKimiPct}{0.0}
\rqSet{rqFFLibBTwoGptFiveNum}{1}
\rqSet{rqFFLibBTwoGptFiveDen}{8}
\rqSet{rqFFLibBTwoGptFivePct}{12.5}
\rqSet{rqFFLibBTwoGeminiNum}{0}
\rqSet{rqFFLibBTwoGeminiDen}{3}
\rqSet{rqFFLibBTwoGeminiPct}{0.0}
\rqSet{rqFFLibPZeroNum}{0}
\rqSet{rqFFLibPZeroDen}{0}
\rqSet{rqFFLibPZeroPct}{n/a}
\rqSet{rqFFLibPZeroKiroNum}{0}
\rqSet{rqFFLibPZeroKiroDen}{0}
\rqSet{rqFFLibPZeroKiroPct}{n/a}
\rqSet{rqFFLibPZeroClaudeNum}{0}
\rqSet{rqFFLibPZeroClaudeDen}{0}
\rqSet{rqFFLibPZeroClaudePct}{n/a}
\rqSet{rqFFLibPZeroCodexNum}{0}
\rqSet{rqFFLibPZeroCodexDen}{0}
\rqSet{rqFFLibPZeroCodexPct}{n/a}
\rqSet{rqFFLibPZeroCTwoRustNum}{0}
\rqSet{rqFFLibPZeroCTwoRustDen}{0}
\rqSet{rqFFLibPZeroCTwoRustPct}{n/a}
\rqSet{rqFFLibPZeroLaertesNum}{0}
\rqSet{rqFFLibPZeroLaertesDen}{0}
\rqSet{rqFFLibPZeroLaertesPct}{n/a}
\rqSet{rqFFLibPZeroCTwoSaferRustNum}{0}
\rqSet{rqFFLibPZeroCTwoSaferRustDen}{0}
\rqSet{rqFFLibPZeroCTwoSaferRustPct}{n/a}
\rqSet{rqFFLibPZeroSmartCTwoRustNum}{0}
\rqSet{rqFFLibPZeroSmartCTwoRustDen}{0}
\rqSet{rqFFLibPZeroSmartCTwoRustPct}{n/a}
\rqSet{rqFFLibPZeroKimiNum}{0}
\rqSet{rqFFLibPZeroKimiDen}{0}
\rqSet{rqFFLibPZeroKimiPct}{n/a}
\rqSet{rqFFLibPZeroGptFiveNum}{0}
\rqSet{rqFFLibPZeroGptFiveDen}{0}
\rqSet{rqFFLibPZeroGptFivePct}{n/a}
\rqSet{rqFFLibPZeroGeminiNum}{0}
\rqSet{rqFFLibPZeroGeminiDen}{0}
\rqSet{rqFFLibPZeroGeminiPct}{n/a}
\rqSet{rqFFLibPOneNum}{0}
\rqSet{rqFFLibPOneDen}{56}
\rqSet{rqFFLibPOnePct}{0.0}
\rqSet{rqFFLibPOneKiroNum}{0}
\rqSet{rqFFLibPOneKiroDen}{0}
\rqSet{rqFFLibPOneKiroPct}{n/a}
\rqSet{rqFFLibPOneClaudeNum}{0}
\rqSet{rqFFLibPOneClaudeDen}{11}
\rqSet{rqFFLibPOneClaudePct}{0.0}
\rqSet{rqFFLibPOneCodexNum}{0}
\rqSet{rqFFLibPOneCodexDen}{45}
\rqSet{rqFFLibPOneCodexPct}{0.0}
\rqSet{rqFFLibPOneCTwoRustNum}{0}
\rqSet{rqFFLibPOneCTwoRustDen}{0}
\rqSet{rqFFLibPOneCTwoRustPct}{n/a}
\rqSet{rqFFLibPOneLaertesNum}{0}
\rqSet{rqFFLibPOneLaertesDen}{0}
\rqSet{rqFFLibPOneLaertesPct}{n/a}
\rqSet{rqFFLibPOneCTwoSaferRustNum}{0}
\rqSet{rqFFLibPOneCTwoSaferRustDen}{0}
\rqSet{rqFFLibPOneCTwoSaferRustPct}{n/a}
\rqSet{rqFFLibPOneSmartCTwoRustNum}{0}
\rqSet{rqFFLibPOneSmartCTwoRustDen}{0}
\rqSet{rqFFLibPOneSmartCTwoRustPct}{n/a}
\rqSet{rqFFLibPOneKimiNum}{0}
\rqSet{rqFFLibPOneKimiDen}{0}
\rqSet{rqFFLibPOneKimiPct}{n/a}
\rqSet{rqFFLibPOneGptFiveNum}{0}
\rqSet{rqFFLibPOneGptFiveDen}{0}
\rqSet{rqFFLibPOneGptFivePct}{n/a}
\rqSet{rqFFLibPOneGeminiNum}{0}
\rqSet{rqFFLibPOneGeminiDen}{0}
\rqSet{rqFFLibPOneGeminiPct}{n/a}
\rqSet{rqFFDisagreeLibAllNum}{110}
\rqSet{rqFFDisagreeLibAllDen}{111}
\rqSet{rqFFDisagreeLibAllPct}{99.1}
\rqSet{rqFFDisagreeLibAllKiroNum}{0}
\rqSet{rqFFDisagreeLibAllKiroDen}{0}
\rqSet{rqFFDisagreeLibAllKiroPct}{n/a}
\rqSet{rqFFDisagreeLibAllClaudeNum}{11}
\rqSet{rqFFDisagreeLibAllClaudeDen}{11}
\rqSet{rqFFDisagreeLibAllClaudePct}{100.0}
\rqSet{rqFFDisagreeLibAllCodexNum}{45}
\rqSet{rqFFDisagreeLibAllCodexDen}{45}
\rqSet{rqFFDisagreeLibAllCodexPct}{100.0}
\rqSet{rqFFDisagreeLibAllCTwoRustNum}{0}
\rqSet{rqFFDisagreeLibAllCTwoRustDen}{0}
\rqSet{rqFFDisagreeLibAllCTwoRustPct}{n/a}
\rqSet{rqFFDisagreeLibAllLaertesNum}{0}
\rqSet{rqFFDisagreeLibAllLaertesDen}{0}
\rqSet{rqFFDisagreeLibAllLaertesPct}{n/a}
\rqSet{rqFFDisagreeLibAllCTwoSaferRustNum}{23}
\rqSet{rqFFDisagreeLibAllCTwoSaferRustDen}{23}
\rqSet{rqFFDisagreeLibAllCTwoSaferRustPct}{100.0}
\rqSet{rqFFDisagreeLibAllSmartCTwoRustNum}{0}
\rqSet{rqFFDisagreeLibAllSmartCTwoRustDen}{0}
\rqSet{rqFFDisagreeLibAllSmartCTwoRustPct}{n/a}
\rqSet{rqFFDisagreeLibAllKimiNum}{15}
\rqSet{rqFFDisagreeLibAllKimiDen}{15}
\rqSet{rqFFDisagreeLibAllKimiPct}{100.0}
\rqSet{rqFFDisagreeLibAllGptFiveNum}{12}
\rqSet{rqFFDisagreeLibAllGptFiveDen}{13}
\rqSet{rqFFDisagreeLibAllGptFivePct}{92.3}
\rqSet{rqFFDisagreeLibAllGeminiNum}{4}
\rqSet{rqFFDisagreeLibAllGeminiDen}{4}
\rqSet{rqFFDisagreeLibAllGeminiPct}{100.0}
\rqSet{rqFFDisagreeLibBOneNum}{31}
\rqSet{rqFFDisagreeLibBOneDen}{31}
\rqSet{rqFFDisagreeLibBOnePct}{100.0}
\rqSet{rqFFDisagreeLibBOneKiroNum}{0}
\rqSet{rqFFDisagreeLibBOneKiroDen}{0}
\rqSet{rqFFDisagreeLibBOneKiroPct}{n/a}
\rqSet{rqFFDisagreeLibBOneClaudeNum}{0}
\rqSet{rqFFDisagreeLibBOneClaudeDen}{0}
\rqSet{rqFFDisagreeLibBOneClaudePct}{n/a}
\rqSet{rqFFDisagreeLibBOneCodexNum}{0}
\rqSet{rqFFDisagreeLibBOneCodexDen}{0}
\rqSet{rqFFDisagreeLibBOneCodexPct}{n/a}
\rqSet{rqFFDisagreeLibBOneCTwoRustNum}{0}
\rqSet{rqFFDisagreeLibBOneCTwoRustDen}{0}
\rqSet{rqFFDisagreeLibBOneCTwoRustPct}{n/a}
\rqSet{rqFFDisagreeLibBOneLaertesNum}{0}
\rqSet{rqFFDisagreeLibBOneLaertesDen}{0}
\rqSet{rqFFDisagreeLibBOneLaertesPct}{n/a}
\rqSet{rqFFDisagreeLibBOneCTwoSaferRustNum}{16}
\rqSet{rqFFDisagreeLibBOneCTwoSaferRustDen}{16}
\rqSet{rqFFDisagreeLibBOneCTwoSaferRustPct}{100.0}
\rqSet{rqFFDisagreeLibBOneSmartCTwoRustNum}{0}
\rqSet{rqFFDisagreeLibBOneSmartCTwoRustDen}{0}
\rqSet{rqFFDisagreeLibBOneSmartCTwoRustPct}{n/a}
\rqSet{rqFFDisagreeLibBOneKimiNum}{9}
\rqSet{rqFFDisagreeLibBOneKimiDen}{9}
\rqSet{rqFFDisagreeLibBOneKimiPct}{100.0}
\rqSet{rqFFDisagreeLibBOneGptFiveNum}{5}
\rqSet{rqFFDisagreeLibBOneGptFiveDen}{5}
\rqSet{rqFFDisagreeLibBOneGptFivePct}{100.0}
\rqSet{rqFFDisagreeLibBOneGeminiNum}{1}
\rqSet{rqFFDisagreeLibBOneGeminiDen}{1}
\rqSet{rqFFDisagreeLibBOneGeminiPct}{100.0}
\rqSet{rqFFDisagreeLibBTwoNum}{23}
\rqSet{rqFFDisagreeLibBTwoDen}{24}
\rqSet{rqFFDisagreeLibBTwoPct}{95.8}
\rqSet{rqFFDisagreeLibBTwoKiroNum}{0}
\rqSet{rqFFDisagreeLibBTwoKiroDen}{0}
\rqSet{rqFFDisagreeLibBTwoKiroPct}{n/a}
\rqSet{rqFFDisagreeLibBTwoClaudeNum}{0}
\rqSet{rqFFDisagreeLibBTwoClaudeDen}{0}
\rqSet{rqFFDisagreeLibBTwoClaudePct}{n/a}
\rqSet{rqFFDisagreeLibBTwoCodexNum}{0}
\rqSet{rqFFDisagreeLibBTwoCodexDen}{0}
\rqSet{rqFFDisagreeLibBTwoCodexPct}{n/a}
\rqSet{rqFFDisagreeLibBTwoCTwoRustNum}{0}
\rqSet{rqFFDisagreeLibBTwoCTwoRustDen}{0}
\rqSet{rqFFDisagreeLibBTwoCTwoRustPct}{n/a}
\rqSet{rqFFDisagreeLibBTwoLaertesNum}{0}
\rqSet{rqFFDisagreeLibBTwoLaertesDen}{0}
\rqSet{rqFFDisagreeLibBTwoLaertesPct}{n/a}
\rqSet{rqFFDisagreeLibBTwoCTwoSaferRustNum}{7}
\rqSet{rqFFDisagreeLibBTwoCTwoSaferRustDen}{7}
\rqSet{rqFFDisagreeLibBTwoCTwoSaferRustPct}{100.0}
\rqSet{rqFFDisagreeLibBTwoSmartCTwoRustNum}{0}
\rqSet{rqFFDisagreeLibBTwoSmartCTwoRustDen}{0}
\rqSet{rqFFDisagreeLibBTwoSmartCTwoRustPct}{n/a}
\rqSet{rqFFDisagreeLibBTwoKimiNum}{6}
\rqSet{rqFFDisagreeLibBTwoKimiDen}{6}
\rqSet{rqFFDisagreeLibBTwoKimiPct}{100.0}
\rqSet{rqFFDisagreeLibBTwoGptFiveNum}{7}
\rqSet{rqFFDisagreeLibBTwoGptFiveDen}{8}
\rqSet{rqFFDisagreeLibBTwoGptFivePct}{87.5}
\rqSet{rqFFDisagreeLibBTwoGeminiNum}{3}
\rqSet{rqFFDisagreeLibBTwoGeminiDen}{3}
\rqSet{rqFFDisagreeLibBTwoGeminiPct}{100.0}
\rqSet{rqFFDisagreeLibPZeroNum}{0}
\rqSet{rqFFDisagreeLibPZeroDen}{0}
\rqSet{rqFFDisagreeLibPZeroPct}{n/a}
\rqSet{rqFFDisagreeLibPZeroKiroNum}{0}
\rqSet{rqFFDisagreeLibPZeroKiroDen}{0}
\rqSet{rqFFDisagreeLibPZeroKiroPct}{n/a}
\rqSet{rqFFDisagreeLibPZeroClaudeNum}{0}
\rqSet{rqFFDisagreeLibPZeroClaudeDen}{0}
\rqSet{rqFFDisagreeLibPZeroClaudePct}{n/a}
\rqSet{rqFFDisagreeLibPZeroCodexNum}{0}
\rqSet{rqFFDisagreeLibPZeroCodexDen}{0}
\rqSet{rqFFDisagreeLibPZeroCodexPct}{n/a}
\rqSet{rqFFDisagreeLibPZeroCTwoRustNum}{0}
\rqSet{rqFFDisagreeLibPZeroCTwoRustDen}{0}
\rqSet{rqFFDisagreeLibPZeroCTwoRustPct}{n/a}
\rqSet{rqFFDisagreeLibPZeroLaertesNum}{0}
\rqSet{rqFFDisagreeLibPZeroLaertesDen}{0}
\rqSet{rqFFDisagreeLibPZeroLaertesPct}{n/a}
\rqSet{rqFFDisagreeLibPZeroCTwoSaferRustNum}{0}
\rqSet{rqFFDisagreeLibPZeroCTwoSaferRustDen}{0}
\rqSet{rqFFDisagreeLibPZeroCTwoSaferRustPct}{n/a}
\rqSet{rqFFDisagreeLibPZeroSmartCTwoRustNum}{0}
\rqSet{rqFFDisagreeLibPZeroSmartCTwoRustDen}{0}
\rqSet{rqFFDisagreeLibPZeroSmartCTwoRustPct}{n/a}
\rqSet{rqFFDisagreeLibPZeroKimiNum}{0}
\rqSet{rqFFDisagreeLibPZeroKimiDen}{0}
\rqSet{rqFFDisagreeLibPZeroKimiPct}{n/a}
\rqSet{rqFFDisagreeLibPZeroGptFiveNum}{0}
\rqSet{rqFFDisagreeLibPZeroGptFiveDen}{0}
\rqSet{rqFFDisagreeLibPZeroGptFivePct}{n/a}
\rqSet{rqFFDisagreeLibPZeroGeminiNum}{0}
\rqSet{rqFFDisagreeLibPZeroGeminiDen}{0}
\rqSet{rqFFDisagreeLibPZeroGeminiPct}{n/a}
\rqSet{rqFFDisagreeLibPOneNum}{56}
\rqSet{rqFFDisagreeLibPOneDen}{56}
\rqSet{rqFFDisagreeLibPOnePct}{100.0}
\rqSet{rqFFDisagreeLibPOneKiroNum}{0}
\rqSet{rqFFDisagreeLibPOneKiroDen}{0}
\rqSet{rqFFDisagreeLibPOneKiroPct}{n/a}
\rqSet{rqFFDisagreeLibPOneClaudeNum}{11}
\rqSet{rqFFDisagreeLibPOneClaudeDen}{11}
\rqSet{rqFFDisagreeLibPOneClaudePct}{100.0}
\rqSet{rqFFDisagreeLibPOneCodexNum}{45}
\rqSet{rqFFDisagreeLibPOneCodexDen}{45}
\rqSet{rqFFDisagreeLibPOneCodexPct}{100.0}
\rqSet{rqFFDisagreeLibPOneCTwoRustNum}{0}
\rqSet{rqFFDisagreeLibPOneCTwoRustDen}{0}
\rqSet{rqFFDisagreeLibPOneCTwoRustPct}{n/a}
\rqSet{rqFFDisagreeLibPOneLaertesNum}{0}
\rqSet{rqFFDisagreeLibPOneLaertesDen}{0}
\rqSet{rqFFDisagreeLibPOneLaertesPct}{n/a}
\rqSet{rqFFDisagreeLibPOneCTwoSaferRustNum}{0}
\rqSet{rqFFDisagreeLibPOneCTwoSaferRustDen}{0}
\rqSet{rqFFDisagreeLibPOneCTwoSaferRustPct}{n/a}
\rqSet{rqFFDisagreeLibPOneSmartCTwoRustNum}{0}
\rqSet{rqFFDisagreeLibPOneSmartCTwoRustDen}{0}
\rqSet{rqFFDisagreeLibPOneSmartCTwoRustPct}{n/a}
\rqSet{rqFFDisagreeLibPOneKimiNum}{0}
\rqSet{rqFFDisagreeLibPOneKimiDen}{0}
\rqSet{rqFFDisagreeLibPOneKimiPct}{n/a}
\rqSet{rqFFDisagreeLibPOneGptFiveNum}{0}
\rqSet{rqFFDisagreeLibPOneGptFiveDen}{0}
\rqSet{rqFFDisagreeLibPOneGptFivePct}{n/a}
\rqSet{rqFFDisagreeLibPOneGeminiNum}{0}
\rqSet{rqFFDisagreeLibPOneGeminiDen}{0}
\rqSet{rqFFDisagreeLibPOneGeminiPct}{n/a}
\rqSet{rqFFNotAdjLibAllNum}{0}
\rqSet{rqFFNotAdjLibAllDen}{111}
\rqSet{rqFFNotAdjLibAllPct}{0.0}
\rqSet{rqFFNotAdjLibAllKiroNum}{0}
\rqSet{rqFFNotAdjLibAllKiroDen}{0}
\rqSet{rqFFNotAdjLibAllKiroPct}{n/a}
\rqSet{rqFFNotAdjLibAllClaudeNum}{0}
\rqSet{rqFFNotAdjLibAllClaudeDen}{11}
\rqSet{rqFFNotAdjLibAllClaudePct}{0.0}
\rqSet{rqFFNotAdjLibAllCodexNum}{0}
\rqSet{rqFFNotAdjLibAllCodexDen}{45}
\rqSet{rqFFNotAdjLibAllCodexPct}{0.0}
\rqSet{rqFFNotAdjLibAllCTwoRustNum}{0}
\rqSet{rqFFNotAdjLibAllCTwoRustDen}{0}
\rqSet{rqFFNotAdjLibAllCTwoRustPct}{n/a}
\rqSet{rqFFNotAdjLibAllLaertesNum}{0}
\rqSet{rqFFNotAdjLibAllLaertesDen}{0}
\rqSet{rqFFNotAdjLibAllLaertesPct}{n/a}
\rqSet{rqFFNotAdjLibAllCTwoSaferRustNum}{0}
\rqSet{rqFFNotAdjLibAllCTwoSaferRustDen}{23}
\rqSet{rqFFNotAdjLibAllCTwoSaferRustPct}{0.0}
\rqSet{rqFFNotAdjLibAllSmartCTwoRustNum}{0}
\rqSet{rqFFNotAdjLibAllSmartCTwoRustDen}{0}
\rqSet{rqFFNotAdjLibAllSmartCTwoRustPct}{n/a}
\rqSet{rqFFNotAdjLibAllKimiNum}{0}
\rqSet{rqFFNotAdjLibAllKimiDen}{15}
\rqSet{rqFFNotAdjLibAllKimiPct}{0.0}
\rqSet{rqFFNotAdjLibAllGptFiveNum}{0}
\rqSet{rqFFNotAdjLibAllGptFiveDen}{13}
\rqSet{rqFFNotAdjLibAllGptFivePct}{0.0}
\rqSet{rqFFNotAdjLibAllGeminiNum}{0}
\rqSet{rqFFNotAdjLibAllGeminiDen}{4}
\rqSet{rqFFNotAdjLibAllGeminiPct}{0.0}
\rqSet{rqFFNotAdjLibBOneNum}{0}
\rqSet{rqFFNotAdjLibBOneDen}{31}
\rqSet{rqFFNotAdjLibBOnePct}{0.0}
\rqSet{rqFFNotAdjLibBOneKiroNum}{0}
\rqSet{rqFFNotAdjLibBOneKiroDen}{0}
\rqSet{rqFFNotAdjLibBOneKiroPct}{n/a}
\rqSet{rqFFNotAdjLibBOneClaudeNum}{0}
\rqSet{rqFFNotAdjLibBOneClaudeDen}{0}
\rqSet{rqFFNotAdjLibBOneClaudePct}{n/a}
\rqSet{rqFFNotAdjLibBOneCodexNum}{0}
\rqSet{rqFFNotAdjLibBOneCodexDen}{0}
\rqSet{rqFFNotAdjLibBOneCodexPct}{n/a}
\rqSet{rqFFNotAdjLibBOneCTwoRustNum}{0}
\rqSet{rqFFNotAdjLibBOneCTwoRustDen}{0}
\rqSet{rqFFNotAdjLibBOneCTwoRustPct}{n/a}
\rqSet{rqFFNotAdjLibBOneLaertesNum}{0}
\rqSet{rqFFNotAdjLibBOneLaertesDen}{0}
\rqSet{rqFFNotAdjLibBOneLaertesPct}{n/a}
\rqSet{rqFFNotAdjLibBOneCTwoSaferRustNum}{0}
\rqSet{rqFFNotAdjLibBOneCTwoSaferRustDen}{16}
\rqSet{rqFFNotAdjLibBOneCTwoSaferRustPct}{0.0}
\rqSet{rqFFNotAdjLibBOneSmartCTwoRustNum}{0}
\rqSet{rqFFNotAdjLibBOneSmartCTwoRustDen}{0}
\rqSet{rqFFNotAdjLibBOneSmartCTwoRustPct}{n/a}
\rqSet{rqFFNotAdjLibBOneKimiNum}{0}
\rqSet{rqFFNotAdjLibBOneKimiDen}{9}
\rqSet{rqFFNotAdjLibBOneKimiPct}{0.0}
\rqSet{rqFFNotAdjLibBOneGptFiveNum}{0}
\rqSet{rqFFNotAdjLibBOneGptFiveDen}{5}
\rqSet{rqFFNotAdjLibBOneGptFivePct}{0.0}
\rqSet{rqFFNotAdjLibBOneGeminiNum}{0}
\rqSet{rqFFNotAdjLibBOneGeminiDen}{1}
\rqSet{rqFFNotAdjLibBOneGeminiPct}{0.0}
\rqSet{rqFFNotAdjLibBTwoNum}{0}
\rqSet{rqFFNotAdjLibBTwoDen}{24}
\rqSet{rqFFNotAdjLibBTwoPct}{0.0}
\rqSet{rqFFNotAdjLibBTwoKiroNum}{0}
\rqSet{rqFFNotAdjLibBTwoKiroDen}{0}
\rqSet{rqFFNotAdjLibBTwoKiroPct}{n/a}
\rqSet{rqFFNotAdjLibBTwoClaudeNum}{0}
\rqSet{rqFFNotAdjLibBTwoClaudeDen}{0}
\rqSet{rqFFNotAdjLibBTwoClaudePct}{n/a}
\rqSet{rqFFNotAdjLibBTwoCodexNum}{0}
\rqSet{rqFFNotAdjLibBTwoCodexDen}{0}
\rqSet{rqFFNotAdjLibBTwoCodexPct}{n/a}
\rqSet{rqFFNotAdjLibBTwoCTwoRustNum}{0}
\rqSet{rqFFNotAdjLibBTwoCTwoRustDen}{0}
\rqSet{rqFFNotAdjLibBTwoCTwoRustPct}{n/a}
\rqSet{rqFFNotAdjLibBTwoLaertesNum}{0}
\rqSet{rqFFNotAdjLibBTwoLaertesDen}{0}
\rqSet{rqFFNotAdjLibBTwoLaertesPct}{n/a}
\rqSet{rqFFNotAdjLibBTwoCTwoSaferRustNum}{0}
\rqSet{rqFFNotAdjLibBTwoCTwoSaferRustDen}{7}
\rqSet{rqFFNotAdjLibBTwoCTwoSaferRustPct}{0.0}
\rqSet{rqFFNotAdjLibBTwoSmartCTwoRustNum}{0}
\rqSet{rqFFNotAdjLibBTwoSmartCTwoRustDen}{0}
\rqSet{rqFFNotAdjLibBTwoSmartCTwoRustPct}{n/a}
\rqSet{rqFFNotAdjLibBTwoKimiNum}{0}
\rqSet{rqFFNotAdjLibBTwoKimiDen}{6}
\rqSet{rqFFNotAdjLibBTwoKimiPct}{0.0}
\rqSet{rqFFNotAdjLibBTwoGptFiveNum}{0}
\rqSet{rqFFNotAdjLibBTwoGptFiveDen}{8}
\rqSet{rqFFNotAdjLibBTwoGptFivePct}{0.0}
\rqSet{rqFFNotAdjLibBTwoGeminiNum}{0}
\rqSet{rqFFNotAdjLibBTwoGeminiDen}{3}
\rqSet{rqFFNotAdjLibBTwoGeminiPct}{0.0}
\rqSet{rqFFNotAdjLibPZeroNum}{0}
\rqSet{rqFFNotAdjLibPZeroDen}{0}
\rqSet{rqFFNotAdjLibPZeroPct}{n/a}
\rqSet{rqFFNotAdjLibPZeroKiroNum}{0}
\rqSet{rqFFNotAdjLibPZeroKiroDen}{0}
\rqSet{rqFFNotAdjLibPZeroKiroPct}{n/a}
\rqSet{rqFFNotAdjLibPZeroClaudeNum}{0}
\rqSet{rqFFNotAdjLibPZeroClaudeDen}{0}
\rqSet{rqFFNotAdjLibPZeroClaudePct}{n/a}
\rqSet{rqFFNotAdjLibPZeroCodexNum}{0}
\rqSet{rqFFNotAdjLibPZeroCodexDen}{0}
\rqSet{rqFFNotAdjLibPZeroCodexPct}{n/a}
\rqSet{rqFFNotAdjLibPZeroCTwoRustNum}{0}
\rqSet{rqFFNotAdjLibPZeroCTwoRustDen}{0}
\rqSet{rqFFNotAdjLibPZeroCTwoRustPct}{n/a}
\rqSet{rqFFNotAdjLibPZeroLaertesNum}{0}
\rqSet{rqFFNotAdjLibPZeroLaertesDen}{0}
\rqSet{rqFFNotAdjLibPZeroLaertesPct}{n/a}
\rqSet{rqFFNotAdjLibPZeroCTwoSaferRustNum}{0}
\rqSet{rqFFNotAdjLibPZeroCTwoSaferRustDen}{0}
\rqSet{rqFFNotAdjLibPZeroCTwoSaferRustPct}{n/a}
\rqSet{rqFFNotAdjLibPZeroSmartCTwoRustNum}{0}
\rqSet{rqFFNotAdjLibPZeroSmartCTwoRustDen}{0}
\rqSet{rqFFNotAdjLibPZeroSmartCTwoRustPct}{n/a}
\rqSet{rqFFNotAdjLibPZeroKimiNum}{0}
\rqSet{rqFFNotAdjLibPZeroKimiDen}{0}
\rqSet{rqFFNotAdjLibPZeroKimiPct}{n/a}
\rqSet{rqFFNotAdjLibPZeroGptFiveNum}{0}
\rqSet{rqFFNotAdjLibPZeroGptFiveDen}{0}
\rqSet{rqFFNotAdjLibPZeroGptFivePct}{n/a}
\rqSet{rqFFNotAdjLibPZeroGeminiNum}{0}
\rqSet{rqFFNotAdjLibPZeroGeminiDen}{0}
\rqSet{rqFFNotAdjLibPZeroGeminiPct}{n/a}
\rqSet{rqFFNotAdjLibPOneNum}{0}
\rqSet{rqFFNotAdjLibPOneDen}{56}
\rqSet{rqFFNotAdjLibPOnePct}{0.0}
\rqSet{rqFFNotAdjLibPOneKiroNum}{0}
\rqSet{rqFFNotAdjLibPOneKiroDen}{0}
\rqSet{rqFFNotAdjLibPOneKiroPct}{n/a}
\rqSet{rqFFNotAdjLibPOneClaudeNum}{0}
\rqSet{rqFFNotAdjLibPOneClaudeDen}{11}
\rqSet{rqFFNotAdjLibPOneClaudePct}{0.0}
\rqSet{rqFFNotAdjLibPOneCodexNum}{0}
\rqSet{rqFFNotAdjLibPOneCodexDen}{45}
\rqSet{rqFFNotAdjLibPOneCodexPct}{0.0}
\rqSet{rqFFNotAdjLibPOneCTwoRustNum}{0}
\rqSet{rqFFNotAdjLibPOneCTwoRustDen}{0}
\rqSet{rqFFNotAdjLibPOneCTwoRustPct}{n/a}
\rqSet{rqFFNotAdjLibPOneLaertesNum}{0}
\rqSet{rqFFNotAdjLibPOneLaertesDen}{0}
\rqSet{rqFFNotAdjLibPOneLaertesPct}{n/a}
\rqSet{rqFFNotAdjLibPOneCTwoSaferRustNum}{0}
\rqSet{rqFFNotAdjLibPOneCTwoSaferRustDen}{0}
\rqSet{rqFFNotAdjLibPOneCTwoSaferRustPct}{n/a}
\rqSet{rqFFNotAdjLibPOneSmartCTwoRustNum}{0}
\rqSet{rqFFNotAdjLibPOneSmartCTwoRustDen}{0}
\rqSet{rqFFNotAdjLibPOneSmartCTwoRustPct}{n/a}
\rqSet{rqFFNotAdjLibPOneKimiNum}{0}
\rqSet{rqFFNotAdjLibPOneKimiDen}{0}
\rqSet{rqFFNotAdjLibPOneKimiPct}{n/a}
\rqSet{rqFFNotAdjLibPOneGptFiveNum}{0}
\rqSet{rqFFNotAdjLibPOneGptFiveDen}{0}
\rqSet{rqFFNotAdjLibPOneGptFivePct}{n/a}
\rqSet{rqFFNotAdjLibPOneGeminiNum}{0}
\rqSet{rqFFNotAdjLibPOneGeminiDen}{0}
\rqSet{rqFFNotAdjLibPOneGeminiPct}{n/a}
\rqSet{rqFFNotReexLibAllNum}{0}
\rqSet{rqFFNotReexLibAllDen}{111}
\rqSet{rqFFNotReexLibAllPct}{0.0}
\rqSet{rqFFNotReexLibAllKiroNum}{0}
\rqSet{rqFFNotReexLibAllKiroDen}{0}
\rqSet{rqFFNotReexLibAllKiroPct}{n/a}
\rqSet{rqFFNotReexLibAllClaudeNum}{0}
\rqSet{rqFFNotReexLibAllClaudeDen}{11}
\rqSet{rqFFNotReexLibAllClaudePct}{0.0}
\rqSet{rqFFNotReexLibAllCodexNum}{0}
\rqSet{rqFFNotReexLibAllCodexDen}{45}
\rqSet{rqFFNotReexLibAllCodexPct}{0.0}
\rqSet{rqFFNotReexLibAllCTwoRustNum}{0}
\rqSet{rqFFNotReexLibAllCTwoRustDen}{0}
\rqSet{rqFFNotReexLibAllCTwoRustPct}{n/a}
\rqSet{rqFFNotReexLibAllLaertesNum}{0}
\rqSet{rqFFNotReexLibAllLaertesDen}{0}
\rqSet{rqFFNotReexLibAllLaertesPct}{n/a}
\rqSet{rqFFNotReexLibAllCTwoSaferRustNum}{0}
\rqSet{rqFFNotReexLibAllCTwoSaferRustDen}{23}
\rqSet{rqFFNotReexLibAllCTwoSaferRustPct}{0.0}
\rqSet{rqFFNotReexLibAllSmartCTwoRustNum}{0}
\rqSet{rqFFNotReexLibAllSmartCTwoRustDen}{0}
\rqSet{rqFFNotReexLibAllSmartCTwoRustPct}{n/a}
\rqSet{rqFFNotReexLibAllKimiNum}{0}
\rqSet{rqFFNotReexLibAllKimiDen}{15}
\rqSet{rqFFNotReexLibAllKimiPct}{0.0}
\rqSet{rqFFNotReexLibAllGptFiveNum}{0}
\rqSet{rqFFNotReexLibAllGptFiveDen}{13}
\rqSet{rqFFNotReexLibAllGptFivePct}{0.0}
\rqSet{rqFFNotReexLibAllGeminiNum}{0}
\rqSet{rqFFNotReexLibAllGeminiDen}{4}
\rqSet{rqFFNotReexLibAllGeminiPct}{0.0}
\rqSet{rqFFNotReexLibBOneNum}{0}
\rqSet{rqFFNotReexLibBOneDen}{31}
\rqSet{rqFFNotReexLibBOnePct}{0.0}
\rqSet{rqFFNotReexLibBTwoNum}{0}
\rqSet{rqFFNotReexLibBTwoDen}{24}
\rqSet{rqFFNotReexLibBTwoPct}{0.0}
\rqSet{rqFFNotReexLibPZeroNum}{0}
\rqSet{rqFFNotReexLibPZeroDen}{0}
\rqSet{rqFFNotReexLibPZeroPct}{n/a}
\rqSet{rqFFNotReexLibPOneNum}{0}
\rqSet{rqFFNotReexLibPOneDen}{56}
\rqSet{rqFFNotReexLibPOnePct}{0.0}
\rqSet{rqFFHarnessErrorLibAllNum}{0}
\rqSet{rqFFHarnessErrorLibAllDen}{111}
\rqSet{rqFFHarnessErrorLibAllPct}{0.0}
\rqSet{rqFFHarnessErrorLibAllKiroNum}{0}
\rqSet{rqFFHarnessErrorLibAllKiroDen}{0}
\rqSet{rqFFHarnessErrorLibAllKiroPct}{n/a}
\rqSet{rqFFHarnessErrorLibAllClaudeNum}{0}
\rqSet{rqFFHarnessErrorLibAllClaudeDen}{11}
\rqSet{rqFFHarnessErrorLibAllClaudePct}{0.0}
\rqSet{rqFFHarnessErrorLibAllCodexNum}{0}
\rqSet{rqFFHarnessErrorLibAllCodexDen}{45}
\rqSet{rqFFHarnessErrorLibAllCodexPct}{0.0}
\rqSet{rqFFHarnessErrorLibAllCTwoRustNum}{0}
\rqSet{rqFFHarnessErrorLibAllCTwoRustDen}{0}
\rqSet{rqFFHarnessErrorLibAllCTwoRustPct}{n/a}
\rqSet{rqFFHarnessErrorLibAllLaertesNum}{0}
\rqSet{rqFFHarnessErrorLibAllLaertesDen}{0}
\rqSet{rqFFHarnessErrorLibAllLaertesPct}{n/a}
\rqSet{rqFFHarnessErrorLibAllCTwoSaferRustNum}{0}
\rqSet{rqFFHarnessErrorLibAllCTwoSaferRustDen}{23}
\rqSet{rqFFHarnessErrorLibAllCTwoSaferRustPct}{0.0}
\rqSet{rqFFHarnessErrorLibAllSmartCTwoRustNum}{0}
\rqSet{rqFFHarnessErrorLibAllSmartCTwoRustDen}{0}
\rqSet{rqFFHarnessErrorLibAllSmartCTwoRustPct}{n/a}
\rqSet{rqFFHarnessErrorLibAllKimiNum}{0}
\rqSet{rqFFHarnessErrorLibAllKimiDen}{15}
\rqSet{rqFFHarnessErrorLibAllKimiPct}{0.0}
\rqSet{rqFFHarnessErrorLibAllGptFiveNum}{0}
\rqSet{rqFFHarnessErrorLibAllGptFiveDen}{13}
\rqSet{rqFFHarnessErrorLibAllGptFivePct}{0.0}
\rqSet{rqFFHarnessErrorLibAllGeminiNum}{0}
\rqSet{rqFFHarnessErrorLibAllGeminiDen}{4}
\rqSet{rqFFHarnessErrorLibAllGeminiPct}{0.0}
\rqSet{rqFFHarnessErrorLibBOneNum}{0}
\rqSet{rqFFHarnessErrorLibBOneDen}{31}
\rqSet{rqFFHarnessErrorLibBOnePct}{0.0}
\rqSet{rqFFHarnessErrorLibBTwoNum}{0}
\rqSet{rqFFHarnessErrorLibBTwoDen}{24}
\rqSet{rqFFHarnessErrorLibBTwoPct}{0.0}
\rqSet{rqFFHarnessErrorLibPZeroNum}{0}
\rqSet{rqFFHarnessErrorLibPZeroDen}{0}
\rqSet{rqFFHarnessErrorLibPZeroPct}{n/a}
\rqSet{rqFFHarnessErrorLibPOneNum}{0}
\rqSet{rqFFHarnessErrorLibPOneDen}{56}
\rqSet{rqFFHarnessErrorLibPOnePct}{0.0}
\rqSet{rqFFRrunRecLibAllNum}{111}
\rqSet{rqFFRrunRecLibAllDen}{111}
\rqSet{rqFFRrunRecLibAllPct}{100.0}
\rqSet{rqFFRrunRecLibAllKiroNum}{0}
\rqSet{rqFFRrunRecLibAllKiroDen}{0}
\rqSet{rqFFRrunRecLibAllKiroPct}{n/a}
\rqSet{rqFFRrunRecLibAllClaudeNum}{11}
\rqSet{rqFFRrunRecLibAllClaudeDen}{11}
\rqSet{rqFFRrunRecLibAllClaudePct}{100.0}
\rqSet{rqFFRrunRecLibAllCodexNum}{45}
\rqSet{rqFFRrunRecLibAllCodexDen}{45}
\rqSet{rqFFRrunRecLibAllCodexPct}{100.0}
\rqSet{rqFFRrunRecLibAllCTwoRustNum}{0}
\rqSet{rqFFRrunRecLibAllCTwoRustDen}{0}
\rqSet{rqFFRrunRecLibAllCTwoRustPct}{n/a}
\rqSet{rqFFRrunRecLibAllLaertesNum}{0}
\rqSet{rqFFRrunRecLibAllLaertesDen}{0}
\rqSet{rqFFRrunRecLibAllLaertesPct}{n/a}
\rqSet{rqFFRrunRecLibAllCTwoSaferRustNum}{23}
\rqSet{rqFFRrunRecLibAllCTwoSaferRustDen}{23}
\rqSet{rqFFRrunRecLibAllCTwoSaferRustPct}{100.0}
\rqSet{rqFFRrunRecLibAllSmartCTwoRustNum}{0}
\rqSet{rqFFRrunRecLibAllSmartCTwoRustDen}{0}
\rqSet{rqFFRrunRecLibAllSmartCTwoRustPct}{n/a}
\rqSet{rqFFRrunRecLibAllKimiNum}{15}
\rqSet{rqFFRrunRecLibAllKimiDen}{15}
\rqSet{rqFFRrunRecLibAllKimiPct}{100.0}
\rqSet{rqFFRrunRecLibAllGptFiveNum}{13}
\rqSet{rqFFRrunRecLibAllGptFiveDen}{13}
\rqSet{rqFFRrunRecLibAllGptFivePct}{100.0}
\rqSet{rqFFRrunRecLibAllGeminiNum}{4}
\rqSet{rqFFRrunRecLibAllGeminiDen}{4}
\rqSet{rqFFRrunRecLibAllGeminiPct}{100.0}
\rqSet{rqFFRrunRecLibBOneNum}{31}
\rqSet{rqFFRrunRecLibBOneDen}{31}
\rqSet{rqFFRrunRecLibBOnePct}{100.0}
\rqSet{rqFFRrunRecLibBTwoNum}{24}
\rqSet{rqFFRrunRecLibBTwoDen}{24}
\rqSet{rqFFRrunRecLibBTwoPct}{100.0}
\rqSet{rqFFRrunRecLibPZeroNum}{0}
\rqSet{rqFFRrunRecLibPZeroDen}{0}
\rqSet{rqFFRrunRecLibPZeroPct}{n/a}
\rqSet{rqFFRrunRecLibPOneNum}{56}
\rqSet{rqFFRrunRecLibPOneDen}{56}
\rqSet{rqFFRrunRecLibPOnePct}{100.0}
\rqSet{rqRZeroTestRRunLibAllNum}{111}
\rqSet{rqRZeroTestRRunLibAllDen}{138}
\rqSet{rqRZeroTestRRunLibAllPct}{80.4}
\rqSet{rqRZeroTestRRunLibAllKiroNum}{0}
\rqSet{rqRZeroTestRRunLibAllKiroDen}{2}
\rqSet{rqRZeroTestRRunLibAllKiroPct}{0.0}
\rqSet{rqRZeroTestRRunLibAllClaudeNum}{11}
\rqSet{rqRZeroTestRRunLibAllClaudeDen}{13}
\rqSet{rqRZeroTestRRunLibAllClaudePct}{84.6}
\rqSet{rqRZeroTestRRunLibAllCodexNum}{45}
\rqSet{rqRZeroTestRRunLibAllCodexDen}{47}
\rqSet{rqRZeroTestRRunLibAllCodexPct}{95.7}
\rqSet{rqRZeroTestRRunLibAllCTwoRustNum}{0}
\rqSet{rqRZeroTestRRunLibAllCTwoRustDen}{1}
\rqSet{rqRZeroTestRRunLibAllCTwoRustPct}{0.0}
\rqSet{rqRZeroTestRRunLibAllLaertesNum}{0}
\rqSet{rqRZeroTestRRunLibAllLaertesDen}{1}
\rqSet{rqRZeroTestRRunLibAllLaertesPct}{0.0}
\rqSet{rqRZeroTestRRunLibAllCTwoSaferRustNum}{23}
\rqSet{rqRZeroTestRRunLibAllCTwoSaferRustDen}{25}
\rqSet{rqRZeroTestRRunLibAllCTwoSaferRustPct}{92.0}
\rqSet{rqRZeroTestRRunLibAllSmartCTwoRustNum}{0}
\rqSet{rqRZeroTestRRunLibAllSmartCTwoRustDen}{0}
\rqSet{rqRZeroTestRRunLibAllSmartCTwoRustPct}{n/a}
\rqSet{rqRZeroTestRRunLibAllKimiNum}{15}
\rqSet{rqRZeroTestRRunLibAllKimiDen}{21}
\rqSet{rqRZeroTestRRunLibAllKimiPct}{71.4}
\rqSet{rqRZeroTestRRunLibAllGptFiveNum}{13}
\rqSet{rqRZeroTestRRunLibAllGptFiveDen}{15}
\rqSet{rqRZeroTestRRunLibAllGptFivePct}{86.7}
\rqSet{rqRZeroTestRRunLibAllGeminiNum}{4}
\rqSet{rqRZeroTestRRunLibAllGeminiDen}{13}
\rqSet{rqRZeroTestRRunLibAllGeminiPct}{30.8}
\rqSet{rqRZeroTestRRunLibBOneNum}{31}
\rqSet{rqRZeroTestRRunLibBOneDen}{38}
\rqSet{rqRZeroTestRRunLibBOnePct}{81.6}
\rqSet{rqRZeroTestRRunLibBTwoNum}{24}
\rqSet{rqRZeroTestRRunLibBTwoDen}{44}
\rqSet{rqRZeroTestRRunLibBTwoPct}{54.5}
\rqSet{rqRZeroTestRRunLibPZeroNum}{0}
\rqSet{rqRZeroTestRRunLibPZeroDen}{0}
\rqSet{rqRZeroTestRRunLibPZeroPct}{n/a}
\rqSet{rqRZeroTestRRunLibPOneNum}{56}
\rqSet{rqRZeroTestRRunLibPOneDen}{56}
\rqSet{rqRZeroTestRRunLibPOnePct}{100.0}
\rqSet{rqRZeroTestRNonrepLibAllNum}{15}
\rqSet{rqRZeroTestRNonrepLibAllDen}{138}
\rqSet{rqRZeroTestRNonrepLibAllPct}{10.9}
\rqSet{rqRZeroTestRNonrepLibAllKiroNum}{2}
\rqSet{rqRZeroTestRNonrepLibAllKiroDen}{2}
\rqSet{rqRZeroTestRNonrepLibAllKiroPct}{100.0}
\rqSet{rqRZeroTestRNonrepLibAllClaudeNum}{2}
\rqSet{rqRZeroTestRNonrepLibAllClaudeDen}{13}
\rqSet{rqRZeroTestRNonrepLibAllClaudePct}{15.4}
\rqSet{rqRZeroTestRNonrepLibAllCodexNum}{1}
\rqSet{rqRZeroTestRNonrepLibAllCodexDen}{47}
\rqSet{rqRZeroTestRNonrepLibAllCodexPct}{2.1}
\rqSet{rqRZeroTestRNonrepLibAllCTwoRustNum}{1}
\rqSet{rqRZeroTestRNonrepLibAllCTwoRustDen}{1}
\rqSet{rqRZeroTestRNonrepLibAllCTwoRustPct}{100.0}
\rqSet{rqRZeroTestRNonrepLibAllLaertesNum}{1}
\rqSet{rqRZeroTestRNonrepLibAllLaertesDen}{1}
\rqSet{rqRZeroTestRNonrepLibAllLaertesPct}{100.0}
\rqSet{rqRZeroTestRNonrepLibAllCTwoSaferRustNum}{2}
\rqSet{rqRZeroTestRNonrepLibAllCTwoSaferRustDen}{25}
\rqSet{rqRZeroTestRNonrepLibAllCTwoSaferRustPct}{8.0}
\rqSet{rqRZeroTestRNonrepLibAllSmartCTwoRustNum}{0}
\rqSet{rqRZeroTestRNonrepLibAllSmartCTwoRustDen}{0}
\rqSet{rqRZeroTestRNonrepLibAllSmartCTwoRustPct}{n/a}
\rqSet{rqRZeroTestRNonrepLibAllKimiNum}{2}
\rqSet{rqRZeroTestRNonrepLibAllKimiDen}{21}
\rqSet{rqRZeroTestRNonrepLibAllKimiPct}{9.5}
\rqSet{rqRZeroTestRNonrepLibAllGptFiveNum}{2}
\rqSet{rqRZeroTestRNonrepLibAllGptFiveDen}{15}
\rqSet{rqRZeroTestRNonrepLibAllGptFivePct}{13.3}
\rqSet{rqRZeroTestRNonrepLibAllGeminiNum}{2}
\rqSet{rqRZeroTestRNonrepLibAllGeminiDen}{13}
\rqSet{rqRZeroTestRNonrepLibAllGeminiPct}{15.4}
\rqSet{rqRZeroTestRNonrepLibBOneNum}{5}
\rqSet{rqRZeroTestRNonrepLibBOneDen}{38}
\rqSet{rqRZeroTestRNonrepLibBOnePct}{13.2}
\rqSet{rqRZeroTestRNonrepLibBTwoNum}{10}
\rqSet{rqRZeroTestRNonrepLibBTwoDen}{44}
\rqSet{rqRZeroTestRNonrepLibBTwoPct}{22.7}
\rqSet{rqRZeroTestRNonrepLibPZeroNum}{0}
\rqSet{rqRZeroTestRNonrepLibPZeroDen}{0}
\rqSet{rqRZeroTestRNonrepLibPZeroPct}{n/a}
\rqSet{rqRZeroTestRNonrepLibPOneNum}{0}
\rqSet{rqRZeroTestRNonrepLibPOneDen}{56}
\rqSet{rqRZeroTestRNonrepLibPOnePct}{0.0}
\rqSet{rqRZeroTestRNobuildLibAllNum}{12}
\rqSet{rqRZeroTestRNobuildLibAllDen}{138}
\rqSet{rqRZeroTestRNobuildLibAllPct}{8.7}
\rqSet{rqRZeroTestRNobuildLibAllKiroNum}{0}
\rqSet{rqRZeroTestRNobuildLibAllKiroDen}{2}
\rqSet{rqRZeroTestRNobuildLibAllKiroPct}{0.0}
\rqSet{rqRZeroTestRNobuildLibAllClaudeNum}{0}
\rqSet{rqRZeroTestRNobuildLibAllClaudeDen}{13}
\rqSet{rqRZeroTestRNobuildLibAllClaudePct}{0.0}
\rqSet{rqRZeroTestRNobuildLibAllCodexNum}{1}
\rqSet{rqRZeroTestRNobuildLibAllCodexDen}{47}
\rqSet{rqRZeroTestRNobuildLibAllCodexPct}{2.1}
\rqSet{rqRZeroTestRNobuildLibAllCTwoRustNum}{0}
\rqSet{rqRZeroTestRNobuildLibAllCTwoRustDen}{1}
\rqSet{rqRZeroTestRNobuildLibAllCTwoRustPct}{0.0}
\rqSet{rqRZeroTestRNobuildLibAllLaertesNum}{0}
\rqSet{rqRZeroTestRNobuildLibAllLaertesDen}{1}
\rqSet{rqRZeroTestRNobuildLibAllLaertesPct}{0.0}
\rqSet{rqRZeroTestRNobuildLibAllCTwoSaferRustNum}{0}
\rqSet{rqRZeroTestRNobuildLibAllCTwoSaferRustDen}{25}
\rqSet{rqRZeroTestRNobuildLibAllCTwoSaferRustPct}{0.0}
\rqSet{rqRZeroTestRNobuildLibAllSmartCTwoRustNum}{0}
\rqSet{rqRZeroTestRNobuildLibAllSmartCTwoRustDen}{0}
\rqSet{rqRZeroTestRNobuildLibAllSmartCTwoRustPct}{n/a}
\rqSet{rqRZeroTestRNobuildLibAllKimiNum}{4}
\rqSet{rqRZeroTestRNobuildLibAllKimiDen}{21}
\rqSet{rqRZeroTestRNobuildLibAllKimiPct}{19.0}
\rqSet{rqRZeroTestRNobuildLibAllGptFiveNum}{0}
\rqSet{rqRZeroTestRNobuildLibAllGptFiveDen}{15}
\rqSet{rqRZeroTestRNobuildLibAllGptFivePct}{0.0}
\rqSet{rqRZeroTestRNobuildLibAllGeminiNum}{7}
\rqSet{rqRZeroTestRNobuildLibAllGeminiDen}{13}
\rqSet{rqRZeroTestRNobuildLibAllGeminiPct}{53.8}
\rqSet{rqRZeroTestRNobuildLibBOneNum}{2}
\rqSet{rqRZeroTestRNobuildLibBOneDen}{38}
\rqSet{rqRZeroTestRNobuildLibBOnePct}{5.3}
\rqSet{rqRZeroTestRNobuildLibBTwoNum}{10}
\rqSet{rqRZeroTestRNobuildLibBTwoDen}{44}
\rqSet{rqRZeroTestRNobuildLibBTwoPct}{22.7}
\rqSet{rqRZeroTestRNobuildLibPZeroNum}{0}
\rqSet{rqRZeroTestRNobuildLibPZeroDen}{0}
\rqSet{rqRZeroTestRNobuildLibPZeroPct}{n/a}
\rqSet{rqRZeroTestRNobuildLibPOneNum}{0}
\rqSet{rqRZeroTestRNobuildLibPOneDen}{56}
\rqSet{rqRZeroTestRNobuildLibPOnePct}{0.0}
\rqSet{rqRZeroTestROtherLibAllNum}{0}
\rqSet{rqRZeroTestROtherLibAllDen}{138}
\rqSet{rqRZeroTestROtherLibAllPct}{0.0}
\rqSet{rqRZeroTestROtherLibAllKiroNum}{0}
\rqSet{rqRZeroTestROtherLibAllKiroDen}{2}
\rqSet{rqRZeroTestROtherLibAllKiroPct}{0.0}
\rqSet{rqRZeroTestROtherLibAllClaudeNum}{0}
\rqSet{rqRZeroTestROtherLibAllClaudeDen}{13}
\rqSet{rqRZeroTestROtherLibAllClaudePct}{0.0}
\rqSet{rqRZeroTestROtherLibAllCodexNum}{0}
\rqSet{rqRZeroTestROtherLibAllCodexDen}{47}
\rqSet{rqRZeroTestROtherLibAllCodexPct}{0.0}
\rqSet{rqRZeroTestROtherLibAllCTwoRustNum}{0}
\rqSet{rqRZeroTestROtherLibAllCTwoRustDen}{1}
\rqSet{rqRZeroTestROtherLibAllCTwoRustPct}{0.0}
\rqSet{rqRZeroTestROtherLibAllLaertesNum}{0}
\rqSet{rqRZeroTestROtherLibAllLaertesDen}{1}
\rqSet{rqRZeroTestROtherLibAllLaertesPct}{0.0}
\rqSet{rqRZeroTestROtherLibAllCTwoSaferRustNum}{0}
\rqSet{rqRZeroTestROtherLibAllCTwoSaferRustDen}{25}
\rqSet{rqRZeroTestROtherLibAllCTwoSaferRustPct}{0.0}
\rqSet{rqRZeroTestROtherLibAllSmartCTwoRustNum}{0}
\rqSet{rqRZeroTestROtherLibAllSmartCTwoRustDen}{0}
\rqSet{rqRZeroTestROtherLibAllSmartCTwoRustPct}{n/a}
\rqSet{rqRZeroTestROtherLibAllKimiNum}{0}
\rqSet{rqRZeroTestROtherLibAllKimiDen}{21}
\rqSet{rqRZeroTestROtherLibAllKimiPct}{0.0}
\rqSet{rqRZeroTestROtherLibAllGptFiveNum}{0}
\rqSet{rqRZeroTestROtherLibAllGptFiveDen}{15}
\rqSet{rqRZeroTestROtherLibAllGptFivePct}{0.0}
\rqSet{rqRZeroTestROtherLibAllGeminiNum}{0}
\rqSet{rqRZeroTestROtherLibAllGeminiDen}{13}
\rqSet{rqRZeroTestROtherLibAllGeminiPct}{0.0}
\rqSet{rqRZeroTestROtherLibBOneNum}{0}
\rqSet{rqRZeroTestROtherLibBOneDen}{38}
\rqSet{rqRZeroTestROtherLibBOnePct}{0.0}
\rqSet{rqRZeroTestROtherLibBTwoNum}{0}
\rqSet{rqRZeroTestROtherLibBTwoDen}{44}
\rqSet{rqRZeroTestROtherLibBTwoPct}{0.0}
\rqSet{rqRZeroTestROtherLibPZeroNum}{0}
\rqSet{rqRZeroTestROtherLibPZeroDen}{0}
\rqSet{rqRZeroTestROtherLibPZeroPct}{n/a}
\rqSet{rqRZeroTestROtherLibPOneNum}{0}
\rqSet{rqRZeroTestROtherLibPOneDen}{56}
\rqSet{rqRZeroTestROtherLibPOnePct}{0.0}
\rqSet{rqLrLibAllNum}{27}
\rqSet{rqLrLibAllDen}{138}
\rqSet{rqLrLibAllPct}{19.6}
\rqSet{rqLrLibAllKiroNum}{2}
\rqSet{rqLrLibAllKiroDen}{2}
\rqSet{rqLrLibAllKiroPct}{100.0}
\rqSet{rqLrLibAllClaudeNum}{2}
\rqSet{rqLrLibAllClaudeDen}{13}
\rqSet{rqLrLibAllClaudePct}{15.4}
\rqSet{rqLrLibAllCodexNum}{2}
\rqSet{rqLrLibAllCodexDen}{47}
\rqSet{rqLrLibAllCodexPct}{4.3}
\rqSet{rqLrLibAllCTwoRustNum}{1}
\rqSet{rqLrLibAllCTwoRustDen}{1}
\rqSet{rqLrLibAllCTwoRustPct}{100.0}
\rqSet{rqLrLibAllLaertesNum}{1}
\rqSet{rqLrLibAllLaertesDen}{1}
\rqSet{rqLrLibAllLaertesPct}{100.0}
\rqSet{rqLrLibAllCTwoSaferRustNum}{2}
\rqSet{rqLrLibAllCTwoSaferRustDen}{25}
\rqSet{rqLrLibAllCTwoSaferRustPct}{8.0}
\rqSet{rqLrLibAllSmartCTwoRustNum}{0}
\rqSet{rqLrLibAllSmartCTwoRustDen}{0}
\rqSet{rqLrLibAllSmartCTwoRustPct}{n/a}
\rqSet{rqLrLibAllKimiNum}{6}
\rqSet{rqLrLibAllKimiDen}{21}
\rqSet{rqLrLibAllKimiPct}{28.6}
\rqSet{rqLrLibAllGptFiveNum}{2}
\rqSet{rqLrLibAllGptFiveDen}{15}
\rqSet{rqLrLibAllGptFivePct}{13.3}
\rqSet{rqLrLibAllGeminiNum}{9}
\rqSet{rqLrLibAllGeminiDen}{13}
\rqSet{rqLrLibAllGeminiPct}{69.2}
\rqSet{rqLrLibBOneNum}{7}
\rqSet{rqLrLibBOneDen}{38}
\rqSet{rqLrLibBOnePct}{18.4}
\rqSet{rqLrLibBOneKiroNum}{0}
\rqSet{rqLrLibBOneKiroDen}{0}
\rqSet{rqLrLibBOneKiroPct}{n/a}
\rqSet{rqLrLibBOneClaudeNum}{0}
\rqSet{rqLrLibBOneClaudeDen}{0}
\rqSet{rqLrLibBOneClaudePct}{n/a}
\rqSet{rqLrLibBOneCodexNum}{0}
\rqSet{rqLrLibBOneCodexDen}{0}
\rqSet{rqLrLibBOneCodexPct}{n/a}
\rqSet{rqLrLibBOneCTwoRustNum}{0}
\rqSet{rqLrLibBOneCTwoRustDen}{0}
\rqSet{rqLrLibBOneCTwoRustPct}{n/a}
\rqSet{rqLrLibBOneLaertesNum}{0}
\rqSet{rqLrLibBOneLaertesDen}{0}
\rqSet{rqLrLibBOneLaertesPct}{n/a}
\rqSet{rqLrLibBOneCTwoSaferRustNum}{1}
\rqSet{rqLrLibBOneCTwoSaferRustDen}{17}
\rqSet{rqLrLibBOneCTwoSaferRustPct}{5.9}
\rqSet{rqLrLibBOneSmartCTwoRustNum}{0}
\rqSet{rqLrLibBOneSmartCTwoRustDen}{0}
\rqSet{rqLrLibBOneSmartCTwoRustPct}{n/a}
\rqSet{rqLrLibBOneKimiNum}{2}
\rqSet{rqLrLibBOneKimiDen}{11}
\rqSet{rqLrLibBOneKimiPct}{18.2}
\rqSet{rqLrLibBOneGptFiveNum}{2}
\rqSet{rqLrLibBOneGptFiveDen}{7}
\rqSet{rqLrLibBOneGptFivePct}{28.6}
\rqSet{rqLrLibBOneGeminiNum}{2}
\rqSet{rqLrLibBOneGeminiDen}{3}
\rqSet{rqLrLibBOneGeminiPct}{66.7}
\rqSet{rqLrLibBTwoNum}{20}
\rqSet{rqLrLibBTwoDen}{44}
\rqSet{rqLrLibBTwoPct}{45.5}
\rqSet{rqLrLibBTwoKiroNum}{2}
\rqSet{rqLrLibBTwoKiroDen}{2}
\rqSet{rqLrLibBTwoKiroPct}{100.0}
\rqSet{rqLrLibBTwoClaudeNum}{2}
\rqSet{rqLrLibBTwoClaudeDen}{2}
\rqSet{rqLrLibBTwoClaudePct}{100.0}
\rqSet{rqLrLibBTwoCodexNum}{2}
\rqSet{rqLrLibBTwoCodexDen}{2}
\rqSet{rqLrLibBTwoCodexPct}{100.0}
\rqSet{rqLrLibBTwoCTwoRustNum}{1}
\rqSet{rqLrLibBTwoCTwoRustDen}{1}
\rqSet{rqLrLibBTwoCTwoRustPct}{100.0}
\rqSet{rqLrLibBTwoLaertesNum}{1}
\rqSet{rqLrLibBTwoLaertesDen}{1}
\rqSet{rqLrLibBTwoLaertesPct}{100.0}
\rqSet{rqLrLibBTwoCTwoSaferRustNum}{1}
\rqSet{rqLrLibBTwoCTwoSaferRustDen}{8}
\rqSet{rqLrLibBTwoCTwoSaferRustPct}{12.5}
\rqSet{rqLrLibBTwoSmartCTwoRustNum}{0}
\rqSet{rqLrLibBTwoSmartCTwoRustDen}{0}
\rqSet{rqLrLibBTwoSmartCTwoRustPct}{n/a}
\rqSet{rqLrLibBTwoKimiNum}{4}
\rqSet{rqLrLibBTwoKimiDen}{10}
\rqSet{rqLrLibBTwoKimiPct}{40.0}
\rqSet{rqLrLibBTwoGptFiveNum}{0}
\rqSet{rqLrLibBTwoGptFiveDen}{8}
\rqSet{rqLrLibBTwoGptFivePct}{0.0}
\rqSet{rqLrLibBTwoGeminiNum}{7}
\rqSet{rqLrLibBTwoGeminiDen}{10}
\rqSet{rqLrLibBTwoGeminiPct}{70.0}
\rqSet{rqLrLibPZeroNum}{0}
\rqSet{rqLrLibPZeroDen}{0}
\rqSet{rqLrLibPZeroPct}{n/a}
\rqSet{rqLrLibPZeroKiroNum}{0}
\rqSet{rqLrLibPZeroKiroDen}{0}
\rqSet{rqLrLibPZeroKiroPct}{n/a}
\rqSet{rqLrLibPZeroClaudeNum}{0}
\rqSet{rqLrLibPZeroClaudeDen}{0}
\rqSet{rqLrLibPZeroClaudePct}{n/a}
\rqSet{rqLrLibPZeroCodexNum}{0}
\rqSet{rqLrLibPZeroCodexDen}{0}
\rqSet{rqLrLibPZeroCodexPct}{n/a}
\rqSet{rqLrLibPZeroCTwoRustNum}{0}
\rqSet{rqLrLibPZeroCTwoRustDen}{0}
\rqSet{rqLrLibPZeroCTwoRustPct}{n/a}
\rqSet{rqLrLibPZeroLaertesNum}{0}
\rqSet{rqLrLibPZeroLaertesDen}{0}
\rqSet{rqLrLibPZeroLaertesPct}{n/a}
\rqSet{rqLrLibPZeroCTwoSaferRustNum}{0}
\rqSet{rqLrLibPZeroCTwoSaferRustDen}{0}
\rqSet{rqLrLibPZeroCTwoSaferRustPct}{n/a}
\rqSet{rqLrLibPZeroSmartCTwoRustNum}{0}
\rqSet{rqLrLibPZeroSmartCTwoRustDen}{0}
\rqSet{rqLrLibPZeroSmartCTwoRustPct}{n/a}
\rqSet{rqLrLibPZeroKimiNum}{0}
\rqSet{rqLrLibPZeroKimiDen}{0}
\rqSet{rqLrLibPZeroKimiPct}{n/a}
\rqSet{rqLrLibPZeroGptFiveNum}{0}
\rqSet{rqLrLibPZeroGptFiveDen}{0}
\rqSet{rqLrLibPZeroGptFivePct}{n/a}
\rqSet{rqLrLibPZeroGeminiNum}{0}
\rqSet{rqLrLibPZeroGeminiDen}{0}
\rqSet{rqLrLibPZeroGeminiPct}{n/a}
\rqSet{rqLrLibPOneNum}{0}
\rqSet{rqLrLibPOneDen}{56}
\rqSet{rqLrLibPOnePct}{0.0}
\rqSet{rqLrLibPOneKiroNum}{0}
\rqSet{rqLrLibPOneKiroDen}{0}
\rqSet{rqLrLibPOneKiroPct}{n/a}
\rqSet{rqLrLibPOneClaudeNum}{0}
\rqSet{rqLrLibPOneClaudeDen}{11}
\rqSet{rqLrLibPOneClaudePct}{0.0}
\rqSet{rqLrLibPOneCodexNum}{0}
\rqSet{rqLrLibPOneCodexDen}{45}
\rqSet{rqLrLibPOneCodexPct}{0.0}
\rqSet{rqLrLibPOneCTwoRustNum}{0}
\rqSet{rqLrLibPOneCTwoRustDen}{0}
\rqSet{rqLrLibPOneCTwoRustPct}{n/a}
\rqSet{rqLrLibPOneLaertesNum}{0}
\rqSet{rqLrLibPOneLaertesDen}{0}
\rqSet{rqLrLibPOneLaertesPct}{n/a}
\rqSet{rqLrLibPOneCTwoSaferRustNum}{0}
\rqSet{rqLrLibPOneCTwoSaferRustDen}{0}
\rqSet{rqLrLibPOneCTwoSaferRustPct}{n/a}
\rqSet{rqLrLibPOneSmartCTwoRustNum}{0}
\rqSet{rqLrLibPOneSmartCTwoRustDen}{0}
\rqSet{rqLrLibPOneSmartCTwoRustPct}{n/a}
\rqSet{rqLrLibPOneKimiNum}{0}
\rqSet{rqLrLibPOneKimiDen}{0}
\rqSet{rqLrLibPOneKimiPct}{n/a}
\rqSet{rqLrLibPOneGptFiveNum}{0}
\rqSet{rqLrLibPOneGptFiveDen}{0}
\rqSet{rqLrLibPOneGptFivePct}{n/a}
\rqSet{rqLrLibPOneGeminiNum}{0}
\rqSet{rqLrLibPOneGeminiDen}{0}
\rqSet{rqLrLibPOneGeminiPct}{n/a}
\rqSet{rqFFReasonNotAdjudicableUbInCLibAllNum}{0}
\rqSet{rqFFReasonNotAdjudicableUbInCLibAllDen}{0}
\rqSet{rqFFReasonNotAdjudicableUbInCLibAllPct}{n/a}
\rqSet{rqFFReasonNotAdjudicableUbInCLibBOneNum}{0}
\rqSet{rqFFReasonNotAdjudicableUbInCLibBOneDen}{0}
\rqSet{rqFFReasonNotAdjudicableUbInCLibBOnePct}{n/a}
\rqSet{rqFFReasonNotAdjudicableUbInCLibBTwoNum}{0}
\rqSet{rqFFReasonNotAdjudicableUbInCLibBTwoDen}{0}
\rqSet{rqFFReasonNotAdjudicableUbInCLibBTwoPct}{n/a}
\rqSet{rqFFReasonNotAdjudicableUbInCLibPZeroNum}{0}
\rqSet{rqFFReasonNotAdjudicableUbInCLibPZeroDen}{0}
\rqSet{rqFFReasonNotAdjudicableUbInCLibPZeroPct}{n/a}
\rqSet{rqFFReasonNotAdjudicableUbInCLibPOneNum}{0}
\rqSet{rqFFReasonNotAdjudicableUbInCLibPOneDen}{0}
\rqSet{rqFFReasonNotAdjudicableUbInCLibPOnePct}{n/a}
\rqSet{rqFFReasonNotAdjudicableBothBudgetLibAllNum}{0}
\rqSet{rqFFReasonNotAdjudicableBothBudgetLibAllDen}{0}
\rqSet{rqFFReasonNotAdjudicableBothBudgetLibAllPct}{n/a}
\rqSet{rqFFReasonNotAdjudicableBothBudgetLibBOneNum}{0}
\rqSet{rqFFReasonNotAdjudicableBothBudgetLibBOneDen}{0}
\rqSet{rqFFReasonNotAdjudicableBothBudgetLibBOnePct}{n/a}
\rqSet{rqFFReasonNotAdjudicableBothBudgetLibBTwoNum}{0}
\rqSet{rqFFReasonNotAdjudicableBothBudgetLibBTwoDen}{0}
\rqSet{rqFFReasonNotAdjudicableBothBudgetLibBTwoPct}{n/a}
\rqSet{rqFFReasonNotAdjudicableBothBudgetLibPZeroNum}{0}
\rqSet{rqFFReasonNotAdjudicableBothBudgetLibPZeroDen}{0}
\rqSet{rqFFReasonNotAdjudicableBothBudgetLibPZeroPct}{n/a}
\rqSet{rqFFReasonNotAdjudicableBothBudgetLibPOneNum}{0}
\rqSet{rqFFReasonNotAdjudicableBothBudgetLibPOneDen}{0}
\rqSet{rqFFReasonNotAdjudicableBothBudgetLibPOnePct}{n/a}
\rqSet{rqFFReasonNotAdjudicableSanitizerInconclusiveLibAllNum}{0}
\rqSet{rqFFReasonNotAdjudicableSanitizerInconclusiveLibAllDen}{0}
\rqSet{rqFFReasonNotAdjudicableSanitizerInconclusiveLibAllPct}{n/a}
\rqSet{rqFFReasonNotAdjudicableSanitizerInconclusiveLibBOneNum}{0}
\rqSet{rqFFReasonNotAdjudicableSanitizerInconclusiveLibBOneDen}{0}
\rqSet{rqFFReasonNotAdjudicableSanitizerInconclusiveLibBOnePct}{n/a}
\rqSet{rqFFReasonNotAdjudicableSanitizerInconclusiveLibBTwoNum}{0}
\rqSet{rqFFReasonNotAdjudicableSanitizerInconclusiveLibBTwoDen}{0}
\rqSet{rqFFReasonNotAdjudicableSanitizerInconclusiveLibBTwoPct}{n/a}
\rqSet{rqFFReasonNotAdjudicableSanitizerInconclusiveLibPZeroNum}{0}
\rqSet{rqFFReasonNotAdjudicableSanitizerInconclusiveLibPZeroDen}{0}
\rqSet{rqFFReasonNotAdjudicableSanitizerInconclusiveLibPZeroPct}{n/a}
\rqSet{rqFFReasonNotAdjudicableSanitizerInconclusiveLibPOneNum}{0}
\rqSet{rqFFReasonNotAdjudicableSanitizerInconclusiveLibPOneDen}{0}
\rqSet{rqFFReasonNotAdjudicableSanitizerInconclusiveLibPOnePct}{n/a}
\rqSet{rqFFReasonNotAdjudicableCUnstableLibAllNum}{0}
\rqSet{rqFFReasonNotAdjudicableCUnstableLibAllDen}{0}
\rqSet{rqFFReasonNotAdjudicableCUnstableLibAllPct}{n/a}
\rqSet{rqFFReasonNotAdjudicableCUnstableLibBOneNum}{0}
\rqSet{rqFFReasonNotAdjudicableCUnstableLibBOneDen}{0}
\rqSet{rqFFReasonNotAdjudicableCUnstableLibBOnePct}{n/a}
\rqSet{rqFFReasonNotAdjudicableCUnstableLibBTwoNum}{0}
\rqSet{rqFFReasonNotAdjudicableCUnstableLibBTwoDen}{0}
\rqSet{rqFFReasonNotAdjudicableCUnstableLibBTwoPct}{n/a}
\rqSet{rqFFReasonNotAdjudicableCUnstableLibPZeroNum}{0}
\rqSet{rqFFReasonNotAdjudicableCUnstableLibPZeroDen}{0}
\rqSet{rqFFReasonNotAdjudicableCUnstableLibPZeroPct}{n/a}
\rqSet{rqFFReasonNotAdjudicableCUnstableLibPOneNum}{0}
\rqSet{rqFFReasonNotAdjudicableCUnstableLibPOneDen}{0}
\rqSet{rqFFReasonNotAdjudicableCUnstableLibPOnePct}{n/a}
\rqSet{rqFFReasonNotAdjudicableUnstableComparisonLibAllNum}{0}
\rqSet{rqFFReasonNotAdjudicableUnstableComparisonLibAllDen}{0}
\rqSet{rqFFReasonNotAdjudicableUnstableComparisonLibAllPct}{n/a}
\rqSet{rqFFReasonNotAdjudicableUnstableComparisonLibBOneNum}{0}
\rqSet{rqFFReasonNotAdjudicableUnstableComparisonLibBOneDen}{0}
\rqSet{rqFFReasonNotAdjudicableUnstableComparisonLibBOnePct}{n/a}
\rqSet{rqFFReasonNotAdjudicableUnstableComparisonLibBTwoNum}{0}
\rqSet{rqFFReasonNotAdjudicableUnstableComparisonLibBTwoDen}{0}
\rqSet{rqFFReasonNotAdjudicableUnstableComparisonLibBTwoPct}{n/a}
\rqSet{rqFFReasonNotAdjudicableUnstableComparisonLibPZeroNum}{0}
\rqSet{rqFFReasonNotAdjudicableUnstableComparisonLibPZeroDen}{0}
\rqSet{rqFFReasonNotAdjudicableUnstableComparisonLibPZeroPct}{n/a}
\rqSet{rqFFReasonNotAdjudicableUnstableComparisonLibPOneNum}{0}
\rqSet{rqFFReasonNotAdjudicableUnstableComparisonLibPOneDen}{0}
\rqSet{rqFFReasonNotAdjudicableUnstableComparisonLibPOnePct}{n/a}
\rqSet{rqFFReasonNotReExaminedNoFailingVectorLibAllNum}{0}
\rqSet{rqFFReasonNotReExaminedNoFailingVectorLibAllDen}{0}
\rqSet{rqFFReasonNotReExaminedNoFailingVectorLibAllPct}{n/a}
\rqSet{rqFFReasonNotReExaminedNoFailingVectorLibBOneNum}{0}
\rqSet{rqFFReasonNotReExaminedNoFailingVectorLibBOneDen}{0}
\rqSet{rqFFReasonNotReExaminedNoFailingVectorLibBOnePct}{n/a}
\rqSet{rqFFReasonNotReExaminedNoFailingVectorLibBTwoNum}{0}
\rqSet{rqFFReasonNotReExaminedNoFailingVectorLibBTwoDen}{0}
\rqSet{rqFFReasonNotReExaminedNoFailingVectorLibBTwoPct}{n/a}
\rqSet{rqFFReasonNotReExaminedNoFailingVectorLibPZeroNum}{0}
\rqSet{rqFFReasonNotReExaminedNoFailingVectorLibPZeroDen}{0}
\rqSet{rqFFReasonNotReExaminedNoFailingVectorLibPZeroPct}{n/a}
\rqSet{rqFFReasonNotReExaminedNoFailingVectorLibPOneNum}{0}
\rqSet{rqFFReasonNotReExaminedNoFailingVectorLibPOneDen}{0}
\rqSet{rqFFReasonNotReExaminedNoFailingVectorLibPOnePct}{n/a}
\rqSet{rqFFReasonNotReExaminedReportMissingLibAllNum}{0}
\rqSet{rqFFReasonNotReExaminedReportMissingLibAllDen}{0}
\rqSet{rqFFReasonNotReExaminedReportMissingLibAllPct}{n/a}
\rqSet{rqFFReasonNotReExaminedReportMissingLibBOneNum}{0}
\rqSet{rqFFReasonNotReExaminedReportMissingLibBOneDen}{0}
\rqSet{rqFFReasonNotReExaminedReportMissingLibBOnePct}{n/a}
\rqSet{rqFFReasonNotReExaminedReportMissingLibBTwoNum}{0}
\rqSet{rqFFReasonNotReExaminedReportMissingLibBTwoDen}{0}
\rqSet{rqFFReasonNotReExaminedReportMissingLibBTwoPct}{n/a}
\rqSet{rqFFReasonNotReExaminedReportMissingLibPZeroNum}{0}
\rqSet{rqFFReasonNotReExaminedReportMissingLibPZeroDen}{0}
\rqSet{rqFFReasonNotReExaminedReportMissingLibPZeroPct}{n/a}
\rqSet{rqFFReasonNotReExaminedReportMissingLibPOneNum}{0}
\rqSet{rqFFReasonNotReExaminedReportMissingLibPOneDen}{0}
\rqSet{rqFFReasonNotReExaminedReportMissingLibPOnePct}{n/a}
\rqSet{rqFFReasonNotReExaminedTimeoutOnlyRulePendingLibAllNum}{0}
\rqSet{rqFFReasonNotReExaminedTimeoutOnlyRulePendingLibAllDen}{0}
\rqSet{rqFFReasonNotReExaminedTimeoutOnlyRulePendingLibAllPct}{n/a}
\rqSet{rqFFReasonNotReExaminedTimeoutOnlyRulePendingLibBOneNum}{0}
\rqSet{rqFFReasonNotReExaminedTimeoutOnlyRulePendingLibBOneDen}{0}
\rqSet{rqFFReasonNotReExaminedTimeoutOnlyRulePendingLibBOnePct}{n/a}
\rqSet{rqFFReasonNotReExaminedTimeoutOnlyRulePendingLibBTwoNum}{0}
\rqSet{rqFFReasonNotReExaminedTimeoutOnlyRulePendingLibBTwoDen}{0}
\rqSet{rqFFReasonNotReExaminedTimeoutOnlyRulePendingLibBTwoPct}{n/a}
\rqSet{rqFFReasonNotReExaminedTimeoutOnlyRulePendingLibPZeroNum}{0}
\rqSet{rqFFReasonNotReExaminedTimeoutOnlyRulePendingLibPZeroDen}{0}
\rqSet{rqFFReasonNotReExaminedTimeoutOnlyRulePendingLibPZeroPct}{n/a}
\rqSet{rqFFReasonNotReExaminedTimeoutOnlyRulePendingLibPOneNum}{0}
\rqSet{rqFFReasonNotReExaminedTimeoutOnlyRulePendingLibPOneDen}{0}
\rqSet{rqFFReasonNotReExaminedTimeoutOnlyRulePendingLibPOnePct}{n/a}
\rqSet{rqFFReasonNotReExaminedBuildFailedInHarnessLibAllNum}{0}
\rqSet{rqFFReasonNotReExaminedBuildFailedInHarnessLibAllDen}{0}
\rqSet{rqFFReasonNotReExaminedBuildFailedInHarnessLibAllPct}{n/a}
\rqSet{rqFFReasonNotReExaminedBuildFailedInHarnessLibBOneNum}{0}
\rqSet{rqFFReasonNotReExaminedBuildFailedInHarnessLibBOneDen}{0}
\rqSet{rqFFReasonNotReExaminedBuildFailedInHarnessLibBOnePct}{n/a}
\rqSet{rqFFReasonNotReExaminedBuildFailedInHarnessLibBTwoNum}{0}
\rqSet{rqFFReasonNotReExaminedBuildFailedInHarnessLibBTwoDen}{0}
\rqSet{rqFFReasonNotReExaminedBuildFailedInHarnessLibBTwoPct}{n/a}
\rqSet{rqFFReasonNotReExaminedBuildFailedInHarnessLibPZeroNum}{0}
\rqSet{rqFFReasonNotReExaminedBuildFailedInHarnessLibPZeroDen}{0}
\rqSet{rqFFReasonNotReExaminedBuildFailedInHarnessLibPZeroPct}{n/a}
\rqSet{rqFFReasonNotReExaminedBuildFailedInHarnessLibPOneNum}{0}
\rqSet{rqFFReasonNotReExaminedBuildFailedInHarnessLibPOneDen}{0}
\rqSet{rqFFReasonNotReExaminedBuildFailedInHarnessLibPOnePct}{n/a}
\rqSet{rqFFReasonNotReExaminedIncompleteHarnessRunVectorUnreadableLibAllNum}{0}
\rqSet{rqFFReasonNotReExaminedIncompleteHarnessRunVectorUnreadableLibAllDen}{0}
\rqSet{rqFFReasonNotReExaminedIncompleteHarnessRunVectorUnreadableLibAllPct}{n/a}
\rqSet{rqFFReasonNotReExaminedIncompleteHarnessRunVectorUnreadableLibBOneNum}{0}
\rqSet{rqFFReasonNotReExaminedIncompleteHarnessRunVectorUnreadableLibBOneDen}{0}
\rqSet{rqFFReasonNotReExaminedIncompleteHarnessRunVectorUnreadableLibBOnePct}{n/a}
\rqSet{rqFFReasonNotReExaminedIncompleteHarnessRunVectorUnreadableLibBTwoNum}{0}
\rqSet{rqFFReasonNotReExaminedIncompleteHarnessRunVectorUnreadableLibBTwoDen}{0}
\rqSet{rqFFReasonNotReExaminedIncompleteHarnessRunVectorUnreadableLibBTwoPct}{n/a}
\rqSet{rqFFReasonNotReExaminedIncompleteHarnessRunVectorUnreadableLibPZeroNum}{0}
\rqSet{rqFFReasonNotReExaminedIncompleteHarnessRunVectorUnreadableLibPZeroDen}{0}
\rqSet{rqFFReasonNotReExaminedIncompleteHarnessRunVectorUnreadableLibPZeroPct}{n/a}
\rqSet{rqFFReasonNotReExaminedIncompleteHarnessRunVectorUnreadableLibPOneNum}{0}
\rqSet{rqFFReasonNotReExaminedIncompleteHarnessRunVectorUnreadableLibPOneDen}{0}
\rqSet{rqFFReasonNotReExaminedIncompleteHarnessRunVectorUnreadableLibPOnePct}{n/a}
\rqSet{rqFFReasonNotReExaminedIncompleteHarnessRunInvalidInputLibAllNum}{0}
\rqSet{rqFFReasonNotReExaminedIncompleteHarnessRunInvalidInputLibAllDen}{0}
\rqSet{rqFFReasonNotReExaminedIncompleteHarnessRunInvalidInputLibAllPct}{n/a}
\rqSet{rqFFReasonNotReExaminedIncompleteHarnessRunInvalidInputLibBOneNum}{0}
\rqSet{rqFFReasonNotReExaminedIncompleteHarnessRunInvalidInputLibBOneDen}{0}
\rqSet{rqFFReasonNotReExaminedIncompleteHarnessRunInvalidInputLibBOnePct}{n/a}
\rqSet{rqFFReasonNotReExaminedIncompleteHarnessRunInvalidInputLibBTwoNum}{0}
\rqSet{rqFFReasonNotReExaminedIncompleteHarnessRunInvalidInputLibBTwoDen}{0}
\rqSet{rqFFReasonNotReExaminedIncompleteHarnessRunInvalidInputLibBTwoPct}{n/a}
\rqSet{rqFFReasonNotReExaminedIncompleteHarnessRunInvalidInputLibPZeroNum}{0}
\rqSet{rqFFReasonNotReExaminedIncompleteHarnessRunInvalidInputLibPZeroDen}{0}
\rqSet{rqFFReasonNotReExaminedIncompleteHarnessRunInvalidInputLibPZeroPct}{n/a}
\rqSet{rqFFReasonNotReExaminedIncompleteHarnessRunInvalidInputLibPOneNum}{0}
\rqSet{rqFFReasonNotReExaminedIncompleteHarnessRunInvalidInputLibPOneDen}{0}
\rqSet{rqFFReasonNotReExaminedIncompleteHarnessRunInvalidInputLibPOnePct}{n/a}
\rqSet{rqFFReasonNotReExaminedIncompleteHarnessRunRunUnobservedLibAllNum}{0}
\rqSet{rqFFReasonNotReExaminedIncompleteHarnessRunRunUnobservedLibAllDen}{0}
\rqSet{rqFFReasonNotReExaminedIncompleteHarnessRunRunUnobservedLibAllPct}{n/a}
\rqSet{rqFFReasonNotReExaminedIncompleteHarnessRunRunUnobservedLibBOneNum}{0}
\rqSet{rqFFReasonNotReExaminedIncompleteHarnessRunRunUnobservedLibBOneDen}{0}
\rqSet{rqFFReasonNotReExaminedIncompleteHarnessRunRunUnobservedLibBOnePct}{n/a}
\rqSet{rqFFReasonNotReExaminedIncompleteHarnessRunRunUnobservedLibBTwoNum}{0}
\rqSet{rqFFReasonNotReExaminedIncompleteHarnessRunRunUnobservedLibBTwoDen}{0}
\rqSet{rqFFReasonNotReExaminedIncompleteHarnessRunRunUnobservedLibBTwoPct}{n/a}
\rqSet{rqFFReasonNotReExaminedIncompleteHarnessRunRunUnobservedLibPZeroNum}{0}
\rqSet{rqFFReasonNotReExaminedIncompleteHarnessRunRunUnobservedLibPZeroDen}{0}
\rqSet{rqFFReasonNotReExaminedIncompleteHarnessRunRunUnobservedLibPZeroPct}{n/a}
\rqSet{rqFFReasonNotReExaminedIncompleteHarnessRunRunUnobservedLibPOneNum}{0}
\rqSet{rqFFReasonNotReExaminedIncompleteHarnessRunRunUnobservedLibPOneDen}{0}
\rqSet{rqFFReasonNotReExaminedIncompleteHarnessRunRunUnobservedLibPOnePct}{n/a}
\rqSet{rqFFReasonNotReExaminedIncompleteHarnessRunSanBuildFailedLibAllNum}{0}
\rqSet{rqFFReasonNotReExaminedIncompleteHarnessRunSanBuildFailedLibAllDen}{0}
\rqSet{rqFFReasonNotReExaminedIncompleteHarnessRunSanBuildFailedLibAllPct}{n/a}
\rqSet{rqFFReasonNotReExaminedIncompleteHarnessRunSanBuildFailedLibBOneNum}{0}
\rqSet{rqFFReasonNotReExaminedIncompleteHarnessRunSanBuildFailedLibBOneDen}{0}
\rqSet{rqFFReasonNotReExaminedIncompleteHarnessRunSanBuildFailedLibBOnePct}{n/a}
\rqSet{rqFFReasonNotReExaminedIncompleteHarnessRunSanBuildFailedLibBTwoNum}{0}
\rqSet{rqFFReasonNotReExaminedIncompleteHarnessRunSanBuildFailedLibBTwoDen}{0}
\rqSet{rqFFReasonNotReExaminedIncompleteHarnessRunSanBuildFailedLibBTwoPct}{n/a}
\rqSet{rqFFReasonNotReExaminedIncompleteHarnessRunSanBuildFailedLibPZeroNum}{0}
\rqSet{rqFFReasonNotReExaminedIncompleteHarnessRunSanBuildFailedLibPZeroDen}{0}
\rqSet{rqFFReasonNotReExaminedIncompleteHarnessRunSanBuildFailedLibPZeroPct}{n/a}
\rqSet{rqFFReasonNotReExaminedIncompleteHarnessRunSanBuildFailedLibPOneNum}{0}
\rqSet{rqFFReasonNotReExaminedIncompleteHarnessRunSanBuildFailedLibPOneDen}{0}
\rqSet{rqFFReasonNotReExaminedIncompleteHarnessRunSanBuildFailedLibPOnePct}{n/a}
\rqSet{rqFFReasonNotReExaminedIncompleteHarnessRunSanInfraLibAllNum}{0}
\rqSet{rqFFReasonNotReExaminedIncompleteHarnessRunSanInfraLibAllDen}{0}
\rqSet{rqFFReasonNotReExaminedIncompleteHarnessRunSanInfraLibAllPct}{n/a}
\rqSet{rqFFReasonNotReExaminedIncompleteHarnessRunSanInfraLibBOneNum}{0}
\rqSet{rqFFReasonNotReExaminedIncompleteHarnessRunSanInfraLibBOneDen}{0}
\rqSet{rqFFReasonNotReExaminedIncompleteHarnessRunSanInfraLibBOnePct}{n/a}
\rqSet{rqFFReasonNotReExaminedIncompleteHarnessRunSanInfraLibBTwoNum}{0}
\rqSet{rqFFReasonNotReExaminedIncompleteHarnessRunSanInfraLibBTwoDen}{0}
\rqSet{rqFFReasonNotReExaminedIncompleteHarnessRunSanInfraLibBTwoPct}{n/a}
\rqSet{rqFFReasonNotReExaminedIncompleteHarnessRunSanInfraLibPZeroNum}{0}
\rqSet{rqFFReasonNotReExaminedIncompleteHarnessRunSanInfraLibPZeroDen}{0}
\rqSet{rqFFReasonNotReExaminedIncompleteHarnessRunSanInfraLibPZeroPct}{n/a}
\rqSet{rqFFReasonNotReExaminedIncompleteHarnessRunSanInfraLibPOneNum}{0}
\rqSet{rqFFReasonNotReExaminedIncompleteHarnessRunSanInfraLibPOneDen}{0}
\rqSet{rqFFReasonNotReExaminedIncompleteHarnessRunSanInfraLibPOnePct}{n/a}
\rqSet{rqFFChanStderrLibAllNum}{0}
\rqSet{rqFFChanStderrLibAllDen}{1}
\rqSet{rqFFChanStderrLibAllPct}{0.0}
\rqSet{rqFFChanStderrLibBOneNum}{0}
\rqSet{rqFFChanStderrLibBOneDen}{0}
\rqSet{rqFFChanStderrLibBOnePct}{n/a}
\rqSet{rqFFChanStderrLibBTwoNum}{0}
\rqSet{rqFFChanStderrLibBTwoDen}{1}
\rqSet{rqFFChanStderrLibBTwoPct}{0.0}
\rqSet{rqFFChanStderrLibPZeroNum}{0}
\rqSet{rqFFChanStderrLibPZeroDen}{0}
\rqSet{rqFFChanStderrLibPZeroPct}{n/a}
\rqSet{rqFFChanStderrLibPOneNum}{0}
\rqSet{rqFFChanStderrLibPOneDen}{0}
\rqSet{rqFFChanStderrLibPOnePct}{n/a}
\rqSet{rqFFChanStdoutLibAllNum}{0}
\rqSet{rqFFChanStdoutLibAllDen}{1}
\rqSet{rqFFChanStdoutLibAllPct}{0.0}
\rqSet{rqFFChanStdoutLibBOneNum}{0}
\rqSet{rqFFChanStdoutLibBOneDen}{0}
\rqSet{rqFFChanStdoutLibBOnePct}{n/a}
\rqSet{rqFFChanStdoutLibBTwoNum}{0}
\rqSet{rqFFChanStdoutLibBTwoDen}{1}
\rqSet{rqFFChanStdoutLibBTwoPct}{0.0}
\rqSet{rqFFChanStdoutLibPZeroNum}{0}
\rqSet{rqFFChanStdoutLibPZeroDen}{0}
\rqSet{rqFFChanStdoutLibPZeroPct}{n/a}
\rqSet{rqFFChanStdoutLibPOneNum}{0}
\rqSet{rqFFChanStdoutLibPOneDen}{0}
\rqSet{rqFFChanStdoutLibPOnePct}{n/a}
\rqSet{rqFFChanStateFloatLibAllNum}{1}
\rqSet{rqFFChanStateFloatLibAllDen}{1}
\rqSet{rqFFChanStateFloatLibAllPct}{100.0}
\rqSet{rqFFChanStateFloatLibBOneNum}{0}
\rqSet{rqFFChanStateFloatLibBOneDen}{0}
\rqSet{rqFFChanStateFloatLibBOnePct}{n/a}
\rqSet{rqFFChanStateFloatLibBTwoNum}{1}
\rqSet{rqFFChanStateFloatLibBTwoDen}{1}
\rqSet{rqFFChanStateFloatLibBTwoPct}{100.0}
\rqSet{rqFFChanStateFloatLibPZeroNum}{0}
\rqSet{rqFFChanStateFloatLibPZeroDen}{0}
\rqSet{rqFFChanStateFloatLibPZeroPct}{n/a}
\rqSet{rqFFChanStateFloatLibPOneNum}{0}
\rqSet{rqFFChanStateFloatLibPOneDen}{0}
\rqSet{rqFFChanStateFloatLibPOnePct}{n/a}
\rqSet{rqFFChanStateOtherLibAllNum}{0}
\rqSet{rqFFChanStateOtherLibAllDen}{1}
\rqSet{rqFFChanStateOtherLibAllPct}{0.0}
\rqSet{rqFFChanStateOtherLibBOneNum}{0}
\rqSet{rqFFChanStateOtherLibBOneDen}{0}
\rqSet{rqFFChanStateOtherLibBOnePct}{n/a}
\rqSet{rqFFChanStateOtherLibBTwoNum}{0}
\rqSet{rqFFChanStateOtherLibBTwoDen}{1}
\rqSet{rqFFChanStateOtherLibBTwoPct}{0.0}
\rqSet{rqFFChanStateOtherLibPZeroNum}{0}
\rqSet{rqFFChanStateOtherLibPZeroDen}{0}
\rqSet{rqFFChanStateOtherLibPZeroPct}{n/a}
\rqSet{rqFFChanStateOtherLibPOneNum}{0}
\rqSet{rqFFChanStateOtherLibPOneDen}{0}
\rqSet{rqFFChanStateOtherLibPOnePct}{n/a}
\rqSet{rqFFChanStateLibAllNum}{0}
\rqSet{rqFFChanStateLibAllDen}{1}
\rqSet{rqFFChanStateLibAllPct}{0.0}
\rqSet{rqFFChanStateLibBOneNum}{0}
\rqSet{rqFFChanStateLibBOneDen}{0}
\rqSet{rqFFChanStateLibBOnePct}{n/a}
\rqSet{rqFFChanStateLibBTwoNum}{0}
\rqSet{rqFFChanStateLibBTwoDen}{1}
\rqSet{rqFFChanStateLibBTwoPct}{0.0}
\rqSet{rqFFChanStateLibPZeroNum}{0}
\rqSet{rqFFChanStateLibPZeroDen}{0}
\rqSet{rqFFChanStateLibPZeroPct}{n/a}
\rqSet{rqFFChanStateLibPOneNum}{0}
\rqSet{rqFFChanStateLibPOneDen}{0}
\rqSet{rqFFChanStateLibPOnePct}{n/a}
\rqSet{rqFFChanResultPanicLibAllNum}{0}
\rqSet{rqFFChanResultPanicLibAllDen}{1}
\rqSet{rqFFChanResultPanicLibAllPct}{0.0}
\rqSet{rqFFChanResultPanicLibBOneNum}{0}
\rqSet{rqFFChanResultPanicLibBOneDen}{0}
\rqSet{rqFFChanResultPanicLibBOnePct}{n/a}
\rqSet{rqFFChanResultPanicLibBTwoNum}{0}
\rqSet{rqFFChanResultPanicLibBTwoDen}{1}
\rqSet{rqFFChanResultPanicLibBTwoPct}{0.0}
\rqSet{rqFFChanResultPanicLibPZeroNum}{0}
\rqSet{rqFFChanResultPanicLibPZeroDen}{0}
\rqSet{rqFFChanResultPanicLibPZeroPct}{n/a}
\rqSet{rqFFChanResultPanicLibPOneNum}{0}
\rqSet{rqFFChanResultPanicLibPOneDen}{0}
\rqSet{rqFFChanResultPanicLibPOnePct}{n/a}
\rqSet{rqFFChanResultSegmentationFaultLibAllNum}{0}
\rqSet{rqFFChanResultSegmentationFaultLibAllDen}{1}
\rqSet{rqFFChanResultSegmentationFaultLibAllPct}{0.0}
\rqSet{rqFFChanResultSegmentationFaultLibBOneNum}{0}
\rqSet{rqFFChanResultSegmentationFaultLibBOneDen}{0}
\rqSet{rqFFChanResultSegmentationFaultLibBOnePct}{n/a}
\rqSet{rqFFChanResultSegmentationFaultLibBTwoNum}{0}
\rqSet{rqFFChanResultSegmentationFaultLibBTwoDen}{1}
\rqSet{rqFFChanResultSegmentationFaultLibBTwoPct}{0.0}
\rqSet{rqFFChanResultSegmentationFaultLibPZeroNum}{0}
\rqSet{rqFFChanResultSegmentationFaultLibPZeroDen}{0}
\rqSet{rqFFChanResultSegmentationFaultLibPZeroPct}{n/a}
\rqSet{rqFFChanResultSegmentationFaultLibPOneNum}{0}
\rqSet{rqFFChanResultSegmentationFaultLibPOneDen}{0}
\rqSet{rqFFChanResultSegmentationFaultLibPOnePct}{n/a}
\rqSet{rqFFChanResultUnknownFailureLibAllNum}{0}
\rqSet{rqFFChanResultUnknownFailureLibAllDen}{1}
\rqSet{rqFFChanResultUnknownFailureLibAllPct}{0.0}
\rqSet{rqFFChanResultUnknownFailureLibBOneNum}{0}
\rqSet{rqFFChanResultUnknownFailureLibBOneDen}{0}
\rqSet{rqFFChanResultUnknownFailureLibBOnePct}{n/a}
\rqSet{rqFFChanResultUnknownFailureLibBTwoNum}{0}
\rqSet{rqFFChanResultUnknownFailureLibBTwoDen}{1}
\rqSet{rqFFChanResultUnknownFailureLibBTwoPct}{0.0}
\rqSet{rqFFChanResultUnknownFailureLibPZeroNum}{0}
\rqSet{rqFFChanResultUnknownFailureLibPZeroDen}{0}
\rqSet{rqFFChanResultUnknownFailureLibPZeroPct}{n/a}
\rqSet{rqFFChanResultUnknownFailureLibPOneNum}{0}
\rqSet{rqFFChanResultUnknownFailureLibPOneDen}{0}
\rqSet{rqFFChanResultUnknownFailureLibPOnePct}{n/a}
\rqSet{rqFFChanResultTimeoutLibAllNum}{0}
\rqSet{rqFFChanResultTimeoutLibAllDen}{1}
\rqSet{rqFFChanResultTimeoutLibAllPct}{0.0}
\rqSet{rqFFChanResultTimeoutLibBOneNum}{0}
\rqSet{rqFFChanResultTimeoutLibBOneDen}{0}
\rqSet{rqFFChanResultTimeoutLibBOnePct}{n/a}
\rqSet{rqFFChanResultTimeoutLibBTwoNum}{0}
\rqSet{rqFFChanResultTimeoutLibBTwoDen}{1}
\rqSet{rqFFChanResultTimeoutLibBTwoPct}{0.0}
\rqSet{rqFFChanResultTimeoutLibPZeroNum}{0}
\rqSet{rqFFChanResultTimeoutLibPZeroDen}{0}
\rqSet{rqFFChanResultTimeoutLibPZeroPct}{n/a}
\rqSet{rqFFChanResultTimeoutLibPOneNum}{0}
\rqSet{rqFFChanResultTimeoutLibPOneDen}{0}
\rqSet{rqFFChanResultTimeoutLibPOnePct}{n/a}
\rqSet{rqFFCTermAnyLibAllNum}{0}
\rqSet{rqFFCTermAnyLibAllDen}{303}
\rqSet{rqFFCTermAnyLibAllPct}{0.0}
\rqSet{rqFFCTermAnyLibBOneNum}{0}
\rqSet{rqFFCTermAnyLibBOneDen}{155}
\rqSet{rqFFCTermAnyLibBOnePct}{0.0}
\rqSet{rqFFCTermAnyLibBTwoNum}{0}
\rqSet{rqFFCTermAnyLibBTwoDen}{92}
\rqSet{rqFFCTermAnyLibBTwoPct}{0.0}
\rqSet{rqFFCTermAnyLibPZeroNum}{0}
\rqSet{rqFFCTermAnyLibPZeroDen}{0}
\rqSet{rqFFCTermAnyLibPZeroPct}{n/a}
\rqSet{rqFFCTermAnyLibPOneNum}{0}
\rqSet{rqFFCTermAnyLibPOneDen}{56}
\rqSet{rqFFCTermAnyLibPOnePct}{0.0}
\rqSet{rqFFCTermExitOneZeroOneLibAllNum}{0}
\rqSet{rqFFCTermExitOneZeroOneLibAllDen}{303}
\rqSet{rqFFCTermExitOneZeroOneLibAllPct}{0.0}
\rqSet{rqFFCTermExitOneZeroOneLibBOneNum}{0}
\rqSet{rqFFCTermExitOneZeroOneLibBOneDen}{155}
\rqSet{rqFFCTermExitOneZeroOneLibBOnePct}{0.0}
\rqSet{rqFFCTermExitOneZeroOneLibBTwoNum}{0}
\rqSet{rqFFCTermExitOneZeroOneLibBTwoDen}{92}
\rqSet{rqFFCTermExitOneZeroOneLibBTwoPct}{0.0}
\rqSet{rqFFCTermExitOneZeroOneLibPZeroNum}{0}
\rqSet{rqFFCTermExitOneZeroOneLibPZeroDen}{0}
\rqSet{rqFFCTermExitOneZeroOneLibPZeroPct}{n/a}
\rqSet{rqFFCTermExitOneZeroOneLibPOneNum}{0}
\rqSet{rqFFCTermExitOneZeroOneLibPOneDen}{56}
\rqSet{rqFFCTermExitOneZeroOneLibPOnePct}{0.0}
\rqSet{rqFFCTermSignalLibAllNum}{0}
\rqSet{rqFFCTermSignalLibAllDen}{303}
\rqSet{rqFFCTermSignalLibAllPct}{0.0}
\rqSet{rqFFCTermSignalLibBOneNum}{0}
\rqSet{rqFFCTermSignalLibBOneDen}{155}
\rqSet{rqFFCTermSignalLibBOnePct}{0.0}
\rqSet{rqFFCTermSignalLibBTwoNum}{0}
\rqSet{rqFFCTermSignalLibBTwoDen}{92}
\rqSet{rqFFCTermSignalLibBTwoPct}{0.0}
\rqSet{rqFFCTermSignalLibPZeroNum}{0}
\rqSet{rqFFCTermSignalLibPZeroDen}{0}
\rqSet{rqFFCTermSignalLibPZeroPct}{n/a}
\rqSet{rqFFCTermSignalLibPOneNum}{0}
\rqSet{rqFFCTermSignalLibPOneDen}{56}
\rqSet{rqFFCTermSignalLibPOnePct}{0.0}
\rqSet{rqFFCTermFFAnyLibAllNum}{0}
\rqSet{rqFFCTermFFAnyLibAllDen}{2}
\rqSet{rqFFCTermFFAnyLibAllPct}{0.0}
\rqSet{rqFFCTermFFAnyLibBOneNum}{0}
\rqSet{rqFFCTermFFAnyLibBOneDen}{0}
\rqSet{rqFFCTermFFAnyLibBOnePct}{n/a}
\rqSet{rqFFCTermFFAnyLibBTwoNum}{0}
\rqSet{rqFFCTermFFAnyLibBTwoDen}{2}
\rqSet{rqFFCTermFFAnyLibBTwoPct}{0.0}
\rqSet{rqFFCTermFFAnyLibPZeroNum}{0}
\rqSet{rqFFCTermFFAnyLibPZeroDen}{0}
\rqSet{rqFFCTermFFAnyLibPZeroPct}{n/a}
\rqSet{rqFFCTermFFAnyLibPOneNum}{0}
\rqSet{rqFFCTermFFAnyLibPOneDen}{0}
\rqSet{rqFFCTermFFAnyLibPOnePct}{n/a}
\rqSet{rqFFCTermFFExitOneZeroOneLibAllNum}{0}
\rqSet{rqFFCTermFFExitOneZeroOneLibAllDen}{2}
\rqSet{rqFFCTermFFExitOneZeroOneLibAllPct}{0.0}
\rqSet{rqFFCTermFFExitOneZeroOneLibBOneNum}{0}
\rqSet{rqFFCTermFFExitOneZeroOneLibBOneDen}{0}
\rqSet{rqFFCTermFFExitOneZeroOneLibBOnePct}{n/a}
\rqSet{rqFFCTermFFExitOneZeroOneLibBTwoNum}{0}
\rqSet{rqFFCTermFFExitOneZeroOneLibBTwoDen}{2}
\rqSet{rqFFCTermFFExitOneZeroOneLibBTwoPct}{0.0}
\rqSet{rqFFCTermFFExitOneZeroOneLibPZeroNum}{0}
\rqSet{rqFFCTermFFExitOneZeroOneLibPZeroDen}{0}
\rqSet{rqFFCTermFFExitOneZeroOneLibPZeroPct}{n/a}
\rqSet{rqFFCTermFFExitOneZeroOneLibPOneNum}{0}
\rqSet{rqFFCTermFFExitOneZeroOneLibPOneDen}{0}
\rqSet{rqFFCTermFFExitOneZeroOneLibPOnePct}{n/a}
\rqSet{rqFFCTermFFSignalLibAllNum}{0}
\rqSet{rqFFCTermFFSignalLibAllDen}{2}
\rqSet{rqFFCTermFFSignalLibAllPct}{0.0}
\rqSet{rqFFCTermFFSignalLibBOneNum}{0}
\rqSet{rqFFCTermFFSignalLibBOneDen}{0}
\rqSet{rqFFCTermFFSignalLibBOnePct}{n/a}
\rqSet{rqFFCTermFFSignalLibBTwoNum}{0}
\rqSet{rqFFCTermFFSignalLibBTwoDen}{2}
\rqSet{rqFFCTermFFSignalLibBTwoPct}{0.0}
\rqSet{rqFFCTermFFSignalLibPZeroNum}{0}
\rqSet{rqFFCTermFFSignalLibPZeroDen}{0}
\rqSet{rqFFCTermFFSignalLibPZeroPct}{n/a}
\rqSet{rqFFCTermFFSignalLibPOneNum}{0}
\rqSet{rqFFCTermFFSignalLibPOneDen}{0}
\rqSet{rqFFCTermFFSignalLibPOnePct}{n/a}
\rqSet{rqFFVecSupportsRejectionLibAllNum}{300}
\rqSet{rqFFVecSupportsRejectionLibAllDen}{303}
\rqSet{rqFFVecSupportsRejectionLibAllPct}{99.0}
\rqSet{rqFFVecSupportsRejectionLibBOneNum}{155}
\rqSet{rqFFVecSupportsRejectionLibBOneDen}{155}
\rqSet{rqFFVecSupportsRejectionLibBOnePct}{100.0}
\rqSet{rqFFVecSupportsRejectionLibBTwoNum}{89}
\rqSet{rqFFVecSupportsRejectionLibBTwoDen}{92}
\rqSet{rqFFVecSupportsRejectionLibBTwoPct}{96.7}
\rqSet{rqFFVecSupportsRejectionLibPZeroNum}{0}
\rqSet{rqFFVecSupportsRejectionLibPZeroDen}{0}
\rqSet{rqFFVecSupportsRejectionLibPZeroPct}{n/a}
\rqSet{rqFFVecSupportsRejectionLibPOneNum}{56}
\rqSet{rqFFVecSupportsRejectionLibPOneDen}{56}
\rqSet{rqFFVecSupportsRejectionLibPOnePct}{100.0}
\rqSet{rqFFVecAgreesLibAllNum}{2}
\rqSet{rqFFVecAgreesLibAllDen}{303}
\rqSet{rqFFVecAgreesLibAllPct}{0.7}
\rqSet{rqFFVecAgreesLibBOneNum}{0}
\rqSet{rqFFVecAgreesLibBOneDen}{155}
\rqSet{rqFFVecAgreesLibBOnePct}{0.0}
\rqSet{rqFFVecAgreesLibBTwoNum}{2}
\rqSet{rqFFVecAgreesLibBTwoDen}{92}
\rqSet{rqFFVecAgreesLibBTwoPct}{2.2}
\rqSet{rqFFVecAgreesLibPZeroNum}{0}
\rqSet{rqFFVecAgreesLibPZeroDen}{0}
\rqSet{rqFFVecAgreesLibPZeroPct}{n/a}
\rqSet{rqFFVecAgreesLibPOneNum}{0}
\rqSet{rqFFVecAgreesLibPOneDen}{56}
\rqSet{rqFFVecAgreesLibPOnePct}{0.0}
\rqSet{rqFFVecNotAdjudicableUbInCLibAllNum}{0}
\rqSet{rqFFVecNotAdjudicableUbInCLibAllDen}{303}
\rqSet{rqFFVecNotAdjudicableUbInCLibAllPct}{0.0}
\rqSet{rqFFVecNotAdjudicableUbInCLibBOneNum}{0}
\rqSet{rqFFVecNotAdjudicableUbInCLibBOneDen}{155}
\rqSet{rqFFVecNotAdjudicableUbInCLibBOnePct}{0.0}
\rqSet{rqFFVecNotAdjudicableUbInCLibBTwoNum}{0}
\rqSet{rqFFVecNotAdjudicableUbInCLibBTwoDen}{92}
\rqSet{rqFFVecNotAdjudicableUbInCLibBTwoPct}{0.0}
\rqSet{rqFFVecNotAdjudicableUbInCLibPZeroNum}{0}
\rqSet{rqFFVecNotAdjudicableUbInCLibPZeroDen}{0}
\rqSet{rqFFVecNotAdjudicableUbInCLibPZeroPct}{n/a}
\rqSet{rqFFVecNotAdjudicableUbInCLibPOneNum}{0}
\rqSet{rqFFVecNotAdjudicableUbInCLibPOneDen}{56}
\rqSet{rqFFVecNotAdjudicableUbInCLibPOnePct}{0.0}
\rqSet{rqFFVecNotAdjudicableBothBudgetLibAllNum}{0}
\rqSet{rqFFVecNotAdjudicableBothBudgetLibAllDen}{303}
\rqSet{rqFFVecNotAdjudicableBothBudgetLibAllPct}{0.0}
\rqSet{rqFFVecNotAdjudicableBothBudgetLibBOneNum}{0}
\rqSet{rqFFVecNotAdjudicableBothBudgetLibBOneDen}{155}
\rqSet{rqFFVecNotAdjudicableBothBudgetLibBOnePct}{0.0}
\rqSet{rqFFVecNotAdjudicableBothBudgetLibBTwoNum}{0}
\rqSet{rqFFVecNotAdjudicableBothBudgetLibBTwoDen}{92}
\rqSet{rqFFVecNotAdjudicableBothBudgetLibBTwoPct}{0.0}
\rqSet{rqFFVecNotAdjudicableBothBudgetLibPZeroNum}{0}
\rqSet{rqFFVecNotAdjudicableBothBudgetLibPZeroDen}{0}
\rqSet{rqFFVecNotAdjudicableBothBudgetLibPZeroPct}{n/a}
\rqSet{rqFFVecNotAdjudicableBothBudgetLibPOneNum}{0}
\rqSet{rqFFVecNotAdjudicableBothBudgetLibPOneDen}{56}
\rqSet{rqFFVecNotAdjudicableBothBudgetLibPOnePct}{0.0}
\rqSet{rqFFVecNotAdjudicableSanitizerInconclusiveLibAllNum}{0}
\rqSet{rqFFVecNotAdjudicableSanitizerInconclusiveLibAllDen}{303}
\rqSet{rqFFVecNotAdjudicableSanitizerInconclusiveLibAllPct}{0.0}
\rqSet{rqFFVecNotAdjudicableSanitizerInconclusiveLibBOneNum}{0}
\rqSet{rqFFVecNotAdjudicableSanitizerInconclusiveLibBOneDen}{155}
\rqSet{rqFFVecNotAdjudicableSanitizerInconclusiveLibBOnePct}{0.0}
\rqSet{rqFFVecNotAdjudicableSanitizerInconclusiveLibBTwoNum}{0}
\rqSet{rqFFVecNotAdjudicableSanitizerInconclusiveLibBTwoDen}{92}
\rqSet{rqFFVecNotAdjudicableSanitizerInconclusiveLibBTwoPct}{0.0}
\rqSet{rqFFVecNotAdjudicableSanitizerInconclusiveLibPZeroNum}{0}
\rqSet{rqFFVecNotAdjudicableSanitizerInconclusiveLibPZeroDen}{0}
\rqSet{rqFFVecNotAdjudicableSanitizerInconclusiveLibPZeroPct}{n/a}
\rqSet{rqFFVecNotAdjudicableSanitizerInconclusiveLibPOneNum}{0}
\rqSet{rqFFVecNotAdjudicableSanitizerInconclusiveLibPOneDen}{56}
\rqSet{rqFFVecNotAdjudicableSanitizerInconclusiveLibPOnePct}{0.0}
\rqSet{rqFFVecNotAdjudicableCUnstableLibAllNum}{0}
\rqSet{rqFFVecNotAdjudicableCUnstableLibAllDen}{303}
\rqSet{rqFFVecNotAdjudicableCUnstableLibAllPct}{0.0}
\rqSet{rqFFVecNotAdjudicableCUnstableLibBOneNum}{0}
\rqSet{rqFFVecNotAdjudicableCUnstableLibBOneDen}{155}
\rqSet{rqFFVecNotAdjudicableCUnstableLibBOnePct}{0.0}
\rqSet{rqFFVecNotAdjudicableCUnstableLibBTwoNum}{0}
\rqSet{rqFFVecNotAdjudicableCUnstableLibBTwoDen}{92}
\rqSet{rqFFVecNotAdjudicableCUnstableLibBTwoPct}{0.0}
\rqSet{rqFFVecNotAdjudicableCUnstableLibPZeroNum}{0}
\rqSet{rqFFVecNotAdjudicableCUnstableLibPZeroDen}{0}
\rqSet{rqFFVecNotAdjudicableCUnstableLibPZeroPct}{n/a}
\rqSet{rqFFVecNotAdjudicableCUnstableLibPOneNum}{0}
\rqSet{rqFFVecNotAdjudicableCUnstableLibPOneDen}{56}
\rqSet{rqFFVecNotAdjudicableCUnstableLibPOnePct}{0.0}
\rqSet{rqFFVecNotAdjudicableUnstableComparisonLibAllNum}{1}
\rqSet{rqFFVecNotAdjudicableUnstableComparisonLibAllDen}{303}
\rqSet{rqFFVecNotAdjudicableUnstableComparisonLibAllPct}{0.3}
\rqSet{rqFFVecNotAdjudicableUnstableComparisonLibBOneNum}{0}
\rqSet{rqFFVecNotAdjudicableUnstableComparisonLibBOneDen}{155}
\rqSet{rqFFVecNotAdjudicableUnstableComparisonLibBOnePct}{0.0}
\rqSet{rqFFVecNotAdjudicableUnstableComparisonLibBTwoNum}{1}
\rqSet{rqFFVecNotAdjudicableUnstableComparisonLibBTwoDen}{92}
\rqSet{rqFFVecNotAdjudicableUnstableComparisonLibBTwoPct}{1.1}
\rqSet{rqFFVecNotAdjudicableUnstableComparisonLibPZeroNum}{0}
\rqSet{rqFFVecNotAdjudicableUnstableComparisonLibPZeroDen}{0}
\rqSet{rqFFVecNotAdjudicableUnstableComparisonLibPZeroPct}{n/a}
\rqSet{rqFFVecNotAdjudicableUnstableComparisonLibPOneNum}{0}
\rqSet{rqFFVecNotAdjudicableUnstableComparisonLibPOneDen}{56}
\rqSet{rqFFVecNotAdjudicableUnstableComparisonLibPOnePct}{0.0}
\rqSet{rqFFVecIncompleteVectorUnreadableLibAllNum}{0}
\rqSet{rqFFVecIncompleteVectorUnreadableLibAllDen}{303}
\rqSet{rqFFVecIncompleteVectorUnreadableLibAllPct}{0.0}
\rqSet{rqFFVecIncompleteVectorUnreadableLibBOneNum}{0}
\rqSet{rqFFVecIncompleteVectorUnreadableLibBOneDen}{155}
\rqSet{rqFFVecIncompleteVectorUnreadableLibBOnePct}{0.0}
\rqSet{rqFFVecIncompleteVectorUnreadableLibBTwoNum}{0}
\rqSet{rqFFVecIncompleteVectorUnreadableLibBTwoDen}{92}
\rqSet{rqFFVecIncompleteVectorUnreadableLibBTwoPct}{0.0}
\rqSet{rqFFVecIncompleteVectorUnreadableLibPZeroNum}{0}
\rqSet{rqFFVecIncompleteVectorUnreadableLibPZeroDen}{0}
\rqSet{rqFFVecIncompleteVectorUnreadableLibPZeroPct}{n/a}
\rqSet{rqFFVecIncompleteVectorUnreadableLibPOneNum}{0}
\rqSet{rqFFVecIncompleteVectorUnreadableLibPOneDen}{56}
\rqSet{rqFFVecIncompleteVectorUnreadableLibPOnePct}{0.0}
\rqSet{rqFFVecIncompleteInvalidInputLibAllNum}{0}
\rqSet{rqFFVecIncompleteInvalidInputLibAllDen}{303}
\rqSet{rqFFVecIncompleteInvalidInputLibAllPct}{0.0}
\rqSet{rqFFVecIncompleteInvalidInputLibBOneNum}{0}
\rqSet{rqFFVecIncompleteInvalidInputLibBOneDen}{155}
\rqSet{rqFFVecIncompleteInvalidInputLibBOnePct}{0.0}
\rqSet{rqFFVecIncompleteInvalidInputLibBTwoNum}{0}
\rqSet{rqFFVecIncompleteInvalidInputLibBTwoDen}{92}
\rqSet{rqFFVecIncompleteInvalidInputLibBTwoPct}{0.0}
\rqSet{rqFFVecIncompleteInvalidInputLibPZeroNum}{0}
\rqSet{rqFFVecIncompleteInvalidInputLibPZeroDen}{0}
\rqSet{rqFFVecIncompleteInvalidInputLibPZeroPct}{n/a}
\rqSet{rqFFVecIncompleteInvalidInputLibPOneNum}{0}
\rqSet{rqFFVecIncompleteInvalidInputLibPOneDen}{56}
\rqSet{rqFFVecIncompleteInvalidInputLibPOnePct}{0.0}
\rqSet{rqFFVecIncompleteRunUnobservedLibAllNum}{0}
\rqSet{rqFFVecIncompleteRunUnobservedLibAllDen}{303}
\rqSet{rqFFVecIncompleteRunUnobservedLibAllPct}{0.0}
\rqSet{rqFFVecIncompleteRunUnobservedLibBOneNum}{0}
\rqSet{rqFFVecIncompleteRunUnobservedLibBOneDen}{155}
\rqSet{rqFFVecIncompleteRunUnobservedLibBOnePct}{0.0}
\rqSet{rqFFVecIncompleteRunUnobservedLibBTwoNum}{0}
\rqSet{rqFFVecIncompleteRunUnobservedLibBTwoDen}{92}
\rqSet{rqFFVecIncompleteRunUnobservedLibBTwoPct}{0.0}
\rqSet{rqFFVecIncompleteRunUnobservedLibPZeroNum}{0}
\rqSet{rqFFVecIncompleteRunUnobservedLibPZeroDen}{0}
\rqSet{rqFFVecIncompleteRunUnobservedLibPZeroPct}{n/a}
\rqSet{rqFFVecIncompleteRunUnobservedLibPOneNum}{0}
\rqSet{rqFFVecIncompleteRunUnobservedLibPOneDen}{56}
\rqSet{rqFFVecIncompleteRunUnobservedLibPOnePct}{0.0}
\rqSet{rqFFVecIncompleteSanBuildFailedLibAllNum}{0}
\rqSet{rqFFVecIncompleteSanBuildFailedLibAllDen}{303}
\rqSet{rqFFVecIncompleteSanBuildFailedLibAllPct}{0.0}
\rqSet{rqFFVecIncompleteSanBuildFailedLibBOneNum}{0}
\rqSet{rqFFVecIncompleteSanBuildFailedLibBOneDen}{155}
\rqSet{rqFFVecIncompleteSanBuildFailedLibBOnePct}{0.0}
\rqSet{rqFFVecIncompleteSanBuildFailedLibBTwoNum}{0}
\rqSet{rqFFVecIncompleteSanBuildFailedLibBTwoDen}{92}
\rqSet{rqFFVecIncompleteSanBuildFailedLibBTwoPct}{0.0}
\rqSet{rqFFVecIncompleteSanBuildFailedLibPZeroNum}{0}
\rqSet{rqFFVecIncompleteSanBuildFailedLibPZeroDen}{0}
\rqSet{rqFFVecIncompleteSanBuildFailedLibPZeroPct}{n/a}
\rqSet{rqFFVecIncompleteSanBuildFailedLibPOneNum}{0}
\rqSet{rqFFVecIncompleteSanBuildFailedLibPOneDen}{56}
\rqSet{rqFFVecIncompleteSanBuildFailedLibPOnePct}{0.0}
\rqSet{rqFFVecIncompleteSanInfraLibAllNum}{0}
\rqSet{rqFFVecIncompleteSanInfraLibAllDen}{303}
\rqSet{rqFFVecIncompleteSanInfraLibAllPct}{0.0}
\rqSet{rqFFVecIncompleteSanInfraLibBOneNum}{0}
\rqSet{rqFFVecIncompleteSanInfraLibBOneDen}{155}
\rqSet{rqFFVecIncompleteSanInfraLibBOnePct}{0.0}
\rqSet{rqFFVecIncompleteSanInfraLibBTwoNum}{0}
\rqSet{rqFFVecIncompleteSanInfraLibBTwoDen}{92}
\rqSet{rqFFVecIncompleteSanInfraLibBTwoPct}{0.0}
\rqSet{rqFFVecIncompleteSanInfraLibPZeroNum}{0}
\rqSet{rqFFVecIncompleteSanInfraLibPZeroDen}{0}
\rqSet{rqFFVecIncompleteSanInfraLibPZeroPct}{n/a}
\rqSet{rqFFVecIncompleteSanInfraLibPOneNum}{0}
\rqSet{rqFFVecIncompleteSanInfraLibPOneDen}{56}
\rqSet{rqFFVecIncompleteSanInfraLibPOnePct}{0.0}
\rqSet{rqRobFFNoStderrLibAllNum}{1}
\rqSet{rqRobFFNoStderrLibAllDen}{111}
\rqSet{rqRobFFNoStderrLibAllPct}{0.9}
\rqSet{rqRobFFNoStderrLibAllKiroNum}{0}
\rqSet{rqRobFFNoStderrLibAllKiroDen}{0}
\rqSet{rqRobFFNoStderrLibAllKiroPct}{n/a}
\rqSet{rqRobFFNoStderrLibAllClaudeNum}{0}
\rqSet{rqRobFFNoStderrLibAllClaudeDen}{11}
\rqSet{rqRobFFNoStderrLibAllClaudePct}{0.0}
\rqSet{rqRobFFNoStderrLibAllCodexNum}{0}
\rqSet{rqRobFFNoStderrLibAllCodexDen}{45}
\rqSet{rqRobFFNoStderrLibAllCodexPct}{0.0}
\rqSet{rqRobFFNoStderrLibAllCTwoRustNum}{0}
\rqSet{rqRobFFNoStderrLibAllCTwoRustDen}{0}
\rqSet{rqRobFFNoStderrLibAllCTwoRustPct}{n/a}
\rqSet{rqRobFFNoStderrLibAllLaertesNum}{0}
\rqSet{rqRobFFNoStderrLibAllLaertesDen}{0}
\rqSet{rqRobFFNoStderrLibAllLaertesPct}{n/a}
\rqSet{rqRobFFNoStderrLibAllCTwoSaferRustNum}{0}
\rqSet{rqRobFFNoStderrLibAllCTwoSaferRustDen}{23}
\rqSet{rqRobFFNoStderrLibAllCTwoSaferRustPct}{0.0}
\rqSet{rqRobFFNoStderrLibAllSmartCTwoRustNum}{0}
\rqSet{rqRobFFNoStderrLibAllSmartCTwoRustDen}{0}
\rqSet{rqRobFFNoStderrLibAllSmartCTwoRustPct}{n/a}
\rqSet{rqRobFFNoStderrLibAllKimiNum}{0}
\rqSet{rqRobFFNoStderrLibAllKimiDen}{15}
\rqSet{rqRobFFNoStderrLibAllKimiPct}{0.0}
\rqSet{rqRobFFNoStderrLibAllGptFiveNum}{1}
\rqSet{rqRobFFNoStderrLibAllGptFiveDen}{13}
\rqSet{rqRobFFNoStderrLibAllGptFivePct}{7.7}
\rqSet{rqRobFFNoStderrLibAllGeminiNum}{0}
\rqSet{rqRobFFNoStderrLibAllGeminiDen}{4}
\rqSet{rqRobFFNoStderrLibAllGeminiPct}{0.0}
\rqSet{rqRobFFNoStderrLibBOneNum}{0}
\rqSet{rqRobFFNoStderrLibBOneDen}{31}
\rqSet{rqRobFFNoStderrLibBOnePct}{0.0}
\rqSet{rqRobFFNoStderrLibBTwoNum}{1}
\rqSet{rqRobFFNoStderrLibBTwoDen}{24}
\rqSet{rqRobFFNoStderrLibBTwoPct}{4.2}
\rqSet{rqRobFFNoStderrLibPZeroNum}{0}
\rqSet{rqRobFFNoStderrLibPZeroDen}{0}
\rqSet{rqRobFFNoStderrLibPZeroPct}{n/a}
\rqSet{rqRobFFNoStderrLibPOneNum}{0}
\rqSet{rqRobFFNoStderrLibPOneDen}{56}
\rqSet{rqRobFFNoStderrLibPOnePct}{0.0}
\rqSet{rqRobFFNoStderrNotAdjLibAllNum}{0}
\rqSet{rqRobFFNoStderrNotAdjLibAllDen}{111}
\rqSet{rqRobFFNoStderrNotAdjLibAllPct}{0.0}
\rqSet{rqRobFFNoStderrNotAdjLibBOneNum}{0}
\rqSet{rqRobFFNoStderrNotAdjLibBOneDen}{31}
\rqSet{rqRobFFNoStderrNotAdjLibBOnePct}{0.0}
\rqSet{rqRobFFNoStderrNotAdjLibBTwoNum}{0}
\rqSet{rqRobFFNoStderrNotAdjLibBTwoDen}{24}
\rqSet{rqRobFFNoStderrNotAdjLibBTwoPct}{0.0}
\rqSet{rqRobFFNoStderrNotAdjLibPZeroNum}{0}
\rqSet{rqRobFFNoStderrNotAdjLibPZeroDen}{0}
\rqSet{rqRobFFNoStderrNotAdjLibPZeroPct}{n/a}
\rqSet{rqRobFFNoStderrNotAdjLibPOneNum}{0}
\rqSet{rqRobFFNoStderrNotAdjLibPOneDen}{56}
\rqSet{rqRobFFNoStderrNotAdjLibPOnePct}{0.0}
\rqSet{rqRobFFNoStderrNotReexLibAllNum}{0}
\rqSet{rqRobFFNoStderrNotReexLibAllDen}{111}
\rqSet{rqRobFFNoStderrNotReexLibAllPct}{0.0}
\rqSet{rqRobFFNoStderrNotReexLibBOneNum}{0}
\rqSet{rqRobFFNoStderrNotReexLibBOneDen}{31}
\rqSet{rqRobFFNoStderrNotReexLibBOnePct}{0.0}
\rqSet{rqRobFFNoStderrNotReexLibBTwoNum}{0}
\rqSet{rqRobFFNoStderrNotReexLibBTwoDen}{24}
\rqSet{rqRobFFNoStderrNotReexLibBTwoPct}{0.0}
\rqSet{rqRobFFNoStderrNotReexLibPZeroNum}{0}
\rqSet{rqRobFFNoStderrNotReexLibPZeroDen}{0}
\rqSet{rqRobFFNoStderrNotReexLibPZeroPct}{n/a}
\rqSet{rqRobFFNoStderrNotReexLibPOneNum}{0}
\rqSet{rqRobFFNoStderrNotReexLibPOneDen}{56}
\rqSet{rqRobFFNoStderrNotReexLibPOnePct}{0.0}
\rqSet{rqRobFFNoStderrChangedLibAllNum}{0}
\rqSet{rqRobFFNoStderrChangedLibAllDen}{111}
\rqSet{rqRobFFNoStderrChangedLibAllPct}{0.0}
\rqSet{rqRobFFNoStderrChangedLibBOneNum}{0}
\rqSet{rqRobFFNoStderrChangedLibBOneDen}{31}
\rqSet{rqRobFFNoStderrChangedLibBOnePct}{0.0}
\rqSet{rqRobFFNoStderrChangedLibBTwoNum}{0}
\rqSet{rqRobFFNoStderrChangedLibBTwoDen}{24}
\rqSet{rqRobFFNoStderrChangedLibBTwoPct}{0.0}
\rqSet{rqRobFFNoStderrChangedLibPZeroNum}{0}
\rqSet{rqRobFFNoStderrChangedLibPZeroDen}{0}
\rqSet{rqRobFFNoStderrChangedLibPZeroPct}{n/a}
\rqSet{rqRobFFNoStderrChangedLibPOneNum}{0}
\rqSet{rqRobFFNoStderrChangedLibPOneDen}{56}
\rqSet{rqRobFFNoStderrChangedLibPOnePct}{0.0}
\rqSet{rqRobFFExactLibAllNum}{0}
\rqSet{rqRobFFExactLibAllDen}{111}
\rqSet{rqRobFFExactLibAllPct}{0.0}
\rqSet{rqRobFFExactLibAllKiroNum}{0}
\rqSet{rqRobFFExactLibAllKiroDen}{0}
\rqSet{rqRobFFExactLibAllKiroPct}{n/a}
\rqSet{rqRobFFExactLibAllClaudeNum}{0}
\rqSet{rqRobFFExactLibAllClaudeDen}{11}
\rqSet{rqRobFFExactLibAllClaudePct}{0.0}
\rqSet{rqRobFFExactLibAllCodexNum}{0}
\rqSet{rqRobFFExactLibAllCodexDen}{45}
\rqSet{rqRobFFExactLibAllCodexPct}{0.0}
\rqSet{rqRobFFExactLibAllCTwoRustNum}{0}
\rqSet{rqRobFFExactLibAllCTwoRustDen}{0}
\rqSet{rqRobFFExactLibAllCTwoRustPct}{n/a}
\rqSet{rqRobFFExactLibAllLaertesNum}{0}
\rqSet{rqRobFFExactLibAllLaertesDen}{0}
\rqSet{rqRobFFExactLibAllLaertesPct}{n/a}
\rqSet{rqRobFFExactLibAllCTwoSaferRustNum}{0}
\rqSet{rqRobFFExactLibAllCTwoSaferRustDen}{23}
\rqSet{rqRobFFExactLibAllCTwoSaferRustPct}{0.0}
\rqSet{rqRobFFExactLibAllSmartCTwoRustNum}{0}
\rqSet{rqRobFFExactLibAllSmartCTwoRustDen}{0}
\rqSet{rqRobFFExactLibAllSmartCTwoRustPct}{n/a}
\rqSet{rqRobFFExactLibAllKimiNum}{0}
\rqSet{rqRobFFExactLibAllKimiDen}{15}
\rqSet{rqRobFFExactLibAllKimiPct}{0.0}
\rqSet{rqRobFFExactLibAllGptFiveNum}{0}
\rqSet{rqRobFFExactLibAllGptFiveDen}{13}
\rqSet{rqRobFFExactLibAllGptFivePct}{0.0}
\rqSet{rqRobFFExactLibAllGeminiNum}{0}
\rqSet{rqRobFFExactLibAllGeminiDen}{4}
\rqSet{rqRobFFExactLibAllGeminiPct}{0.0}
\rqSet{rqRobFFExactLibBOneNum}{0}
\rqSet{rqRobFFExactLibBOneDen}{31}
\rqSet{rqRobFFExactLibBOnePct}{0.0}
\rqSet{rqRobFFExactLibBTwoNum}{0}
\rqSet{rqRobFFExactLibBTwoDen}{24}
\rqSet{rqRobFFExactLibBTwoPct}{0.0}
\rqSet{rqRobFFExactLibPZeroNum}{0}
\rqSet{rqRobFFExactLibPZeroDen}{0}
\rqSet{rqRobFFExactLibPZeroPct}{n/a}
\rqSet{rqRobFFExactLibPOneNum}{0}
\rqSet{rqRobFFExactLibPOneDen}{56}
\rqSet{rqRobFFExactLibPOnePct}{0.0}
\rqSet{rqRobFFExactNotAdjLibAllNum}{0}
\rqSet{rqRobFFExactNotAdjLibAllDen}{111}
\rqSet{rqRobFFExactNotAdjLibAllPct}{0.0}
\rqSet{rqRobFFExactNotAdjLibBOneNum}{0}
\rqSet{rqRobFFExactNotAdjLibBOneDen}{31}
\rqSet{rqRobFFExactNotAdjLibBOnePct}{0.0}
\rqSet{rqRobFFExactNotAdjLibBTwoNum}{0}
\rqSet{rqRobFFExactNotAdjLibBTwoDen}{24}
\rqSet{rqRobFFExactNotAdjLibBTwoPct}{0.0}
\rqSet{rqRobFFExactNotAdjLibPZeroNum}{0}
\rqSet{rqRobFFExactNotAdjLibPZeroDen}{0}
\rqSet{rqRobFFExactNotAdjLibPZeroPct}{n/a}
\rqSet{rqRobFFExactNotAdjLibPOneNum}{0}
\rqSet{rqRobFFExactNotAdjLibPOneDen}{56}
\rqSet{rqRobFFExactNotAdjLibPOnePct}{0.0}
\rqSet{rqRobFFExactNotReexLibAllNum}{0}
\rqSet{rqRobFFExactNotReexLibAllDen}{111}
\rqSet{rqRobFFExactNotReexLibAllPct}{0.0}
\rqSet{rqRobFFExactNotReexLibBOneNum}{0}
\rqSet{rqRobFFExactNotReexLibBOneDen}{31}
\rqSet{rqRobFFExactNotReexLibBOnePct}{0.0}
\rqSet{rqRobFFExactNotReexLibBTwoNum}{0}
\rqSet{rqRobFFExactNotReexLibBTwoDen}{24}
\rqSet{rqRobFFExactNotReexLibBTwoPct}{0.0}
\rqSet{rqRobFFExactNotReexLibPZeroNum}{0}
\rqSet{rqRobFFExactNotReexLibPZeroDen}{0}
\rqSet{rqRobFFExactNotReexLibPZeroPct}{n/a}
\rqSet{rqRobFFExactNotReexLibPOneNum}{0}
\rqSet{rqRobFFExactNotReexLibPOneDen}{56}
\rqSet{rqRobFFExactNotReexLibPOnePct}{0.0}
\rqSet{rqRobFFExactChangedLibAllNum}{1}
\rqSet{rqRobFFExactChangedLibAllDen}{111}
\rqSet{rqRobFFExactChangedLibAllPct}{0.9}
\rqSet{rqRobFFExactChangedLibBOneNum}{0}
\rqSet{rqRobFFExactChangedLibBOneDen}{31}
\rqSet{rqRobFFExactChangedLibBOnePct}{0.0}
\rqSet{rqRobFFExactChangedLibBTwoNum}{1}
\rqSet{rqRobFFExactChangedLibBTwoDen}{24}
\rqSet{rqRobFFExactChangedLibBTwoPct}{4.2}
\rqSet{rqRobFFExactChangedLibPZeroNum}{0}
\rqSet{rqRobFFExactChangedLibPZeroDen}{0}
\rqSet{rqRobFFExactChangedLibPZeroPct}{n/a}
\rqSet{rqRobFFExactChangedLibPOneNum}{0}
\rqSet{rqRobFFExactChangedLibPOneDen}{56}
\rqSet{rqRobFFExactChangedLibPOnePct}{0.0}
\rqSet{rqRobFFOrigFloatLibAllNum}{0}
\rqSet{rqRobFFOrigFloatLibAllDen}{111}
\rqSet{rqRobFFOrigFloatLibAllPct}{0.0}
\rqSet{rqRobFFOrigFloatLibAllKiroNum}{0}
\rqSet{rqRobFFOrigFloatLibAllKiroDen}{0}
\rqSet{rqRobFFOrigFloatLibAllKiroPct}{n/a}
\rqSet{rqRobFFOrigFloatLibAllClaudeNum}{0}
\rqSet{rqRobFFOrigFloatLibAllClaudeDen}{11}
\rqSet{rqRobFFOrigFloatLibAllClaudePct}{0.0}
\rqSet{rqRobFFOrigFloatLibAllCodexNum}{0}
\rqSet{rqRobFFOrigFloatLibAllCodexDen}{45}
\rqSet{rqRobFFOrigFloatLibAllCodexPct}{0.0}
\rqSet{rqRobFFOrigFloatLibAllCTwoRustNum}{0}
\rqSet{rqRobFFOrigFloatLibAllCTwoRustDen}{0}
\rqSet{rqRobFFOrigFloatLibAllCTwoRustPct}{n/a}
\rqSet{rqRobFFOrigFloatLibAllLaertesNum}{0}
\rqSet{rqRobFFOrigFloatLibAllLaertesDen}{0}
\rqSet{rqRobFFOrigFloatLibAllLaertesPct}{n/a}
\rqSet{rqRobFFOrigFloatLibAllCTwoSaferRustNum}{0}
\rqSet{rqRobFFOrigFloatLibAllCTwoSaferRustDen}{23}
\rqSet{rqRobFFOrigFloatLibAllCTwoSaferRustPct}{0.0}
\rqSet{rqRobFFOrigFloatLibAllSmartCTwoRustNum}{0}
\rqSet{rqRobFFOrigFloatLibAllSmartCTwoRustDen}{0}
\rqSet{rqRobFFOrigFloatLibAllSmartCTwoRustPct}{n/a}
\rqSet{rqRobFFOrigFloatLibAllKimiNum}{0}
\rqSet{rqRobFFOrigFloatLibAllKimiDen}{15}
\rqSet{rqRobFFOrigFloatLibAllKimiPct}{0.0}
\rqSet{rqRobFFOrigFloatLibAllGptFiveNum}{0}
\rqSet{rqRobFFOrigFloatLibAllGptFiveDen}{13}
\rqSet{rqRobFFOrigFloatLibAllGptFivePct}{0.0}
\rqSet{rqRobFFOrigFloatLibAllGeminiNum}{0}
\rqSet{rqRobFFOrigFloatLibAllGeminiDen}{4}
\rqSet{rqRobFFOrigFloatLibAllGeminiPct}{0.0}
\rqSet{rqRobFFOrigFloatLibBOneNum}{0}
\rqSet{rqRobFFOrigFloatLibBOneDen}{31}
\rqSet{rqRobFFOrigFloatLibBOnePct}{0.0}
\rqSet{rqRobFFOrigFloatLibBTwoNum}{0}
\rqSet{rqRobFFOrigFloatLibBTwoDen}{24}
\rqSet{rqRobFFOrigFloatLibBTwoPct}{0.0}
\rqSet{rqRobFFOrigFloatLibPZeroNum}{0}
\rqSet{rqRobFFOrigFloatLibPZeroDen}{0}
\rqSet{rqRobFFOrigFloatLibPZeroPct}{n/a}
\rqSet{rqRobFFOrigFloatLibPOneNum}{0}
\rqSet{rqRobFFOrigFloatLibPOneDen}{56}
\rqSet{rqRobFFOrigFloatLibPOnePct}{0.0}
\rqSet{rqRobFFOrigFloatNotAdjLibAllNum}{0}
\rqSet{rqRobFFOrigFloatNotAdjLibAllDen}{111}
\rqSet{rqRobFFOrigFloatNotAdjLibAllPct}{0.0}
\rqSet{rqRobFFOrigFloatNotAdjLibBOneNum}{0}
\rqSet{rqRobFFOrigFloatNotAdjLibBOneDen}{31}
\rqSet{rqRobFFOrigFloatNotAdjLibBOnePct}{0.0}
\rqSet{rqRobFFOrigFloatNotAdjLibBTwoNum}{0}
\rqSet{rqRobFFOrigFloatNotAdjLibBTwoDen}{24}
\rqSet{rqRobFFOrigFloatNotAdjLibBTwoPct}{0.0}
\rqSet{rqRobFFOrigFloatNotAdjLibPZeroNum}{0}
\rqSet{rqRobFFOrigFloatNotAdjLibPZeroDen}{0}
\rqSet{rqRobFFOrigFloatNotAdjLibPZeroPct}{n/a}
\rqSet{rqRobFFOrigFloatNotAdjLibPOneNum}{0}
\rqSet{rqRobFFOrigFloatNotAdjLibPOneDen}{56}
\rqSet{rqRobFFOrigFloatNotAdjLibPOnePct}{0.0}
\rqSet{rqRobFFOrigFloatNotReexLibAllNum}{0}
\rqSet{rqRobFFOrigFloatNotReexLibAllDen}{111}
\rqSet{rqRobFFOrigFloatNotReexLibAllPct}{0.0}
\rqSet{rqRobFFOrigFloatNotReexLibBOneNum}{0}
\rqSet{rqRobFFOrigFloatNotReexLibBOneDen}{31}
\rqSet{rqRobFFOrigFloatNotReexLibBOnePct}{0.0}
\rqSet{rqRobFFOrigFloatNotReexLibBTwoNum}{0}
\rqSet{rqRobFFOrigFloatNotReexLibBTwoDen}{24}
\rqSet{rqRobFFOrigFloatNotReexLibBTwoPct}{0.0}
\rqSet{rqRobFFOrigFloatNotReexLibPZeroNum}{0}
\rqSet{rqRobFFOrigFloatNotReexLibPZeroDen}{0}
\rqSet{rqRobFFOrigFloatNotReexLibPZeroPct}{n/a}
\rqSet{rqRobFFOrigFloatNotReexLibPOneNum}{0}
\rqSet{rqRobFFOrigFloatNotReexLibPOneDen}{56}
\rqSet{rqRobFFOrigFloatNotReexLibPOnePct}{0.0}
\rqSet{rqRobFFOrigFloatChangedLibAllNum}{1}
\rqSet{rqRobFFOrigFloatChangedLibAllDen}{111}
\rqSet{rqRobFFOrigFloatChangedLibAllPct}{0.9}
\rqSet{rqRobFFOrigFloatChangedLibBOneNum}{0}
\rqSet{rqRobFFOrigFloatChangedLibBOneDen}{31}
\rqSet{rqRobFFOrigFloatChangedLibBOnePct}{0.0}
\rqSet{rqRobFFOrigFloatChangedLibBTwoNum}{1}
\rqSet{rqRobFFOrigFloatChangedLibBTwoDen}{24}
\rqSet{rqRobFFOrigFloatChangedLibBTwoPct}{4.2}
\rqSet{rqRobFFOrigFloatChangedLibPZeroNum}{0}
\rqSet{rqRobFFOrigFloatChangedLibPZeroDen}{0}
\rqSet{rqRobFFOrigFloatChangedLibPZeroPct}{n/a}
\rqSet{rqRobFFOrigFloatChangedLibPOneNum}{0}
\rqSet{rqRobFFOrigFloatChangedLibPOneDen}{56}
\rqSet{rqRobFFOrigFloatChangedLibPOnePct}{0.0}
\rqSet{rqSRLibAllNum}{1281}
\rqSet{rqSRLibAllDen}{2190}
\rqSet{rqSRLibAllPct}{58.5}
\rqSet{rqSRLibAllKiroNum}{217}
\rqSet{rqSRLibAllKiroDen}{219}
\rqSet{rqSRLibAllKiroPct}{99.1}
\rqSet{rqSRLibAllClaudeNum}{206}
\rqSet{rqSRLibAllClaudeDen}{219}
\rqSet{rqSRLibAllClaudePct}{94.1}
\rqSet{rqSRLibAllCodexNum}{171}
\rqSet{rqSRLibAllCodexDen}{219}
\rqSet{rqSRLibAllCodexPct}{78.1}
\rqSet{rqSRLibAllCTwoRustNum}{141}
\rqSet{rqSRLibAllCTwoRustDen}{219}
\rqSet{rqSRLibAllCTwoRustPct}{64.4}
\rqSet{rqSRLibAllLaertesNum}{137}
\rqSet{rqSRLibAllLaertesDen}{219}
\rqSet{rqSRLibAllLaertesPct}{62.6}
\rqSet{rqSRLibAllCTwoSaferRustNum}{105}
\rqSet{rqSRLibAllCTwoSaferRustDen}{219}
\rqSet{rqSRLibAllCTwoSaferRustPct}{47.9}
\rqSet{rqSRLibAllSmartCTwoRustNum}{0}
\rqSet{rqSRLibAllSmartCTwoRustDen}{219}
\rqSet{rqSRLibAllSmartCTwoRustPct}{0.0}
\rqSet{rqSRLibAllKimiNum}{85}
\rqSet{rqSRLibAllKimiDen}{219}
\rqSet{rqSRLibAllKimiPct}{38.8}
\rqSet{rqSRLibAllGptFiveNum}{111}
\rqSet{rqSRLibAllGptFiveDen}{219}
\rqSet{rqSRLibAllGptFivePct}{50.7}
\rqSet{rqSRLibAllGeminiNum}{108}
\rqSet{rqSRLibAllGeminiDen}{219}
\rqSet{rqSRLibAllGeminiPct}{49.3}
\rqSet{rqSRLibBOneNum}{637}
\rqSet{rqSRLibBOneDen}{770}
\rqSet{rqSRLibBOnePct}{82.7}
\rqSet{rqSRLibBOneKiroNum}{77}
\rqSet{rqSRLibBOneKiroDen}{77}
\rqSet{rqSRLibBOneKiroPct}{100.0}
\rqSet{rqSRLibBOneClaudeNum}{77}
\rqSet{rqSRLibBOneClaudeDen}{77}
\rqSet{rqSRLibBOneClaudePct}{100.0}
\rqSet{rqSRLibBOneCodexNum}{77}
\rqSet{rqSRLibBOneCodexDen}{77}
\rqSet{rqSRLibBOneCodexPct}{100.0}
\rqSet{rqSRLibBOneCTwoRustNum}{77}
\rqSet{rqSRLibBOneCTwoRustDen}{77}
\rqSet{rqSRLibBOneCTwoRustPct}{100.0}
\rqSet{rqSRLibBOneLaertesNum}{74}
\rqSet{rqSRLibBOneLaertesDen}{77}
\rqSet{rqSRLibBOneLaertesPct}{96.1}
\rqSet{rqSRLibBOneCTwoSaferRustNum}{60}
\rqSet{rqSRLibBOneCTwoSaferRustDen}{77}
\rqSet{rqSRLibBOneCTwoSaferRustPct}{77.9}
\rqSet{rqSRLibBOneSmartCTwoRustNum}{0}
\rqSet{rqSRLibBOneSmartCTwoRustDen}{77}
\rqSet{rqSRLibBOneSmartCTwoRustPct}{0.0}
\rqSet{rqSRLibBOneKimiNum}{57}
\rqSet{rqSRLibBOneKimiDen}{77}
\rqSet{rqSRLibBOneKimiPct}{74.0}
\rqSet{rqSRLibBOneGptFiveNum}{67}
\rqSet{rqSRLibBOneGptFiveDen}{77}
\rqSet{rqSRLibBOneGptFivePct}{87.0}
\rqSet{rqSRLibBOneGeminiNum}{71}
\rqSet{rqSRLibBOneGeminiDen}{77}
\rqSet{rqSRLibBOneGeminiPct}{92.2}
\rqSet{rqSRLibBTwoNum}{472}
\rqSet{rqSRLibBTwoDen}{660}
\rqSet{rqSRLibBTwoPct}{71.5}
\rqSet{rqSRLibBTwoKiroNum}{64}
\rqSet{rqSRLibBTwoKiroDen}{66}
\rqSet{rqSRLibBTwoKiroPct}{97.0}
\rqSet{rqSRLibBTwoClaudeNum}{64}
\rqSet{rqSRLibBTwoClaudeDen}{66}
\rqSet{rqSRLibBTwoClaudePct}{97.0}
\rqSet{rqSRLibBTwoCodexNum}{63}
\rqSet{rqSRLibBTwoCodexDen}{66}
\rqSet{rqSRLibBTwoCodexPct}{95.5}
\rqSet{rqSRLibBTwoCTwoRustNum}{64}
\rqSet{rqSRLibBTwoCTwoRustDen}{66}
\rqSet{rqSRLibBTwoCTwoRustPct}{97.0}
\rqSet{rqSRLibBTwoLaertesNum}{63}
\rqSet{rqSRLibBTwoLaertesDen}{66}
\rqSet{rqSRLibBTwoLaertesPct}{95.5}
\rqSet{rqSRLibBTwoCTwoSaferRustNum}{45}
\rqSet{rqSRLibBTwoCTwoSaferRustDen}{66}
\rqSet{rqSRLibBTwoCTwoSaferRustPct}{68.2}
\rqSet{rqSRLibBTwoSmartCTwoRustNum}{0}
\rqSet{rqSRLibBTwoSmartCTwoRustDen}{66}
\rqSet{rqSRLibBTwoSmartCTwoRustPct}{0.0}
\rqSet{rqSRLibBTwoKimiNum}{28}
\rqSet{rqSRLibBTwoKimiDen}{66}
\rqSet{rqSRLibBTwoKimiPct}{42.4}
\rqSet{rqSRLibBTwoGptFiveNum}{44}
\rqSet{rqSRLibBTwoGptFiveDen}{66}
\rqSet{rqSRLibBTwoGptFivePct}{66.7}
\rqSet{rqSRLibBTwoGeminiNum}{37}
\rqSet{rqSRLibBTwoGeminiDen}{66}
\rqSet{rqSRLibBTwoGeminiPct}{56.1}
\rqSet{rqSRLibPZeroNum}{0}
\rqSet{rqSRLibPZeroDen}{0}
\rqSet{rqSRLibPZeroPct}{n/a}
\rqSet{rqSRLibPZeroKiroNum}{0}
\rqSet{rqSRLibPZeroKiroDen}{0}
\rqSet{rqSRLibPZeroKiroPct}{n/a}
\rqSet{rqSRLibPZeroClaudeNum}{0}
\rqSet{rqSRLibPZeroClaudeDen}{0}
\rqSet{rqSRLibPZeroClaudePct}{n/a}
\rqSet{rqSRLibPZeroCodexNum}{0}
\rqSet{rqSRLibPZeroCodexDen}{0}
\rqSet{rqSRLibPZeroCodexPct}{n/a}
\rqSet{rqSRLibPZeroCTwoRustNum}{0}
\rqSet{rqSRLibPZeroCTwoRustDen}{0}
\rqSet{rqSRLibPZeroCTwoRustPct}{n/a}
\rqSet{rqSRLibPZeroLaertesNum}{0}
\rqSet{rqSRLibPZeroLaertesDen}{0}
\rqSet{rqSRLibPZeroLaertesPct}{n/a}
\rqSet{rqSRLibPZeroCTwoSaferRustNum}{0}
\rqSet{rqSRLibPZeroCTwoSaferRustDen}{0}
\rqSet{rqSRLibPZeroCTwoSaferRustPct}{n/a}
\rqSet{rqSRLibPZeroSmartCTwoRustNum}{0}
\rqSet{rqSRLibPZeroSmartCTwoRustDen}{0}
\rqSet{rqSRLibPZeroSmartCTwoRustPct}{n/a}
\rqSet{rqSRLibPZeroKimiNum}{0}
\rqSet{rqSRLibPZeroKimiDen}{0}
\rqSet{rqSRLibPZeroKimiPct}{n/a}
\rqSet{rqSRLibPZeroGptFiveNum}{0}
\rqSet{rqSRLibPZeroGptFiveDen}{0}
\rqSet{rqSRLibPZeroGptFivePct}{n/a}
\rqSet{rqSRLibPZeroGeminiNum}{0}
\rqSet{rqSRLibPZeroGeminiDen}{0}
\rqSet{rqSRLibPZeroGeminiPct}{n/a}
\rqSet{rqSRLibPOneNum}{172}
\rqSet{rqSRLibPOneDen}{760}
\rqSet{rqSRLibPOnePct}{22.6}
\rqSet{rqSRLibPOneKiroNum}{76}
\rqSet{rqSRLibPOneKiroDen}{76}
\rqSet{rqSRLibPOneKiroPct}{100.0}
\rqSet{rqSRLibPOneClaudeNum}{65}
\rqSet{rqSRLibPOneClaudeDen}{76}
\rqSet{rqSRLibPOneClaudePct}{85.5}
\rqSet{rqSRLibPOneCodexNum}{31}
\rqSet{rqSRLibPOneCodexDen}{76}
\rqSet{rqSRLibPOneCodexPct}{40.8}
\rqSet{rqSRLibPOneCTwoRustNum}{0}
\rqSet{rqSRLibPOneCTwoRustDen}{76}
\rqSet{rqSRLibPOneCTwoRustPct}{0.0}
\rqSet{rqSRLibPOneLaertesNum}{0}
\rqSet{rqSRLibPOneLaertesDen}{76}
\rqSet{rqSRLibPOneLaertesPct}{0.0}
\rqSet{rqSRLibPOneCTwoSaferRustNum}{0}
\rqSet{rqSRLibPOneCTwoSaferRustDen}{76}
\rqSet{rqSRLibPOneCTwoSaferRustPct}{0.0}
\rqSet{rqSRLibPOneSmartCTwoRustNum}{0}
\rqSet{rqSRLibPOneSmartCTwoRustDen}{76}
\rqSet{rqSRLibPOneSmartCTwoRustPct}{0.0}
\rqSet{rqSRLibPOneKimiNum}{0}
\rqSet{rqSRLibPOneKimiDen}{76}
\rqSet{rqSRLibPOneKimiPct}{0.0}
\rqSet{rqSRLibPOneGptFiveNum}{0}
\rqSet{rqSRLibPOneGptFiveDen}{76}
\rqSet{rqSRLibPOneGptFivePct}{0.0}
\rqSet{rqSRLibPOneGeminiNum}{0}
\rqSet{rqSRLibPOneGeminiDen}{76}
\rqSet{rqSRLibPOneGeminiPct}{0.0}
\rqSet{rqSRLoLibAllNum}{1024}
\rqSet{rqSRLoLibAllDen}{2190}
\rqSet{rqSRLoLibAllPct}{46.8}
\rqSet{rqSRLoLibAllKiroNum}{185}
\rqSet{rqSRLoLibAllKiroDen}{219}
\rqSet{rqSRLoLibAllKiroPct}{84.5}
\rqSet{rqSRLoLibAllClaudeNum}{175}
\rqSet{rqSRLoLibAllClaudeDen}{219}
\rqSet{rqSRLoLibAllClaudePct}{79.9}
\rqSet{rqSRLoLibAllCodexNum}{141}
\rqSet{rqSRLoLibAllCodexDen}{219}
\rqSet{rqSRLoLibAllCodexPct}{64.4}
\rqSet{rqSRLoLibAllCTwoRustNum}{113}
\rqSet{rqSRLoLibAllCTwoRustDen}{219}
\rqSet{rqSRLoLibAllCTwoRustPct}{51.6}
\rqSet{rqSRLoLibAllLaertesNum}{111}
\rqSet{rqSRLoLibAllLaertesDen}{219}
\rqSet{rqSRLoLibAllLaertesPct}{50.7}
\rqSet{rqSRLoLibAllCTwoSaferRustNum}{80}
\rqSet{rqSRLoLibAllCTwoSaferRustDen}{219}
\rqSet{rqSRLoLibAllCTwoSaferRustPct}{36.5}
\rqSet{rqSRLoLibAllSmartCTwoRustNum}{0}
\rqSet{rqSRLoLibAllSmartCTwoRustDen}{219}
\rqSet{rqSRLoLibAllSmartCTwoRustPct}{0.0}
\rqSet{rqSRLoLibAllKimiNum}{57}
\rqSet{rqSRLoLibAllKimiDen}{219}
\rqSet{rqSRLoLibAllKimiPct}{26.0}
\rqSet{rqSRLoLibAllGptFiveNum}{78}
\rqSet{rqSRLoLibAllGptFiveDen}{219}
\rqSet{rqSRLoLibAllGptFivePct}{35.6}
\rqSet{rqSRLoLibAllGeminiNum}{84}
\rqSet{rqSRLoLibAllGeminiDen}{219}
\rqSet{rqSRLoLibAllGeminiPct}{38.4}
\rqSet{rqSRLoLibBOneNum}{502}
\rqSet{rqSRLoLibBOneDen}{770}
\rqSet{rqSRLoLibBOnePct}{65.2}
\rqSet{rqSRLoLibBOneKiroNum}{63}
\rqSet{rqSRLoLibBOneKiroDen}{77}
\rqSet{rqSRLoLibBOneKiroPct}{81.8}
\rqSet{rqSRLoLibBOneClaudeNum}{63}
\rqSet{rqSRLoLibBOneClaudeDen}{77}
\rqSet{rqSRLoLibBOneClaudePct}{81.8}
\rqSet{rqSRLoLibBOneCodexNum}{63}
\rqSet{rqSRLoLibBOneCodexDen}{77}
\rqSet{rqSRLoLibBOneCodexPct}{81.8}
\rqSet{rqSRLoLibBOneCTwoRustNum}{62}
\rqSet{rqSRLoLibBOneCTwoRustDen}{77}
\rqSet{rqSRLoLibBOneCTwoRustPct}{80.5}
\rqSet{rqSRLoLibBOneLaertesNum}{60}
\rqSet{rqSRLoLibBOneLaertesDen}{77}
\rqSet{rqSRLoLibBOneLaertesPct}{77.9}
\rqSet{rqSRLoLibBOneCTwoSaferRustNum}{45}
\rqSet{rqSRLoLibBOneCTwoSaferRustDen}{77}
\rqSet{rqSRLoLibBOneCTwoSaferRustPct}{58.4}
\rqSet{rqSRLoLibBOneSmartCTwoRustNum}{0}
\rqSet{rqSRLoLibBOneSmartCTwoRustDen}{77}
\rqSet{rqSRLoLibBOneSmartCTwoRustPct}{0.0}
\rqSet{rqSRLoLibBOneKimiNum}{40}
\rqSet{rqSRLoLibBOneKimiDen}{77}
\rqSet{rqSRLoLibBOneKimiPct}{51.9}
\rqSet{rqSRLoLibBOneGptFiveNum}{48}
\rqSet{rqSRLoLibBOneGptFiveDen}{77}
\rqSet{rqSRLoLibBOneGptFivePct}{62.3}
\rqSet{rqSRLoLibBOneGeminiNum}{58}
\rqSet{rqSRLoLibBOneGeminiDen}{77}
\rqSet{rqSRLoLibBOneGeminiPct}{75.3}
\rqSet{rqSRLoLibBTwoNum}{353}
\rqSet{rqSRLoLibBTwoDen}{660}
\rqSet{rqSRLoLibBTwoPct}{53.5}
\rqSet{rqSRLoLibBTwoKiroNum}{46}
\rqSet{rqSRLoLibBTwoKiroDen}{66}
\rqSet{rqSRLoLibBTwoKiroPct}{69.7}
\rqSet{rqSRLoLibBTwoClaudeNum}{47}
\rqSet{rqSRLoLibBTwoClaudeDen}{66}
\rqSet{rqSRLoLibBTwoClaudePct}{71.2}
\rqSet{rqSRLoLibBTwoCodexNum}{50}
\rqSet{rqSRLoLibBTwoCodexDen}{66}
\rqSet{rqSRLoLibBTwoCodexPct}{75.8}
\rqSet{rqSRLoLibBTwoCTwoRustNum}{51}
\rqSet{rqSRLoLibBTwoCTwoRustDen}{66}
\rqSet{rqSRLoLibBTwoCTwoRustPct}{77.3}
\rqSet{rqSRLoLibBTwoLaertesNum}{51}
\rqSet{rqSRLoLibBTwoLaertesDen}{66}
\rqSet{rqSRLoLibBTwoLaertesPct}{77.3}
\rqSet{rqSRLoLibBTwoCTwoSaferRustNum}{35}
\rqSet{rqSRLoLibBTwoCTwoSaferRustDen}{66}
\rqSet{rqSRLoLibBTwoCTwoSaferRustPct}{53.0}
\rqSet{rqSRLoLibBTwoSmartCTwoRustNum}{0}
\rqSet{rqSRLoLibBTwoSmartCTwoRustDen}{66}
\rqSet{rqSRLoLibBTwoSmartCTwoRustPct}{0.0}
\rqSet{rqSRLoLibBTwoKimiNum}{17}
\rqSet{rqSRLoLibBTwoKimiDen}{66}
\rqSet{rqSRLoLibBTwoKimiPct}{25.8}
\rqSet{rqSRLoLibBTwoGptFiveNum}{30}
\rqSet{rqSRLoLibBTwoGptFiveDen}{66}
\rqSet{rqSRLoLibBTwoGptFivePct}{45.5}
\rqSet{rqSRLoLibBTwoGeminiNum}{26}
\rqSet{rqSRLoLibBTwoGeminiDen}{66}
\rqSet{rqSRLoLibBTwoGeminiPct}{39.4}
\rqSet{rqSRLoLibPZeroNum}{0}
\rqSet{rqSRLoLibPZeroDen}{0}
\rqSet{rqSRLoLibPZeroPct}{n/a}
\rqSet{rqSRLoLibPZeroKiroNum}{0}
\rqSet{rqSRLoLibPZeroKiroDen}{0}
\rqSet{rqSRLoLibPZeroKiroPct}{n/a}
\rqSet{rqSRLoLibPZeroClaudeNum}{0}
\rqSet{rqSRLoLibPZeroClaudeDen}{0}
\rqSet{rqSRLoLibPZeroClaudePct}{n/a}
\rqSet{rqSRLoLibPZeroCodexNum}{0}
\rqSet{rqSRLoLibPZeroCodexDen}{0}
\rqSet{rqSRLoLibPZeroCodexPct}{n/a}
\rqSet{rqSRLoLibPZeroCTwoRustNum}{0}
\rqSet{rqSRLoLibPZeroCTwoRustDen}{0}
\rqSet{rqSRLoLibPZeroCTwoRustPct}{n/a}
\rqSet{rqSRLoLibPZeroLaertesNum}{0}
\rqSet{rqSRLoLibPZeroLaertesDen}{0}
\rqSet{rqSRLoLibPZeroLaertesPct}{n/a}
\rqSet{rqSRLoLibPZeroCTwoSaferRustNum}{0}
\rqSet{rqSRLoLibPZeroCTwoSaferRustDen}{0}
\rqSet{rqSRLoLibPZeroCTwoSaferRustPct}{n/a}
\rqSet{rqSRLoLibPZeroSmartCTwoRustNum}{0}
\rqSet{rqSRLoLibPZeroSmartCTwoRustDen}{0}
\rqSet{rqSRLoLibPZeroSmartCTwoRustPct}{n/a}
\rqSet{rqSRLoLibPZeroKimiNum}{0}
\rqSet{rqSRLoLibPZeroKimiDen}{0}
\rqSet{rqSRLoLibPZeroKimiPct}{n/a}
\rqSet{rqSRLoLibPZeroGptFiveNum}{0}
\rqSet{rqSRLoLibPZeroGptFiveDen}{0}
\rqSet{rqSRLoLibPZeroGptFivePct}{n/a}
\rqSet{rqSRLoLibPZeroGeminiNum}{0}
\rqSet{rqSRLoLibPZeroGeminiDen}{0}
\rqSet{rqSRLoLibPZeroGeminiPct}{n/a}
\rqSet{rqSRLoLibPOneNum}{169}
\rqSet{rqSRLoLibPOneDen}{760}
\rqSet{rqSRLoLibPOnePct}{22.2}
\rqSet{rqSRLoLibPOneKiroNum}{76}
\rqSet{rqSRLoLibPOneKiroDen}{76}
\rqSet{rqSRLoLibPOneKiroPct}{100.0}
\rqSet{rqSRLoLibPOneClaudeNum}{65}
\rqSet{rqSRLoLibPOneClaudeDen}{76}
\rqSet{rqSRLoLibPOneClaudePct}{85.5}
\rqSet{rqSRLoLibPOneCodexNum}{28}
\rqSet{rqSRLoLibPOneCodexDen}{76}
\rqSet{rqSRLoLibPOneCodexPct}{36.8}
\rqSet{rqSRLoLibPOneCTwoRustNum}{0}
\rqSet{rqSRLoLibPOneCTwoRustDen}{76}
\rqSet{rqSRLoLibPOneCTwoRustPct}{0.0}
\rqSet{rqSRLoLibPOneLaertesNum}{0}
\rqSet{rqSRLoLibPOneLaertesDen}{76}
\rqSet{rqSRLoLibPOneLaertesPct}{0.0}
\rqSet{rqSRLoLibPOneCTwoSaferRustNum}{0}
\rqSet{rqSRLoLibPOneCTwoSaferRustDen}{76}
\rqSet{rqSRLoLibPOneCTwoSaferRustPct}{0.0}
\rqSet{rqSRLoLibPOneSmartCTwoRustNum}{0}
\rqSet{rqSRLoLibPOneSmartCTwoRustDen}{76}
\rqSet{rqSRLoLibPOneSmartCTwoRustPct}{0.0}
\rqSet{rqSRLoLibPOneKimiNum}{0}
\rqSet{rqSRLoLibPOneKimiDen}{76}
\rqSet{rqSRLoLibPOneKimiPct}{0.0}
\rqSet{rqSRLoLibPOneGptFiveNum}{0}
\rqSet{rqSRLoLibPOneGptFiveDen}{76}
\rqSet{rqSRLoLibPOneGptFivePct}{0.0}
\rqSet{rqSRLoLibPOneGeminiNum}{0}
\rqSet{rqSRLoLibPOneGeminiDen}{76}
\rqSet{rqSRLoLibPOneGeminiPct}{0.0}
\rqSet{rqSRHiLibAllNum}{1244}
\rqSet{rqSRHiLibAllDen}{2190}
\rqSet{rqSRHiLibAllPct}{56.8}
\rqSet{rqSRHiLibAllKiroNum}{212}
\rqSet{rqSRHiLibAllKiroDen}{219}
\rqSet{rqSRHiLibAllKiroPct}{96.8}
\rqSet{rqSRHiLibAllClaudeNum}{202}
\rqSet{rqSRHiLibAllClaudeDen}{219}
\rqSet{rqSRHiLibAllClaudePct}{92.2}
\rqSet{rqSRHiLibAllCodexNum}{171}
\rqSet{rqSRHiLibAllCodexDen}{219}
\rqSet{rqSRHiLibAllCodexPct}{78.1}
\rqSet{rqSRHiLibAllCTwoRustNum}{139}
\rqSet{rqSRHiLibAllCTwoRustDen}{219}
\rqSet{rqSRHiLibAllCTwoRustPct}{63.5}
\rqSet{rqSRHiLibAllLaertesNum}{135}
\rqSet{rqSRHiLibAllLaertesDen}{219}
\rqSet{rqSRHiLibAllLaertesPct}{61.6}
\rqSet{rqSRHiLibAllCTwoSaferRustNum}{100}
\rqSet{rqSRHiLibAllCTwoSaferRustDen}{219}
\rqSet{rqSRHiLibAllCTwoSaferRustPct}{45.7}
\rqSet{rqSRHiLibAllSmartCTwoRustNum}{0}
\rqSet{rqSRHiLibAllSmartCTwoRustDen}{219}
\rqSet{rqSRHiLibAllSmartCTwoRustPct}{0.0}
\rqSet{rqSRHiLibAllKimiNum}{75}
\rqSet{rqSRHiLibAllKimiDen}{219}
\rqSet{rqSRHiLibAllKimiPct}{34.2}
\rqSet{rqSRHiLibAllGptFiveNum}{99}
\rqSet{rqSRHiLibAllGptFiveDen}{219}
\rqSet{rqSRHiLibAllGptFivePct}{45.2}
\rqSet{rqSRHiLibAllGeminiNum}{111}
\rqSet{rqSRHiLibAllGeminiDen}{219}
\rqSet{rqSRHiLibAllGeminiPct}{50.7}
\rqSet{rqSRHiLibBOneNum}{606}
\rqSet{rqSRHiLibBOneDen}{770}
\rqSet{rqSRHiLibBOnePct}{78.7}
\rqSet{rqSRHiLibBOneKiroNum}{75}
\rqSet{rqSRHiLibBOneKiroDen}{77}
\rqSet{rqSRHiLibBOneKiroPct}{97.4}
\rqSet{rqSRHiLibBOneClaudeNum}{75}
\rqSet{rqSRHiLibBOneClaudeDen}{77}
\rqSet{rqSRHiLibBOneClaudePct}{97.4}
\rqSet{rqSRHiLibBOneCodexNum}{75}
\rqSet{rqSRHiLibBOneCodexDen}{77}
\rqSet{rqSRHiLibBOneCodexPct}{97.4}
\rqSet{rqSRHiLibBOneCTwoRustNum}{74}
\rqSet{rqSRHiLibBOneCTwoRustDen}{77}
\rqSet{rqSRHiLibBOneCTwoRustPct}{96.1}
\rqSet{rqSRHiLibBOneLaertesNum}{71}
\rqSet{rqSRHiLibBOneLaertesDen}{77}
\rqSet{rqSRHiLibBOneLaertesPct}{92.2}
\rqSet{rqSRHiLibBOneCTwoSaferRustNum}{56}
\rqSet{rqSRHiLibBOneCTwoSaferRustDen}{77}
\rqSet{rqSRHiLibBOneCTwoSaferRustPct}{72.7}
\rqSet{rqSRHiLibBOneSmartCTwoRustNum}{0}
\rqSet{rqSRHiLibBOneSmartCTwoRustDen}{77}
\rqSet{rqSRHiLibBOneSmartCTwoRustPct}{0.0}
\rqSet{rqSRHiLibBOneKimiNum}{50}
\rqSet{rqSRHiLibBOneKimiDen}{77}
\rqSet{rqSRHiLibBOneKimiPct}{64.9}
\rqSet{rqSRHiLibBOneGptFiveNum}{60}
\rqSet{rqSRHiLibBOneGptFiveDen}{77}
\rqSet{rqSRHiLibBOneGptFivePct}{77.9}
\rqSet{rqSRHiLibBOneGeminiNum}{70}
\rqSet{rqSRHiLibBOneGeminiDen}{77}
\rqSet{rqSRHiLibBOneGeminiPct}{90.9}
\rqSet{rqSRHiLibBTwoNum}{466}
\rqSet{rqSRHiLibBTwoDen}{660}
\rqSet{rqSRHiLibBTwoPct}{70.6}
\rqSet{rqSRHiLibBTwoKiroNum}{61}
\rqSet{rqSRHiLibBTwoKiroDen}{66}
\rqSet{rqSRHiLibBTwoKiroPct}{92.4}
\rqSet{rqSRHiLibBTwoClaudeNum}{62}
\rqSet{rqSRHiLibBTwoClaudeDen}{66}
\rqSet{rqSRHiLibBTwoClaudePct}{93.9}
\rqSet{rqSRHiLibBTwoCodexNum}{65}
\rqSet{rqSRHiLibBTwoCodexDen}{66}
\rqSet{rqSRHiLibBTwoCodexPct}{98.5}
\rqSet{rqSRHiLibBTwoCTwoRustNum}{65}
\rqSet{rqSRHiLibBTwoCTwoRustDen}{66}
\rqSet{rqSRHiLibBTwoCTwoRustPct}{98.5}
\rqSet{rqSRHiLibBTwoLaertesNum}{64}
\rqSet{rqSRHiLibBTwoLaertesDen}{66}
\rqSet{rqSRHiLibBTwoLaertesPct}{97.0}
\rqSet{rqSRHiLibBTwoCTwoSaferRustNum}{44}
\rqSet{rqSRHiLibBTwoCTwoSaferRustDen}{66}
\rqSet{rqSRHiLibBTwoCTwoSaferRustPct}{66.7}
\rqSet{rqSRHiLibBTwoSmartCTwoRustNum}{0}
\rqSet{rqSRHiLibBTwoSmartCTwoRustDen}{66}
\rqSet{rqSRHiLibBTwoSmartCTwoRustPct}{0.0}
\rqSet{rqSRHiLibBTwoKimiNum}{25}
\rqSet{rqSRHiLibBTwoKimiDen}{66}
\rqSet{rqSRHiLibBTwoKimiPct}{37.9}
\rqSet{rqSRHiLibBTwoGptFiveNum}{39}
\rqSet{rqSRHiLibBTwoGptFiveDen}{66}
\rqSet{rqSRHiLibBTwoGptFivePct}{59.1}
\rqSet{rqSRHiLibBTwoGeminiNum}{41}
\rqSet{rqSRHiLibBTwoGeminiDen}{66}
\rqSet{rqSRHiLibBTwoGeminiPct}{62.1}
\rqSet{rqSRHiLibPZeroNum}{0}
\rqSet{rqSRHiLibPZeroDen}{0}
\rqSet{rqSRHiLibPZeroPct}{n/a}
\rqSet{rqSRHiLibPZeroKiroNum}{0}
\rqSet{rqSRHiLibPZeroKiroDen}{0}
\rqSet{rqSRHiLibPZeroKiroPct}{n/a}
\rqSet{rqSRHiLibPZeroClaudeNum}{0}
\rqSet{rqSRHiLibPZeroClaudeDen}{0}
\rqSet{rqSRHiLibPZeroClaudePct}{n/a}
\rqSet{rqSRHiLibPZeroCodexNum}{0}
\rqSet{rqSRHiLibPZeroCodexDen}{0}
\rqSet{rqSRHiLibPZeroCodexPct}{n/a}
\rqSet{rqSRHiLibPZeroCTwoRustNum}{0}
\rqSet{rqSRHiLibPZeroCTwoRustDen}{0}
\rqSet{rqSRHiLibPZeroCTwoRustPct}{n/a}
\rqSet{rqSRHiLibPZeroLaertesNum}{0}
\rqSet{rqSRHiLibPZeroLaertesDen}{0}
\rqSet{rqSRHiLibPZeroLaertesPct}{n/a}
\rqSet{rqSRHiLibPZeroCTwoSaferRustNum}{0}
\rqSet{rqSRHiLibPZeroCTwoSaferRustDen}{0}
\rqSet{rqSRHiLibPZeroCTwoSaferRustPct}{n/a}
\rqSet{rqSRHiLibPZeroSmartCTwoRustNum}{0}
\rqSet{rqSRHiLibPZeroSmartCTwoRustDen}{0}
\rqSet{rqSRHiLibPZeroSmartCTwoRustPct}{n/a}
\rqSet{rqSRHiLibPZeroKimiNum}{0}
\rqSet{rqSRHiLibPZeroKimiDen}{0}
\rqSet{rqSRHiLibPZeroKimiPct}{n/a}
\rqSet{rqSRHiLibPZeroGptFiveNum}{0}
\rqSet{rqSRHiLibPZeroGptFiveDen}{0}
\rqSet{rqSRHiLibPZeroGptFivePct}{n/a}
\rqSet{rqSRHiLibPZeroGeminiNum}{0}
\rqSet{rqSRHiLibPZeroGeminiDen}{0}
\rqSet{rqSRHiLibPZeroGeminiPct}{n/a}
\rqSet{rqSRHiLibPOneNum}{172}
\rqSet{rqSRHiLibPOneDen}{760}
\rqSet{rqSRHiLibPOnePct}{22.6}
\rqSet{rqSRHiLibPOneKiroNum}{76}
\rqSet{rqSRHiLibPOneKiroDen}{76}
\rqSet{rqSRHiLibPOneKiroPct}{100.0}
\rqSet{rqSRHiLibPOneClaudeNum}{65}
\rqSet{rqSRHiLibPOneClaudeDen}{76}
\rqSet{rqSRHiLibPOneClaudePct}{85.5}
\rqSet{rqSRHiLibPOneCodexNum}{31}
\rqSet{rqSRHiLibPOneCodexDen}{76}
\rqSet{rqSRHiLibPOneCodexPct}{40.8}
\rqSet{rqSRHiLibPOneCTwoRustNum}{0}
\rqSet{rqSRHiLibPOneCTwoRustDen}{76}
\rqSet{rqSRHiLibPOneCTwoRustPct}{0.0}
\rqSet{rqSRHiLibPOneLaertesNum}{0}
\rqSet{rqSRHiLibPOneLaertesDen}{76}
\rqSet{rqSRHiLibPOneLaertesPct}{0.0}
\rqSet{rqSRHiLibPOneCTwoSaferRustNum}{0}
\rqSet{rqSRHiLibPOneCTwoSaferRustDen}{76}
\rqSet{rqSRHiLibPOneCTwoSaferRustPct}{0.0}
\rqSet{rqSRHiLibPOneSmartCTwoRustNum}{0}
\rqSet{rqSRHiLibPOneSmartCTwoRustDen}{76}
\rqSet{rqSRHiLibPOneSmartCTwoRustPct}{0.0}
\rqSet{rqSRHiLibPOneKimiNum}{0}
\rqSet{rqSRHiLibPOneKimiDen}{76}
\rqSet{rqSRHiLibPOneKimiPct}{0.0}
\rqSet{rqSRHiLibPOneGptFiveNum}{0}
\rqSet{rqSRHiLibPOneGptFiveDen}{76}
\rqSet{rqSRHiLibPOneGptFivePct}{0.0}
\rqSet{rqSRHiLibPOneGeminiNum}{0}
\rqSet{rqSRHiLibPOneGeminiDen}{76}
\rqSet{rqSRHiLibPOneGeminiPct}{0.0}
\rqSet{rqSRWidthLibAllNum}{220}
\rqSet{rqSRWidthLibAllDen}{2190}
\rqSet{rqSRWidthLibAllPct}{10.0}
\rqSet{rqSRWidthLibAllKiroNum}{27}
\rqSet{rqSRWidthLibAllKiroDen}{219}
\rqSet{rqSRWidthLibAllKiroPct}{12.3}
\rqSet{rqSRWidthLibAllClaudeNum}{27}
\rqSet{rqSRWidthLibAllClaudeDen}{219}
\rqSet{rqSRWidthLibAllClaudePct}{12.3}
\rqSet{rqSRWidthLibAllCodexNum}{30}
\rqSet{rqSRWidthLibAllCodexDen}{219}
\rqSet{rqSRWidthLibAllCodexPct}{13.7}
\rqSet{rqSRWidthLibAllCTwoRustNum}{26}
\rqSet{rqSRWidthLibAllCTwoRustDen}{219}
\rqSet{rqSRWidthLibAllCTwoRustPct}{11.9}
\rqSet{rqSRWidthLibAllLaertesNum}{24}
\rqSet{rqSRWidthLibAllLaertesDen}{219}
\rqSet{rqSRWidthLibAllLaertesPct}{11.0}
\rqSet{rqSRWidthLibAllCTwoSaferRustNum}{20}
\rqSet{rqSRWidthLibAllCTwoSaferRustDen}{219}
\rqSet{rqSRWidthLibAllCTwoSaferRustPct}{9.1}
\rqSet{rqSRWidthLibAllSmartCTwoRustNum}{0}
\rqSet{rqSRWidthLibAllSmartCTwoRustDen}{219}
\rqSet{rqSRWidthLibAllSmartCTwoRustPct}{0.0}
\rqSet{rqSRWidthLibAllKimiNum}{18}
\rqSet{rqSRWidthLibAllKimiDen}{219}
\rqSet{rqSRWidthLibAllKimiPct}{8.2}
\rqSet{rqSRWidthLibAllGptFiveNum}{21}
\rqSet{rqSRWidthLibAllGptFiveDen}{219}
\rqSet{rqSRWidthLibAllGptFivePct}{9.6}
\rqSet{rqSRWidthLibAllGeminiNum}{27}
\rqSet{rqSRWidthLibAllGeminiDen}{219}
\rqSet{rqSRWidthLibAllGeminiPct}{12.3}
\rqSet{rqSRWidthLibBOneNum}{104}
\rqSet{rqSRWidthLibBOneDen}{770}
\rqSet{rqSRWidthLibBOnePct}{13.5}
\rqSet{rqSRWidthLibBOneKiroNum}{12}
\rqSet{rqSRWidthLibBOneKiroDen}{77}
\rqSet{rqSRWidthLibBOneKiroPct}{15.6}
\rqSet{rqSRWidthLibBOneClaudeNum}{12}
\rqSet{rqSRWidthLibBOneClaudeDen}{77}
\rqSet{rqSRWidthLibBOneClaudePct}{15.6}
\rqSet{rqSRWidthLibBOneCodexNum}{12}
\rqSet{rqSRWidthLibBOneCodexDen}{77}
\rqSet{rqSRWidthLibBOneCodexPct}{15.6}
\rqSet{rqSRWidthLibBOneCTwoRustNum}{12}
\rqSet{rqSRWidthLibBOneCTwoRustDen}{77}
\rqSet{rqSRWidthLibBOneCTwoRustPct}{15.6}
\rqSet{rqSRWidthLibBOneLaertesNum}{11}
\rqSet{rqSRWidthLibBOneLaertesDen}{77}
\rqSet{rqSRWidthLibBOneLaertesPct}{14.3}
\rqSet{rqSRWidthLibBOneCTwoSaferRustNum}{11}
\rqSet{rqSRWidthLibBOneCTwoSaferRustDen}{77}
\rqSet{rqSRWidthLibBOneCTwoSaferRustPct}{14.3}
\rqSet{rqSRWidthLibBOneSmartCTwoRustNum}{0}
\rqSet{rqSRWidthLibBOneSmartCTwoRustDen}{77}
\rqSet{rqSRWidthLibBOneSmartCTwoRustPct}{0.0}
\rqSet{rqSRWidthLibBOneKimiNum}{10}
\rqSet{rqSRWidthLibBOneKimiDen}{77}
\rqSet{rqSRWidthLibBOneKimiPct}{13.0}
\rqSet{rqSRWidthLibBOneGptFiveNum}{12}
\rqSet{rqSRWidthLibBOneGptFiveDen}{77}
\rqSet{rqSRWidthLibBOneGptFivePct}{15.6}
\rqSet{rqSRWidthLibBOneGeminiNum}{12}
\rqSet{rqSRWidthLibBOneGeminiDen}{77}
\rqSet{rqSRWidthLibBOneGeminiPct}{15.6}
\rqSet{rqSRWidthLibBTwoNum}{113}
\rqSet{rqSRWidthLibBTwoDen}{660}
\rqSet{rqSRWidthLibBTwoPct}{17.1}
\rqSet{rqSRWidthLibBTwoKiroNum}{15}
\rqSet{rqSRWidthLibBTwoKiroDen}{66}
\rqSet{rqSRWidthLibBTwoKiroPct}{22.7}
\rqSet{rqSRWidthLibBTwoClaudeNum}{15}
\rqSet{rqSRWidthLibBTwoClaudeDen}{66}
\rqSet{rqSRWidthLibBTwoClaudePct}{22.7}
\rqSet{rqSRWidthLibBTwoCodexNum}{15}
\rqSet{rqSRWidthLibBTwoCodexDen}{66}
\rqSet{rqSRWidthLibBTwoCodexPct}{22.7}
\rqSet{rqSRWidthLibBTwoCTwoRustNum}{14}
\rqSet{rqSRWidthLibBTwoCTwoRustDen}{66}
\rqSet{rqSRWidthLibBTwoCTwoRustPct}{21.2}
\rqSet{rqSRWidthLibBTwoLaertesNum}{13}
\rqSet{rqSRWidthLibBTwoLaertesDen}{66}
\rqSet{rqSRWidthLibBTwoLaertesPct}{19.7}
\rqSet{rqSRWidthLibBTwoCTwoSaferRustNum}{9}
\rqSet{rqSRWidthLibBTwoCTwoSaferRustDen}{66}
\rqSet{rqSRWidthLibBTwoCTwoSaferRustPct}{13.6}
\rqSet{rqSRWidthLibBTwoSmartCTwoRustNum}{0}
\rqSet{rqSRWidthLibBTwoSmartCTwoRustDen}{66}
\rqSet{rqSRWidthLibBTwoSmartCTwoRustPct}{0.0}
\rqSet{rqSRWidthLibBTwoKimiNum}{8}
\rqSet{rqSRWidthLibBTwoKimiDen}{66}
\rqSet{rqSRWidthLibBTwoKimiPct}{12.1}
\rqSet{rqSRWidthLibBTwoGptFiveNum}{9}
\rqSet{rqSRWidthLibBTwoGptFiveDen}{66}
\rqSet{rqSRWidthLibBTwoGptFivePct}{13.6}
\rqSet{rqSRWidthLibBTwoGeminiNum}{15}
\rqSet{rqSRWidthLibBTwoGeminiDen}{66}
\rqSet{rqSRWidthLibBTwoGeminiPct}{22.7}
\rqSet{rqSRWidthLibPZeroNum}{0}
\rqSet{rqSRWidthLibPZeroDen}{0}
\rqSet{rqSRWidthLibPZeroPct}{n/a}
\rqSet{rqSRWidthLibPZeroKiroNum}{0}
\rqSet{rqSRWidthLibPZeroKiroDen}{0}
\rqSet{rqSRWidthLibPZeroKiroPct}{n/a}
\rqSet{rqSRWidthLibPZeroClaudeNum}{0}
\rqSet{rqSRWidthLibPZeroClaudeDen}{0}
\rqSet{rqSRWidthLibPZeroClaudePct}{n/a}
\rqSet{rqSRWidthLibPZeroCodexNum}{0}
\rqSet{rqSRWidthLibPZeroCodexDen}{0}
\rqSet{rqSRWidthLibPZeroCodexPct}{n/a}
\rqSet{rqSRWidthLibPZeroCTwoRustNum}{0}
\rqSet{rqSRWidthLibPZeroCTwoRustDen}{0}
\rqSet{rqSRWidthLibPZeroCTwoRustPct}{n/a}
\rqSet{rqSRWidthLibPZeroLaertesNum}{0}
\rqSet{rqSRWidthLibPZeroLaertesDen}{0}
\rqSet{rqSRWidthLibPZeroLaertesPct}{n/a}
\rqSet{rqSRWidthLibPZeroCTwoSaferRustNum}{0}
\rqSet{rqSRWidthLibPZeroCTwoSaferRustDen}{0}
\rqSet{rqSRWidthLibPZeroCTwoSaferRustPct}{n/a}
\rqSet{rqSRWidthLibPZeroSmartCTwoRustNum}{0}
\rqSet{rqSRWidthLibPZeroSmartCTwoRustDen}{0}
\rqSet{rqSRWidthLibPZeroSmartCTwoRustPct}{n/a}
\rqSet{rqSRWidthLibPZeroKimiNum}{0}
\rqSet{rqSRWidthLibPZeroKimiDen}{0}
\rqSet{rqSRWidthLibPZeroKimiPct}{n/a}
\rqSet{rqSRWidthLibPZeroGptFiveNum}{0}
\rqSet{rqSRWidthLibPZeroGptFiveDen}{0}
\rqSet{rqSRWidthLibPZeroGptFivePct}{n/a}
\rqSet{rqSRWidthLibPZeroGeminiNum}{0}
\rqSet{rqSRWidthLibPZeroGeminiDen}{0}
\rqSet{rqSRWidthLibPZeroGeminiPct}{n/a}
\rqSet{rqSRWidthLibPOneNum}{3}
\rqSet{rqSRWidthLibPOneDen}{760}
\rqSet{rqSRWidthLibPOnePct}{0.4}
\rqSet{rqSRWidthLibPOneKiroNum}{0}
\rqSet{rqSRWidthLibPOneKiroDen}{76}
\rqSet{rqSRWidthLibPOneKiroPct}{0.0}
\rqSet{rqSRWidthLibPOneClaudeNum}{0}
\rqSet{rqSRWidthLibPOneClaudeDen}{76}
\rqSet{rqSRWidthLibPOneClaudePct}{0.0}
\rqSet{rqSRWidthLibPOneCodexNum}{3}
\rqSet{rqSRWidthLibPOneCodexDen}{76}
\rqSet{rqSRWidthLibPOneCodexPct}{3.9}
\rqSet{rqSRWidthLibPOneCTwoRustNum}{0}
\rqSet{rqSRWidthLibPOneCTwoRustDen}{76}
\rqSet{rqSRWidthLibPOneCTwoRustPct}{0.0}
\rqSet{rqSRWidthLibPOneLaertesNum}{0}
\rqSet{rqSRWidthLibPOneLaertesDen}{76}
\rqSet{rqSRWidthLibPOneLaertesPct}{0.0}
\rqSet{rqSRWidthLibPOneCTwoSaferRustNum}{0}
\rqSet{rqSRWidthLibPOneCTwoSaferRustDen}{76}
\rqSet{rqSRWidthLibPOneCTwoSaferRustPct}{0.0}
\rqSet{rqSRWidthLibPOneSmartCTwoRustNum}{0}
\rqSet{rqSRWidthLibPOneSmartCTwoRustDen}{76}
\rqSet{rqSRWidthLibPOneSmartCTwoRustPct}{0.0}
\rqSet{rqSRWidthLibPOneKimiNum}{0}
\rqSet{rqSRWidthLibPOneKimiDen}{76}
\rqSet{rqSRWidthLibPOneKimiPct}{0.0}
\rqSet{rqSRWidthLibPOneGptFiveNum}{0}
\rqSet{rqSRWidthLibPOneGptFiveDen}{76}
\rqSet{rqSRWidthLibPOneGptFivePct}{0.0}
\rqSet{rqSRWidthLibPOneGeminiNum}{0}
\rqSet{rqSRWidthLibPOneGeminiDen}{76}
\rqSet{rqSRWidthLibPOneGeminiPct}{0.0}
\rqSet{rqTermCPassLibAllNum}{4484947}
\rqSet{rqTermCPassLibAllDen}{6554751}
\rqSet{rqTermCPassLibAllPct}{68.4}
\rqSet{rqTermCPassLibBOneNum}{1794487}
\rqSet{rqTermCPassLibBOneDen}{3259995}
\rqSet{rqTermCPassLibBOnePct}{55.0}
\rqSet{rqTermCPassLibBTwoNum}{1675953}
\rqSet{rqTermCPassLibBTwoDen}{2280249}
\rqSet{rqTermCPassLibBTwoPct}{73.5}
\rqSet{rqTermCPassLibPZeroNum}{0}
\rqSet{rqTermCPassLibPZeroDen}{0}
\rqSet{rqTermCPassLibPZeroPct}{n/a}
\rqSet{rqTermCPassLibPOneNum}{1014507}
\rqSet{rqTermCPassLibPOneDen}{1014507}
\rqSet{rqTermCPassLibPOnePct}{100.0}
\rqSet{rqTermCVectorComparisonFailedLibAllNum}{1940123}
\rqSet{rqTermCVectorComparisonFailedLibAllDen}{6554751}
\rqSet{rqTermCVectorComparisonFailedLibAllPct}{29.6}
\rqSet{rqTermCVectorComparisonFailedLibBOneNum}{1379334}
\rqSet{rqTermCVectorComparisonFailedLibBOneDen}{3259995}
\rqSet{rqTermCVectorComparisonFailedLibBOnePct}{42.3}
\rqSet{rqTermCVectorComparisonFailedLibBTwoNum}{560789}
\rqSet{rqTermCVectorComparisonFailedLibBTwoDen}{2280249}
\rqSet{rqTermCVectorComparisonFailedLibBTwoPct}{24.6}
\rqSet{rqTermCVectorComparisonFailedLibPZeroNum}{0}
\rqSet{rqTermCVectorComparisonFailedLibPZeroDen}{0}
\rqSet{rqTermCVectorComparisonFailedLibPZeroPct}{n/a}
\rqSet{rqTermCVectorComparisonFailedLibPOneNum}{0}
\rqSet{rqTermCVectorComparisonFailedLibPOneDen}{1014507}
\rqSet{rqTermCVectorComparisonFailedLibPOnePct}{0.0}
\rqSet{rqTermCPanicLibAllNum}{0}
\rqSet{rqTermCPanicLibAllDen}{6554751}
\rqSet{rqTermCPanicLibAllPct}{0.0}
\rqSet{rqTermCPanicLibBOneNum}{0}
\rqSet{rqTermCPanicLibBOneDen}{3259995}
\rqSet{rqTermCPanicLibBOnePct}{0.0}
\rqSet{rqTermCPanicLibBTwoNum}{0}
\rqSet{rqTermCPanicLibBTwoDen}{2280249}
\rqSet{rqTermCPanicLibBTwoPct}{0.0}
\rqSet{rqTermCPanicLibPZeroNum}{0}
\rqSet{rqTermCPanicLibPZeroDen}{0}
\rqSet{rqTermCPanicLibPZeroPct}{n/a}
\rqSet{rqTermCPanicLibPOneNum}{0}
\rqSet{rqTermCPanicLibPOneDen}{1014507}
\rqSet{rqTermCPanicLibPOnePct}{0.0}
\rqSet{rqTermCSegmentationFaultLibAllNum}{92839}
\rqSet{rqTermCSegmentationFaultLibAllDen}{6554751}
\rqSet{rqTermCSegmentationFaultLibAllPct}{1.4}
\rqSet{rqTermCSegmentationFaultLibBOneNum}{66290}
\rqSet{rqTermCSegmentationFaultLibBOneDen}{3259995}
\rqSet{rqTermCSegmentationFaultLibBOnePct}{2.0}
\rqSet{rqTermCSegmentationFaultLibBTwoNum}{26549}
\rqSet{rqTermCSegmentationFaultLibBTwoDen}{2280249}
\rqSet{rqTermCSegmentationFaultLibBTwoPct}{1.2}
\rqSet{rqTermCSegmentationFaultLibPZeroNum}{0}
\rqSet{rqTermCSegmentationFaultLibPZeroDen}{0}
\rqSet{rqTermCSegmentationFaultLibPZeroPct}{n/a}
\rqSet{rqTermCSegmentationFaultLibPOneNum}{0}
\rqSet{rqTermCSegmentationFaultLibPOneDen}{1014507}
\rqSet{rqTermCSegmentationFaultLibPOnePct}{0.0}
\rqSet{rqTermCUnknownFailureLibAllNum}{23447}
\rqSet{rqTermCUnknownFailureLibAllDen}{6554751}
\rqSet{rqTermCUnknownFailureLibAllPct}{0.4}
\rqSet{rqTermCUnknownFailureLibBOneNum}{19857}
\rqSet{rqTermCUnknownFailureLibBOneDen}{3259995}
\rqSet{rqTermCUnknownFailureLibBOnePct}{0.6}
\rqSet{rqTermCUnknownFailureLibBTwoNum}{3590}
\rqSet{rqTermCUnknownFailureLibBTwoDen}{2280249}
\rqSet{rqTermCUnknownFailureLibBTwoPct}{0.2}
\rqSet{rqTermCUnknownFailureLibPZeroNum}{0}
\rqSet{rqTermCUnknownFailureLibPZeroDen}{0}
\rqSet{rqTermCUnknownFailureLibPZeroPct}{n/a}
\rqSet{rqTermCUnknownFailureLibPOneNum}{0}
\rqSet{rqTermCUnknownFailureLibPOneDen}{1014507}
\rqSet{rqTermCUnknownFailureLibPOnePct}{0.0}
\rqSet{rqTermCTimeoutLibAllNum}{5514}
\rqSet{rqTermCTimeoutLibAllDen}{6554751}
\rqSet{rqTermCTimeoutLibAllPct}{0.1}
\rqSet{rqTermCTimeoutLibBOneNum}{27}
\rqSet{rqTermCTimeoutLibBOneDen}{3259995}
\rqSet{rqTermCTimeoutLibBOnePct}{0.0}
\rqSet{rqTermCTimeoutLibBTwoNum}{5487}
\rqSet{rqTermCTimeoutLibBTwoDen}{2280249}
\rqSet{rqTermCTimeoutLibBTwoPct}{0.2}
\rqSet{rqTermCTimeoutLibPZeroNum}{0}
\rqSet{rqTermCTimeoutLibPZeroDen}{0}
\rqSet{rqTermCTimeoutLibPZeroPct}{n/a}
\rqSet{rqTermCTimeoutLibPOneNum}{0}
\rqSet{rqTermCTimeoutLibPOneDen}{1014507}
\rqSet{rqTermCTimeoutLibPOnePct}{0.0}
\rqSet{rqTermCOutputLimitLibAllNum}{0}
\rqSet{rqTermCOutputLimitLibAllDen}{6554751}
\rqSet{rqTermCOutputLimitLibAllPct}{0.0}
\rqSet{rqTermCOutputLimitLibBOneNum}{0}
\rqSet{rqTermCOutputLimitLibBOneDen}{3259995}
\rqSet{rqTermCOutputLimitLibBOnePct}{0.0}
\rqSet{rqTermCOutputLimitLibBTwoNum}{0}
\rqSet{rqTermCOutputLimitLibBTwoDen}{2280249}
\rqSet{rqTermCOutputLimitLibBTwoPct}{0.0}
\rqSet{rqTermCOutputLimitLibPZeroNum}{0}
\rqSet{rqTermCOutputLimitLibPZeroDen}{0}
\rqSet{rqTermCOutputLimitLibPZeroPct}{n/a}
\rqSet{rqTermCOutputLimitLibPOneNum}{0}
\rqSet{rqTermCOutputLimitLibPOneDen}{1014507}
\rqSet{rqTermCOutputLimitLibPOnePct}{0.0}
\rqSet{rqTermCMemLimitLibAllNum}{7881}
\rqSet{rqTermCMemLimitLibAllDen}{6554751}
\rqSet{rqTermCMemLimitLibAllPct}{0.1}
\rqSet{rqTermCMemLimitLibBOneNum}{0}
\rqSet{rqTermCMemLimitLibBOneDen}{3259995}
\rqSet{rqTermCMemLimitLibBOnePct}{0.0}
\rqSet{rqTermCMemLimitLibBTwoNum}{7881}
\rqSet{rqTermCMemLimitLibBTwoDen}{2280249}
\rqSet{rqTermCMemLimitLibBTwoPct}{0.3}
\rqSet{rqTermCMemLimitLibPZeroNum}{0}
\rqSet{rqTermCMemLimitLibPZeroDen}{0}
\rqSet{rqTermCMemLimitLibPZeroPct}{n/a}
\rqSet{rqTermCMemLimitLibPOneNum}{0}
\rqSet{rqTermCMemLimitLibPOneDen}{1014507}
\rqSet{rqTermCMemLimitLibPOnePct}{0.0}
\rqSet{rqTermCInfraLibAllNum}{0}
\rqSet{rqTermCInfraLibAllDen}{6554751}
\rqSet{rqTermCInfraLibAllPct}{0.0}
\rqSet{rqTermCInfraLibBOneNum}{0}
\rqSet{rqTermCInfraLibBOneDen}{3259995}
\rqSet{rqTermCInfraLibBOnePct}{0.0}
\rqSet{rqTermCInfraLibBTwoNum}{0}
\rqSet{rqTermCInfraLibBTwoDen}{2280249}
\rqSet{rqTermCInfraLibBTwoPct}{0.0}
\rqSet{rqTermCInfraLibPZeroNum}{0}
\rqSet{rqTermCInfraLibPZeroDen}{0}
\rqSet{rqTermCInfraLibPZeroPct}{n/a}
\rqSet{rqTermCInfraLibPOneNum}{0}
\rqSet{rqTermCInfraLibPOneDen}{1014507}
\rqSet{rqTermCInfraLibPOnePct}{0.0}
\rqSet{rqTermCInvalidInputLibAllNum}{0}
\rqSet{rqTermCInvalidInputLibAllDen}{6554751}
\rqSet{rqTermCInvalidInputLibAllPct}{0.0}
\rqSet{rqTermCInvalidInputLibBOneNum}{0}
\rqSet{rqTermCInvalidInputLibBOneDen}{3259995}
\rqSet{rqTermCInvalidInputLibBOnePct}{0.0}
\rqSet{rqTermCInvalidInputLibBTwoNum}{0}
\rqSet{rqTermCInvalidInputLibBTwoDen}{2280249}
\rqSet{rqTermCInvalidInputLibBTwoPct}{0.0}
\rqSet{rqTermCInvalidInputLibPZeroNum}{0}
\rqSet{rqTermCInvalidInputLibPZeroDen}{0}
\rqSet{rqTermCInvalidInputLibPZeroPct}{n/a}
\rqSet{rqTermCInvalidInputLibPOneNum}{0}
\rqSet{rqTermCInvalidInputLibPOneDen}{1014507}
\rqSet{rqTermCInvalidInputLibPOnePct}{0.0}
\rqSet{rqTermRPassLibAllNum}{4464435}
\rqSet{rqTermRPassLibAllDen}{6554751}
\rqSet{rqTermRPassLibAllPct}{68.1}
\rqSet{rqTermRPassLibBOneNum}{1777315}
\rqSet{rqTermRPassLibBOneDen}{3259995}
\rqSet{rqTermRPassLibBOnePct}{54.5}
\rqSet{rqTermRPassLibBTwoNum}{1672613}
\rqSet{rqTermRPassLibBTwoDen}{2280249}
\rqSet{rqTermRPassLibBTwoPct}{73.4}
\rqSet{rqTermRPassLibPZeroNum}{0}
\rqSet{rqTermRPassLibPZeroDen}{0}
\rqSet{rqTermRPassLibPZeroPct}{n/a}
\rqSet{rqTermRPassLibPOneNum}{1014507}
\rqSet{rqTermRPassLibPOneDen}{1014507}
\rqSet{rqTermRPassLibPOnePct}{100.0}
\rqSet{rqTermRVectorComparisonFailedLibAllNum}{1949555}
\rqSet{rqTermRVectorComparisonFailedLibAllDen}{6554751}
\rqSet{rqTermRVectorComparisonFailedLibAllPct}{29.7}
\rqSet{rqTermRVectorComparisonFailedLibBOneNum}{1388679}
\rqSet{rqTermRVectorComparisonFailedLibBOneDen}{3259995}
\rqSet{rqTermRVectorComparisonFailedLibBOnePct}{42.6}
\rqSet{rqTermRVectorComparisonFailedLibBTwoNum}{560876}
\rqSet{rqTermRVectorComparisonFailedLibBTwoDen}{2280249}
\rqSet{rqTermRVectorComparisonFailedLibBTwoPct}{24.6}
\rqSet{rqTermRVectorComparisonFailedLibPZeroNum}{0}
\rqSet{rqTermRVectorComparisonFailedLibPZeroDen}{0}
\rqSet{rqTermRVectorComparisonFailedLibPZeroPct}{n/a}
\rqSet{rqTermRVectorComparisonFailedLibPOneNum}{0}
\rqSet{rqTermRVectorComparisonFailedLibPOneDen}{1014507}
\rqSet{rqTermRVectorComparisonFailedLibPOnePct}{0.0}
\rqSet{rqTermRPanicLibAllNum}{36726}
\rqSet{rqTermRPanicLibAllDen}{6554751}
\rqSet{rqTermRPanicLibAllPct}{0.6}
\rqSet{rqTermRPanicLibBOneNum}{30006}
\rqSet{rqTermRPanicLibBOneDen}{3259995}
\rqSet{rqTermRPanicLibBOnePct}{0.9}
\rqSet{rqTermRPanicLibBTwoNum}{6720}
\rqSet{rqTermRPanicLibBTwoDen}{2280249}
\rqSet{rqTermRPanicLibBTwoPct}{0.3}
\rqSet{rqTermRPanicLibPZeroNum}{0}
\rqSet{rqTermRPanicLibPZeroDen}{0}
\rqSet{rqTermRPanicLibPZeroPct}{n/a}
\rqSet{rqTermRPanicLibPOneNum}{0}
\rqSet{rqTermRPanicLibPOneDen}{1014507}
\rqSet{rqTermRPanicLibPOnePct}{0.0}
\rqSet{rqTermRSegmentationFaultLibAllNum}{63178}
\rqSet{rqTermRSegmentationFaultLibAllDen}{6554751}
\rqSet{rqTermRSegmentationFaultLibAllPct}{1.0}
\rqSet{rqTermRSegmentationFaultLibBOneNum}{41596}
\rqSet{rqTermRSegmentationFaultLibBOneDen}{3259995}
\rqSet{rqTermRSegmentationFaultLibBOnePct}{1.3}
\rqSet{rqTermRSegmentationFaultLibBTwoNum}{21582}
\rqSet{rqTermRSegmentationFaultLibBTwoDen}{2280249}
\rqSet{rqTermRSegmentationFaultLibBTwoPct}{0.9}
\rqSet{rqTermRSegmentationFaultLibPZeroNum}{0}
\rqSet{rqTermRSegmentationFaultLibPZeroDen}{0}
\rqSet{rqTermRSegmentationFaultLibPZeroPct}{n/a}
\rqSet{rqTermRSegmentationFaultLibPOneNum}{0}
\rqSet{rqTermRSegmentationFaultLibPOneDen}{1014507}
\rqSet{rqTermRSegmentationFaultLibPOnePct}{0.0}
\rqSet{rqTermRUnknownFailureLibAllNum}{23073}
\rqSet{rqTermRUnknownFailureLibAllDen}{6554751}
\rqSet{rqTermRUnknownFailureLibAllPct}{0.4}
\rqSet{rqTermRUnknownFailureLibBOneNum}{18591}
\rqSet{rqTermRUnknownFailureLibBOneDen}{3259995}
\rqSet{rqTermRUnknownFailureLibBOnePct}{0.6}
\rqSet{rqTermRUnknownFailureLibBTwoNum}{4482}
\rqSet{rqTermRUnknownFailureLibBTwoDen}{2280249}
\rqSet{rqTermRUnknownFailureLibBTwoPct}{0.2}
\rqSet{rqTermRUnknownFailureLibPZeroNum}{0}
\rqSet{rqTermRUnknownFailureLibPZeroDen}{0}
\rqSet{rqTermRUnknownFailureLibPZeroPct}{n/a}
\rqSet{rqTermRUnknownFailureLibPOneNum}{0}
\rqSet{rqTermRUnknownFailureLibPOneDen}{1014507}
\rqSet{rqTermRUnknownFailureLibPOnePct}{0.0}
\rqSet{rqTermRTimeoutLibAllNum}{5465}
\rqSet{rqTermRTimeoutLibAllDen}{6554751}
\rqSet{rqTermRTimeoutLibAllPct}{0.1}
\rqSet{rqTermRTimeoutLibBOneNum}{0}
\rqSet{rqTermRTimeoutLibBOneDen}{3259995}
\rqSet{rqTermRTimeoutLibBOnePct}{0.0}
\rqSet{rqTermRTimeoutLibBTwoNum}{5465}
\rqSet{rqTermRTimeoutLibBTwoDen}{2280249}
\rqSet{rqTermRTimeoutLibBTwoPct}{0.2}
\rqSet{rqTermRTimeoutLibPZeroNum}{0}
\rqSet{rqTermRTimeoutLibPZeroDen}{0}
\rqSet{rqTermRTimeoutLibPZeroPct}{n/a}
\rqSet{rqTermRTimeoutLibPOneNum}{0}
\rqSet{rqTermRTimeoutLibPOneDen}{1014507}
\rqSet{rqTermRTimeoutLibPOnePct}{0.0}
\rqSet{rqTermROutputLimitLibAllNum}{0}
\rqSet{rqTermROutputLimitLibAllDen}{6554751}
\rqSet{rqTermROutputLimitLibAllPct}{0.0}
\rqSet{rqTermROutputLimitLibBOneNum}{0}
\rqSet{rqTermROutputLimitLibBOneDen}{3259995}
\rqSet{rqTermROutputLimitLibBOnePct}{0.0}
\rqSet{rqTermROutputLimitLibBTwoNum}{0}
\rqSet{rqTermROutputLimitLibBTwoDen}{2280249}
\rqSet{rqTermROutputLimitLibBTwoPct}{0.0}
\rqSet{rqTermROutputLimitLibPZeroNum}{0}
\rqSet{rqTermROutputLimitLibPZeroDen}{0}
\rqSet{rqTermROutputLimitLibPZeroPct}{n/a}
\rqSet{rqTermROutputLimitLibPOneNum}{0}
\rqSet{rqTermROutputLimitLibPOneDen}{1014507}
\rqSet{rqTermROutputLimitLibPOnePct}{0.0}
\rqSet{rqTermRMemLimitLibAllNum}{12319}
\rqSet{rqTermRMemLimitLibAllDen}{6554751}
\rqSet{rqTermRMemLimitLibAllPct}{0.2}
\rqSet{rqTermRMemLimitLibBOneNum}{3808}
\rqSet{rqTermRMemLimitLibBOneDen}{3259995}
\rqSet{rqTermRMemLimitLibBOnePct}{0.1}
\rqSet{rqTermRMemLimitLibBTwoNum}{8511}
\rqSet{rqTermRMemLimitLibBTwoDen}{2280249}
\rqSet{rqTermRMemLimitLibBTwoPct}{0.4}
\rqSet{rqTermRMemLimitLibPZeroNum}{0}
\rqSet{rqTermRMemLimitLibPZeroDen}{0}
\rqSet{rqTermRMemLimitLibPZeroPct}{n/a}
\rqSet{rqTermRMemLimitLibPOneNum}{0}
\rqSet{rqTermRMemLimitLibPOneDen}{1014507}
\rqSet{rqTermRMemLimitLibPOnePct}{0.0}
\rqSet{rqTermRInfraLibAllNum}{0}
\rqSet{rqTermRInfraLibAllDen}{6554751}
\rqSet{rqTermRInfraLibAllPct}{0.0}
\rqSet{rqTermRInfraLibBOneNum}{0}
\rqSet{rqTermRInfraLibBOneDen}{3259995}
\rqSet{rqTermRInfraLibBOnePct}{0.0}
\rqSet{rqTermRInfraLibBTwoNum}{0}
\rqSet{rqTermRInfraLibBTwoDen}{2280249}
\rqSet{rqTermRInfraLibBTwoPct}{0.0}
\rqSet{rqTermRInfraLibPZeroNum}{0}
\rqSet{rqTermRInfraLibPZeroDen}{0}
\rqSet{rqTermRInfraLibPZeroPct}{n/a}
\rqSet{rqTermRInfraLibPOneNum}{0}
\rqSet{rqTermRInfraLibPOneDen}{1014507}
\rqSet{rqTermRInfraLibPOnePct}{0.0}
\rqSet{rqTermRInvalidInputLibAllNum}{0}
\rqSet{rqTermRInvalidInputLibAllDen}{6554751}
\rqSet{rqTermRInvalidInputLibAllPct}{0.0}
\rqSet{rqTermRInvalidInputLibBOneNum}{0}
\rqSet{rqTermRInvalidInputLibBOneDen}{3259995}
\rqSet{rqTermRInvalidInputLibBOnePct}{0.0}
\rqSet{rqTermRInvalidInputLibBTwoNum}{0}
\rqSet{rqTermRInvalidInputLibBTwoDen}{2280249}
\rqSet{rqTermRInvalidInputLibBTwoPct}{0.0}
\rqSet{rqTermRInvalidInputLibPZeroNum}{0}
\rqSet{rqTermRInvalidInputLibPZeroDen}{0}
\rqSet{rqTermRInvalidInputLibPZeroPct}{n/a}
\rqSet{rqTermRInvalidInputLibPOneNum}{0}
\rqSet{rqTermRInvalidInputLibPOneDen}{1014507}
\rqSet{rqTermRInvalidInputLibPOnePct}{0.0}
\rqSet{rqFFComboStateFloatLibAllNum}{1}
\rqSet{rqFFComboStateFloatLibAllDen}{1}
\rqSet{rqFFComboStateFloatLibAllPct}{100.0}
\rqSet{rqFFComboStateFloatLibBOneNum}{0}
\rqSet{rqFFComboStateFloatLibBOneDen}{0}
\rqSet{rqFFComboStateFloatLibBOnePct}{n/a}
\rqSet{rqFFComboStateFloatLibBTwoNum}{1}
\rqSet{rqFFComboStateFloatLibBTwoDen}{1}
\rqSet{rqFFComboStateFloatLibBTwoPct}{100.0}
\rqSet{rqFFComboStateFloatLibPZeroNum}{0}
\rqSet{rqFFComboStateFloatLibPZeroDen}{0}
\rqSet{rqFFComboStateFloatLibPZeroPct}{n/a}
\rqSet{rqFFComboStateFloatLibPOneNum}{0}
\rqSet{rqFFComboStateFloatLibPOneDen}{0}
\rqSet{rqFFComboStateFloatLibPOnePct}{n/a}
}{}
\IfFileExists{generated/v3-numbers.tex}{% SINH TU DONG boi writing/scripts/gen-v3-macros.py. KHONG SUA TAY.
% Repo commit: d46010271517095fe135f15a336a18c2bc9b23b3. Nguon: generated/rq-numbers.tex, generated/sanbound-numbers.tex (exp 38), generated/census-numbers.tex, experiments/41-lp-star/results (exp 41).
% Dinh nghia: decision-log 2026-11-12 (dau tren chat L_R*, L_P*); v3/03-methodology.tex muc ve khoang RQ3.
\providecommand{\rqSet}[2]{\expandafter\def\csname #1\endcsname{#2}}
\rqSet{rqSRUpExeAllPct}{35.0}
\rqSet{rqSRUpExeAllNum}{371}
\rqSet{rqSRUpExeAllKiroPct}{62.3}
\rqSet{rqSRUpExeAllKiroNum}{66}
\rqSet{rqSRUpExeAllClaudePct}{70.8}
\rqSet{rqSRUpExeAllClaudeNum}{75}
\rqSet{rqSRUpExeAllCodexPct}{50.9}
\rqSet{rqSRUpExeAllCodexNum}{54}
\rqSet{rqSRUpExeAllCTwoRustPct}{48.1}
\rqSet{rqSRUpExeAllCTwoRustNum}{51}
\rqSet{rqSRUpExeAllLaertesPct}{50.0}
\rqSet{rqSRUpExeAllLaertesNum}{53}
\rqSet{rqSRUpExeAllCTwoSaferRustPct}{26.4}
\rqSet{rqSRUpExeAllCTwoSaferRustNum}{28}
\rqSet{rqSRUpExeAllSmartCTwoRustPct}{10.4}
\rqSet{rqSRUpExeAllSmartCTwoRustNum}{11}
\rqSet{rqSRUpExeAllKimiPct}{8.5}
\rqSet{rqSRUpExeAllKimiNum}{9}
\rqSet{rqSRUpExeAllGptFivePct}{9.4}
\rqSet{rqSRUpExeAllGptFiveNum}{10}
\rqSet{rqSRUpExeAllGeminiPct}{13.2}
\rqSet{rqSRUpExeAllGeminiNum}{14}
\rqSet{rqSRUpLibAllPct}{56.8}
\rqSet{rqSRUpLibAllNum}{1243}
\rqSet{rqSRUpLibAllKiroPct}{96.8}
\rqSet{rqSRUpLibAllKiroNum}{212}
\rqSet{rqSRUpLibAllClaudePct}{92.2}
\rqSet{rqSRUpLibAllClaudeNum}{202}
\rqSet{rqSRUpLibAllCodexPct}{77.6}
\rqSet{rqSRUpLibAllCodexNum}{170}
\rqSet{rqSRUpLibAllCTwoRustPct}{63.5}
\rqSet{rqSRUpLibAllCTwoRustNum}{139}
\rqSet{rqSRUpLibAllLaertesPct}{61.6}
\rqSet{rqSRUpLibAllLaertesNum}{135}
\rqSet{rqSRUpLibAllCTwoSaferRustPct}{45.7}
\rqSet{rqSRUpLibAllCTwoSaferRustNum}{100}
\rqSet{rqSRUpLibAllSmartCTwoRustPct}{0.0}
\rqSet{rqSRUpLibAllSmartCTwoRustNum}{0}
\rqSet{rqSRUpLibAllKimiPct}{34.2}
\rqSet{rqSRUpLibAllKimiNum}{75}
\rqSet{rqSRUpLibAllGptFivePct}{45.2}
\rqSet{rqSRUpLibAllGptFiveNum}{99}
\rqSet{rqSRUpLibAllGeminiPct}{50.7}
\rqSet{rqSRUpLibAllGeminiNum}{111}
\rqSet{rqSRUpExeAllDen}{1060}
\rqSet{rqSRUpNoLrStarExeAllNum}{407}
\rqSet{rqSRUpNoLrStarExeAllPct}{38.4}
\rqSet{rqSRUpExeAllKiroDen}{106}
\rqSet{rqSRUpExeAllClaudeDen}{106}
\rqSet{rqSRUpExeAllCodexDen}{106}
\rqSet{rqSRUpNoLrStarExeAllCodexNum}{90}
\rqSet{rqSRUpNoLrStarExeAllCodexPct}{84.9}
\rqSet{rqSRUpExeAllCTwoRustDen}{106}
\rqSet{rqSRUpExeAllLaertesDen}{106}
\rqSet{rqSRUpExeAllCTwoSaferRustDen}{106}
\rqSet{rqSRUpExeAllSmartCTwoRustDen}{106}
\rqSet{rqSRUpExeAllKimiDen}{106}
\rqSet{rqSRUpExeAllGptFiveDen}{106}
\rqSet{rqSRUpExeAllGeminiDen}{106}
\rqSet{rqLpStarRetExeAllNum}{12}
\rqSet{rqSRUpLibAllDen}{2190}
\rqSet{rqSRUpLibAllKiroDen}{219}
\rqSet{rqSRUpLibAllClaudeDen}{219}
\rqSet{rqSRUpLibAllCodexDen}{219}
\rqSet{rqSRUpLibAllCTwoRustDen}{219}
\rqSet{rqSRUpLibAllLaertesDen}{219}
\rqSet{rqSRUpLibAllCTwoSaferRustDen}{219}
\rqSet{rqSRUpLibAllSmartCTwoRustDen}{219}
\rqSet{rqSRUpLibAllKimiDen}{219}
\rqSet{rqSRUpLibAllGptFiveDen}{219}
\rqSet{rqSRUpLibAllGeminiDen}{219}
\rqSet{rqLpStarRetLibAllNum}{1}
\rqSet{rqLpStarExeAllNum}{71}
\rqSet{rqFPBoundLoTightExeAllNum}{188}
\rqSet{rqFPBoundLoTightExeAllDen}{559}
\rqSet{rqFPBoundLoTightExeAllPct}{33.6}
\rqSet{rqLossTotalExeAllNum}{119}
\rqSet{rqOpenExeAllNum}{71}
\rqSet{rqLrStarExeAllNum}{0}
\rqSet{rqSRWidthTightExeAllPct}{6.7}
\rqSet{rqSRDecrMinExeAllPct}{17.7}
\rqSet{rqSRDecrMaxExeAllPct}{24.4}
\rqSet{rqUnfinishedExeAllNum}{72}
\rqSet{rqUnfinishedOpenExeAllNum}{60}
\rqSet{rqLpStarLibAllNum}{192}
\rqSet{rqFPBoundLoTightLibAllNum}{66}
\rqSet{rqFPBoundLoTightLibAllDen}{1281}
\rqSet{rqFPBoundLoTightLibAllPct}{5.2}
\rqSet{rqLossTotalLibAllNum}{220}
\rqSet{rqOpenLibAllNum}{219}
\rqSet{rqLrStarLibAllNum}{27}
\rqSet{rqSRWidthTightLibAllPct}{10.0}
\rqSet{rqSRDecrMinLibAllPct}{1.7}
\rqSet{rqSRDecrMaxLibAllPct}{11.7}
\rqSet{rqUnfinishedLibAllNum}{91}
\rqSet{rqUnfinishedOpenLibAllNum}{90}
\rqSet{rqInsensPtestExeAllNum}{132}
\rqSet{rqInsensDExeAllNum}{0}
\rqSet{rqSensFPExeAllNum}{176}
\rqSet{rqSensFPExeAllDen}{344}
\rqSet{rqSensFPExeAllPct}{51.2}
\rqSet{rqLabInspN}{140}
\rqSet{rqLabInspAgree}{128}
\rqSet{rqLabRqTwoN}{61}
\rqSet{rqLabRqTwoAgree}{60}
\rqSet{rqLabThemeN}{65}
\rqSet{rqLabThemeAgree}{60}
\rqSet{rqLabNonrepN}{24}
\rqSet{rqLabNonrepAgree}{24}
\rqSet{nTotalPilot}{13}
\rqSet{nTotalPrimary}{325}
\rqSet{nTotalPzero}{3250}
\rqSet{nTotalKimiMissing}{123}
}{}
\thispagestyle{plain}
\begin{center}
{\fontsize{13}{16}\selectfont\bfseries ABSTRACT}
\end{center}
\vspace{1em}

% Phu luc 06: noi dung abstract va keywords dung 12pt, khac co chu toan cuc 13pt
\begingroup
\fontsize{12}{15}\selectfont
\noindent\textbf{Abstract:}
Systems that translate C programs into Rust are usually evaluated by compiling each translation and running an existing test suite. The result for each translation is recorded as a PASS or FAIL decision, and the decisions are aggregated into the success rate of a system. How far such a rate can be trusted depends on whether the tests reveal differences in behavior and on whether the expectations used for grading are accurate, and studies of related tasks have shown that accepted outputs often fail when they are tested further. This thesis re-evaluates the published decisions of ten translation systems on the TRACTOR Public Test Corpus without rerunning any translator. It builds an independent differential testing harness that runs each Rust translation and its C program on the same inputs and compares their termination, their output and, for libraries, the state after a call. For accepted translations, the harness searches for a behavioral divergence on the original test vectors and on generated inputs; for translations rejected at the test stage, it replays the vectors that failed. An observation counts as evidence only if it is stable under repeated runs and the C execution completes a sanitizer check for undefined behavior without a report. Among the accepted translations that were fully re-evaluated, \rqFPExeAllNum/\rqFPExeAllDen{} executable (\rqFPExeAllPct\%) and \rqFPLibAllNum/\rqFPLibAllDen{} library translations (\rqFPLibAllPct\%) have at least one retained divergence, found almost entirely on generated inputs. Among the test-stage rejections that reached an outcome, \rqFFExeAllNum/\rqFFExeAllDen{} executable and \rqFFLibAllNum/\rqFFLibAllDen{} library rejections lack supporting evidence. When these reclassifications and the decisions that remain without a result are carried back into the original population, the success rate of executables under the oracle of this study lies between \rqSRLoExeAllPct\% and \rqSRUpExeAllPct\%, against \rqSRExeAllPct\% under the published decisions, and that of libraries between \rqSRLoLibAllPct\% and \rqSRUpLibAllPct\%, against \rqSRLibAllPct\%. On this corpus, the evidence that changes the reading of an evaluation lies mainly on the side of acceptances. The results indicate that accepted translations need behavioral evidence beyond the original tests, and that a reported success rate should be accompanied by its comparison rule, its full denominator and the uncertainty left by the decisions that were not examined. The findings are limited to one corpus, and a retained divergence records an observed difference whose status as a semantic error remains open.

\vspace{1em}
\noindent\textit{\textbf{Keywords:}} \textit{C-to-Rust translation, differential testing, test oracle, undefined behavior}
\endgroup
\clearpage
